%%%%%%%%%%%%%%%%%%%%%%%%%%%%%%%%%%%%%%%%%
% Nuo Li's Doctoral Thesis



\documentclass[11pt,a4paper,UKenglish,twoside,openright]{Thesis}
\usepackage{etex}

\PassOptionsToPackage{utf8x}{inputenc}

%\DeclareUnicodeCharacter{00A0}{ }

\graphicspath{{Pictures/}} % Specifies the directory where pictures are stored


\usepackage[all]{xy}
%agda literal file
\usepackage{agda}

\usepackage{multicol}
\usepackage{dsfont} % math symbols
\usepackage{amsthm}
\usepackage{relsize}
\usepackage{color}
\usepackage{amsmath}
\usepackage{amsfonts}
\usepackage{amssymb}
\usepackage{bbm}
\usepackage[greek,english]{babel}
\usepackage{stmaryrd}
\usepackage{textgreek}

%\usepackage{unicode-math}

%\usepackage[T1]{fontenc}
\usepackage[utf8x]{inputenx}
\usepackage{autofe} % font for code environment


\usepackage{booktabs}% horizontal lines

\usepackage{proof} % For type derivation
\setlength{\inferLineSkip}{4pt}

\usepackage{diagxy} % commutative diagram 

\usepackage{bcprules}

\usepackage[square, numbers, comma, sort&compress]{natbib} 
% Use the natbib reference package - read up on this to edit the reference style; if you want text (e.g. Smith et al., 2012) for the in-text references (instead of numbers), remove 'numbers' 

%\usepackage[toc,page]{appendix}


\usepackage{xkvltxp}
\usepackage[]{fixme}
\fxsetup{
    status=draft,
    author=,
    layout=inline,
    theme=color
}

\definecolor{fxnote}{rgb}{0.8000,0.0000,0.0000}
% define the background colour:
\colorlet{fxnotebg}{yellow}

% refedine the layout macro:
\makeatletter
\renewcommand*\FXLayoutInline[3]{%
  \@fxdocolon {#3}{%
    \@fxuseface {inline}%
    \colorbox{fx#1bg}{\color {fx#1}\ignorespaces #3\@fxcolon #2}}}
\makeatother

%\usepackage{glossaries}


\usepackage[textwidth=3cm]{todonotes}
\reversemarginpar
\setlength{\marginparwidth}{3cm}


%%% symbols and macros, mostly from the HoTT book



\newcommand{\prd}[1]{\Pi_{#1}}
\newcommand{\sm}[1]{\Sigma_{#1}}

% \newcommand{\smmultiline}[1]{\Sigma \; \left #1 \right}

\newcommand{\lam}[1]{\lambda #1 .}

%own
\newcommand{\fa}[2]{\forall (#1 : #2) . \,}
% \newcommand{\fa}[1]{\forall (#1). \,}
\newcommand{\fasimple}[1]{\forall #1 . \,}
\newcommand{\fapairs}[3]{\fa {#1 \, #2} #3}
\newcommand{\fatypes}[1]{\fa #1 \type} % forall types

%%% Definitional equality (used infix) %%%
\newcommand{\jdeq}{\equiv}      % An equality judgment
\newcommand{\defeq}{\vcentcolon\equiv}  % A judgmental equality currently being defined

%%% Homotopies (written infix) %%%
\newcommand{\eqvsym}{\simeq}
\newcommand{\eqv}[2]{\ensuremath{#1 \eqvsym #2}\xspace}
\newcommand{\eqvauto}[1]{\eqv{#1}{#1}}


%%% own notation: equivalences
\newcommand{\iseqown}[1]{\ensuremath{\mathsf{e}_{#1}}\xspace}
\newcommand{\eqvsymspace}{\enspace \eqvsym \enspace}
\newcommand{\eqvspace}[2]{\ensuremath{#1 \eqvsymspace #2}\xspace}
\newcommand{\eqvautospace}[1]{\eqvspace{#1}{#1}}


%%% Identity functions %%%
\newcommand{\idfunc}[1][]{\ensuremath{\mathsf{id}_{#1}}\xspace}

%%% Identity types %%%
\newcommand{\idsym}{{=}}
\newcommand{\id}[3][]{\ensuremath{#2 =_{#1} #3}\xspace}
\newcommand{\idspace}[3][]{\ensuremath{#2 \, =_{#1} \, #3}\xspace}
\newcommand{\idtype}[3][]{\ensuremath{\mathsf{Id}_{#1}(#2,#3)}\xspace}
\newcommand{\idtypevar}[1]{\ensuremath{\mathsf{Id}_{#1}}\xspace}
\newcommand{\idauto}[2][]{\id[#1]{#2}{#2}}


%%% Reflexivity terms %%%
\newcommand{\refl}[1]{\ensuremath{\mathsf{refl}_{#1}}\xspace}

%%% Path concatenation (used infix, in diagrammatic order) %%%
\newcommand{\ct}{%
  \mathchoice{\mathbin{\raisebox{0.5ex}{$\displaystyle\centerdot$}}}%
             {\mathbin{\raisebox{0.5ex}{$\centerdot$}}}%
             {\mathbin{\raisebox{0.25ex}{$\scriptstyle\,\centerdot\,$}}}%
             {\mathbin{\raisebox{0.1ex}{$\scriptscriptstyle\,\centerdot\,$}}}
}

%%% Path reversal %%%
\newcommand{\opp}[1]{\mathord{{#1}^{-1}}}
\let\rev\opp
\let\sym\opp

%%% Transport (covariant) %%%
\newcommand{\trans}[2]{\ensuremath{{#1}_{*}\mathopen{}\left({#2}\right)\mathclose{}}\xspace}

\newcommand{\isequiv}{\ensuremath{\mathsf{isequiv}}}
\newcommand{\idequiv}{\ensuremath{\textsf{id-equiv}}}


%%% Universe types %%%
\newcommand{\UU}{\ensuremath{\mathcal{U}}\xspace}
\let\bbU\UU
\let\type\UU
% own
\newcommand{\UUp}{{\UU_\bullet}}
% \newcommand{\UUtt}[2]{\ensuremath{ \UU_{#1}^{\mathsmaller{\leq} #2}}}
% \newcommand{\UUt}[1]{\ensuremath{ \UU^{\mathsmaller{\leq} #1}}}
%%% Natural numbers
\newcommand{\N}{\ensuremath{\mathbb{N}}\xspace}
%\newcommand{\N}{\textbf{N}}
\let\nat\N
\newcommand{\suc}{\mathsf{succ}}
\newcommand{\add}{\mathsf{add}}
% \newcommand{\ack}{\mathsf{ack}}
\newcommand{\ass}{\mathsf{ass}}
\newcommand{\ite}{\mathsf{iter}}

%own
\newcommand{\Q}{\ensuremath{\mathbb{Q}}\xspace}
% \newcommand{\R}{\ensuremath{\mathbb{R}}\xspace}



\newcommand{\dbl}{\ensuremath{\mathsf{double}}}


%%% The empty type
\newcommand{\emptyt}{\ensuremath{\mathbf{0}}\xspace}

%%% The unit type
\newcommand{\unit}{\ensuremath{\mathbf{1}}\xspace}
\newcommand{\ttt}{\ensuremath{\star}\xspace}

%%% The two-element type
\newcommand{\bool}{\ensuremath{\mathbf{2}}\xspace}
\newcommand{\btrue}{{1_{\bool}}}
\newcommand{\bfalse}{{0_{\bool}}}
\let\two\bool
\let\trueb\btrue
\let\falseb\bfalse

\newcommand{\swap}{\ensuremath{\textsf{swap}}}
\newcommand{\swapequiv}{\ensuremath{\textsf{swap-equiv}}}


%%% Truncation levels
\newcommand{\iscontr}{\ensuremath{\mathsf{isContr}}}
\newcommand{\isprop}{\ensuremath{\mathsf{isProp}}}
\newcommand{\isset}{\ensuremath{\mathsf{isSet}}}


\def\compare#1#2#3#4{\if#1#3\if#2#41\else0\fi\else0\fi}

\newcommand{\istype}[1]{
  \edef\a{\compare-2#1\empty\empty}
  \if\a1 \iscontr \else
  \edef\b{\compare-1#1\empty\empty}
  \if\b1 \isprop \else
  \edef\c{#1}
  \if0\c \isset \else
  \mathsf{is}\mbox{-}{#1}\mbox{-}\mathsf{type} \fi\fi\fi
}



%%% Univalence
\newcommand{\ua}{\ensuremath{\mathsf{ua}}\xspace} % the inverse of idtoeqv
\newcommand{\idtoeqv}{\ensuremath{\mathsf{idtoeqv}}\xspace}
\newcommand{\univalence}{\ensuremath{\mathsf{univalence}}\xspace} % the full axiom


%%% Circles and spheres
\newcommand{\Sn}{\mathbb{S}}
\newcommand{\base}{\ensuremath{\mathsf{base}}\xspace}
\newcommand{\lloop}{\ensuremath{\mathsf{loop}}\xspace}
\newcommand{\surf}{\ensuremath{\mathsf{surf}}\xspace}

%%% Map on paths %%%
\newcommand{\mapfunc}[1]{\ensuremath{\mathsf{ap}_{#1}}\xspace} % Without argument
\newcommand{\map}[2]{\ensuremath{{#1}\mathopen{}\left({#2}\right)\mathclose{}}\xspace}
\let\Ap\map
%\newcommand{\Ap}[2]{\ensuremath{{#1}\left({#2}\right)}\xspace}
\newcommand{\mapdepfunc}[1]{\ensuremath{\mathsf{apd}_{#1}}\xspace} % Without argument
% \newcommand{\mapdep}[2]{\ensuremath{{#1}\llparenthesis{#2}\rrparenthesis}\xspace}
\newcommand{\mapdep}[2]{\ensuremath{\mapdepfunc{#1}\mathopen{}\left(#2\right)\mathclose{}}\xspace}
\let\apfunc\mapfunc
\let\ap\map
\let\apdfunc\mapdepfunc
\let\apd\mapdep

%own
\newcommand{\happly}{\ensuremath{\mathsf{happly}}}

%% OWN: UIP
\newcommand{\uipplain}{\mathsf{UIP}}
\newcommand{\uipext}[1]{\uipplain_{#1}}

\newcommand{\dec}{\mathsf{decidable}}
\newcommand{\deceq}{\mathsf{discrete}}


% % \newcommand{\trunc}[2]{\left\lVert #1 \right\rVert_{#2}} % note: \| is not symmetric, is it?
% % \newcommand{\truncmap}[2]{\left\lvert #1 \right\rvert_{#2}} % note: \| is not symmetric, is it?
% \newcommand{\trunc}[1]{\left\lVert #1 \right\rVert} % note: \| is not symmetric, is it?
% \newcommand{\truncmap}[1]{\left\lvert #1 \right\rvert} % note: \| is not symmetric, is it?

%%% Bracket/squash/truncation types %%%
\newcommand{\trunc}[2]{\mathopen{}\left\Vert #2\right\Vert_{#1}\mathclose{}}
\newcommand{\ttrunc}[2]{\bigl\Vert #2\bigr\Vert_{#1}}
\newcommand{\Trunc}[2]{\Bigl\Vert #2\Bigr\Vert_{#1}}
\newcommand{\truncf}[1]{\Vert \blank \Vert_{#1}}
\newcommand{\tproj}[3][]{\mathopen{}\left|#3\right|_{#2}^{#1}\mathclose{}}
\newcommand{\tprojf}[2][]{|\blank|_{#2}^{#1}}
\def\pizero{\trunc0}
%\newcommand{\brck}[1]{\trunc{-1}{#1}}
%\newcommand{\Brck}[1]{\Trunc{-1}{#1}}
%\newcommand{\bproj}[1]{\tproj{-1}{#1}}
%\newcommand{\bprojf}{\tprojf{-1}}



\newcommand{\bbbrck}[1]{\Bigl\Vert #1 \Bigr\Vert}
\newcommand{\brck}[1]{\trunc{}{#1}}
\newcommand{\bbrck}[1]{\ttrunc{}{#1}}
\let\bracket\bbrck
\newcommand{\Brck}[1]{\Trunc{}{#1}}
\let\Bracket\Brck
\newcommand{\bproj}[1]{\tproj{}{#1}}
\newcommand{\bprojf}{\tprojf{}}

% OWN
\newcommand{\populated}[1]{\langle \! \langle #1 \rangle \! \rangle}


%OWN
\newcommand{\fst}{\pi_1 \!}
\newcommand{\snd}{\pi_2 \!}
\newcommand{\inl}{\mathsf{inl}\xspace}
\newcommand{\inr}{\mathsf{inr}\xspace}
\newcommand{\inone}{\mathsf{in_1}\xspace}
\newcommand{\intwo}{\mathsf{in_2}\xspace}
\newcommand{\inthree}{\mathsf{in_3}\xspace}


%\newcommand{\hproptype}{\textbf{hProp}}


\newcommand{\LEM}[1]{\ensuremath{\mathsf{LEM}_{#1}}\xspace}
\newcommand{\choice}[1]{\ensuremath{\mathsf{AC}_{#1}}\xspace}

\usepackage{mypack}

\usepackage{tikz}

\hypersetup{urlcolor=blue, colorlinks=true}
% Colors hyperlinks in blue - change to black if annoying


\title{\ttitle}
% Defines the thesis title - don't touch this



\newcommand{\HRule}{\rule{\linewidth}{0.5mm}} % New command to make the lines in the title page





\begin{document}

\frontmatter % Use roman page numbering style (i, ii, iii, iv...) for the pre-content pages

\setstretch{1.3} % Line spacing of 1.3

% Define the page headers using the FancyHdr package and set up for one-sided printing
\fancyhead{} % Clears all page headers and footers
\rhead{\thepage} % Sets the right side header to show the page number
\lhead{} % Clears the left side page header

\pagestyle{fancy} % Finally, use the "fancy" page style to implement the FancyHdr headers

%\newcommand{\HRule}{\rule{\linewidth}{0.5mm}} % New command to make the lines in the title page

% PDF meta-data
\hypersetup{pdftitle={\ttitle}}
\hypersetup{pdfsubject=\subjectname}
\hypersetup{pdfauthor=\authornames}
\hypersetup{pdfkeywords=\keywordnames}


\listoftodos

%----------------------------------------------------------------------------------------
%	TITLE PAGE
%----------------------------------------------------------------------------------------

\begin{titlepage}

\begin{center}

\includegraphics[height=1.5cm,scale=0.1]{Pictures/2000px-University_of_Nottingham.png}% Logo instead of name
%\textsc{\LARGE \univname}\\[1.5cm] % University name
\textsc{\Large Doctoral Thesis}\\[0.5cm] % Thesis type

\HRule \\[0.4cm] % Horizontal line
{\huge \bfseries \ttitle}\\[0.4cm] % Thesis title
\HRule \\[1.5cm] % Horizontal line
 
\begin{minipage}{0.4\textwidth}
\begin{flushleft} \large
\emph{Author:}\\
{\authornames} % Author name - remove the \href bracket to remove the link
\end{flushleft}
\end{minipage}
\begin{minipage}{0.4\textwidth}
\begin{flushright} \large
\emph{Supervisor:} \\
{\supname} % Supervisor name - remove the \href bracket to remove the link  
\end{flushright}
\end{minipage}\\[3cm]
 
\large \textit{A thesis submitted in fulfillment of the requirements\\ for the degree of \degreename}\\[0.3cm] % University requirement text
\textit{in the}\\[0.4cm]
\groupname\\\deptname\\[2cm] % Research group name and department name
 
{\large \today}\\[4cm] % Date
%\includegraphics{Logo} % University/department logo - uncomment to place it
 
\vfill
\end{center}

\end{titlepage}


\clearpage % Start a new page

%----------------------------------------------------------------------------------------
%	ABSTRACT PAGE
%----------------------------------------------------------------------------------------

\addtotoc{Abstract} % Add the "Abstract" page entry to the Contents

\abstract{\addtocontents{toc}{\vspace{1em}} % Add a gap in the Contents, for aesthetics

The Thesis Abstract is written here (and usually kept to just this page). This thesis mainly covered the quotient types in type theory
}


\clearpage % Start a new page

%----------------------------------------------------------------------------------------
%	ACKNOWLEDGEMENTS
%----------------------------------------------------------------------------------------

\setstretch{1.3} % Reset the line-spacing to 1.3 for body text (if it has changed)

\acknowledgements{\addtocontents{toc}{\vspace{1em}} % Add a gap in the Contents, for aesthetics

The acknowledgments and the people to thank go here, don't forget to include your project advisor\ldots
}
\clearpage % Start a new page

%----------------------------------------------------------------------------------------
%	LIST OF CONTENTS/FIGURES/TABLES PAGES
%----------------------------------------------------------------------------------------

\pagestyle{fancy} % The page style headers have been "empty" all this time, now use the "fancy" headers as defined before to bring them back

\lhead{\emph{Contents}} % Set the left side page header to "Contents"
\tableofcontents % Write out the Table of Contents

\lhead{\emph{List of Figures}} % Set the left side page header to "List of Figures"
\listoffigures % Write out the List of Figures

\lhead{\emph{List of Tables}} % Set the left side page header to "List of Tables"
\listoftables % Write out the List of Tables




%----------------------------------------------------------------------------------------
%	THESIS CONTENT - CHAPTERS
%----------------------------------------------------------------------------------------

\mainmatter % Begin numeric (1,2,3...) page numbering

\fancyfoot{}
\fancyhead[RO,LE]{\thepage}
\fancyhead[LO]{\leftmark}
\fancyhead[RE]{\rightmark}

\pagestyle{fancy} % Return the page headers back to the "fancy" style

% Include the chapters of the thesis as separate files from the Chapters folder
% Uncomment the lines as you write the chapters


\chapter{Introduction}


\todo{what we are going to write}




\section{Organisation}

\todo{short introductions to each section}
%%

\chapter{Type Theory}
\label{bg}

Type theory is usually used to refer to a formal systems in which every term always has its type. It is initially invented as a foundation for mathematics as an alternative to set theory, and later works well in computer science as programming languages which we can write certified programs with type system. There are various type theories, but here we mainly study Per Martin-L\"{o}f's intuitionistic type theory (N.B.\ we will use "type theory'' specially for it if not ambiguous). It is one of the most widely studied type theory in computer science and has been implemented as programming languages like Agda, Coq, NuPRL, Epigram etc. Another type theory will play important role in this thesis is Homotopy Type Theory which is a comparatively new field between mathematics and computer science which interprets type theory using notions in homotopy type theory.

In this chapter we will introduce type theory, especially \itt, from history to rules. We will also introduce one of the implementation of intensional type theory -- Agda briefly. We will also discuss the topic of extensional concepts which are missing from \itt as a prerequisite for the topic of our thesis -- quotient types.



\section{A brief history of type theory}

The concept of \emph{sets} has been used since thounsands of years ago. However mathematician have not study the theory of sets until 1870s when George Cantor and Richard Dedekind founded set theory as a branch of mathematical logic. Since then set theory is used as a language to describe definitions of most mathematical objects, namely it is a foundational system for mathematics.

In the 1900s, Bertrand Russell discovered a paradox in this system. In this naive set theory, there was no distinction between small sets like the set of natural numbers or the set of real numbers and "larger" sets like the set of all sets. This lead to Russell's paradox.

\begin{example}[Russell's Paradox]
Let $R$ be the set of all sets which do not contain themselves
$R = \{x ~| ~x \not\in  x\}$
Then we got a cotradiction
$R \in R \iff R \not\in R$
\end{example}

To avoid this paradox, Russell
proposed the theory of types \cite{rus:1903} as an alternative to
naïve set theory. Each mathematical object is assigned a type. This is done in a hierarchical structure such that "larger" sets and small sets reside in different levels. The set of all
sets is no longer on the same level as its elements and the
paradox disappears.

The elementary notion of type theory is \emph{type} which plays a similar role to set in set theory, but differs fundamentally. Every object in type theory comes with its unique type, while an object in set theory can appear in multiple sets and we can talk about an object without knowing which set it belongs to.
To explain the difference, we use the the number $\mathsf{2}$ as an example. In set theory, $2$ is not an element of only one specific set, it belongs to the set of natural numbers $\N$ and also the set of integers $\Z$. While in type
theory, it is impossible to avoid mentioning the type of object $\mathsf{2}$. Usually the term $\AgdaInductiveConstructor{suc} \, \AgdaSymbol{(}\AgdaInductiveConstructor{suc} \, \AgdaInductiveConstructor{zero}\AgdaSymbol{)}$ stands for $\mathsf{2}$ of type $\N$ and we have a different term of type $\Z$ constructed by constructors of $\Z$ which is a different object to the one of $\N$. 



Since Russell's type theory, a variety of type theories have been developed by mathematicians and computer scientists, for example Gödel's System T \cite{gdl:1931}. There are two families of famous type theories building the bridges between mathematics and computer science, \emph{lambda calculus} and \emph{\mltt}.


\subsection{Lambda Calculus}

Alonzo Church introduced lambda calculus in the 1930s. He first introduced
an untyped lambda calculus which turned out to be inconsistent due to
the Kleene-Rosser paradox\cite{kleene1935inconsistency}.

\begin{example}[Kleene-Rosser paradox]
Suppose we have a function $f = \lambda x . \neg (x ~ x) $, then we can deduce a contradiction by applying it to itself:

$f f = (\lambda x . \neg (x ~ x)) f = \neg (f ~ f)$
\end{example}

Then he refined it with adding types. This theory is also called
Church's theory of types or simply typed lambda calculus is
introduced as a foundation of mathematics. An important change is that functions become primitive objects which means functions are types defined inductively using the $\rightarrow$ type former. It is widely applied to
various fields especially computer science. Some languages are extentions of lambda calculus, for example Haskell. Haskell belongs to one of the variants of lambda calculus called System F, although it has evolved into System FC recently. There are also other refinements of lambda calculus which is illustrated by the $\lambda$-cube \cite{barendregt1991introduction}.



\subsection{Per Martin-L\"{o}f's Type Theory}

% also cite http://www.cs.ru.nl/~herman/onderwijs/provingwithCA/paper-lncs.pdf}
% cite http://plato.stanford.edu/entries/type-theory/

In 1970s, Per Martin-L\"{o}f \cite{per:71,per:82}  developed his profound intuitionistic type theory. 


\textbf{Universe and impredicativity} The notion of a \emph{universe of types}  was first used by Martin-L\"{o}f \cite{Martin-Lof-1973} to describe the type of all types. 
His 1971's formulation has an impredicative universe i.e.\ for the universe $\mathsf{U}$, there is an axiom $\mathsf{U} : \mathsf{U}$. 
It was later proved to be inconsistent due to Girard's paradox \cite{hurkens1995simplification}.
In his later versions, a predicative hierarchy of universes are adopted and is more widely used. Briefly speaking, it starts from a universe of small types called $U_0$ and for each $n : \N$ we have $U_n : U_{n+1}$ which forms a cumulative hierarchy of universe. There is a more detailed introduction to the notion of universe written by Erik Palmgren \cite{Palmgren98onuniverses}.



It serves as a foundation of constructive mathematics \cite{martin1984intuitionistic} and also as a functional programming languages \cite{DBLP:dblp_journals/tcs/Troelstra99}. Different to set theory whose axioms are based on first-order logic, \mltt provides a means of implementing intuitionistic logic. This is achieved by the Curry-Howard
isomorphism\emph{``propositions can be interpreted as types and their
  proofs are inhabitants of these types''}. It relates computer programs with mathematical proofs. Technically, a proposition can be encoded as a type, and then a proof of it can be given by constructing a program (or term) of it. It internalises the  Brouwer–Heyting–Kolmogorov (BHK) interpretation of intuitionistic logic, for example a proof of $P \wedge Q$ is a pair of $p : P$ and $q : Q$.
Started from these, implications are functions, negations are functions into empty types, modus ponens is function application etc. This makes it a programming makes it a programming languages in which we can certify programs within the same language.
Compared to set theory, type theory has less axioms and gives us the possiblity of construct mathematics in computers so that we can do theorem proving and verification with the help of computers.

Another important feature of \mltt is dependent types. The type systems of simply lambda calculus like Haskell are not expressive enough to encode predicate logic without dependent types.

\begin{definition}\label{dpty}
\textit{Dependent type}. Dependent types are types that depends on values of other types \cite{dtw}. 
\end{definition}

As a example, it is possible to write a type for a list of natural number of length 3 as $List~\N ~3$. 

With dependent types, the quantifiers like $\forall$ and $\exists$ can be encoded.
The Curry-Howard isomorphism is then extended to predicate logic. 
A predicate on $X$ can be written as a dependent type $P ~x$ where $x : X$. 


\textbf{Equalities in Type Theory}

%cite http://adam.chlipala.net/cpdt/html/Equality.html

In regular mathematics the notion of equality are usually used to describe sameness and taken as a given.
However in type theory it is one of the most contentious topic.
There are different classifications of different notions of equalities.

Usually there are at least two important notions of equality, definitional equality and propositional equality.

\begin{definition}
\textit{Definitional equality} is a judgement-level equality, which holds when two objects have the same normal forms\cite{nor:90}.
\end{definition}

With dependent types, it is possible to write a type to encode the equality of objects.

\begin{definition}
\textit{Propositional equality} is a type which represents a propostion that two objects of the same type are equal.
\end{definition}

Intuitively if two objects are definitionally equal, they must be propositionally equal.

\begin{equation*}
\infer[\text{Id-intro}]{a = b}{a \equiv b}
\end{equation*}

But how about the other way around (we call it the \emph{equality reflecition rule})? Are two propositional equal objects definitional equal?
\begin{equation}
\label{reflection}
\infer[\text{Reflection}]{a \equiv b}{a = b}
\end{equation}
 Actually
the treatment of equality creates two different versions of \mltt, \emph{intensional} one and \emph{extensional} one.

In \itt, the answer is no. Propositional equality (also called intensional equality  \cite{nor:90}) is different to definitional equality. 
The definitional equality is always decidable hence type checking that depends on definitional equality is
decidable as well~\cite{alti:lics99}. Therefore \itt has better computational behaviors.
Types like $\mathsf{a = b}$ which stands for a
propostion that $\mathsf{a}$ equals $\mathsf{b}$ are propositional equalities. They are some types which we need to prove
or disprove by construction. Each of them has an unique element $\mathsf{refl}$ which only exists if $\mathsf{a}$ and $\mathsf{b}$ are
definitionally equal in all cases. However it is not enough for other extensional equalities, for example the equality of functions.


\ett adopted the equality reflection rule which means the propositional equality is extensional and is undistinguished with definitional equality, in other words, two propositional equal objects are judgementally equal.

There are some extensional concepts exists in \ett. For example the functional extensionality.

\begin{definition}[Functional extensionality]\label{fun-ext}
Objects of different normal forms, for example point-wise equal functions or different proofs of the same proposition, may be definitionally equal. This is called functional extensionality.

\begin{equation}
\infer[\text{fun-ext}]{f = g}{f g : A \rightarrow B, \forall a : A, f \, a = g \, a}
\end{equation}

\end{definition}



\begin{lemma}
\label{functional extensionality is available in ett}
Functional extensionality \ref{fun-ext} is derivable from equality refleciton rule \ref{reflection} (in \ett).
\end{lemma}
 
\begin{proof}
Suppose $\Gamma \vdash f \,a = g \,a$, with reflection rule we have $\Gamma \vdash f \,a \equiv g \,a$.
Then using $\xi$-rule, we know that $\Gamma \vdash \lambda a . f \,a \equiv \lambda a . g \,a$.
From $\eta$-equivalence, we know that $\Gamma \vdash f \equiv g$. We can conclude that $\Gamma \vdash f = g$.
\end{proof}

In \itt, this is not provable. If we add it as an axiom, we will lose the canonicity since we can construct a natural number with this extensionality and substitution for propositional equality.

\begin{example}[non-canonical construction]
Suppose we define $id := \lambda x . x$ and $p0 := \lambda x . x + 0$, recursively it is provable that $p': \forall x , id~ x = p0~x$, from functional extensionality we obtain a propositional equality $p: id = p0$ such that $p \equiv ext ~p'$. We can use substitution to construct non-canonical natural numbers, for instance $subst~ (\lambda f. \N) ~p ~0$. This is non-canonical because the equality is non-canonical so that the expression cannot be normalised.
\end{example}


The non-canonicity makes the definitional equality undecidable and so is type checking. The termination of type checker is not assured.


The choice whether adopt equality reflection rule or not decides the version of \mltt. Agda chooses the intensional one while NuPRL chooses the extensional one.  However we consider intensional version more as a better choice because it ihas better computational behaviors. Even though we do not get the extensional concepts automatically, we still want them in \itt and people are keeping on studying the way to introduce these into \itt.

Altenkirch and McBride introduced a variant of \ett called
\emph{Observational Type Theory}  \cite{alt:06} in which definitional equality is
decidable and propositional equality is extensional.


We already know that \mltt can be encoded as programming languages in
which the evaluation of a well-typed program always terminates \cite{nor:90}.


\section{Agda}

Agda is a dependently typed functional programming language which is designed based on intensional version
of \mltt \cite{agdawiki:main}. It is used as the main tool to study the topic of this thesis.

As we have seen, \mltt is based on the Curry-Howard
isomorphism: types are identified with propositions and terms (or programs) are identified with proofs. It turns Agda is into a proof assistant like Coq, which allows users to do mathematical reasoning and also computer program reasoning. 
Usually to prove the correctness of programs, we need to state some theorems of programming languages on the meta-level, but in Agda we can prove and use these theorems alongwith writing programs.  As Nordström et al. \cite{nps} pointed out that we could express both specifications and programs at the same time when using the type theory to construct proofs using programs.


\todo{following paragraph needs to be rewritten or added with code}
We can prove a proposition following steps below:
First we give the name of the proposition and encode it as the type. Then we can gradually refine the goal to formalise a type-correct program namely the proof. As long as we have the proof, it can be used as a lemma in other proofs or programs. Usually, there are no tactics like in Coq (it may be implemented in the future). But with the gradually refinement mechanism, the process of building proofs is very similar to conceiving proofs in regular mathematics.


There are more features of Agda as follows:

\begin{itemize}
\item \textit{Dependent type}. 
As mentioned in \ref{dpty}, dependent types are types that depends on values of other types \cite{dtw}. They enable us to write more expressive types as program speficication or propositions in order to reduce bugs. In Haskell and other Hindley-Milner style languages, types and values are clearly distinct \cite{tutorial}, In Agda, we can define types depending on values which means the border between types and values is vague. To illustrate what this means, the most common example is $\mathsf{Vector A n}$ where we can length-explicit lists called vectors. It is a data type which represents a vector containing elements of type \textbf{A} and depends on a natural number \textbf{n} which is the length of the list. We can specify types with more constraints such that the we can express what programs we can better and leave the checking work to the type chekcer. For instance, to use the length-explicit vector, we will not encounter exceptions like out of bounds in Java, since it is impossible to define such functions before compiling.

\item \textit{Functional programming language}. As the name indicates that, functional programming languages emphasizes the application of functions rather than changing data in the imperative style like C{}\verb!++! and Java. The base of functional programming is lambda calculus. The key motivation to develop functional programmming language is to eliminating the side effects which means we can ensure the result will be the same no matter how many times we input the same data. There are several generations of functional programming languages, for example Lisp, Erlang, Haskell etc. Most of the applications of them are currently in the academic fields, however as the functional programming developed, more applications will be explored.

% \item \textit{Per Martin-Löf Type Theory}. It has different names like Intuitionistic type theory or Constructive type theory and is developed by Per Martin-Löf in 1980s. It associated functional programs with proofs of mathematical propositions written as dependent types. That means we can now represent propositions we want to prove as types in Agda by dependent types and Curry-Howard isomorphism \cite{aboa}. Then we only need to construct a program of the corresponding type to prove that propostion. For example:

\end{itemize}

As a functional programming languages, Agda also has some nice features for theorem proving,

\begin{itemize}

\item \textit{Pattern matching}. The mechanism for dependently typed pattern matching is very powerful \cite{alti:pisigma-new}. We could prove propositions case by case. In fact it is similar to the approach to prove propositions case by case in regular mathematics. Pattern match is a more intuitive way to use the eliminators of types. For example, to define the negation on booleans $neg : Bool \to Bool$, we can write the program in two cases with respect to the value of the argument.

\item \textit{Inductive \& Recursive definition}. In Agda, types are often defined inductively, for example, natural numbers is defined as

\begin{code}\>\<%
\>\AgdaKeyword{data} \AgdaDatatype{ℕ} \AgdaSymbol{:} \AgdaPrimitiveType{Set} \AgdaKeyword{where}\<%
\\
\>[0]\AgdaIndent{2}{}\<[2]%
\>[2]\AgdaInductiveConstructor{zero} \AgdaSymbol{:} \AgdaDatatype{ℕ}\<%
\\
\>[0]\AgdaIndent{2}{}\<[2]%
\>[2]\AgdaInductiveConstructor{suc} \<[7]%
\>[7]\AgdaSymbol{:} \AgdaSymbol{(}\AgdaBound{n} \AgdaSymbol{:} \AgdaDatatype{ℕ}\AgdaSymbol{)} \AgdaSymbol{→} \AgdaDatatype{ℕ}\<%
\>\<\end{code}

The function for inductive types are usually written in recursive style, for example, the double function for natural numbers,

\begin{code}\>\<%
\\
\>\AgdaFunction{double} \AgdaSymbol{:} \AgdaDatatype{ℕ} \AgdaSymbol{→} \AgdaDatatype{ℕ}\<%
\\
\>\AgdaFunction{double} \AgdaInductiveConstructor{zero} \AgdaSymbol{=} \AgdaInductiveConstructor{zero}\<%
\\
\>\AgdaFunction{double} \AgdaSymbol{(}\AgdaInductiveConstructor{suc} \AgdaBound{n}\AgdaSymbol{)} \AgdaSymbol{=} \AgdaInductiveConstructor{suc} \AgdaSymbol{(}\AgdaInductiveConstructor{suc} \AgdaSymbol{(}\AgdaFunction{double} \AgdaBound{n}\AgdaSymbol{))}\<%
\>\<\end{code}

The availability of recursive definition enables programmers to prove propositions in the same manner of mathematical induction. 

\item \textit{Construction of functions}. One of the advantage of using a functional programming language as a theorem prover is the construction of functions which makes the proving more flexible.

In functional programming languages, complicated programs are commonly built gradually using aunxiliary functions and frequently used functions in the library.

Described as a proof assistant, complicated theorems are commonly proved gradually using lemmas and other theorems we have proved.

This decreases the difficulty of interpreting proofs in mathematics into Agda.

\item \textit{Lazy evaluation}. Lazy evaluation could eliminate unecessary operation because Agda is lazy to delay a computation until we need its result. It is often used to handle infinite data structures. \cite{wiki:Lazy_evaluation}

\end{itemize}

Agda also has some special functions in its interactive emacs interface beyond simple functional programming languages which enhance the ease and convenience of this language.

\begin{itemize}
\item \textit{Type Checker}. Type checker is an essential part of Agda. You can use to to type check a file without compiling it. It is the type checker that detect type mismatch problem and for theorem proving, it means the proof is incorrect. It interactively shows the goals, assumptions and variables when buiding a proof. 

The \emph{coverage checker} makes sure that the patterns cover all possible cases \cite{aboa}. 

The \emph{termination checker} will warn possiblily non-terminated error. The missing cases error will be reported by type checker. The suspected non-terminated definition can not be used by other ones. All programs must terminate in Agda so that it will not crash \cite{tutorial}.  The type checker then ensures that the proof is complete and not been proved by itself. 

In Agda, type signatures for functions are essential due to the presence of type checker (which is different to Haskell).
 
\item \textit{Interactive interface}. It has a Emacs-based interface for interactively writing and verifying proofs.  With type checker we can refine our proofs step by step \cite{aboa}. It also has some convenient functions and emacs means the potential to be extended.

\item \textit{Unicode support}. In Haskell and Coq, unicode support is not an essential part. However in Agda, to be a better theorem prover, it reads unicode symbols like: $\beta$, $\forall$ and $\exists$ and supports mixfix operators like: $+$ and $-$, which are very common for mathematics. It provides more meaningful names for types and lemmas and more flexible way to define operators. This also improve the readablity of the Agda proofs. For example, the commutativity of plus for natural numbers can be encoded as follows

\begin{code}
\>\AgdaFunction{comm} \AgdaSymbol{:} \AgdaSymbol{∀} \AgdaSymbol{(}\AgdaBound{a} \AgdaBound{b} \AgdaSymbol{:} \AgdaDatatype{ℕ}\AgdaSymbol{)} \AgdaSymbol{→} \AgdaBound{a} \AgdaFunction{+} \AgdaBound{b} \AgdaDatatype{≡} \AgdaBound{b} \AgdaFunction{+} \AgdaBound{a}\<%
\end{code}

We can use symbols we are familiar in regular mathematics.



Secondly we could use symbols to replace some common-used properties to simply the proofs a lot. The following code was simplied using several symbols,


Finally, we could use some other languages characters to define functions such as Chinese characters.

\item \textit{Code navigation}. As long as a program is loaded, it provides shortcut keys to move to the original definitions of certain object and move back. In real life programming it alleviates a great deal of work of programmers to look up the library.

\item \textit{Implicit arguments}. Sometime it is unnessary to write an argument since it can be inferred from other arguments by the type checker. It can simplify the application of functions and make the programs more concise. For example, to define a polymorphic function $\mathsf{id}$,

\begin{code}\>\<%
\\
\>\AgdaFunction{id} \AgdaSymbol{:} \AgdaSymbol{\{}\AgdaBound{A} \AgdaSymbol{:} \AgdaPrimitiveType{Set}\AgdaSymbol{\}} \AgdaSymbol{→} \AgdaBound{A} \AgdaSymbol{→} \AgdaBound{A}\<%
\\
\>\AgdaFunction{id} \AgdaBound{a} \AgdaSymbol{=} \AgdaBound{a}\<%
\end{code}

Whenever we give an argument $\mathsf{a}$,  its type $\mathsf{A}$ must be inferable.

\item \textit{Module system}. The mechanism of parametrised modules makes it possible to define generic operations and prove a whole set of generic properties.

\item \textit{Coinduction}. We can define coinductive types like streams in Agda which are typically infinite data structures. Coinductive occurences must be labelled with $\infty$ and coninductive types do not need to terminate but has to be productive. It is often used in conjunction with lazy evaluation. \cite{wiki:Coinduction}
	
\end{itemize} 

With these helpful features, Agda is a very powerful proof assisstant. It does not magically prove theorems for people, but it really helps mathematicians and computer scientists to do formalised reasoning with verification by high-performance computers. 


\subsection{Basic syntax}

Agda is a functional programming language and at the same time a theorem prover for mathematician. Its syntax has some similarities with Haskell but there are also many differences. In the other perspective, even though it is an implementation of \mltt, it adopts a syntax more closely to functional programming languages.

\begin{itemize}
\item Since Agda has a type checker, the programmer use typing judngement more often, so the designer decide to use single colon $\AgdaSymbol{:}$ for typing judgement, for example $\AgdaFunction{a} \AgdaSymbol{:} \AgdaDatatype{A}$ means that $\AgdaFunction{a}$ is of type $\AgdaDatatype{A}$, while double colons $::$ for the definition of the \emph{cons} constructor for list.

\item  The symbols for different equalities are contrary to the conventions in text. Like in some other programming languages e.g.\ Java Haskell, the equality symbol "$\AgdaSymbol{=}$" is reserved for function definition. Instead the cogruence symbol "$\AgdaDatatype{≡}$ is chosen for identity type which internalise propositional equality. 
% This is inconsistent with our conventional choices of symbols in articles, but it follows the conventions in Haskell and other programming languages that "$\AgdaSymbol{=}$" is used for definition.

\item Agda provides a more flexible way to define mixfix operators. With the unicode support, it is possible to define infix plus operator $\_+\_$, where the underscore marks the spaces for the explicit arguments in non-prefix operators. Underscores in expressions like $f ~ x ~ \_ z$ represent wildcards in function application for unnecessary  arguments (namely if can be inferred by type checker).

\item We use \textbf{data} to define constructors for inductive and coinductive datatypes. 

\item We have universe levels parameters in a lot of definitions which makes code looks unnecessarily cumbersome. 
The universe of small types is encoded as $\Set_{0}$ or $\Set$ rather than $\Type$, even though it is not a set in set-theoretical sense.
We will follow the \textbf{typical ambiguity} in this thesis which says that we write $\AgdaBound{A} \AgdaSymbol{:} \AgdaPrimitiveType{Set}$ for $\AgdaBound{A} \AgdaSymbol{:} \AgdaPrimitiveType{Set} \AgdaBound{a}$ and $\AgdaPrimitiveType{Set} \AgdaSymbol{:} \AgdaPrimitiveType{Set}$ which stands for $\AgdaPrimitiveType{Set}\AgdaBound{i} \AgdaSymbol{:} \AgdaPrimitiveType{Set}\AgdaBound{(i+1)}$.
The universe of propositions $\Prop$ ($\Prop \subset \Set$) does not exists in Agda because there is no proof-irrelevance in \itt as we will see later \ref{extensionality}. For ease of reading, we will use $\Set$ in replace of $\Prop$ in Agda code. Only necessary, we will explicitly add the proof-irrelevance property for a given proposition $P : \Set$, i.e.\ for all $p~ q : P$, $p = q$.


\item Agda has a more liberal way to define $\Pi$-types. They are often written as special cases of function types, for example $\Pi x : A. B$ can be written as $(x : A) \to B$. $\Sigma$-types are defined in Agda standard library. There is also a generalised $\Sigma$-types called \emph{dependent record type} which can be defined by keyword \textbf{record}.

\item For coinductive types and more generally mixed inductive/coinductive types \cite{txa:mpc2010g}, we adopt a set of operators which are defined in module \textbf{Coinduction}. A infinite list (or stream) can be defined as:

\begin{code}
\>\AgdaKeyword{data} \AgdaDatatype{Stream} \AgdaSymbol{(}\AgdaBound{A} \AgdaSymbol{:} \AgdaPrimitiveType{Set}\AgdaSymbol{)} \AgdaSymbol{:} \AgdaPrimitiveType{Set} \AgdaKeyword{where}\<%
\\
\>[0]\AgdaIndent{2}{}\<[2]%
\>[2]\AgdaInductiveConstructor{\_∷\_} \AgdaSymbol{:} \AgdaBound{A} \AgdaSymbol{→} \AgdaDatatype{∞} \AgdaSymbol{(}\AgdaDatatype{Stream} \AgdaBound{A}\AgdaSymbol{)} \AgdaSymbol{→} \AgdaDatatype{Stream} \AgdaBound{A}\<%
\end{code}

The delay operator $\infty$ denotes an coinductive argument. Given $a:A$, the expression with a delay function $\sharp~a$ is a element of type $\infty~A$. $\beta~x$ will force computation in $x : \infty~A$.

\item If type checker can infer the type for some arguments, we can use implicit arguments which are indicated by curly brackets.
For example a function $f:\{x:A\} \to B$ allows us to omit unnecessary argument which makes the code more readable.




\end{itemize}



\subsection{Identity Type}

Identity type is the type introduced by Martin-L\"{o}f to encode the propositional equality for definitionally equal terms \cite{nor:90}. For any two terms of $\mathsf{a\,b : A}$, we have the type $\mathsf{Id (A , a , b)}$ which is inhabitted when $\mathsf{a}$ and $\mathsf{b}$ are definitionally equal. Here we use an alternative equivalent version named after Paulin-Mohring which is parameterized with the left side of the identity. %This also includes the identity type for arbitrary universe level.

\begin{code}%
\\
\>\AgdaKeyword{data} \AgdaDatatype{\_≡\_} \AgdaSymbol{\{}\AgdaBound{A} \AgdaSymbol{:} \AgdaPrimitiveType{Set}\AgdaSymbol{\}} \AgdaSymbol{(}\AgdaBound{x} \AgdaSymbol{:} \AgdaBound{A}\AgdaSymbol{)} \AgdaSymbol{:} \AgdaBound{A} \AgdaSymbol{→} \AgdaPrimitiveType{Set} \AgdaKeyword{where}\<%
\\
\>[0]\AgdaIndent{2}{}\<[2]%
\>[2]\AgdaInductiveConstructor{refl} \AgdaSymbol{:} \AgdaBound{x} \AgdaDatatype{≡} \AgdaBound{x}\<%
\\
%
\end{code}

In Agda eliminators are not automatically derived for the types defined. Instead we have pattern matching generally which is sometimes stronger than eliminators.
As long as we pattern match on a variable of an identity type with the unique inhabitant $\AgdaInductiveConstructor{refl}$, all occurences of both variables become the same.
It is stronger and it provides the eliminator J.

\begin{code}
%
\\
\>\AgdaFunction{J} \AgdaSymbol{:} \AgdaSymbol{(}\AgdaBound{A} \AgdaSymbol{:} \AgdaPrimitiveType{Set}\AgdaSymbol{)(}\AgdaBound{a} \AgdaSymbol{:} \AgdaBound{A}\AgdaSymbol{)} \AgdaSymbol{→} \AgdaSymbol{(}\AgdaBound{P} \AgdaSymbol{:} \AgdaSymbol{(}\AgdaBound{b} \AgdaSymbol{:} \AgdaBound{A}\AgdaSymbol{)} \AgdaSymbol{→} \AgdaBound{a} \AgdaDatatype{≡} \AgdaBound{b} \AgdaSymbol{→} \AgdaPrimitiveType{Set}\AgdaSymbol{)}\<%
\\
\>[0]\AgdaIndent{2}{}\<[2]%
\>[2]\AgdaSymbol{→} \AgdaBound{P} \AgdaBound{a} \AgdaInductiveConstructor{refl}\<%
\\
\>[0]\AgdaIndent{2}{}\<[2]%
\>[2]\AgdaSymbol{→} \AgdaSymbol{(}\AgdaBound{b} \AgdaSymbol{:} \AgdaBound{A}\AgdaSymbol{)(}\AgdaBound{p} \AgdaSymbol{:} \AgdaBound{a} \AgdaDatatype{≡} \AgdaBound{b}\AgdaSymbol{)} \AgdaSymbol{→} \AgdaBound{P} \AgdaBound{b} \AgdaBound{p}\<%
\\
\>\AgdaFunction{J} \AgdaBound{A} \AgdaSymbol{.}\AgdaBound{b} \AgdaBound{P} \AgdaBound{m} \AgdaBound{b} \AgdaInductiveConstructor{refl} \AgdaSymbol{=} \AgdaBound{m}\<%
\\
%
\end{code}

\subsection{Extensionality}
\label{extensionality}

In regular mathematics, equality does not only exists between intensionally equal terms. Objects are also equal if they have the same extensional external properties, like functions and propositions. 

As Martin Hofmann summarises in \cite{hof:phd}, there are several important extensional concepts which we will also encounter in this thesis, \emph{Functional extensionality}, \emph{Uniqueness of identity}, \emph{Proof-irrelevance}, \emph{Propositional extensionality}, \emph{Quotient types}. 
These principles are not currently available in \itt, but they are valid extensions of Type Theory and worth interpreting to help both Mathematics and programs constructions. Intuitively speaking, the extensionally equal terms can be distingushed by any other terms, hence the extensionality is justifiable. 

\begin{itemize}
\item \textit{Functional extensionality} It has been introduced in \ref{fun-ext}.

\item \textit{Uniqueness of identity} In type theory we define a notion of \emph{set} as follows:

\begin{definition}\label{UIP}
A type $A$ is a \textbf{set} if for all $x,y:A$ and all $p,q:x=y$, we have $p=q$.
\end{definition}

This property is usually called uniqueness of identity (UIP). From the usual definition of identity types, UIP is not a result for every type. UIP is also equivalent to Streicher’s “Axiom K”.

\begin{axiom}[K]
For all $x:A$ and $p: x=x$ we have $p=\text{refl}_{x}$.
\end{axiom}

Another characterisation of a set in type theory is given by Hedberg's Theorem.
\begin{theorem}[Hedberg]
If $A$ has decidable equality, then $A$ is a set.
\end{theorem}

\item \textit{Proof-irrelevance} All proofs of the same proposition are propositionally equal.

\item \textit{Propositional extensionality} Two logically equivalent propositions are propositionally equal.

\item \textit{Quotient types} A quotient type is a type formed by redefing equality on a underlying type with a given equivalence relation on it. 

\item \textit{Univalence Axiom}
In \hott, univalence is a extensional principle which states that ismorphic types are propositionally equal.
\end{itemize}

It is interesting to extend Type Theory with these extensional principles. However, it only makes sense if the type-checking decidability and terms canonicity are not sacrificed. This thesis mainly focus on the extension of quotient types.


\section{Homotopy Type Theory}

\hott is a variant of intensional \mltt{} which is a new branch developed between theoretical computer science
and mathematics. Vladimir Voevodsky found a surprising connection between homotopy theory and type theory \cite{voe:06}. He proposed the univalence axiom, which identifies isomorphic structures, as a univalent foundation for mathematics. 
In \hott, there is an observation that notions in type theory can be interpreted by homotopy-theoretical terms. A type is regarded as a \emph{space} and a term of this type is a \emph{point} of this space. Functions between types are \emph{continous maps} and identity types are usually considered as \emph{paths}. Identity types of identity types are \emph{homotopies}. Although these notions are originally defined with topological bases, we only employ them as homotopical notions on a higher level. 


As univalence axiom states, equality is equivalent to equivalence. Acutally it can be seen as an formal acceptance of the ''common sense'' in Mathematics that isomorphic structures can be identitfied. The higher structures of the equivalence also allows us to study the different ways of identification. Therefore it is more appropriate to interpret types as higher groupoids. People is trying to implement \hott in \itt and one possible way is to interpret \wog first in Agda. The author has done some work in this direction which can be found in Chapter~\ref{wog}.


I will not explain this topic in detail here, but in \autoref{HITs}
we give a more detailed introduction to it.



 % Something without a type makes no sense to us because we are not sure what it stands for and how do we use it. The type definition describe the syntax so that some symbol makes sense and the semantic meaning may be revealed from the construction.

% There is another question, whether mathematics is a collection of patterns and laws which is observed, or it is a system created and built by people to explain the patterns and laws in the world. I think people prefer the second answer usually accept the type theory more easily, although most people (probably 99.9 percent) prefer the first one. When we learn what is natural numbers, we learn it as "numbers like 1, 2, 3 ,4 and perhaps 0", the commutative law, associate law are axioms because there is no way to prove it if we introduce it in this manner. We are convinced by some examples like "2 + 3 = 3 + 2" and we find it works for most of the cases then we accept it by observations. It is some methods physicians used a lot -- to conclude some laws from a number of facts. It is a proper method for physicians because what they research on is world can only be observed. However for mathematics, even though it is applied to the real world, it is a system completely created by people. People used their fingers to count, wrote symbols for results, even though it was very shallow it is obviously a aritificial system. People extend 


% Type theory is strongly connected with computation theory.
% Set


% Type theory has fewer axioms, simpler model than set theory which has mutual foudations: logic and axioms.

\section{Summary}


In mathematics, set theory is still a more popular choice over type theory. However in computer science, instuitionistic type theory is worth more studying. It is more close to program construction and from a computer scientist's point of view, it is very natural to accept intuitionistic logic. 

It provides a foundation of mathematics which can be implemented as a programming language so that proving is just programming and verification is just type checking. The aid of computers saves a lot of work from mathematicians and reduces the chance of making mistakes, although the absence of \emph{principle of excluded middle} in intuitionistic logic makes some mathematicians hard to accept. There is a very good talk given by Andrej Bauer in IAS called "Five Stages of Accepting Constructive Mathematics'' online \footnote{available on Youtube}.

\paragraph{Summary of different versions of Per Martin L\"{o}f's intuintionsitc type theory}

%The basic notion in Martin-L¨of’s type theory is the notion of type. A type
%is explained by saying what an object of the type is and what it means for
%two objects of the type to be identical. This means that we can make the
%judgement -- http://www.cse.chalmers.se/research/group/logic/TypesSS05/Extra/nordstrom.pdf


\begin{enumerate}

\item The 71's formulation \cite{per:71} is impredicative and turned our to be inconsistent due to Girard's paradox.

\item The 72's version which is published in 1996 \cite{Martin-Lof-1972} abandons the impredicative universe and all later versions are predicative. It does not have inductive identity types but recursively defines equalities for given types.

\item The 73's version \cite{Martin-Lof-1973} firstly introduced inductively defined identity types to interalising equality as types. 

\item The 80's version which is summarised by Giovanni Sambin in 1984 \cite{martin1984intuitionistic} is extensional. It adopts equality reflection, namely an inhabitant of an identity type implies definitionally equality.

\item In the homotopic version \cite{hott}, Vladimir Voevodsky extends it with univalence axiom and provides a homotopic interpretation of it.

\end{enumerate}








%\input{Chapters/Extensionality}

\chapter{Quotient Types}
\label{qt}
% Quotient type or quotient types?


%Quotient type is a type generated by redefining equality on a given type. 

%Technically we cannot redifine equality

 %Generally speaking, given a setoid $A$ and an 
 %equivalence relation $\sim$ on it, a quotient type denoted as
%$\qset{A}$ is a type generated by redefining equality of $A$ by
%$\sim$. In \itt, we usually use setoids to represent quotients but it
%is not satisfactory. It is a very important topic to extend \itt with
%quotient types.

%From a programmming persepective, there are more tools to define data
%types. 
%In the chapter, we first introduce the quotients in Mathematics and
%then quotient types in type theory with categorically interpretations.


%\section{Quotients in Mathematics}

% "M"athematics must be capital


%structures  used also extended to other branches, such as set
%theory, group theory, topology etc.

%we usually use the same operators for abstractly similar 
%operations. For instance, the product is extended to the cartesian product in set theory.


Quotient is one of the primitive notions in \maths. 
In arithmetic, quotient is the result of division

$$8 \div 4 = 2 ~~ \text{or}~~ 8/4 = 2$$


In abstract algebra quotient is the result of ``dividing'' a set, group, space or another algebraic structure (dividend) with respect to an given equivalence relation on it. We use quotient set as an example here.

\begin{definition}
\textbf{Equivalence relation}.
An equivalence relation is a binary
relation which is reflexive, symmetric and transitive.
\end{definition}

Intuitively, given any equivalence relation, a set can be partitioned into
some equivalence classes,

\begin{definition}
\textbf{Equivalence class}.
The equivalence class of an element $a$ is a set whose elements are
all equivalent to $a$
\begin{equation}
\class a = \{x : A \;| \; a \sim x \}
\end{equation}
\end{definition}

The collection of these equivalence classes is called
the quotient set.

\begin{definition}
\textbf{Quotient set}.
Given a set A equipped with an equivalence relation $\sim$, a quotient
set is denoted as $\qset{A}$,
\begin{equation}
\qset{A} = \{ [ a ] \;|\; a : A \}
\end{equation}
\end{definition}

Similarly in topology, group theory or other branches of \maths there are also quotient space, quotient group and other quotients. 

Naturally one would also expect quotient types (or quotients) in Type Theory. Intuitively speaking, a \emph{quotient type} $\qset{A}$ is a type $A$ whose equality is redefined by an equivalence relation. 
Unlike \ett, there is no quotient operator in traditional \itt, but setoids are usually used to represent quotients.

\begin{definition}
\textbf{Setoid}.
\noindent A setoid $(A,\sim,\text{eqv}_{\sim})$ (usually written as just $(A,\sim)$) consists of
\begin{enumerate}
\item a set (type) $A : \Set$,
\item a binary relation $\sim : A \to A \to \Prop$, and
\item it is an equivalence, i.e.\ the proofs that it is reflexive, symmetric and transitive.
\end{enumerate}
\end{definition}

%Notice that $\Prop$ here is not necessarily a universe of propositions. In an \itt like Agda, we use $\Set$ instead (in standard library) which does not reflect the fact that the relation is propositional. Indeed we can fix it by adding additional condition that the relation is propositional.

\begin{remark}
This notion is also called a \emph{total setoid}. If the relation of a total setoid is not required to be reflexive, then it becomes a \emph{partial setoid}. Barthe Capretta and Pons have studied different definitions of setoids and categories of setoids in \cite{bar:03}. In this thesis, the term "setoid" refers to total setoids in general.
\end{remark}

However using setoids to represents quotients is not an ideal solution. Since it is an alternative representation of sets, everything defined on $\Set$ has to be redefined on $\Setoid$ again. 
For example, functions on setoids, equalities on setoids,
products on setoids etc. Moreover, considering a setoid $(A,\sim)$ whose $A$ is also encoded as a setoid, 
it is essential to define higher setoids.
From a programming perspective, setoids are also
unsafe because we have access to the underlying
 sets \cite{aan}. We can perform operations which do not respect
 equivalence relation and the result does not make any sense.
Therefore, it is essential to use a quotient of type $\Set$ which can be uniformly treated with other types. It is also the case in various branches of mathematics that the quotient is the same kind of object as the base one.

In this chapter, we present a definition of quotient type in an \itt extended with a proof-irrelevant universe of propositions in the sense of \autoref{ittprop}. We also explain the categorical meanings of quotient types.


\section{Quotients in Type Theory}

\subsection{Rules for quotients}\label{iqs}

The quotient types are defined by the following rules as described in \cite{Jacobs94quotientsin,hof:95:sm}. 


\infrule[Q-\bf{Form}]
{ \Gamma \vdash A  \andalso \Gamma ,x : A , y : A \vdash x \sim y : \bf{Prop} \andalso \sim~\text{is an equivalence}}
{\Gamma \vdash \qset{A}}

Given a type $A$ and with a binary equivalence relation $\sim$ on $A$, we can form the quotient $\qset{A}$. Here we use an infix notation for readability.

The equivalence properties includes


\begin{multicols}{2}
\infrule[$\sim$-\bf{refl}]
{ \Gamma \vdash a : A}
{\Gamma \vdash \text{ref}_{\sim}(a) : a \sim a}
\columnbreak
\infrule[$\sim$-\bf{sym}]
{ \Gamma \vdash a : A \andalso \Gamma \vdash b : A \andalso \Gamma \vdash p : a \sim b}
{\Gamma \vdash \text{sym}_{\sim}(p) : b \sim a}
\end{multicols}




\infrule[$\sim$-\bf{trans}]
{ \Gamma \vdash a : A \andalso \Gamma \vdash b : A \andalso \Gamma \vdash c : A \\
\Gamma \vdash p : a \sim b \andalso \Gamma \vdash q : b \sim c}
{\Gamma \vdash \text{trn}_{\sim}(p,q) : a \sim c}

\begin{remark}
Notice that the formation rule is different to Hofmann's version \cite{hof:95:sm} where $\sim$ is not required to be only equivalence relation. In fact his version is just more general which accepts non equivalence relation $R : A \to A \to \Prop$, but $A/R$ has to be understood as the quotient of $A$ by the equivalence closure of $R$.
\end{remark}

\begin{multicols}{2}
\infrule[Q-\bf{Intro}]{\Gamma \vdash a : A}{\Gamma \vdash [ a ] : \qset{A}}
\columnbreak
\infrule[Q-\bf{Ax}]
{\Gamma \vdash a , b :  A  \andalso  \Gamma \vdash  p : a \sim b}
{ \Gamma \vdash \text{Qax}(p) : [a]=_{\qset{A}} [b]}
\end{multicols}


We introduce an ``equivalence class'' for each element of $A$. It is usually denoted as $[ a ]$, or $[ a ]_{\sim}$ for $\sim$ if it is unclear which relation it refers to. 
 Qax states that the ``equivalence classes'' of two terms which are related by $\sim$ are (propositionally) equal.

 Notice that the notation of terms $[a]$ can be confused with the substitution notation like $B[a]$ or $B [[ a ]/x]$. So for a $\Pi$-type $B : (x : A) \to \Set$ and $a : A$, we write $B(a) : \Set$ for $B [[ a ]/x]$ where order of the arguments in brackets matches its definition.

% $for any possibly ambiguous occurrence, we use full notations i.e.\ $[a]_{\sim}$  for terms and $B[a/x]$ for substitution.$


In Hofmann's \cite{hof:95:sm} definition, it comes with an eliminator (also called \emph{lifting}) with its computation rule ($\beta$-rule) and an induction principle (equivalent to a $\eta$-rule): \footnote{We use shorthand notation $\hat{}$ for lifting here} 


\infrule[Q-\bf{elim}]
{\Gamma \vdash  B \andalso \Gamma \vdash f : A \to B \\
\Gamma, a: A, b : A, p : a \sim b \vdash  f^{\sim}(a,b,p) : f(a) =_{B}  f(b) \andalso \Gamma \vdash  q : \qset{A}}
{\Gamma \vdash  \hat{f}(q) : B}

\infrule[Q-\bf{comp}]{\Gamma \vdash  a : A}{\Gamma \vdash  \text{Qcomp}(a)  : \hat{f}([ a ]) \equiv f(a) }


\infrule[Q-\bf{ind}]
{\Gamma,  x : \qset{A} \vdash P : \Prop \andalso \Gamma, a : A \vdash h(a) : P([a]) \andalso \Gamma \vdash  q : \qset{A}}
{\Gamma \vdash \text{Qind}(h,q) : P(q)}

Given a function $f : A \to B$ which respects $\sim$, we can lift it to be a function on $\qset{A}$ as $\hat{f} : \qset{A} \to B$ such that for any element $a : A$, $\hat{f}([ a ])$ computes to the same value as $f(a)$. It allows us to define functions on quotient types by functions on base types (representatives).
%Notice that \emph{function application} is written like "$f ~x ~ y$'' where $x$ and $y$ are two arguments for function $f$. 
Notice that We omit $f^=$ since the computation rule already implies that it is proof-irrelevant.

The induction principle states that for any proposition $P : \qset{A} \to \Prop$. it is enough to just consider cases $ P([ a ])$ for all $a : A$. In other words, $\qset{A}$ only consists of "equivalence classes" i.e.\ $[ a ]$.


Alternatively, a \emph{dependent} eliminator (dependent lifting) serves the same purpose:

\infrule[Q-\bf{dep-elim}]
{\Gamma, x : \qset{A} \vdash B \andalso \Gamma \vdash f : (a : A) \to B([ a ]) \\
\Gamma, a : A, b : A, p : a \sim b \vdash f^= (a,b,p) : f(a)
\stackrel{p}{=}  f(b) \andalso \Gamma \vdash q :
\qset{A}}
{\Gamma \vdash \hat{f}(q) : B(q)}

\infrule[Q-\bf{dep-comp}]{\Gamma \vdash a : A}
{\Gamma \vdash \text{Qdcomp} (a)  : \hat{f}([ a ]) \equiv f(a) }


Notice that $\stackrel{p}{=}$ is an abbreviation for propositional equality which requires substitution in the type of the left hand side by $\text{Qax}(p)$ so that both sides have the same type.
We use the same notation for the two versions of eliminators because they are in fact equivalent (see \autoref{elim-equiv}).

\begin{proposition}\label{elim-equiv}
The \emph{non-dependent eliminator} with the \emph{induction principle} is equivalent to the \emph{dependent eliminator}.
\end{proposition}
\begin{proof}
1. Assume we have the non-dependent eliminator and the induction principle, $B$ is a dependent type on $\qset{A}$, $f$ is a dependent function of type $(a : A) \to B([ a ])$ and it respects $\sim$ under substitution (i.e.\ $f^=$), $q$ is an element of $\qset{A}$.

Set $B'$ as a dependent product $\Sigma (r : \qset{A}) ~B[r/x]$,

Then a non-dependent version of $f$ which is of type $A \to B'$ can be defined as

$$f' a \defeq [ a ] , f ~a$$

Given $p : a \sim b$, we can conclude that $f'(a) =_{B'} f'(b)$ is inhabited from Qax and $f^=$.

It allows us to lift the non-dependent function $f'$ as $\hat{f'}$ such that 

\begin{equation}\label{f'comp}
\hat{f'} ([ a ]) \equiv [ a ] , f ~a
\end{equation}
Applying first projection on both sides of \ref{f'comp}, the following propositional equality is inhabited:

 $$\pi_1 (\hat{f'}([ a ])) = [ a ]$$

By induction principle, the predicate $P : \qset{A} \to \Prop$ defined as

$$P(q) \defeq \pi_1 ~(\hat{f'}(q)) =_{\qset{A}} q$$

is inhabited for all $q : \qset{A}$.

Finally, to complete the dependent eliminator, we can construct an element of type $B[q/x]$ by

$$\pi_2 ~(\hat{f'}(q))$$

which has the correct type because $P(q)$ holds. The computation rule is simply derivable from \ref{f'comp}.

2. It is easy to find out that the non-dependent eliminator and induction principle are just special cases of dependent eliminator.

A constructive version of this proof in Agda can be found in \autoref{app:dq}.
\end{proof}


Additionally, a quotient is effective (or exact) if an "equivalence class" only contains terms that are related by $\sim$.

\infrule[Q-\bf{effective}]
{\Gamma \vdash a :  A \andalso \Gamma \vdash b :  A  \andalso p : [a] =_{\qset{A}} [b] }
{\text{eff}(p) : a \sim b}

In fact propositional extensionality implies that all quotients defined with \emph{equivalence} relation are effective.

\begin{proposition}
With propositional extensionality, we can prove that all quotients\footnote{we assume $\sim$ is an equivalence relation} are effective.
\end{proposition}

\begin{proof}\label{PUEF}
%Suppose we have a set $A$ with an equivalence relation $\sim : A \to A \to \Prop$, a quotient set is $\qset{A}$.

Suppose we have a quotient type $\qset{A}$, two elements of $a, b : A$ and $[ a ] =_{\qset{A}} [ b ]$


Set a predicate $P_a : A \to \Prop$ as 
$$P_a (x) \defeq a \sim x$$

$P_a$ is compatible with $\sim$ since

$x \sim y$

$\Rightarrow a \sim x \iff a \sim y$ (symmetry and transitivity)

$\equiv P_a(x) \iff P_a(y)$ (propositional extensionality)

$\Rightarrow P_a(x) = P_a(y)$

Therefore we can lift $P_a$\footnote{The elimination rule applies to large types} such that for any $x : A$

$$\hat{P} ([ x ]) \equiv a \sim x$$


We can simply deduce $\hat{P} ([ a ]) = \hat{P} ([ b ])$ from assumption $[ a ] = [ b ]$ which by definition
is just

$$a \sim a  = a \sim b$$ 

Finally with eliminator $J$ and $\text{ref}(a) : a \sim a$ we can easily prove

$$a \sim b$$.

\end{proof}

%Alternatively, we can prove it as follows:

%\begin{proof}

%Firstly we prove the equivalence relation is well-defined on the
%quotient types, namely it respects the equivalence relation:

%Suppose we have $a \sim b$ and $c \sim d$, we can deduce $a \sim c \iff
%b \sim d$. Then applying the propositional univalence axiom, we know
%that $a \sim c = b \sim d$, hence the equivalence relation is
%well-defined.

%Because it is well-defined, we can lift it such that

%$[ a ] ~\hat{\sim}~ [ b ] \equiv a \sim b$


%From reflexivity of the equivalence relation, $\forall x : A, x \sim x$, 
%we know that $\forall x : A, [x]~\hat{\sim}~[x]$.

%Assume $[a]=[b]$, using J-eliminator in $[a]~\hat{\sim}~[a]$
%(reflexivity), $[a]~\hat{\sim}~[b]$ which is definitionally equal to $a \sim b$, Hence the quotient is
%effective.
%\end{proof}


\section{Quotients are coequalizers}

%\begin{remark}\label{equivalencerelationiso}

%\end{remark}

The rules of quotient types can be characterised in a category-theoretical way.

Categorically speaking, a quotient is a \textbf{coequalizer}. Let us recall the definition.


\begin{definition}
\textbf{Coequalizer}.
Given two objects $X$ and $Y$ and two parallel morphisms $f, g : \morph{X}{Y}$ , a coequalizer is an object $Q$ with a morphism $q : \morph{Y}{Q}$ such that $q \circ f = q \circ g$ and it is universal: any pair $(Q' , q')$ satisfying $q' \circ f = q' \circ g$ has a unique factorisation $u$ such that $q' = u \circ q$:
\begin{displaymath}
    \xymatrix{X \ar@<0.5ex>[rr]^f \ar@<-0.5ex>[rr]_g && Y \ar[rr]^q
      \ar[ddrr]_{q'} && Q
      \ar@{.>}[dd]^u \\ \\
& &&& Q' }
\end{displaymath}
\end{definition}


In a cartesian closed category, e.g.\ $\Set$, assume

$$R \defeq \Sigma (a_1,a_2 : A) ~ a_1 \sim a_2$$

%$$R = \{(a_1,a_2) : A \times A ~|~ a_1 \sim a_2\}$$

and the two projections are two parallel morphisms $\pi_1,\pi_2 : \morph{R}{A}$,

A quotient is just the coequalizer $(\qset{A},[\_])$ such that the following diagram commutes:
\begin{displaymath}
    \xymatrix{R \ar@<0.5ex>[rr]^{\pi_1} \ar@<-0.5ex>[rr]_{\pi_2} && A \ar[rr]^{ [\_]}
      \ar[ddrr]_{f} && \qset{A}
      \ar@{.>}[dd]^{\hat{f}} \\ \\
& &&& B}
\end{displaymath}

The factorisation $\hat{\_}$ is just the eliminator, the computation rule and induction principle corresponds to the universal property of it. 

\begin{proposition}
The induction principle implies uniqueness, and is also derivable from the definition of coequalizer.
\end{proposition}

\begin{proof}
It is easy to see that induction principle implies the uniqueness of $\hat{f}$:

Given any $g : \qset{A} \to B$ fulfils the same property as $\hat{f}$, applying induction principle on 

$$\forall (a :A) \to g([a]) = \hat{f}([a])$$

we can deduce that $$\forall (q : \qset{A}) \to g(q) = \hat{f}(q)$$ hence $g = \hat{f}$.

The other way is more difficult:

Given $P : A \to \Prop$, $h : (x : A) \to P (x)$ define

$$P' \defeq \Sigma(x : A) ~P(x) ~\text{and}~ h'(x) = ([x],h(x))$$

we can observe that 

\begin{equation}\label{eq1}
\pi_1 \circ h' = [\_]
\end{equation}

By universal property, there is a unique $\hat{h'}$ s.t. 

\begin{equation}\label{eq2}
\hat{h'} \circ [\_] = h'
\end{equation}

By replacing \ref{eq2} in \ref{eq1}

\begin{equation}\label{eq3}
\pi_1 \circ \hat{h'} \circ [\_] = [\_]
\end{equation}

From \emph{uniqueness} we can easily prove that $[\_]$ is an \emph{epimorphism} as follows:

given $f_1,f_2 : \qset{A} \to B$ s.t. $f_1 \circ [\_] = f_2 \circ [\_]$

Then $f_1 \circ [\_] \circ \pi_1 = f_1 \circ [\_] \circ \pi_2$ and $f_2 \circ [\_] \circ \pi_1 = f_2 \circ [\_] \circ \pi_2$ there is a unique $g : \qset{A} \to B$ such that

$g \circ [\_] = f_1 \circ [\_] = f_2 \circ [\_]$

Because both $f_1$ and $f_2$ fulfil this property, $f_1 = f_2$, hence $[\_]$ is epi.


Thus from \ref{eq3}, we prove that

\begin{equation}\label{eq4}
\pi_1 \circ \hat{h'} = id
\end{equation}

Thus for any $q : \qset{A}$, the type of $\pi_2(\hat{h'}(q))$ is

$P(\pi_1(\hat{h'}(q))) = P(q)$

as expected, hence we derived the induction principle. 
In fact, following the same procedure, a dependent eliminator is also derivable.
\end{proof}

The coequalizer (quotient) is effective if the following diagram is a \emph{pullback}

\[\xymatrix{
R\pullbackcorner\ar[rr]^{\pi_1}\ar[dd]_{\pi_2} && A\ar[dd]^{[\_]} \\ \\
A\ar[rr]_{[\_]} && Q
}\]

\begin{proof}
Assume we have two points $a, b : \textbf{1} \to A$ satisfying $[ a ] = [ b ]$.

From the pullback property, there is a unique point $r : \textbf{1} \to R$ such that

$$\pi_1(r) = a$$ and $$\pi_2(r) = b$$ 

Hence $(a, b)$ is an element of $R$, by definition it means

$$a \sim b$$
\end{proof}

In \autoref{dq}, we also introduce two other notions: prequotient and definable quotient.

Cateogorically speaking, a \emph{prequotient} is just a \emph{fork}

\[\xymatrix{
R\ar@<0.5ex>[rr]^{\pi_0}\ar@<-0.5ex>[rr]_{\pi_1}& & A\ar[rr]^{[\_]}
& & Q
}\]

and a \emph{definable quotient} corresponds to a \emph{split coequalizer} which is a \emph{fork} with two morphisms $\emb : Q \to A$ and $t : A \to R$ such that 

$[\_] \circ \emb = 1_Q$

$\emb \circ [\_]  = \pi_2 \circ t$ and 

$\pi_1 \circ t = 1_A$.

In fact, $t$ corresponds to the "complete" property of definable quotients as $t(a) = (\emb [ a ] , a)$.

\section{Quotient is a left adjoint functor}

As Jacobs\cite{Jacobs94quotientsin} suggests, quotient can be described as a left adjoint to an equality functor.

Let us recall the definition first.
\begin{definition}
\textbf{Adjunction}.
Given two categories $A$ $B$, a functor $F : A \to B$ is left adjoint
to $G : B \to A$ if we have a natural isomorphism
$\Phi : \text{hom}_{B}(F ~\_ , \_) \to hom_{A}(\_, G ~\_)$
\end{definition}


Given the category of setoids $\textbf{Setoid}$ and category of sets $\textbf{Set}$, there is an equality functor $\nabla : \textbf{Set} \to \textbf{Setoid}$ defined as
$$\nabla A \defeq (A , =_{A})$$

where the morphism part is trivial embedding.

Quotients can be seen as a functor $\textbf{Q} : \textbf{Setoid} \to \textbf{Set}$ which is
left-adjoint to a equality functor $\nabla A \defeq (A , =_{A})$

The object part of this functor corresponds to the formation rule of quotients, hence we can use $\qset{B}$
to represent $\textbf{Q} ~ (B , \sim)$.

The adjunction can be described by a natural isomorphism

$\Phi : \text{hom}_{\text{Set}}(\textbf{Q}  ~\_ , \_) \to \text{hom}_{\text{Setoid}}(\_ , \nabla ~\_)$

or a diagram for each $(Y, \sim) : \textbf{Setoid}$ and $X : \textbf{Set}$:

\begin{equation*}
\begin{aligned}
\qset{Y} & \to X \\
\midrule
\midrule
(Y, \sim) & \to (X , =_{X})
\end{aligned}
\end{equation*}

%\begin{equation*}
%\begin{aligned}
%\text{hom}_{\text{Sets}}(\qset{Y},X) \\
%\midrule
%\midrule
%\text{hom}_{\text{Setoids}}((Y, \sim), (X , =))
%\end{aligned}
%\end{equation*}


which consists of $\Phi_{(Y, \sim),X}$ and its inverse $\Phi^{-1}_{(Y, \sim),X}$ (the subscripts are omitted later).


Given an identity morphism $\text{id} : \qset{A} \to \qset{A}$, 

$$\Phi(\text{id}) : (A, \sim) \to (\qset{A}, =_{\qset{A}})$$

is just the introduction rule $[\_] : A \to \qset{A}$ with the property that it is compatible with $\sim$.
It is also called \emph{unit} written as $\eta_{(A,\sim)}$.

Given a morphism $f : (A, \sim) \to (B , =_{B})$ which is a function compatible with $\sim$,

$$\Phi^{-1}(f) : \qset{A} \to B$$ 

which corresponds to the elimination rule.

The computation rule $\hat{f} \circ [\_] \equiv f$ corresponds to the following digram in the category of setoids:


\[\xymatrix{
(A,\sim)\ar@<-0.5ex>[dd]_{\eta_{(A,\sim)}}\ar@<0.5ex>[ddrr]^{f} \\ \\
(\qset{A}, =_{\qset{A}})\ar@<-0.5ex>[rr]_{\nabla(\Phi^{-1}(f))} & & (B , =_{B})
}\]

which is commutative because

$\nabla (\Phi^{-1}(f)) \circ \eta_{(A,\sim)}$

$=\Phi(\Phi^{-1}(f))$ by adjunction law $G(f) \circ \eta_{Y} = \Phi(f)$

$=f$



We can also recover the adjunction from the definition of quotients. Define

$$\textbf{Q} ~ (Y , \sim) \defeq \qset{Y}$$

The adjunction is given by

$$\Phi(f) \defeq f \circ [\_] ~\text{and} ~ \Phi^{-1}  (g, g^{\sim}) \defeq \hat{g}$$

The computation rule and induction principle just express that these two mapping are each other inverses.

%which is compatible with $\sim$ because $[\_]$ is compatible with $\sim$.

%We can define its inverse as

%: assume $f : \qset{Y} \to X$, there is a mapping

%It is an isomorphism because it has an inverse: assume $(g, g^{\sim}) : (Y, \sim) \to (X, =_{X})$


%such that

%$\Phi^{-1} (\Phi(f))= \dlift{f \circ [\_]} = f$

%because for all $y : Y$, $\dlift{f \circ [\_]} ([ y ]) \equiv (f \circ [\_]) (y) \equiv f ([ y ])$




%where $g$ is a function compatible with $\sim$

%To prove that this is an isomorphism, assume $f : \qset{B} \to A$, there exists a mapping

%$\Omega ~ f = f \circ [\_]$ and its inverse

%$\Omega^{-1} ~ g = \hat{g}$ (we omit the property that $g$ respects $\sim$)

%The isomorphic properties can be verified as follows,

%$\Omega^{-1} (\Omega~f) = \dlift{f \circ [\_]} = f$ by the uniqueness from lifting.

%$\Omega (\Omega^{-1}~g) = \hat{g} \circ [\_] = g$ by definition of lifting.


\section{Quotients in \hott}\label{qthott}

As we mentioned before, quotient types (in the sense of \ref{iqs}) are available in \hott. Because of the different interpretations of types, it makes less confusion to call them \emph{quotients} or \emph{set quotients} here.

First, let us recall that

\begin{itemize}
\item an h-proposition (hProp) $P$ is a type which has the property $\forall(a, b : A) \to a =_{A} b$, and

\item an hSet $S$ is a type such that forall $x, y : S$, $x =_{S} y$ are h-propositions.
\end{itemize}

For simplicity, we use "sets" for h-sets and "propositions" for h-propositions here. Also $\Prop$ is \emph{not} the built-in universe of propositions in Coq, but the universe of h-propositions.

\subsection{An impredicative encoding of quotient sets}\label{impredicative}

Vladimir Voevodsky has introduced an impredicative definition of quotients which has been encoded in Coq \cite{voe:hset}. 

%\footnote{\url{http://www.cse.chalmers.se/~coquand/cirm.pdf}} 

Assume we have a type $A$ (which is not necessarily a set) and an equivalence relation $\sim : A \to A \to \Prop$. 

\begin{definition}
An \textbf{equivalence class} is a predicate $P : A \to \Prop$ such that

it is inhabited $\exists (a : A) ~P (a)$

For all $x, y : A$,

$P(x) \to P(y) \to x \sim y$

$P(x) \to x \sim y \to P(y) $

The properties can be encoded as 

$$\text{EqClass}(P) \defeq (\exists (a : A) ~ P (a)) \wedge (\forall (x, y : A) \to P(x) \to (x \sim y \iff P(y)))$$
\end{definition}

\begin{definition}
We define the \textbf{set quotient} as

$\qset{A} \defeq \Sigma (P : A \to \Prop) ~ \text{EqClass}(P)$
\end{definition}

$\qset{A}$ is a set because $A \to \Prop$ is a set and $\text{EqClass}(P)$ is a proposition. $\wedge$ is the non-dependent $\Sigma$-type for propositions and $\forall$ is the $\Pi$-type for propositions. The encoding of $\exists(a : A)~ P(a)$ is given by a truncated $\Sigma$-type: $\| \Sigma(a : A)~P(a) \|$.
The (-1)-truncation $\|-\|$ is defined impredicatively as

$$\|X\| \defeq \forall (P : \Prop) \to (X \to P) \to P$$

with a trivial embedding function $|\_| : X \to \| X \|$:

$$| x | \defeq \lambda P ~ f \to f(x)$$

Note that $\|X\|$ is in the universe of $\Set_1$, but with \textbf{resizing rules} proposed by Voevodsky \cite{Universe-poly,RR}, $\|X\|$ is moved to the universe of $\Set$. We can apply the resizing rule for proposition because $\|X\|$ behaves like a proposition.
It also has to be noticed that it is impossible to extract an element of $A$ from an proof of $\text{EqClass} ~ P$ because of the impredicative encoding. 


There is a canonical function $[\_] : A \to \qset{A}$ defined as

$$[ a ] \defeq (\lambda x \to a \sim x, | a, \text{ref}(a) | , \lambda x~y ~p \to (\lambda q \to\text{trn}(p,q), \lambda q \to \text{trn}(\text{sym}(p),q)))$$

which is compatible with $\sim$, or in another word, it respects $\sim$.

\begin{definition}
We say a function $f : A \to B$ is \textbf{compatible} with $\sim$ if

$$\forall (x, y : A) \to x \sim y \to f(x) =_{B} f(y)$$
\end{definition}

 The verification of compatibility requires propositional extensionality and functional extensionality which are available in \hott.


Given a function $f : A \to B$ into a \emph{set} $B$ which is compatible with $\sim$, there exists a unique function $\hat{f} = \qset{A} \to B$ such that $\hat{f} ([ a ]) \equiv f(a)$. Since we can not extract a element of $A$, it has to be defined in a more complicated way.


Assume we have an element of set $\qset{A}$ which can be written as

$$(P, p, q) : \qset{A}$$

where $P$ is the predicate, $p$ is the truncated witness of $P$ is inhabited, and $q$ contains the proofs of the logical equivalence.

We can define a function $f_{P} : (\Sigma (x : A) ~P(x)) \to B$ simply by 

$$f_{P} \defeq f \circ \pi_{1}$$

but our witness $p : \| \Sigma (x : A) ~P(x) \|$ is truncated which can not be applied to $f_{P}$. However we can generate a function $\bar{f_{P}} : \| (\Sigma (x : A) ~P(x)) \| \to B$ with the following lemma \cite{krausEscardoEtAll_existence}.

\begin{lemma}\label{constantT}
Given a constant function $g : X \to Y$ where $Y$ is a set, hence it satisfies

$$\forall(x,y : X) \to g(x) = g(y)$$

there exists a function $\bar{g} : \|X\| \to Y$ such that $\bar{g}(| x |) \equiv g(x)$.
\end{lemma}

\begin{proof}

Define the subset

$$Y' \defeq \Sigma (y : Y)~ \|\Sigma (x : X) ~g(x) = y \|$$

Intuitively it contains only the image of the constant function i.e.\ \textbf{$Y'$
is propositional}:

For any $(y_1,p_1) : Y'$ and $(y_2,p_2): Y'$, 

we can first generate the proofs 

$p_1 ((g(x) = y_1),\pi_2) : g(x) = y_1$ and 

$p_2 ((g(x) = y_2),\pi_2) : g(x) = y_2$

By symmetry and transitivity we can prove that $y_1 = y_2$.

From the fact that truncated type is always propositional, we can also deduce that $p1 = p2$, then $(y_1,p_1) = (y_2,p_2)$. Hence we can conclude that $Y'$ is propositional.

We can simply define a function $g' : X \to Y'$ from $g$ as

$$g'(x) \defeq (g(x), \lambda Q ~f \to f (x , \text{refl}_{=}(g(x))))$$


Because $Y'$ is propositional, it is possible to lift $g'$ to a function $\hat{g'} : \| X\| \to Y'$ which is technically defined as

$$\hat{g'}(x) \defeq x (Y',g')$$


Finally we define

$\bar{g} \defeq \pi_1 \circ \hat{g'}$

Which fulfils the computation rule

$\bar{g}(|x|) \equiv  \pi_1 (|x|(Y',g'))  \equiv \pi_1 (g' (x)) \equiv g(x)$
\end{proof}

\begin{proposition}
$f_P$ is a constant function.
\end{proposition}

\begin{proof}
Given two elements of $\Sigma (x : A) ~P(x)$, $(x_1,p_1)$ and $(x_2,p_2)$.

From the properties of equivalence class $P$, given $p_1 : P(x_1)$ and $p_2 : P(x_2)$,
we know that 

$$x_1 \sim x_2$$

Because $f$ is compatible with $\sim$,

$$f(x_1) = f(x_2)$$


By definition of $f_P$,

$$f_P(x_1,p_1) \equiv f (x_1) = f(x_2) \equiv f_P(x_2,p_2)$$

hence $f_P$ is a constant function.
\end{proof} 

Now apply \autoref{constantT} to constant function $f_{P}$, we obtain the expected function $\bar{f_{P}} : \| (\Sigma (x : A) ~P(x)) \| \to B$ which allows us to generate $\bar{f_{P}}(p) : B$.

To conclude, we can lift a function $f : A \to B$ which is compatible with $\sim$. The computational rule can be verified easily as follows:

$$\hat{f}([a]) \equiv \bar{f}_{\lambda x \to a \sim x}(|a|) \equiv f(a)$$



\subsection{Quotient inductive types}


%As we mentioned before, quotient types (in the sense of \ref{iqs}) are available in \hott. Because of the different interpretations of types, it makes less confusion to call them \emph{quotients} here.

An alternative way to define quotients in \hott is by using higher inductive types.


Assume $A$ is a set and $\_\sim\_ : A \to A \to \Prop$ is an equivalence relation. 
To build a quotient, we can simply impose level-$1$ morphisms in the structure of given sets according to the equivalence relation.
Thus, a quotient $\qset{A}$ can be defined as a higher inductive type:

\begin{itemize}
\item $[\_] : A \rightarrow \qset{A}$
\item $eqv : (a,b : A) \rightarrow a \sim b \rightarrow  [a] = [b]$
\item $isSet : (x,y:\qset{A}) \rightarrow (p_1,p_2 : x = y) \rightarrow p_1 = p_2$
\end{itemize}

It is also a set so we call it \textbf{set-quotient} or \textbf{quotient inductive types} (QITs).

\begin{remark}
Categorically, there is no difference between a set and a setoid because a skeleton of a setoid
is just a discrete category which is equivalent to a set.
\end{remark}

Some examples suggests that QITs are more powerful than quotient types.

One of the example is the definition of real numbers $\R$ which will be discussed in \autoref{rl}. Our construction of reals by Cauchy sequences of rational numbers is not Cauchy complete because there is no limit of each equivalence class. However, the Cauchy approximation approach in \cite{hott} (see subsection 11.3.1) using quotient inductive types is Cauchy complete due to the fact that the equivalence relation and limits are included in its definition.

Another example is unordered trees (rooted tree) which are trees connected to a multiset of subtrees, hence there is no ordering on subtrees.

Firstly we define ordered trees as:

\begin{itemize}
\item A leaf $l: \mathsf{Tree}$, or
\item An ordered list of subtrees indexed by $I$, $st : (I \rightarrow \mathsf{Tree}) \rightarrow \mathsf{Tree}$,
\end{itemize}

With the following equivalence relation:

\begin{itemize}
\item $l_{eq} : l \sim l$,
\item $st_{eq} : (f , g : I \to \mathsf{Tree}) \to f \sim_{p} g \to sp(f) \sim sp(g)$
\end{itemize}

where $f \sim_{p} g$ stands for $f$ is a permutation of $g$. The permutation can be defined using a bijective map $p : I \to I$ which relates equivalent subtrees recursively.

%However, if we use quotient type $\mathsf{Tree^{\sim}} := \qset{\mathsf{Tree}}$, the resulting trees have unordered subtrees which are themselves ordered.

If we define unordered trees as a quotient type $\mathsf{Tree^{\sim}} := \qset{\mathsf{Tree}}$, 
it is problematic to lift $st$.
For trees with finite subtrees such as \emph{binary trees} where $I \defeq \bool$, $st$ can be lifted by nesting lifting functions,

$$\hat{st}(a,b) = \dlift{\hat{st}(a)}(b)$$

because its type is isomorphic to $\mathsf{BTree} \rightarrow \mathsf{BTree} \rightarrow \mathsf{BTree}$.
Intuitively this approach can be applied to trees with finite subtrees. However it fails if have infinite subtrees, for example when $I \defeq \N$.

However we can avoid the problem by using QITs:

\begin{itemize}
\item $l: \mathsf{Tree}$, 
\item $st : (I \rightarrow \mathsf{Tree}) \rightarrow \mathsf{Tree}$

and a set of paths relates two permuted trees:

\item $l_{eq} : l  =_{\mathsf{Tree}} l $

\item$st_{eq} : \forall (f, g : I \to \mathsf{Tree}) \rightarrow  f \sim_{p} g \rightarrow  st(f) =_{\mathsf{Tree}} st(g)$
\end{itemize}

Similarly the cumulative hierarchy of all sets introduced in \cite{hott} (see section 10.5) also suggests that quotient types have some weaknesses compared to quotient inductive types.

A cumulative hierarchy can be given by constructors,

$$\{\_\} : (I : \Set) \to (I \to M_0) \to M_0$$

along with a subset relation,

$$\_\in\_ : M_0 \to M_0 \to \Prop$$

which is inhabited if $ f(i) \in \{ I , f \}$.

Then we can easily define the equivalence relation on "sets" using set-theoretical definition

$A \sim B \defeq \forall m : M_0, m \in A \iff m \in B$

Similar to unordered trees, we can not obtain the constructor $\dlift{\{\_\}}$ because the index set $I$ can be infinite.

To summarise, it seems that quotient inductive types are more powerful than quotient types due to the ability of defining term constructors and equivalence relations simultaneously. 
However, quotient inductive types are not available in type theories other than \hott and the computational interpretation of them is still an open problem.
Moreover, there can be more general quotients in \hott, for example a quotient of a type by a $1$-groupoid (See section 9.9 in \cite{hott}). It is interesting to investigate \emph{real} \emph{quotient types} in \hott, but it is beyond this thesis.

\section{Related work}

The introduction of quotient types in Type Theory has been studied by several authors in different versions of \mltt and using various approaches.

\begin{itemize}

\item In \cite{cab}, Mendler et al. have first considered building new types from a
given type using a quotient operator $//$. Their work is done in an
implementation of \ett, NuPRL. 
In NuPRL, every type comes with its own equality relation, so the quotient operator can be
seen as a way of redefining equality in a type. But it is not all
about building new types. They also discuss problems that arise from
defining functions on the new type which can be illustrated using a simple example. 

Assume the base type is $A$ and the new equivalence relation is $E$, the new
type can be formed as $A//E$. 

When we want to define a function $f \,\colon\, A//E \to \2$,  $f(a) \not= f(b$ may
exists for $a, b \,\colon A$ such that $E\,a\,b$. This will lead to
inconsistency since $E(a,b)$ implies $a$ converts to $b$ in \ett, hence
the left hand side $f(a)$ can be converted to $f(b)$, namely we get $f(b) \not= f(b)$
which is contradicted with the equality reflection rule. 

Therefore a function is said to be well-defined \cite{cab} on the new type only
if it respects the equivalence relation $E$, namely

$$\forall a,b : A \to E(a,b) \to f(a) = f(b)$$


 After the introduction of quotient types, Mendler further investigates
 this topic from a categorical perspective in ~\cite{men:90}. He uses
 the correspondence between quotient types in \mltt{} and coequalizers
 in a category of types to define a notion called \emph{squash types},
 which is further discussed by Nogin \cite{nog:02}.

\item Nogin \cite{nog:02} considers a modular approach to axiomatizing quotient types in NuPRL as well. Despite the ease of constructing new types
from base types, he also discusses some
problems about quotient types. For example, since the equality is
extensional, we cannot recover the
witness of the equality.  He suggests including more axioms to
conceptualise quotients. He decomposes the formalisation of quotient type
into several smaller primitives such that they can be handled much
simpler.


\item Jacobs \cite{Jacobs94quotientsin} introduces a syntax for quotient types based on a predicate logic
within simple type theory. He discusses quotient types from a categorical theoretical perspective. In fact the syntax of quotient types arises from an adjunction as we mentioned before.

\item To add quotient types to \mltt, Hofmann proposes three models for
quotient types in \cite{hof:phd}. The first one is a setoid model for
quotient types. In this model all types are attached with partial
equivalence relations, namely all types are partial setoids rather than
sets. It does not provide dependency at the level of types but only at the level of the relations. 
%This is similar to \ett in some sense but retains computational property in \itt. 
The second one is groupoid model which supports most features required 
but it is not definable in \itt. 
He also proposes a third model to as an attempt to overcome problems in the previous two models. More type dependency is provided and quotient types are believed to be definable in this model.
However it also has some disadvantages. 
he also shows that \ett is conservative over \itt extended with quotient types \cite{hof:95:con}.


\item Altenkirch \cite{alti:lics99} also provides a different setoid model which is built in an \itt
extended with proof-irrelevant universes of propositions and $\eta$-rules for $\Pi$-types and $\Sigma$-types. It is decidable, adequate and permits large eliminations. We have implemented this setoid model and interpret quotient types in it (see \autoref{models}).

\item Homeier \cite{hom} also axiomatises quotient types in Higher Order Logic
(HOL), which is also a theorem prover. He creates a tool package to
construct quotient types as a conservative extension of HOL such that
users are able to define new types in HOL. Next he defines the
normalisation functions and proves several properties of
these. Finally he discussed the issues when quotienting on the
aggregate types such as lists and pairs.


\item Courtieu \cite{cou:01} shows an extension of Calculus of Inductive Constructions
with \emph{Normalised Types} which are similar to quotient types, but equivalence relations are replaced by normalisation functions which select a canonical element for each equivalence class. 
In fact normalised types can be seen as a proper subset of quotient types. We can easily recover a quotient type from a normalised type as below

$$a \sim b \defeq [ a ] = [ b ]$$

However not all quotient types have normal forms, for example, the set of real numbers (see \autoref{rl}).
The notion \emph{definable quotients} we proposed in \autoref{dq} is also similar to it, but does not provide a new type automatically.

\item Barthe and Geuvers \cite{bar:96} also propose a new notion called
\emph{congruence types}, which is also a special class of quotient
types, in which the base type are inductively defined and with a set
of reduction rules called the term-rewriting system. The idea behind
it is the $\beta$-equivalence is replaced by a set of
$\beta$-conversion rules. Congruence types can be treated as an
alternative to the pattern matching introduced in \cite{coq:92}. The main
purpose of introducing congruence types is to solve problems in
term rewriting systems rather than to implement quotient types.
Congruence types are not inductive but have good computational behaviour because we can use term-rewriting system to link a term of base type to a unique term of congruence type as its normal form.
However this approach has some problems in termination criteria and interaction between rewriting systems \cite{cou:01}.

\item Barthe and Capretta \cite{bar:03} compare different ways to setoids in Type Theory.
The setoid is classified as partial setoid or total setoid depending
on whether the equality relation is reflexive or not. They also
consider obtain quotients with different kinds of setoids, especially
the ones from partial setoids are difficult to define because the lack
of reflexivity.

\item Abbott, Altenkirch et al. \cite{abb:04} provides the basis for
programming with quotient datatypes polymorphically based on their
works on containers which are datatypes whose instances are
collections of objects, such as arrays, trees and so on. Generalising
the notion of container, they define quotient containers as the
containers quotiented by a collection of isomorphisms on the positions
within the containers.

\item Voevodsky \cite{voe:hset} implements quotients in Coq based on a set
of axioms of \hott. He firstly implement
equivalence class and use it to implement quotients which is an
analogy to the construction of quotient sets in set theory. The detail is given in \autoref{impredicative}.
\end{itemize}





\section{Summary}

We have given the syntax of quotient types in this chapter. We have provided a non-dependent eliminator with an induction principle and a dependent eliminator and proved they are equivalent. We have also shown that propositional extensionality actually implies the effectiveness of quotients. The characterisation of quotients in category theory has been given. It does not only corresponds to coequalizers but also can be generated from a left adjoint functor to the equality functor $\nabla : \textbf{Set} \to \textbf{Setoid}$. We have also briefly reviewed the research of quotient types by other authors.








\chapter{Definable Quotients}
\label{dq}

\todo{1. Distributivity 2. Rational numbers 3. other definable quotients}

Quotient types may be necessary for defining some types in \itt, but
it is not always the case. The set of integers and the set of rational
numbers can be defined without quotients types, but there are more
properties revealed if we view them as quotients. 
However there are some good
properties if we relate them with the base types and equivalence
relation, for example we can lift functions and their properties from base types to
quotient types. Moreover, if the base types are simpler to manipulate,
it is worthwhile using the base type to define functions and reasoning
and then lifting them. We can achieve more convenience by manipulating base types and then lifting the operators and propositions according to the relation between quotient types and base types. The main things we
are going to discuss in this chapter is called  definable quotient
structures which does not require quotient types
to be added into the \itt.

In this Chapter we will show this using one of the examples, the set
of integers. Some of the work is is conducted by Thorsten Altenkirch,
Thomas Anberr\'{e}e and the author together, and summarised in \cite{aan} .


\section{Integers}

From the usual symbols to represent integers, we can easily figure out
one inductive definition for integers,

\begin{code}
\\
\>\AgdaKeyword{data} \AgdaDatatype{ℤ} \AgdaSymbol{:} \AgdaPrimitiveType{Set} \AgdaKeyword{where}\<%
\\
\>[0]\AgdaIndent{2}{}\<[2]%
\>[2]\AgdaInductiveConstructor{+\_} \AgdaSymbol{:} \AgdaDatatype{ℕ} \AgdaSymbol{→} \AgdaDatatype{ℤ}\<%
\\
\>[0]\AgdaIndent{2}{}\<[2]%
\>[2]\AgdaInductiveConstructor{zero} \AgdaSymbol{:} \AgdaDatatype{ℤ}\<%
\\
\>[0]\AgdaIndent{2}{}\<[2]%
\>[2]\AgdaInductiveConstructor{-\_} \AgdaSymbol{:} \AgdaDatatype{ℕ} \AgdaSymbol{→} \AgdaDatatype{ℤ}\<%
\\
\end{code}

However we face a trade-off: three different representation for zero or
to use code $+0$ for number $+1$. Usually the principle is to not
losing canonicity because it requires unnecessary checking for whether
some functions respect the equivalence or not. Therefore, the second
choice makes more sense and we refine it a bit as:

\begin{code}
\\
\>\AgdaKeyword{data} \AgdaDatatype{ℤ} \AgdaSymbol{:} \AgdaPrimitiveType{Set} \AgdaKeyword{where}\<%
\\
\>[0]\AgdaIndent{2}{}\<[2]%
\>[2]\AgdaInductiveConstructor{+suc\_} \AgdaSymbol{:} \AgdaDatatype{ℕ} \AgdaSymbol{→} \AgdaDatatype{ℤ}\<%
\\
\>[0]\AgdaIndent{2}{}\<[2]%
\>[2]\AgdaInductiveConstructor{zero} \<[8]%
\>[8]\AgdaSymbol{:} \AgdaDatatype{ℤ}\<%
\\
\>[0]\AgdaIndent{2}{}\<[2]%
\>[2]\AgdaInductiveConstructor{-suc\_} \AgdaSymbol{:} \AgdaDatatype{ℕ} \AgdaSymbol{→} \AgdaDatatype{ℤ}\<%
\\
\>\<\end{code}

This is better, but in practice it is expected to have more cases if
we use pattern matching. Every time we use pattern matching, a case will be
split into three. This becomes worse and worse when we have mutiple integer
arguments and we have to do case analysis on all of them. A simple
refinement is combining the first two constructors:

\begin{code}
\\
\>\AgdaKeyword{data} \AgdaDatatype{ℤ} \AgdaSymbol{:} \AgdaPrimitiveType{Set} \AgdaKeyword{where}\<%
\\
\>[0]\AgdaIndent{2}{}\<[2]%
\>[2]\AgdaInductiveConstructor{+\_} \<[8]%
\>[8]\AgdaSymbol{:} \AgdaDatatype{ℕ} \AgdaSymbol{→} \AgdaDatatype{ℤ}\<%
\\
\>[0]\AgdaIndent{2}{}\<[2]%
\>[2]\AgdaInductiveConstructor{-suc\_} \AgdaSymbol{:} \AgdaDatatype{ℕ} \AgdaSymbol{→} \AgdaDatatype{ℤ}\<%
\\
\>\<\end{code}

This is the most proper version we decided to use for the set of integers. It is
inductively defined and is readable because it is just an
intepretation of the usual symbols for integers in regular mathematics.

Usually we believe that the reason of inventing integers is the lack of symbols to
represent the results of subtraction between two
natural numbers. Integers are used to represent these results, and
vice versa,
every integer can be represented as a pair of natural numbers and the
choice is not unique.
For example, from the equation $1 - 4 = - 3$, it is clear that the
integer $- 3$ can be represented the pair $(1,4)$. 
Therefore we can use the paired natural numbers as an alternative
definition for integers.

$$\Z_0=\N \times \N$$

However since there are different pairs for one integer, we have to
quotient it with an equivalence relation. For any two pairs of natural
numbers $(n_1, n_2)$ and $(n_3, n_4)$, we know they represent the same
integer if

$$ n_1 - n_2 = n_3 - n_4$$

Technically, this does not work because the subtraction defined for natural
numbers only returns zero if the pair is for negative number. We only
need to do some small modifications:

$$ n_1 + n_4 = n_3 + n_2$$

This helps us define a relation but it is not enough. This is an
equation in mathamtics, but in Type Theory we have to prove that it is an
equivalence relation, namely, it is reflexive, symmetric and transitive.

Combining the carrier (the pair of natural numbers), the equivalence
relation and its proof, we have a setoid.

\begin{code}
\\
\>\AgdaFunction{ℤ-Setoid} \AgdaSymbol{:} \AgdaRecord{Setoid}\<%
\\
\>\AgdaFunction{ℤ-Setoid} \AgdaSymbol{=} \AgdaKeyword{record}\<%
\\
\>[2]\AgdaIndent{3}{}\<[3]%
\>[3]\AgdaSymbol{\{} \AgdaField{Carrier} \<[19]%
\>[19]\AgdaSymbol{=} \AgdaFunction{ℤ₀}\<%
\\
\>[2]\AgdaIndent{3}{}\<[3]%
\>[3]\AgdaSymbol{;} \AgdaField{\_≈\_} \<[19]%
\>[19]\AgdaSymbol{=} \AgdaPostulate{\_∼\_}\<%
\\
\>[2]\AgdaIndent{3}{}\<[3]%
\>[3]\AgdaSymbol{;} \AgdaField{isEquivalence} \AgdaSymbol{=} \AgdaPostulate{\_∼\_isEquivalence}\<%
\\
\>[2]\AgdaIndent{3}{}\<[3]%
\>[3]\AgdaSymbol{\}}\<%
\\
\end{code}

Since the set of integers is definable as we discussed before, they can
be seen as the normal forms of the equivalent classes. The normalisation
function can be defined as follows:
 
\begin{code}
\\
\>\AgdaFunction{[\_]} \<[18]%
\>[18]\AgdaSymbol{:} \AgdaFunction{ℤ₀} \AgdaSymbol{→} \AgdaDatatype{ℤ}\<%
\\
\>\AgdaFunction{[} \AgdaBound{m} \AgdaInductiveConstructor{,} \AgdaInductiveConstructor{0} \AgdaFunction{]} \<[18]%
\>[18]\AgdaSymbol{=} \AgdaInductiveConstructor{+} \AgdaBound{m}\<%
\\
\>\AgdaFunction{[} \AgdaInductiveConstructor{0} \AgdaInductiveConstructor{,} \AgdaInductiveConstructor{suc} \AgdaBound{n} \AgdaFunction{]} \<[18]%
\>[18]\AgdaSymbol{=} \AgdaInductiveConstructor{-suc} \AgdaBound{n}\<%
\\
\>\AgdaFunction{[} \AgdaInductiveConstructor{suc} \AgdaBound{m} \AgdaInductiveConstructor{,} \AgdaInductiveConstructor{suc} \AgdaBound{n} \AgdaFunction{]} \AgdaSymbol{=} \AgdaFunction{[} \AgdaBound{m} \AgdaInductiveConstructor{,} \AgdaBound{n} \AgdaFunction{]}\<%
\\
\end{code}

The function should be proved well-defined on the setoid, namely it
has to respect the equivalence relation. We call it soundness here. 
It is not trivial but easy to observe that the function is
sound\footnote{the formal proof can be found in appendix (we cheat a
  bit by defining embedding function to make it simpler)}.

A setoid and a function respects this equivalence (not necessary to be
a normalisation function) constitute a prequotient.

\begin{definition}
Prequotient.

\noindent
Given a setoid $(A,\sim)$,  a \emph{prequotient} $(Q,[\_],\sound)$ over that setoid consists in
\begin{enumerate}
\item \label{enum:Q} a set $Q$,
\item \label{enum:box}a function $[\_]: A \to Q$,
\item \label{enum:sound} a proof \emph{sound} that  the function $[\_]$ is compatible with the relation $\sim$,
that is \[\sound\colon (a,b : A) \to a\sim b \to [a] = [b],\]
\end{enumerate}


Prequotient only includes the formalisation rules and introduction
rules. To complete a \emph{quotient}, we also need the elimination rule added into such a prequotient

\begin{enumerate}
\setcounter{enumi}{3}
\item \label{enum:elim}
for any $B: Q\to\Set$, an eliminator
 \begin{align*}
 \qelim_B\,\,:\,\,\,&(f\colon (a:A) \to B\,\class a) \\
        {\to}\, &((p:a\sim b) \to f\,a \simeq_{\sound\,p}f\,b)\\
        {\to}\, &((q:Q) \to B\,q)
 \end{align*}
such that $\qelimbeta : \qelim_B\,f \,p\,\class a\equiv f\,a$.
\end{enumerate}

This eliminator is also called dependent lifting function because it
actually lifts a function which is well-defined on the setoid to a
function defined on the quotient type. The result type is also
dependent on the quotient type. There is an equivalent definition
given by Martin Hofmann
which has a non-dependent eliminator with an induction principle
instead.

\[\lift\colon (f\colon A \to B) \to (\forall a,b\cdot a\sim b \to f\,a
\equiv f\,b) \to (Q \to B)\]

Suppose $B$ is a predicate, 
\[\qind \colon((a: A)\to B \,\class a)\to ((q : Q)\to B\,q)\]

However, it is oberservable that given any non-empty set $Q$, all
constant functions fit in this definition. Any element of $Q$ has to
be mapped from at most one equivalence class of the setoid
$(A,\sim)$. This property is called \emph{exact} here 

\begin{enumerate}
\setcounter{enumi}{4}
\item $exact :(\forall a,b : A) \to  \class a \equiv \class b \to a \sim b$.

\end{enumerate}
\end{definition}

The quotient is exact if exactly one equivalence class corresponds to
an element of $Q$.

We already know that the integer is definable and it is plausible to
find a representative in each equivalence classes. Since
we treat elements of $Q$ as the name for the equivalence classes, the
selection function can be defined as an embedding function from  the
quotient type $Q$ to base type $A$. This is an alternative and more
flexible way to eliminate the quotient type $Q$ and if a
prequotient $(Q, \class{\dotph}, \sound)$ on a setoid $(A,\sim)$ has
an embedding function which is specified as
\begin{align*}
\emb &: Q \to A\\
\complete &: (a : A) \to \emb {\class a} \sim a\\
\stable &: (q:Q) \to \class{\emb\,q} \equiv q.\\
\end{align*}

Then it is a \emph{definable quotient}. Composing the ``normalisation'' function $\class\dotph$ with the
embedding function, we obtain the real normalisation function. A
definable quotient is an \emph{exact} quotient which is proved in \cite{aan}.



\paragraph{Operations}

For a definable quotient, we can lift an operation by mixing the
normalisation and embedding functions. For example, given an unary
operator 

\begin{code}
\\
\>\AgdaFunction{lift₁} \AgdaSymbol{:} \AgdaSymbol{(}\AgdaBound{op} \AgdaSymbol{:} \AgdaFunction{ℤ₀} \AgdaSymbol{→} \AgdaFunction{ℤ₀}\AgdaSymbol{)} \AgdaSymbol{→} \AgdaDatatype{ℤ} \AgdaSymbol{→} \AgdaDatatype{ℤ}\<%
\\
\>\AgdaFunction{lift₁} \AgdaBound{op} \AgdaSymbol{=} \AgdaFunction{[\_]} \AgdaFunction{∘} \AgdaBound{op} \AgdaFunction{∘} \AgdaFunction{⌜\_⌝}\<%
\\
\end{code}

To lift binary or n-ary operators, we only need to apply the operator
to the representative for each "equivalence class" and "normalise" the
result so that it becomes a function defined on the set of
"equivalence classes".

But there is a unavoidable problem: not all operations defined on
$\Z_0$ is defind on the setoid, namely respect the equivalence
relation. Therefore it is reasonable to verify if the function is
well-defined on the setoid:

$$a \sim b → op \, a \sim op \, b$$

We will show how to define the addition for the quotient
integers. Given two numbers $(a_1 , b_1)$ and $(a_2 , b_2)$ We
only need to add them pair-wisely together, and it can be verified
easily because we know that 

$(a_1 - b_1) + (a_2 - b_2) = (a_1 + a_2) - (b_1) - (b_2)$

The verification is not necessary here but should be important in
other cases.

\begin{code}
\\
\>\AgdaFunction{\_+\_} \AgdaSymbol{:} \AgdaFunction{ℤ₀} \AgdaSymbol{→} \AgdaFunction{ℤ₀} \AgdaSymbol{→} \AgdaFunction{ℤ₀}\<%
\\
\>\AgdaSymbol{(}\AgdaBound{x+} \AgdaInductiveConstructor{,} \AgdaBound{x-}\AgdaSymbol{)} \AgdaFunction{+} \AgdaSymbol{(}\AgdaBound{y+} \AgdaInductiveConstructor{,} \AgdaBound{y-}\AgdaSymbol{)} \AgdaSymbol{=} \AgdaSymbol{(}\AgdaBound{x+} \AgdaFunction{ℕ+} \AgdaBound{y+}\AgdaSymbol{)} \AgdaInductiveConstructor{,} \AgdaSymbol{(}\AgdaBound{x-} \AgdaFunction{ℕ+} \AgdaBound{y-}\AgdaSymbol{)}\<%
\\
\end{code}


\paragraph{Properties}

We can also define the ring of $\Z$. It contains a lot of properties
to prove which are seemed as axioms in classic mathematics. In
constructive mathematics, the only axioms for integers are the
constructors and the elimination rules.

As what we have menetioned, even though the definition of integers
only has two constructors, it gradually increase the difficulty of
proving when doing case analyses on more and more integers. One
example is the proving of distributivity.

\paragraph{An efficient utilization of quotient structure: the proving of distributivity}

One of the most important motivations of using setoid integers is that
the setoid definition reduces the complexity of programs involving
integers. We have shown it is simpler to define some operators, but it
will be much more evident when proving properties of the ring of integers.

Most simple laws of the ring of integers are not as unbearably
complicated as the distributivity laws. An attempt of the right
distributivity for $\Z$ ($ (y + z) \times x = y \times x + z \times x$)
is shown as below. The multiplication is not defined with pattern
matching, but in the arithmetic approach.

\begin{code}
\\
\>\AgdaFunction{\_ℤ*\_} \AgdaSymbol{:} \AgdaDatatype{ℤ} \AgdaSymbol{→} \AgdaDatatype{ℤ} \AgdaSymbol{→} \AgdaDatatype{ℤ}\<%
\\
\>\AgdaBound{i} \AgdaFunction{ℤ*} \AgdaBound{j} \AgdaSymbol{=} \AgdaFunction{sign} \AgdaBound{i} \AgdaFunction{S*} \AgdaFunction{sign} \AgdaBound{j} \AgdaFunction{◃} \AgdaFunction{∣} \AgdaBound{i} \AgdaFunction{∣} \AgdaPrimitive{ℕ*} \AgdaFunction{∣} \AgdaBound{j} \AgdaFunction{∣}\<%
\\
%
\\
\end{code}

The first idea which I came up
with is to use the right distributivity law of natural
numbers. It is observable that if all three variables have the same
signs, it is easy to apply the right distributivity law of natural
numbers. We can relax the constraint a bit, only y and z have the
same symbol, it is still plausible. To prove this part we need three
lemmas,

\begin{code}
\\
\>\AgdaFunction{lem1} \AgdaSymbol{:} \AgdaSymbol{∀} \AgdaBound{x} \AgdaBound{y} \AgdaSymbol{→} \AgdaFunction{sign} \AgdaBound{x} \AgdaDatatype{≡} \AgdaFunction{sign} \AgdaBound{y} \AgdaSymbol{→} \AgdaFunction{∣} \AgdaBound{x} \AgdaFunction{ℤ+} \AgdaBound{y} \AgdaFunction{∣} \AgdaDatatype{≡} \AgdaFunction{∣} \AgdaBound{x} \AgdaFunction{∣} \AgdaPrimitive{ℕ+} \AgdaFunction{∣} \AgdaBound{y} \AgdaFunction{∣} \<[61]%
\>[61]\<%
\\
\>\AgdaFunction{lem1} \AgdaSymbol{(}\AgdaInductiveConstructor{-suc} \AgdaBound{x}\AgdaSymbol{)} \AgdaSymbol{(}\AgdaInductiveConstructor{-suc} \AgdaBound{y}\AgdaSymbol{)} \AgdaBound{e} \AgdaSymbol{=} \AgdaFunction{cong} \AgdaInductiveConstructor{suc} \AgdaSymbol{(}\AgdaFunction{sym} \AgdaSymbol{(}\AgdaFunction{m+1+n≡1+m+n} \AgdaBound{x} \AgdaBound{y}\AgdaSymbol{))}\<%
\\
\>\AgdaFunction{lem1} \AgdaSymbol{(}\AgdaInductiveConstructor{-suc} \AgdaBound{x}\AgdaSymbol{)} \AgdaSymbol{(}\AgdaInductiveConstructor{+} \AgdaBound{y}\AgdaSymbol{)} \AgdaSymbol{()}\<%
\\
\>\AgdaFunction{lem1} \AgdaSymbol{(}\AgdaInductiveConstructor{+} \AgdaBound{x}\AgdaSymbol{)} \AgdaSymbol{(}\AgdaInductiveConstructor{-suc} \AgdaBound{y}\AgdaSymbol{)} \AgdaSymbol{()}\<%
\\
\>\AgdaFunction{lem1} \AgdaSymbol{(}\AgdaInductiveConstructor{+} \AgdaBound{x}\AgdaSymbol{)} \AgdaSymbol{(}\AgdaInductiveConstructor{+} \AgdaBound{y}\AgdaSymbol{)} \AgdaBound{e} \AgdaSymbol{=} \AgdaInductiveConstructor{refl}\<%
\\
%
\\
\>\AgdaFunction{lem2} \AgdaSymbol{:} \AgdaSymbol{∀} \AgdaBound{x} \AgdaBound{y} \AgdaSymbol{→} \AgdaFunction{sign} \AgdaBound{x} \AgdaDatatype{≡} \AgdaFunction{sign} \AgdaBound{y} \AgdaSymbol{→} \AgdaFunction{sign} \AgdaSymbol{(}\AgdaBound{x} \AgdaFunction{ℤ+} \AgdaBound{y}\AgdaSymbol{)} \AgdaDatatype{≡} \AgdaFunction{sign} \AgdaBound{y}\<%
\\
\>\AgdaFunction{lem2} \AgdaSymbol{(}\AgdaInductiveConstructor{-suc} \AgdaBound{x}\AgdaSymbol{)} \AgdaSymbol{(} \AgdaInductiveConstructor{-suc} \AgdaBound{y}\AgdaSymbol{)} \AgdaBound{e} \AgdaSymbol{=} \AgdaInductiveConstructor{refl}\<%
\\
\>\AgdaFunction{lem2} \AgdaSymbol{(}\AgdaInductiveConstructor{-suc} \AgdaBound{x}\AgdaSymbol{)} \AgdaSymbol{(}\AgdaInductiveConstructor{+} \AgdaBound{y}\AgdaSymbol{)} \AgdaSymbol{()}\<%
\\
\>\AgdaFunction{lem2} \AgdaSymbol{(}\AgdaInductiveConstructor{+} \AgdaBound{x}\AgdaSymbol{)} \AgdaSymbol{(}\AgdaInductiveConstructor{-suc} \AgdaBound{y}\AgdaSymbol{)} \AgdaSymbol{()}\<%
\\
\>\AgdaFunction{lem2} \AgdaSymbol{(}\AgdaInductiveConstructor{+} \AgdaBound{x}\AgdaSymbol{)} \AgdaSymbol{(}\AgdaInductiveConstructor{+} \AgdaBound{y}\AgdaSymbol{)} \AgdaBound{e} \AgdaSymbol{=} \AgdaInductiveConstructor{refl}\<%
\\
%
\\
%
\\
\>\AgdaFunction{lem3} \AgdaSymbol{:} \AgdaSymbol{∀} \AgdaBound{x} \AgdaBound{y} \AgdaBound{s} \AgdaSymbol{→} \AgdaBound{s} \AgdaFunction{◃} \AgdaSymbol{(}\AgdaBound{x} \AgdaPrimitive{ℕ+} \AgdaBound{y}\AgdaSymbol{)} \AgdaDatatype{≡} \AgdaSymbol{(}\AgdaBound{s} \AgdaFunction{◃} \AgdaBound{x}\AgdaSymbol{)} \AgdaFunction{ℤ+} \AgdaSymbol{(}\AgdaBound{s} \AgdaFunction{◃} \AgdaBound{y}\AgdaSymbol{)}\<%
\\
\>\AgdaFunction{lem3} \AgdaInductiveConstructor{0} \AgdaInductiveConstructor{0} \AgdaBound{s} \AgdaSymbol{=} \AgdaInductiveConstructor{refl}\<%
\\
\>\AgdaFunction{lem3} \AgdaInductiveConstructor{0} \AgdaSymbol{(}\AgdaInductiveConstructor{suc} \AgdaBound{y}\AgdaSymbol{)} \AgdaBound{s} \AgdaSymbol{=} \AgdaFunction{sym} \AgdaSymbol{(}\AgdaFunction{ℤ-id-l} \AgdaSymbol{\_)}\<%
\\
\>\AgdaFunction{lem3} \AgdaSymbol{(}\AgdaInductiveConstructor{suc} \AgdaBound{x}\AgdaSymbol{)} \AgdaBound{y} \AgdaBound{s} \AgdaSymbol{=} \AgdaFunction{trans} \AgdaSymbol{(}\AgdaFunction{h} \AgdaBound{s} \AgdaSymbol{(}\AgdaBound{x} \AgdaPrimitive{ℕ+} \AgdaBound{y}\AgdaSymbol{))} \AgdaSymbol{(}\<%
\\
\>[4]\AgdaIndent{19}{}\<[19]%
\>[19]\AgdaFunction{trans} \AgdaSymbol{(}\AgdaFunction{cong} \AgdaSymbol{(λ} \AgdaBound{n} \AgdaSymbol{→} \AgdaSymbol{(}\AgdaBound{s} \AgdaFunction{◃} \AgdaInductiveConstructor{suc} \AgdaNumber{0}\AgdaSymbol{)} \AgdaFunction{ℤ+} \AgdaBound{n}\AgdaSymbol{)} \AgdaSymbol{(}\AgdaFunction{lem3} \AgdaBound{x} \AgdaBound{y} \AgdaBound{s}\AgdaSymbol{))} \AgdaSymbol{(}\<%
\\
\>[4]\AgdaIndent{19}{}\<[19]%
\>[19]\AgdaFunction{trans} \AgdaSymbol{(}\AgdaFunction{sym} \AgdaSymbol{(}\AgdaPostulate{ℤ-+-assoc} \AgdaSymbol{(}\AgdaBound{s} \AgdaFunction{◃} \AgdaInductiveConstructor{suc} \AgdaNumber{0}\AgdaSymbol{)} \AgdaSymbol{(}\AgdaBound{s} \AgdaFunction{◃} \AgdaBound{x}\AgdaSymbol{)} \AgdaSymbol{(}\AgdaBound{s} \AgdaFunction{◃} \AgdaBound{y}\AgdaSymbol{)))} \AgdaSymbol{(}\<%
\\
\>[4]\AgdaIndent{19}{}\<[19]%
\>[19]\AgdaFunction{cong} \AgdaSymbol{(λ} \AgdaBound{n} \AgdaSymbol{→} \AgdaBound{n} \AgdaFunction{ℤ+} \AgdaSymbol{(}\AgdaBound{s} \AgdaFunction{◃} \AgdaBound{y}\AgdaSymbol{))} \AgdaSymbol{(}\AgdaFunction{sym} \AgdaSymbol{(}\AgdaFunction{h} \AgdaBound{s} \AgdaBound{x}\AgdaSymbol{)))))}\<%
\\
\>[13]\AgdaIndent{2}{}\<[2]%
\>[2]\AgdaKeyword{where}\<%
\\
\>[0]\AgdaIndent{4}{}\<[4]%
\>[4]\AgdaFunction{h} \AgdaSymbol{:} \AgdaSymbol{∀} \AgdaBound{s} \AgdaBound{y} \AgdaSymbol{→} \AgdaBound{s} \AgdaFunction{◃} \AgdaInductiveConstructor{suc} \AgdaBound{y} \AgdaDatatype{≡} \AgdaSymbol{(}\AgdaBound{s} \AgdaFunction{◃} \AgdaSymbol{(}\AgdaInductiveConstructor{suc} \AgdaNumber{0}\AgdaSymbol{))} \AgdaFunction{ℤ+} \AgdaSymbol{(}\AgdaBound{s} \AgdaFunction{◃} \AgdaBound{y}\AgdaSymbol{)}\<%
\\
\>[0]\AgdaIndent{4}{}\<[4]%
\>[4]\AgdaFunction{h} \AgdaBound{s} \AgdaInductiveConstructor{0} \AgdaSymbol{=} \AgdaFunction{sym} \AgdaSymbol{(}\AgdaPostulate{ℤ-id-r} \AgdaSymbol{\_)}\<%
\\
\>[0]\AgdaIndent{4}{}\<[4]%
\>[4]\AgdaFunction{h} \AgdaInductiveConstructor{Sign.-} \AgdaSymbol{(}\AgdaInductiveConstructor{suc} \AgdaBound{y}\AgdaSymbol{)} \AgdaSymbol{=} \AgdaInductiveConstructor{refl}\<%
\\
\>[0]\AgdaIndent{4}{}\<[4]%
\>[4]\AgdaFunction{h} \AgdaInductiveConstructor{Sign.+} \AgdaSymbol{(}\AgdaInductiveConstructor{suc} \AgdaBound{y}\AgdaSymbol{)} \AgdaSymbol{=} \AgdaInductiveConstructor{refl}\<%
\\
\end{code}

The following is a partial definition with the first case that 

\begin{code}
\\
\>\AgdaFunction{distʳ} \<[7]%
\>[7]\AgdaSymbol{:} \AgdaFunction{\_ℤ*\_} \AgdaFunction{DistributesOverʳ} \AgdaFunction{\_ℤ+\_}\<%
\\
\>\AgdaFunction{distʳ} \AgdaBound{x} \AgdaBound{y} \AgdaBound{z} \AgdaKeyword{with} \AgdaFunction{sign} \AgdaBound{y} \AgdaFunction{S≟} \AgdaFunction{sign} \AgdaBound{z}\<%
\\
\>\AgdaFunction{distʳ} \AgdaBound{x} \AgdaBound{y} \AgdaBound{z} \<[15]%
\>[15]\AgdaSymbol{|} \AgdaInductiveConstructor{yes} \AgdaBound{p} \<[23]%
\>[23]\<%
\\
\>[4]\AgdaIndent{9}{}\<[9]%
\>[9]\AgdaKeyword{rewrite} \AgdaFunction{lem1} \AgdaBound{y} \AgdaBound{z} \AgdaBound{p} \AgdaSymbol{|} \AgdaBound{p} \<[32]%
\>[32]\<%
\\
\>[9]\AgdaIndent{15}{}\<[15]%
\>[15]\AgdaSymbol{|} \AgdaFunction{lem2} \AgdaBound{y} \AgdaBound{z} \AgdaBound{p} \AgdaSymbol{=} \<[30]%
\>[30]\<%
\\
\>[0]\AgdaIndent{6}{}\<[6]%
\>[6]\AgdaFunction{trans} \AgdaSymbol{(}\AgdaFunction{cong} \AgdaSymbol{(λ} \AgdaBound{n} \AgdaSymbol{→} \AgdaFunction{sign} \AgdaBound{z} \AgdaFunction{S*} \AgdaFunction{sign} \AgdaBound{x} \AgdaFunction{◃} \AgdaBound{n}\AgdaSymbol{)} \<[47]%
\>[47]\<%
\\
\>[0]\AgdaIndent{12}{}\<[12]%
\>[12]\AgdaSymbol{(}\AgdaFunction{distribʳ} \AgdaFunction{∣} \AgdaBound{x} \AgdaFunction{∣} \AgdaFunction{∣} \AgdaBound{y} \AgdaFunction{∣} \AgdaSymbol{(}\AgdaFunction{∣} \AgdaBound{z} \AgdaFunction{∣}\AgdaSymbol{)))} \<[44]%
\>[44]\<%
\\
\>[0]\AgdaIndent{6}{}\<[6]%
\>[6]\AgdaSymbol{(}\AgdaFunction{lem3} \AgdaSymbol{(}\AgdaFunction{∣} \AgdaBound{y} \AgdaFunction{∣} \AgdaPrimitive{ℕ*} \AgdaFunction{∣} \AgdaBound{x} \AgdaFunction{∣}\AgdaSymbol{)} \AgdaSymbol{(}\AgdaFunction{∣} \AgdaBound{z} \AgdaFunction{∣} \AgdaPrimitive{ℕ*} \AgdaFunction{∣} \AgdaBound{x} \AgdaFunction{∣}\AgdaSymbol{)} \<[46]%
\>[46]\<%
\\
\>[0]\AgdaIndent{12}{}\<[12]%
\>[12]\AgdaSymbol{(}\AgdaFunction{sign} \AgdaBound{z} \AgdaFunction{S*} \AgdaFunction{sign} \AgdaBound{x}\AgdaSymbol{))}\<%
\\
\>\AgdaFunction{distʳ} \AgdaBound{x} \AgdaBound{y} \AgdaBound{z} \AgdaSymbol{|} \AgdaInductiveConstructor{no} \AgdaBound{¬p} \AgdaSymbol{=} \AgdaBound{...}\<%
\\
\end{code}

If y and z have different
signs, it is impossible to apply the right distributivity law for
natural numbers. We have to prove it from scratch. Even though it
seems that there are only two cases to prove, it is conceivable that how
many lemmas we need as prerequisites. The proving is as complicated as
using pattern matching on each variable. This is not the best solution
we want. We should benefit more from the observation that any equation
of integers can be turned into equation of natural numbers\ref{obs-int-nat}.










However, if we prove
the laws for quotient integers, it is much simpler since there is
only one case to prove.

However, if we prove it for the quotient integers, it is much
easier. In fact, it is in generally automatically provable. Since the
equality of any two quotient integers is essentially the equality of two
natural numbers after normalising. To prove the distributivity, the
simpliest way is to use semiring
solver for natural numbers. $\AgdaFunction{DistributesOverˡ}$ means
that the distributivity of the first operators over the second one.

\begin{remark}
A ring solver is an automatic equation checker for rings e.g.\ the
ring of integers. It is implemented based on the theory described in
"Proving Equalities in a Commutative Ring Done Right in Coq" by
Grégoire and Mahboubi \cite{gregoire2005proving}.

\end{remark}

\begin{code}
\\
\>\AgdaFunction{distˡ} \AgdaSymbol{:} \<[9]%
\>[9]\AgdaFunction{\_*\_} \AgdaFunction{DistributesOverˡ} \AgdaFunction{\_+\_}\<%
\\
\>\AgdaFunction{distˡ} \AgdaSymbol{(}\AgdaBound{a} \AgdaInductiveConstructor{,} \AgdaBound{b}\AgdaSymbol{)} \AgdaSymbol{(}\AgdaBound{c} \AgdaInductiveConstructor{,} \AgdaBound{d}\AgdaSymbol{)} \AgdaSymbol{(}\AgdaBound{e} \AgdaInductiveConstructor{,} \AgdaBound{f}\AgdaSymbol{)} \AgdaSymbol{=} \AgdaFunction{solve} \AgdaNumber{6} \<[40]%
\>[40]\<%
\\
\>[0]\AgdaIndent{2}{}\<[2]%
\>[2]\AgdaSymbol{(λ} \AgdaBound{a} \AgdaBound{b} \AgdaBound{c} \AgdaBound{d} \AgdaBound{e} \AgdaBound{f} \AgdaSymbol{→} \AgdaBound{a} \AgdaFunction{:*} \AgdaSymbol{(}\AgdaBound{c} \AgdaFunction{:+} \AgdaBound{e}\AgdaSymbol{)} \AgdaFunction{:+} \AgdaBound{b} \AgdaFunction{:*} \AgdaSymbol{(}\AgdaBound{d} \AgdaFunction{:+} \AgdaBound{f}\AgdaSymbol{)} \AgdaFunction{:+}\<%
\\
\>[2]\AgdaIndent{6}{}\<[6]%
\>[6]\AgdaSymbol{(}\AgdaBound{a} \AgdaFunction{:*} \AgdaBound{d} \AgdaFunction{:+} \AgdaBound{b} \AgdaFunction{:*} \AgdaBound{c} \AgdaFunction{:+} \AgdaSymbol{(}\AgdaBound{a} \AgdaFunction{:*} \AgdaBound{f} \AgdaFunction{:+} \AgdaBound{b} \AgdaFunction{:*} \AgdaBound{e}\AgdaSymbol{))}\<%
\\
\>[2]\AgdaIndent{6}{}\<[6]%
\>[6]\AgdaFunction{:=}\<%
\\
\>[2]\AgdaIndent{6}{}\<[6]%
\>[6]\AgdaBound{a} \AgdaFunction{:*} \AgdaBound{c} \AgdaFunction{:+} \AgdaBound{b} \AgdaFunction{:*} \AgdaBound{d} \AgdaFunction{:+} \AgdaSymbol{(}\AgdaBound{a} \AgdaFunction{:*} \AgdaBound{e} \AgdaFunction{:+} \AgdaBound{b} \AgdaFunction{:*} \AgdaBound{f}\AgdaSymbol{)} \AgdaFunction{:+}\<%
\\
\>[2]\AgdaIndent{6}{}\<[6]%
\>[6]\AgdaSymbol{(}\AgdaBound{a} \AgdaFunction{:*} \AgdaSymbol{(}\AgdaBound{d} \AgdaFunction{:+} \AgdaBound{f}\AgdaSymbol{)} \AgdaFunction{:+} \AgdaBound{b} \AgdaFunction{:*} \AgdaSymbol{(}\AgdaBound{c} \AgdaFunction{:+} \AgdaBound{e}\AgdaSymbol{)))} \AgdaInductiveConstructor{refl} \AgdaBound{a} \AgdaBound{b} \AgdaBound{c} \AgdaBound{d} \AgdaBound{e} \AgdaBound{f}\<%
\\
\end{code}

The utilization of ring solver can be simplified even further by
adopting "reflection". It helps us quote the type of the goal so
that we can define a function that automatically do it without
explicitly providing the equations. There is already some work
done by van der Walt \cite{van2012reflection}. It can be seen as an
analogy of the "ring" tactic from Coq.



The main drawback of this method is the type verification of the
terms automatically generated requires more computations than the terms
we manually construct. The optimization of Agda already shows a big
improvement in this technical efficiency issue. However it will affect
other functions that use this proof. It may heavily slow the type
checking.
Therefore although it is still very convenient to use the ring solver to prove any proposition for the quotient
integers, we decide to prove commonly used properties for the commutative
ring of the integers by hand. Luckily it is still much simpler than the ones
for the set of integers $\Z$.


This is a good example of how the definable quotient structure helps
simplifying definiting and proving about new types based on existing
functions and theorems. 


\section{Rational numbers}

\todo{add this paragraph:}
From a discussion on Agda mailing list, a better way to manipulate the
rational numbers is required. The set definition of rational numbers
contains a proof of the coprime property which means that whenther you
want to define a function, you have to reduce the fraction until the
enumerator and denominator are comprime. However from the experience
of using fracitons in arithematics, we can leave them without reducing
and it is still correct. Therefore, it is better for us to have the
flexibility of using the un-reduced fractions or reduced ones in
operations. With the help of definable quotient structures, it is
feasible. We can define operations and prove theorems about rational
numbers using the setoid definition and lift them to the set
version. It makes the definition much easier and whenever we want the
reduced fraction, we just need to apply the normalisation function
within the structure. It also improves the efficiency of programs
using rational numbers, even some people claims that it will be too
large such that it is less efficient.




The quotient of rational numbers is better known than the previous
quotient of integers. We usually write two integers $m$ and $n$ ($n$ is not zero) in
fractional form $\frac{m}{n}$ to represent a rational number. Alternatively we
can use an integer and a positive natural number such that it is
simpler to exclude 0 in the denominator. Two fractions are equal if
they are reduced to the same irreducible term. If the numerator and
denominator of a fraction are coprime, it is said to be an irreducible
fraction. Based on this observation, it is naturally to form a definable quotient, where the base type is 

$$\Q_0 = \Z \times \N$$

The integer stands for the \emph{numerator} and the natural number is
\emph{denominator-1} (We use $N$ for $N^{+}$ to avoid invalid
fractions from construction rather than from zero test)

In Agda, to make the terms more meaningful we define it as

\begin{code}
\\
\>\AgdaKeyword{data} \AgdaDatatype{ℚ₀} \AgdaSymbol{:} \AgdaPrimitiveType{Set} \AgdaKeyword{where}\<%
\\
\>[-1]\AgdaIndent{2}{}\<[2]%
\>[2]\AgdaInductiveConstructor{\_/suc\_} \AgdaSymbol{:} \AgdaSymbol{(}\AgdaBound{n} \AgdaSymbol{:} \AgdaDatatype{ℤ}\AgdaSymbol{)} \AgdaSymbol{→} \AgdaSymbol{(}\AgdaBound{d} \AgdaSymbol{:} \AgdaDatatype{ℕ}\AgdaSymbol{)} \AgdaSymbol{→} \AgdaDatatype{ℚ₀}\<%
\\
\end{code}

In mathematics, to judge the equality of two fractions, it is easier to conduct the following conversion,

$$ \frac{a}{b} = \frac{c}{d} \iff a \times d = c \times b $$

Therefore the equivalence relation can be defined as,

\begin{code}
\\
\>\AgdaFunction{\_*\_} \AgdaSymbol{:} \AgdaDatatype{ℤ} \AgdaSymbol{→} \AgdaDatatype{ℕ} \AgdaSymbol{→} \AgdaDatatype{ℤ}\<%
\\
\>\AgdaSymbol{(}\AgdaInductiveConstructor{+} \AgdaBound{x}\AgdaSymbol{)} \AgdaFunction{*} \AgdaBound{d} \AgdaSymbol{=} \AgdaInductiveConstructor{+} \AgdaSymbol{(}\AgdaBound{x} \AgdaPrimitive{ℕ*} \AgdaBound{d}\AgdaSymbol{)}\<%
\\
\>\AgdaSymbol{(}\AgdaInductiveConstructor{-suc} \AgdaBound{x}\AgdaSymbol{)} \AgdaFunction{*} \AgdaInductiveConstructor{0} \AgdaSymbol{=} \AgdaInductiveConstructor{+} \AgdaNumber{0}\<%
\\
\>\AgdaSymbol{(}\AgdaInductiveConstructor{-suc} \AgdaBound{x}\AgdaSymbol{)} \AgdaFunction{*} \AgdaInductiveConstructor{suc} \AgdaBound{d} \AgdaSymbol{=} \AgdaInductiveConstructor{-suc} \AgdaSymbol{(}\AgdaBound{x} \AgdaPrimitive{ℕ+} \AgdaInductiveConstructor{suc} \AgdaBound{x} \AgdaPrimitive{ℕ*} \AgdaBound{d}\AgdaSymbol{)}\<%
\\
%
\\
\>\AgdaFunction{\_∼\_} \<[35]%
\>[35]\AgdaSymbol{:} \AgdaFunction{Rel} \AgdaDatatype{ℚ₀} \AgdaSymbol{\_}\<%
\\
\>\AgdaSymbol{(}\AgdaBound{n1} \AgdaInductiveConstructor{/suc} \AgdaBound{d1}\AgdaSymbol{)} \AgdaFunction{∼} \AgdaSymbol{(}\AgdaBound{n2} \AgdaInductiveConstructor{/suc} \AgdaBound{d2}\AgdaSymbol{)} \AgdaSymbol{=} \<[31]%
\>[31]\AgdaBound{n1} \AgdaFunction{*} \AgdaInductiveConstructor{suc} \AgdaBound{d2} \AgdaDatatype{≡} \AgdaBound{n2} \AgdaFunction{*} \AgdaInductiveConstructor{suc} \AgdaBound{d1}\<%
\\
\end{code}

The normal form of rational numbers, namely the quotient type in this quotient is the set of irreducible fractions. We only need to add a restriction that the numerator and denominator is coprime,

$$\Q = \Sigma (n \colon \Z). \Sigma (d \colon \N). \mathsf{coprime} \,n \,(d +1)$$

It can be defined as follows which is available in standard library,

\begin{code}
\\
\>\AgdaKeyword{record} \AgdaRecord{ℚ} \AgdaSymbol{:} \AgdaPrimitiveType{Set} \AgdaKeyword{where}\<%
\\
\>[0]\AgdaIndent{2}{}\<[2]%
\>[2]\AgdaKeyword{field}\<%
\\
\>[2]\AgdaIndent{4}{}\<[4]%
\>[4]\AgdaField{numerator} \<[18]%
\>[18]\AgdaSymbol{:} \AgdaDatatype{ℤ}\<%
\\
\>[2]\AgdaIndent{4}{}\<[4]%
\>[4]\AgdaField{denominator-1} \AgdaSymbol{:} \AgdaDatatype{ℕ}\<%
\\
\>[2]\AgdaIndent{4}{}\<[4]%
\>[4]\AgdaField{isCoprime} \<[18]%
\>[18]\AgdaSymbol{:} \AgdaFunction{True} \AgdaSymbol{(}\AgdaFunction{C.coprime?} \AgdaFunction{∣} \AgdaBound{numerator} \AgdaFunction{∣} \AgdaSymbol{(}\AgdaInductiveConstructor{suc} \AgdaBound{denominator-1}\AgdaSymbol{))}\<%
\\
%
\\
\>[0]\AgdaIndent{2}{}\<[2]%
\>[2]\AgdaFunction{denominator} \AgdaSymbol{:} \AgdaDatatype{ℤ}\<%
\\
\>[0]\AgdaIndent{2}{}\<[2]%
\>[2]\AgdaFunction{denominator} \AgdaSymbol{=} \AgdaInductiveConstructor{+} \AgdaInductiveConstructor{suc} \AgdaFunction{denominator-1}\<%
\\
%
\\
\>[0]\AgdaIndent{2}{}\<[2]%
\>[2]\AgdaFunction{coprime} \AgdaSymbol{:} \AgdaFunction{Coprime} \AgdaFunction{numerator} \AgdaFunction{denominator}\<%
\\
\>[0]\AgdaIndent{2}{}\<[2]%
\>[2]\AgdaFunction{coprime} \AgdaSymbol{=} \AgdaFunction{toWitness} \AgdaFunction{isCoprime}\<%
\\
\end{code}

The normalisation function is an implementation of the reducing
process. But first we need to  the |gcd| function which calculates the greatest common
divisor can help us reduce the fraction and give us the proof of
coprime. First we need to define the conversion from the results of GCD
to normal rational numbers (the full definition can be found in
Appendix \autoref{rational-gcd}),

\begin{code}
\\
\>\AgdaFunction{GCD′→ℚ} \AgdaSymbol{:} \AgdaSymbol{∀} \AgdaBound{x} \AgdaBound{y} \AgdaBound{di} \AgdaSymbol{→} \AgdaBound{y} \AgdaFunction{≢} \AgdaNumber{0} \AgdaSymbol{→} \AgdaDatatype{C.GCD′} \AgdaBound{x} \AgdaBound{y} \AgdaBound{di} \AgdaSymbol{→} \AgdaRecord{ℚ}\<%
\\
\end{code}

To normalise a fractional, we split it into 3 cases with respect to
the numerator. The idea is to calculate the "gcd'' and then use the
above function to get the normalised rational number.

\begin{code}
\\
\>\AgdaFunction{[\_]} \AgdaSymbol{:} \AgdaDatatype{ℚ₀} \AgdaSymbol{→} \AgdaRecord{ℚ}\<%
\\
\>\AgdaFunction{[} \AgdaSymbol{(}\AgdaInductiveConstructor{+} \AgdaInductiveConstructor{0}\AgdaSymbol{)} \AgdaInductiveConstructor{/suc} \AgdaBound{d} \AgdaFunction{]} \AgdaSymbol{=} \AgdaInductiveConstructor{ℤ.+\_} \AgdaNumber{0} \AgdaFunction{÷} \AgdaNumber{1}\<%
\\
\>\AgdaFunction{[} \AgdaSymbol{(}\AgdaInductiveConstructor{+} \AgdaSymbol{(}\AgdaInductiveConstructor{suc} \AgdaBound{n}\AgdaSymbol{))} \AgdaInductiveConstructor{/suc} \AgdaBound{d} \AgdaFunction{]} \AgdaKeyword{with} \AgdaFunction{gcd} \AgdaSymbol{(}\AgdaInductiveConstructor{suc} \AgdaBound{n}\AgdaSymbol{)} \AgdaSymbol{(}\AgdaInductiveConstructor{suc} \AgdaBound{d}\AgdaSymbol{)}\<%
\\
\>\AgdaFunction{[} \AgdaSymbol{(}\AgdaInductiveConstructor{+} \AgdaInductiveConstructor{suc} \AgdaBound{n}\AgdaSymbol{)} \AgdaInductiveConstructor{/suc} \AgdaBound{d} \AgdaFunction{]} \AgdaSymbol{|} \AgdaBound{di} \AgdaInductiveConstructor{,} \AgdaBound{g} \AgdaSymbol{=} \AgdaFunction{GCD′→ℚ} \AgdaSymbol{(}\AgdaInductiveConstructor{suc} \AgdaBound{n}\AgdaSymbol{)} \AgdaSymbol{(}\AgdaInductiveConstructor{suc} \AgdaBound{d}\AgdaSymbol{)} \AgdaBound{di} \AgdaSymbol{(λ} \AgdaSymbol{())} \AgdaSymbol{(}\AgdaFunction{C.gcd-gcd′} \AgdaBound{g}\AgdaSymbol{)}\<%
\\
\end{code}


The embedding function is simple. We only need to forget the coprime proof in the normal form,

\begin{code}
\\
\>\AgdaFunction{⌜\_⌝} \AgdaSymbol{:} \AgdaRecord{ℚ} \AgdaSymbol{→} \AgdaDatatype{ℚ₀}\<%
\\
\>\AgdaFunction{⌜} \AgdaBound{x} \AgdaFunction{⌝} \AgdaSymbol{=} \AgdaSymbol{(}\AgdaFunction{ℤcon} \AgdaSymbol{(}\AgdaFunction{ℚ.numerator} \AgdaBound{x}\AgdaSymbol{))} \AgdaInductiveConstructor{/suc} \AgdaSymbol{(}\AgdaFunction{ℚ.denominator-1} \AgdaBound{x}\AgdaSymbol{)}\<%
\\
\end{code} 

Similarly, we are able to construct the setoid, the prequotient and
then the definable quotient of rational numbers. We can benefit from
the ease of defining operators and proving theorems on setoids while
still using the normal form of rational numbers, the lifted operators
and properties which are safer.

The same approach works here as well. Since we can easily embed the natural numbers into integers, the
equations of the quotient rational numbers are degraded to equations of the
integers.  The commutative ring of integers also enable us to prove
all properties of rational numbers automatically.

 

\chapter{Real numbers and other undefinable quotients}
\label{rl}

If you mean something specific, always write "the"

In the previous chapter, only definable quotient types are
investigated. But some other
types are undefinable in \itt{} without quotients. In this chapter, the real numbers,
the multisets and
the partiality monad and also the proofs that
the undefinablity of them will be disccussed in detail.


\section{Real numbers}

We have several choices to represent real numbers. We choose Cauchy
sequences of rational numbers to represent real numbers \cite{bis:85}.

$$\R_{0} = \set{s : \N\to\Q \mid \forall\varepsilon
  :\Q,\varepsilon>0\to\exists m:\N, \forall i:\N, i>m\to \vert  s_i -
  s_m \vert  <\varepsilon}$$

It is implementable in Type Theory, but there is a problem of the
choice of type of the second part of the cauchy sequence, namely the property
that it is a cauchy sequence. Do we distinguish the same sequences
with different proofs? Logically speaking we should not. It means that
we need to truncate it to proposition but we will lose the important
tool to guess what the real number is. To avoid this problem, we could
use an alternative equivalent definition which is a subset of $\R_{0}$:

$$\R_0' = \set{f: \N\to\Q \mid \forall k
  :\N,\forall m , n > k, \to \vert  f_m -
  f_n \vert  < \frac{1}{k}}$$

With the definition, the condition part is propositional and we can
guess the number by applying any number k to the sequence and we know
the interval where it should be located.

Different cauchy sequences can represent the same number. Therefore an equivalence relation\footnote{
The Agda version is in Appendix} is expected. In mathematics two Cauchy sequences $\R_0$ are said to be
equal if their pointwise difference converges to zero,

$$r \sim s = \forall\varepsilon :\Q,\varepsilon>0\to\exists m:\N,
\forall i:\N, i>m\to \vert  r_i - s_i \vert <\varepsilon$$

\subsection{Non-normalizability of Cauchy Sequences}

To prove that it is impossible to give a full definable quotient
structure of real numbers with the setoid of cauchy sequences, we
could show it by proving that it is impossible to define a
normalisation function for the cauchy sequences.

\begin{definition}\label{def:nor}
We say that a quotient structure $\qset{A}$ is definable via a
normalisation, if we have a normalisation function
 \begin{equation}
  nf : A \to A
 \end{equation}
 with the property that it respects $\sim$
\begin{equation}
 p : \prd{c_1,c_2 : A} c_1 \sim c_2 \to nf(c_1) = nf(c_2).
\end{equation}
such that
 \begin{equation}
 q : \prd{c : A}  {nf(c)} \sim c.
 \end{equation}
\end{definition}

It is equivalent to say that we have a definable quotient structure in
the sense of \cite{aan}, because we can form the set of equivalence
classes as

\begin{equation*}
Q \defeq \Sigma_{c : \R_0} nf(c) = c
\end{equation*}

where the second part is propositional, and the "normalisation''
function can be defined as

\begin{equation*}
[c] \defeq nf(c) , p(q)
\end{equation*}

and the embedding function is just the first projection. The properties can be verified
easily.

In the other way around, the true normalisation function is just

\begin{equation*}
n \defeq emb \circ [\_]
\end{equation*}

and the properties hold as well.


We have made an attempt in the original version of our \cite{aan}
draft, but there is something important problem pointed out by Martin Escardo. Laterly, Nicolai
Kraus suggests to fix the proof by proving it as a meta-theoretical
property. We will show an adaption of his proof here.

\textbf{Some preliminaries}

In fact the proof is mainly conducted using topological tools. The
following definitions are helpful for someone who are not so familiar
with topological concepts.


\begin{definition}
Metric space. In mathematics, a metric space is a set where a notion
of distance (called a metric) between elements of the set is defined.

A metric space is an ordered pair $(M , d)$ where M is a set and d is a metric on M:
\begin{enumerate}
\item M is a set,
\item and $d : M \times M \rightarrow \R$ s.t.
\item $d (x , y) = 0 \iff x = y$
\item $d(x,y)=d(y,x)$
\item $d(x,y)+d(y,z) \ge d(x,z)$ 
\end{enumerate}
\end{definition}

We usually define a standard topological structure for discrete types.

\begin{itemize}

\item $(\bool , h)$ where 
$h(m,n) =
\left\{
	\begin{array}{ll}
		0  & \mbox{if } m = n \\
		1 & \mbox{if } m \neq n
	\end{array}
\right.
$

\item $(\N , d)$ where 
$d(m,n) =
\left\{
	\begin{array}{ll}
		0  & \mbox{if } m = n \\
		1 & \mbox{if } m \neq n
	\end{array}
\right.
$

\item $(\Q , e)$ where 
$e(m,n) =
\left\{
	\begin{array}{ll}
		0  & \mbox{if } m = n \\
		1 & \mbox{if } m \neq n
	\end{array}
\right.
$

\end{itemize}

We use a slightly different definition of cauchy sequences of rational
numbers here:

\begin{definition}
 We call a function $f : \PN \to \Q$\footnote{we
   use $\PN$ instead of $\N$ because $n$ must be positive in $\frac 1 n$} a \emph{Cauchy Sequence} if it satisfies
 \begin{equation}
  \iscauchy (f) \defeq \fa{n}{\PN} \fa{m} {\PN} m > n \to |f(n) - f(m)| < \frac 1 n. \label{eq:prop-property-short}
 \end{equation}
 The type of Cauchy Sequences is thus
 \begin{equation*}
  \R_0 \defeq \sm{f : \PN \to \Q} \iscauchy (f).
 \end{equation*}
\end{definition}

And the standard metric space for the sequences $\PN \rightarrow \Q $
is defined by the distance function

\begin{equation}
g(f_1,f_2) = 2^{-\mathsf{inf}\{k \in \N \; | \; f_1(k) \not= f_2(k)\}}
\end{equation}

Among all these standard metric spaces, It is a folklore that all
definable functions are continous.

\begin{theorem}\label{defcon}
All definable functions are continous.
\end{theorem}

Let us introduce the following auxiliary definition:
\begin{definition}
 For a sequence $f : \N \to \Q$, we say that $f$ is \emph{Cauchy with
   factor $k$}, written as $\iscauchy_k$, for some $k \in \Q$, $k > 0$, if
 \begin{equation}
  \iscauchy_k(f) \; \defeq \; \fapairs{n}{m}{\N} m > n \to |f(n) - f(m)| < \frac 1 {k \cdot n}. \label{eq:cauchy-aux}
 \end{equation}
 The usual Cauchy condition $\iscauchy$ is therefore ``Cauchy with factor $1$''.
\end{definition}

\begin{remark} If we claim a function $f$ is defined on $\R_0$ that
  respects $\sim$, it means that we have a proof
\begin{equation}
 p : \prd{c_1,c_2 : \R_0} c_1 \sim c_2 \to f(c_1) = f(c_2).
\end{equation}

\end{remark}

Now we have enough tools to prove the following proposition.

\begin{proposition} \label{prop:main}
 \textbf{"$\R_0 / \sim$" is connected.} It means that any continuous function $f$ 
 \begin{equation}
  f : \R_0 \to \bool
 \end{equation} that respects $\sim$ is constant. We prove that it is impossible to find $c_1, c_2:
\R_0$ such that $f(c_1) \not= f(c_2)$ meta-theoretically, instead of
deriving a proof term of this in type theory.

The definability of function implies that it is a continous function \ref{defcon}
between the standard metric spaces for $\R_0$ and $\bool$\footnote{The
  metric of $\R$ comes from the first component. Technically, if
  $\R_0$ is defined by \ref{eq:usedbytxa}, this would not make it a
  metric space (as the distance between non-equal elements could be
  $0$); however, this would not matter, and for our definition, there
  is no problem anyway.}.

\end{proposition}



\begin{proof}

The general idea is to interpret our definition in classical mathematics,
assume we have a non-constant function and deduce a
contradiction.
 
Consider the ``naive'' set model (with ``classical standard mathematics'' as meta-theory). This clearly works if we are in a minimalisitic type theory with $\Pi$, $\Sigma$, $\mathsf{W}$, $=$, $\mathbb N$; however, if we restrict ourselves to the types in the lowest universe of homotopy type theory (which is enough), it also works for HoTT.
 We use $\tometa \cdot$ as an interpreation function; 
 for example, we write $\R$ for the field of real numbers which can be defined as $\tometa {\R_0} / \tometa{\sim}$. 
 By abuse of notation, we write $\tometa {\R_0}$ for the set of Cauchy sequences in the model that fulfill the Cauchy condition, without the actual proof thereof. This is justified as this property is propositional. 
 
 For readability, we use the symbol $=$ for equality in the theory as well in the model, and we do not use the semantic brackets for natural numbers such as $2$ or $4$.
 In the model, we use $\clim \cdot : \tometa {\R_0} \to \R$ as the
 function that gives us the limit of a Cauchy sequence (Not all
 functions in the ``naive'' set model have to be continuous).
 Thus, for $r : \R_0$, we write $\climeta r \in \R$ for the real number it represents.

Assume $f, p$ are given. We prove that $\tometa f : \tometa {\R_0} \to \tometa \bool$ is constant in the model, which implies the statement of Proposition \ref{prop:main}.

Thus, assume $\tometa f$ is non-constant,  there are $c_1, c_2 : \tometa {\R_0}$ with $\tometa f(c_1) \not= \tometa f(c_2)$. 

 Define
 \begin{align}
  m_1 & \defeq \mathsf{sup} \{\clim d \in \R \; | \; d \in \tometa {\R_0}, \clim d \leq \mathsf{max}(\clim{c_1}, \clim{c_2}), \tometa f(d) = \tometa \btrue \}\\
  m_2 & \defeq \mathsf{sup} \{\clim d \in \R \; | \; d \in \tometa {\R_0}, \clim d \leq \mathsf{max}(\clim{c_1}, \clim{c_2}), \tometa f(d) = \tometa \bfalse \}
 \end{align}
 (note that one of these two necessarily has to be $\clim{c_1}$ or $\clim{c_2}$, whichever is bigger).
 Set $m \defeq \mathsf{min}(m_1,m_2)$, and we can observe that in
 \emph{every} neighborhood $U$ of $m$, given any $t$,
we can always find another point $x \in U$ such that $x = \clim e$
(for some $e$) with $\tometa f(e) \not= \tometa f(t)$.
 
 Let $c \in \tometa {\R_0}$ be a Cauchy sequence such that $\clim c$ is equal to $m$. 
 We may assume that $c$ satisfies the condition
 $\tometa{\iscauchy_5}$.\footnote{the factor 5 is chosen due to the
   need of a later proof.} 
%  The slight difference ($\frac 1 n$ is replaced by $\frac 1 {4n}$) will be important.
%  Without loss of generality, assume $\tometa{f}(c) = \tometa \btrue$. % not needed.
 
 As $f$ (and thereby $\tometa f$) is continuous (remember the metric
 spaces), there exists $n_0 \in \tometa {\N}$ such that for any Cauchy
 sequence $c' \in \tometa {\R_0}$, if the first $n_0$ sequence
 elements of $c'$ coincide with those of $c$ (namely the distance
 $g(c, c') = 2^{-\mathsf{inf}\{k \in \N \; | \; c(k) \not= c'(k)\}}
 \leq 2^{-{n_0}}$ ), then $\tometa f(c') = \tometa f(c)$. 
 Write $U \subset \tometa {\R_0}$ for the set of Cauchy sequences which fulfill this property, and $\clim U \defeq \{\clim d \; | \; d \in U\}$ for the set of reals that $U$ corresponds to.
 
 We claim that $\clim U$ is a neighborhood of $m$. 
 More precisely, we prove: 
 The interval $I \defeq (m - \frac 1 {2n_0} , m+ \frac 1 {2n_0})$ is contained in $\clim U$.
 Let $x \in \R$ be in $I$. There is a sequence $t : \tometa {\N \to \Q}$ such that $\tometa {\iscauchy_{5n_0}}(t)$ and $\clim t = x$. 
 Let us now ``concatenate'' the first $n_0$ elements of the sequence $c$ with $t$, that is, define
 \begin{align}
  &g : \tometa{\N \to \Q} \\
  &g (n) = \begin{cases}
            c(n) & \text{if $n \leq n_0$} \\
            t(n-n_0) & \text{else}.
           \end{cases}
 \end{align}

 Observe that $g$ is also a Cauchy sequence, i.e. $\tometa \iscauchy(g)$: The only thing that needs to be checked is whether the two ``parts'' of $g$ work well together. Let $0 < n \leq n_0$ and $m > n_0$ be two natural numbers. We need to show that
 \begin{equation}
  |g(n) - g(m)| < \frac 1 n.
 \end{equation}
 Calculate
 \begin{align}
  & \quad |g(n) - g(m)| \\
  = & \quad|c(n) - t(m - n_0)| \\
  = & \quad |c(n) - \clim c + \clim c - \clim t + \clim t - t(m - n_0)| \\
  \leq & \quad | c(n) - \clim c | + |\clim c - \clim t| + |\clim t - t(m - n_0)| \\
  \leq &  \quad  \frac{1}{5n}  + \frac{1}{2n_0} + \frac{1}{5n_0 \cdot (m-n_0)} \\
  \leq &  \quad  \frac{1}{5n}  + \frac{1}{2n} + \frac{1}{5n_0} \\
  <  & \quad \frac 1 n.
 \end{align}

From the continuity property of $f$ and the definition of
$g$ we know that $\tometa f (g) = \tometa f(c)$. Clearly, $\clim g =
\clim t = x \in I$. Therefore \emph{all}  $s \in \tometa {\R_0}$ with $\clim
s \in I$, we can use the "concatenation'' approach to find a $g$
satisfies $s \, \tometa \sim \,  g$ (namely $\clim s = \clim g$), and by the condition that $f$ (and
thereby $\tometa f$) repects $\sim$, we can conclude that $\tometa f
(s) = \tometa f (g) = \tometa f(c)$.
 
%  It is easy to see that $g$ is a Cauchy sequence, i.e. $\tometa \iscauchy(g)$,\footnote{Of course, this is why we needed $\iscauchy_5$ instead of $\iscauchy$ earlier. The $5$ is somewhat arbitrary (todo: think about what is actually needed).} with $\clim g = x$. Further, the first $n_0$ elements of $g$ coincide with those of $c$, and thus, $\tometa f(g) = \tometa f(c)$.
 
However, as we have seen, in \emph{every} neighborhood of $m$, and thus in particular in $(m - \frac 1 {2n_0} , m+ \frac 1 {2n_0})$, there is an $x$ such that $x = \clim e$ (for some $e$) with $\tometa f(e) \not= \tometa f(c)$, in contradiction to the just established statement.
\end{proof}

The proposition that "$\R_0 / \sim$ is connected'' implies the
following corollary:

\begin{corollary}\label{dis:con}
Any continous function from $\R_0$ to any discrete type that respects $\sim$ is constant.
\end{corollary}



\begin{theorem}
Any continuous endofunction $f$ on $\R_0$ that respects $\sim$ which means
 \begin{equation}
  p : \prd{c_1,c_2 : \R_0} c_1 \sim c_2 \to f(c_1) = f(c_2).
 \end{equation}
 is constant (in the sense of corollary \ref{dis:con}).
\end{theorem}
\begin{proof}

Assume we have $f,p$ as required. 

To prove $f$ is constant, it is enough to show that the sequence part is
constant because the proof part is propositional. Again, by slight abuse of notation, we write $\tometa f : \tometa{\R_0} \to \tometa {\R_0}$, omitting the proof part of $f$.
 

Given a postive natural number $n : \tometa \PN$, we have a projection
function $\pi_{n} : \tometa{ \R_0} \to \tometa \Q$. Define a function
$h_n : \tometa {\R_0} \to \tometa \Q$ as

 \begin{equation*}
 h_n(c) \defeq \pi_{n} \circ f
 \end{equation*}

By corollary \ref{dis:con} we know that $h_n$ is constant, hence $f$
is constant everywhere, it is enough to show that $f$ is constant.


% By Proposition \ref{prop:main} we know it is a constant function and
% we can observe that $h (\mathsf 0) = true$, can we infer that $\forall
% c : \tometa {\R_0}, \tometa f c = \tometa f (\mathsf 0)$


%  We only need to show that $\pi_1 \circ f$ (the actual sequence) is constant as the proof of being a Cauchy sequence is propositional.\footnote{Even if $\iscauchy$ is not a propositional predicate (as in \ref{eq:usedbytxa}), it will still be true that $m$ is constant. This is simply because $\sim$ is defined only in terms of the actual sequence part.} Again, by slight abuse of notation, we write $\tometa f : \tometa{\R_0} \to \tometa {\R_0}$, omitting the proof part of $f$.
 
%  Given $c : \tometa {\R_0}$, we want to show $\tometa f (c) = \tometa f (\mathsf 0)$, where $\mathsf 0$ is the sequence that is constantly $0$. 
%  To do so, it is enough to show that, for a given $k : \tometa \N$, we have $\tometa{f}(c)(k) = \tometa{f}(\mathsf 0)(k)$. 
%  If this was not true, we would have a function $\tometa \R \to \tometa \bool$, defined by 
%  \begin{equation*}
%   \lambda c . \mathsf{isEquval }\left(\tometa f (c)(k)\right) \left(\tometa f (\mathsf 0)(k)\right)
%  \end{equation*}
% that is not constant, contradicting Proposition \ref{prop:main}.
\end{proof}


\begin{corollary}
 There is no definable normalisation function on $\R_0$ in the sense
 of \ref{def:nor}
\end{corollary}

\begin{corollary}
 $\R_0 / \sim$ is not definable in the sense of \cite{aan}.
\end{corollary}


However, it doesn't imply that we cannot define the set of real numbers in minimalisitic type theory with $\Pi, \Sigma, \mathsf{W}, =,
\N$. The meaning of definability of real numbers is still not clear enough. To make it more precise, we define
it as whether there is a type $A$ in $\mathsf{TT}$ (minimalisitic type
theory) such that its embedding $\tometa A$ in $\mathsf{TT} + \mathsf{Q}$ (type theory
extended with quotient types) is isomorphic to $\tometa \R_0
/ \tometa \sim$ (where it is a valid type). We conjecture that it is still not
definable.

\begin{proof}
Assume the set of real numbers is definable, we have a type $A$ and
its embedding in $\mathsf{TT} + \mathsf{Q}$  is $\tometa A$. It also
gives us a normalisation function and a representative function
between $\tometa \R_0$ and $\tometa A$.

\todo{P(T) -> Connected (T) -> Contractible(T), $\R/ \sim$ is connected but
  not contractible?}

\end{proof}

\subsection{Cauchy completeness of the Cauchy reals}

\todo{expand Cauchy completeness of Cauchy reals: it should rely on
  the axiom of countable choice}

Whether our definition of Cauchy sequence is Cauchy complete? In other
words, is there a representative Cauchy sequence as a limit for every
equivalence class? The answer is no.

In the HoTT book \cite{hott}, an alternative definition is used
instead which is called cauchy approximation. Because every
approximation has been proven to have a limit, it is Cauchy complete.
The definition uses the higher inductive types which will be discussed
in later Chapter \ref{HITs}.

\section{Multisets}

\begin{definition}

Multiset. A multiset (or bag) is a set without the constraint that there is no repetitive elements.

\end{definition}

A set is just a special case of multiset when the
\emph{multiplicity}(the number of the occurences) of every element is one.
Multisets (or bags) are believed to be used in ancient times, but it
is only 
studied by mathematicans from 20th century.

In set theory, a multiset is defined as a pair of a set $A$ and a
occurences counting function $m : A \rightarrow \N$.

However in type theory, the set-theoretical \emph{"set"} is not a
primitive notion and we need to define multisets in a type-theoretic
way. 
In type theoretical language, a multiset can be seen as a unordered
list. It can be encoded as lists which identifies
permutations. To make it simpler we use length-explict list
\emph{Vector} ("Vector A 3" stands for a list of type A of length 3).


The equivalence relation required is as follows:

Given two lists of type $A$ and has length n, $p \, q : Vector\,A\,n$,
they are equivalent if we have an isomorphism between them.

$$ p \sim q = \Sigma \phi : Fin\,n \rightarrow Fin\,n, \phi\,is\,a\,
bijection \wedge \forall x, p_x = q_{\phi(x)}$$


\todo{decidable order -> of course definable, A-> N definable, give a
  explicit proof}

\todo{definable -> split quotient because definable is too general}


\section{Partiality monad}

The partiality monad is a coninductive type which is available in the
standard library of Agda. It is used to encode delayed computation.

\begin{code}
\\
\>\AgdaKeyword{data} \AgdaDatatype{Delay} \AgdaSymbol{(}\AgdaBound{A} \AgdaSymbol{:} \AgdaPrimitiveType{Set}\AgdaSymbol{)} \AgdaSymbol{:} \AgdaPrimitiveType{Set} \AgdaKeyword{where}\<%
\\
\>[0]\AgdaIndent{2}{}\<[2]%
\>[2]\AgdaInductiveConstructor{now} \AgdaSymbol{:} \AgdaBound{A} \AgdaSymbol{→} \AgdaDatatype{Delay} \AgdaBound{A}\<%
\\
\>[0]\AgdaIndent{2}{}\<[2]%
\>[2]\AgdaInductiveConstructor{later} \AgdaSymbol{:} \AgdaDatatype{∞} \AgdaSymbol{(}\AgdaDatatype{Delay} \AgdaBound{A}\AgdaSymbol{)} \AgdaSymbol{→} \AgdaDatatype{Delay} \AgdaBound{A}\<%
\\
\end{code}

A non-terminating program can be defined in a coinductive way.

\begin{code}
\\
\>\AgdaFunction{never} \AgdaSymbol{:} \AgdaSymbol{\{}\AgdaBound{A} \AgdaSymbol{:} \AgdaPrimitiveType{Set}\AgdaSymbol{\}} \AgdaSymbol{→} \AgdaDatatype{Delay} \AgdaBound{A}\<%
\\
\>\AgdaFunction{never} \AgdaSymbol{=} \AgdaInductiveConstructor{later} \AgdaSymbol{(}\AgdaCoinductiveConstructor{♯} \AgdaFunction{never}\AgdaSymbol{)}\<%
\\
\end{code}

We have two equality for the Delay type: strong bisimilarity and weak
bisimilarity. 
Two computation are strongly bisimilar if they are the
same after the same number of steps delay (there can be infinite
steps).

\begin{code}
\\
\>\AgdaKeyword{data} \AgdaDatatype{\_∼\_} \AgdaSymbol{\{}\AgdaBound{A} \AgdaSymbol{:} \AgdaPrimitiveType{Set}\AgdaSymbol{\}} \AgdaSymbol{:} \AgdaDatatype{Delay} \AgdaBound{A} \AgdaSymbol{→} \AgdaDatatype{Delay} \AgdaBound{A} \AgdaSymbol{→} \AgdaPrimitiveType{Set} \AgdaKeyword{where}\<%
\\
\>[2]\AgdaIndent{4}{}\<[4]%
\>[4]\AgdaInductiveConstructor{now} \<[11]%
\>[11]\AgdaSymbol{:} \AgdaSymbol{∀} \AgdaSymbol{\{}\AgdaBound{x}\AgdaSymbol{\}} \AgdaSymbol{→} \AgdaSymbol{(}\AgdaInductiveConstructor{now} \AgdaBound{x}\AgdaSymbol{)} \AgdaDatatype{∼} \AgdaSymbol{(}\AgdaInductiveConstructor{now} \AgdaBound{x}\AgdaSymbol{)}\<%
\\
\>[2]\AgdaIndent{4}{}\<[4]%
\>[4]\AgdaInductiveConstructor{later} \<[11]%
\>[11]\AgdaSymbol{:} \AgdaSymbol{∀} \AgdaSymbol{\{}\AgdaBound{x} \AgdaBound{y}\AgdaSymbol{\}} \AgdaSymbol{(}\AgdaBound{x∼y} \AgdaSymbol{:} \AgdaDatatype{∞} \AgdaSymbol{((}\AgdaFunction{♭} \AgdaBound{x}\AgdaSymbol{)} \AgdaDatatype{∼} \AgdaSymbol{(}\AgdaFunction{♭} \AgdaBound{y}\AgdaSymbol{)))} \AgdaSymbol{→} \AgdaSymbol{(}\AgdaInductiveConstructor{later} \AgdaBound{x}\AgdaSymbol{)} \AgdaDatatype{∼} \AgdaSymbol{(}\AgdaInductiveConstructor{later} \AgdaBound{y}\AgdaSymbol{)}\<%
\\
\end{code}

We inductively define an operator which states that "x terminates with
y" if we write $\AgdaBound{x}\AgdaSymbol{))} \AgdaDatatype{↓} \AgdaBound{y}$.

\begin{code}
\\
\>\AgdaKeyword{infix} \AgdaNumber{4} \_↓\_\<%
\\
%
\\
\>\AgdaKeyword{data} \AgdaDatatype{\_↓\_} \AgdaSymbol{\{}\AgdaBound{A} \AgdaSymbol{:} \AgdaPrimitiveType{Set}\AgdaSymbol{\}} \AgdaSymbol{:} \AgdaDatatype{Delay} \AgdaBound{A} \AgdaSymbol{→} \AgdaBound{A} \AgdaSymbol{→} \AgdaPrimitiveType{Set} \AgdaKeyword{where}\<%
\\
\>[0]\AgdaIndent{2}{}\<[2]%
\>[2]\AgdaInductiveConstructor{nowT} \<[9]%
\>[9]\AgdaSymbol{:} \AgdaSymbol{∀\{}\AgdaBound{a}\AgdaSymbol{\}} \AgdaSymbol{→} \AgdaSymbol{(}\AgdaInductiveConstructor{now} \AgdaBound{a}\AgdaSymbol{)} \AgdaDatatype{↓} \AgdaBound{a}\<%
\\
\>[0]\AgdaIndent{2}{}\<[2]%
\>[2]\AgdaInductiveConstructor{laterT} \AgdaSymbol{:} \AgdaSymbol{∀\{}\AgdaBound{d} \AgdaBound{a}\AgdaSymbol{\}} \AgdaSymbol{→} \AgdaBound{d} \AgdaDatatype{↓} \AgdaBound{a} \AgdaSymbol{→} \AgdaSymbol{(}\AgdaInductiveConstructor{later} \AgdaSymbol{(}\AgdaCoinductiveConstructor{♯} \AgdaBound{d}\AgdaSymbol{))} \AgdaDatatype{↓} \AgdaBound{a}\<%
\\
\end{code}

And two computation are weakly bisimilar if they terminates with the
same value.

\begin{code}
\\
\>\AgdaKeyword{data} \AgdaDatatype{\_≈\_} \AgdaSymbol{\{}\AgdaBound{A} \AgdaSymbol{:} \AgdaPrimitiveType{Set}\AgdaSymbol{\}} \AgdaSymbol{:} \AgdaDatatype{Delay} \AgdaBound{A} \AgdaSymbol{→} \AgdaDatatype{Delay} \AgdaBound{A} \AgdaSymbol{→} \AgdaPrimitiveType{Set} \AgdaKeyword{where}\<%
\\
\>[2]\AgdaIndent{4}{}\<[4]%
\>[4]\AgdaInductiveConstructor{now} \<[11]%
\>[11]\AgdaSymbol{:} \AgdaSymbol{∀} \AgdaSymbol{\{}\AgdaBound{x} \AgdaBound{y} \AgdaBound{a}\AgdaSymbol{\}} \AgdaSymbol{→} \AgdaBound{x} \AgdaDatatype{↓} \AgdaBound{a} \AgdaSymbol{→} \AgdaBound{y} \AgdaDatatype{↓} \AgdaBound{a} \AgdaSymbol{→} \AgdaBound{x} \AgdaDatatype{≈} \AgdaBound{y}\<%
\\
\>[2]\AgdaIndent{4}{}\<[4]%
\>[4]\AgdaInductiveConstructor{later} \<[11]%
\>[11]\AgdaSymbol{:} \AgdaSymbol{∀} \AgdaSymbol{\{}\AgdaBound{x} \AgdaBound{y}\AgdaSymbol{\}} \AgdaSymbol{(}\AgdaBound{x∼y} \AgdaSymbol{:} \AgdaDatatype{∞} \AgdaSymbol{((}\AgdaFunction{♭} \AgdaBound{x}\AgdaSymbol{)} \AgdaDatatype{≈} \AgdaSymbol{(}\AgdaFunction{♭} \AgdaBound{y}\AgdaSymbol{)))} \AgdaSymbol{→} \AgdaSymbol{(}\AgdaInductiveConstructor{later} \AgdaBound{x}\AgdaSymbol{)} \AgdaDatatype{≈} \AgdaSymbol{(}\AgdaInductiveConstructor{later} \AgdaBound{y}\AgdaSymbol{)}\<%
\\
\end{code}


The quotient derived from the equivalent relation (weak bisimilarity) which represent
the set of all computations can also be a good
example of undefinable quotient.

It is equivalent to another definition inductively on either left side
or right side.


\section{Not all connected type is contractible}

Martin's example is that we can define a type which is proved
different but we cannot find a function from it to \emph{2}.

\todo{Martin Escardo's
  \url{http://www.cs.bham.ac.uk/~mhe/agda/FailureOfTotalSeparatedness.html}
should be considered and discussed here, couterexample of Connected ->
contractible(?)}

conjecture: every quotient definable in pure type theory without
quotient is split

one way to prove it is any connected type is contractible.

the question is still open.


This will implies that reals are not only non-split but also
undefianble in general.


1. why quotients simplifies reasoning. (writened up)
2. Undefinablity
3. omega-groupoids model


%\chapter{The Setoid Model}
\label{sm}
% always forgets the "the"

\todo{the Setoid Model}

Quotient types are one of the extensional concepts in Type Theory \cite{hof:phd}. There are several existing intensional models for extensional
concepts. The first one we are going to work with is Altenkirch's
setoid model. To introduce an extensional propositional equality in \itt{}, 
Altenkirch \cite{alti:lics99} proposes an intensional setoid model
with a proof-irrelevant universe of propositions \textbf{Prop}.


\begin{equation}[proof-irr]
\begin{aligned}
\Gamma \vdash P : \textbf{Prop} & & \Gamma \vdash p,q : P & \\
\midrule
& \Gamma \vdash p=q : P & \\
\end{aligned}
\end{equation}

It only contains "propositional'' sets which has at most one
inhabitant. Notice that it is not a definition of types, which means
that we cannot conclude a type is of type \textbf{Prop} if we have a
proof that all
inhabitants are definitionally equal.

The propositional universe is closed under "$\Pi$" and "$\Sigma$", namely dependent functions
and dependent products.

\begin{equation}[\Pi-Prop]
\begin{aligned}
\Gamma \vdash A : \textbf{Set} & & \Gamma,x : A \vdash P \in \textbf{Prop} & \\
\midrule
& \Gamma \vdash \Pi\, x : A.P & \\
\end{aligned}
\end{equation}



\begin{equation}[\Sigma-Prop]
\begin{aligned}
\Gamma \vdash P : \textbf{Prop} & & \Gamma,x : P \vdash Q \in \textbf{Prop} & \\
\midrule
& \Gamma \vdash \Sigma\, x : P.Q & \\
\end{aligned}
\end{equation}



 It is called a setoid model since types are interpreted as setoids.
The solution to introduce the extensional equality is an object type theory defined inside the setoid model which serves as the metatheory. He also proved that the extended type theory generated from the metatheory is decidable and adequate, functional extensionality is
inhabited and it permits large elimination (defining a dependent type by recursion). Within this type theory,
introduction of quotient types is straightforward.

This model is different to a setoid model as an E-category, for instance
the one introduced by Hofmann \cite{hofmann1995interpretation} . An E-category is a category equipped with
an equivalence relation for homsets. To distinguish them, we call this
category \textbf{E-setoids}.  All morphisms of \textbf{E-setoids}
gives rise to types and they are cartesian closed, namely it is a a locally
cartesian closed category (LCCC). Not all morphisms in our category of setoids give
rise to types and it is not an LCCC. Every LCCC can serve as a model for categories with
families but not every category with families has to be an
LCCC. 

\todo{write why this model is not lccc explicitly. refer to Nicolais's
result}



\paragraph{The category of setoids is not a LCCC}

The pullback functor.

\begin{displaymath}
    \xymatrix{X' \ar[r]^{p} \ar[d]_{f^{*}(a)} & Y' \ar[d]^a \\
      X \ar[r]^f& Y }
\end{displaymath}


Observe that $X \rightarrow 1 \cong X$, therefore the pullback of y which is
$X/1 \rightarrow X \times Y / Y$ can be seen as a pullback of X of type $X \rightarrow
X/Y$.

The left adjoint to the pullback functor $f*$ is just the post
composition of $f$ written as $f \circ\_$ or $\Sigma_f$.

\begin{displaymath}
    \xymatrix{X' \ar@{=}[r] \ar[d]_{a} & X' \ar[d]^{\Sigma_f a} \\
      X \ar[r]^f& Y }
\end{displaymath}


However this setoid model is still a model for Type
Theory just like the groupoid model which is a generalisation of it.
To develop this model of type theory in Agda, we have implemented the
categories with families of setoids. 
We build a category with families of setoids to accommodate the types theory described in
\cite{alti:lics99}  so that it is possible to define quotient types
following Martin Hofmann's Paper \cite{hof:95:sm}.  Only necessary
part for the Setoid model will be present here.


\section{An implementation of categories with Families in Agda}

Following the work in \cite{alti:99}, we first define a
proof-irrelevant universe of propositions. We name it as \textbf{hProp}
since \textbf{Prop} is a  reserved word which can't be used and
\textbf{hProp} is a notion from Homotopy Type Theory which we will introduce later.

\section{hProp}

\AgdaHide{

\begin{code}\>\<%
\\
%
\\
\>\AgdaKeyword{open} \AgdaKeyword{import} \AgdaModule{Level}\<%
\\
\>\AgdaKeyword{open} \AgdaKeyword{import} \AgdaModule{Relation.Binary.PropositionalEquality}\<%
\\
%
\\
\>\AgdaKeyword{module} \AgdaModule{hProp} \AgdaSymbol{(}\AgdaBound{ext} \AgdaSymbol{:} \AgdaFunction{Extensionality} \AgdaPrimitive{zero} \AgdaPrimitive{zero}\AgdaSymbol{)} \AgdaKeyword{where}\<%
\\
%
\\
\>\AgdaKeyword{open} \AgdaKeyword{import} \AgdaModule{Relation.Nullary}\<%
\\
\>\AgdaKeyword{open} \AgdaKeyword{import} \AgdaModule{Data.Unit}\<%
\\
\>\AgdaKeyword{open} \AgdaKeyword{import} \AgdaModule{Data.Empty}\<%
\\
\>\AgdaKeyword{open} \AgdaKeyword{import} \AgdaModule{Data.Nat}\<%
\\
\>\AgdaKeyword{open} \AgdaKeyword{import} \AgdaModule{Data.Product}\<%
\\
%
\\
\>\AgdaKeyword{infixr} \AgdaNumber{2} \_⇒\_\<%
\\
%
\\
\>\AgdaKeyword{infixr} \AgdaNumber{3} \_∧\_\<%
\\
%
\\
%
\\
\>\<\end{code}
}

A proof-irrelvant universe only contains sets with at most one inhabitant. 

\begin{code}\>\<%
\\
%
\\
\>\AgdaKeyword{record} \AgdaRecord{hProp} \AgdaSymbol{:} \AgdaPrimitiveType{Set₁} \AgdaKeyword{where}\<%
\\
\>[0]\AgdaIndent{2}{}\<[2]%
\>[2]\AgdaKeyword{constructor} \AgdaInductiveConstructor{hp}\<%
\\
\>[0]\AgdaIndent{2}{}\<[2]%
\>[2]\AgdaKeyword{field}\<%
\\
\>[2]\AgdaIndent{4}{}\<[4]%
\>[4]\AgdaField{prf} \AgdaSymbol{:} \AgdaPrimitiveType{Set}\<%
\\
\>[2]\AgdaIndent{4}{}\<[4]%
\>[4]\AgdaField{Uni} \AgdaSymbol{:} \AgdaSymbol{\{}\AgdaBound{p} \AgdaBound{q} \AgdaSymbol{:} \AgdaBound{prf}\AgdaSymbol{\}} \AgdaSymbol{→} \AgdaBound{p} \AgdaDatatype{≡} \AgdaBound{q}\<%
\\
%
\\
\>\AgdaKeyword{open} \AgdaModule{hProp} \AgdaKeyword{public} \AgdaKeyword{renaming} \AgdaSymbol{(}prf \AgdaSymbol{to} <\_>\AgdaSymbol{)}\<%
\\
%
\\
\>\<\end{code}

We can extract the proof of any propostion $A : hProp$ by using $<>$ and there is always a proof that all inhabitants of it are the same, in other words, if there is any proof of it, the proof is unique. This is not exactly the same as the $Prop$ universe in Altenkirch's approach which is judgemental. It is just a judgement whether a set behaves like a $Proposition$. The $hProp$ we define above is propositional since we can extract the proof of uniqueness.

We would like to have some basic propositions $\top$ and $\bot$. To distinguish them with the ones for non-proof irrelevant propositions which are already available in Agda library, we add a prime to all similar symbols.

\begin{code}\>\<%
\\
%
\\
\>\AgdaFunction{⊤'} \AgdaSymbol{:} \AgdaRecord{hProp}\<%
\\
\>\AgdaFunction{⊤'} \AgdaSymbol{=} \AgdaInductiveConstructor{hp} \AgdaRecord{⊤} \AgdaInductiveConstructor{refl}\<%
\\
%
\\
\>\AgdaFunction{⊥'} \AgdaSymbol{:} \AgdaRecord{hProp}\<%
\\
\>\AgdaFunction{⊥'} \AgdaSymbol{=} \AgdaInductiveConstructor{hp} \AgdaDatatype{⊥} \AgdaSymbol{(λ} \AgdaSymbol{\{}\AgdaBound{p}\AgdaSymbol{\}} \AgdaSymbol{→} \AgdaFunction{⊥-elim} \AgdaBound{p}\AgdaSymbol{)}\<%
\\
%
\\
\>\<\end{code}

We also want the universal and existential quantifier for $hProp$, namely it is closed under $\Pi$-types and $\Sigma$-types.
The universal quantifier of $hProp$ can be axiomitised but we decide to explicitly state that we 
require the functional extensionality to use this module. The reason is that functional extensionality is actually equivalent to the closure under $\Pi$-types.

\begin{code}\>\<%
\\
%
\\
\>\AgdaFunction{∀'} \AgdaSymbol{:} \AgdaSymbol{(}\AgdaBound{A} \AgdaSymbol{:} \AgdaPrimitiveType{Set}\AgdaSymbol{)(}\AgdaBound{P} \AgdaSymbol{:} \AgdaBound{A} \AgdaSymbol{→} \AgdaRecord{hProp}\AgdaSymbol{)} \AgdaSymbol{→} \AgdaRecord{hProp}\<%
\\
\>\AgdaFunction{∀'} \AgdaBound{A} \AgdaBound{P} \AgdaSymbol{=} \AgdaInductiveConstructor{hp} \AgdaSymbol{((}\AgdaBound{x} \AgdaSymbol{:} \AgdaBound{A}\AgdaSymbol{)} \AgdaSymbol{→} \AgdaFunction{<} \AgdaBound{P} \AgdaBound{x} \AgdaFunction{>}\AgdaSymbol{)} \AgdaSymbol{(}\AgdaBound{ext} \AgdaSymbol{(λ} \AgdaBound{x} \AgdaSymbol{→} \AgdaFunction{Uni} \AgdaSymbol{(}\AgdaBound{P} \AgdaBound{x}\AgdaSymbol{)))}\<%
\\
%
\\
\>\<\end{code}


\AgdaHide{
\begin{code}\>\<%
\\
%
\\
\>\AgdaFunction{sig-eq} \AgdaSymbol{:} \AgdaSymbol{\{}\AgdaBound{A} \AgdaSymbol{:} \AgdaPrimitiveType{Set}\AgdaSymbol{\}\{}\AgdaBound{B} \AgdaSymbol{:} \AgdaBound{A} \AgdaSymbol{→} \AgdaPrimitiveType{Set}\AgdaSymbol{\}\{}\AgdaBound{a} \AgdaBound{b} \AgdaSymbol{:} \AgdaBound{A}\AgdaSymbol{\}} \AgdaSymbol{→} \<[43]%
\>[43]\<%
\\
\>[4]\AgdaIndent{9}{}\<[9]%
\>[9]\AgdaSymbol{(}\AgdaBound{p} \AgdaSymbol{:} \AgdaBound{a} \AgdaDatatype{≡} \AgdaBound{b}\AgdaSymbol{)} \AgdaSymbol{→} \<[23]%
\>[23]\<%
\\
\>[4]\AgdaIndent{9}{}\<[9]%
\>[9]\AgdaSymbol{\{}\AgdaBound{c} \AgdaSymbol{:} \AgdaBound{B} \AgdaBound{a}\AgdaSymbol{\}\{}\AgdaBound{d} \AgdaSymbol{:} \AgdaBound{B} \AgdaBound{b}\AgdaSymbol{\}} \AgdaSymbol{→} \<[30]%
\>[30]\<%
\\
\>[4]\AgdaIndent{9}{}\<[9]%
\>[9]\AgdaSymbol{(}\AgdaFunction{subst} \AgdaSymbol{(λ} \AgdaBound{x} \AgdaSymbol{→} \AgdaBound{B} \AgdaBound{x}\AgdaSymbol{)} \AgdaBound{p} \AgdaBound{c} \AgdaDatatype{≡} \AgdaBound{d}\AgdaSymbol{)} \<[37]%
\>[37]\<%
\\
\>[4]\AgdaIndent{9}{}\<[9]%
\>[9]\AgdaSymbol{→} \AgdaDatatype{\_≡\_} \AgdaSymbol{\{\_\}} \AgdaSymbol{\{}\AgdaRecord{Σ} \AgdaBound{A} \AgdaBound{B}\AgdaSymbol{\}} \AgdaSymbol{(}\AgdaBound{a} \AgdaInductiveConstructor{,} \AgdaBound{c}\AgdaSymbol{)} \AgdaSymbol{(}\AgdaBound{b} \AgdaInductiveConstructor{,} \AgdaBound{d}\AgdaSymbol{)}\<%
\\
\>\AgdaFunction{sig-eq} \AgdaInductiveConstructor{refl} \AgdaInductiveConstructor{refl} \AgdaSymbol{=} \AgdaInductiveConstructor{refl}\<%
\\
%
\\
\>\<\end{code}
}

\begin{code}\>\<%
\\
%
\\
%
\\
\>\AgdaFunction{Σ'} \AgdaSymbol{:} \AgdaSymbol{(}\AgdaBound{P} \AgdaSymbol{:} \AgdaRecord{hProp}\AgdaSymbol{)(}\AgdaBound{Q} \AgdaSymbol{:} \AgdaFunction{<} \AgdaBound{P} \AgdaFunction{>} \AgdaSymbol{→} \AgdaRecord{hProp}\AgdaSymbol{)} \AgdaSymbol{→} \AgdaRecord{hProp}\<%
\\
\>\AgdaFunction{Σ'} \AgdaBound{P} \AgdaBound{Q} \AgdaSymbol{=} \AgdaInductiveConstructor{hp} \AgdaSymbol{(}\AgdaRecord{Σ} \AgdaFunction{<} \AgdaBound{P} \AgdaFunction{>} \AgdaSymbol{(λ} \AgdaBound{x} \AgdaSymbol{→} \AgdaFunction{<} \AgdaBound{Q} \AgdaBound{x} \AgdaFunction{>}\AgdaSymbol{))} \<[38]%
\>[38]\<%
\\
\>[9]\AgdaIndent{12}{}\<[12]%
\>[12]\AgdaSymbol{(λ} \AgdaSymbol{\{}\AgdaBound{p}\AgdaSymbol{\}} \AgdaSymbol{\{}\AgdaBound{q}\AgdaSymbol{\}} \AgdaSymbol{→} \<[25]%
\>[25]\<%
\\
\>[9]\AgdaIndent{12}{}\<[12]%
\>[12]\AgdaFunction{sig-eq} \AgdaSymbol{(}\AgdaFunction{Uni} \AgdaBound{P}\AgdaSymbol{)} \AgdaSymbol{(}\AgdaFunction{Uni} \AgdaSymbol{(}\AgdaBound{Q} \AgdaSymbol{(}\AgdaFunction{proj₁} \AgdaBound{q}\AgdaSymbol{))))}\<%
\\
%
\\
\>\<\end{code}

Implication and conjuction which are independent ones of them follow simply.

\begin{code}\>\<%
\\
%
\\
\>\AgdaFunction{\_⇒\_} \AgdaSymbol{:} \AgdaSymbol{(}\AgdaBound{P} \AgdaBound{Q} \AgdaSymbol{:} \AgdaRecord{hProp}\AgdaSymbol{)} \AgdaSymbol{→} \AgdaRecord{hProp}\<%
\\
\>\AgdaBound{P} \AgdaFunction{⇒} \AgdaBound{Q} \AgdaSymbol{=} \<[9]%
\>[9]\AgdaFunction{∀'} \AgdaFunction{<} \AgdaBound{P} \AgdaFunction{>} \AgdaSymbol{(λ} \AgdaBound{\_} \AgdaSymbol{→} \AgdaBound{Q}\AgdaSymbol{)}\<%
\\
%
\\
\>\AgdaFunction{\_∧\_} \AgdaSymbol{:} \AgdaSymbol{(}\AgdaBound{P} \AgdaBound{Q} \AgdaSymbol{:} \AgdaRecord{hProp}\AgdaSymbol{)} \AgdaSymbol{→} \AgdaRecord{hProp}\<%
\\
\>\AgdaBound{P} \AgdaFunction{∧} \AgdaBound{Q} \AgdaSymbol{=} \AgdaFunction{Σ'} \AgdaBound{P} \AgdaSymbol{(λ} \AgdaBound{\_} \AgdaSymbol{→} \AgdaBound{Q}\AgdaSymbol{)}\<%
\\
%
\\
\>\<\end{code}

\AgdaHide{
\begin{code}\>\<%
\\
\>\AgdaKeyword{syntax} ∀' A \AgdaSymbol{(λ} x \AgdaSymbol{→} B\AgdaSymbol{)} \AgdaSymbol{=} ∀'[ x ∶ A ] B


\AgdaKeyword{syntax} Σ' A \AgdaSymbol{(λ} x \AgdaSymbol{→} B\AgdaSymbol{)} \AgdaSymbol{=} Σ'[ x ∶ A ] B

\<\end{code}
}

As long as we have implication and conjuction, more operators on proposition can be defined, for instances negation and logical equivalence.

\begin{code}\>\<%
\\
%
\\
\>\AgdaFunction{¬} \AgdaSymbol{:} \AgdaRecord{hProp} \AgdaSymbol{→} \AgdaRecord{hProp}\<%
\\
\>\AgdaFunction{¬} \AgdaBound{P} \AgdaSymbol{=} \AgdaBound{P} \AgdaFunction{⇒} \AgdaFunction{⊥'} \<[13]%
\>[13]\<%
\\
%
\\
\>\AgdaFunction{\_↔\_} \<[6]%
\>[6]\AgdaSymbol{:} \AgdaSymbol{(}\AgdaBound{P} \AgdaBound{Q} \AgdaSymbol{:} \AgdaRecord{hProp}\AgdaSymbol{)} \AgdaSymbol{→} \AgdaRecord{hProp}\<%
\\
\>\AgdaBound{P} \AgdaFunction{↔} \AgdaBound{Q} \AgdaSymbol{=} \AgdaSymbol{(}\AgdaBound{P} \AgdaFunction{⇒} \AgdaBound{Q}\AgdaSymbol{)} \AgdaFunction{∧} \AgdaSymbol{(}\AgdaBound{Q} \AgdaFunction{⇒} \AgdaBound{P}\AgdaSymbol{)}\<%
\\
%
\\
\>\<\end{code}

\section{Category}

To define category of setoids we should define category first.


\AgdaHide{
\begin{code}\>\<%
\\
%
\\
\>\AgdaSymbol{\{-\#} \AgdaKeyword{OPTIONS} --type-in-type \AgdaSymbol{\#-\}}\<%
\\
%
\\
\>\AgdaKeyword{module} \AgdaModule{Category} \AgdaKeyword{where}\<%
\\
%
\\
\>\AgdaKeyword{open} \AgdaKeyword{import} \AgdaModule{Data.Product}\<%
\\
\>\AgdaKeyword{open} \AgdaKeyword{import} \AgdaModule{Relation.Binary.PropositionalEquality}\<%
\\
%
\\
\>\AgdaKeyword{open} \AgdaKeyword{import} \AgdaModule{Level}\<%
\\
%
\\
\>\<\end{code}
}

\AgdaHide{
\begin{code}\>\<%
\\
%
\\
\>\AgdaKeyword{record} \AgdaRecord{IsCategory}\<%
\\
\>[0]\AgdaIndent{2}{}\<[2]%
\>[2]\AgdaSymbol{(}\AgdaBound{obj} \<[12]%
\>[12]\AgdaSymbol{:} \AgdaPrimitiveType{Set}\AgdaSymbol{)}\<%
\\
%
\\
\>[0]\AgdaIndent{2}{}\<[2]%
\>[2]\AgdaSymbol{(}\AgdaBound{hom} \<[12]%
\>[12]\AgdaSymbol{:} \AgdaBound{obj} \AgdaSymbol{→} \AgdaBound{obj} \AgdaSymbol{→} \AgdaPrimitiveType{Set}\AgdaSymbol{)}\<%
\\
%
\\
\>[0]\AgdaIndent{2}{}\<[2]%
\>[2]\AgdaSymbol{(}\AgdaBound{id} \<[12]%
\>[12]\AgdaSymbol{:} \AgdaSymbol{∀} \AgdaBound{α} \AgdaSymbol{→} \AgdaBound{hom} \AgdaBound{α} \AgdaBound{α}\AgdaSymbol{)}\<%
\\
%
\\
\>[0]\AgdaIndent{2}{}\<[2]%
\>[2]\AgdaSymbol{(}\AgdaBound{[\_⇒\_]\_∘\_} \AgdaSymbol{:} \AgdaSymbol{∀} \AgdaBound{α} \AgdaSymbol{\{}\AgdaBound{β}\AgdaSymbol{\}} \AgdaBound{γ}\<%
\\
\>[2]\AgdaIndent{12}{}\<[12]%
\>[12]\AgdaSymbol{→} \AgdaBound{hom} \AgdaBound{β} \AgdaBound{γ}\<%
\\
\>[2]\AgdaIndent{12}{}\<[12]%
\>[12]\AgdaSymbol{→} \AgdaBound{hom} \AgdaBound{α} \AgdaBound{β}\<%
\\
\>[2]\AgdaIndent{12}{}\<[12]%
\>[12]\AgdaSymbol{→} \AgdaBound{hom} \AgdaBound{α} \AgdaBound{γ}\AgdaSymbol{)}\<%
\\
\>[0]\AgdaIndent{2}{}\<[2]%
\>[2]\AgdaSymbol{:} \AgdaPrimitiveType{Set}\<%
\\
\>[0]\AgdaIndent{2}{}\<[2]%
\>[2]\AgdaKeyword{where}\<%
\\
\>[2]\AgdaIndent{4}{}\<[4]%
\>[4]\AgdaKeyword{constructor} \AgdaInductiveConstructor{IsCatC}\<%
\\
\>[2]\AgdaIndent{4}{}\<[4]%
\>[4]\AgdaKeyword{field}\<%
\\
\>[4]\AgdaIndent{6}{}\<[6]%
\>[6]\AgdaField{id₁} \<[11]%
\>[11]\AgdaSymbol{:} \AgdaSymbol{∀} \AgdaBound{α} \AgdaBound{β} \AgdaSymbol{(}\AgdaBound{f} \AgdaSymbol{:} \AgdaBound{hom} \AgdaBound{α} \AgdaBound{β}\AgdaSymbol{)}\<%
\\
\>[6]\AgdaIndent{11}{}\<[11]%
\>[11]\AgdaSymbol{→} \AgdaBound{[} \AgdaBound{α} \AgdaBound{⇒} \AgdaBound{β} \AgdaBound{]} \AgdaBound{f} \AgdaBound{∘} \AgdaSymbol{(}\AgdaBound{id} \AgdaBound{α}\AgdaSymbol{)} \AgdaDatatype{≡} \AgdaBound{f}\<%
\\
%
\\
\>[0]\AgdaIndent{6}{}\<[6]%
\>[6]\AgdaField{id₂} \<[11]%
\>[11]\AgdaSymbol{:} \AgdaSymbol{∀} \AgdaBound{α} \AgdaBound{β} \AgdaSymbol{(}\AgdaBound{f} \AgdaSymbol{:} \AgdaBound{hom} \AgdaBound{α} \AgdaBound{β}\AgdaSymbol{)}\<%
\\
\>[0]\AgdaIndent{11}{}\<[11]%
\>[11]\AgdaSymbol{→} \AgdaBound{[} \AgdaBound{α} \AgdaBound{⇒} \AgdaBound{β} \AgdaBound{]} \AgdaSymbol{(}\AgdaBound{id} \AgdaBound{β}\AgdaSymbol{)} \AgdaBound{∘} \AgdaBound{f} \AgdaDatatype{≡} \AgdaBound{f}\<%
\\
%
\\
\>[0]\AgdaIndent{6}{}\<[6]%
\>[6]\AgdaField{comp} \AgdaSymbol{:} \AgdaSymbol{∀} \AgdaBound{α} \AgdaSymbol{\{}\AgdaBound{β} \AgdaBound{γ}\AgdaSymbol{\}} \AgdaBound{δ} \AgdaSymbol{(}\AgdaBound{f} \AgdaSymbol{:} \AgdaBound{hom} \AgdaBound{α} \AgdaBound{β}\AgdaSymbol{)} \AgdaSymbol{(}\AgdaBound{g} \AgdaSymbol{:} \AgdaBound{hom} \AgdaBound{β} \AgdaBound{γ}\AgdaSymbol{)} \AgdaSymbol{(}\AgdaBound{h} \AgdaSymbol{:} \AgdaBound{hom} \AgdaBound{γ} \AgdaBound{δ}\AgdaSymbol{)}\<%
\\
\>[0]\AgdaIndent{11}{}\<[11]%
\>[11]\AgdaSymbol{→} \AgdaBound{[} \AgdaBound{α} \AgdaBound{⇒} \AgdaBound{δ} \AgdaBound{]} \AgdaBound{[} \AgdaBound{β} \AgdaBound{⇒} \AgdaBound{δ} \AgdaBound{]} \AgdaBound{h} \AgdaBound{∘} \AgdaBound{g} \AgdaBound{∘} \AgdaBound{f} \AgdaDatatype{≡} \AgdaBound{[} \AgdaBound{α} \AgdaBound{⇒} \AgdaBound{δ} \AgdaBound{]} \AgdaBound{h} \AgdaBound{∘} \AgdaSymbol{(}\AgdaBound{[} \AgdaBound{α} \AgdaBound{⇒} \AgdaBound{γ} \AgdaBound{]} \AgdaBound{g} \AgdaBound{∘} \AgdaBound{f}\AgdaSymbol{)}\<%
\\
%
\\
%
\\
\>\<\end{code}
}


\begin{code}\>\<%
\\
\>\AgdaKeyword{record} \AgdaRecord{Category} \AgdaSymbol{:} \AgdaPrimitiveType{Set} \AgdaKeyword{where}\<%
\\
\>[0]\AgdaIndent{2}{}\<[2]%
\>[2]\AgdaKeyword{constructor} \AgdaInductiveConstructor{CatC}\<%
\\
\>[0]\AgdaIndent{2}{}\<[2]%
\>[2]\AgdaKeyword{field}\<%
\\
\>[2]\AgdaIndent{4}{}\<[4]%
\>[4]\AgdaField{obj} \<[15]%
\>[15]\AgdaSymbol{:} \AgdaPrimitiveType{Set}\<%
\\
%
\\
\>[2]\AgdaIndent{4}{}\<[4]%
\>[4]\AgdaField{hom} \<[15]%
\>[15]\AgdaSymbol{:} \AgdaBound{obj} \AgdaSymbol{→} \AgdaBound{obj} \AgdaSymbol{→} \AgdaPrimitiveType{Set}\<%
\\
%
\\
\>[2]\AgdaIndent{4}{}\<[4]%
\>[4]\AgdaField{id} \<[15]%
\>[15]\AgdaSymbol{:} \AgdaSymbol{∀} \AgdaBound{α}\<%
\\
\>[4]\AgdaIndent{15}{}\<[15]%
\>[15]\AgdaSymbol{→} \AgdaBound{hom} \AgdaBound{α} \AgdaBound{α}\<%
\\
%
\\
\>[0]\AgdaIndent{4}{}\<[4]%
\>[4]\AgdaField{[\_⇒\_]\_∘\_} \<[15]%
\>[15]\AgdaSymbol{:} \AgdaSymbol{∀} \AgdaBound{α} \AgdaSymbol{\{}\AgdaBound{β}\AgdaSymbol{\}} \AgdaBound{γ}\<%
\\
\>[0]\AgdaIndent{15}{}\<[15]%
\>[15]\AgdaSymbol{→} \AgdaBound{hom} \AgdaBound{β} \AgdaBound{γ}\<%
\\
\>[0]\AgdaIndent{15}{}\<[15]%
\>[15]\AgdaSymbol{→} \AgdaBound{hom} \AgdaBound{α} \AgdaBound{β}\<%
\\
\>[0]\AgdaIndent{15}{}\<[15]%
\>[15]\AgdaSymbol{→} \AgdaBound{hom} \AgdaBound{α} \AgdaBound{γ}\<%
\\
%
\\
\>[0]\AgdaIndent{4}{}\<[4]%
\>[4]\AgdaField{isCategory} \AgdaSymbol{:} \AgdaRecord{IsCategory} \AgdaBound{obj} \AgdaBound{hom} \AgdaBound{id} \AgdaBound{[\_⇒\_]\_∘\_}\<%
\\
%
\\
%
\\
\>\<\end{code}

\AgdaHide{
\begin{code}\>\<%
\\
%
\\
\>[0]\AgdaIndent{2}{}\<[2]%
\>[2]\AgdaKeyword{open} \AgdaModule{IsCategory} \AgdaKeyword{public}\<%
\\
%
\\
\>\<\end{code}
}

$isCategory$ contains all the laws for this structure to be a category, for instance the
associativity laws for composition.

\section{Category of setoids}


\AgdaHide{
\begin{code}\>\<%
\\
%
\\
\>\AgdaSymbol{\{-\#} \AgdaKeyword{OPTIONS} --type-in-type \AgdaSymbol{\#-\}}\<%
\\
%
\\
\>\AgdaKeyword{open} \AgdaKeyword{import} \AgdaModule{Level}\<%
\\
\>\AgdaKeyword{open} \AgdaKeyword{import} \AgdaModule{Relation.Binary.PropositionalEquality} \AgdaSymbol{as} \AgdaModule{PE} \AgdaKeyword{hiding} \AgdaSymbol{(}refl\AgdaSymbol{;} sym \AgdaSymbol{;} trans\AgdaSymbol{;} isEquivalence\AgdaSymbol{)}\<%
\\
%
\\
\>\AgdaKeyword{module} \AgdaModule{CategoryOfSetoid} \<[25]%
\>[25]\AgdaSymbol{(}\AgdaBound{ext} \AgdaSymbol{:} \AgdaFunction{Extensionality} \AgdaPrimitive{zero} \AgdaPrimitive{zero}\AgdaSymbol{)} \AgdaKeyword{where}\<%
\\
%
\\
\>\AgdaKeyword{open} \AgdaKeyword{import} \AgdaModule{Cats.Category}\<%
\\
\>\AgdaKeyword{open} \AgdaKeyword{import} \AgdaModule{Function}\<%
\\
\>\AgdaKeyword{open} \AgdaKeyword{import} \AgdaModule{Relation.Binary.Core} \AgdaKeyword{using} \AgdaSymbol{(}\_⇒\_\AgdaSymbol{)}\<%
\\
\>\AgdaKeyword{open} \AgdaKeyword{import} \AgdaModule{Data.Empty}\<%
\\
\>\AgdaKeyword{open} \AgdaKeyword{import} \AgdaModule{HProp} \AgdaBound{ext} \AgdaKeyword{public}\<%
\\
\>\AgdaKeyword{open} \AgdaKeyword{import} \AgdaModule{Data.Unit}\<%
\\
%
\\
\>\<\end{code}
}

\section{Category of setoids}

\begin{code}\>\<%
\\
%
\\
\>\AgdaKeyword{record} \AgdaRecord{HSetoid} \AgdaSymbol{:} \AgdaPrimitiveType{Set₁} \AgdaKeyword{where}\<%
\\
\>[0]\AgdaIndent{2}{}\<[2]%
\>[2]\AgdaKeyword{constructor} \AgdaInductiveConstructor{\_,\_,\_,\_,\_}\<%
\\
\>[0]\AgdaIndent{2}{}\<[2]%
\>[2]\AgdaKeyword{infix} \AgdaNumber{4} \_≈h\_ \_≈\_\<%
\\
\>[0]\AgdaIndent{2}{}\<[2]%
\>[2]\AgdaKeyword{field}\<%
\\
\>[2]\AgdaIndent{4}{}\<[4]%
\>[4]\AgdaField{Carrier} \AgdaSymbol{:} \AgdaPrimitiveType{Set}\<%
\\
\>[2]\AgdaIndent{4}{}\<[4]%
\>[4]\AgdaField{\_≈h\_} \<[13]%
\>[13]\AgdaSymbol{:} \AgdaBound{Carrier} \AgdaSymbol{→} \AgdaBound{Carrier} \AgdaSymbol{→} \AgdaRecord{HProp}\<%
\\
\>[2]\AgdaIndent{4}{}\<[4]%
\>[4]\AgdaField{refl} \<[12]%
\>[12]\AgdaSymbol{:} \AgdaSymbol{\{}\AgdaBound{x} \AgdaSymbol{:} \AgdaBound{Carrier}\AgdaSymbol{\}} \AgdaSymbol{→} \AgdaFunction{<} \AgdaBound{x} \AgdaBound{≈h} \AgdaBound{x} \AgdaFunction{>}\<%
\\
\>[2]\AgdaIndent{4}{}\<[4]%
\>[4]\AgdaField{sym} \<[12]%
\>[12]\AgdaSymbol{:} \AgdaSymbol{\{}\AgdaBound{x} \AgdaBound{y} \AgdaSymbol{:} \AgdaBound{Carrier}\AgdaSymbol{\}} \AgdaSymbol{→} \AgdaFunction{<} \AgdaBound{x} \AgdaBound{≈h} \AgdaBound{y} \AgdaFunction{>} \AgdaSymbol{→} \AgdaFunction{<} \AgdaBound{y} \AgdaBound{≈h} \AgdaBound{x} \AgdaFunction{>}\<%
\\
\>[2]\AgdaIndent{4}{}\<[4]%
\>[4]\AgdaField{trans} \<[12]%
\>[12]\AgdaSymbol{:} \AgdaSymbol{\{}\AgdaBound{x} \AgdaBound{y} \AgdaBound{z} \AgdaSymbol{:} \AgdaBound{Carrier}\AgdaSymbol{\}} \AgdaSymbol{→} \AgdaFunction{<} \AgdaBound{x} \AgdaBound{≈h} \AgdaBound{y} \AgdaFunction{>} \AgdaSymbol{→} \AgdaFunction{<} \AgdaBound{y} \AgdaBound{≈h} \AgdaBound{z} \AgdaFunction{>} \AgdaSymbol{→} \AgdaFunction{<} \AgdaBound{x} \AgdaBound{≈h} \AgdaBound{z} \AgdaFunction{>}\<%
\\
\>[0]\AgdaIndent{2}{}\<[2]%
\>[2]\<%
\\
\>[0]\AgdaIndent{2}{}\<[2]%
\>[2]\AgdaFunction{\_≈\_} \AgdaSymbol{:} \AgdaFunction{Carrier} \AgdaSymbol{→} \AgdaFunction{Carrier} \AgdaSymbol{→} \AgdaPrimitiveType{Set}\<%
\\
\>[0]\AgdaIndent{2}{}\<[2]%
\>[2]\AgdaBound{a} \AgdaFunction{≈} \AgdaBound{b} \AgdaSymbol{=} \AgdaFunction{<} \AgdaBound{a} \AgdaFunction{≈h} \AgdaBound{b} \AgdaFunction{>}\<%
\\
%
\\
\>[0]\AgdaIndent{2}{}\<[2]%
\>[2]\AgdaFunction{PI} \AgdaSymbol{:} \AgdaSymbol{\{}\AgdaBound{x} \AgdaBound{y} \AgdaSymbol{:} \AgdaFunction{Carrier}\AgdaSymbol{\}\{}\AgdaBound{B} \AgdaSymbol{:} \AgdaPrimitiveType{Set}\AgdaSymbol{\}(}\AgdaBound{A} \AgdaSymbol{:} \AgdaBound{x} \AgdaFunction{≈} \AgdaBound{y} \AgdaSymbol{→} \AgdaBound{B}\AgdaSymbol{)\{}\AgdaBound{p} \AgdaBound{q} \AgdaSymbol{:} \AgdaBound{x} \AgdaFunction{≈} \AgdaBound{y}\AgdaSymbol{\}} \AgdaSymbol{→} \AgdaBound{A} \AgdaBound{p} \AgdaDatatype{≡} \AgdaBound{A} \AgdaBound{q}\<%
\\
\>[0]\AgdaIndent{2}{}\<[2]%
\>[2]\AgdaFunction{PI} \AgdaSymbol{\{}\AgdaBound{x}\AgdaSymbol{\}} \AgdaSymbol{\{}\AgdaBound{y}\AgdaSymbol{\}} \AgdaBound{A} \AgdaSymbol{\{}\AgdaBound{p}\AgdaSymbol{\}} \AgdaSymbol{\{}\AgdaBound{q}\AgdaSymbol{\}} \AgdaKeyword{with} \AgdaFunction{Uni} \AgdaSymbol{(}\AgdaBound{x} \AgdaFunction{≈h} \AgdaBound{y}\AgdaSymbol{)} \AgdaSymbol{\{}\AgdaBound{p}\AgdaSymbol{\}} \AgdaSymbol{\{}\AgdaBound{q}\AgdaSymbol{\}}\<%
\\
\>[0]\AgdaIndent{2}{}\<[2]%
\>[2]\AgdaFunction{PI} \AgdaBound{A} \AgdaSymbol{|} \AgdaInductiveConstructor{PE.refl} \AgdaSymbol{=} \AgdaInductiveConstructor{PE.refl}\<%
\\
%
\\
\>[0]\AgdaIndent{2}{}\<[2]%
\>[2]\AgdaFunction{reflexive} \AgdaSymbol{:} \AgdaDatatype{\_≡\_} \AgdaFunction{⇒} \AgdaFunction{\_≈\_}\<%
\\
\>[0]\AgdaIndent{2}{}\<[2]%
\>[2]\AgdaFunction{reflexive} \AgdaInductiveConstructor{PE.refl} \AgdaSymbol{=} \AgdaFunction{refl}\<%
\\
%
\\
\>\AgdaKeyword{open} \AgdaModule{HSetoid} \AgdaKeyword{public} \AgdaKeyword{renaming} \AgdaSymbol{(}refl \AgdaSymbol{to} [\_]refl\AgdaSymbol{;} sym \AgdaSymbol{to} [\_]sym\AgdaSymbol{;} 
  \_≈\_ \AgdaSymbol{to} [\_]\_≈\_ \AgdaSymbol{;} \_≈h\_ \AgdaSymbol{to} [\_]\_≈h\_ \AgdaSymbol{;} Carrier \AgdaSymbol{to} ∣\_∣ \AgdaSymbol{;} trans \AgdaSymbol{to} [\_]trans\AgdaSymbol{)}\<%
\\
%
\\
\>\AgdaFunction{rel} \AgdaSymbol{:} \AgdaSymbol{(}\AgdaBound{A} \AgdaSymbol{:} \AgdaRecord{HSetoid}\AgdaSymbol{)} \AgdaSymbol{→} \AgdaFunction{∣} \AgdaBound{A} \AgdaFunction{∣} \AgdaSymbol{→} \AgdaFunction{∣} \AgdaBound{A} \AgdaFunction{∣} \AgdaSymbol{→} \AgdaRecord{HProp}\<%
\\
\>\AgdaFunction{rel} \AgdaBound{A} \AgdaBound{a} \AgdaBound{b} \AgdaSymbol{=} \AgdaFunction{[} \AgdaBound{A} \AgdaFunction{]} \AgdaBound{a} \AgdaFunction{≈h} \AgdaBound{b}\<%
\\
%
\\
\>\AgdaFunction{[\_]uip} \AgdaSymbol{:} \AgdaSymbol{∀(}\AgdaBound{Γ} \AgdaSymbol{:} \AgdaRecord{HSetoid}\AgdaSymbol{)\{}\AgdaBound{a} \AgdaBound{b} \AgdaSymbol{:} \AgdaFunction{∣} \AgdaBound{Γ} \AgdaFunction{∣}\AgdaSymbol{\}\{}\AgdaBound{p} \AgdaBound{q} \AgdaSymbol{:} \AgdaFunction{[} \AgdaBound{Γ} \AgdaFunction{]} \AgdaBound{a} \AgdaFunction{≈} \AgdaBound{b}\AgdaSymbol{\}} \AgdaSymbol{→} \AgdaBound{p} \AgdaDatatype{≡} \AgdaBound{q}\<%
\\
\>\AgdaFunction{[} \AgdaBound{Γ} \AgdaFunction{]uip} \AgdaSymbol{\{}\AgdaBound{a}\AgdaSymbol{\}} \AgdaSymbol{\{}\AgdaBound{b}\AgdaSymbol{\}} \AgdaSymbol{=} \AgdaFunction{Uni} \AgdaSymbol{(}\AgdaFunction{[} \AgdaBound{Γ} \AgdaFunction{]} \AgdaBound{a} \AgdaFunction{≈h} \AgdaBound{b}\AgdaSymbol{)}\<%
\\
%
\\
\>\<\end{code}

Morphisms between HSetoids (Functors)

\begin{code}\>\<%
\\
%
\\
\>\AgdaKeyword{infix} \AgdaNumber{5} \_⇉\_\<%
\\
%
\\
\>\AgdaKeyword{record} \AgdaRecord{\_⇉\_} \AgdaSymbol{(}\AgdaBound{A} \AgdaBound{B} \AgdaSymbol{:} \AgdaRecord{HSetoid}\AgdaSymbol{)} \AgdaSymbol{:} \AgdaPrimitiveType{Set} \AgdaKeyword{where}\<%
\\
\>[0]\AgdaIndent{2}{}\<[2]%
\>[2]\AgdaKeyword{constructor} \AgdaInductiveConstructor{fn:\_resp:\_}\<%
\\
\>[0]\AgdaIndent{2}{}\<[2]%
\>[2]\AgdaKeyword{field}\<%
\\
\>[2]\AgdaIndent{4}{}\<[4]%
\>[4]\AgdaField{fn} \<[9]%
\>[9]\AgdaSymbol{:} \AgdaFunction{∣} \AgdaBound{A} \AgdaFunction{∣} \AgdaSymbol{→} \AgdaFunction{∣} \AgdaBound{B} \AgdaFunction{∣}\<%
\\
\>[2]\AgdaIndent{4}{}\<[4]%
\>[4]\AgdaField{resp} \AgdaSymbol{:} \AgdaSymbol{\{}\AgdaBound{x} \AgdaBound{y} \AgdaSymbol{:} \AgdaFunction{∣} \AgdaBound{A} \AgdaFunction{∣}\AgdaSymbol{\}} \AgdaSymbol{→} \<[27]%
\>[27]\<%
\\
\>[4]\AgdaIndent{11}{}\<[11]%
\>[11]\AgdaFunction{[} \AgdaBound{A} \AgdaFunction{]} \AgdaBound{x} \AgdaFunction{≈} \AgdaBound{y} \AgdaSymbol{→} \<[25]%
\>[25]\<%
\\
\>[4]\AgdaIndent{11}{}\<[11]%
\>[11]\AgdaFunction{[} \AgdaBound{B} \AgdaFunction{]} \AgdaBound{fn} \AgdaBound{x} \AgdaFunction{≈} \AgdaBound{fn} \AgdaBound{y}\<%
\\
\>\AgdaKeyword{open} \AgdaModule{\_⇉\_} \AgdaKeyword{public} \AgdaKeyword{renaming} \AgdaSymbol{(}fn \AgdaSymbol{to} [\_]fn \AgdaSymbol{;} resp \AgdaSymbol{to} [\_]resp\AgdaSymbol{)}\<%
\\
%
\\
\>\<\end{code}

Identity

\begin{code}\>\<%
\\
%
\\
\>\AgdaFunction{idCH} \AgdaSymbol{:} \AgdaSymbol{\{}\AgdaBound{Γ} \AgdaSymbol{:} \AgdaRecord{HSetoid}\AgdaSymbol{\}} \AgdaSymbol{→} \AgdaBound{Γ} \AgdaRecord{⇉} \AgdaBound{Γ} \<[29]%
\>[29]\<%
\\
\>\AgdaFunction{idCH} \AgdaSymbol{=} \AgdaKeyword{record} \AgdaSymbol{\{} \AgdaField{fn} \AgdaSymbol{=} \AgdaFunction{id}\AgdaSymbol{;} \AgdaField{resp} \AgdaSymbol{=} \AgdaFunction{id}\AgdaSymbol{\}}\<%
\\
%
\\
\>\<\end{code}

Composition

\begin{code}\>\<%
\\
%
\\
\>\AgdaKeyword{infixl} \AgdaNumber{5} \_∘c\_\<%
\\
%
\\
\>\AgdaFunction{\_∘c\_} \AgdaSymbol{:} \AgdaSymbol{∀\{}\AgdaBound{Γ} \AgdaBound{Δ} \AgdaBound{Z}\AgdaSymbol{\}} \AgdaSymbol{→} \AgdaBound{Δ} \AgdaRecord{⇉} \AgdaBound{Z} \AgdaSymbol{→} \AgdaBound{Γ} \AgdaRecord{⇉} \AgdaBound{Δ} \AgdaSymbol{→} \AgdaBound{Γ} \AgdaRecord{⇉} \AgdaBound{Z}\<%
\\
\>\AgdaBound{yz} \AgdaFunction{∘c} \AgdaBound{xy} \AgdaSymbol{=} \AgdaKeyword{record} \<[18]%
\>[18]\<%
\\
\>[4]\AgdaIndent{11}{}\<[11]%
\>[11]\AgdaSymbol{\{} \AgdaField{fn} \AgdaSymbol{=} \AgdaFunction{[} \AgdaBound{yz} \AgdaFunction{]fn} \AgdaFunction{∘} \AgdaFunction{[} \AgdaBound{xy} \AgdaFunction{]fn}\<%
\\
\>[4]\AgdaIndent{11}{}\<[11]%
\>[11]\AgdaSymbol{;} \AgdaField{resp} \AgdaSymbol{=} \AgdaFunction{[} \AgdaBound{yz} \AgdaFunction{]resp} \AgdaFunction{∘} \AgdaFunction{[} \AgdaBound{xy} \AgdaFunction{]resp}\<%
\\
\>[4]\AgdaIndent{11}{}\<[11]%
\>[11]\AgdaSymbol{\}}\<%
\\
\>\<\end{code}

Categorical laws

\begin{code}\>\<%
\\
%
\\
\>\AgdaFunction{id₁} \AgdaSymbol{:} \AgdaSymbol{∀} \AgdaSymbol{\{}\AgdaBound{Γ} \AgdaBound{Δ}\AgdaSymbol{\}(}\AgdaBound{ch} \AgdaSymbol{:} \AgdaBound{Γ} \AgdaRecord{⇉} \AgdaBound{Δ}\AgdaSymbol{)} \AgdaSymbol{→} \AgdaBound{ch} \AgdaFunction{∘c} \AgdaFunction{idCH} \AgdaDatatype{≡} \AgdaBound{ch}\<%
\\
\>\AgdaFunction{id₁} \AgdaBound{ch} \AgdaSymbol{=} \AgdaInductiveConstructor{PE.refl}\<%
\\
%
\\
\>\AgdaFunction{id₂} \AgdaSymbol{:} \AgdaSymbol{∀} \AgdaSymbol{\{}\AgdaBound{Γ} \AgdaBound{Δ}\AgdaSymbol{\}(}\AgdaBound{ch} \AgdaSymbol{:} \AgdaBound{Γ} \AgdaRecord{⇉} \AgdaBound{Δ}\AgdaSymbol{)} \AgdaSymbol{→} \AgdaFunction{idCH} \AgdaFunction{∘c} \AgdaBound{ch} \AgdaDatatype{≡} \AgdaBound{ch}\<%
\\
\>\AgdaFunction{id₂} \AgdaBound{ch} \AgdaSymbol{=} \AgdaInductiveConstructor{PE.refl}\<%
\\
%
\\
\>\AgdaFunction{comp} \AgdaSymbol{:} \AgdaSymbol{∀} \AgdaSymbol{\{}\AgdaBound{Γ} \AgdaBound{Δ} \AgdaBound{Φ} \AgdaBound{Ψ}\AgdaSymbol{\}} \AgdaSymbol{(}\AgdaBound{f} \AgdaSymbol{:} \AgdaBound{Γ} \AgdaRecord{⇉} \AgdaBound{Δ}\AgdaSymbol{)} \AgdaSymbol{(}\AgdaBound{g} \AgdaSymbol{:} \AgdaBound{Δ} \AgdaRecord{⇉} \AgdaBound{Φ}\AgdaSymbol{)} \AgdaSymbol{(}\AgdaBound{h} \AgdaSymbol{:} \AgdaBound{Φ} \AgdaRecord{⇉} \AgdaBound{Ψ}\AgdaSymbol{)}\<%
\\
\>[0]\AgdaIndent{7}{}\<[7]%
\>[7]\AgdaSymbol{→} \AgdaBound{h} \AgdaFunction{∘c} \AgdaBound{g} \AgdaFunction{∘c} \AgdaBound{f} \AgdaDatatype{≡} \AgdaBound{h} \AgdaFunction{∘c} \AgdaSymbol{(}\AgdaBound{g} \AgdaFunction{∘c} \AgdaBound{f}\AgdaSymbol{)}\<%
\\
\>\AgdaFunction{comp} \AgdaBound{f} \AgdaBound{g} \AgdaBound{h} \AgdaSymbol{=} \AgdaInductiveConstructor{PE.refl}\<%
\\
%
\\
\>\AgdaFunction{\_f≈\_} \AgdaSymbol{:} \<[8]%
\>[8]\AgdaSymbol{∀\{}\AgdaBound{Γ} \AgdaBound{Δ} \AgdaSymbol{:} \AgdaRecord{HSetoid}\AgdaSymbol{\}} \AgdaSymbol{→} \AgdaSymbol{(}\AgdaBound{f} \AgdaBound{g} \AgdaSymbol{:} \AgdaBound{Γ} \AgdaRecord{⇉} \AgdaBound{Δ}\AgdaSymbol{)} \AgdaSymbol{→} \AgdaRecord{HProp}\<%
\\
\>\AgdaFunction{\_f≈\_} \AgdaSymbol{\{}\AgdaBound{Γ} \AgdaInductiveConstructor{,} \AgdaBound{\_≈h\_} \AgdaInductiveConstructor{,} \AgdaBound{refl} \AgdaInductiveConstructor{,} \AgdaBound{sym} \AgdaInductiveConstructor{,} \AgdaBound{trans}\AgdaSymbol{\}} \<[37]%
\>[37]\<%
\\
\>[0]\AgdaIndent{5}{}\<[5]%
\>[5]\AgdaSymbol{\{}\AgdaBound{Δ} \AgdaInductiveConstructor{,} \AgdaBound{\_≈h₁\_} \AgdaInductiveConstructor{,} \AgdaBound{refl₁} \AgdaInductiveConstructor{,} \AgdaBound{sym₁} \AgdaInductiveConstructor{,} \AgdaBound{trans₁}\AgdaSymbol{\}}\<%
\\
\>[0]\AgdaIndent{5}{}\<[5]%
\>[5]\AgdaSymbol{(}\AgdaInductiveConstructor{fn:} \AgdaBound{fn} \AgdaInductiveConstructor{resp:} \AgdaBound{fresp}\AgdaSymbol{)} \AgdaSymbol{(}\AgdaInductiveConstructor{fn:} \AgdaBound{gn} \AgdaInductiveConstructor{resp:} \AgdaBound{gresp}\AgdaSymbol{)} \<[47]%
\>[47]\<%
\\
\>[0]\AgdaIndent{2}{}\<[2]%
\>[2]\AgdaSymbol{=} \AgdaKeyword{record} \<[11]%
\>[11]\<%
\\
\>[0]\AgdaIndent{11}{}\<[11]%
\>[11]\AgdaSymbol{\{} \AgdaField{prf} \AgdaSymbol{=} \AgdaSymbol{(}\AgdaBound{g} \AgdaSymbol{:} \AgdaBound{Γ}\AgdaSymbol{)} \AgdaSymbol{→} \AgdaFunction{<} \AgdaBound{fn} \AgdaBound{g} \AgdaBound{≈h₁} \AgdaBound{gn} \AgdaBound{g} \AgdaFunction{>}\<%
\\
\>[0]\AgdaIndent{11}{}\<[11]%
\>[11]\AgdaSymbol{;} \AgdaField{Uni} \AgdaSymbol{=} \AgdaBound{ext} \AgdaSymbol{(λ} \AgdaBound{g} \AgdaSymbol{→} \AgdaFunction{Uni} \AgdaSymbol{(}\AgdaBound{fn} \AgdaBound{g} \AgdaBound{≈h₁} \AgdaBound{gn} \AgdaBound{g}\AgdaSymbol{))}\<%
\\
\>[0]\AgdaIndent{11}{}\<[11]%
\>[11]\AgdaSymbol{\}}\<%
\\
%
\\
\>\<\end{code}

Category of Setoids

\begin{code}\>\<%
\\
%
\\
%
\\
\>\AgdaFunction{Std} \AgdaSymbol{:} \AgdaRecord{Category}\<%
\\
\>\AgdaFunction{Std} \AgdaSymbol{=} \AgdaInductiveConstructor{CatC} \AgdaRecord{HSetoid} \AgdaRecord{\_⇉\_} \AgdaSymbol{(λ} \AgdaBound{\_} \AgdaSymbol{→} \AgdaFunction{idCH}\AgdaSymbol{)} \AgdaSymbol{(λ} \AgdaBound{\_} \AgdaBound{\_} \AgdaSymbol{→} \AgdaFunction{\_∘c\_}\AgdaSymbol{)} \<[51]%
\>[51]\<%
\\
\>[0]\AgdaIndent{6}{}\<[6]%
\>[6]\AgdaSymbol{(}\AgdaInductiveConstructor{IsCatC} \AgdaSymbol{(λ} \AgdaBound{α} \AgdaBound{β} \AgdaBound{f} \AgdaSymbol{→} \AgdaInductiveConstructor{PE.refl}\AgdaSymbol{)} \<[34]%
\>[34]\<%
\\
\>[0]\AgdaIndent{6}{}\<[6]%
\>[6]\AgdaSymbol{(λ} \AgdaBound{α} \AgdaBound{β} \AgdaBound{f} \AgdaSymbol{→} \AgdaInductiveConstructor{PE.refl}\AgdaSymbol{)} \<[26]%
\>[26]\<%
\\
\>[0]\AgdaIndent{6}{}\<[6]%
\>[6]\AgdaSymbol{(λ} \AgdaBound{α} \AgdaBound{δ} \AgdaBound{f} \AgdaBound{g} \AgdaBound{h} \AgdaSymbol{→} \AgdaInductiveConstructor{PE.refl}\AgdaSymbol{))}\<%
\\
%
\\
\>\<\end{code}

Terminal object

\begin{code}\>\<%
\\
%
\\
\>\AgdaFunction{●} \AgdaSymbol{:} \AgdaRecord{HSetoid}\<%
\\
\>\AgdaFunction{●} \<[4]%
\>[4]\AgdaSymbol{=} \AgdaKeyword{record} \AgdaSymbol{\{}\<%
\\
\>[0]\AgdaIndent{6}{}\<[6]%
\>[6]\AgdaField{Carrier} \AgdaSymbol{=} \AgdaRecord{⊤}\AgdaSymbol{;}\<%
\\
\>[0]\AgdaIndent{6}{}\<[6]%
\>[6]\AgdaField{\_≈h\_} \<[14]%
\>[14]\AgdaSymbol{=} \AgdaSymbol{λ} \AgdaBound{\_} \AgdaBound{\_} \AgdaSymbol{→} \AgdaFunction{⊤'}\AgdaSymbol{;}\<%
\\
\>[0]\AgdaIndent{6}{}\<[6]%
\>[6]\AgdaField{refl} \<[14]%
\>[14]\AgdaSymbol{=} \AgdaInductiveConstructor{tt}\AgdaSymbol{;}\<%
\\
\>[0]\AgdaIndent{6}{}\<[6]%
\>[6]\AgdaField{sym} \<[14]%
\>[14]\AgdaSymbol{=} \AgdaSymbol{λ} \AgdaBound{\_} \AgdaSymbol{→} \AgdaInductiveConstructor{tt}\AgdaSymbol{;}\<%
\\
\>[0]\AgdaIndent{6}{}\<[6]%
\>[6]\AgdaField{trans} \<[14]%
\>[14]\AgdaSymbol{=} \AgdaSymbol{λ} \AgdaBound{\_} \AgdaBound{\_} \AgdaSymbol{→} \AgdaInductiveConstructor{tt} \AgdaSymbol{\}}\<%
\\
%
\\
%
\\
\>\AgdaFunction{⋆} \AgdaSymbol{:} \AgdaSymbol{\{}\AgdaBound{Δ} \AgdaSymbol{:} \AgdaRecord{HSetoid}\AgdaSymbol{\}} \AgdaSymbol{→} \AgdaBound{Δ} \AgdaRecord{⇉} \AgdaFunction{●}\<%
\\
\>\AgdaFunction{⋆} \AgdaSymbol{=} \AgdaKeyword{record} \<[11]%
\>[11]\<%
\\
\>[0]\AgdaIndent{6}{}\<[6]%
\>[6]\AgdaSymbol{\{} \AgdaField{fn} \AgdaSymbol{=} \AgdaSymbol{λ} \AgdaBound{\_} \AgdaSymbol{→} \AgdaInductiveConstructor{tt}\<%
\\
\>[0]\AgdaIndent{6}{}\<[6]%
\>[6]\AgdaSymbol{;} \AgdaField{resp} \AgdaSymbol{=} \AgdaSymbol{λ} \AgdaBound{\_} \AgdaSymbol{→} \AgdaInductiveConstructor{tt} \AgdaSymbol{\}}\<%
\\
%
\\
\>\AgdaFunction{uniqueHom} \AgdaSymbol{:} \AgdaSymbol{∀} \AgdaSymbol{(}\AgdaBound{Δ} \AgdaSymbol{:} \AgdaRecord{HSetoid}\AgdaSymbol{)} \AgdaSymbol{→} \AgdaSymbol{(}\AgdaBound{f} \AgdaSymbol{:} \AgdaBound{Δ} \AgdaRecord{⇉} \AgdaFunction{●}\AgdaSymbol{)} \AgdaSymbol{→} \AgdaBound{f} \AgdaDatatype{≡} \AgdaFunction{⋆}\<%
\\
\>\AgdaFunction{uniqueHom} \AgdaBound{Δ} \AgdaBound{f} \AgdaSymbol{=} \AgdaInductiveConstructor{PE.refl}\<%
\\
%
\\
\>\<\end{code}


\section{categories with families of setoids}


A Category with families consists of a base category and a functor
\cite{clairambault2005categories}. We firstly define the categories with
families of sets in Agda  as a guidance for the one for setoids. We
would present the setoid one here since it is relevant.

\AgdaHide{
\begin{code}\>\<%
\\
%
\\
\>\AgdaKeyword{module} \AgdaModule{CwF-setoid} \AgdaKeyword{where}\<%
\\
%
\\
\>\AgdaKeyword{open} \AgdaKeyword{import} \AgdaModule{CategoryofSetoid}\<%
\\
\>\AgdaKeyword{open} \AgdaKeyword{import} \AgdaModule{Data.Product}\<%
\\
\>\AgdaKeyword{open} \AgdaKeyword{import} \AgdaModule{Function}\<%
\\
\>\AgdaKeyword{open} \AgdaKeyword{import} \AgdaModule{Data.Nat}\<%
\\
\>\AgdaKeyword{open} \AgdaKeyword{import} \AgdaModule{Data.Empty}\<%
\\
\>\AgdaKeyword{open} \AgdaKeyword{import} \AgdaModule{Data.Unit}\<%
\\
%
\\
\>\<\end{code}
}

\section{Categories with families}

Context are interpreted as setoids

\begin{code}\>\<%
\\
\>\AgdaFunction{Con} \AgdaSymbol{=} \AgdaRecord{Setoid}\<%
\\
\>\<\end{code}

Semantic Types

\begin{code}\>\<%
\\
\>\AgdaKeyword{record} \AgdaRecord{Ty} \AgdaSymbol{(}\AgdaBound{Γ} \AgdaSymbol{:} \AgdaRecord{Setoid}\AgdaSymbol{)} \AgdaSymbol{:} \AgdaPrimitiveType{Set₁} \AgdaKeyword{where}\<%
\\
\>[0]\AgdaIndent{2}{}\<[2]%
\>[2]\AgdaKeyword{field}\<%
\\
\>[2]\AgdaIndent{4}{}\<[4]%
\>[4]\AgdaField{fm} \<[11]%
\>[11]\AgdaSymbol{:} \AgdaFunction{∣} \AgdaBound{Γ} \AgdaFunction{∣} \AgdaSymbol{→} \AgdaRecord{Setoid}\<%
\\
\>[2]\AgdaIndent{4}{}\<[4]%
\>[4]\AgdaField{substT} \AgdaSymbol{:} \AgdaSymbol{\{}\AgdaBound{x} \AgdaBound{y} \AgdaSymbol{:} \AgdaFunction{∣} \AgdaBound{Γ} \AgdaFunction{∣}\AgdaSymbol{\}} \AgdaSymbol{→} \<[29]%
\>[29]\<%
\\
\>[4]\AgdaIndent{13}{}\<[13]%
\>[13]\AgdaSymbol{.(}\AgdaFunction{[} \AgdaBound{Γ} \AgdaFunction{]} \AgdaBound{x} \AgdaFunction{≈} \AgdaBound{y}\AgdaSymbol{)} \AgdaSymbol{→}\<%
\\
\>[4]\AgdaIndent{13}{}\<[13]%
\>[13]\AgdaFunction{∣} \AgdaBound{fm} \AgdaBound{x} \AgdaFunction{∣} \AgdaSymbol{→} \<[24]%
\>[24]\<%
\\
\>[4]\AgdaIndent{13}{}\<[13]%
\>[13]\AgdaFunction{∣} \AgdaBound{fm} \AgdaBound{y} \AgdaFunction{∣}\<%
\\
\>[0]\AgdaIndent{4}{}\<[4]%
\>[4]\AgdaSymbol{.}\AgdaField{subst*} \AgdaSymbol{:} \AgdaSymbol{∀\{}\AgdaBound{x} \AgdaBound{y} \AgdaSymbol{:} \AgdaFunction{∣} \AgdaBound{Γ} \AgdaFunction{∣}\AgdaSymbol{\}}\<%
\\
\>[0]\AgdaIndent{13}{}\<[13]%
\>[13]\AgdaSymbol{(}\AgdaBound{p} \AgdaSymbol{:} \AgdaSymbol{(}\AgdaFunction{[} \AgdaBound{Γ} \AgdaFunction{]} \AgdaBound{x} \AgdaFunction{≈} \AgdaBound{y}\AgdaSymbol{))}\<%
\\
\>[0]\AgdaIndent{13}{}\<[13]%
\>[13]\AgdaSymbol{\{}\AgdaBound{a} \AgdaBound{b} \AgdaSymbol{:} \AgdaFunction{∣} \AgdaBound{fm} \AgdaBound{x} \AgdaFunction{∣}\AgdaSymbol{\}} \AgdaSymbol{→}\<%
\\
\>[0]\AgdaIndent{13}{}\<[13]%
\>[13]\AgdaSymbol{.(}\AgdaFunction{[} \AgdaBound{fm} \AgdaBound{x} \AgdaFunction{]} \AgdaBound{a} \AgdaFunction{≈} \AgdaBound{b}\AgdaSymbol{)} \AgdaSymbol{→}\<%
\\
\>[0]\AgdaIndent{13}{}\<[13]%
\>[13]\AgdaSymbol{(}\AgdaFunction{[} \AgdaBound{fm} \AgdaBound{y} \AgdaFunction{]} \AgdaBound{substT} \AgdaBound{p} \AgdaBound{a} \AgdaFunction{≈} \AgdaBound{substT} \AgdaBound{p} \AgdaBound{b}\AgdaSymbol{)}\<%
\\
%
\\
\>[0]\AgdaIndent{4}{}\<[4]%
\>[4]\AgdaSymbol{.}\AgdaField{refl*} \<[12]%
\>[12]\AgdaSymbol{:} \AgdaSymbol{∀\{}\AgdaBound{x} \AgdaSymbol{:} \AgdaFunction{∣} \AgdaBound{Γ} \AgdaFunction{∣}\AgdaSymbol{\}\{}\AgdaBound{a} \AgdaSymbol{:} \AgdaFunction{∣} \AgdaBound{fm} \AgdaBound{x} \AgdaFunction{∣}\AgdaSymbol{\}} \AgdaSymbol{→} \<[43]%
\>[43]\<%
\\
\>[0]\AgdaIndent{13}{}\<[13]%
\>[13]\AgdaFunction{[} \AgdaBound{fm} \AgdaBound{x} \AgdaFunction{]} \AgdaBound{substT} \AgdaSymbol{(}\AgdaFunction{[} \AgdaBound{Γ} \AgdaFunction{]refl}\AgdaSymbol{)} \AgdaBound{a} \AgdaFunction{≈} \AgdaBound{a}\<%
\\
\>[0]\AgdaIndent{4}{}\<[4]%
\>[4]\AgdaSymbol{.}\AgdaField{trans*} \AgdaSymbol{:} \AgdaSymbol{∀\{}\AgdaBound{x} \AgdaBound{y} \AgdaBound{z} \AgdaSymbol{:} \AgdaFunction{∣} \AgdaBound{Γ} \AgdaFunction{∣}\AgdaSymbol{\}}\<%
\\
\>[0]\AgdaIndent{13}{}\<[13]%
\>[13]\AgdaSymbol{\{}\AgdaBound{p} \AgdaSymbol{:} \AgdaFunction{[} \AgdaBound{Γ} \AgdaFunction{]} \AgdaBound{x} \AgdaFunction{≈} \AgdaBound{y}\AgdaSymbol{\}}\<%
\\
\>[0]\AgdaIndent{13}{}\<[13]%
\>[13]\AgdaSymbol{\{}\AgdaBound{q} \AgdaSymbol{:} \AgdaFunction{[} \AgdaBound{Γ} \AgdaFunction{]} \AgdaBound{y} \AgdaFunction{≈} \AgdaBound{z}\AgdaSymbol{\}}\<%
\\
\>[0]\AgdaIndent{13}{}\<[13]%
\>[13]\AgdaSymbol{(}\AgdaBound{a} \AgdaSymbol{:} \AgdaFunction{∣} \AgdaBound{fm} \AgdaBound{x} \AgdaFunction{∣}\AgdaSymbol{)} \AgdaSymbol{→} \<[30]%
\>[30]\<%
\\
\>[0]\AgdaIndent{13}{}\<[13]%
\>[13]\AgdaFunction{[} \AgdaBound{fm} \AgdaBound{z} \AgdaFunction{]} \AgdaBound{substT} \AgdaBound{q} \AgdaSymbol{(}\AgdaBound{substT} \AgdaBound{p} \AgdaBound{a}\AgdaSymbol{)} \<[44]%
\>[44]\<%
\\
\>[0]\AgdaIndent{12}{}\<[12]%
\>[12]\AgdaFunction{≈} \AgdaBound{substT} \AgdaSymbol{(}\AgdaFunction{[} \AgdaBound{Γ} \AgdaFunction{]trans} \AgdaBound{p} \AgdaBound{q}\AgdaSymbol{)} \AgdaBound{a}\<%
\\
%
\\
\>[0]\AgdaIndent{2}{}\<[2]%
\>[2]\AgdaSymbol{.}\AgdaFunction{tr*} \AgdaSymbol{:} \AgdaSymbol{∀\{}\AgdaBound{x} \AgdaBound{y} \AgdaSymbol{:} \AgdaFunction{∣} \AgdaBound{Γ} \AgdaFunction{∣}\AgdaSymbol{\}}\<%
\\
\>[2]\AgdaIndent{10}{}\<[10]%
\>[10]\AgdaSymbol{\{}\AgdaBound{p} \AgdaSymbol{:} \AgdaFunction{[} \AgdaBound{Γ} \AgdaFunction{]} \AgdaBound{y} \AgdaFunction{≈} \AgdaBound{x}\AgdaSymbol{\}}\<%
\\
\>[2]\AgdaIndent{10}{}\<[10]%
\>[10]\AgdaSymbol{\{}\AgdaBound{q} \AgdaSymbol{:} \AgdaFunction{[} \AgdaBound{Γ} \AgdaFunction{]} \AgdaBound{x} \AgdaFunction{≈} \AgdaBound{y}\AgdaSymbol{\}}\<%
\\
\>[2]\AgdaIndent{10}{}\<[10]%
\>[10]\AgdaSymbol{\{}\AgdaBound{a} \AgdaSymbol{:} \AgdaFunction{∣} \AgdaFunction{fm} \AgdaBound{x} \AgdaFunction{∣}\AgdaSymbol{\}} \AgdaSymbol{→} \<[27]%
\>[27]\<%
\\
\>[2]\AgdaIndent{10}{}\<[10]%
\>[10]\AgdaFunction{[} \AgdaFunction{fm} \AgdaBound{x} \AgdaFunction{]} \AgdaFunction{substT} \AgdaBound{p} \AgdaSymbol{(}\AgdaFunction{substT} \AgdaBound{q} \AgdaBound{a}\AgdaSymbol{)} \AgdaFunction{≈} \AgdaBound{a}\<%
\\
\>[0]\AgdaIndent{2}{}\<[2]%
\>[2]\AgdaFunction{tr*} \AgdaSymbol{=} \AgdaFunction{[} \AgdaFunction{fm} \AgdaSymbol{\_} \AgdaFunction{]trans} \AgdaSymbol{(}\AgdaFunction{trans*} \AgdaSymbol{\_)} \AgdaFunction{refl*}\<%
\\
%
\\
\>[0]\AgdaIndent{2}{}\<[2]%
\>[2]\AgdaFunction{substT-inv} \AgdaSymbol{:} \AgdaSymbol{\{}\AgdaBound{x} \AgdaBound{y} \AgdaSymbol{:} \AgdaFunction{∣} \AgdaBound{Γ} \AgdaFunction{∣}\AgdaSymbol{\}} \AgdaSymbol{→} \<[31]%
\>[31]\<%
\\
\>[2]\AgdaIndent{13}{}\<[13]%
\>[13]\AgdaSymbol{.(}\AgdaFunction{[} \AgdaBound{Γ} \AgdaFunction{]} \AgdaBound{x} \AgdaFunction{≈} \AgdaBound{y}\AgdaSymbol{)} \AgdaSymbol{→}\<%
\\
\>[2]\AgdaIndent{13}{}\<[13]%
\>[13]\AgdaFunction{∣} \AgdaFunction{fm} \AgdaBound{y} \AgdaFunction{∣} \AgdaSymbol{→} \<[24]%
\>[24]\<%
\\
\>[2]\AgdaIndent{13}{}\<[13]%
\>[13]\AgdaFunction{∣} \AgdaFunction{fm} \AgdaBound{x} \AgdaFunction{∣}\<%
\\
\>[0]\AgdaIndent{2}{}\<[2]%
\>[2]\AgdaFunction{substT-inv} \AgdaBound{p} \AgdaBound{y} \AgdaSymbol{=} \AgdaFunction{substT} \AgdaSymbol{(}\AgdaFunction{[} \AgdaBound{Γ} \AgdaFunction{]sym} \AgdaBound{p}\AgdaSymbol{)} \AgdaBound{y}\<%
\\
\>\<\end{code}

\AgdaHide{
\begin{code}\>\<%
\\
\>[0]\AgdaIndent{2}{}\<[2]%
\>[2]\AgdaFunction{subst-mir1} \AgdaSymbol{:} \AgdaSymbol{∀\{}\AgdaBound{x} \AgdaBound{y} \AgdaSymbol{:} \AgdaFunction{∣} \AgdaBound{Γ} \AgdaFunction{∣}\AgdaSymbol{\}\{}\AgdaBound{P} \AgdaSymbol{:} \AgdaPrimitiveType{Set}\AgdaSymbol{\}}\<%
\\
\>[2]\AgdaIndent{14}{}\<[14]%
\>[14]\AgdaSymbol{\{}\AgdaBound{p} \AgdaSymbol{:} \AgdaFunction{[} \AgdaBound{Γ} \AgdaFunction{]} \AgdaBound{x} \AgdaFunction{≈} \AgdaBound{y}\AgdaSymbol{\}\{}\AgdaBound{q} \AgdaSymbol{:} \AgdaFunction{[} \AgdaBound{Γ} \AgdaFunction{]} \AgdaBound{y} \AgdaFunction{≈} \AgdaBound{x}\AgdaSymbol{\}}\<%
\\
\>[2]\AgdaIndent{14}{}\<[14]%
\>[14]\AgdaSymbol{\{}\AgdaBound{a} \AgdaSymbol{:} \AgdaFunction{∣} \AgdaFunction{fm} \AgdaBound{x} \AgdaFunction{∣}\AgdaSymbol{\}\{}\AgdaBound{b} \AgdaSymbol{:} \AgdaFunction{∣} \AgdaFunction{fm} \AgdaBound{y} \AgdaFunction{∣}\AgdaSymbol{\}} \AgdaSymbol{→} \<[45]%
\>[45]\<%
\\
\>[2]\AgdaIndent{14}{}\<[14]%
\>[14]\AgdaSymbol{(.(}\AgdaFunction{[} \AgdaFunction{fm} \AgdaBound{y} \AgdaFunction{]} \AgdaFunction{substT} \AgdaBound{p} \AgdaBound{a} \AgdaFunction{≈} \AgdaBound{b}\AgdaSymbol{)} \AgdaSymbol{→} \AgdaBound{P}\AgdaSymbol{)} \<[47]%
\>[47]\<%
\\
\>[-1]\AgdaIndent{13}{}\<[13]%
\>[13]\AgdaSymbol{→} \AgdaSymbol{(.(}\AgdaFunction{[} \AgdaFunction{fm} \AgdaBound{x} \AgdaFunction{]} \AgdaBound{a} \AgdaFunction{≈} \AgdaFunction{substT} \AgdaBound{q} \AgdaBound{b}\AgdaSymbol{)} \AgdaSymbol{→} \AgdaBound{P}\AgdaSymbol{)}\<%
\\
\>[0]\AgdaIndent{2}{}\<[2]%
\>[2]\AgdaFunction{subst-mir1} \AgdaBound{eq} \AgdaBound{p} \AgdaSymbol{=} \AgdaBound{eq} \AgdaSymbol{(}\AgdaFunction{[} \AgdaFunction{fm} \AgdaSymbol{\_} \AgdaFunction{]trans} \AgdaSymbol{(}\AgdaFunction{subst*} \AgdaSymbol{\_} \AgdaBound{p}\AgdaSymbol{)} \AgdaFunction{tr*}\AgdaSymbol{)}\<%
\\
%
\\
\>\AgdaKeyword{open} \AgdaModule{Ty} \AgdaKeyword{public} \<[15]%
\>[15]\<%
\\
\>[0]\AgdaIndent{2}{}\<[2]%
\>[2]\AgdaKeyword{renaming} \AgdaSymbol{(}substT \AgdaSymbol{to} [\_]subst\AgdaSymbol{;} substT-inv \AgdaSymbol{to} [\_]subst-inv
           \AgdaSymbol{;} subst* \AgdaSymbol{to} [\_]subst*\AgdaSymbol{;} fm \AgdaSymbol{to} [\_]fm \AgdaSymbol{;}
            refl* \AgdaSymbol{to} [\_]refl* \AgdaSymbol{;} trans* \AgdaSymbol{to} [\_]trans*\AgdaSymbol{;} tr* \AgdaSymbol{to} [\_]tr*\AgdaSymbol{)}\<%
\\
%
\\
\>\<\end{code}
}

Type substitution

\begin{code}\>\<%
\\
\>\AgdaFunction{\_[\_]T} \AgdaSymbol{:} \AgdaSymbol{∀} \AgdaSymbol{\{}\AgdaBound{Γ} \AgdaBound{Δ} \AgdaSymbol{:} \AgdaRecord{Setoid}\AgdaSymbol{\}} \AgdaSymbol{→} \AgdaRecord{Ty} \AgdaBound{Δ} \AgdaSymbol{→} \AgdaBound{Γ} \AgdaRecord{⇉} \AgdaBound{Δ} \AgdaSymbol{→} \AgdaRecord{Ty} \AgdaBound{Γ}\<%
\\
\>\AgdaFunction{\_[\_]T} \AgdaSymbol{\{}\AgdaBound{Γ}\AgdaSymbol{\}} \AgdaSymbol{\{}\AgdaBound{Δ}\AgdaSymbol{\}} \AgdaBound{A} \AgdaBound{f}\<%
\\
\>[2]\AgdaIndent{5}{}\<[5]%
\>[5]\AgdaSymbol{=} \AgdaKeyword{record}\<%
\\
\>[2]\AgdaIndent{5}{}\<[5]%
\>[5]\AgdaSymbol{\{} \AgdaField{fm} \<[14]%
\>[14]\AgdaSymbol{=} \AgdaSymbol{λ} \AgdaBound{x} \AgdaSymbol{→} \AgdaFunction{fm} \AgdaSymbol{(}\AgdaFunction{fn} \AgdaBound{x}\AgdaSymbol{)}\<%
\\
\>[2]\AgdaIndent{5}{}\<[5]%
\>[5]\AgdaSymbol{;} \AgdaField{substT} \AgdaSymbol{=} \AgdaSymbol{λ} \AgdaBound{p} \AgdaSymbol{→} \AgdaFunction{substT} \AgdaSymbol{\_}\<%
\\
\>[2]\AgdaIndent{5}{}\<[5]%
\>[5]\AgdaSymbol{;} \AgdaField{subst*} \AgdaSymbol{=} \AgdaSymbol{λ} \AgdaBound{p} \AgdaSymbol{→} \AgdaFunction{subst*} \AgdaSymbol{(}\AgdaFunction{resp} \AgdaBound{p}\AgdaSymbol{)}\<%
\\
\>[2]\AgdaIndent{5}{}\<[5]%
\>[5]\AgdaSymbol{;} \AgdaField{refl*} \<[14]%
\>[14]\AgdaSymbol{=} \AgdaFunction{refl*}\<%
\\
\>[2]\AgdaIndent{5}{}\<[5]%
\>[5]\AgdaSymbol{;} \AgdaField{trans*} \AgdaSymbol{=} \AgdaFunction{trans*}\<%
\\
\>[2]\AgdaIndent{5}{}\<[5]%
\>[5]\AgdaSymbol{\}}\<%
\\
\>[2]\AgdaIndent{5}{}\<[5]%
\>[5]\AgdaKeyword{where} \<[11]%
\>[11]\<%
\\
\>[5]\AgdaIndent{7}{}\<[7]%
\>[7]\AgdaKeyword{open} \AgdaModule{Ty} \AgdaBound{A}\<%
\\
\>[5]\AgdaIndent{7}{}\<[7]%
\>[7]\AgdaKeyword{open} \AgdaModule{\_⇉\_} \AgdaBound{f}\<%
\\
%
\\
%
\\
\>\<\end{code}


Semantic Terms

\begin{code}\>\<%
\\
\>\AgdaKeyword{record} \AgdaRecord{Tm} \AgdaSymbol{\{}\AgdaBound{Γ} \AgdaSymbol{:} \AgdaFunction{Con}\AgdaSymbol{\}(}\AgdaBound{A} \AgdaSymbol{:} \AgdaRecord{Ty} \AgdaBound{Γ}\AgdaSymbol{)} \AgdaSymbol{:} \AgdaPrimitiveType{Set} \AgdaKeyword{where}\<%
\\
\>[0]\AgdaIndent{2}{}\<[2]%
\>[2]\AgdaKeyword{constructor} \AgdaInductiveConstructor{tm:\_resp:\_}\<%
\\
\>[0]\AgdaIndent{2}{}\<[2]%
\>[2]\AgdaKeyword{field}\<%
\\
\>[2]\AgdaIndent{4}{}\<[4]%
\>[4]\AgdaField{tm} \<[10]%
\>[10]\AgdaSymbol{:} \AgdaSymbol{(}\AgdaBound{x} \AgdaSymbol{:} \AgdaFunction{∣} \AgdaBound{Γ} \AgdaFunction{∣}\AgdaSymbol{)} \AgdaSymbol{→} \AgdaFunction{∣} \AgdaFunction{[} \AgdaBound{A} \AgdaFunction{]fm} \AgdaBound{x} \AgdaFunction{∣}\<%
\\
\>[2]\AgdaIndent{4}{}\<[4]%
\>[4]\AgdaSymbol{.}\AgdaField{respt} \AgdaSymbol{:} \AgdaSymbol{∀} \AgdaSymbol{\{}\AgdaBound{x} \AgdaBound{y} \AgdaSymbol{:} \AgdaFunction{∣} \AgdaBound{Γ} \AgdaFunction{∣}\AgdaSymbol{\}} \AgdaSymbol{→} \<[31]%
\>[31]\<%
\\
\>[4]\AgdaIndent{10}{}\<[10]%
\>[10]\AgdaSymbol{(}\AgdaBound{p} \AgdaSymbol{:} \AgdaFunction{[} \AgdaBound{Γ} \AgdaFunction{]} \AgdaBound{x} \AgdaFunction{≈} \AgdaBound{y}\AgdaSymbol{)} \AgdaSymbol{→} \<[30]%
\>[30]\<%
\\
\>[4]\AgdaIndent{10}{}\<[10]%
\>[10]\AgdaFunction{[} \AgdaFunction{[} \AgdaBound{A} \AgdaFunction{]fm} \AgdaBound{y} \AgdaFunction{]} \AgdaFunction{[} \AgdaBound{A} \AgdaFunction{]subst} \AgdaBound{p} \AgdaSymbol{(}\AgdaBound{tm} \AgdaBound{x}\AgdaSymbol{)} \AgdaFunction{≈} \AgdaBound{tm} \AgdaBound{y}\<%
\\
%
\\
\>\AgdaKeyword{open} \AgdaModule{Tm} \AgdaKeyword{public} \AgdaKeyword{renaming} \AgdaSymbol{(}tm \AgdaSymbol{to} [\_]tm \AgdaSymbol{;} respt \AgdaSymbol{to} [\_]respt\AgdaSymbol{)}\<%
\\
\>\<\end{code}

Term substitution

\begin{code}\>\<%
\\
\>\AgdaFunction{\_[\_]m} \AgdaSymbol{:} \AgdaSymbol{∀} \AgdaSymbol{\{}\AgdaBound{Γ} \AgdaBound{Δ} \AgdaSymbol{:} \AgdaFunction{Con}\AgdaSymbol{\}\{}\AgdaBound{A} \AgdaSymbol{:} \AgdaRecord{Ty} \AgdaBound{Δ}\AgdaSymbol{\}} \AgdaSymbol{→} \AgdaRecord{Tm} \AgdaBound{A} \<[39]%
\>[39]\<%
\\
\>[0]\AgdaIndent{6}{}\<[6]%
\>[6]\AgdaSymbol{→} \AgdaSymbol{(}\AgdaBound{f} \AgdaSymbol{:} \AgdaBound{Γ} \AgdaRecord{⇉} \AgdaBound{Δ}\AgdaSymbol{)} \AgdaSymbol{→} \AgdaRecord{Tm} \AgdaSymbol{(}\AgdaBound{A} \AgdaFunction{[} \AgdaBound{f} \AgdaFunction{]T}\AgdaSymbol{)}\<%
\\
\>\AgdaFunction{\_[\_]m} \AgdaBound{t} \AgdaBound{f} \AgdaSymbol{=} \AgdaKeyword{record} \<[19]%
\>[19]\<%
\\
\>[0]\AgdaIndent{10}{}\<[10]%
\>[10]\AgdaSymbol{\{} \AgdaField{tm} \AgdaSymbol{=} \AgdaFunction{[} \AgdaBound{t} \AgdaFunction{]tm} \AgdaFunction{∘} \AgdaFunction{[} \AgdaBound{f} \AgdaFunction{]fn}\<%
\\
\>[0]\AgdaIndent{10}{}\<[10]%
\>[10]\AgdaSymbol{;} \AgdaField{respt} \AgdaSymbol{=} \AgdaFunction{[} \AgdaBound{t} \AgdaFunction{]respt} \AgdaFunction{∘} \AgdaFunction{[} \AgdaBound{f} \AgdaFunction{]resp} \<[43]%
\>[43]\<%
\\
\>[0]\AgdaIndent{10}{}\<[10]%
\>[10]\AgdaSymbol{\}}\<%
\\
\>\<\end{code}

Context comprehension

\begin{code}\>\<%
\\
\>\AgdaFunction{\_\&\_} \AgdaSymbol{:} \AgdaSymbol{(}\AgdaBound{Γ} \AgdaSymbol{:} \AgdaRecord{Setoid}\AgdaSymbol{)} \AgdaSymbol{→} \AgdaRecord{Ty} \AgdaBound{Γ} \AgdaSymbol{→} \AgdaRecord{Setoid}\<%
\\
\>\AgdaBound{Γ} \AgdaFunction{\&} \AgdaBound{A} \AgdaSymbol{=} \<[8]%
\>[8]\<%
\\
\>[0]\AgdaIndent{2}{}\<[2]%
\>[2]\AgdaKeyword{record} \AgdaSymbol{\{} \AgdaField{Carrier} \AgdaSymbol{=} \AgdaRecord{Σ[} \AgdaBound{x} \AgdaRecord{∶} \AgdaFunction{∣} \AgdaBound{Γ} \AgdaFunction{∣} \AgdaRecord{]} \AgdaFunction{∣} \AgdaFunction{fm} \AgdaBound{x} \AgdaFunction{∣} \<[45]%
\>[45]\<%
\\
\>[0]\AgdaIndent{9}{}\<[9]%
\>[9]\AgdaSymbol{;} \AgdaField{\_≈\_} \AgdaSymbol{=} \AgdaSymbol{λ\{(}\AgdaBound{x} \AgdaInductiveConstructor{,} \AgdaBound{a}\AgdaSymbol{)} \AgdaSymbol{(}\AgdaBound{y} \AgdaInductiveConstructor{,} \AgdaBound{b}\AgdaSymbol{)} \AgdaSymbol{→} \<[37]%
\>[37]\<%
\\
\>[9]\AgdaIndent{13}{}\<[13]%
\>[13]\AgdaRecord{Σ[} \AgdaBound{p} \AgdaRecord{∶} \AgdaBound{x} \AgdaFunction{≈} \AgdaBound{y} \AgdaRecord{]} \AgdaFunction{[} \AgdaFunction{fm} \AgdaBound{y} \AgdaFunction{]} \AgdaSymbol{(}\AgdaFunction{substT} \AgdaBound{p} \AgdaBound{a}\AgdaSymbol{)} \AgdaFunction{≈} \AgdaBound{b} \AgdaSymbol{\}}\<%
\\
\>[0]\AgdaIndent{9}{}\<[9]%
\>[9]\AgdaSymbol{;} \AgdaField{refl} \AgdaSymbol{=} \AgdaFunction{refl} \AgdaInductiveConstructor{,} \AgdaFunction{refl*}\<%
\\
\>[0]\AgdaIndent{9}{}\<[9]%
\>[9]\AgdaSymbol{;} \AgdaField{sym} \AgdaSymbol{=} \<[18]%
\>[18]\AgdaSymbol{λ} \AgdaSymbol{\{(}\AgdaBound{p} \AgdaInductiveConstructor{,} \AgdaBound{q}\AgdaSymbol{)} \AgdaSymbol{→} \AgdaSymbol{(}\AgdaFunction{sym} \AgdaBound{p}\AgdaSymbol{)} \AgdaInductiveConstructor{,} \<[41]%
\>[41]\<%
\\
\>[9]\AgdaIndent{13}{}\<[13]%
\>[13]\AgdaFunction{[} \AgdaFunction{fm} \AgdaSymbol{\_} \AgdaFunction{]trans} \AgdaSymbol{(}\AgdaFunction{subst*} \AgdaSymbol{\_} \AgdaSymbol{(}\AgdaFunction{[} \AgdaFunction{fm} \AgdaSymbol{\_} \AgdaFunction{]sym} \AgdaBound{q}\AgdaSymbol{))} \AgdaFunction{tr*} \AgdaSymbol{\}}\<%
\\
\>[0]\AgdaIndent{9}{}\<[9]%
\>[9]\AgdaSymbol{;} \AgdaField{trans} \AgdaSymbol{=} \AgdaSymbol{λ} \AgdaSymbol{\{(}\AgdaBound{p} \AgdaInductiveConstructor{,} \AgdaBound{q}\AgdaSymbol{)} \AgdaSymbol{(}\AgdaBound{m} \AgdaInductiveConstructor{,} \AgdaBound{n}\AgdaSymbol{)} \AgdaSymbol{→} \AgdaFunction{trans} \AgdaBound{p} \AgdaBound{m} \AgdaInductiveConstructor{,}\<%
\\
\>[0]\AgdaIndent{19}{}\<[19]%
\>[19]\AgdaFunction{[} \AgdaFunction{fm} \AgdaSymbol{\_} \AgdaFunction{]trans} \AgdaSymbol{(}\AgdaFunction{[} \AgdaFunction{fm} \AgdaSymbol{\_} \AgdaFunction{]trans} \<[48]%
\>[48]\<%
\\
\>[0]\AgdaIndent{19}{}\<[19]%
\>[19]\AgdaSymbol{(}\AgdaFunction{[} \AgdaFunction{fm} \AgdaSymbol{\_} \AgdaFunction{]sym} \AgdaSymbol{(}\AgdaFunction{trans*} \AgdaSymbol{\_))} \AgdaSymbol{(}\AgdaFunction{subst*} \AgdaSymbol{\_} \AgdaBound{q}\AgdaSymbol{))} \AgdaBound{n}\AgdaSymbol{\}}\<%
\\
\>[0]\AgdaIndent{9}{}\<[9]%
\>[9]\AgdaSymbol{\}}\<%
\\
\>[0]\AgdaIndent{4}{}\<[4]%
\>[4]\AgdaKeyword{where}\<%
\\
\>[4]\AgdaIndent{6}{}\<[6]%
\>[6]\AgdaKeyword{open} \AgdaModule{Setoid} \AgdaBound{Γ}\<%
\\
\>[4]\AgdaIndent{6}{}\<[6]%
\>[6]\AgdaKeyword{open} \AgdaModule{Ty} \AgdaBound{A}\<%
\\
%
\\
%
\\
\>\AgdaKeyword{infixl} \AgdaNumber{5} \_\&\_\<%
\\
%
\\
\>\AgdaFunction{fst\&} \AgdaSymbol{:} \AgdaSymbol{\{}\AgdaBound{Γ} \AgdaSymbol{:} \AgdaFunction{Con}\AgdaSymbol{\}\{}\AgdaBound{A} \AgdaSymbol{:} \AgdaRecord{Ty} \AgdaBound{Γ}\AgdaSymbol{\}} \AgdaSymbol{→} \AgdaBound{Γ} \AgdaFunction{\&} \AgdaBound{A} \AgdaRecord{⇉} \AgdaBound{Γ}\<%
\\
\>\AgdaFunction{fst\&} \AgdaSymbol{=} \AgdaKeyword{record} \<[14]%
\>[14]\<%
\\
\>[6]\AgdaIndent{9}{}\<[9]%
\>[9]\AgdaSymbol{\{} \AgdaField{fn} \AgdaSymbol{=} \AgdaFunction{proj₁}\<%
\\
\>[6]\AgdaIndent{9}{}\<[9]%
\>[9]\AgdaSymbol{;} \AgdaField{resp} \AgdaSymbol{=} \AgdaFunction{proj₁}\<%
\\
\>[6]\AgdaIndent{9}{}\<[9]%
\>[9]\AgdaSymbol{\}}\<%
\\
\>\<\end{code}

Pairing operation

\begin{code}\>\<%
\\
\>\AgdaFunction{\_,,\_} \AgdaSymbol{:} \AgdaSymbol{\{}\AgdaBound{Γ} \AgdaBound{Δ} \AgdaSymbol{:} \AgdaFunction{Con}\AgdaSymbol{\}\{}\AgdaBound{A} \AgdaSymbol{:} \AgdaRecord{Ty} \AgdaBound{Δ}\AgdaSymbol{\}(}\AgdaBound{f} \AgdaSymbol{:} \AgdaBound{Γ} \AgdaRecord{⇉} \AgdaBound{Δ}\AgdaSymbol{)} \<[40]%
\>[40]\<%
\\
\>[-2]\AgdaIndent{5}{}\<[5]%
\>[5]\AgdaSymbol{→} \AgdaSymbol{(}\AgdaRecord{Tm} \AgdaSymbol{(}\AgdaBound{A} \AgdaFunction{[} \AgdaBound{f} \AgdaFunction{]T}\AgdaSymbol{))} \AgdaSymbol{→} \AgdaBound{Γ} \AgdaRecord{⇉} \AgdaSymbol{(}\AgdaBound{Δ} \AgdaFunction{\&} \AgdaBound{A}\AgdaSymbol{)}\<%
\\
\>\AgdaBound{f} \AgdaFunction{,,} \AgdaBound{t} \AgdaSymbol{=} \AgdaKeyword{record} \<[16]%
\>[16]\<%
\\
\>[0]\AgdaIndent{9}{}\<[9]%
\>[9]\AgdaSymbol{\{} \AgdaField{fn} \AgdaSymbol{=} \AgdaFunction{<} \AgdaFunction{[} \AgdaBound{f} \AgdaFunction{]fn} \AgdaFunction{,} \AgdaFunction{[} \AgdaBound{t} \AgdaFunction{]tm} \AgdaFunction{>}\<%
\\
\>[0]\AgdaIndent{9}{}\<[9]%
\>[9]\AgdaSymbol{;} \AgdaField{resp} \AgdaSymbol{=} \AgdaFunction{<} \AgdaFunction{[} \AgdaBound{f} \AgdaFunction{]resp} \AgdaFunction{,} \AgdaFunction{[} \AgdaBound{t} \AgdaFunction{]respt} \AgdaFunction{>}\<%
\\
\>[0]\AgdaIndent{9}{}\<[9]%
\>[9]\AgdaSymbol{\}}\<%
\\
\>\<\end{code}

Projections

\begin{code}\>\<%
\\
\>\AgdaFunction{fst} \AgdaSymbol{:} \AgdaSymbol{\{}\AgdaBound{Γ} \AgdaBound{Δ} \AgdaSymbol{:} \AgdaFunction{Con}\AgdaSymbol{\}\{}\AgdaBound{A} \AgdaSymbol{:} \AgdaRecord{Ty} \AgdaBound{Δ}\AgdaSymbol{\}} \AgdaSymbol{→} \AgdaBound{Γ} \AgdaRecord{⇉} \AgdaSymbol{(}\AgdaBound{Δ} \AgdaFunction{\&} \AgdaBound{A}\AgdaSymbol{)} \AgdaSymbol{→} \AgdaBound{Γ} \AgdaRecord{⇉} \AgdaBound{Δ}\<%
\\
\>\AgdaFunction{fst} \AgdaBound{f} \AgdaSymbol{=} \AgdaKeyword{record} \<[15]%
\>[15]\<%
\\
\>[0]\AgdaIndent{8}{}\<[8]%
\>[8]\AgdaSymbol{\{} \AgdaField{fn} \AgdaSymbol{=} \AgdaFunction{proj₁} \AgdaFunction{∘} \AgdaFunction{[} \AgdaBound{f} \AgdaFunction{]fn}\<%
\\
\>[0]\AgdaIndent{8}{}\<[8]%
\>[8]\AgdaSymbol{;} \AgdaField{resp} \AgdaSymbol{=} \AgdaFunction{proj₁} \AgdaFunction{∘} \AgdaFunction{[} \AgdaBound{f} \AgdaFunction{]resp} \<[35]%
\>[35]\<%
\\
\>[0]\AgdaIndent{8}{}\<[8]%
\>[8]\AgdaSymbol{\}}\<%
\\
%
\\
\>\AgdaFunction{snd} \AgdaSymbol{:} \AgdaSymbol{\{}\AgdaBound{Γ} \AgdaBound{Δ} \AgdaSymbol{:} \AgdaFunction{Con}\AgdaSymbol{\}\{}\AgdaBound{A} \AgdaSymbol{:} \AgdaRecord{Ty} \AgdaBound{Δ}\AgdaSymbol{\}} \AgdaSymbol{→} \AgdaSymbol{(}\AgdaBound{f} \AgdaSymbol{:} \AgdaBound{Γ} \AgdaRecord{⇉} \AgdaSymbol{(}\AgdaBound{Δ} \AgdaFunction{\&} \AgdaBound{A}\AgdaSymbol{))}\<%
\\
\>[0]\AgdaIndent{4}{}\<[4]%
\>[4]\AgdaSymbol{→} \AgdaRecord{Tm} \AgdaSymbol{(}\AgdaBound{A} \AgdaFunction{[} \AgdaFunction{fst} \AgdaSymbol{\{}A \AgdaSymbol{=} \AgdaBound{A}\AgdaSymbol{\}} \AgdaBound{f} \AgdaFunction{]T}\AgdaSymbol{)}\<%
\\
\>\AgdaFunction{snd} \AgdaBound{f} \AgdaSymbol{=} \AgdaKeyword{record} \<[15]%
\>[15]\<%
\\
\>[0]\AgdaIndent{8}{}\<[8]%
\>[8]\AgdaSymbol{\{} \AgdaField{tm} \AgdaSymbol{=} \AgdaFunction{proj₂} \AgdaFunction{∘} \AgdaFunction{[} \AgdaBound{f} \AgdaFunction{]fn}\<%
\\
\>[0]\AgdaIndent{8}{}\<[8]%
\>[8]\AgdaSymbol{;} \AgdaField{respt} \AgdaSymbol{=} \AgdaFunction{proj₂} \AgdaFunction{∘} \AgdaFunction{[} \AgdaBound{f} \AgdaFunction{]resp} \<[36]%
\>[36]\<%
\\
\>[0]\AgdaIndent{8}{}\<[8]%
\>[8]\AgdaSymbol{\}}\<%
\\
%
\\
\>\AgdaFunction{\_\textasciicircum\_} \AgdaSymbol{:} \AgdaSymbol{\{}\AgdaBound{Γ} \AgdaBound{Δ} \AgdaSymbol{:} \AgdaFunction{Con}\AgdaSymbol{\}(}\AgdaBound{f} \AgdaSymbol{:} \AgdaBound{Γ} \AgdaRecord{⇉} \AgdaBound{Δ}\AgdaSymbol{)(}\AgdaBound{A} \AgdaSymbol{:} \AgdaRecord{Ty} \AgdaBound{Δ}\AgdaSymbol{)} \<[39]%
\>[39]\<%
\\
\>[0]\AgdaIndent{4}{}\<[4]%
\>[4]\AgdaSymbol{→} \AgdaBound{Γ} \AgdaFunction{\&} \AgdaBound{A} \AgdaFunction{[} \AgdaBound{f} \AgdaFunction{]T} \AgdaRecord{⇉} \AgdaBound{Δ} \AgdaFunction{\&} \AgdaBound{A}\<%
\\
\>\AgdaBound{f} \AgdaFunction{\textasciicircum} \AgdaBound{A} \AgdaSymbol{=} \AgdaKeyword{record} \<[15]%
\>[15]\<%
\\
\>[0]\AgdaIndent{8}{}\<[8]%
\>[8]\AgdaSymbol{\{} \AgdaField{fn} \AgdaSymbol{=} \AgdaFunction{<} \AgdaFunction{[} \AgdaBound{f} \AgdaFunction{]fn} \AgdaFunction{∘} \AgdaFunction{proj₁} \AgdaFunction{,} \AgdaFunction{proj₂} \AgdaFunction{>}\<%
\\
\>[0]\AgdaIndent{8}{}\<[8]%
\>[8]\AgdaSymbol{;} \AgdaField{resp} \AgdaSymbol{=} \AgdaFunction{<} \AgdaFunction{[} \AgdaBound{f} \AgdaFunction{]resp} \AgdaFunction{∘} \AgdaFunction{proj₁} \AgdaFunction{,} \AgdaFunction{proj₂} \AgdaFunction{>}\<%
\\
\>[0]\AgdaIndent{8}{}\<[8]%
\>[8]\AgdaSymbol{\}}\<%
\\
\>\<\end{code}

$\Pi$-types (object level)

\begin{code}\>\<%
\\
\>\AgdaFunction{Π} \AgdaSymbol{:} \AgdaSymbol{\{}\AgdaBound{Γ} \AgdaSymbol{:} \AgdaRecord{Setoid}\AgdaSymbol{\}(}\AgdaBound{A} \AgdaSymbol{:} \AgdaRecord{Ty} \AgdaBound{Γ}\AgdaSymbol{)(}\AgdaBound{B} \AgdaSymbol{:} \AgdaRecord{Ty} \AgdaSymbol{(}\AgdaBound{Γ} \AgdaFunction{\&} \AgdaBound{A}\AgdaSymbol{))} \AgdaSymbol{→} \AgdaRecord{Ty} \AgdaBound{Γ}\<%
\\
\>\AgdaFunction{Π} \AgdaSymbol{\{}\AgdaBound{Γ}\AgdaSymbol{\}} \AgdaBound{A} \AgdaBound{B} \AgdaSymbol{=} \AgdaKeyword{record} \<[19]%
\>[19]\<%
\\
\>[0]\AgdaIndent{2}{}\<[2]%
\>[2]\AgdaSymbol{\{} \AgdaField{fm} \AgdaSymbol{=} \AgdaSymbol{λ} \AgdaBound{x} \AgdaSymbol{→} \AgdaKeyword{let} \AgdaBound{Ax} \AgdaSymbol{=} \AgdaFunction{[} \AgdaBound{A} \AgdaFunction{]fm} \AgdaBound{x} \AgdaKeyword{in}\<%
\\
\>[0]\AgdaIndent{15}{}\<[15]%
\>[15]\AgdaKeyword{let} \AgdaBound{Bx} \AgdaSymbol{=} \AgdaSymbol{λ} \AgdaBound{a} \AgdaSymbol{→} \AgdaFunction{[} \AgdaBound{B} \AgdaFunction{]fm} \AgdaSymbol{(}\AgdaBound{x} \AgdaInductiveConstructor{,} \AgdaBound{a}\AgdaSymbol{)} \AgdaKeyword{in}\<%
\\
\>[0]\AgdaIndent{9}{}\<[9]%
\>[9]\AgdaKeyword{record}\<%
\\
\>[0]\AgdaIndent{9}{}\<[9]%
\>[9]\AgdaSymbol{\{} \AgdaField{Carrier} \AgdaSymbol{=} \AgdaRecord{Subset} \AgdaSymbol{((}\AgdaBound{a} \AgdaSymbol{:} \AgdaFunction{∣} \AgdaBound{Ax} \AgdaFunction{∣}\AgdaSymbol{)} \AgdaSymbol{→} \AgdaFunction{∣} \AgdaBound{Bx} \AgdaBound{a} \AgdaFunction{∣}\AgdaSymbol{)} \AgdaSymbol{(λ} \AgdaBound{fn} \AgdaSymbol{→} \<[62]%
\>[62]\<%
\\
\>[9]\AgdaIndent{21}{}\<[21]%
\>[21]\AgdaSymbol{(}\AgdaBound{a} \AgdaBound{b} \AgdaSymbol{:} \AgdaFunction{∣} \AgdaBound{Ax} \AgdaFunction{∣}\AgdaSymbol{)}\<%
\\
\>[9]\AgdaIndent{21}{}\<[21]%
\>[21]\AgdaSymbol{(}\AgdaBound{p} \AgdaSymbol{:} \AgdaFunction{[} \AgdaBound{Ax} \AgdaFunction{]} \AgdaBound{a} \AgdaFunction{≈} \AgdaBound{b}\AgdaSymbol{)} \AgdaSymbol{→}\<%
\\
\>[9]\AgdaIndent{21}{}\<[21]%
\>[21]\AgdaFunction{[} \AgdaBound{Bx} \AgdaBound{b} \AgdaFunction{]} \AgdaFunction{[} \AgdaBound{B} \AgdaFunction{]subst} \AgdaSymbol{(}\AgdaFunction{[} \AgdaBound{Γ} \AgdaFunction{]refl} \AgdaInductiveConstructor{,} \<[54]%
\>[54]\<%
\\
\>[9]\AgdaIndent{21}{}\<[21]%
\>[21]\AgdaFunction{[} \AgdaBound{Ax} \AgdaFunction{]trans} \AgdaFunction{[} \AgdaBound{A} \AgdaFunction{]refl*} \AgdaBound{p}\AgdaSymbol{)} \AgdaSymbol{(}\AgdaBound{fn} \AgdaBound{a}\AgdaSymbol{)} \AgdaFunction{≈} \AgdaBound{fn} \AgdaBound{b}\AgdaSymbol{)}\<%
\\
%
\\
%
\\
\>[0]\AgdaIndent{9}{}\<[9]%
\>[9]\AgdaSymbol{;} \AgdaField{\_≈\_} \<[18]%
\>[18]\AgdaSymbol{=} \AgdaSymbol{λ\{(}\AgdaBound{f} \AgdaInductiveConstructor{,} \AgdaSymbol{\_)} \AgdaSymbol{(}\AgdaBound{g} \AgdaInductiveConstructor{,} \AgdaSymbol{\_)} \AgdaSymbol{→} \AgdaSymbol{∀} \AgdaBound{a} \AgdaSymbol{→} \AgdaFunction{[} \AgdaBound{Bx} \AgdaBound{a} \AgdaFunction{]} \AgdaBound{f} \AgdaBound{a} \AgdaFunction{≈} \AgdaBound{g} \AgdaBound{a} \AgdaSymbol{\}}\<%
\\
\>[0]\AgdaIndent{9}{}\<[9]%
\>[9]\AgdaSymbol{;} \AgdaField{refl} \<[19]%
\>[19]\AgdaSymbol{=} \AgdaSymbol{λ} \AgdaBound{a} \AgdaSymbol{→} \AgdaFunction{[} \AgdaBound{Bx} \AgdaSymbol{\_} \AgdaFunction{]refl} \<[40]%
\>[40]\<%
\\
\>[0]\AgdaIndent{9}{}\<[9]%
\>[9]\AgdaSymbol{;} \AgdaField{sym} \<[19]%
\>[19]\AgdaSymbol{=} \AgdaSymbol{λ} \AgdaBound{f} \AgdaBound{a} \AgdaSymbol{→} \AgdaFunction{[} \AgdaBound{Bx} \AgdaSymbol{\_} \AgdaFunction{]sym} \AgdaSymbol{(}\AgdaBound{f} \AgdaBound{a}\AgdaSymbol{)}\<%
\\
\>[0]\AgdaIndent{9}{}\<[9]%
\>[9]\AgdaSymbol{;} \AgdaField{trans} \<[19]%
\>[19]\AgdaSymbol{=} \AgdaSymbol{λ} \AgdaBound{f} \AgdaBound{g} \AgdaBound{a} \AgdaSymbol{→} \AgdaFunction{[} \AgdaBound{Bx} \AgdaSymbol{\_} \AgdaFunction{]trans} \AgdaSymbol{(}\AgdaBound{f} \AgdaBound{a}\AgdaSymbol{)} \AgdaSymbol{(}\AgdaBound{g} \AgdaBound{a}\AgdaSymbol{)}\<%
\\
\>[0]\AgdaIndent{9}{}\<[9]%
\>[9]\AgdaSymbol{\}}\<%
\\
%
\\
\>[0]\AgdaIndent{2}{}\<[2]%
\>[2]\AgdaSymbol{;} \AgdaField{substT} \AgdaSymbol{=} \AgdaSymbol{λ} \AgdaSymbol{\{}\AgdaBound{x}\AgdaSymbol{\}} \AgdaSymbol{\{}\AgdaBound{y}\AgdaSymbol{\}} \AgdaBound{p} \AgdaSymbol{→} \AgdaSymbol{λ} \AgdaSymbol{\{(}\AgdaBound{f} \AgdaInductiveConstructor{,} \AgdaBound{rsp}\AgdaSymbol{)} \AgdaSymbol{→}\<%
\\
\>[0]\AgdaIndent{19}{}\<[19]%
\>[19]\AgdaKeyword{let} \AgdaBound{y2x} \AgdaSymbol{=} \AgdaSymbol{λ} \AgdaBound{a} \AgdaSymbol{→} \AgdaFunction{[} \AgdaBound{A} \AgdaFunction{]subst} \AgdaSymbol{(}\AgdaFunction{[} \AgdaBound{Γ} \AgdaFunction{]sym} \AgdaBound{p}\AgdaSymbol{)} \AgdaBound{a} \AgdaKeyword{in}\<%
\\
\>[0]\AgdaIndent{19}{}\<[19]%
\>[19]\AgdaKeyword{let} \AgdaBound{x2y} \AgdaSymbol{=} \AgdaSymbol{λ} \AgdaBound{a} \AgdaSymbol{→} \AgdaFunction{[} \AgdaBound{A} \AgdaFunction{]subst} \AgdaBound{p} \AgdaBound{a} \AgdaKeyword{in}\<%
\\
\>[0]\AgdaIndent{13}{}\<[13]%
\>[13]\AgdaSymbol{(λ} \AgdaBound{a} \AgdaSymbol{→} \AgdaFunction{[} \AgdaBound{B} \AgdaFunction{]subst} \AgdaSymbol{(}\AgdaBound{p} \AgdaInductiveConstructor{,} \AgdaFunction{[} \AgdaBound{A} \AgdaFunction{]tr*}\AgdaSymbol{)} \<[46]%
\>[46]\<%
\\
\>[0]\AgdaIndent{13}{}\<[13]%
\>[13]\AgdaSymbol{(}\AgdaBound{f} \AgdaSymbol{(}\AgdaBound{y2x} \AgdaBound{a}\AgdaSymbol{)))} \AgdaInductiveConstructor{,} \<[28]%
\>[28]\<%
\\
\>[0]\AgdaIndent{13}{}\<[13]%
\>[13]\AgdaSymbol{(λ} \AgdaBound{a} \AgdaBound{b} \AgdaBound{q} \AgdaSymbol{→}\<%
\\
\>[13]\AgdaIndent{16}{}\<[16]%
\>[16]\AgdaKeyword{let} \AgdaBound{a'} \AgdaSymbol{=} \AgdaBound{y2x} \AgdaBound{a} \AgdaKeyword{in} \<[34]%
\>[34]\<%
\\
\>[13]\AgdaIndent{16}{}\<[16]%
\>[16]\AgdaKeyword{let} \AgdaBound{b'} \AgdaSymbol{=} \AgdaBound{y2x} \AgdaBound{b} \AgdaKeyword{in}\<%
\\
\>[13]\AgdaIndent{16}{}\<[16]%
\>[16]\AgdaKeyword{let} \AgdaBound{q'} \AgdaSymbol{=} \AgdaFunction{[} \AgdaBound{A} \AgdaFunction{]subst*} \AgdaSymbol{(}\AgdaFunction{[} \AgdaBound{Γ} \AgdaFunction{]sym} \AgdaBound{p}\AgdaSymbol{)} \AgdaBound{q} \AgdaKeyword{in}\<%
\\
\>[13]\AgdaIndent{16}{}\<[16]%
\>[16]\AgdaKeyword{let} \AgdaBound{H} \AgdaSymbol{=} \AgdaBound{rsp} \AgdaBound{a'} \AgdaBound{b'} \AgdaSymbol{(}\AgdaFunction{[} \AgdaBound{A} \AgdaFunction{]subst*} \AgdaSymbol{(}\AgdaFunction{[} \AgdaBound{Γ} \AgdaFunction{]sym} \AgdaBound{p}\AgdaSymbol{)} \AgdaBound{q}\AgdaSymbol{)} \AgdaKeyword{in}\<%
\\
\>[13]\AgdaIndent{16}{}\<[16]%
\>[16]\AgdaKeyword{let} \AgdaBound{r} \AgdaSymbol{:} \AgdaFunction{[} \AgdaBound{Γ} \AgdaFunction{\&} \AgdaBound{A} \AgdaFunction{]} \AgdaSymbol{(}\AgdaBound{x} \AgdaInductiveConstructor{,} \AgdaBound{b'}\AgdaSymbol{)} \AgdaFunction{≈} \AgdaSymbol{(}\AgdaBound{y} \AgdaInductiveConstructor{,} \AgdaBound{b}\AgdaSymbol{)}
                    r \AgdaSymbol{=} \AgdaSymbol{(}\AgdaBound{p} \AgdaInductiveConstructor{,} \AgdaFunction{[} \AgdaBound{A} \AgdaFunction{]tr*}\AgdaSymbol{)} \AgdaKeyword{in}\<%
\\
\>[13]\AgdaIndent{16}{}\<[16]%
\>[16]\AgdaKeyword{let} \AgdaBound{pre} \AgdaSymbol{=} \AgdaFunction{[} \AgdaBound{B} \AgdaFunction{]subst*} \AgdaBound{r} \<[40]%
\>[40]\<%
\\
\>[16]\AgdaIndent{24}{}\<[24]%
\>[24]\AgdaSymbol{(}\AgdaBound{rsp} \AgdaBound{a'} \AgdaBound{b'} \AgdaSymbol{(}\AgdaFunction{[} \AgdaBound{A} \AgdaFunction{]subst*} \AgdaSymbol{(}\AgdaFunction{[} \AgdaBound{Γ} \AgdaFunction{]sym} \AgdaBound{p}\AgdaSymbol{)} \AgdaBound{q}\AgdaSymbol{))} \AgdaKeyword{in} \<[68]%
\>[68]\<%
\\
\>[0]\AgdaIndent{16}{}\<[16]%
\>[16]\AgdaFunction{[} \AgdaFunction{[} \AgdaBound{B} \AgdaFunction{]fm} \AgdaSymbol{(}\AgdaBound{y} \AgdaInductiveConstructor{,} \AgdaBound{b}\AgdaSymbol{)} \AgdaFunction{]trans} \<[41]%
\>[41]\<%
\\
\>[0]\AgdaIndent{16}{}\<[16]%
\>[16]\AgdaSymbol{(}\AgdaFunction{[} \AgdaBound{B} \AgdaFunction{]trans*} \AgdaSymbol{\_)} \<[32]%
\>[32]\<%
\\
\>[0]\AgdaIndent{16}{}\<[16]%
\>[16]\AgdaSymbol{(}\AgdaFunction{[} \AgdaFunction{[} \AgdaBound{B} \AgdaFunction{]fm} \AgdaSymbol{(}\AgdaBound{y} \AgdaInductiveConstructor{,} \AgdaBound{b}\AgdaSymbol{)} \AgdaFunction{]trans} \<[42]%
\>[42]\<%
\\
\>[0]\AgdaIndent{16}{}\<[16]%
\>[16]\AgdaSymbol{(}\AgdaFunction{[} \AgdaFunction{[} \AgdaBound{B} \AgdaFunction{]fm} \AgdaSymbol{(}\AgdaBound{y} \AgdaInductiveConstructor{,} \AgdaBound{b}\AgdaSymbol{)} \AgdaFunction{]sym} \AgdaSymbol{(}\AgdaFunction{[} \AgdaBound{B} \AgdaFunction{]trans*} \AgdaSymbol{\_))} \<[57]%
\>[57]\<%
\\
\>[0]\AgdaIndent{16}{}\<[16]%
\>[16]\AgdaBound{pre}\AgdaSymbol{))\}}\<%
\\
%
\\
%
\\
\>[0]\AgdaIndent{2}{}\<[2]%
\>[2]\AgdaSymbol{;} \AgdaField{subst*} \AgdaSymbol{=} \AgdaSymbol{λ} \AgdaBound{\_} \AgdaBound{q} \AgdaBound{\_} \AgdaSymbol{→} \AgdaFunction{[} \AgdaBound{B} \AgdaFunction{]subst*} \AgdaSymbol{\_} \AgdaSymbol{(}\AgdaBound{q} \AgdaSymbol{\_)}\<%
\\
\>[0]\AgdaIndent{2}{}\<[2]%
\>[2]\AgdaSymbol{;} \AgdaField{refl*} \AgdaSymbol{=} \AgdaSymbol{λ} \AgdaSymbol{\{}\AgdaBound{x}\AgdaSymbol{\}} \AgdaSymbol{\{}\AgdaBound{a}\AgdaSymbol{\}} \AgdaBound{ax} \<[25]%
\>[25]\<%
\\
\>[2]\AgdaIndent{10}{}\<[10]%
\>[10]\AgdaSymbol{→} \AgdaKeyword{let} \AgdaBound{rsp} \AgdaSymbol{=} \AgdaFunction{prj₂} \AgdaBound{a} \AgdaKeyword{in} \AgdaSymbol{(}\AgdaBound{rsp} \AgdaSymbol{\_} \AgdaSymbol{\_} \AgdaFunction{[} \AgdaBound{A} \AgdaFunction{]refl*}\AgdaSymbol{)}\<%
\\
\>[0]\AgdaIndent{2}{}\<[2]%
\>[2]\AgdaSymbol{;} \AgdaField{trans*} \AgdaSymbol{=} \<[14]%
\>[14]\AgdaSymbol{λ} \AgdaSymbol{\{(}\AgdaBound{f} \AgdaInductiveConstructor{,} \AgdaBound{rsp}\AgdaSymbol{)} \AgdaBound{a} \AgdaSymbol{→}\<%
\\
\>[0]\AgdaIndent{13}{}\<[13]%
\>[13]\AgdaFunction{[} \AgdaFunction{[} \AgdaBound{B} \AgdaFunction{]fm} \AgdaSymbol{\_} \AgdaFunction{]trans} \<[32]%
\>[32]\<%
\\
\>[0]\AgdaIndent{13}{}\<[13]%
\>[13]\AgdaSymbol{(}\AgdaFunction{[} \AgdaFunction{[} \AgdaBound{B} \AgdaFunction{]fm} \AgdaSymbol{\_} \AgdaFunction{]trans} \<[33]%
\>[33]\<%
\\
\>[0]\AgdaIndent{13}{}\<[13]%
\>[13]\AgdaSymbol{(}\AgdaFunction{[} \AgdaBound{B} \AgdaFunction{]trans*} \AgdaSymbol{\_)} \<[29]%
\>[29]\<%
\\
\>[0]\AgdaIndent{13}{}\<[13]%
\>[13]\AgdaSymbol{(}\AgdaFunction{[} \AgdaFunction{[} \AgdaBound{B} \AgdaFunction{]fm} \AgdaSymbol{\_} \AgdaFunction{]sym} \AgdaSymbol{(}\AgdaFunction{[} \AgdaBound{B} \AgdaFunction{]trans*} \AgdaSymbol{\_)))}\<%
\\
\>[0]\AgdaIndent{13}{}\<[13]%
\>[13]\AgdaSymbol{(}\AgdaFunction{[} \AgdaBound{B} \AgdaFunction{]subst*} \AgdaSymbol{\_} \AgdaSymbol{(}\AgdaBound{rsp} \AgdaSymbol{\_} \AgdaSymbol{\_} \AgdaSymbol{(}\AgdaFunction{[} \AgdaBound{A} \AgdaFunction{]trans*} \AgdaSymbol{\_)} \AgdaSymbol{))} \AgdaSymbol{\}}\<%
\\
\>[0]\AgdaIndent{2}{}\<[2]%
\>[2]\AgdaSymbol{\}}\<%
\\
%
\\
%
\\
\>\AgdaFunction{lam} \AgdaSymbol{:} \AgdaSymbol{\{}\AgdaBound{Γ} \AgdaSymbol{:} \AgdaFunction{Con}\AgdaSymbol{\}\{}\AgdaBound{A} \AgdaSymbol{:} \AgdaRecord{Ty} \AgdaBound{Γ}\AgdaSymbol{\}\{}\AgdaBound{B} \AgdaSymbol{:} \AgdaRecord{Ty} \AgdaSymbol{(}\AgdaBound{Γ} \AgdaFunction{\&} \AgdaBound{A}\AgdaSymbol{)\}} \AgdaSymbol{→} \AgdaRecord{Tm} \AgdaBound{B} \AgdaSymbol{→} \AgdaRecord{Tm} \AgdaSymbol{(}\AgdaFunction{Π} \AgdaBound{A} \AgdaBound{B}\AgdaSymbol{)}\<%
\\
\>\AgdaFunction{lam} \AgdaSymbol{\{}\AgdaBound{Γ}\AgdaSymbol{\}} \AgdaSymbol{\{}\AgdaBound{A}\AgdaSymbol{\}} \AgdaSymbol{(}\AgdaInductiveConstructor{tm:} \AgdaBound{tm} \AgdaInductiveConstructor{resp:} \AgdaBound{respt}\AgdaSymbol{)} \AgdaSymbol{=} \<[35]%
\>[35]\<%
\\
\>[0]\AgdaIndent{2}{}\<[2]%
\>[2]\AgdaKeyword{record} \AgdaSymbol{\{} \AgdaField{tm} \AgdaSymbol{=} \AgdaSymbol{λ} \AgdaBound{x} \AgdaSymbol{→} \AgdaSymbol{(λ} \AgdaBound{a} \AgdaSymbol{→} \AgdaBound{tm} \AgdaSymbol{(}\AgdaBound{x} \AgdaInductiveConstructor{,} \AgdaBound{a}\AgdaSymbol{))} \<[41]%
\>[41]\<%
\\
\>[2]\AgdaIndent{16}{}\<[16]%
\>[16]\AgdaInductiveConstructor{,} \AgdaSymbol{(λ} \AgdaBound{a} \AgdaBound{b} \AgdaBound{p} \AgdaSymbol{→} \AgdaBound{respt} \AgdaSymbol{(}\AgdaFunction{[} \AgdaBound{Γ} \AgdaFunction{]refl} \AgdaInductiveConstructor{,} \<[48]%
\>[48]\<%
\\
\>[2]\AgdaIndent{16}{}\<[16]%
\>[16]\AgdaFunction{[} \AgdaFunction{[} \AgdaBound{A} \AgdaFunction{]fm} \AgdaBound{x} \AgdaFunction{]trans} \AgdaFunction{[} \AgdaBound{A} \AgdaFunction{]refl*} \AgdaBound{p}\AgdaSymbol{))}\<%
\\
\>[5]\AgdaIndent{9}{}\<[9]%
\>[9]\AgdaSymbol{;} \AgdaField{respt} \AgdaSymbol{=} \AgdaSymbol{λ} \AgdaBound{p} \AgdaBound{\_} \AgdaSymbol{→} \AgdaBound{respt} \AgdaSymbol{(}\AgdaBound{p} \AgdaInductiveConstructor{,} \AgdaFunction{[} \AgdaBound{A} \AgdaFunction{]tr*}\AgdaSymbol{)}\<%
\\
\>[0]\AgdaIndent{9}{}\<[9]%
\>[9]\AgdaSymbol{\}}\<%
\\
%
\\
%
\\
\>\AgdaFunction{app} \AgdaSymbol{:} \AgdaSymbol{\{}\AgdaBound{Γ} \AgdaSymbol{:} \AgdaFunction{Con}\AgdaSymbol{\}\{}\AgdaBound{A} \AgdaSymbol{:} \AgdaRecord{Ty} \AgdaBound{Γ}\AgdaSymbol{\}\{}\AgdaBound{B} \AgdaSymbol{:} \AgdaRecord{Ty} \AgdaSymbol{(}\AgdaBound{Γ} \AgdaFunction{\&} \AgdaBound{A}\AgdaSymbol{)\}} \AgdaSymbol{→} \AgdaRecord{Tm} \AgdaSymbol{(}\AgdaFunction{Π} \AgdaBound{A} \AgdaBound{B}\AgdaSymbol{)} \AgdaSymbol{→} \AgdaRecord{Tm} \AgdaBound{B}\<%
\\
\>\AgdaFunction{app} \AgdaSymbol{\{}\AgdaBound{Γ}\AgdaSymbol{\}} \AgdaSymbol{\{}\AgdaBound{A}\AgdaSymbol{\}} \AgdaSymbol{\{}\AgdaBound{B}\AgdaSymbol{\}} \AgdaSymbol{(}\AgdaInductiveConstructor{tm:} \AgdaBound{tm} \AgdaInductiveConstructor{resp:} \AgdaBound{respt}\AgdaSymbol{)} \AgdaSymbol{=} \<[39]%
\>[39]\<%
\\
\>[0]\AgdaIndent{2}{}\<[2]%
\>[2]\AgdaKeyword{record} \AgdaSymbol{\{} \AgdaField{tm} \AgdaSymbol{=} \AgdaSymbol{λ} \AgdaSymbol{\{(}\AgdaBound{x} \AgdaInductiveConstructor{,} \AgdaBound{a}\AgdaSymbol{)} \AgdaSymbol{→} \AgdaFunction{prj₁} \AgdaSymbol{(}\AgdaBound{tm} \AgdaBound{x}\AgdaSymbol{)} \AgdaBound{a}\AgdaSymbol{\}}\<%
\\
\>[0]\AgdaIndent{9}{}\<[9]%
\>[9]\AgdaSymbol{;} \AgdaField{respt} \AgdaSymbol{=} \AgdaSymbol{λ} \AgdaSymbol{\{}\AgdaBound{x}\AgdaSymbol{\}} \AgdaSymbol{\{}\AgdaBound{y}\AgdaSymbol{\}} \AgdaSymbol{→} \AgdaSymbol{λ} \AgdaSymbol{\{(}\AgdaBound{p} \AgdaInductiveConstructor{,} \AgdaBound{tr}\AgdaSymbol{)} \AgdaSymbol{→} \<[45]%
\>[45]\<%
\\
\>[9]\AgdaIndent{17}{}\<[17]%
\>[17]\AgdaKeyword{let} \AgdaBound{fresp} \AgdaSymbol{=} \AgdaFunction{prj₂} \AgdaSymbol{(}\AgdaBound{tm} \AgdaSymbol{(}\AgdaFunction{proj₁} \AgdaBound{x}\AgdaSymbol{))} \AgdaKeyword{in}\<%
\\
\>[9]\AgdaIndent{17}{}\<[17]%
\>[17]\AgdaFunction{[} \AgdaFunction{[} \AgdaBound{B} \AgdaFunction{]fm} \AgdaSymbol{\_} \AgdaFunction{]trans} \<[36]%
\>[36]\<%
\\
\>[9]\AgdaIndent{17}{}\<[17]%
\>[17]\AgdaSymbol{(}\AgdaFunction{[} \AgdaBound{B} \AgdaFunction{]subst*} \AgdaSymbol{(}\AgdaBound{p} \AgdaInductiveConstructor{,} \AgdaBound{tr}\AgdaSymbol{)} \<[39]%
\>[39]\<%
\\
\>[9]\AgdaIndent{17}{}\<[17]%
\>[17]\AgdaSymbol{(}\AgdaFunction{[} \AgdaFunction{[} \AgdaBound{B} \AgdaFunction{]fm} \AgdaSymbol{\_} \AgdaFunction{]sym} \AgdaFunction{[} \AgdaBound{B} \AgdaFunction{]refl*}\AgdaSymbol{))}\<%
\\
\>[9]\AgdaIndent{17}{}\<[17]%
\>[17]\AgdaSymbol{(}\AgdaFunction{[} \AgdaFunction{[} \AgdaBound{B} \AgdaFunction{]fm} \AgdaSymbol{\_} \AgdaFunction{]trans}\<%
\\
\>[9]\AgdaIndent{17}{}\<[17]%
\>[17]\AgdaSymbol{(}\AgdaFunction{[} \AgdaBound{B} \AgdaFunction{]trans*} \AgdaSymbol{\{}p \AgdaSymbol{=} \AgdaSymbol{(}\AgdaFunction{[} \AgdaBound{Γ} \AgdaFunction{]refl} \AgdaInductiveConstructor{,} \AgdaFunction{[} \AgdaBound{A} \AgdaFunction{]refl*}\AgdaSymbol{)\}} \AgdaSymbol{\_)}\<%
\\
\>[0]\AgdaIndent{8}{}\<[8]%
\>[8]\<%
\\
\>[0]\AgdaIndent{17}{}\<[17]%
\>[17]\AgdaSymbol{(}\AgdaFunction{[} \AgdaFunction{[} \AgdaBound{B} \AgdaFunction{]fm} \AgdaSymbol{\_} \AgdaFunction{]trans} \<[37]%
\>[37]\<%
\\
\>[0]\AgdaIndent{17}{}\<[17]%
\>[17]\AgdaSymbol{(}\AgdaFunction{[} \AgdaFunction{[} \AgdaBound{B} \AgdaFunction{]fm} \AgdaSymbol{\_} \AgdaFunction{]sym} \<[35]%
\>[35]\<%
\\
\>[0]\AgdaIndent{17}{}\<[17]%
\>[17]\AgdaSymbol{(}\AgdaFunction{[} \AgdaBound{B} \AgdaFunction{]trans*} \AgdaSymbol{\{}q \AgdaSymbol{=} \AgdaSymbol{(}\AgdaBound{p} \AgdaInductiveConstructor{,} \AgdaFunction{[} \AgdaBound{A} \AgdaFunction{]tr*}\AgdaSymbol{)\}} \AgdaSymbol{\_))}\<%
\\
\>[0]\AgdaIndent{17}{}\<[17]%
\>[17]\AgdaSymbol{(}\AgdaFunction{[} \AgdaFunction{[} \AgdaBound{B} \AgdaFunction{]fm} \AgdaSymbol{\_} \AgdaFunction{]trans} \<[37]%
\>[37]\<%
\\
\>[0]\AgdaIndent{17}{}\<[17]%
\>[17]\AgdaSymbol{(}\AgdaFunction{[} \AgdaBound{B} \AgdaFunction{]subst*} \AgdaSymbol{\_} \AgdaSymbol{(}\AgdaBound{fresp} \AgdaSymbol{\_} \AgdaSymbol{\_} \<[43]%
\>[43]\<%
\\
\>[17]\AgdaIndent{19}{}\<[19]%
\>[19]\AgdaSymbol{(}\AgdaFunction{[} \AgdaFunction{[} \AgdaBound{A} \AgdaFunction{]fm} \AgdaSymbol{\_} \AgdaFunction{]trans} \AgdaSymbol{(}\AgdaFunction{[} \AgdaFunction{[} \AgdaBound{A} \AgdaFunction{]fm} \AgdaSymbol{\_} \AgdaFunction{]sym} \AgdaFunction{[} \AgdaBound{A} \AgdaFunction{]tr*}\AgdaSymbol{)} \<[67]%
\>[67]\<%
\\
\>[17]\AgdaIndent{19}{}\<[19]%
\>[19]\AgdaSymbol{(}\AgdaFunction{[} \AgdaBound{A} \AgdaFunction{]subst*} \AgdaSymbol{(}\AgdaFunction{[} \AgdaBound{Γ} \AgdaFunction{]sym} \AgdaBound{p}\AgdaSymbol{)} \AgdaBound{tr}\AgdaSymbol{))))}\<%
\\
\>[0]\AgdaIndent{17}{}\<[17]%
\>[17]\AgdaSymbol{(}\AgdaBound{respt} \AgdaBound{p} \AgdaSymbol{\_))))} \AgdaSymbol{\}}\<%
\\
\>[0]\AgdaIndent{9}{}\<[9]%
\>[9]\AgdaSymbol{\}}\<%
\\
%
\\
\>\AgdaFunction{\_⇒\_} \AgdaSymbol{:} \AgdaSymbol{\{}\AgdaBound{Γ} \AgdaSymbol{:} \AgdaFunction{Con}\AgdaSymbol{\}(}\AgdaBound{A} \AgdaBound{B} \AgdaSymbol{:} \AgdaRecord{Ty} \AgdaBound{Γ}\AgdaSymbol{)} \AgdaSymbol{→} \AgdaRecord{Ty} \AgdaBound{Γ}\<%
\\
\>\AgdaBound{A} \AgdaFunction{⇒} \AgdaBound{B} \AgdaSymbol{=} \AgdaFunction{Π} \AgdaBound{A} \AgdaSymbol{(}\AgdaBound{B} \AgdaFunction{[} \AgdaFunction{fst\&} \AgdaSymbol{\{}A \AgdaSymbol{=} \AgdaBound{A}\AgdaSymbol{\}} \AgdaFunction{]T}\AgdaSymbol{)}\<%
\\
%
\\
\>\AgdaKeyword{infixr} \AgdaNumber{6} \_⇒\_\<%
\\
\>\<\end{code}

Simpler definition for functions

\begin{code}\>\<%
\\
\>\AgdaFunction{[\_,\_]\_⇒fm\_} \AgdaSymbol{:} \AgdaSymbol{(}\AgdaBound{Γ} \AgdaSymbol{:} \AgdaFunction{Con}\AgdaSymbol{)(}\AgdaBound{x} \AgdaSymbol{:} \AgdaFunction{∣} \AgdaBound{Γ} \AgdaFunction{∣}\AgdaSymbol{)} \<[34]%
\>[34]\<%
\\
\>[9]\AgdaIndent{11}{}\<[11]%
\>[11]\AgdaSymbol{→} \AgdaRecord{Setoid} \AgdaSymbol{→} \AgdaRecord{Setoid} \AgdaSymbol{→} \AgdaRecord{Setoid}\<%
\\
\>\AgdaFunction{[} \AgdaBound{Γ} \AgdaFunction{,} \AgdaBound{x} \AgdaFunction{]} \AgdaBound{Ax} \AgdaFunction{⇒fm} \AgdaBound{Bx} \<[20]%
\>[20]\<%
\\
\>[0]\AgdaIndent{2}{}\<[2]%
\>[2]\AgdaSymbol{=} \AgdaKeyword{record}\<%
\\
\>[0]\AgdaIndent{6}{}\<[6]%
\>[6]\AgdaSymbol{\{} \AgdaField{Carrier} \AgdaSymbol{=} \AgdaRecord{Σ[} \AgdaBound{fn} \AgdaRecord{∶} \AgdaSymbol{(}\AgdaFunction{∣} \AgdaBound{Ax} \AgdaFunction{∣} \AgdaSymbol{→} \AgdaFunction{∣} \AgdaBound{Bx} \AgdaFunction{∣}\AgdaSymbol{)} \AgdaRecord{]} \AgdaSymbol{((}\AgdaBound{a} \AgdaBound{b} \AgdaSymbol{:} \AgdaFunction{∣} \AgdaBound{Ax} \AgdaFunction{∣}\AgdaSymbol{)}\<%
\\
\>[6]\AgdaIndent{18}{}\<[18]%
\>[18]\AgdaSymbol{(}\AgdaBound{p} \AgdaSymbol{:} \AgdaFunction{[} \AgdaBound{Ax} \AgdaFunction{]} \AgdaBound{a} \AgdaFunction{≈} \AgdaBound{b}\AgdaSymbol{)} \AgdaSymbol{→} \AgdaFunction{[} \AgdaBound{Bx} \AgdaFunction{]} \AgdaBound{fn} \AgdaBound{a} \AgdaFunction{≈} \AgdaBound{fn} \AgdaBound{b}\AgdaSymbol{)}\<%
\\
\>[0]\AgdaIndent{6}{}\<[6]%
\>[6]\AgdaSymbol{;} \AgdaField{\_≈\_} \<[15]%
\>[15]\AgdaSymbol{=} \AgdaSymbol{λ\{(}\AgdaBound{f} \AgdaInductiveConstructor{,} \AgdaSymbol{\_)} \AgdaSymbol{(}\AgdaBound{g} \AgdaInductiveConstructor{,} \AgdaSymbol{\_)} \AgdaSymbol{→} \AgdaSymbol{∀} \AgdaBound{a} \AgdaSymbol{→} \AgdaFunction{[} \AgdaBound{Bx} \AgdaFunction{]} \AgdaBound{f} \AgdaBound{a} \AgdaFunction{≈} \AgdaBound{g} \AgdaBound{a} \AgdaSymbol{\}}\<%
\\
\>[0]\AgdaIndent{6}{}\<[6]%
\>[6]\AgdaSymbol{;} \AgdaField{refl} \<[16]%
\>[16]\AgdaSymbol{=} \AgdaSymbol{λ} \AgdaBound{\_} \AgdaSymbol{→} \AgdaFunction{[} \AgdaBound{Bx} \AgdaFunction{]refl} \<[35]%
\>[35]\<%
\\
\>[0]\AgdaIndent{6}{}\<[6]%
\>[6]\AgdaSymbol{;} \AgdaField{sym} \<[16]%
\>[16]\AgdaSymbol{=} \AgdaSymbol{λ} \AgdaBound{f} \AgdaBound{a} \AgdaSymbol{→} \AgdaFunction{[} \AgdaBound{Bx} \AgdaFunction{]sym} \AgdaSymbol{(}\AgdaBound{f} \AgdaBound{a}\AgdaSymbol{)}\<%
\\
\>[0]\AgdaIndent{6}{}\<[6]%
\>[6]\AgdaSymbol{;} \AgdaField{trans} \<[16]%
\>[16]\AgdaSymbol{=} \AgdaSymbol{λ} \AgdaBound{f} \AgdaBound{g} \AgdaBound{a} \AgdaSymbol{→} \AgdaFunction{[} \AgdaBound{Bx} \AgdaFunction{]trans} \AgdaSymbol{(}\AgdaBound{f} \AgdaBound{a}\AgdaSymbol{)} \AgdaSymbol{(}\AgdaBound{g} \AgdaBound{a}\AgdaSymbol{)}\<%
\\
\>[0]\AgdaIndent{6}{}\<[6]%
\>[6]\AgdaSymbol{\}}\<%
\\
\>\<\end{code}

$\Sigma$-types (object level)

\begin{code}\>\<%
\\
%
\\
\>\AgdaFunction{Σ'} \AgdaSymbol{:} \AgdaSymbol{\{}\AgdaBound{Γ} \AgdaSymbol{:} \AgdaFunction{Con}\AgdaSymbol{\}(}\AgdaBound{A} \AgdaSymbol{:} \AgdaRecord{Ty} \AgdaBound{Γ}\AgdaSymbol{)(}\AgdaBound{B} \AgdaSymbol{:} \AgdaRecord{Ty} \AgdaSymbol{(}\AgdaBound{Γ} \AgdaFunction{\&} \AgdaBound{A}\AgdaSymbol{))} \AgdaSymbol{→} \AgdaRecord{Ty} \AgdaBound{Γ}\<%
\\
\>\AgdaFunction{Σ'} \AgdaSymbol{\{}\AgdaBound{Γ}\AgdaSymbol{\}} \AgdaBound{A} \AgdaBound{B} \AgdaSymbol{=} \AgdaKeyword{record} \<[20]%
\>[20]\<%
\\
\>[6]\AgdaIndent{8}{}\<[8]%
\>[8]\AgdaSymbol{\{} \AgdaField{fm} \AgdaSymbol{=} \AgdaSymbol{λ} \AgdaBound{x} \AgdaSymbol{→} \AgdaKeyword{let} \AgdaBound{Ax} \AgdaSymbol{=} \AgdaFunction{[} \AgdaBound{A} \AgdaFunction{]fm} \AgdaBound{x} \AgdaKeyword{in}\<%
\\
\>[8]\AgdaIndent{15}{}\<[15]%
\>[15]\AgdaKeyword{let} \AgdaBound{Bx} \AgdaSymbol{=} \AgdaSymbol{λ} \AgdaBound{a} \AgdaSymbol{→} \AgdaFunction{[} \AgdaBound{B} \AgdaFunction{]fm} \AgdaSymbol{(}\AgdaBound{x} \AgdaInductiveConstructor{,} \AgdaBound{a}\AgdaSymbol{)} \AgdaKeyword{in}\<%
\\
\>[-2]\AgdaIndent{9}{}\<[9]%
\>[9]\AgdaKeyword{record}\<%
\\
\>[0]\AgdaIndent{11}{}\<[11]%
\>[11]\AgdaSymbol{\{} \AgdaField{Carrier} \AgdaSymbol{=} \AgdaRecord{Σ[} \AgdaBound{a} \AgdaRecord{∶} \AgdaFunction{∣} \AgdaBound{Ax} \AgdaFunction{∣} \AgdaRecord{]} \AgdaFunction{∣} \AgdaBound{Bx} \AgdaBound{a} \AgdaFunction{∣}\<%
\\
%
\\
\>[0]\AgdaIndent{11}{}\<[11]%
\>[11]\AgdaSymbol{;} \AgdaField{\_≈\_} \<[20]%
\>[20]\AgdaSymbol{=} \AgdaSymbol{λ\{(}\AgdaBound{a₁} \AgdaInductiveConstructor{,} \AgdaBound{b₁}\AgdaSymbol{)} \AgdaSymbol{(}\AgdaBound{a₂} \AgdaInductiveConstructor{,} \AgdaBound{b₂}\AgdaSymbol{)} \AgdaSymbol{→} \<[46]%
\>[46]\<%
\\
\>[11]\AgdaIndent{22}{}\<[22]%
\>[22]\AgdaRecord{Subset} \AgdaSymbol{(}\AgdaFunction{[} \AgdaBound{Ax} \AgdaFunction{]} \AgdaBound{a₁} \AgdaFunction{≈} \AgdaBound{a₂}\AgdaSymbol{)} \<[46]%
\>[46]\<%
\\
\>[11]\AgdaIndent{22}{}\<[22]%
\>[22]\AgdaSymbol{(λ} \AgdaBound{eq₁} \AgdaSymbol{→} \AgdaFunction{[} \AgdaBound{Bx} \AgdaSymbol{\_} \AgdaFunction{]} \AgdaFunction{[} \AgdaBound{B} \AgdaFunction{]subst} \<[51]%
\>[51]\<%
\\
\>[22]\AgdaIndent{23}{}\<[23]%
\>[23]\AgdaSymbol{(}\AgdaFunction{[} \AgdaBound{Γ} \AgdaFunction{]refl} \AgdaInductiveConstructor{,} \AgdaFunction{[} \AgdaFunction{[} \AgdaBound{A} \AgdaFunction{]fm} \AgdaBound{x} \AgdaFunction{]trans} \<[55]%
\>[55]\<%
\\
\>[22]\AgdaIndent{23}{}\<[23]%
\>[23]\AgdaFunction{[} \AgdaBound{A} \AgdaFunction{]refl*} \AgdaBound{eq₁}\AgdaSymbol{)} \AgdaBound{b₁} \AgdaFunction{≈} \AgdaBound{b₂}\AgdaSymbol{)}\<%
\\
\>[-10]\AgdaIndent{11}{}\<[11]%
\>[11]\AgdaSymbol{\}}\<%
\\
%
\\
\>[0]\AgdaIndent{11}{}\<[11]%
\>[11]\AgdaSymbol{;} \AgdaField{refl} \<[21]%
\>[21]\AgdaSymbol{=} \AgdaSymbol{λ} \AgdaSymbol{\{}\AgdaBound{t}\AgdaSymbol{\}} \AgdaSymbol{→} \AgdaFunction{[} \AgdaBound{Ax} \AgdaFunction{]refl} \AgdaInductiveConstructor{,} \AgdaFunction{[} \AgdaBound{B} \AgdaFunction{]refl*}\<%
\\
%
\\
\>[0]\AgdaIndent{11}{}\<[11]%
\>[11]\AgdaSymbol{;} \AgdaField{sym} \<[21]%
\>[21]\AgdaSymbol{=} \AgdaSymbol{λ} \AgdaSymbol{\{(}\AgdaBound{p} \AgdaInductiveConstructor{,} \AgdaBound{q}\AgdaSymbol{)} \AgdaSymbol{→} \AgdaSymbol{(}\AgdaFunction{[} \AgdaBound{Ax} \AgdaFunction{]sym} \AgdaBound{p}\AgdaSymbol{)} \AgdaInductiveConstructor{,} \<[52]%
\>[52]\<%
\\
\>[11]\AgdaIndent{23}{}\<[23]%
\>[23]\AgdaFunction{[} \AgdaBound{Bx} \AgdaSymbol{\_} \AgdaFunction{]trans} \AgdaSymbol{(}\AgdaFunction{[} \AgdaBound{B} \AgdaFunction{]subst*} \AgdaSymbol{\_} \<[52]%
\>[52]\<%
\\
\>[11]\AgdaIndent{23}{}\<[23]%
\>[23]\AgdaSymbol{(}\AgdaFunction{[} \AgdaBound{Bx} \AgdaSymbol{\_} \AgdaFunction{]sym} \AgdaBound{q}\AgdaSymbol{))} \AgdaFunction{[} \AgdaBound{B} \AgdaFunction{]tr*}\AgdaSymbol{\}}\<%
\\
%
\\
\>[0]\AgdaIndent{11}{}\<[11]%
\>[11]\AgdaSymbol{;} \AgdaField{trans} \<[21]%
\>[21]\AgdaSymbol{=} \AgdaSymbol{λ} \AgdaSymbol{\{(}\AgdaBound{p} \AgdaInductiveConstructor{,} \AgdaBound{q}\AgdaSymbol{)} \AgdaSymbol{(}\AgdaBound{r} \AgdaInductiveConstructor{,} \AgdaBound{s}\AgdaSymbol{)} \AgdaSymbol{→} \AgdaSymbol{(}\AgdaFunction{[} \AgdaBound{Ax} \AgdaFunction{]trans} \AgdaBound{p} \AgdaBound{r}\AgdaSymbol{)} \AgdaInductiveConstructor{,}\<%
\\
\>[0]\AgdaIndent{23}{}\<[23]%
\>[23]\AgdaFunction{[} \AgdaBound{Bx} \AgdaSymbol{\_} \AgdaFunction{]trans} \AgdaSymbol{(}\AgdaFunction{[} \AgdaBound{Bx} \AgdaSymbol{\_} \AgdaFunction{]trans} \<[52]%
\>[52]\<%
\\
\>[0]\AgdaIndent{23}{}\<[23]%
\>[23]\AgdaSymbol{(}\AgdaFunction{[} \AgdaBound{Bx} \AgdaSymbol{\_} \AgdaFunction{]sym} \AgdaSymbol{(}\AgdaFunction{[} \AgdaBound{B} \AgdaFunction{]trans*} \AgdaSymbol{\_))} \<[53]%
\>[53]\<%
\\
\>[0]\AgdaIndent{23}{}\<[23]%
\>[23]\AgdaSymbol{(}\AgdaFunction{[} \AgdaBound{B} \AgdaFunction{]subst*} \AgdaSymbol{\_} \AgdaBound{q}\AgdaSymbol{))} \AgdaBound{s}\AgdaSymbol{\}}\<%
\\
\>[0]\AgdaIndent{11}{}\<[11]%
\>[11]\AgdaSymbol{\}}\<%
\\
%
\\
\>[0]\AgdaIndent{8}{}\<[8]%
\>[8]\AgdaSymbol{;} \AgdaField{substT} \AgdaSymbol{=} \AgdaSymbol{λ} \AgdaBound{x≈y} \AgdaSymbol{→} \AgdaSymbol{λ} \AgdaSymbol{\{(}\AgdaBound{p} \AgdaInductiveConstructor{,} \AgdaBound{q}\AgdaSymbol{)} \AgdaSymbol{→} \<[40]%
\>[40]\<%
\\
\>[8]\AgdaIndent{19}{}\<[19]%
\>[19]\AgdaSymbol{(}\AgdaFunction{[} \AgdaBound{A} \AgdaFunction{]subst} \AgdaBound{x≈y} \AgdaBound{p}\AgdaSymbol{)} \AgdaInductiveConstructor{,} \AgdaFunction{[} \AgdaBound{B} \AgdaFunction{]subst} \AgdaSymbol{(}\AgdaBound{x≈y} \AgdaInductiveConstructor{,} \<[58]%
\>[58]\<%
\\
\>[8]\AgdaIndent{19}{}\<[19]%
\>[19]\AgdaFunction{[} \AgdaFunction{[} \AgdaBound{A} \AgdaFunction{]fm} \AgdaSymbol{\_} \AgdaFunction{]refl}\AgdaSymbol{)} \AgdaBound{q}\AgdaSymbol{\}}\<%
\\
%
\\
\>[0]\AgdaIndent{8}{}\<[8]%
\>[8]\AgdaSymbol{;} \AgdaField{subst*} \AgdaSymbol{=} \AgdaSymbol{λ} \AgdaBound{x≈y} \AgdaSymbol{→} \AgdaSymbol{λ} \AgdaSymbol{\{(}\AgdaBound{p} \AgdaInductiveConstructor{,} \AgdaBound{q}\AgdaSymbol{)} \AgdaSymbol{→} \AgdaFunction{[} \AgdaBound{A} \AgdaFunction{]subst*} \AgdaBound{x≈y} \AgdaBound{p} \AgdaInductiveConstructor{,} \<[60]%
\>[60]\<%
\\
\>[0]\AgdaIndent{19}{}\<[19]%
\>[19]\AgdaFunction{[} \AgdaFunction{[} \AgdaBound{B} \AgdaFunction{]fm} \AgdaSymbol{\_} \AgdaFunction{]trans} \AgdaSymbol{(}\AgdaFunction{[} \AgdaFunction{[} \AgdaBound{B} \AgdaFunction{]fm} \AgdaSymbol{\_} \AgdaFunction{]trans} \<[58]%
\>[58]\<%
\\
\>[0]\AgdaIndent{19}{}\<[19]%
\>[19]\AgdaSymbol{(}\AgdaFunction{[} \AgdaBound{B} \AgdaFunction{]trans*} \AgdaSymbol{\_)}\<%
\\
\>[0]\AgdaIndent{19}{}\<[19]%
\>[19]\AgdaSymbol{(}\AgdaFunction{[} \AgdaFunction{[} \AgdaBound{B} \AgdaFunction{]fm} \AgdaSymbol{\_} \AgdaFunction{]sym} \AgdaSymbol{(}\AgdaFunction{[} \AgdaBound{B} \AgdaFunction{]trans*} \AgdaSymbol{\_)))} \<[55]%
\>[55]\<%
\\
\>[0]\AgdaIndent{19}{}\<[19]%
\>[19]\AgdaSymbol{(}\AgdaFunction{[} \AgdaBound{B} \AgdaFunction{]subst*} \AgdaSymbol{(}\AgdaBound{x≈y} \AgdaInductiveConstructor{,} \AgdaFunction{[} \AgdaFunction{[} \AgdaBound{A} \AgdaFunction{]fm} \AgdaSymbol{\_} \AgdaFunction{]refl}\AgdaSymbol{)} \AgdaBound{q}\AgdaSymbol{)} \AgdaSymbol{\}}\<%
\\
\>[0]\AgdaIndent{8}{}\<[8]%
\>[8]\AgdaSymbol{;} \AgdaField{refl*} \AgdaSymbol{=} \AgdaSymbol{λ} \AgdaSymbol{\{}\AgdaBound{x}\AgdaSymbol{\}} \AgdaSymbol{\{}\AgdaBound{a}\AgdaSymbol{\}} \AgdaSymbol{→} \<[30]%
\>[30]\<%
\\
\>[0]\AgdaIndent{18}{}\<[18]%
\>[18]\AgdaKeyword{let} \AgdaSymbol{(}\AgdaBound{p} \AgdaInductiveConstructor{,} \AgdaBound{q}\AgdaSymbol{)} \AgdaSymbol{=} \AgdaBound{a} \AgdaKeyword{in} \AgdaFunction{[} \AgdaBound{A} \AgdaFunction{]refl*} \AgdaInductiveConstructor{,} \AgdaFunction{[} \AgdaBound{B} \AgdaFunction{]tr*}\<%
\\
\>[0]\AgdaIndent{8}{}\<[8]%
\>[8]\AgdaSymbol{;} \AgdaField{trans*} \AgdaSymbol{=} \<[20]%
\>[20]\AgdaSymbol{λ} \AgdaSymbol{\{(}\AgdaBound{p} \AgdaInductiveConstructor{,} \AgdaBound{q}\AgdaSymbol{)} \<[32]%
\>[32]\AgdaSymbol{→} \AgdaSymbol{(}\AgdaFunction{[} \AgdaBound{A} \AgdaFunction{]trans*} \AgdaSymbol{\_)} \AgdaInductiveConstructor{,} \<[52]%
\>[52]\<%
\\
\>[0]\AgdaIndent{20}{}\<[20]%
\>[20]\AgdaSymbol{(}\AgdaFunction{[} \AgdaFunction{[} \AgdaBound{B} \AgdaFunction{]fm} \AgdaSymbol{\_} \AgdaFunction{]trans}\<%
\\
\>[0]\AgdaIndent{20}{}\<[20]%
\>[20]\AgdaSymbol{(}\AgdaFunction{[} \AgdaBound{B} \AgdaFunction{]trans*} \AgdaSymbol{\_)} \AgdaSymbol{(}\AgdaFunction{[} \AgdaBound{B} \AgdaFunction{]trans*} \AgdaSymbol{\_))} \AgdaSymbol{\}}\<%
\\
\>[0]\AgdaIndent{8}{}\<[8]%
\>[8]\AgdaSymbol{\}}\<%
\\
%
\\
\>\<\end{code}

Binary relation

\begin{code}\>\<%
\\
\>\AgdaFunction{Rel} \AgdaSymbol{:} \AgdaSymbol{\{}\AgdaBound{Γ} \AgdaSymbol{:} \AgdaFunction{Con}\AgdaSymbol{\}} \AgdaSymbol{→} \AgdaRecord{Ty} \AgdaBound{Γ} \AgdaSymbol{→} \AgdaPrimitiveType{Set₁}\<%
\\
\>\AgdaFunction{Rel} \AgdaSymbol{\{}\AgdaBound{Γ}\AgdaSymbol{\}} \AgdaBound{A} \AgdaSymbol{=} \AgdaRecord{Ty} \AgdaSymbol{(}\AgdaBound{Γ} \AgdaFunction{\&} \AgdaBound{A} \AgdaFunction{\&} \AgdaBound{A} \AgdaFunction{[} \AgdaFunction{fst\&} \AgdaSymbol{\{}A \AgdaSymbol{=} \AgdaBound{A}\AgdaSymbol{\}} \AgdaFunction{]T}\AgdaSymbol{)}\<%
\\
\>\<\end{code}

Natural numbers

Axiom: irrelevant:

\begin{code}\>\<%
\\
\>\AgdaKeyword{postulate}\<%
\\
\>[0]\AgdaIndent{4}{}\<[4]%
\>[4]\AgdaSymbol{.}\AgdaPostulate{irrelevant} \AgdaSymbol{:} \AgdaSymbol{\{}\AgdaBound{A} \AgdaSymbol{:} \AgdaPrimitiveType{Set}\AgdaSymbol{\}} \AgdaSymbol{→} \AgdaSymbol{.}\AgdaBound{A} \AgdaSymbol{→} \AgdaBound{A}\<%
\\
\>\<\end{code}


\begin{code}\>\<%
\\
\>\AgdaKeyword{module} \AgdaModule{Natural} \AgdaSymbol{(}\AgdaBound{Γ} \AgdaSymbol{:} \AgdaFunction{Con}\AgdaSymbol{)} \AgdaKeyword{where}\<%
\\
%
\\
\>[0]\AgdaIndent{2}{}\<[2]%
\>[2]\AgdaFunction{\_≈nat\_} \AgdaSymbol{:} \AgdaDatatype{ℕ} \AgdaSymbol{→} \AgdaDatatype{ℕ} \AgdaSymbol{→} \AgdaPrimitiveType{Set}\<%
\\
\>[0]\AgdaIndent{2}{}\<[2]%
\>[2]\AgdaInductiveConstructor{zero} \AgdaFunction{≈nat} \AgdaInductiveConstructor{zero} \AgdaSymbol{=} \AgdaRecord{⊤}\<%
\\
\>[0]\AgdaIndent{2}{}\<[2]%
\>[2]\AgdaInductiveConstructor{zero} \AgdaFunction{≈nat} \AgdaInductiveConstructor{suc} \AgdaBound{n} \AgdaSymbol{=} \AgdaDatatype{⊥}\<%
\\
\>[0]\AgdaIndent{2}{}\<[2]%
\>[2]\AgdaInductiveConstructor{suc} \AgdaBound{m} \AgdaFunction{≈nat} \AgdaInductiveConstructor{zero} \AgdaSymbol{=} \AgdaDatatype{⊥}\<%
\\
\>[0]\AgdaIndent{2}{}\<[2]%
\>[2]\AgdaInductiveConstructor{suc} \AgdaBound{m} \AgdaFunction{≈nat} \AgdaInductiveConstructor{suc} \AgdaBound{n} \AgdaSymbol{=} \AgdaBound{m} \AgdaFunction{≈nat} \AgdaBound{n}\<%
\\
\>[0]\AgdaIndent{2}{}\<[2]%
\>[2]\<%
\\
\>[0]\AgdaIndent{2}{}\<[2]%
\>[2]\AgdaFunction{reflNat} \AgdaSymbol{:} \AgdaSymbol{\{}\AgdaBound{x} \AgdaSymbol{:} \AgdaDatatype{ℕ}\AgdaSymbol{\}} \AgdaSymbol{→} \AgdaBound{x} \AgdaFunction{≈nat} \AgdaBound{x}\<%
\\
\>[0]\AgdaIndent{2}{}\<[2]%
\>[2]\AgdaFunction{reflNat} \AgdaSymbol{\{}\AgdaInductiveConstructor{zero}\AgdaSymbol{\}} \AgdaSymbol{=} \AgdaInductiveConstructor{tt}\<%
\\
\>[0]\AgdaIndent{2}{}\<[2]%
\>[2]\AgdaFunction{reflNat} \AgdaSymbol{\{}\AgdaInductiveConstructor{suc} \AgdaBound{n}\AgdaSymbol{\}} \AgdaSymbol{=} \AgdaFunction{reflNat} \AgdaSymbol{\{}\AgdaBound{n}\AgdaSymbol{\}}\<%
\\
%
\\
\>[0]\AgdaIndent{2}{}\<[2]%
\>[2]\AgdaFunction{symNat} \AgdaSymbol{:} \AgdaSymbol{\{}\AgdaBound{x} \AgdaBound{y} \AgdaSymbol{:} \AgdaDatatype{ℕ}\AgdaSymbol{\}} \AgdaSymbol{→} \AgdaBound{x} \AgdaFunction{≈nat} \AgdaBound{y} \AgdaSymbol{→} \AgdaBound{y} \AgdaFunction{≈nat} \AgdaBound{x}\<%
\\
\>[0]\AgdaIndent{2}{}\<[2]%
\>[2]\AgdaFunction{symNat} \AgdaSymbol{\{}\AgdaInductiveConstructor{zero}\AgdaSymbol{\}} \AgdaSymbol{\{}\AgdaInductiveConstructor{zero}\AgdaSymbol{\}} \AgdaBound{eq} \AgdaSymbol{=} \AgdaInductiveConstructor{tt}\<%
\\
\>[0]\AgdaIndent{2}{}\<[2]%
\>[2]\AgdaFunction{symNat} \AgdaSymbol{\{}\AgdaInductiveConstructor{zero}\AgdaSymbol{\}} \AgdaSymbol{\{}\AgdaInductiveConstructor{suc} \AgdaSymbol{\_\}} \AgdaBound{eq} \AgdaSymbol{=} \AgdaBound{eq}\<%
\\
\>[0]\AgdaIndent{2}{}\<[2]%
\>[2]\AgdaFunction{symNat} \AgdaSymbol{\{}\AgdaInductiveConstructor{suc} \AgdaSymbol{\_\}} \AgdaSymbol{\{}\AgdaInductiveConstructor{zero}\AgdaSymbol{\}} \AgdaBound{eq} \AgdaSymbol{=} \AgdaBound{eq}\<%
\\
\>[0]\AgdaIndent{2}{}\<[2]%
\>[2]\AgdaFunction{symNat} \AgdaSymbol{\{}\AgdaInductiveConstructor{suc} \AgdaBound{x}\AgdaSymbol{\}} \AgdaSymbol{\{}\AgdaInductiveConstructor{suc} \AgdaBound{y}\AgdaSymbol{\}} \AgdaBound{eq} \AgdaSymbol{=} \AgdaFunction{symNat} \AgdaSymbol{\{}\AgdaBound{x}\AgdaSymbol{\}} \AgdaSymbol{\{}\AgdaBound{y}\AgdaSymbol{\}} \AgdaBound{eq}\<%
\\
%
\\
\>[0]\AgdaIndent{2}{}\<[2]%
\>[2]\AgdaFunction{transNat} \AgdaSymbol{:} \AgdaSymbol{\{}\AgdaBound{x} \AgdaBound{y} \AgdaBound{z} \AgdaSymbol{:} \AgdaDatatype{ℕ}\AgdaSymbol{\}}\<%
\\
\>[2]\AgdaIndent{11}{}\<[11]%
\>[11]\AgdaSymbol{→} \AgdaBound{x} \AgdaFunction{≈nat} \AgdaBound{y} \AgdaSymbol{→} \AgdaBound{y} \AgdaFunction{≈nat} \AgdaBound{z} \AgdaSymbol{→} \AgdaBound{x} \AgdaFunction{≈nat} \AgdaBound{z}\<%
\\
\>[0]\AgdaIndent{2}{}\<[2]%
\>[2]\AgdaFunction{transNat} \AgdaSymbol{\{}\AgdaInductiveConstructor{zero}\AgdaSymbol{\}} \AgdaSymbol{\{}\AgdaInductiveConstructor{zero}\AgdaSymbol{\}} \AgdaBound{xy} \AgdaBound{yz} \AgdaSymbol{=} \AgdaBound{yz}\<%
\\
\>[0]\AgdaIndent{2}{}\<[2]%
\>[2]\AgdaFunction{transNat} \AgdaSymbol{\{}\AgdaInductiveConstructor{zero}\AgdaSymbol{\}} \AgdaSymbol{\{}\AgdaInductiveConstructor{suc} \AgdaSymbol{\_\}} \AgdaSymbol{()} \AgdaBound{yz}\<%
\\
\>[0]\AgdaIndent{2}{}\<[2]%
\>[2]\AgdaFunction{transNat} \AgdaSymbol{\{}\AgdaInductiveConstructor{suc} \AgdaSymbol{\_\}} \AgdaSymbol{\{}\AgdaInductiveConstructor{zero}\AgdaSymbol{\}} \AgdaSymbol{()} \AgdaBound{yz}\<%
\\
\>[0]\AgdaIndent{2}{}\<[2]%
\>[2]\AgdaFunction{transNat} \AgdaSymbol{\{}\AgdaInductiveConstructor{suc} \AgdaSymbol{\_\}} \AgdaSymbol{\{}\AgdaInductiveConstructor{suc} \AgdaSymbol{\_\}} \AgdaSymbol{\{}\AgdaInductiveConstructor{zero}\AgdaSymbol{\}} \AgdaBound{xy} \AgdaBound{yz} \AgdaSymbol{=} \AgdaBound{yz}\<%
\\
\>[0]\AgdaIndent{2}{}\<[2]%
\>[2]\AgdaFunction{transNat} \AgdaSymbol{\{}\AgdaInductiveConstructor{suc} \AgdaBound{x}\AgdaSymbol{\}} \AgdaSymbol{\{}\AgdaInductiveConstructor{suc} \AgdaBound{y}\AgdaSymbol{\}} \AgdaSymbol{\{}\AgdaInductiveConstructor{suc} \AgdaBound{z}\AgdaSymbol{\}} \AgdaBound{xy} \AgdaBound{yz} \AgdaSymbol{=}\<%
\\
\>[2]\AgdaIndent{11}{}\<[11]%
\>[11]\AgdaFunction{transNat} \AgdaSymbol{\{}\AgdaBound{x}\AgdaSymbol{\}} \AgdaSymbol{\{}\AgdaBound{y}\AgdaSymbol{\}} \AgdaSymbol{\{}\AgdaBound{z}\AgdaSymbol{\}} \AgdaBound{xy} \AgdaBound{yz}\<%
\\
\>[6]\AgdaIndent{3}{}\<[3]%
\>[3]\<%
\\
%
\\
\>[0]\AgdaIndent{2}{}\<[2]%
\>[2]\AgdaFunction{⟦Nat⟧} \AgdaSymbol{:} \AgdaRecord{Ty} \AgdaBound{Γ}\<%
\\
\>[0]\AgdaIndent{2}{}\<[2]%
\>[2]\AgdaFunction{⟦Nat⟧} \AgdaSymbol{=} \AgdaKeyword{record} \<[17]%
\>[17]\<%
\\
\>[2]\AgdaIndent{4}{}\<[4]%
\>[4]\AgdaSymbol{\{} \AgdaField{fm} \AgdaSymbol{=} \AgdaSymbol{λ} \AgdaBound{γ} \AgdaSymbol{→} \AgdaKeyword{record}\<%
\\
\>[4]\AgdaIndent{9}{}\<[9]%
\>[9]\AgdaSymbol{\{} \AgdaField{Carrier} \AgdaSymbol{=} \AgdaDatatype{ℕ}\<%
\\
\>[4]\AgdaIndent{9}{}\<[9]%
\>[9]\AgdaSymbol{;} \AgdaField{\_≈\_} \AgdaSymbol{=} \AgdaFunction{\_≈nat\_}\<%
\\
\>[4]\AgdaIndent{9}{}\<[9]%
\>[9]\AgdaSymbol{;} \AgdaField{refl} \AgdaSymbol{=} \AgdaSymbol{λ} \AgdaSymbol{\{}\AgdaBound{n}\AgdaSymbol{\}} \AgdaSymbol{→} \AgdaFunction{reflNat} \AgdaSymbol{\{}\AgdaBound{n}\AgdaSymbol{\}}\<%
\\
\>[4]\AgdaIndent{9}{}\<[9]%
\>[9]\AgdaSymbol{;} \AgdaField{sym} \AgdaSymbol{=} \AgdaSymbol{λ} \AgdaSymbol{\{}\AgdaBound{x}\AgdaSymbol{\}} \AgdaSymbol{\{}\AgdaBound{y}\AgdaSymbol{\}} \AgdaSymbol{→} \AgdaFunction{symNat} \AgdaSymbol{\{}\AgdaBound{x}\AgdaSymbol{\}} \AgdaSymbol{\{}\AgdaBound{y}\AgdaSymbol{\}}\<%
\\
\>[4]\AgdaIndent{9}{}\<[9]%
\>[9]\AgdaSymbol{;} \AgdaField{trans} \AgdaSymbol{=} \AgdaSymbol{λ} \AgdaSymbol{\{}\AgdaBound{x}\AgdaSymbol{\}} \AgdaSymbol{\{}\AgdaBound{y}\AgdaSymbol{\}} \AgdaSymbol{\{}\AgdaBound{z}\AgdaSymbol{\}} \AgdaSymbol{→} \AgdaFunction{transNat} \AgdaSymbol{\{}\AgdaBound{x}\AgdaSymbol{\}} \AgdaSymbol{\{}\AgdaBound{y}\AgdaSymbol{\}} \AgdaSymbol{\{}\AgdaBound{z}\AgdaSymbol{\}}\<%
\\
\>[4]\AgdaIndent{9}{}\<[9]%
\>[9]\AgdaSymbol{\}}\<%
\\
\>[0]\AgdaIndent{4}{}\<[4]%
\>[4]\AgdaSymbol{;} \AgdaField{substT} \AgdaSymbol{=} \AgdaSymbol{λ} \AgdaBound{\_} \AgdaBound{n} \AgdaSymbol{→} \AgdaBound{n}\<%
\\
\>[0]\AgdaIndent{4}{}\<[4]%
\>[4]\AgdaSymbol{;} \AgdaField{subst*} \AgdaSymbol{=} \AgdaSymbol{λ} \AgdaBound{\_} \AgdaBound{x} \AgdaSymbol{→} \AgdaPostulate{irrelevant} \AgdaBound{x}\<%
\\
\>[0]\AgdaIndent{4}{}\<[4]%
\>[4]\AgdaSymbol{;} \AgdaField{refl*} \AgdaSymbol{=} \AgdaSymbol{λ} \AgdaSymbol{\{}\AgdaBound{x}\AgdaSymbol{\}} \AgdaSymbol{\{}\AgdaBound{a}\AgdaSymbol{\}} \AgdaSymbol{→} \AgdaFunction{reflNat} \AgdaSymbol{\{}\AgdaBound{a}\AgdaSymbol{\}}\<%
\\
\>[0]\AgdaIndent{4}{}\<[4]%
\>[4]\AgdaSymbol{;} \AgdaField{trans*} \AgdaSymbol{=} \AgdaSymbol{λ} \AgdaBound{a} \AgdaSymbol{→} \AgdaFunction{reflNat} \AgdaSymbol{\{}\AgdaBound{a}\AgdaSymbol{\}} \<[33]%
\>[33]\<%
\\
\>[0]\AgdaIndent{4}{}\<[4]%
\>[4]\AgdaSymbol{\}}\<%
\\
%
\\
\>[0]\AgdaIndent{2}{}\<[2]%
\>[2]\AgdaFunction{⟦0⟧} \AgdaSymbol{:} \AgdaRecord{Tm} \AgdaFunction{⟦Nat⟧}\<%
\\
\>[0]\AgdaIndent{2}{}\<[2]%
\>[2]\AgdaFunction{⟦0⟧} \AgdaSymbol{=} \AgdaKeyword{record}\<%
\\
\>[2]\AgdaIndent{6}{}\<[6]%
\>[6]\AgdaSymbol{\{} \AgdaField{tm} \AgdaSymbol{=} \AgdaSymbol{λ} \AgdaBound{\_} \AgdaSymbol{→} \AgdaNumber{0}\<%
\\
\>[2]\AgdaIndent{6}{}\<[6]%
\>[6]\AgdaSymbol{;} \AgdaField{respt} \AgdaSymbol{=} \AgdaSymbol{λ} \AgdaBound{p} \AgdaSymbol{→} \AgdaInductiveConstructor{tt}\<%
\\
\>[2]\AgdaIndent{6}{}\<[6]%
\>[6]\AgdaSymbol{\}}\<%
\\
%
\\
\>[0]\AgdaIndent{2}{}\<[2]%
\>[2]\AgdaFunction{⟦s⟧} \AgdaSymbol{:} \AgdaRecord{Tm} \AgdaFunction{⟦Nat⟧} \AgdaSymbol{→} \AgdaRecord{Tm} \AgdaFunction{⟦Nat⟧}\<%
\\
\>[0]\AgdaIndent{2}{}\<[2]%
\>[2]\AgdaFunction{⟦s⟧} \AgdaSymbol{(}\AgdaInductiveConstructor{tm:} \AgdaBound{t} \AgdaInductiveConstructor{resp:} \AgdaBound{respt}\AgdaSymbol{)} \<[26]%
\>[26]\<%
\\
\>[2]\AgdaIndent{6}{}\<[6]%
\>[6]\AgdaSymbol{=} \AgdaKeyword{record}\<%
\\
\>[2]\AgdaIndent{6}{}\<[6]%
\>[6]\AgdaSymbol{\{} \AgdaField{tm} \AgdaSymbol{=} \AgdaInductiveConstructor{suc} \AgdaFunction{∘} \AgdaBound{t}\<%
\\
\>[2]\AgdaIndent{6}{}\<[6]%
\>[6]\AgdaSymbol{;} \AgdaField{respt} \AgdaSymbol{=} \AgdaBound{respt}\<%
\\
\>[2]\AgdaIndent{6}{}\<[6]%
\>[6]\AgdaSymbol{\}}\<%
\\
\>\<\end{code}

Simply typed universe

\AgdaHide{
\begin{code}\>\<%
\\
\>[0]\AgdaIndent{2}{}\<[2]%
\>[2]\AgdaKeyword{data} \AgdaDatatype{⟦U⟧⁰} \AgdaSymbol{:} \AgdaPrimitiveType{Set} \AgdaKeyword{where}\<%
\\
\>[0]\AgdaIndent{4}{}\<[4]%
\>[4]\AgdaInductiveConstructor{nat} \AgdaSymbol{:} \AgdaDatatype{⟦U⟧⁰}\<%
\\
\>[0]\AgdaIndent{4}{}\<[4]%
\>[4]\AgdaInductiveConstructor{arr<\_,\_>} \AgdaSymbol{:} \AgdaSymbol{(}\AgdaBound{a} \AgdaBound{b} \AgdaSymbol{:} \AgdaDatatype{⟦U⟧⁰}\AgdaSymbol{)} \AgdaSymbol{→} \AgdaDatatype{⟦U⟧⁰}\<%
\\
%
\\
\>[0]\AgdaIndent{2}{}\<[2]%
\>[2]\AgdaFunction{\_\textasciitilde⟦U⟧\_} \AgdaSymbol{:} \AgdaDatatype{⟦U⟧⁰} \AgdaSymbol{→} \AgdaDatatype{⟦U⟧⁰} \AgdaSymbol{→} \AgdaPrimitiveType{Set} \AgdaComment{-- HProp}\<%
\\
\>[0]\AgdaIndent{2}{}\<[2]%
\>[2]\AgdaInductiveConstructor{nat} \AgdaFunction{\textasciitilde⟦U⟧} \AgdaInductiveConstructor{nat} \AgdaSymbol{=} \AgdaRecord{⊤}\<%
\\
\>[0]\AgdaIndent{2}{}\<[2]%
\>[2]\AgdaInductiveConstructor{nat} \AgdaFunction{\textasciitilde⟦U⟧} \AgdaInductiveConstructor{arr<} \AgdaBound{a} \AgdaInductiveConstructor{,} \AgdaBound{b} \AgdaInductiveConstructor{>} \AgdaSymbol{=} \AgdaDatatype{⊥}\<%
\\
\>[0]\AgdaIndent{2}{}\<[2]%
\>[2]\AgdaInductiveConstructor{arr<} \AgdaBound{a} \AgdaInductiveConstructor{,} \AgdaBound{b} \AgdaInductiveConstructor{>} \AgdaFunction{\textasciitilde⟦U⟧} \AgdaInductiveConstructor{nat} \AgdaSymbol{=} \AgdaDatatype{⊥}\<%
\\
\>[0]\AgdaIndent{2}{}\<[2]%
\>[2]\AgdaInductiveConstructor{arr<} \AgdaBound{a} \AgdaInductiveConstructor{,} \AgdaBound{b} \AgdaInductiveConstructor{>} \AgdaFunction{\textasciitilde⟦U⟧} \AgdaInductiveConstructor{arr<} \AgdaBound{a'} \AgdaInductiveConstructor{,} \AgdaBound{b'} \AgdaInductiveConstructor{>} \AgdaSymbol{=} \AgdaBound{a} \AgdaFunction{\textasciitilde⟦U⟧} \AgdaBound{a'} \AgdaFunction{×} \AgdaBound{b} \AgdaFunction{\textasciitilde⟦U⟧} \AgdaBound{b'}\<%
\\
%
\\
\>[0]\AgdaIndent{2}{}\<[2]%
\>[2]\AgdaFunction{reflU} \AgdaSymbol{:} \<[11]%
\>[11]\AgdaSymbol{\{}\AgdaBound{x} \AgdaSymbol{:} \AgdaDatatype{⟦U⟧⁰}\AgdaSymbol{\}} \AgdaSymbol{→} \AgdaBound{x} \AgdaFunction{\textasciitilde⟦U⟧} \AgdaBound{x}\<%
\\
\>[0]\AgdaIndent{2}{}\<[2]%
\>[2]\AgdaFunction{reflU} \AgdaSymbol{\{}\AgdaInductiveConstructor{nat}\AgdaSymbol{\}} \AgdaSymbol{=} \AgdaInductiveConstructor{tt}\<%
\\
\>[0]\AgdaIndent{2}{}\<[2]%
\>[2]\AgdaFunction{reflU} \AgdaSymbol{\{}\AgdaInductiveConstructor{arr<} \AgdaBound{a} \AgdaInductiveConstructor{,} \AgdaBound{b} \AgdaInductiveConstructor{>}\AgdaSymbol{\}} \AgdaSymbol{=} \AgdaFunction{reflU} \AgdaSymbol{\{}\AgdaBound{a}\AgdaSymbol{\}} \AgdaInductiveConstructor{,} \AgdaFunction{reflU} \AgdaSymbol{\{}\AgdaBound{b}\AgdaSymbol{\}}\<%
\\
%
\\
\>[0]\AgdaIndent{2}{}\<[2]%
\>[2]\AgdaFunction{symU} \AgdaSymbol{:} \AgdaSymbol{\{}\AgdaBound{x} \AgdaBound{y} \AgdaSymbol{:} \AgdaDatatype{⟦U⟧⁰}\AgdaSymbol{\}} \AgdaSymbol{→} \AgdaBound{x} \AgdaFunction{\textasciitilde⟦U⟧} \AgdaBound{y} \AgdaSymbol{→} \AgdaBound{y} \AgdaFunction{\textasciitilde⟦U⟧} \AgdaBound{x}\<%
\\
\>[0]\AgdaIndent{2}{}\<[2]%
\>[2]\AgdaFunction{symU} \AgdaSymbol{\{}\AgdaInductiveConstructor{nat}\AgdaSymbol{\}} \AgdaSymbol{\{}\AgdaInductiveConstructor{nat}\AgdaSymbol{\}} \AgdaBound{eq} \AgdaSymbol{=} \AgdaInductiveConstructor{tt}\<%
\\
\>[0]\AgdaIndent{2}{}\<[2]%
\>[2]\AgdaFunction{symU} \AgdaSymbol{\{}\AgdaInductiveConstructor{nat}\AgdaSymbol{\}} \AgdaSymbol{\{}\AgdaInductiveConstructor{arr<} \AgdaBound{a} \AgdaInductiveConstructor{,} \AgdaBound{b} \AgdaInductiveConstructor{>}\AgdaSymbol{\}} \AgdaBound{eq} \AgdaSymbol{=} \AgdaBound{eq}\<%
\\
\>[0]\AgdaIndent{2}{}\<[2]%
\>[2]\AgdaFunction{symU} \AgdaSymbol{\{}\AgdaInductiveConstructor{arr<} \AgdaBound{a} \AgdaInductiveConstructor{,} \AgdaBound{b} \AgdaInductiveConstructor{>}\AgdaSymbol{\}} \AgdaSymbol{\{}\AgdaInductiveConstructor{nat}\AgdaSymbol{\}} \AgdaBound{eq} \AgdaSymbol{=} \AgdaBound{eq}\<%
\\
\>[0]\AgdaIndent{2}{}\<[2]%
\>[2]\AgdaFunction{symU} \AgdaSymbol{\{}\AgdaInductiveConstructor{arr<} \AgdaBound{a} \AgdaInductiveConstructor{,} \AgdaBound{b} \AgdaInductiveConstructor{>}\AgdaSymbol{\}} \AgdaSymbol{\{}\AgdaInductiveConstructor{arr<} \AgdaBound{a'} \AgdaInductiveConstructor{,} \AgdaBound{b'} \AgdaInductiveConstructor{>}\AgdaSymbol{\}} \AgdaSymbol{(}\AgdaBound{p} \AgdaInductiveConstructor{,} \AgdaBound{q}\AgdaSymbol{)} \AgdaSymbol{=} \AgdaSymbol{(}\AgdaFunction{symU} \AgdaSymbol{\{}\AgdaBound{a}\AgdaSymbol{\}} \AgdaSymbol{\{}\AgdaBound{a'}\AgdaSymbol{\}} \AgdaBound{p}\AgdaSymbol{)} \<[67]%
\>[67]\<%
\\
\>[2]\AgdaIndent{47}{}\<[47]%
\>[47]\AgdaInductiveConstructor{,} \AgdaSymbol{(}\AgdaFunction{symU} \AgdaSymbol{\{}\AgdaBound{b}\AgdaSymbol{\}} \AgdaSymbol{\{}\AgdaBound{b'}\AgdaSymbol{\}} \AgdaBound{q}\AgdaSymbol{)}\<%
\\
%
\\
\>[0]\AgdaIndent{2}{}\<[2]%
\>[2]\AgdaFunction{transU} \AgdaSymbol{:} \AgdaSymbol{\{}\AgdaBound{x} \AgdaBound{y} \AgdaBound{z} \AgdaSymbol{:} \AgdaDatatype{⟦U⟧⁰}\AgdaSymbol{\}} \AgdaSymbol{→} \AgdaBound{x} \AgdaFunction{\textasciitilde⟦U⟧} \AgdaBound{y} \AgdaSymbol{→} \AgdaBound{y} \AgdaFunction{\textasciitilde⟦U⟧} \AgdaBound{z} \AgdaSymbol{→} \AgdaBound{x} \AgdaFunction{\textasciitilde⟦U⟧} \AgdaBound{z}\<%
\\
\>[0]\AgdaIndent{2}{}\<[2]%
\>[2]\AgdaFunction{transU} \AgdaSymbol{\{}\AgdaInductiveConstructor{nat}\AgdaSymbol{\}} \AgdaSymbol{\{}\AgdaInductiveConstructor{nat}\AgdaSymbol{\}} \AgdaBound{eq1} \AgdaBound{eq2} \AgdaSymbol{=} \AgdaBound{eq2}\<%
\\
\>[0]\AgdaIndent{2}{}\<[2]%
\>[2]\AgdaFunction{transU} \AgdaSymbol{\{}\AgdaInductiveConstructor{nat}\AgdaSymbol{\}} \AgdaSymbol{\{}\AgdaInductiveConstructor{arr<} \AgdaBound{a} \AgdaInductiveConstructor{,} \AgdaBound{b} \AgdaInductiveConstructor{>}\AgdaSymbol{\}} \AgdaSymbol{()} \AgdaBound{eq2}\<%
\\
\>[0]\AgdaIndent{2}{}\<[2]%
\>[2]\AgdaFunction{transU} \AgdaSymbol{\{}\AgdaInductiveConstructor{arr<} \AgdaBound{a} \AgdaInductiveConstructor{,} \AgdaBound{b} \AgdaInductiveConstructor{>}\AgdaSymbol{\}} \AgdaSymbol{\{}\AgdaInductiveConstructor{nat}\AgdaSymbol{\}} \AgdaSymbol{()} \AgdaBound{eq2}\<%
\\
\>[0]\AgdaIndent{2}{}\<[2]%
\>[2]\AgdaFunction{transU} \AgdaSymbol{\{}\AgdaInductiveConstructor{arr<} \AgdaBound{a} \AgdaInductiveConstructor{,} \AgdaBound{b} \AgdaInductiveConstructor{>}\AgdaSymbol{\}} \AgdaSymbol{\{}\AgdaInductiveConstructor{arr<} \AgdaBound{a'} \AgdaInductiveConstructor{,} \AgdaBound{b'} \AgdaInductiveConstructor{>}\AgdaSymbol{\}} \AgdaSymbol{\{}\AgdaInductiveConstructor{nat}\AgdaSymbol{\}} \AgdaBound{eq1} \AgdaBound{eq2} \AgdaSymbol{=} \AgdaBound{eq2}\<%
\\
\>[0]\AgdaIndent{2}{}\<[2]%
\>[2]\AgdaFunction{transU} \AgdaSymbol{\{}\AgdaInductiveConstructor{arr<} \AgdaBound{a} \AgdaInductiveConstructor{,} \AgdaBound{b} \AgdaInductiveConstructor{>}\AgdaSymbol{\}} \AgdaSymbol{\{}\AgdaInductiveConstructor{arr<} \AgdaBound{a'} \AgdaInductiveConstructor{,} \AgdaBound{b'} \AgdaInductiveConstructor{>}\AgdaSymbol{\}} \AgdaSymbol{\{}\AgdaInductiveConstructor{arr<} \AgdaBound{a0} \AgdaInductiveConstructor{,} \AgdaBound{b0} \AgdaInductiveConstructor{>}\AgdaSymbol{\}} \AgdaSymbol{(}\AgdaBound{p1} \AgdaInductiveConstructor{,} \AgdaBound{q1}\AgdaSymbol{)} \<[68]%
\>[68]\<%
\\
\>[2]\AgdaIndent{9}{}\<[9]%
\>[9]\AgdaSymbol{(}\AgdaBound{p2} \AgdaInductiveConstructor{,} \AgdaBound{q2}\AgdaSymbol{)} \AgdaSymbol{=} \AgdaSymbol{(}\AgdaFunction{transU} \AgdaSymbol{\{}\AgdaBound{a}\AgdaSymbol{\}} \AgdaSymbol{\{}\AgdaBound{a'}\AgdaSymbol{\}} \AgdaSymbol{\{}\AgdaBound{a0}\AgdaSymbol{\}} \AgdaBound{p1} \AgdaBound{p2}\AgdaSymbol{)} \<[50]%
\>[50]\<%
\\
\>[2]\AgdaIndent{9}{}\<[9]%
\>[9]\AgdaInductiveConstructor{,} \AgdaFunction{transU} \AgdaSymbol{\{}\AgdaBound{b}\AgdaSymbol{\}} \AgdaSymbol{\{}\AgdaBound{b'}\AgdaSymbol{\}} \AgdaSymbol{\{}\AgdaBound{b0}\AgdaSymbol{\}} \AgdaBound{q1} \AgdaBound{q2}\<%
\\
%
\\
\>[0]\AgdaIndent{2}{}\<[2]%
\>[2]\AgdaFunction{⟦U⟧} \AgdaSymbol{:} \AgdaRecord{Ty} \AgdaBound{Γ}\<%
\\
\>[0]\AgdaIndent{2}{}\<[2]%
\>[2]\AgdaFunction{⟦U⟧} \AgdaSymbol{=} \AgdaKeyword{record} \<[15]%
\>[15]\<%
\\
\>[2]\AgdaIndent{4}{}\<[4]%
\>[4]\AgdaSymbol{\{} \AgdaField{fm} \AgdaSymbol{=} \AgdaSymbol{λ} \AgdaBound{γ} \AgdaSymbol{→} \AgdaKeyword{record}\<%
\\
\>[4]\AgdaIndent{9}{}\<[9]%
\>[9]\AgdaSymbol{\{} \AgdaField{Carrier} \AgdaSymbol{=} \AgdaDatatype{⟦U⟧⁰}\<%
\\
\>[4]\AgdaIndent{9}{}\<[9]%
\>[9]\AgdaSymbol{;} \AgdaField{\_≈\_} \AgdaSymbol{=} \AgdaFunction{\_\textasciitilde⟦U⟧\_}\<%
\\
\>[4]\AgdaIndent{9}{}\<[9]%
\>[9]\AgdaSymbol{;} \AgdaField{refl} \AgdaSymbol{=} \AgdaSymbol{λ} \AgdaSymbol{\{}\AgdaBound{x}\AgdaSymbol{\}} \AgdaSymbol{→} \AgdaFunction{reflU} \AgdaSymbol{\{}\AgdaBound{x}\AgdaSymbol{\}}\<%
\\
\>[4]\AgdaIndent{9}{}\<[9]%
\>[9]\AgdaSymbol{;} \AgdaField{sym} \AgdaSymbol{=} \AgdaSymbol{λ} \AgdaSymbol{\{}\AgdaBound{x}\AgdaSymbol{\}} \AgdaSymbol{\{}\AgdaBound{y}\AgdaSymbol{\}} \AgdaSymbol{→} \AgdaFunction{symU} \AgdaSymbol{\{}\AgdaBound{x}\AgdaSymbol{\}} \AgdaSymbol{\{}\AgdaBound{y}\AgdaSymbol{\}}\<%
\\
\>[4]\AgdaIndent{9}{}\<[9]%
\>[9]\AgdaSymbol{;} \AgdaField{trans} \AgdaSymbol{=} \AgdaSymbol{λ} \AgdaSymbol{\{}\AgdaBound{x}\AgdaSymbol{\}} \AgdaSymbol{\{}\AgdaBound{y}\AgdaSymbol{\}} \AgdaSymbol{\{}\AgdaBound{z}\AgdaSymbol{\}} \AgdaSymbol{→} \AgdaFunction{transU} \AgdaSymbol{\{}\AgdaBound{x}\AgdaSymbol{\}} \AgdaSymbol{\{}\AgdaBound{y}\AgdaSymbol{\}} \AgdaSymbol{\{}\AgdaBound{z}\AgdaSymbol{\}}\<%
\\
\>[4]\AgdaIndent{9}{}\<[9]%
\>[9]\AgdaSymbol{\}}\<%
\\
\>[0]\AgdaIndent{4}{}\<[4]%
\>[4]\AgdaSymbol{;} \AgdaField{substT} \AgdaSymbol{=} \AgdaSymbol{λ} \AgdaBound{\_} \AgdaBound{x} \AgdaSymbol{→} \AgdaBound{x}\<%
\\
\>[0]\AgdaIndent{4}{}\<[4]%
\>[4]\AgdaSymbol{;} \AgdaField{subst*} \AgdaSymbol{=} \<[16]%
\>[16]\AgdaSymbol{λ} \AgdaBound{\_} \AgdaBound{x} \AgdaSymbol{→} \AgdaPostulate{irrelevant} \AgdaBound{x}\<%
\\
\>[0]\AgdaIndent{4}{}\<[4]%
\>[4]\AgdaSymbol{;} \AgdaField{refl*} \AgdaSymbol{=} \AgdaSymbol{λ} \AgdaSymbol{\{}\AgdaBound{x}\AgdaSymbol{\}} \AgdaSymbol{\{}\AgdaBound{a}\AgdaSymbol{\}} \AgdaSymbol{→} \AgdaFunction{reflU} \AgdaSymbol{\{}\AgdaBound{a}\AgdaSymbol{\}}\<%
\\
\>[0]\AgdaIndent{4}{}\<[4]%
\>[4]\AgdaSymbol{;} \AgdaField{trans*} \AgdaSymbol{=} \AgdaSymbol{λ} \AgdaBound{a} \AgdaSymbol{→} \AgdaFunction{reflU} \AgdaSymbol{\{}\AgdaBound{a}\AgdaSymbol{\}}\<%
\\
\>[0]\AgdaIndent{4}{}\<[4]%
\>[4]\AgdaSymbol{\}}\<%
\\
%
\\
\>[0]\AgdaIndent{2}{}\<[2]%
\>[2]\AgdaFunction{elfm} \AgdaSymbol{:} \AgdaRecord{Σ} \AgdaFunction{∣} \AgdaBound{Γ} \AgdaFunction{∣} \AgdaSymbol{(λ} \AgdaBound{x} \AgdaSymbol{→} \AgdaDatatype{⟦U⟧⁰}\AgdaSymbol{)} \AgdaSymbol{→} \AgdaRecord{Setoid}\<%
\\
\>[0]\AgdaIndent{2}{}\<[2]%
\>[2]\AgdaFunction{elfm} \AgdaSymbol{(}\AgdaBound{γ} \AgdaInductiveConstructor{,} \AgdaInductiveConstructor{nat}\AgdaSymbol{)} \AgdaSymbol{=} \AgdaFunction{[} \AgdaFunction{⟦Nat⟧} \AgdaFunction{]fm} \AgdaBound{γ}\<%
\\
\>[0]\AgdaIndent{2}{}\<[2]%
\>[2]\AgdaFunction{elfm} \AgdaSymbol{(}\AgdaBound{γ} \AgdaInductiveConstructor{,} \AgdaInductiveConstructor{arr<} \AgdaBound{a} \AgdaInductiveConstructor{,} \AgdaBound{b} \AgdaInductiveConstructor{>}\AgdaSymbol{)} \AgdaSymbol{=} \AgdaFunction{[} \AgdaBound{Γ} \AgdaFunction{,} \AgdaBound{γ} \AgdaFunction{]} \AgdaFunction{elfm} \AgdaSymbol{(}\AgdaBound{γ} \AgdaInductiveConstructor{,} \AgdaBound{a}\AgdaSymbol{)} \AgdaFunction{⇒fm} \AgdaFunction{elfm} \AgdaSymbol{(}\AgdaBound{γ} \AgdaInductiveConstructor{,} \AgdaBound{b}\AgdaSymbol{)}\<%
\\
\>\<\end{code}
}

\section{What we can do in this model}

\section{Examples of types}

\AgdaHide{
\begin{code}\>\<%
\\
%
\\
\>\AgdaSymbol{\{-\#} \AgdaKeyword{OPTIONS} --type-in-type \AgdaSymbol{\#-\}}\<%
\\
%
\\
\>\AgdaKeyword{import} \AgdaModule{Level}\<%
\\
\>\AgdaKeyword{open} \AgdaKeyword{import} \AgdaModule{Relation.Binary.PropositionalEquality} \AgdaSymbol{as} \AgdaModule{PE} \AgdaKeyword{hiding} \AgdaSymbol{(}refl \AgdaSymbol{;} sym \AgdaSymbol{;} trans\AgdaSymbol{;} isEquivalence\AgdaSymbol{;} [\_]\AgdaSymbol{)}\<%
\\
%
\\
\>\AgdaKeyword{module} \AgdaModule{CwF-ctd} \AgdaSymbol{(}\AgdaBound{ext} \AgdaSymbol{:} \AgdaFunction{Extensionality} \AgdaPrimitive{Level.zero} \AgdaPrimitive{Level.zero}\AgdaSymbol{)} \AgdaKeyword{where}\<%
\\
%
\\
\>\AgdaKeyword{open} \AgdaKeyword{import} \AgdaModule{Data.Unit}\<%
\\
\>\AgdaKeyword{open} \AgdaKeyword{import} \AgdaModule{Function}\<%
\\
\>\AgdaKeyword{open} \AgdaKeyword{import} \AgdaModule{Data.Product}\<%
\\
%
\\
\>\AgdaKeyword{open} \AgdaKeyword{import} \AgdaModule{CwF-setoidwo} \AgdaBound{ext} \AgdaKeyword{public}\<%
\\
%
\\
\>\AgdaKeyword{open} \AgdaKeyword{import} \AgdaModule{Data.Nat}\<%
\\
%
\\
\>\<\end{code}
}

Binary relation

\begin{code}\>\<%
\\
\>\AgdaFunction{Rel} \AgdaSymbol{:} \AgdaSymbol{\{}\AgdaBound{Γ} \AgdaSymbol{:} \AgdaFunction{Con}\AgdaSymbol{\}} \AgdaSymbol{→} \AgdaRecord{Ty} \AgdaBound{Γ} \AgdaSymbol{→} \AgdaPrimitiveType{Set₁}\<%
\\
\>\AgdaFunction{Rel} \AgdaSymbol{\{}\AgdaBound{Γ}\AgdaSymbol{\}} \AgdaBound{A} \AgdaSymbol{=} \AgdaRecord{Ty} \AgdaSymbol{(}\AgdaBound{Γ} \AgdaFunction{\&} \AgdaBound{A} \AgdaFunction{\&} \AgdaBound{A} \AgdaFunction{[} \AgdaFunction{fst\&} \AgdaSymbol{\{}A \AgdaSymbol{=} \AgdaBound{A}\AgdaSymbol{\}} \AgdaFunction{]T}\AgdaSymbol{)}\<%
\\
\>\<\end{code}

Natural numbers

\begin{code}\>\<%
\\
\>\AgdaKeyword{module} \AgdaModule{Natural} \AgdaSymbol{(}\AgdaBound{Γ} \AgdaSymbol{:} \AgdaFunction{Con}\AgdaSymbol{)} \AgdaKeyword{where}\<%
\\
%
\\
\>[0]\AgdaIndent{2}{}\<[2]%
\>[2]\AgdaFunction{\_≈nat\_} \AgdaSymbol{:} \AgdaDatatype{ℕ} \AgdaSymbol{→} \AgdaDatatype{ℕ} \AgdaSymbol{→} \AgdaRecord{HProp}\<%
\\
\>[0]\AgdaIndent{2}{}\<[2]%
\>[2]\AgdaInductiveConstructor{zero} \AgdaFunction{≈nat} \AgdaInductiveConstructor{zero} \AgdaSymbol{=} \AgdaFunction{⊤'}\<%
\\
\>[0]\AgdaIndent{2}{}\<[2]%
\>[2]\AgdaInductiveConstructor{zero} \AgdaFunction{≈nat} \AgdaInductiveConstructor{suc} \AgdaBound{n} \AgdaSymbol{=} \AgdaFunction{⊥'}\<%
\\
\>[0]\AgdaIndent{2}{}\<[2]%
\>[2]\AgdaInductiveConstructor{suc} \AgdaBound{m} \AgdaFunction{≈nat} \AgdaInductiveConstructor{zero} \AgdaSymbol{=} \AgdaFunction{⊥'}\<%
\\
\>[0]\AgdaIndent{2}{}\<[2]%
\>[2]\AgdaInductiveConstructor{suc} \AgdaBound{m} \AgdaFunction{≈nat} \AgdaInductiveConstructor{suc} \AgdaBound{n} \AgdaSymbol{=} \AgdaBound{m} \AgdaFunction{≈nat} \AgdaBound{n}\<%
\\
\>[0]\AgdaIndent{2}{}\<[2]%
\>[2]\<%
\\
\>[0]\AgdaIndent{2}{}\<[2]%
\>[2]\AgdaFunction{reflNat} \AgdaSymbol{:} \AgdaSymbol{\{}\AgdaBound{x} \AgdaSymbol{:} \AgdaDatatype{ℕ}\AgdaSymbol{\}} \AgdaSymbol{→} \AgdaFunction{<} \AgdaBound{x} \AgdaFunction{≈nat} \AgdaBound{x} \AgdaFunction{>} \<[35]%
\>[35]\<%
\\
\>[0]\AgdaIndent{2}{}\<[2]%
\>[2]\AgdaFunction{reflNat} \AgdaSymbol{\{}\AgdaInductiveConstructor{zero}\AgdaSymbol{\}} \AgdaSymbol{=} \AgdaInductiveConstructor{tt}\<%
\\
\>[0]\AgdaIndent{2}{}\<[2]%
\>[2]\AgdaFunction{reflNat} \AgdaSymbol{\{}\AgdaInductiveConstructor{suc} \AgdaBound{n}\AgdaSymbol{\}} \AgdaSymbol{=} \AgdaFunction{reflNat} \AgdaSymbol{\{}\AgdaBound{n}\AgdaSymbol{\}}\<%
\\
%
\\
\>[0]\AgdaIndent{2}{}\<[2]%
\>[2]\AgdaFunction{symNat} \AgdaSymbol{:} \AgdaSymbol{\{}\AgdaBound{x} \AgdaBound{y} \AgdaSymbol{:} \AgdaDatatype{ℕ}\AgdaSymbol{\}} \AgdaSymbol{→} \AgdaFunction{<} \AgdaBound{x} \AgdaFunction{≈nat} \AgdaBound{y} \AgdaFunction{>} \AgdaSymbol{→} \AgdaFunction{<} \AgdaBound{y} \AgdaFunction{≈nat} \AgdaBound{x} \AgdaFunction{>}\<%
\\
\>[0]\AgdaIndent{2}{}\<[2]%
\>[2]\AgdaFunction{symNat} \AgdaSymbol{\{}\AgdaInductiveConstructor{zero}\AgdaSymbol{\}} \AgdaSymbol{\{}\AgdaInductiveConstructor{zero}\AgdaSymbol{\}} \AgdaBound{eq} \AgdaSymbol{=} \AgdaInductiveConstructor{tt}\<%
\\
\>[0]\AgdaIndent{2}{}\<[2]%
\>[2]\AgdaFunction{symNat} \AgdaSymbol{\{}\AgdaInductiveConstructor{zero}\AgdaSymbol{\}} \AgdaSymbol{\{}\AgdaInductiveConstructor{suc} \AgdaSymbol{\_\}} \AgdaBound{eq} \AgdaSymbol{=} \AgdaBound{eq}\<%
\\
\>[0]\AgdaIndent{2}{}\<[2]%
\>[2]\AgdaFunction{symNat} \AgdaSymbol{\{}\AgdaInductiveConstructor{suc} \AgdaSymbol{\_\}} \AgdaSymbol{\{}\AgdaInductiveConstructor{zero}\AgdaSymbol{\}} \AgdaBound{eq} \AgdaSymbol{=} \AgdaBound{eq}\<%
\\
\>[0]\AgdaIndent{2}{}\<[2]%
\>[2]\AgdaFunction{symNat} \AgdaSymbol{\{}\AgdaInductiveConstructor{suc} \AgdaBound{x}\AgdaSymbol{\}} \AgdaSymbol{\{}\AgdaInductiveConstructor{suc} \AgdaBound{y}\AgdaSymbol{\}} \AgdaBound{eq} \AgdaSymbol{=} \AgdaFunction{symNat} \AgdaSymbol{\{}\AgdaBound{x}\AgdaSymbol{\}} \AgdaSymbol{\{}\AgdaBound{y}\AgdaSymbol{\}} \AgdaBound{eq}\<%
\\
%
\\
\>[0]\AgdaIndent{2}{}\<[2]%
\>[2]\AgdaFunction{transNat} \AgdaSymbol{:} \AgdaSymbol{\{}\AgdaBound{x} \AgdaBound{y} \AgdaBound{z} \AgdaSymbol{:} \AgdaDatatype{ℕ}\AgdaSymbol{\}} \AgdaSymbol{→} \AgdaFunction{<} \AgdaBound{x} \AgdaFunction{≈nat} \AgdaBound{y} \AgdaFunction{>} \AgdaSymbol{→} \AgdaFunction{<} \AgdaBound{y} \AgdaFunction{≈nat} \AgdaBound{z} \AgdaFunction{>} \AgdaSymbol{→} \AgdaFunction{<} \AgdaBound{x} \AgdaFunction{≈nat} \AgdaBound{z} \AgdaFunction{>}\<%
\\
\>[0]\AgdaIndent{2}{}\<[2]%
\>[2]\AgdaFunction{transNat} \AgdaSymbol{\{}\AgdaInductiveConstructor{zero}\AgdaSymbol{\}} \AgdaSymbol{\{}\AgdaInductiveConstructor{zero}\AgdaSymbol{\}} \AgdaBound{xy} \AgdaBound{yz} \AgdaSymbol{=} \AgdaBound{yz}\<%
\\
\>[0]\AgdaIndent{2}{}\<[2]%
\>[2]\AgdaFunction{transNat} \AgdaSymbol{\{}\AgdaInductiveConstructor{zero}\AgdaSymbol{\}} \AgdaSymbol{\{}\AgdaInductiveConstructor{suc} \AgdaSymbol{\_\}} \AgdaSymbol{()} \AgdaBound{yz}\<%
\\
\>[0]\AgdaIndent{2}{}\<[2]%
\>[2]\AgdaFunction{transNat} \AgdaSymbol{\{}\AgdaInductiveConstructor{suc} \AgdaSymbol{\_\}} \AgdaSymbol{\{}\AgdaInductiveConstructor{zero}\AgdaSymbol{\}} \AgdaSymbol{()} \AgdaBound{yz}\<%
\\
\>[0]\AgdaIndent{2}{}\<[2]%
\>[2]\AgdaFunction{transNat} \AgdaSymbol{\{}\AgdaInductiveConstructor{suc} \AgdaSymbol{\_\}} \AgdaSymbol{\{}\AgdaInductiveConstructor{suc} \AgdaSymbol{\_\}} \AgdaSymbol{\{}\AgdaInductiveConstructor{zero}\AgdaSymbol{\}} \AgdaBound{xy} \AgdaBound{yz} \AgdaSymbol{=} \AgdaBound{yz}\<%
\\
\>[0]\AgdaIndent{2}{}\<[2]%
\>[2]\AgdaFunction{transNat} \AgdaSymbol{\{}\AgdaInductiveConstructor{suc} \AgdaBound{x}\AgdaSymbol{\}} \AgdaSymbol{\{}\AgdaInductiveConstructor{suc} \AgdaBound{y}\AgdaSymbol{\}} \AgdaSymbol{\{}\AgdaInductiveConstructor{suc} \AgdaBound{z}\AgdaSymbol{\}} \AgdaBound{xy} \AgdaBound{yz} \AgdaSymbol{=} \AgdaFunction{transNat} \AgdaSymbol{\{}\AgdaBound{x}\AgdaSymbol{\}} \AgdaSymbol{\{}\AgdaBound{y}\AgdaSymbol{\}} \AgdaSymbol{\{}\AgdaBound{z}\AgdaSymbol{\}} \AgdaBound{xy} \AgdaBound{yz}\<%
\\
%
\\
\>[0]\AgdaIndent{2}{}\<[2]%
\>[2]\AgdaFunction{⟦Nat⟧} \AgdaSymbol{:} \AgdaRecord{Ty} \AgdaBound{Γ}\<%
\\
\>[0]\AgdaIndent{2}{}\<[2]%
\>[2]\AgdaFunction{⟦Nat⟧} \AgdaSymbol{=} \AgdaKeyword{record} \<[17]%
\>[17]\<%
\\
\>[2]\AgdaIndent{4}{}\<[4]%
\>[4]\AgdaSymbol{\{} \AgdaField{fm} \AgdaSymbol{=} \AgdaSymbol{λ} \AgdaBound{γ} \AgdaSymbol{→} \AgdaKeyword{record}\<%
\\
\>[4]\AgdaIndent{9}{}\<[9]%
\>[9]\AgdaSymbol{\{} \AgdaField{Carrier} \AgdaSymbol{=} \AgdaDatatype{ℕ}\<%
\\
\>[4]\AgdaIndent{9}{}\<[9]%
\>[9]\AgdaSymbol{;} \AgdaField{\_≈h\_} \AgdaSymbol{=} \AgdaFunction{\_≈nat\_}\<%
\\
\>[4]\AgdaIndent{9}{}\<[9]%
\>[9]\AgdaSymbol{;} \AgdaField{refl} \AgdaSymbol{=} \AgdaSymbol{λ} \AgdaSymbol{\{}\AgdaBound{n}\AgdaSymbol{\}} \AgdaSymbol{→} \AgdaFunction{reflNat} \AgdaSymbol{\{}\AgdaBound{n}\AgdaSymbol{\}}\<%
\\
\>[4]\AgdaIndent{9}{}\<[9]%
\>[9]\AgdaSymbol{;} \AgdaField{sym} \AgdaSymbol{=} \AgdaSymbol{λ} \AgdaSymbol{\{}\AgdaBound{x}\AgdaSymbol{\}} \AgdaSymbol{\{}\AgdaBound{y}\AgdaSymbol{\}} \AgdaSymbol{→} \AgdaFunction{symNat} \AgdaSymbol{\{}\AgdaBound{x}\AgdaSymbol{\}} \AgdaSymbol{\{}\AgdaBound{y}\AgdaSymbol{\}}\<%
\\
\>[4]\AgdaIndent{9}{}\<[9]%
\>[9]\AgdaSymbol{;} \AgdaField{trans} \AgdaSymbol{=} \AgdaSymbol{λ} \AgdaSymbol{\{}\AgdaBound{x}\AgdaSymbol{\}} \AgdaSymbol{\{}\AgdaBound{y}\AgdaSymbol{\}} \AgdaSymbol{\{}\AgdaBound{z}\AgdaSymbol{\}} \AgdaSymbol{→} \AgdaFunction{transNat} \AgdaSymbol{\{}\AgdaBound{x}\AgdaSymbol{\}} \AgdaSymbol{\{}\AgdaBound{y}\AgdaSymbol{\}} \AgdaSymbol{\{}\AgdaBound{z}\AgdaSymbol{\}}\<%
\\
\>[4]\AgdaIndent{9}{}\<[9]%
\>[9]\AgdaSymbol{\}}\<%
\\
\>[0]\AgdaIndent{4}{}\<[4]%
\>[4]\AgdaSymbol{;} \AgdaField{substT} \AgdaSymbol{=} \AgdaSymbol{λ} \AgdaBound{\_} \AgdaSymbol{→} \AgdaFunction{id}\<%
\\
\>[0]\AgdaIndent{4}{}\<[4]%
\>[4]\AgdaSymbol{;} \AgdaField{subst*} \AgdaSymbol{=} \AgdaSymbol{λ} \AgdaBound{\_} \AgdaSymbol{→} \AgdaFunction{id}\<%
\\
\>[0]\AgdaIndent{4}{}\<[4]%
\>[4]\AgdaSymbol{;} \AgdaField{refl*} \AgdaSymbol{=} \AgdaSymbol{λ} \AgdaBound{x} \AgdaBound{a} \AgdaSymbol{→} \AgdaFunction{reflNat} \AgdaSymbol{\{}\AgdaBound{a}\AgdaSymbol{\}}\<%
\\
\>[0]\AgdaIndent{4}{}\<[4]%
\>[4]\AgdaSymbol{;} \AgdaField{trans*} \AgdaSymbol{=} \AgdaSymbol{λ} \AgdaBound{a} \AgdaSymbol{→} \AgdaFunction{reflNat} \AgdaSymbol{\{}\AgdaBound{a}\AgdaSymbol{\}} \<[33]%
\>[33]\<%
\\
\>[0]\AgdaIndent{4}{}\<[4]%
\>[4]\AgdaSymbol{\}}\<%
\\
%
\\
\>[0]\AgdaIndent{2}{}\<[2]%
\>[2]\AgdaFunction{⟦0⟧} \AgdaSymbol{:} \AgdaRecord{Tm} \AgdaFunction{⟦Nat⟧}\<%
\\
\>[0]\AgdaIndent{2}{}\<[2]%
\>[2]\AgdaFunction{⟦0⟧} \AgdaSymbol{=} \AgdaKeyword{record}\<%
\\
\>[2]\AgdaIndent{6}{}\<[6]%
\>[6]\AgdaSymbol{\{} \AgdaField{tm} \AgdaSymbol{=} \AgdaSymbol{λ} \AgdaBound{\_} \AgdaSymbol{→} \AgdaNumber{0}\<%
\\
\>[2]\AgdaIndent{6}{}\<[6]%
\>[6]\AgdaSymbol{;} \AgdaField{respt} \AgdaSymbol{=} \AgdaSymbol{λ} \AgdaBound{p} \AgdaSymbol{→} \AgdaInductiveConstructor{tt}\<%
\\
\>[2]\AgdaIndent{6}{}\<[6]%
\>[6]\AgdaSymbol{\}}\<%
\\
%
\\
\>[0]\AgdaIndent{2}{}\<[2]%
\>[2]\AgdaFunction{⟦s⟧} \AgdaSymbol{:} \AgdaRecord{Tm} \AgdaFunction{⟦Nat⟧} \AgdaSymbol{→} \AgdaRecord{Tm} \AgdaFunction{⟦Nat⟧}\<%
\\
\>[0]\AgdaIndent{2}{}\<[2]%
\>[2]\AgdaFunction{⟦s⟧} \AgdaSymbol{(}\AgdaInductiveConstructor{tm:} \AgdaBound{t} \AgdaInductiveConstructor{resp:} \AgdaBound{respt}\AgdaSymbol{)} \<[26]%
\>[26]\<%
\\
\>[2]\AgdaIndent{6}{}\<[6]%
\>[6]\AgdaSymbol{=} \AgdaKeyword{record}\<%
\\
\>[2]\AgdaIndent{6}{}\<[6]%
\>[6]\AgdaSymbol{\{} \AgdaField{tm} \AgdaSymbol{=} \AgdaInductiveConstructor{suc} \AgdaFunction{∘} \AgdaBound{t}\<%
\\
\>[2]\AgdaIndent{6}{}\<[6]%
\>[6]\AgdaSymbol{;} \AgdaField{respt} \AgdaSymbol{=} \AgdaBound{respt}\<%
\\
\>[2]\AgdaIndent{6}{}\<[6]%
\>[6]\AgdaSymbol{\}}\<%
\\
\>\<\end{code}

Simply typed universe

\AgdaHide{
\begin{code}\>\<%
\\
\>[0]\AgdaIndent{2}{}\<[2]%
\>[2]\AgdaKeyword{data} \AgdaDatatype{⟦U⟧⁰} \AgdaSymbol{:} \AgdaPrimitiveType{Set} \AgdaKeyword{where}\<%
\\
\>[0]\AgdaIndent{4}{}\<[4]%
\>[4]\AgdaInductiveConstructor{nat} \AgdaSymbol{:} \AgdaDatatype{⟦U⟧⁰}\<%
\\
\>[0]\AgdaIndent{4}{}\<[4]%
\>[4]\AgdaInductiveConstructor{arr<\_,\_>} \AgdaSymbol{:} \AgdaSymbol{(}\AgdaBound{a} \AgdaBound{b} \AgdaSymbol{:} \AgdaDatatype{⟦U⟧⁰}\AgdaSymbol{)} \AgdaSymbol{→} \AgdaDatatype{⟦U⟧⁰}\<%
\\
%
\\
\>[0]\AgdaIndent{2}{}\<[2]%
\>[2]\AgdaFunction{\_\textasciitilde⟦U⟧\_} \AgdaSymbol{:} \AgdaDatatype{⟦U⟧⁰} \AgdaSymbol{→} \AgdaDatatype{⟦U⟧⁰} \AgdaSymbol{→} \AgdaRecord{HProp}\<%
\\
\>[0]\AgdaIndent{2}{}\<[2]%
\>[2]\AgdaInductiveConstructor{nat} \AgdaFunction{\textasciitilde⟦U⟧} \AgdaInductiveConstructor{nat} \AgdaSymbol{=} \AgdaFunction{⊤'}\<%
\\
\>[0]\AgdaIndent{2}{}\<[2]%
\>[2]\AgdaInductiveConstructor{nat} \AgdaFunction{\textasciitilde⟦U⟧} \AgdaInductiveConstructor{arr<} \AgdaBound{a} \AgdaInductiveConstructor{,} \AgdaBound{b} \AgdaInductiveConstructor{>} \AgdaSymbol{=} \AgdaFunction{⊥'}\<%
\\
\>[0]\AgdaIndent{2}{}\<[2]%
\>[2]\AgdaInductiveConstructor{arr<} \AgdaBound{a} \AgdaInductiveConstructor{,} \AgdaBound{b} \AgdaInductiveConstructor{>} \AgdaFunction{\textasciitilde⟦U⟧} \AgdaInductiveConstructor{nat} \AgdaSymbol{=} \AgdaFunction{⊥'}\<%
\\
\>[0]\AgdaIndent{2}{}\<[2]%
\>[2]\AgdaInductiveConstructor{arr<} \AgdaBound{a} \AgdaInductiveConstructor{,} \AgdaBound{b} \AgdaInductiveConstructor{>} \AgdaFunction{\textasciitilde⟦U⟧} \AgdaInductiveConstructor{arr<} \AgdaBound{a'} \AgdaInductiveConstructor{,} \AgdaBound{b'} \AgdaInductiveConstructor{>} \AgdaSymbol{=} \AgdaBound{a} \AgdaFunction{\textasciitilde⟦U⟧} \AgdaBound{a'} \AgdaFunction{∧} \AgdaBound{b} \AgdaFunction{\textasciitilde⟦U⟧} \AgdaBound{b'}\<%
\\
%
\\
\>[0]\AgdaIndent{2}{}\<[2]%
\>[2]\AgdaFunction{reflU} \AgdaSymbol{:} \<[11]%
\>[11]\AgdaSymbol{\{}\AgdaBound{x} \AgdaSymbol{:} \AgdaDatatype{⟦U⟧⁰}\AgdaSymbol{\}} \AgdaSymbol{→} \AgdaFunction{<} \AgdaBound{x} \AgdaFunction{\textasciitilde⟦U⟧} \AgdaBound{x} \AgdaFunction{>}\<%
\\
\>[0]\AgdaIndent{2}{}\<[2]%
\>[2]\AgdaFunction{reflU} \AgdaSymbol{\{}\AgdaInductiveConstructor{nat}\AgdaSymbol{\}} \AgdaSymbol{=} \AgdaInductiveConstructor{tt}\<%
\\
\>[0]\AgdaIndent{2}{}\<[2]%
\>[2]\AgdaFunction{reflU} \AgdaSymbol{\{}\AgdaInductiveConstructor{arr<} \AgdaBound{a} \AgdaInductiveConstructor{,} \AgdaBound{b} \AgdaInductiveConstructor{>}\AgdaSymbol{\}} \AgdaSymbol{=} \AgdaFunction{reflU} \AgdaSymbol{\{}\AgdaBound{a}\AgdaSymbol{\}} \AgdaInductiveConstructor{,} \AgdaFunction{reflU} \AgdaSymbol{\{}\AgdaBound{b}\AgdaSymbol{\}}\<%
\\
%
\\
\>[0]\AgdaIndent{2}{}\<[2]%
\>[2]\AgdaFunction{symU} \AgdaSymbol{:} \AgdaSymbol{\{}\AgdaBound{x} \AgdaBound{y} \AgdaSymbol{:} \AgdaDatatype{⟦U⟧⁰}\AgdaSymbol{\}} \AgdaSymbol{→} \AgdaFunction{<} \AgdaBound{x} \AgdaFunction{\textasciitilde⟦U⟧} \AgdaBound{y} \AgdaFunction{>} \AgdaSymbol{→} \AgdaFunction{<} \AgdaBound{y} \AgdaFunction{\textasciitilde⟦U⟧} \AgdaBound{x} \AgdaFunction{>}\<%
\\
\>[0]\AgdaIndent{2}{}\<[2]%
\>[2]\AgdaFunction{symU} \AgdaSymbol{\{}\AgdaInductiveConstructor{nat}\AgdaSymbol{\}} \AgdaSymbol{\{}\AgdaInductiveConstructor{nat}\AgdaSymbol{\}} \AgdaBound{eq} \AgdaSymbol{=} \AgdaInductiveConstructor{tt}\<%
\\
\>[0]\AgdaIndent{2}{}\<[2]%
\>[2]\AgdaFunction{symU} \AgdaSymbol{\{}\AgdaInductiveConstructor{nat}\AgdaSymbol{\}} \AgdaSymbol{\{}\AgdaInductiveConstructor{arr<} \AgdaBound{a} \AgdaInductiveConstructor{,} \AgdaBound{b} \AgdaInductiveConstructor{>}\AgdaSymbol{\}} \AgdaBound{eq} \AgdaSymbol{=} \AgdaBound{eq}\<%
\\
\>[0]\AgdaIndent{2}{}\<[2]%
\>[2]\AgdaFunction{symU} \AgdaSymbol{\{}\AgdaInductiveConstructor{arr<} \AgdaBound{a} \AgdaInductiveConstructor{,} \AgdaBound{b} \AgdaInductiveConstructor{>}\AgdaSymbol{\}} \AgdaSymbol{\{}\AgdaInductiveConstructor{nat}\AgdaSymbol{\}} \AgdaBound{eq} \AgdaSymbol{=} \AgdaBound{eq}\<%
\\
\>[0]\AgdaIndent{2}{}\<[2]%
\>[2]\AgdaFunction{symU} \AgdaSymbol{\{}\AgdaInductiveConstructor{arr<} \AgdaBound{a} \AgdaInductiveConstructor{,} \AgdaBound{b} \AgdaInductiveConstructor{>}\AgdaSymbol{\}} \AgdaSymbol{\{}\AgdaInductiveConstructor{arr<} \AgdaBound{a'} \AgdaInductiveConstructor{,} \AgdaBound{b'} \AgdaInductiveConstructor{>}\AgdaSymbol{\}} \AgdaSymbol{(}\AgdaBound{p} \AgdaInductiveConstructor{,} \AgdaBound{q}\AgdaSymbol{)} \AgdaSymbol{=} \AgdaSymbol{(}\AgdaFunction{symU} \AgdaSymbol{\{}\AgdaBound{a}\AgdaSymbol{\}} \AgdaSymbol{\{}\AgdaBound{a'}\AgdaSymbol{\}} \AgdaBound{p}\AgdaSymbol{)} \<[67]%
\>[67]\<%
\\
\>[2]\AgdaIndent{47}{}\<[47]%
\>[47]\AgdaInductiveConstructor{,} \AgdaSymbol{(}\AgdaFunction{symU} \AgdaSymbol{\{}\AgdaBound{b}\AgdaSymbol{\}} \AgdaSymbol{\{}\AgdaBound{b'}\AgdaSymbol{\}} \AgdaBound{q}\AgdaSymbol{)}\<%
\\
%
\\
\>[0]\AgdaIndent{2}{}\<[2]%
\>[2]\AgdaFunction{transU} \AgdaSymbol{:} \AgdaSymbol{\{}\AgdaBound{x} \AgdaBound{y} \AgdaBound{z} \AgdaSymbol{:} \AgdaDatatype{⟦U⟧⁰}\AgdaSymbol{\}} \AgdaSymbol{→} \AgdaFunction{<} \AgdaBound{x} \AgdaFunction{\textasciitilde⟦U⟧} \AgdaBound{y} \AgdaFunction{>} \AgdaSymbol{→} \AgdaFunction{<} \AgdaBound{y} \AgdaFunction{\textasciitilde⟦U⟧} \AgdaBound{z} \AgdaFunction{>} \AgdaSymbol{→} \AgdaFunction{<} \AgdaBound{x} \AgdaFunction{\textasciitilde⟦U⟧} \AgdaBound{z} \AgdaFunction{>}\<%
\\
\>[0]\AgdaIndent{2}{}\<[2]%
\>[2]\AgdaFunction{transU} \AgdaSymbol{\{}\AgdaInductiveConstructor{nat}\AgdaSymbol{\}} \AgdaSymbol{\{}\AgdaInductiveConstructor{nat}\AgdaSymbol{\}} \AgdaBound{eq1} \AgdaBound{eq2} \AgdaSymbol{=} \AgdaBound{eq2}\<%
\\
\>[0]\AgdaIndent{2}{}\<[2]%
\>[2]\AgdaFunction{transU} \AgdaSymbol{\{}\AgdaInductiveConstructor{nat}\AgdaSymbol{\}} \AgdaSymbol{\{}\AgdaInductiveConstructor{arr<} \AgdaBound{a} \AgdaInductiveConstructor{,} \AgdaBound{b} \AgdaInductiveConstructor{>}\AgdaSymbol{\}} \AgdaSymbol{()} \AgdaBound{eq2}\<%
\\
\>[0]\AgdaIndent{2}{}\<[2]%
\>[2]\AgdaFunction{transU} \AgdaSymbol{\{}\AgdaInductiveConstructor{arr<} \AgdaBound{a} \AgdaInductiveConstructor{,} \AgdaBound{b} \AgdaInductiveConstructor{>}\AgdaSymbol{\}} \AgdaSymbol{\{}\AgdaInductiveConstructor{nat}\AgdaSymbol{\}} \AgdaSymbol{()} \AgdaBound{eq2}\<%
\\
\>[0]\AgdaIndent{2}{}\<[2]%
\>[2]\AgdaFunction{transU} \AgdaSymbol{\{}\AgdaInductiveConstructor{arr<} \AgdaBound{a} \AgdaInductiveConstructor{,} \AgdaBound{b} \AgdaInductiveConstructor{>}\AgdaSymbol{\}} \AgdaSymbol{\{}\AgdaInductiveConstructor{arr<} \AgdaBound{a'} \AgdaInductiveConstructor{,} \AgdaBound{b'} \AgdaInductiveConstructor{>}\AgdaSymbol{\}} \AgdaSymbol{\{}\AgdaInductiveConstructor{nat}\AgdaSymbol{\}} \AgdaBound{eq1} \AgdaBound{eq2} \AgdaSymbol{=} \AgdaBound{eq2}\<%
\\
\>[0]\AgdaIndent{2}{}\<[2]%
\>[2]\AgdaFunction{transU} \AgdaSymbol{\{}\AgdaInductiveConstructor{arr<} \AgdaBound{a} \AgdaInductiveConstructor{,} \AgdaBound{b} \AgdaInductiveConstructor{>}\AgdaSymbol{\}} \AgdaSymbol{\{}\AgdaInductiveConstructor{arr<} \AgdaBound{a'} \AgdaInductiveConstructor{,} \AgdaBound{b'} \AgdaInductiveConstructor{>}\AgdaSymbol{\}} \AgdaSymbol{\{}\AgdaInductiveConstructor{arr<} \AgdaBound{a0} \AgdaInductiveConstructor{,} \AgdaBound{b0} \AgdaInductiveConstructor{>}\AgdaSymbol{\}} \AgdaSymbol{(}\AgdaBound{p1} \AgdaInductiveConstructor{,} \AgdaBound{q1}\AgdaSymbol{)} \<[68]%
\>[68]\<%
\\
\>[2]\AgdaIndent{9}{}\<[9]%
\>[9]\AgdaSymbol{(}\AgdaBound{p2} \AgdaInductiveConstructor{,} \AgdaBound{q2}\AgdaSymbol{)} \AgdaSymbol{=} \AgdaSymbol{(}\AgdaFunction{transU} \AgdaSymbol{\{}\AgdaBound{a}\AgdaSymbol{\}} \AgdaSymbol{\{}\AgdaBound{a'}\AgdaSymbol{\}} \AgdaSymbol{\{}\AgdaBound{a0}\AgdaSymbol{\}} \AgdaBound{p1} \AgdaBound{p2}\AgdaSymbol{)} \<[50]%
\>[50]\<%
\\
\>[2]\AgdaIndent{9}{}\<[9]%
\>[9]\AgdaInductiveConstructor{,} \AgdaFunction{transU} \AgdaSymbol{\{}\AgdaBound{b}\AgdaSymbol{\}} \AgdaSymbol{\{}\AgdaBound{b'}\AgdaSymbol{\}} \AgdaSymbol{\{}\AgdaBound{b0}\AgdaSymbol{\}} \AgdaBound{q1} \AgdaBound{q2}\<%
\\
%
\\
\>[0]\AgdaIndent{2}{}\<[2]%
\>[2]\AgdaFunction{⟦U⟧} \AgdaSymbol{:} \AgdaRecord{Ty} \AgdaBound{Γ}\<%
\\
\>[0]\AgdaIndent{2}{}\<[2]%
\>[2]\AgdaFunction{⟦U⟧} \AgdaSymbol{=} \AgdaKeyword{record} \<[15]%
\>[15]\<%
\\
\>[2]\AgdaIndent{4}{}\<[4]%
\>[4]\AgdaSymbol{\{} \AgdaField{fm} \AgdaSymbol{=} \AgdaSymbol{λ} \AgdaBound{γ} \AgdaSymbol{→} \AgdaKeyword{record}\<%
\\
\>[4]\AgdaIndent{9}{}\<[9]%
\>[9]\AgdaSymbol{\{} \AgdaField{Carrier} \AgdaSymbol{=} \AgdaDatatype{⟦U⟧⁰}\<%
\\
\>[4]\AgdaIndent{9}{}\<[9]%
\>[9]\AgdaSymbol{;} \AgdaField{\_≈h\_} \AgdaSymbol{=} \AgdaFunction{\_\textasciitilde⟦U⟧\_}\<%
\\
\>[4]\AgdaIndent{9}{}\<[9]%
\>[9]\AgdaSymbol{;} \AgdaField{refl} \AgdaSymbol{=} \AgdaSymbol{λ} \AgdaSymbol{\{}\AgdaBound{x}\AgdaSymbol{\}} \AgdaSymbol{→} \AgdaFunction{reflU} \AgdaSymbol{\{}\AgdaBound{x}\AgdaSymbol{\}}\<%
\\
\>[4]\AgdaIndent{9}{}\<[9]%
\>[9]\AgdaSymbol{;} \AgdaField{sym} \AgdaSymbol{=} \AgdaSymbol{λ} \AgdaSymbol{\{}\AgdaBound{x}\AgdaSymbol{\}} \AgdaSymbol{\{}\AgdaBound{y}\AgdaSymbol{\}} \AgdaSymbol{→} \AgdaFunction{symU} \AgdaSymbol{\{}\AgdaBound{x}\AgdaSymbol{\}} \AgdaSymbol{\{}\AgdaBound{y}\AgdaSymbol{\}}\<%
\\
\>[4]\AgdaIndent{9}{}\<[9]%
\>[9]\AgdaSymbol{;} \AgdaField{trans} \AgdaSymbol{=} \AgdaSymbol{λ} \AgdaSymbol{\{}\AgdaBound{x}\AgdaSymbol{\}} \AgdaSymbol{\{}\AgdaBound{y}\AgdaSymbol{\}} \AgdaSymbol{\{}\AgdaBound{z}\AgdaSymbol{\}} \AgdaSymbol{→} \AgdaFunction{transU} \AgdaSymbol{\{}\AgdaBound{x}\AgdaSymbol{\}} \AgdaSymbol{\{}\AgdaBound{y}\AgdaSymbol{\}} \AgdaSymbol{\{}\AgdaBound{z}\AgdaSymbol{\}}\<%
\\
\>[4]\AgdaIndent{9}{}\<[9]%
\>[9]\AgdaSymbol{\}}\<%
\\
\>[0]\AgdaIndent{4}{}\<[4]%
\>[4]\AgdaSymbol{;} \AgdaField{substT} \AgdaSymbol{=} \AgdaSymbol{λ} \AgdaBound{\_} \AgdaSymbol{→} \AgdaFunction{id}\<%
\\
\>[0]\AgdaIndent{4}{}\<[4]%
\>[4]\AgdaSymbol{;} \AgdaField{subst*} \AgdaSymbol{=} \AgdaSymbol{λ} \AgdaBound{\_} \AgdaSymbol{→} \AgdaFunction{id}\<%
\\
\>[0]\AgdaIndent{4}{}\<[4]%
\>[4]\AgdaSymbol{;} \AgdaField{refl*} \AgdaSymbol{=} \AgdaSymbol{λ} \AgdaBound{x} \AgdaBound{a} \AgdaSymbol{→} \AgdaFunction{reflU} \AgdaSymbol{\{}\AgdaBound{a}\AgdaSymbol{\}}\<%
\\
\>[0]\AgdaIndent{4}{}\<[4]%
\>[4]\AgdaSymbol{;} \AgdaField{trans*} \AgdaSymbol{=} \AgdaSymbol{λ} \AgdaBound{a} \AgdaSymbol{→} \AgdaFunction{reflU} \AgdaSymbol{\{}\AgdaBound{a}\AgdaSymbol{\}}\<%
\\
\>[0]\AgdaIndent{4}{}\<[4]%
\>[4]\AgdaSymbol{\}}\<%
\\
%
\\
\>[0]\AgdaIndent{2}{}\<[2]%
\>[2]\AgdaFunction{elfm} \AgdaSymbol{:} \AgdaRecord{Σ} \AgdaFunction{∣} \AgdaBound{Γ} \AgdaFunction{∣} \AgdaSymbol{(λ} \AgdaBound{x} \AgdaSymbol{→} \AgdaDatatype{⟦U⟧⁰}\AgdaSymbol{)} \AgdaSymbol{→} \AgdaRecord{HSetoid}\<%
\\
\>[0]\AgdaIndent{2}{}\<[2]%
\>[2]\AgdaFunction{elfm} \AgdaSymbol{(}\AgdaBound{γ} \AgdaInductiveConstructor{,} \AgdaInductiveConstructor{nat}\AgdaSymbol{)} \AgdaSymbol{=} \AgdaFunction{[} \AgdaFunction{⟦Nat⟧} \AgdaFunction{]fm} \AgdaBound{γ}\<%
\\
\>[0]\AgdaIndent{2}{}\<[2]%
\>[2]\AgdaFunction{elfm} \AgdaSymbol{(}\AgdaBound{γ} \AgdaInductiveConstructor{,} \AgdaInductiveConstructor{arr<} \AgdaBound{a} \AgdaInductiveConstructor{,} \AgdaBound{b} \AgdaInductiveConstructor{>}\AgdaSymbol{)} \AgdaSymbol{=} \AgdaFunction{[} \AgdaBound{Γ} \AgdaFunction{,} \AgdaBound{γ} \AgdaFunction{]} \AgdaFunction{elfm} \AgdaSymbol{(}\AgdaBound{γ} \AgdaInductiveConstructor{,} \AgdaBound{a}\AgdaSymbol{)} \AgdaFunction{⇒fm} \AgdaFunction{elfm} \AgdaSymbol{(}\AgdaBound{γ} \AgdaInductiveConstructor{,} \AgdaBound{b}\AgdaSymbol{)}\<%
\\
\>\<\end{code}
}

\AgdaHide{
\begin{code}\>\<%
\\
%
\\
\>\AgdaComment{\{- To do : To find the way to extract the substT from ->

  elsubstT : \{x y : Σ ∣ Γ ∣ (λ x' → ⟦U⟧⁰)\} →
      Σ < [ Γ ] proj₁ x ≈h proj₁ y > (λ x' → < proj₂ x \textasciitilde⟦U⟧ proj₂ y >) →
      ∣ elfm x ∣ → ∣ elfm y ∣
  elsubstT \{\_ , nat\} \{\_ , nat\} \_ x' = x'
  elsubstT \{\_ , nat\} \{\_ , arr< a , b >\} (p , ()) x'
  elsubstT \{\_ , arr< a , b >\} \{\_ , nat\} (p , ()) x'
  elsubstT \{γ , arr< a , b >\} \{γ' , arr< a' , b' >\} (p , qa , qb) (s1 , s2) = 
   \{!!\}

  ⟦El⟧ : Ty (Γ \& ⟦U⟧)
  ⟦El⟧ = record 
       \{ fm = elfm
       ; substT = elsubstT
       ; subst* = \{!!\}
       ; refl* = \{!!\}
       ; trans* = \{!!\} 
       \}

-\}}\<%
\\
\>\<\end{code}
}

The equality type

\begin{code}\>\<%
\\
\>\AgdaKeyword{module} \AgdaModule{Equality-Type} \AgdaSymbol{(}\AgdaBound{Γ} \AgdaSymbol{:} \AgdaFunction{Con}\AgdaSymbol{)(}\AgdaBound{A} \AgdaSymbol{:} \AgdaRecord{Ty} \AgdaBound{Γ}\AgdaSymbol{)} \AgdaKeyword{where}\<%
\\
%
\\
\>[0]\AgdaIndent{2}{}\<[2]%
\>[2]\AgdaFunction{⟦Id⟧} \AgdaSymbol{:} \AgdaFunction{Rel} \AgdaBound{A}\<%
\\
\>[0]\AgdaIndent{2}{}\<[2]%
\>[2]\AgdaFunction{⟦Id⟧} \AgdaSymbol{=} \AgdaKeyword{record} \<[16]%
\>[16]\<%
\\
\>[2]\AgdaIndent{4}{}\<[4]%
\>[4]\AgdaSymbol{\{} \AgdaField{fm} \AgdaSymbol{=} \AgdaSymbol{λ} \AgdaSymbol{\{((}\AgdaBound{x} \AgdaInductiveConstructor{,} \AgdaBound{a}\AgdaSymbol{)} \AgdaInductiveConstructor{,} \AgdaBound{b}\AgdaSymbol{)} \AgdaSymbol{→} \<[30]%
\>[30]\<%
\\
\>[4]\AgdaIndent{13}{}\<[13]%
\>[13]\AgdaKeyword{record} \<[20]%
\>[20]\<%
\\
\>[4]\AgdaIndent{13}{}\<[13]%
\>[13]\AgdaSymbol{\{} \AgdaField{Carrier} \AgdaSymbol{=} \AgdaFunction{[} \AgdaFunction{[} \AgdaBound{A} \AgdaFunction{]fm} \AgdaBound{x} \AgdaFunction{]} \AgdaBound{a} \AgdaFunction{≈} \AgdaBound{b}\<%
\\
\>[4]\AgdaIndent{13}{}\<[13]%
\>[13]\AgdaSymbol{;} \AgdaField{\_≈h\_} \AgdaSymbol{=} \AgdaSymbol{λ} \AgdaBound{\_} \AgdaBound{\_} \AgdaSymbol{→} \AgdaKeyword{record} \AgdaSymbol{\{} \AgdaField{prf} \AgdaSymbol{=} \AgdaRecord{⊤} \AgdaSymbol{;} \AgdaField{Uni} \AgdaSymbol{=} \AgdaInductiveConstructor{PE.refl} \AgdaSymbol{\}}\<%
\\
\>[4]\AgdaIndent{13}{}\<[13]%
\>[13]\AgdaSymbol{;} \AgdaField{refl} \AgdaSymbol{=} \AgdaInductiveConstructor{tt} \<[25]%
\>[25]\<%
\\
\>[4]\AgdaIndent{13}{}\<[13]%
\>[13]\AgdaSymbol{;} \AgdaField{sym} \AgdaSymbol{=} \AgdaFunction{id}\<%
\\
\>[4]\AgdaIndent{13}{}\<[13]%
\>[13]\AgdaSymbol{;} \AgdaField{trans} \AgdaSymbol{=} \AgdaSymbol{λ} \AgdaBound{\_} \AgdaBound{\_} \AgdaSymbol{→} \AgdaInductiveConstructor{tt}\<%
\\
\>[4]\AgdaIndent{13}{}\<[13]%
\>[13]\AgdaSymbol{\}}\<%
\\
\>[4]\AgdaIndent{13}{}\<[13]%
\>[13]\AgdaSymbol{\}}\<%
\\
\>[0]\AgdaIndent{4}{}\<[4]%
\>[4]\AgdaSymbol{;} \AgdaField{substT} \AgdaSymbol{=} \AgdaSymbol{λ} \AgdaSymbol{\{((}\AgdaBound{x} \AgdaInductiveConstructor{,} \AgdaBound{a}\AgdaSymbol{)} \AgdaInductiveConstructor{,} \AgdaBound{b}\AgdaSymbol{)} \AgdaBound{x0} \AgdaSymbol{→} \<[37]%
\>[37]\<%
\\
\>[0]\AgdaIndent{15}{}\<[15]%
\>[15]\AgdaFunction{[} \AgdaFunction{[} \AgdaBound{A} \AgdaFunction{]fm} \AgdaSymbol{\_} \AgdaFunction{]trans} \<[34]%
\>[34]\<%
\\
\>[0]\AgdaIndent{15}{}\<[15]%
\>[15]\AgdaSymbol{(}\AgdaFunction{[} \AgdaFunction{[} \AgdaBound{A} \AgdaFunction{]fm} \AgdaSymbol{\_} \AgdaFunction{]sym} \AgdaBound{a}\AgdaSymbol{)} \<[36]%
\>[36]\<%
\\
\>[0]\AgdaIndent{15}{}\<[15]%
\>[15]\AgdaSymbol{(}\AgdaFunction{[} \AgdaFunction{[} \AgdaBound{A} \AgdaFunction{]fm} \AgdaSymbol{\_} \AgdaFunction{]trans} \<[35]%
\>[35]\<%
\\
\>[0]\AgdaIndent{15}{}\<[15]%
\>[15]\AgdaSymbol{(}\AgdaFunction{[} \AgdaBound{A} \AgdaFunction{]subst*} \AgdaSymbol{\_} \AgdaBound{x0}\AgdaSymbol{)} \AgdaBound{b}\AgdaSymbol{)} \<[37]%
\>[37]\<%
\\
\>[0]\AgdaIndent{15}{}\<[15]%
\>[15]\AgdaSymbol{\}}\<%
\\
\>[0]\AgdaIndent{4}{}\<[4]%
\>[4]\AgdaSymbol{;} \AgdaField{subst*} \AgdaSymbol{=} \AgdaSymbol{λ} \AgdaBound{\_} \AgdaBound{\_} \AgdaSymbol{→} \AgdaInductiveConstructor{tt}\<%
\\
\>[0]\AgdaIndent{4}{}\<[4]%
\>[4]\AgdaSymbol{;} \AgdaField{refl*} \AgdaSymbol{=} \AgdaSymbol{λ} \AgdaBound{\_} \AgdaBound{\_} \AgdaSymbol{→} \AgdaInductiveConstructor{tt}\<%
\\
\>[0]\AgdaIndent{4}{}\<[4]%
\>[4]\AgdaSymbol{;} \AgdaField{trans*} \AgdaSymbol{=} \AgdaSymbol{λ} \AgdaBound{\_} \AgdaSymbol{→} \AgdaInductiveConstructor{tt} \<[24]%
\>[24]\<%
\\
\>[0]\AgdaIndent{4}{}\<[4]%
\>[4]\AgdaSymbol{\}}\<%
\\
%
\\
\>[0]\AgdaIndent{2}{}\<[2]%
\>[2]\AgdaFunction{⟦refl⟧⁰} \AgdaSymbol{:} \AgdaRecord{Tm} \AgdaSymbol{\{}\AgdaBound{Γ} \AgdaFunction{\&} \AgdaBound{A}\AgdaSymbol{\}} \AgdaSymbol{(}\AgdaFunction{⟦Id⟧} \AgdaFunction{[} \AgdaKeyword{record} \AgdaSymbol{\{} \AgdaField{fn} \AgdaSymbol{=} \AgdaSymbol{λ} \AgdaBound{x'} \AgdaSymbol{→} \AgdaBound{x'} \AgdaInductiveConstructor{,} \AgdaFunction{proj₂} \AgdaBound{x'} \<[66]%
\>[66]\<%
\\
\>[0]\AgdaIndent{23}{}\<[23]%
\>[23]\AgdaSymbol{;} \AgdaField{resp} \AgdaSymbol{=} \AgdaSymbol{λ} \AgdaBound{x'} \AgdaSymbol{→} \AgdaBound{x'} \AgdaInductiveConstructor{,} \AgdaFunction{proj₂} \AgdaBound{x'} \AgdaSymbol{\}} \AgdaFunction{]T}\AgdaSymbol{)} \<[59]%
\>[59]\<%
\\
\>[0]\AgdaIndent{2}{}\<[2]%
\>[2]\AgdaFunction{⟦refl⟧⁰} \AgdaSymbol{=} \AgdaKeyword{record}\<%
\\
\>[0]\AgdaIndent{11}{}\<[11]%
\>[11]\AgdaSymbol{\{} \AgdaField{tm} \AgdaSymbol{=} \AgdaSymbol{λ} \AgdaSymbol{\{(}\AgdaBound{x} \AgdaInductiveConstructor{,} \AgdaBound{a}\AgdaSymbol{)} \AgdaSymbol{→} \AgdaFunction{[} \AgdaFunction{[} \AgdaBound{A} \AgdaFunction{]fm} \AgdaBound{x} \AgdaFunction{]refl} \AgdaSymbol{\{}\AgdaBound{a}\AgdaSymbol{\}} \AgdaSymbol{\}}\<%
\\
\>[0]\AgdaIndent{11}{}\<[11]%
\>[11]\AgdaSymbol{;} \AgdaField{respt} \AgdaSymbol{=} \AgdaSymbol{λ} \AgdaBound{p} \AgdaSymbol{→} \AgdaInductiveConstructor{tt}\<%
\\
\>[0]\AgdaIndent{11}{}\<[11]%
\>[11]\AgdaSymbol{\}}\<%
\\
%
\\
\>[0]\AgdaIndent{2}{}\<[2]%
\>[2]\AgdaFunction{⟦refl⟧} \AgdaSymbol{=} \<[12]%
\>[12]\AgdaFunction{lam} \AgdaSymbol{\{}\AgdaBound{Γ}\AgdaSymbol{\}} \AgdaSymbol{\{}\AgdaBound{A}\AgdaSymbol{\}} \AgdaFunction{⟦refl⟧⁰}\<%
\\
\>\<\end{code}

Subst using equality types

\begin{code}\>\<%
\\
%
\\
\>[0]\AgdaIndent{2}{}\<[2]%
\>[2]\AgdaKeyword{module} \AgdaModule{substIn} \AgdaSymbol{(}\AgdaBound{B} \AgdaSymbol{:} \AgdaRecord{Ty} \AgdaSymbol{(}\AgdaBound{Γ} \AgdaFunction{\&} \AgdaBound{A}\AgdaSymbol{))} \AgdaKeyword{where}\<%
\\
\>[0]\AgdaIndent{2}{}\<[2]%
\>[2]\<%
\\
\>[2]\AgdaIndent{4}{}\<[4]%
\>[4]\AgdaFunction{⟦subst⟧⁰} \AgdaSymbol{:} \AgdaRecord{Tm} \AgdaSymbol{\{}\AgdaBound{Γ} \AgdaFunction{\&} \AgdaBound{A} \AgdaFunction{\&} \AgdaSymbol{(}\AgdaBound{A} \AgdaFunction{[} \AgdaFunction{fst\&} \AgdaSymbol{\{}A \AgdaSymbol{=} \AgdaBound{A}\AgdaSymbol{\}} \AgdaFunction{]T}\AgdaSymbol{)} \<[49]%
\>[49]\<%
\\
\>[4]\AgdaIndent{15}{}\<[15]%
\>[15]\AgdaFunction{\&} \AgdaFunction{⟦Id⟧} \AgdaFunction{\&} \AgdaBound{B} \AgdaFunction{[} \AgdaFunction{fst\&} \AgdaSymbol{\{}A \AgdaSymbol{=} \AgdaBound{A} \AgdaFunction{[} \AgdaFunction{fst\&} \AgdaSymbol{\{}A \AgdaSymbol{=} \AgdaBound{A}\AgdaSymbol{\}} \AgdaFunction{]T}\AgdaSymbol{\}} \<[60]%
\>[60]\AgdaFunction{]T} \<[63]%
\>[63]\<%
\\
\>[4]\AgdaIndent{15}{}\<[15]%
\>[15]\AgdaFunction{[} \AgdaFunction{fst\&} \AgdaSymbol{\{}A \AgdaSymbol{=} \AgdaFunction{⟦Id⟧}\AgdaSymbol{\}} \AgdaFunction{]T}\AgdaSymbol{\}} \<[37]%
\>[37]\<%
\\
\>[-2]\AgdaIndent{13}{}\<[13]%
\>[13]\AgdaSymbol{(}\AgdaBound{B} \AgdaFunction{[} \AgdaKeyword{record} \AgdaSymbol{\{} \AgdaField{fn} \AgdaSymbol{=} \AgdaSymbol{λ} \AgdaBound{x} \AgdaSymbol{→} \AgdaSymbol{(}\AgdaFunction{proj₁} \AgdaSymbol{(}\AgdaFunction{proj₁} \AgdaSymbol{(}\AgdaFunction{proj₁} \AgdaSymbol{(}\AgdaFunction{proj₁} \AgdaBound{x}\AgdaSymbol{))))} \<[72]%
\>[72]\<%
\\
\>[0]\AgdaIndent{13}{}\<[13]%
\>[13]\AgdaInductiveConstructor{,} \AgdaSymbol{(}\AgdaFunction{proj₂} \AgdaSymbol{(}\AgdaFunction{proj₁} \AgdaSymbol{(}\AgdaFunction{proj₁} \AgdaBound{x}\AgdaSymbol{)))} \<[41]%
\>[41]\<%
\\
\>[0]\AgdaIndent{13}{}\<[13]%
\>[13]\AgdaSymbol{;} \AgdaField{resp} \AgdaSymbol{=} \AgdaSymbol{λ} \AgdaBound{x} \AgdaSymbol{→} \AgdaFunction{proj₁} \AgdaSymbol{(}\AgdaFunction{proj₁} \AgdaSymbol{(}\AgdaFunction{proj₁} \AgdaSymbol{(}\AgdaFunction{proj₁} \AgdaBound{x}\AgdaSymbol{)))} \<[60]%
\>[60]\<%
\\
\>[0]\AgdaIndent{13}{}\<[13]%
\>[13]\AgdaInductiveConstructor{,} \AgdaFunction{proj₂} \AgdaSymbol{(}\AgdaFunction{proj₁} \AgdaSymbol{(}\AgdaFunction{proj₁} \AgdaBound{x}\AgdaSymbol{))} \AgdaSymbol{\}} \AgdaFunction{]T}\AgdaSymbol{)}\<%
\\
%
\\
\>[0]\AgdaIndent{4}{}\<[4]%
\>[4]\AgdaFunction{⟦subst⟧⁰} \AgdaSymbol{=} \AgdaKeyword{record}\<%
\\
\>[0]\AgdaIndent{11}{}\<[11]%
\>[11]\AgdaSymbol{\{} \AgdaField{tm} \AgdaSymbol{=} \AgdaSymbol{λ} \AgdaSymbol{\{((((}\AgdaBound{x} \AgdaInductiveConstructor{,} \AgdaBound{a}\AgdaSymbol{)} \AgdaInductiveConstructor{,} \AgdaBound{b}\AgdaSymbol{)} \AgdaInductiveConstructor{,} \AgdaBound{p}\AgdaSymbol{)} \AgdaInductiveConstructor{,} \AgdaBound{PA}\AgdaSymbol{)} \AgdaSymbol{→} \AgdaFunction{[} \AgdaBound{B} \AgdaFunction{]subst} \<[61]%
\>[61]\<%
\\
\>[11]\AgdaIndent{18}{}\<[18]%
\>[18]\AgdaSymbol{(}\AgdaFunction{[} \AgdaBound{Γ} \AgdaFunction{]refl} \AgdaInductiveConstructor{,} \AgdaFunction{[} \AgdaFunction{[} \AgdaBound{A} \AgdaFunction{]fm} \AgdaSymbol{\_} \AgdaFunction{]trans} \<[50]%
\>[50]\<%
\\
\>[11]\AgdaIndent{18}{}\<[18]%
\>[18]\AgdaSymbol{(}\AgdaFunction{[} \AgdaBound{A} \AgdaFunction{]refl*} \AgdaSymbol{\_} \AgdaSymbol{\_)} \AgdaBound{p}\AgdaSymbol{)} \AgdaBound{PA} \AgdaSymbol{\}}\<%
\\
\>[0]\AgdaIndent{11}{}\<[11]%
\>[11]\AgdaSymbol{;} \AgdaField{respt} \AgdaSymbol{=} \AgdaSymbol{λ} \AgdaSymbol{\{((((}\AgdaBound{m} \AgdaInductiveConstructor{,} \AgdaBound{a}\AgdaSymbol{)} \AgdaInductiveConstructor{,} \AgdaBound{b}\AgdaSymbol{)} \AgdaInductiveConstructor{,} \AgdaBound{p}\AgdaSymbol{)} \AgdaInductiveConstructor{,} \AgdaBound{PA}\AgdaSymbol{)} \AgdaSymbol{→} \<[53]%
\>[53]\<%
\\
\>[0]\AgdaIndent{13}{}\<[13]%
\>[13]\AgdaFunction{[} \AgdaFunction{[} \AgdaBound{B} \AgdaFunction{]fm} \AgdaSymbol{\_} \AgdaFunction{]trans} \<[32]%
\>[32]\<%
\\
\>[0]\AgdaIndent{13}{}\<[13]%
\>[13]\AgdaSymbol{(}\AgdaFunction{[} \AgdaBound{B} \AgdaFunction{]trans*} \AgdaSymbol{\_)} \<[29]%
\>[29]\<%
\\
\>[13]\AgdaIndent{14}{}\<[14]%
\>[14]\AgdaSymbol{(}\AgdaFunction{[} \AgdaFunction{[} \AgdaBound{B} \AgdaFunction{]fm} \AgdaSymbol{\_} \AgdaFunction{]trans} \<[34]%
\>[34]\<%
\\
\>[0]\AgdaIndent{13}{}\<[13]%
\>[13]\AgdaFunction{[} \AgdaBound{B} \AgdaFunction{]subst-pi} \<[27]%
\>[27]\<%
\\
\>[0]\AgdaIndent{13}{}\<[13]%
\>[13]\AgdaSymbol{(}\AgdaFunction{[} \AgdaFunction{[} \AgdaBound{B} \AgdaFunction{]fm} \AgdaSymbol{\_} \AgdaFunction{]trans} \<[33]%
\>[33]\<%
\\
\>[0]\AgdaIndent{13}{}\<[13]%
\>[13]\AgdaSymbol{(}\AgdaFunction{[} \AgdaFunction{[} \AgdaBound{B} \AgdaFunction{]fm} \AgdaSymbol{\_} \AgdaFunction{]sym} \AgdaSymbol{(}\AgdaFunction{[} \AgdaBound{B} \AgdaFunction{]trans*} \AgdaSymbol{\_))}\<%
\\
\>[0]\AgdaIndent{13}{}\<[13]%
\>[13]\AgdaSymbol{(}\AgdaFunction{[} \AgdaBound{B} \AgdaFunction{]subst*} \AgdaSymbol{\_} \AgdaBound{PA}\AgdaSymbol{)} \AgdaSymbol{))} \AgdaSymbol{\}}\<%
\\
\>[0]\AgdaIndent{11}{}\<[11]%
\>[11]\AgdaSymbol{\}}\<%
\\
%
\\
\>[0]\AgdaIndent{4}{}\<[4]%
\>[4]\AgdaFunction{⟦subst⟧} \AgdaSymbol{=} \AgdaFunction{lam} \AgdaSymbol{\{}A \AgdaSymbol{=} \AgdaSymbol{(}\AgdaBound{A} \AgdaFunction{[} \AgdaFunction{fst\&} \AgdaSymbol{\{}A \AgdaSymbol{=} \AgdaBound{A}\AgdaSymbol{\}} \AgdaFunction{]T}\AgdaSymbol{)\}} \<[46]%
\>[46]\<%
\\
\>[4]\AgdaIndent{14}{}\<[14]%
\>[14]\AgdaSymbol{(}\AgdaFunction{lam} \AgdaSymbol{\{}A \AgdaSymbol{=} \AgdaFunction{⟦Id⟧}\AgdaSymbol{\}} \<[30]%
\>[30]\<%
\\
\>[4]\AgdaIndent{14}{}\<[14]%
\>[14]\AgdaSymbol{(}\AgdaFunction{lam} \AgdaSymbol{\{}A \AgdaSymbol{=} \AgdaBound{B} \AgdaFunction{[} \AgdaFunction{fst\&} \AgdaSymbol{\{}A \AgdaSymbol{=} \AgdaBound{A} \AgdaFunction{[} \AgdaFunction{fst\&} \AgdaSymbol{\{}A \AgdaSymbol{=} \AgdaBound{A}\AgdaSymbol{\}} \AgdaFunction{]T}\AgdaSymbol{\}}\AgdaFunction{]T} \<[61]%
\>[61]\<%
\\
\>[14]\AgdaIndent{15}{}\<[15]%
\>[15]\AgdaFunction{[} \AgdaFunction{fst\&} \AgdaSymbol{\{}A \AgdaSymbol{=} \AgdaFunction{⟦Id⟧}\AgdaSymbol{\}} \AgdaFunction{]T}\AgdaSymbol{\}} \AgdaFunction{⟦subst⟧⁰}\AgdaSymbol{))}\<%
\\
%
\\
\>\<\end{code}


\AgdaHide{
\begin{code}\>\<%
\\
\>\AgdaKeyword{import} \AgdaModule{Level}\<%
\\
\>\AgdaKeyword{open} \AgdaKeyword{import} \AgdaModule{Relation.Binary.PropositionalEquality} \AgdaSymbol{as} \AgdaModule{PE} \AgdaKeyword{hiding} \AgdaSymbol{(}refl \AgdaSymbol{;} sym \AgdaSymbol{;} trans\AgdaSymbol{;} isEquivalence\AgdaSymbol{;} [\_]\AgdaSymbol{)}\<%
\\
%
\\
\>\AgdaKeyword{module} \AgdaModule{CwF-quotient} \AgdaSymbol{(}\AgdaBound{ext} \AgdaSymbol{:} \AgdaFunction{Extensionality} \AgdaPrimitive{Level.zero} \AgdaPrimitive{Level.zero}\AgdaSymbol{)} \AgdaKeyword{where}\<%
\\
%
\\
\>\AgdaKeyword{open} \AgdaKeyword{import} \AgdaModule{Data.Unit}\<%
\\
\>\AgdaKeyword{open} \AgdaKeyword{import} \AgdaModule{Function}\<%
\\
\>\AgdaKeyword{open} \AgdaKeyword{import} \AgdaModule{Data.Product}\<%
\\
\>\AgdaKeyword{open} \AgdaKeyword{import} \AgdaModule{Data.Bool}\<%
\\
\>\AgdaKeyword{open} \AgdaKeyword{import} \AgdaModule{CategoryofSetoid}\<%
\\
\>\AgdaKeyword{open} \AgdaKeyword{import} \AgdaModule{CwF-setoid}\<%
\\
%
\\
\>\<\end{code}
}

Quotient types

\begin{code}\>\<%
\\
\>\AgdaKeyword{module} \AgdaModule{Q} \AgdaSymbol{(}\AgdaBound{Γ} \AgdaSymbol{:} \AgdaFunction{Con}\AgdaSymbol{)(}\AgdaBound{A} \AgdaSymbol{:} \AgdaRecord{Ty} \AgdaBound{Γ}\AgdaSymbol{)}\<%
\\
%
\\
\>[0]\AgdaIndent{9}{}\<[9]%
\>[9]\AgdaSymbol{(}\AgdaBound{R} \AgdaSymbol{:} \AgdaSymbol{(}\AgdaBound{γ} \AgdaSymbol{:} \AgdaFunction{∣} \AgdaBound{Γ} \AgdaFunction{∣}\AgdaSymbol{)} \AgdaSymbol{→} \AgdaFunction{∣} \AgdaFunction{[} \AgdaBound{A} \AgdaFunction{]fm} \AgdaBound{γ} \AgdaFunction{∣} \AgdaSymbol{→} \AgdaFunction{∣} \AgdaFunction{[} \AgdaBound{A} \AgdaFunction{]fm} \AgdaBound{γ} \AgdaFunction{∣} \AgdaSymbol{→} \AgdaPrimitiveType{Set}\AgdaSymbol{)}\<%
\\
%
\\
\>[0]\AgdaIndent{9}{}\<[9]%
\>[9]\AgdaSymbol{.(}\AgdaBound{Rrespt} \AgdaSymbol{:} \AgdaSymbol{∀\{}\AgdaBound{γ} \AgdaBound{γ'} \AgdaSymbol{:} \AgdaFunction{∣} \AgdaBound{Γ} \AgdaFunction{∣}\AgdaSymbol{\}}\<%
\\
\>[9]\AgdaIndent{19}{}\<[19]%
\>[19]\AgdaSymbol{(}\AgdaBound{p} \AgdaSymbol{:} \AgdaFunction{[} \AgdaBound{Γ} \AgdaFunction{]} \AgdaBound{γ} \AgdaFunction{≈} \AgdaBound{γ'}\AgdaSymbol{)}\<%
\\
\>[9]\AgdaIndent{19}{}\<[19]%
\>[19]\AgdaSymbol{(}\AgdaBound{a} \AgdaBound{b} \AgdaSymbol{:} \AgdaFunction{∣} \AgdaFunction{[} \AgdaBound{A} \AgdaFunction{]fm} \AgdaBound{γ} \AgdaFunction{∣}\AgdaSymbol{)} \AgdaSymbol{→}\<%
\\
\>[9]\AgdaIndent{19}{}\<[19]%
\>[19]\AgdaSymbol{.(}\AgdaBound{R} \AgdaBound{γ} \AgdaBound{a} \AgdaBound{b}\AgdaSymbol{)} \AgdaSymbol{→} \<[32]%
\>[32]\<%
\\
\>[9]\AgdaIndent{19}{}\<[19]%
\>[19]\AgdaBound{R} \AgdaBound{γ'} \AgdaSymbol{(}\AgdaFunction{[} \AgdaBound{A} \AgdaFunction{]subst} \AgdaBound{p} \AgdaBound{a}\AgdaSymbol{)} \AgdaSymbol{(}\AgdaFunction{[} \AgdaBound{A} \AgdaFunction{]subst} \AgdaBound{p} \AgdaBound{b}\AgdaSymbol{))}\<%
\\
%
\\
\>[0]\AgdaIndent{9}{}\<[9]%
\>[9]\AgdaSymbol{.(}\AgdaBound{Rrsp} \AgdaSymbol{:} \AgdaSymbol{∀} \AgdaSymbol{\{}\AgdaBound{γ} \AgdaBound{a} \AgdaBound{b}\AgdaSymbol{\}} \AgdaSymbol{→} \AgdaSymbol{.(}\AgdaFunction{[} \AgdaFunction{[} \AgdaBound{A} \AgdaFunction{]fm} \AgdaBound{γ} \AgdaFunction{]} \AgdaBound{a} \AgdaFunction{≈} \AgdaBound{b}\AgdaSymbol{)} \AgdaSymbol{→} \AgdaBound{R} \AgdaBound{γ} \AgdaBound{a} \AgdaBound{b}\AgdaSymbol{)}\<%
\\
%
\\
\>[0]\AgdaIndent{9}{}\<[9]%
\>[9]\AgdaSymbol{.(}\AgdaBound{Rref} \AgdaSymbol{:} \AgdaSymbol{∀} \AgdaSymbol{\{}\AgdaBound{γ} \AgdaBound{a}\AgdaSymbol{\}} \AgdaSymbol{→} \AgdaBound{R} \AgdaBound{γ} \AgdaBound{a} \AgdaBound{a}\AgdaSymbol{)}\<%
\\
\>[0]\AgdaIndent{9}{}\<[9]%
\>[9]\AgdaSymbol{.(}\AgdaBound{Rsym} \AgdaSymbol{:} \AgdaSymbol{(∀} \AgdaSymbol{\{}\AgdaBound{γ} \AgdaBound{a} \AgdaBound{b}\AgdaSymbol{\}} \AgdaSymbol{→} \AgdaSymbol{.(}\AgdaBound{R} \AgdaBound{γ} \AgdaBound{a} \AgdaBound{b}\AgdaSymbol{)} \AgdaSymbol{→} \AgdaBound{R} \AgdaBound{γ} \AgdaBound{b} \AgdaBound{a}\AgdaSymbol{))}\<%
\\
\>[0]\AgdaIndent{9}{}\<[9]%
\>[9]\AgdaSymbol{.(}\AgdaBound{Rtrn} \AgdaSymbol{:} \<[19]%
\>[19]\AgdaSymbol{(∀} \AgdaSymbol{\{}\AgdaBound{γ} \AgdaBound{a} \AgdaBound{b} \AgdaBound{c}\AgdaSymbol{\}} \AgdaSymbol{→} \AgdaSymbol{.(}\AgdaBound{R} \AgdaBound{γ} \AgdaBound{a} \AgdaBound{b}\AgdaSymbol{)}\<%
\\
\>[9]\AgdaIndent{17}{}\<[17]%
\>[17]\AgdaSymbol{→} \<[20]%
\>[20]\AgdaSymbol{.(}\AgdaBound{R} \AgdaBound{γ} \AgdaBound{b} \AgdaBound{c}\AgdaSymbol{)} \AgdaSymbol{→} \AgdaBound{R} \AgdaBound{γ} \AgdaBound{a} \AgdaBound{c}\AgdaSymbol{))}\<%
\\
\>[0]\AgdaIndent{9}{}\<[9]%
\>[9]\AgdaKeyword{where}\<%
\\
%
\\
\>[0]\AgdaIndent{2}{}\<[2]%
\>[2]\AgdaFunction{⟦Q⟧₀} \AgdaSymbol{:} \AgdaFunction{∣} \AgdaBound{Γ} \AgdaFunction{∣} \AgdaSymbol{→} \AgdaRecord{Setoid}\<%
\\
\>[0]\AgdaIndent{2}{}\<[2]%
\>[2]\AgdaFunction{⟦Q⟧₀} \AgdaBound{γ} \AgdaSymbol{=} \AgdaKeyword{record}\<%
\\
\>[2]\AgdaIndent{9}{}\<[9]%
\>[9]\AgdaSymbol{\{} \AgdaField{Carrier} \AgdaSymbol{=} \AgdaFunction{∣} \AgdaFunction{[} \AgdaBound{A} \AgdaFunction{]fm} \AgdaBound{γ} \AgdaFunction{∣}\<%
\\
\>[2]\AgdaIndent{9}{}\<[9]%
\>[9]\AgdaSymbol{;} \AgdaField{\_≈\_} \AgdaSymbol{=} \AgdaBound{R} \AgdaBound{γ}\<%
\\
\>[2]\AgdaIndent{9}{}\<[9]%
\>[9]\AgdaSymbol{;} \AgdaField{refl} \AgdaSymbol{=} \AgdaBound{Rref}\<%
\\
\>[2]\AgdaIndent{9}{}\<[9]%
\>[9]\AgdaSymbol{;} \AgdaField{sym} \AgdaSymbol{=} \AgdaBound{Rsym}\<%
\\
\>[2]\AgdaIndent{9}{}\<[9]%
\>[9]\AgdaSymbol{;} \AgdaField{trans} \AgdaSymbol{=} \AgdaBound{Rtrn}\<%
\\
\>[2]\AgdaIndent{9}{}\<[9]%
\>[9]\AgdaSymbol{\}}\<%
\\
%
\\
%
\\
\>[0]\AgdaIndent{2}{}\<[2]%
\>[2]\AgdaFunction{⟦Q⟧} \AgdaSymbol{:} \AgdaRecord{Ty} \AgdaBound{Γ}\<%
\\
\>[0]\AgdaIndent{2}{}\<[2]%
\>[2]\AgdaFunction{⟦Q⟧} \AgdaSymbol{=} \AgdaKeyword{record} \<[15]%
\>[15]\<%
\\
\>[2]\AgdaIndent{4}{}\<[4]%
\>[4]\AgdaSymbol{\{} \AgdaField{fm} \AgdaSymbol{=} \AgdaFunction{⟦Q⟧₀}\<%
\\
\>[2]\AgdaIndent{4}{}\<[4]%
\>[4]\AgdaSymbol{;} \AgdaField{substT} \AgdaSymbol{=} \AgdaFunction{[} \AgdaBound{A} \AgdaFunction{]subst}\<%
\\
\>[2]\AgdaIndent{4}{}\<[4]%
\>[4]\AgdaSymbol{;} \AgdaField{subst*} \AgdaSymbol{=} \AgdaSymbol{λ} \AgdaBound{p} \AgdaBound{q} \AgdaSymbol{→} \AgdaBound{Rrespt} \AgdaBound{p} \AgdaSymbol{\_} \AgdaSymbol{\_} \AgdaBound{q}\<%
\\
\>[2]\AgdaIndent{4}{}\<[4]%
\>[4]\AgdaSymbol{;} \AgdaField{refl*} \AgdaSymbol{=} \AgdaBound{Rrsp} \AgdaFunction{[} \AgdaBound{A} \AgdaFunction{]refl*}\<%
\\
\>[2]\AgdaIndent{4}{}\<[4]%
\>[4]\AgdaSymbol{;} \AgdaField{trans*} \AgdaSymbol{=} \AgdaSymbol{λ} \AgdaBound{a} \AgdaSymbol{→} \AgdaBound{Rrsp} \AgdaSymbol{(}\AgdaFunction{[} \AgdaBound{A} \AgdaFunction{]trans*} \AgdaSymbol{\_)}\<%
\\
\>[2]\AgdaIndent{4}{}\<[4]%
\>[4]\AgdaSymbol{\}}\<%
\\
%
\\
\>[0]\AgdaIndent{2}{}\<[2]%
\>[2]\AgdaFunction{⟦[\_]⟧} \AgdaSymbol{:} \AgdaRecord{Tm} \AgdaBound{A} \AgdaSymbol{→} \AgdaRecord{Tm} \AgdaFunction{⟦Q⟧}\<%
\\
\>[0]\AgdaIndent{2}{}\<[2]%
\>[2]\AgdaFunction{⟦[} \AgdaBound{x} \AgdaFunction{]⟧} \AgdaSymbol{=} \AgdaKeyword{record}\<%
\\
\>[2]\AgdaIndent{11}{}\<[11]%
\>[11]\AgdaSymbol{\{} \AgdaField{tm} \AgdaSymbol{=} \AgdaFunction{[} \AgdaBound{x} \AgdaFunction{]tm}\<%
\\
\>[2]\AgdaIndent{11}{}\<[11]%
\>[11]\AgdaSymbol{;} \AgdaField{respt} \AgdaSymbol{=} \AgdaSymbol{λ} \AgdaBound{p} \AgdaSymbol{→} \AgdaBound{Rrsp} \AgdaSymbol{(}\AgdaFunction{[} \AgdaBound{x} \AgdaFunction{]respt} \AgdaBound{p}\AgdaSymbol{)}\<%
\\
\>[2]\AgdaIndent{11}{}\<[11]%
\>[11]\AgdaSymbol{\}}\<%
\\
\>[0]\AgdaIndent{2}{}\<[2]%
\>[2]\<%
\\
\>[0]\AgdaIndent{2}{}\<[2]%
\>[2]\AgdaFunction{⟦[\_]⟧'} \AgdaSymbol{:} \AgdaRecord{Tm} \AgdaSymbol{(}\AgdaBound{A} \AgdaFunction{⇒} \AgdaFunction{⟦Q⟧}\AgdaSymbol{)}\<%
\\
\>[0]\AgdaIndent{2}{}\<[2]%
\>[2]\AgdaFunction{⟦[\_]⟧'} \AgdaSymbol{=} \AgdaKeyword{record}\<%
\\
\>[2]\AgdaIndent{11}{}\<[11]%
\>[11]\AgdaSymbol{\{} \AgdaField{tm} \AgdaSymbol{=} \AgdaSymbol{λ} \AgdaBound{x} \AgdaSymbol{→} \AgdaSymbol{(λ} \AgdaBound{a} \AgdaSymbol{→} \AgdaBound{a}\AgdaSymbol{)} \AgdaInductiveConstructor{,} \<[36]%
\>[36]\<%
\\
\>[11]\AgdaIndent{18}{}\<[18]%
\>[18]\AgdaSymbol{(λ} \AgdaBound{a} \AgdaBound{b} \AgdaBound{p} \AgdaSymbol{→} \<[29]%
\>[29]\<%
\\
\>[11]\AgdaIndent{18}{}\<[18]%
\>[18]\AgdaBound{Rrsp} \AgdaSymbol{(}\AgdaFunction{[} \AgdaFunction{[} \AgdaBound{A} \AgdaFunction{]fm} \AgdaSymbol{\_} \AgdaFunction{]trans} \AgdaFunction{[} \AgdaBound{A} \AgdaFunction{]refl*} \AgdaBound{p}\AgdaSymbol{))}\<%
\\
\>[0]\AgdaIndent{11}{}\<[11]%
\>[11]\AgdaSymbol{;} \AgdaField{respt} \AgdaSymbol{=} \AgdaSymbol{λ} \AgdaBound{p} \AgdaBound{a} \AgdaSymbol{→} \AgdaBound{Rrsp} \AgdaFunction{[} \AgdaBound{A} \AgdaFunction{]tr*}\<%
\\
\>[0]\AgdaIndent{11}{}\<[11]%
\>[11]\AgdaSymbol{\}}\<%
\\
\>\<\end{code}

\begin{code}\>\>[0]\AgdaIndent{2}{}\<[2]%
\>[2]\<%
\\
\>[0]\AgdaIndent{2}{}\<[2]%
\>[2]\AgdaSymbol{.}\AgdaFunction{Q-Ax} \AgdaSymbol{:} \AgdaSymbol{∀} \AgdaBound{γ} \AgdaBound{a} \AgdaBound{b} \AgdaSymbol{→} \AgdaFunction{[} \AgdaFunction{[} \AgdaBound{A} \AgdaFunction{]fm} \AgdaBound{γ} \AgdaFunction{]} \AgdaBound{a} \AgdaFunction{≈} \AgdaBound{b} \AgdaSymbol{→} \AgdaFunction{[} \AgdaFunction{[} \AgdaFunction{⟦Q⟧} \AgdaFunction{]fm} \AgdaSymbol{\_} \AgdaFunction{]} \AgdaBound{a} \AgdaFunction{≈} \AgdaBound{b}\<%
\\
\>[0]\AgdaIndent{2}{}\<[2]%
\>[2]\AgdaFunction{Q-Ax} \AgdaBound{γ} \AgdaBound{a} \AgdaBound{b} \AgdaSymbol{=} \AgdaBound{Rrsp}\<%
\\
%
\\
\>[0]\AgdaIndent{2}{}\<[2]%
\>[2]\AgdaFunction{Q-elim} \AgdaSymbol{:} \AgdaSymbol{(}\AgdaBound{B} \AgdaSymbol{:} \AgdaRecord{Ty} \AgdaBound{Γ}\AgdaSymbol{)} \<[22]%
\>[22]\<%
\\
\>[2]\AgdaIndent{9}{}\<[9]%
\>[9]\AgdaSymbol{→} \AgdaSymbol{(}\AgdaBound{f} \AgdaSymbol{:} \AgdaRecord{Tm} \AgdaSymbol{(}\AgdaBound{A} \AgdaFunction{⇒} \AgdaBound{B}\AgdaSymbol{))}\<%
\\
\>[2]\AgdaIndent{9}{}\<[9]%
\>[9]\AgdaSymbol{→} \AgdaSymbol{(}\AgdaBound{frespR} \AgdaSymbol{:} \AgdaSymbol{∀} \AgdaBound{γ} \AgdaBound{a} \AgdaBound{b}\<%
\\
\>[9]\AgdaIndent{18}{}\<[18]%
\>[18]\AgdaSymbol{→} \AgdaSymbol{(}\AgdaBound{R} \AgdaBound{γ} \AgdaBound{a} \AgdaBound{b}\AgdaSymbol{)}\<%
\\
\>[9]\AgdaIndent{18}{}\<[18]%
\>[18]\AgdaSymbol{→} \AgdaFunction{[} \AgdaFunction{[} \AgdaBound{B} \AgdaFunction{]fm} \AgdaBound{γ} \AgdaFunction{]} \AgdaFunction{prj₁} \AgdaSymbol{(}\AgdaFunction{[} \AgdaBound{f} \AgdaFunction{]tm} \AgdaBound{γ}\AgdaSymbol{)} \AgdaBound{a} \<[53]%
\>[53]\<%
\\
\>[18]\AgdaIndent{20}{}\<[20]%
\>[20]\AgdaFunction{≈} \<[23]%
\>[23]\AgdaFunction{prj₁} \AgdaSymbol{(}\AgdaFunction{[} \AgdaBound{f} \AgdaFunction{]tm} \AgdaBound{γ}\AgdaSymbol{)} \AgdaBound{b}\AgdaSymbol{)}\<%
\\
\>[-7]\AgdaIndent{9}{}\<[9]%
\>[9]\AgdaSymbol{→} \AgdaRecord{Tm} \AgdaSymbol{(}\AgdaFunction{⟦Q⟧} \AgdaFunction{⇒} \AgdaBound{B}\AgdaSymbol{)}\<%
\\
\>[0]\AgdaIndent{2}{}\<[2]%
\>[2]\AgdaFunction{Q-elim} \AgdaBound{B} \AgdaBound{f} \AgdaBound{frespR} \AgdaSymbol{=} \AgdaKeyword{record}\<%
\\
\>[2]\AgdaIndent{11}{}\<[11]%
\>[11]\AgdaSymbol{\{} \AgdaField{tm} \AgdaSymbol{=} \AgdaSymbol{λ} \AgdaBound{γ} \AgdaSymbol{→} \AgdaFunction{prj₁} \AgdaSymbol{(}\AgdaFunction{[} \AgdaBound{f} \AgdaFunction{]tm} \AgdaBound{γ}\AgdaSymbol{)} \AgdaInductiveConstructor{,} \AgdaSymbol{(λ} \AgdaBound{a} \AgdaBound{b} \AgdaBound{p} \AgdaSymbol{→} \<[54]%
\>[54]\<%
\\
\>[11]\AgdaIndent{18}{}\<[18]%
\>[18]\AgdaFunction{[} \AgdaFunction{[} \AgdaBound{B} \AgdaFunction{]fm} \AgdaSymbol{\_} \AgdaFunction{]trans} \AgdaFunction{[} \AgdaBound{B} \AgdaFunction{]refl*} \AgdaSymbol{(}\AgdaBound{frespR} \AgdaSymbol{\_} \AgdaSymbol{\_} \AgdaSymbol{\_} \AgdaBound{p}\AgdaSymbol{))}\<%
\\
\>[0]\AgdaIndent{11}{}\<[11]%
\>[11]\AgdaSymbol{;} \AgdaField{respt} \AgdaSymbol{=} \AgdaSymbol{λ} \AgdaSymbol{\{}\AgdaBound{γ}\AgdaSymbol{\}} \AgdaSymbol{\{}\AgdaBound{γ'}\AgdaSymbol{\}} \AgdaBound{p} \AgdaBound{a} \AgdaSymbol{→} \AgdaFunction{[} \AgdaBound{f} \AgdaFunction{]respt} \AgdaBound{p} \AgdaBound{a}\<%
\\
\>[0]\AgdaIndent{11}{}\<[11]%
\>[11]\AgdaSymbol{\}}\<%
\\
%
\\
\>[0]\AgdaIndent{2}{}\<[2]%
\>[2]\<%
\\
\>[0]\AgdaIndent{2}{}\<[2]%
\>[2]\AgdaFunction{substQ} \AgdaSymbol{:} \AgdaSymbol{(}\AgdaBound{Γ} \AgdaFunction{\&} \AgdaBound{A}\AgdaSymbol{)} \AgdaRecord{⇉} \AgdaSymbol{(}\AgdaBound{Γ} \AgdaFunction{\&} \AgdaFunction{⟦Q⟧}\AgdaSymbol{)}\<%
\\
\>[0]\AgdaIndent{2}{}\<[2]%
\>[2]\AgdaFunction{substQ} \AgdaSymbol{=} \AgdaKeyword{record}\<%
\\
\>[2]\AgdaIndent{11}{}\<[11]%
\>[11]\AgdaSymbol{\{} \AgdaField{fn} \AgdaSymbol{=} \AgdaSymbol{λ} \AgdaSymbol{\{(}\AgdaBound{x} \AgdaInductiveConstructor{,} \AgdaBound{a}\AgdaSymbol{)} \AgdaSymbol{→} \AgdaBound{x} \AgdaInductiveConstructor{,} \AgdaBound{a}\AgdaSymbol{\}}\<%
\\
\>[2]\AgdaIndent{11}{}\<[11]%
\>[11]\AgdaSymbol{;} \AgdaField{resp} \AgdaSymbol{=} \AgdaSymbol{λ\{} \AgdaSymbol{(}\AgdaBound{p} \AgdaInductiveConstructor{,} \AgdaBound{q}\AgdaSymbol{)} \AgdaSymbol{→} \AgdaBound{p} \AgdaInductiveConstructor{,} \AgdaSymbol{(}\AgdaBound{Rrsp} \AgdaBound{q}\AgdaSymbol{)\}}\<%
\\
\>[2]\AgdaIndent{11}{}\<[11]%
\>[11]\AgdaSymbol{\}}\<%
\\
%
\\
\>[0]\AgdaIndent{2}{}\<[2]%
\>[2]\AgdaFunction{Q-ind} \AgdaSymbol{:} \AgdaSymbol{(}\AgdaBound{P} \AgdaSymbol{:} \AgdaRecord{Ty} \AgdaSymbol{(}\AgdaBound{Γ} \AgdaFunction{\&} \AgdaFunction{⟦Q⟧}\AgdaSymbol{))}\<%
\\
\>[0]\AgdaIndent{9}{}\<[9]%
\>[9]\AgdaSymbol{→} \AgdaSymbol{(}\AgdaBound{isProp} \AgdaSymbol{:} \AgdaSymbol{∀} \AgdaSymbol{\{}\AgdaBound{x} \AgdaBound{a}\AgdaSymbol{\}} \AgdaSymbol{(}\AgdaBound{r} \AgdaBound{s} \AgdaSymbol{:} \AgdaFunction{∣} \AgdaFunction{[} \AgdaBound{P} \AgdaFunction{]fm} \AgdaSymbol{(}\AgdaBound{x} \AgdaInductiveConstructor{,} \AgdaBound{a}\AgdaSymbol{)} \AgdaFunction{∣}\AgdaSymbol{)} \AgdaSymbol{→} \<[59]%
\>[59]\<%
\\
\>[9]\AgdaIndent{22}{}\<[22]%
\>[22]\AgdaFunction{[} \AgdaFunction{[} \AgdaBound{P} \AgdaFunction{]fm} \AgdaSymbol{(}\AgdaBound{x} \AgdaInductiveConstructor{,} \AgdaBound{a}\AgdaSymbol{)} \AgdaFunction{]} \AgdaBound{r} \AgdaFunction{≈} \AgdaBound{s} \AgdaSymbol{)}\<%
\\
\>[0]\AgdaIndent{9}{}\<[9]%
\>[9]\AgdaSymbol{→} \AgdaSymbol{(}\AgdaBound{h} \AgdaSymbol{:} \AgdaRecord{Tm} \AgdaSymbol{(}\AgdaFunction{Π} \AgdaBound{A} \AgdaSymbol{(}\AgdaBound{P} \AgdaFunction{[} \AgdaFunction{substQ} \AgdaFunction{]T}\AgdaSymbol{)))}\<%
\\
\>[0]\AgdaIndent{9}{}\<[9]%
\>[9]\AgdaSymbol{→} \AgdaRecord{Tm} \AgdaSymbol{(}\AgdaFunction{Π} \AgdaFunction{⟦Q⟧} \AgdaBound{P}\AgdaSymbol{)}\<%
\\
\>[0]\AgdaIndent{2}{}\<[2]%
\>[2]\AgdaFunction{Q-ind} \AgdaBound{P} \AgdaBound{isProp} \AgdaBound{h} \AgdaSymbol{=} \AgdaKeyword{record}\<%
\\
\>[0]\AgdaIndent{11}{}\<[11]%
\>[11]\AgdaSymbol{\{} \AgdaField{tm} \AgdaSymbol{=} \AgdaSymbol{λ} \AgdaBound{x} \AgdaSymbol{→} \AgdaSymbol{(}\AgdaFunction{prj₁} \AgdaSymbol{(}\AgdaFunction{[} \AgdaBound{h} \AgdaFunction{]tm} \AgdaBound{x}\AgdaSymbol{))} \AgdaInductiveConstructor{,} \<[45]%
\>[45]\<%
\\
\>[11]\AgdaIndent{18}{}\<[18]%
\>[18]\AgdaSymbol{(λ} \AgdaBound{a} \AgdaBound{b} \AgdaBound{p} \AgdaSymbol{→} \AgdaBound{isProp} \AgdaSymbol{\{}\AgdaBound{x}\AgdaSymbol{\}} \AgdaSymbol{\{}\AgdaBound{b}\AgdaSymbol{\}} \AgdaSymbol{\_} \AgdaSymbol{\_)}\<%
\\
\>[0]\AgdaIndent{11}{}\<[11]%
\>[11]\AgdaSymbol{;} \AgdaField{respt} \AgdaSymbol{=} \AgdaFunction{[} \AgdaBound{h} \AgdaFunction{]respt}\<%
\\
\>[0]\AgdaIndent{11}{}\<[11]%
\>[11]\AgdaSymbol{\}}\<%
\\
\>\<\end{code}






\section{Quotient types in setoid model}





\section{Observational equality}

%definitional distinct types

Later in in \cite{alti:ott-conf}, Altenkirch and McBride further
simplifies the setoid model by adopting McBride's heterogeneous
approach to equality. They identifies values up to observation rather than
  construction which is called observational equality. It is the
  propositional equality induced by the Setoid model.  In general we have a heterogeneous equality which
  compares terms of types which are different in construction. It only
  make sense when we can prove the types are the same. It helps us
  avoids the heavy use of $subst$ which makes formalisation and
  reasoning involved. We could simplify the setoid model by adapting this
  approach and the implementation could be easier.


%\chapter{Homotopy Type Theory and higher inductive types}\label{HITs}


\chapter{Models of Type Theory}\label{models}


To introduce extensional concepts in \itt, one can simply postulate them as axioms but the good computational properties of Type Theory will be lost. Therefore it is crucial that these axioms have 
a computational interpretation. One solution is to construct a model of theses axioms in \itt or in a constructive setting such that extensional concepts like functional extensionality, quotient types, univalence are automatically derivable. Types as usually interpreted as structured objects rather than sets, and the model is named after the structured objects, for example, setoid, groupoid, simplicial sets. 


In this chapter, we discuss several models of extensional concepts, and mainly introduce Altenkirch's setoid model. We build a category with families of setoids to accommodate the types theory described in
\cite{alti:lics99}  so that it is possible to define quotient types following Martin Hofmann's Paper \cite{hof:95:sm}. We also briefly introduce several models of \hott where quotient types are also available.

\section{Setoid model}


One of the intensional models of extensional concepts inside \itt is setoid model where types are interpreted as setoids. Martin Hofmann has studied this approach in \cite{hof:phd}, but a naive version of setoid model does not satisfy all definitional equalities. The version in \cite{DBLP:conf/tlca/Hofmann95} which is a model for quotient types does not allow large eliminations (defining a dependent type by recursion).

Altenkirch then proposes \cite{alti:lics99} a different approach in which the setoid model serves as the metatheory.
He uses an extension of \itt by a universe of propositions $\Prop$ as metatheory, and the $\eta$-rules for $\Pi$-types and $\Sigma$-types hold. 

%In \cite{DBLP:conf/csl/Hofmann94}, Martin Hofmann first proposes a setoid model where types are interpreted as setoids, and the function are equivalence preserving. Cateogically, a setoid is a special groupoid where every isomorphism is unique, and then functions between two setoids are just functors.


\infrule[proof-irr]{\Gamma \vdash P : \Prop \andalso \Gamma \vdash p,q : P}{\Gamma \vdash p \equiv q : P}



 $\Prop$ only contains "propositional'' sets which has at most one
inhabitant. Notice that it is not a definition of types, which means
that we cannot conclude a type is of type \textbf{Prop} if we have a
proof that all inhabitants of it are definitionally equal.

The propositional universe is closed under $\Pi$-types and $\Sigma$-types:



\infrule[$\Pi$-Prop]{\Gamma \vdash A : \Set \andalso \Gamma, x : A \vdash P : \Prop}
{\Gamma \vdash \Pi~ (x : A) \to P : \Prop}



\infrule[$\Sigma$-Prop]{\Gamma \vdash P : \Prop \andalso \Gamma,x : P \vdash Q : \Prop}
{\Gamma \vdash \Sigma ~(x : P) ~ Q : \Prop}


The metatheory is proved:

\begin{itemize}
\item Decidable. The definitional equality is decideable, hence type checking is decidable.

\item Consistent. Not all types inhabited and not all well typed definitionality equality holds. 

\item Adequate. All terms of type $\N$ are reducible to numerals.
\end{itemize}


And then Altenkirch constrct an intensional model within this metatheory which is decidable and adequate, functional extensionality is inhabited and it permits large elimination. 


% Within this type theory, introduction of quotient types is straightforward. 
%The set of functions are naturally quotient types, the hidden information is the definition of the functions and the equivalence relation is the functional extensionality.



\begin{remark}[The setoid model is not LCCC]
This model is different to a setoid model as an E-category, for instance
the one introduced by Hofmann \cite{hofmann1995interpretation}. An E-category is a category equipped with
an equivalence relation for homsets. To distinguish them, we call this
category \textbf{E-setoids}.  All morphisms of \textbf{E-setoids}
gives rise to types and they are cartesian closed, namely it is a a locally
cartesian closed category (LCCC). However not all morphisms in our category of setoids give rise to types which means that it is not an LCCC. Every LCCC can serve as a model for categories with
families but not every category with families has to be a LCCC. 

This setoid model is the category of setoids \textbf{Std} which is a full subcategory of \textbf{Gpd} (the category of small groupoids). Every object of \textbf{Gpd} whose all homsets contain at most one morphism are in this subcategory. Altenkirch and Klaus \cite{Altenkirch12setoidsare} prove that both \textbf{Gpd} and \textbf{Std} are cartesian closed but not locally cartesian closed.
\end{remark}


\section{Category with families}


The setoid model is defined as categories with families as introduced by Dybjer \cite{Dyb:96} and Hofmann
\cite{hof:97}. The object theory is decidable because its definitional equalities are interpreted by definitional equality in the metatheory which is decidable.

We construct the setoid model in Agda together with some basic types including quotient types. We will introduce our implementation with real Agda code. We omit some unnecessary part of the code, but the complete code can be found in Appendix \ref{cwf}.


\subsection{Metatheory}

We first postulate axioms for metatheory. We define \textbf{HProp} which serves as the proof-irrelevant universe of propositions although it is not exactly the same as the \textbf{Prop} universe which is judgemental. A set can only be claimed in $\textbf{Prop}$. However any set behaves like a proposition can be encoded as \textbf{HProp} but not in \textbf{Prop}. Another difference is that the uniqueness of identity proof is propositional which complicates the proving.

\begin{code}\>\<%
\\
\>\AgdaKeyword{record} \AgdaRecord{HProp} \AgdaSymbol{:} \AgdaPrimitiveType{Set₁} \AgdaKeyword{where}\<%
\\
\>[0]\AgdaIndent{2}{}\<[2]%
\>[2]\AgdaKeyword{constructor} \AgdaInductiveConstructor{hProp}\<%
\\
\>[0]\AgdaIndent{2}{}\<[2]%
\>[2]\AgdaKeyword{field}\<%
\\
\>[2]\AgdaIndent{4}{}\<[4]%
\>[4]\AgdaField{prf} \AgdaSymbol{:} \AgdaPrimitiveType{Set}\<%
\\
\>[2]\AgdaIndent{4}{}\<[4]%
\>[4]\AgdaField{Uni} \AgdaSymbol{:} \AgdaSymbol{\{}\AgdaBound{p} \AgdaBound{q} \AgdaSymbol{:} \AgdaBound{prf}\AgdaSymbol{\}} \AgdaSymbol{→} \AgdaBound{p} \AgdaDatatype{≡} \AgdaBound{q}\<%%
\\
\>\AgdaKeyword{open} \AgdaModule{HProp} \AgdaKeyword{public} \AgdaKeyword{renaming} \AgdaSymbol{(}prf \AgdaSymbol{to} <\_>\AgdaSymbol{)}\<%
\end{code}

We also define basic propositions $\top$ and $\bot$ and the universal and existential quantifier, namely it is closed under $\Pi$-types and $\Sigma$-types. Notice that, the $\Pi$-closure requires functional extensionality which can be postulated since $\Pi$-closure is itself an axiom in this model.

\begin{code}\>\<
%
\\
\>\AgdaFunction{∀'} \AgdaSymbol{:} \AgdaSymbol{(}\AgdaBound{A} \AgdaSymbol{:} \AgdaPrimitiveType{Set}\AgdaSymbol{)(}\AgdaBound{P} \AgdaSymbol{:} \AgdaBound{A} \AgdaSymbol{→} \AgdaRecord{HProp}\AgdaSymbol{)} \AgdaSymbol{→} \AgdaRecord{HProp}\<%
\\
\>\AgdaFunction{∀'} \AgdaBound{A} \AgdaBound{P} \AgdaSymbol{=} \AgdaInductiveConstructor{hProp} \AgdaSymbol{((}\AgdaBound{x} \AgdaSymbol{:} \AgdaBound{A}\AgdaSymbol{)} \AgdaSymbol{→} \AgdaFunction{<} \AgdaBound{P} \AgdaBound{x} \AgdaFunction{>}\AgdaSymbol{)} \AgdaSymbol{(}\AgdaBound{ext} \AgdaSymbol{(λ} \AgdaBound{x} \AgdaSymbol{→} \AgdaFunction{Uni} \AgdaSymbol{(}\AgdaBound{P} \AgdaBound{x}\AgdaSymbol{)))}\<%
\>\<\end{code}

Notice that we use a syntactic abbreviation $<\_>$ for projection of proof in $\textbf{HProp}$. Moever to avoid conficts with existing code, we sometimes add a $'$ symbol in names. 


\subsection{Category of setoids}

Then setoids can be defined as follows.

\begin{code}\>\<%
\\
\>\AgdaKeyword{record} \AgdaRecord{HSetoid} \AgdaSymbol{:} \AgdaPrimitiveType{Set₁} \AgdaKeyword{where}\<%
\\
\>[0]\AgdaIndent{2}{}\<[2]%
\>[2]\AgdaKeyword{field}\<%
\\
\>[2]\AgdaIndent{4}{}\<[4]%
\>[4]\AgdaField{Carrier} \AgdaSymbol{:} \AgdaPrimitiveType{Set}\<%
\\
\>[2]\AgdaIndent{4}{}\<[4]%
\>[4]\AgdaField{\_≈h\_} \<[13]%
\>[13]\AgdaSymbol{:} \AgdaBound{Carrier} \AgdaSymbol{→} \AgdaBound{Carrier} \AgdaSymbol{→} \AgdaRecord{HProp}\<%
\\
\>[2]\AgdaIndent{4}{}\<[4]%
\>[4]\AgdaField{refl} \<[12]%
\>[12]\AgdaSymbol{:} \AgdaSymbol{\{}\AgdaBound{x} \AgdaSymbol{:} \AgdaBound{Carrier}\AgdaSymbol{\}} \AgdaSymbol{→} \AgdaFunction{<} \AgdaBound{x} \AgdaBound{≈h} \AgdaBound{x} \AgdaFunction{>}\<%
\\
\>[2]\AgdaIndent{4}{}\<[4]%
\>[4]\AgdaField{sym} \<[12]%
\>[12]\AgdaSymbol{:} \AgdaSymbol{\{}\AgdaBound{x} \AgdaBound{y} \AgdaSymbol{:} \AgdaBound{Carrier}\AgdaSymbol{\}} \AgdaSymbol{→} \AgdaFunction{<} \AgdaBound{x} \AgdaBound{≈h} \AgdaBound{y} \AgdaFunction{>} \AgdaSymbol{→} \AgdaFunction{<} \AgdaBound{y} \AgdaBound{≈h} \AgdaBound{x} \AgdaFunction{>}\<%
\\
\>[2]\AgdaIndent{4}{}\<[4]%
\>[4]\AgdaField{trans} \<[12]%
\>[12]\AgdaSymbol{:} \AgdaSymbol{\{}\AgdaBound{x} \AgdaBound{y} \AgdaBound{z} \AgdaSymbol{:} \AgdaBound{Carrier}\AgdaSymbol{\}} \AgdaSymbol{→} \AgdaFunction{<} \AgdaBound{x} \AgdaBound{≈h} \AgdaBound{y} \AgdaFunction{>} \AgdaSymbol{→} \AgdaFunction{<} \AgdaBound{y} \AgdaBound{≈h} \AgdaBound{z} \AgdaFunction{>} \AgdaSymbol{→} \AgdaFunction{<} \AgdaBound{x} \AgdaBound{≈h} \AgdaBound{z} \AgdaFunction{>}\<%
\\
\>\<\end{code}

Notice that since \textbf{HProp} is not considered as a sort of types, we simpified the definition by defining the laws non-proposionally, e.g.\ reflexivity.



Functors between setoids can be given as

\begin{code}
\\
\>\AgdaKeyword{record} \AgdaRecord{\_⇉\_} \AgdaSymbol{(}\AgdaBound{A} \AgdaBound{B} \AgdaSymbol{:} \AgdaRecord{HSetoid}\AgdaSymbol{)} \AgdaSymbol{:} \AgdaPrimitiveType{Set} \AgdaKeyword{where}\<%
\\
\>[0]\AgdaIndent{2}{}\<[2]%
\>[2]\AgdaKeyword{field}\<%
\\
\>[2]\AgdaIndent{4}{}\<[4]%
\>[4]\AgdaField{fn} \<[9]%
\>[9]\AgdaSymbol{:} \AgdaFunction{∣} \AgdaBound{A} \AgdaFunction{∣} \AgdaSymbol{→} \AgdaFunction{∣} \AgdaBound{B} \AgdaFunction{∣}\<%
\\
\>[2]\AgdaIndent{4}{}\<[4]%
\>[4]\AgdaField{resp} \AgdaSymbol{:} \AgdaSymbol{\{}\AgdaBound{x} \AgdaBound{y} \AgdaSymbol{:} \AgdaFunction{∣} \AgdaBound{A} \AgdaFunction{∣}\AgdaSymbol{\}} \AgdaSymbol{→} \<[27]%
\>[27]\<%
\\
\>[4]\AgdaIndent{11}{}\<[11]%
\>[11]\AgdaFunction{[} \AgdaBound{A} \AgdaFunction{]} \AgdaBound{x} \AgdaFunction{≈} \AgdaBound{y} \AgdaSymbol{→} \<[25]%
\>[25]\<%
\\
\>[4]\AgdaIndent{11}{}\<[11]%
\>[11]\AgdaFunction{[} \AgdaBound{B} \AgdaFunction{]} \AgdaBound{fn} \AgdaBound{x} \AgdaFunction{≈} \AgdaBound{fn} \AgdaBound{y}\<%
\\
\end{code}

%大略介绍thorsten的model和martin hofmann的Quotient types就可以了

With identity morphisms, composition operator and all other ingredients we can build the category of setoids of $\textbf{Std}$.


\subsection{Categories with families}

Usually a a categories with families is defined as follows:

\begin{definition}
\textbf{Categories with families}.

\begin{itemize}
\item A category $\mathsf{C}$ with a terminal object where objects are contexts and morphisms are substitutions.

\item A functor $F : \mathsf{C}^{op} \to Fam$ where $Fam$ is a categories of families whose objects are dependent products and morphisms are naturally pairs of functions. The types and terms are projections of this functor: given an object $\Gamma : \mathsf{C}$, we usually write $F ~ \Gamma \defeq \Sigma A : \text{Ty} ~\Gamma, \text{Tm} ~\Gamma ~ A)$, and the substitution of types and terms are just contained in the morphism part of this functor. 

\item A comprehension of $\Gamma$ and $A : \text{Ty}~ \Gamma$ written as $\Gamma, A$ is a construction of a new object in $\mathsf{C}$ such that it expresses the extension of contexts.
\end{itemize}

\end{definition}

In the setoid model, both the category $\mathsf{C}$ and $Fam$ are not based on sets, but setoids.
The cateogry of contexts is just $\textbf{Std}$,

\begin{code}
%
\\
\>\AgdaFunction{Con} \AgdaSymbol{=} \AgdaRecord{HSetoid}\<%
\\
\end{code}

The second ingredient is the functors, namely types and terms.
Given a context $\Gamma$, types of the context can be defined as functors from $\Gamma \to \textbf{Std}$ since a context is a setoid which are special cases of groupoids. However setoids are not implemented as categories, we manually build \emph{functors} which respects the internal equivalence of setoids:

\begin{code}%
\\
\>\AgdaKeyword{record} \AgdaRecord{Ty} \AgdaSymbol{(}\AgdaBound{Γ} \AgdaSymbol{:} \AgdaFunction{Con}\AgdaSymbol{)} \AgdaSymbol{:} \AgdaPrimitiveType{Set₁} \AgdaKeyword{where}\<%
\\
\>[0]\AgdaIndent{2}{}\<[2]%
\>[2]\AgdaKeyword{field}\<%
\\
\>[2]\AgdaIndent{4}{}\<[4]%
\>[4]\AgdaField{fm} \<[11]%
\>[11]\AgdaSymbol{:} \AgdaFunction{∣} \AgdaBound{Γ} \AgdaFunction{∣} \AgdaSymbol{→} \AgdaRecord{HSetoid}\<%
\\
%
\\
\>[2]\AgdaIndent{4}{}\<[4]%
\>[4]\AgdaField{substT} \AgdaSymbol{:} \AgdaSymbol{\{}\AgdaBound{x} \AgdaBound{y} \AgdaSymbol{:} \AgdaFunction{∣} \AgdaBound{Γ} \AgdaFunction{∣}\AgdaSymbol{\}} \AgdaSymbol{→} \<[29]%
\>[29]\<%
\\
\>[4]\AgdaIndent{13}{}\<[13]%
\>[13]\AgdaFunction{[} \AgdaBound{Γ} \AgdaFunction{]} \AgdaBound{x} \AgdaFunction{≈} \AgdaBound{y} \AgdaSymbol{→}\<%
\\
\>[4]\AgdaIndent{13}{}\<[13]%
\>[13]\AgdaFunction{∣} \AgdaBound{fm} \AgdaBound{x} \AgdaFunction{∣} \AgdaSymbol{→} \AgdaFunction{∣} \AgdaBound{fm} \AgdaBound{y} \AgdaFunction{∣}\<%
\\
%
\\
\>[0]\AgdaIndent{4}{}\<[4]%
\>[4]\AgdaField{subst*} \AgdaSymbol{:} \AgdaSymbol{∀\{}\AgdaBound{x} \AgdaBound{y} \AgdaSymbol{:} \AgdaFunction{∣} \AgdaBound{Γ} \AgdaFunction{∣}\AgdaSymbol{\}}\<%
\\
\>[0]\AgdaIndent{13}{}\<[13]%
\>[13]\AgdaSymbol{(}\AgdaBound{p} \AgdaSymbol{:} \AgdaFunction{[} \AgdaBound{Γ} \AgdaFunction{]} \AgdaBound{x} \AgdaFunction{≈} \AgdaBound{y}\AgdaSymbol{)}\<%
\\
\>[0]\AgdaIndent{13}{}\<[13]%
\>[13]\AgdaSymbol{\{}\AgdaBound{a} \AgdaBound{b} \AgdaSymbol{:} \AgdaFunction{∣} \AgdaBound{fm} \AgdaBound{x} \AgdaFunction{∣}\AgdaSymbol{\}} \AgdaSymbol{→}\<%
\\
\>[0]\AgdaIndent{13}{}\<[13]%
\>[13]\AgdaFunction{[} \AgdaBound{fm} \AgdaBound{x} \AgdaFunction{]} \AgdaBound{a} \AgdaFunction{≈} \AgdaBound{b} \AgdaSymbol{→}\<%
\\
\>[0]\AgdaIndent{13}{}\<[13]%
\>[13]\AgdaFunction{[} \AgdaBound{fm} \AgdaBound{y} \AgdaFunction{]} \AgdaBound{substT} \AgdaBound{p} \AgdaBound{a} \AgdaFunction{≈} \AgdaBound{substT} \AgdaBound{p} \AgdaBound{b}\<%
\\
%
\\
\>[0]\AgdaIndent{4}{}\<[4]%
\>[4]\AgdaField{refl*} \<[11]%
\>[11]\AgdaSymbol{:} \AgdaSymbol{∀(}\AgdaBound{x} \AgdaSymbol{:} \AgdaFunction{∣} \AgdaBound{Γ} \AgdaFunction{∣}\AgdaSymbol{)(}\AgdaBound{a} \AgdaSymbol{:} \AgdaFunction{∣} \AgdaBound{fm} \AgdaBound{x} \AgdaFunction{∣}\AgdaSymbol{)} \AgdaSymbol{→} \<[42]%
\>[42]\<%
\\
\>[0]\AgdaIndent{13}{}\<[13]%
\>[13]\AgdaFunction{[} \AgdaBound{fm} \AgdaBound{x} \AgdaFunction{]} \AgdaBound{substT} \AgdaFunction{[} \AgdaBound{Γ} \AgdaFunction{]refl} \AgdaBound{a} \AgdaFunction{≈} \AgdaBound{a}\<%
\\
%
\\
\>[0]\AgdaIndent{4}{}\<[4]%
\>[4]\AgdaField{trans*} \AgdaSymbol{:} \AgdaSymbol{∀\{}\AgdaBound{x} \AgdaBound{y} \AgdaBound{z} \AgdaSymbol{:} \AgdaFunction{∣} \AgdaBound{Γ} \AgdaFunction{∣}\AgdaSymbol{\}}\<%
\\
\>[0]\AgdaIndent{13}{}\<[13]%
\>[13]\AgdaSymbol{\{}\AgdaBound{p} \AgdaSymbol{:} \AgdaFunction{[} \AgdaBound{Γ} \AgdaFunction{]} \AgdaBound{x} \AgdaFunction{≈} \AgdaBound{y}\AgdaSymbol{\}\{}\AgdaBound{q} \AgdaSymbol{:} \AgdaFunction{[} \AgdaBound{Γ} \AgdaFunction{]} \AgdaBound{y} \AgdaFunction{≈} \AgdaBound{z}\AgdaSymbol{\}}\<%
\\
\>[0]\AgdaIndent{13}{}\<[13]%
\>[13]\AgdaSymbol{(}\AgdaBound{a} \AgdaSymbol{:} \AgdaFunction{∣} \AgdaBound{fm} \AgdaBound{x} \AgdaFunction{∣}\AgdaSymbol{)} \AgdaSymbol{→} \<[30]%
\>[30]\<%
\\
\>[0]\AgdaIndent{13}{}\<[13]%
\>[13]\AgdaFunction{[} \AgdaBound{fm} \AgdaBound{z} \AgdaFunction{]} \AgdaBound{substT} \AgdaBound{q} \AgdaSymbol{(}\AgdaBound{substT} \AgdaBound{p} \AgdaBound{a}\AgdaSymbol{)} \AgdaFunction{≈} \AgdaBound{substT} \AgdaSymbol{(}\AgdaFunction{[} \AgdaBound{Γ} \AgdaFunction{]trans} \AgdaBound{p} \AgdaBound{q}\AgdaSymbol{)} \AgdaBound{a}\<%
\\
\end{code}

Notice that there is a problem arising from the fact that $HProp$ is not definitionally proof-irrelevant which means that the substitution here is not automatically proof-irrelevant. Therefore we have to define a series of proof-irrelevance lemmas to help construction and reasoning which are omitted here.

Terms follows naturally as families of elements in the underlying set of types indexed by $\gamma : \Gamma$ which respects the equality as well:


\begin{code}
%
\\
\>\AgdaKeyword{record} \AgdaRecord{Tm} \AgdaSymbol{\{}\AgdaBound{Γ} \AgdaSymbol{:} \AgdaFunction{Con}\AgdaSymbol{\}(}\AgdaBound{A} \AgdaSymbol{:} \AgdaRecord{Ty} \AgdaBound{Γ}\AgdaSymbol{)} \AgdaSymbol{:} \AgdaPrimitiveType{Set} \AgdaKeyword{where}\<%
\\
\>[0]\AgdaIndent{2}{}\<[2]%
\>[2]\AgdaKeyword{constructor} \AgdaInductiveConstructor{tm:\_resp:\_}\<%
\\
\>[0]\AgdaIndent{2}{}\<[2]%
\>[2]\AgdaKeyword{field}\<%
\\
\>[2]\AgdaIndent{4}{}\<[4]%
\>[4]\AgdaField{tm} \<[10]%
\>[10]\AgdaSymbol{:} \AgdaSymbol{(}\AgdaBound{x} \AgdaSymbol{:} \AgdaFunction{∣} \AgdaBound{Γ} \AgdaFunction{∣}\AgdaSymbol{)} \AgdaSymbol{→} \AgdaFunction{∣} \AgdaFunction{[} \AgdaBound{A} \AgdaFunction{]fm} \AgdaBound{x} \AgdaFunction{∣}\<%
\\
\>[2]\AgdaIndent{4}{}\<[4]%
\>[4]\AgdaField{respt} \AgdaSymbol{:} \AgdaSymbol{∀} \AgdaSymbol{\{}\AgdaBound{x} \AgdaBound{y} \AgdaSymbol{:} \AgdaFunction{∣} \AgdaBound{Γ} \AgdaFunction{∣}\AgdaSymbol{\}} \AgdaSymbol{→} \<[30]%
\>[30]\<%
\\
\>[4]\AgdaIndent{10}{}\<[10]%
\>[10]\AgdaSymbol{(}\AgdaBound{p} \AgdaSymbol{:} \AgdaFunction{[} \AgdaBound{Γ} \AgdaFunction{]} \AgdaBound{x} \AgdaFunction{≈} \AgdaBound{y}\AgdaSymbol{)} \AgdaSymbol{→} \<[30]%
\>[30]\<%
\\
\>[4]\AgdaIndent{10}{}\<[10]%
\>[10]\AgdaFunction{[} \AgdaFunction{[} \AgdaBound{A} \AgdaFunction{]fm} \AgdaBound{y} \AgdaFunction{]} \AgdaFunction{[} \AgdaBound{A} \AgdaFunction{]subst} \AgdaBound{p} \AgdaSymbol{(}\AgdaBound{tm} \AgdaBound{x}\AgdaSymbol{)} \AgdaFunction{≈} \AgdaBound{tm} \AgdaBound{y}\<%
\\
\end{code}


Due to the record type mechanism in Agda, given a type $A : \text{Ty} ~\Gamma$, one has to write $\text{substT} ~A$ to use the field. For readability, we adopt a better syntax using renaming of fields to highlight the precise types or terms in the style of $\text{[} ~A ~\text{]substT}$ , but the detailed code is also omitted here. Indeed, an example can be found in the above code where we already adopt the syntax stype for internal equivalence of setoids.


The substitution of types can be defined simply by composing the underlying objects of types and context morphisms:

\begin{code}
%
\\
\>\AgdaFunction{\_[\_]T} \AgdaSymbol{:} \AgdaSymbol{∀} \AgdaSymbol{\{}\AgdaBound{Γ} \AgdaBound{Δ} \AgdaSymbol{:} \AgdaFunction{Con}\AgdaSymbol{\}} \AgdaSymbol{→} \AgdaRecord{Ty} \AgdaBound{Δ} \AgdaSymbol{→} \AgdaBound{Γ} \AgdaRecord{⇉} \AgdaBound{Δ} \AgdaSymbol{→} \AgdaRecord{Ty} \AgdaBound{Γ}\<%
\\
\>\AgdaFunction{\_[\_]T} \AgdaSymbol{\{}\AgdaBound{Γ}\AgdaSymbol{\}} \AgdaSymbol{\{}\AgdaBound{Δ}\AgdaSymbol{\}} \AgdaBound{A} \AgdaBound{f}\<%
\\
\>[2]\AgdaIndent{5}{}\<[5]%
\>[5]\AgdaSymbol{=} \AgdaKeyword{record}\<%
\\
\>[2]\AgdaIndent{5}{}\<[5]%
\>[5]\AgdaSymbol{\{} \AgdaField{fm} \<[14]%
\>[14]\AgdaSymbol{=} \AgdaFunction{fm} \AgdaFunction{∘} \AgdaFunction{fn}\<%
\\
\>[2]\AgdaIndent{5}{}\<[5]%
\>[5]\AgdaSymbol{;} \AgdaField{substT} \AgdaSymbol{=} \AgdaFunction{substT} \AgdaFunction{∘} \AgdaFunction{resp}\<%
\\
\>[2]\AgdaIndent{5}{}\<[5]%
\>[5]\AgdaSymbol{;} \AgdaField{subst*} \AgdaSymbol{=} \AgdaFunction{subst*} \AgdaFunction{∘} \AgdaFunction{resp}\<%
\\
\>[2]\AgdaIndent{5}{}\<[5]%
\>[5]\AgdaSymbol{;} \AgdaField{refl*} \<[14]%
\>[14]\AgdaSymbol{=} \AgdaSymbol{λ} \AgdaBound{\_} \AgdaBound{\_} \AgdaSymbol{→} \AgdaFunction{subst-pi'}\<%
\\
\>[2]\AgdaIndent{5}{}\<[5]%
\>[5]\AgdaSymbol{;} \AgdaField{trans*} \AgdaSymbol{=} \AgdaSymbol{λ} \AgdaBound{\_} \AgdaSymbol{→} \AgdaFunction{[} \AgdaFunction{fm} \AgdaSymbol{(}\AgdaFunction{fn} \AgdaSymbol{\_)} \AgdaFunction{]trans} \AgdaSymbol{(}\AgdaFunction{trans*} \AgdaSymbol{\_)} \AgdaFunction{subst-pi}\<%
\\
\>[2]\AgdaIndent{5}{}\<[5]%
\>[5]\AgdaSymbol{\}}\<%
\\
\>[2]\AgdaIndent{5}{}\<[5]%
\>[5]\AgdaKeyword{where} \<[11]%
\>[11]\<%
\\
\>[5]\AgdaIndent{7}{}\<[7]%
\>[7]\AgdaKeyword{open} \AgdaModule{Ty} \AgdaBound{A}\<%
\\
\>[5]\AgdaIndent{7}{}\<[7]%
\>[7]\AgdaKeyword{open} \AgdaModule{\_⇉\_} \AgdaBound{f}\<%
\\
\end{code}

The substitution of terms are similarly natural so that we omit it here.


To complete the categories with families, the terminal object is just empty context
\begin{code}
%
\\
\>\AgdaFunction{●} \AgdaSymbol{:} \AgdaRecord{HSetoid}\<%
\\
\>\AgdaFunction{●} \<[4]%
\>[4]\AgdaSymbol{=} \AgdaKeyword{record} \AgdaSymbol{\{}\<%
\\
\>[0]\AgdaIndent{6}{}\<[6]%
\>[6]\AgdaField{Carrier} \AgdaSymbol{=} \AgdaRecord{⊤}\AgdaSymbol{;}\<%
\\
\>[0]\AgdaIndent{6}{}\<[6]%
\>[6]\AgdaField{\_≈h\_} \<[14]%
\>[14]\AgdaSymbol{=} \AgdaSymbol{λ} \AgdaBound{\_} \AgdaBound{\_} \AgdaSymbol{→} \AgdaFunction{⊤'}\AgdaSymbol{;}\<%
\\
\>[0]\AgdaIndent{6}{}\<[6]%
\>[6]\AgdaField{refl} \<[14]%
\>[14]\AgdaSymbol{=} \AgdaInductiveConstructor{tt}\AgdaSymbol{;}\<%
\\
\>[0]\AgdaIndent{6}{}\<[6]%
\>[6]\AgdaField{sym} \<[14]%
\>[14]\AgdaSymbol{=} \AgdaSymbol{λ} \AgdaBound{\_} \AgdaSymbol{→} \AgdaInductiveConstructor{tt}\AgdaSymbol{;}\<%
\\
\>[0]\AgdaIndent{6}{}\<[6]%
\>[6]\AgdaField{trans} \<[14]%
\>[14]\AgdaSymbol{=} \AgdaSymbol{λ} \AgdaBound{\_} \AgdaBound{\_} \AgdaSymbol{→} \AgdaInductiveConstructor{tt} \AgdaSymbol{\}}\<%
\\
\end{code}

and context comprehension is encoded as follows:

\begin{code}
%
\\
\>\AgdaFunction{\_\&\_} \AgdaSymbol{:} \AgdaSymbol{(}\AgdaBound{Γ} \AgdaSymbol{:} \AgdaFunction{Con}\AgdaSymbol{)} \AgdaSymbol{→} \AgdaRecord{Ty} \AgdaBound{Γ} \AgdaSymbol{→} \AgdaRecord{HSetoid}\<%
\\
\>\AgdaBound{Γ} \AgdaFunction{\&} \AgdaBound{A} \AgdaSymbol{=} \AgdaKeyword{record} \<[15]%
\>[15]\<%
\\
\>[0]\AgdaIndent{7}{}\<[7]%
\>[7]\AgdaSymbol{\{} \AgdaField{Carrier} \AgdaSymbol{=} \AgdaRecord{Σ[} \AgdaBound{x} \AgdaRecord{∶} \AgdaFunction{∣} \AgdaBound{Γ} \AgdaFunction{∣} \AgdaRecord{]} \AgdaFunction{∣} \AgdaFunction{fm} \AgdaBound{x} \AgdaFunction{∣}\<%
\\
\>[0]\AgdaIndent{7}{}\<[7]%
\>[7]\AgdaSymbol{;} \AgdaField{\_≈h\_} \<[15]%
\>[15]\AgdaSymbol{=} \AgdaSymbol{λ\{(}\AgdaBound{x} \AgdaInductiveConstructor{,} \AgdaBound{a}\AgdaSymbol{)} \AgdaSymbol{(}\AgdaBound{y} \AgdaInductiveConstructor{,} \AgdaBound{b}\AgdaSymbol{)} \AgdaSymbol{→} \<[37]%
\>[37]\<%
\\
\>[7]\AgdaIndent{17}{}\<[17]%
\>[17]\AgdaFunction{Σ'[} \AgdaBound{p} \AgdaFunction{∶} \AgdaBound{x} \AgdaFunction{≈h} \AgdaBound{y} \AgdaFunction{]} \AgdaFunction{[} \AgdaFunction{fm} \AgdaBound{y} \AgdaFunction{]} \AgdaFunction{substT} \AgdaBound{p} \AgdaBound{a} \AgdaFunction{≈h} \AgdaBound{b}\AgdaSymbol{\}}\<%
\\
\>[0]\AgdaIndent{7}{}\<[7]%
\>[7]\AgdaSymbol{;} \AgdaField{refl} \<[15]%
\>[15]\AgdaSymbol{=} \AgdaFunction{refl} \AgdaInductiveConstructor{,} \AgdaSymbol{(}\AgdaFunction{refl*} \AgdaSymbol{\_} \AgdaSymbol{\_)}\<%
\\
\>[0]\AgdaIndent{7}{}\<[7]%
\>[7]\AgdaSymbol{;} \AgdaField{sym} \<[15]%
\>[15]\AgdaSymbol{=} \AgdaSymbol{λ} \AgdaSymbol{\{(}\AgdaBound{p} \AgdaInductiveConstructor{,} \AgdaBound{q}\AgdaSymbol{)} \AgdaSymbol{→} \AgdaSymbol{(}\AgdaFunction{sym} \AgdaBound{p}\AgdaSymbol{)} \AgdaInductiveConstructor{,} \<[40]%
\>[40]\<%
\\
\>[7]\AgdaIndent{17}{}\<[17]%
\>[17]\AgdaFunction{[} \AgdaFunction{fm} \AgdaSymbol{\_} \AgdaFunction{]trans} \AgdaSymbol{(}\AgdaFunction{subst*} \AgdaSymbol{\_} \AgdaSymbol{(}\AgdaFunction{[} \AgdaFunction{fm} \AgdaSymbol{\_} \AgdaFunction{]sym} \AgdaBound{q}\AgdaSymbol{))} \AgdaFunction{trans-refl} \AgdaSymbol{\}}\<%
\\
\>[0]\AgdaIndent{7}{}\<[7]%
\>[7]\AgdaSymbol{;} \AgdaField{trans} \AgdaSymbol{=} \AgdaSymbol{λ} \AgdaSymbol{\{(}\AgdaBound{p} \AgdaInductiveConstructor{,} \AgdaBound{q}\AgdaSymbol{)} \AgdaSymbol{(}\AgdaBound{m} \AgdaInductiveConstructor{,} \AgdaBound{n}\AgdaSymbol{)} \AgdaSymbol{→}\<%
\\
\>[0]\AgdaIndent{17}{}\<[17]%
\>[17]\AgdaFunction{trans} \AgdaBound{p} \AgdaBound{m} \AgdaInductiveConstructor{,} \<[29]%
\>[29]\<%
\\
\>[0]\AgdaIndent{17}{}\<[17]%
\>[17]\AgdaFunction{[} \AgdaFunction{fm} \AgdaSymbol{\_} \AgdaFunction{]trans} \AgdaSymbol{(}\AgdaFunction{[} \AgdaFunction{fm} \AgdaSymbol{\_} \AgdaFunction{]trans} \AgdaSymbol{(}\AgdaFunction{[} \AgdaFunction{fm} \AgdaSymbol{\_} \AgdaFunction{]sym} \AgdaSymbol{(}\AgdaFunction{trans*} \AgdaSymbol{\_))}\<%
\\
\>[0]\AgdaIndent{17}{}\<[17]%
\>[17]\AgdaSymbol{(}\AgdaFunction{subst*} \AgdaSymbol{\_} \AgdaBound{q}\AgdaSymbol{))} \AgdaBound{n} \AgdaSymbol{\}}\<%
\\
\>[0]\AgdaIndent{7}{}\<[7]%
\>[7]\AgdaSymbol{\}}\<%
\\
\>[0]\AgdaIndent{7}{}\<[7]%
\>[7]\AgdaKeyword{where} \<[13]%
\>[13]\<%
\\
\>[7]\AgdaIndent{9}{}\<[9]%
\>[9]\AgdaKeyword{open} \AgdaModule{HSetoid} \AgdaBound{Γ}\<%
\\
\>[7]\AgdaIndent{9}{}\<[9]%
\>[9]\AgdaKeyword{open} \AgdaModule{Ty} \AgdaBound{A} \<[23]%
\>[23]\<%
\\
\end{code}

Since the context $\Gamma$ and type $A$ as record types are opened in the scope, we can unambiguously use fields such as $fm$ and $subst*$.

\subsection{$\Pi$-types and some basic types}

Dependent function types $\Pi$-types are essential in a depenent theory and serves as a fundamental test case of the type construction in the setoid model. 
The following code is an example of the complicated and nearly unreadable encoding of types in this model, but in fact it just follows the reasoning and construction in \cite{alti:lics99} with minor adaption. Later explanations where detailed real codes is not shown can be found in Appendix \ref{cwf}.

\begin{code}
%
\\
\>\AgdaFunction{Π} \AgdaSymbol{:} \AgdaSymbol{\{}\AgdaBound{Γ} \AgdaSymbol{:} \AgdaFunction{Con}\AgdaSymbol{\}(}\AgdaBound{A} \AgdaSymbol{:} \AgdaRecord{Ty} \AgdaBound{Γ}\AgdaSymbol{)(}\AgdaBound{B} \AgdaSymbol{:} \AgdaRecord{Ty} \AgdaSymbol{(}\AgdaBound{Γ} \AgdaFunction{\&} \AgdaBound{A}\AgdaSymbol{))} \AgdaSymbol{→} \AgdaRecord{Ty} \AgdaBound{Γ}\<%
\\
\>\AgdaFunction{Π} \AgdaSymbol{\{}\AgdaBound{Γ}\AgdaSymbol{\}} \AgdaBound{A} \AgdaBound{B} \AgdaSymbol{=} \AgdaKeyword{record} \<[19]%
\>[19]\<%
\\
\>[0]\AgdaIndent{2}{}\<[2]%
\>[2]\AgdaSymbol{\{} \AgdaField{fm} \AgdaSymbol{=} \AgdaSymbol{λ} \AgdaBound{x} \AgdaSymbol{→} \AgdaKeyword{let} \AgdaBound{Ax} \AgdaSymbol{=} \AgdaFunction{[} \AgdaBound{A} \AgdaFunction{]fm} \AgdaBound{x} \AgdaKeyword{in}\<%
\\
\>[0]\AgdaIndent{15}{}\<[15]%
\>[15]\AgdaKeyword{let} \AgdaBound{Bx} \AgdaSymbol{=} \AgdaSymbol{λ} \AgdaBound{a} \AgdaSymbol{→} \AgdaFunction{[} \AgdaBound{B} \AgdaFunction{]fm} \AgdaSymbol{(}\AgdaBound{x} \AgdaInductiveConstructor{,} \AgdaBound{a}\AgdaSymbol{)} \AgdaKeyword{in}\<%
\\
\>[0]\AgdaIndent{9}{}\<[9]%
\>[9]\AgdaKeyword{record}\<%
\\
\>[0]\AgdaIndent{9}{}\<[9]%
\>[9]\AgdaSymbol{\{} \AgdaField{Carrier} \AgdaSymbol{=} \AgdaRecord{Σ[} \AgdaBound{fn} \AgdaRecord{∶} \AgdaSymbol{((}\AgdaBound{a} \AgdaSymbol{:} \AgdaFunction{∣} \AgdaBound{Ax} \AgdaFunction{∣}\AgdaSymbol{)} \AgdaSymbol{→} \AgdaFunction{∣} \AgdaBound{Bx} \AgdaBound{a} \AgdaFunction{∣}\AgdaSymbol{)} \AgdaRecord{]}\<%
\\
\>[9]\AgdaIndent{21}{}\<[21]%
\>[21]\AgdaSymbol{((}\AgdaBound{a} \AgdaBound{b} \AgdaSymbol{:} \AgdaFunction{∣} \AgdaBound{Ax} \AgdaFunction{∣}\AgdaSymbol{)}\<%
\\
\>[9]\AgdaIndent{21}{}\<[21]%
\>[21]\AgdaSymbol{(}\AgdaBound{p} \AgdaSymbol{:} \AgdaFunction{[} \AgdaBound{Ax} \AgdaFunction{]} \AgdaBound{a} \AgdaFunction{≈} \AgdaBound{b}\AgdaSymbol{)} \AgdaSymbol{→}\<%
\\
\>[9]\AgdaIndent{21}{}\<[21]%
\>[21]\AgdaFunction{[} \AgdaBound{Bx} \AgdaBound{b} \AgdaFunction{]} \AgdaFunction{[} \AgdaBound{B} \AgdaFunction{]subst} \AgdaSymbol{(}\AgdaFunction{[} \AgdaBound{Γ} \AgdaFunction{]refl} \AgdaInductiveConstructor{,} \<[54]%
\>[54]\<%
\\
\>[9]\AgdaIndent{21}{}\<[21]%
\>[21]\AgdaFunction{[} \AgdaBound{Ax} \AgdaFunction{]trans} \AgdaSymbol{(}\AgdaFunction{[} \AgdaBound{A} \AgdaFunction{]refl*} \AgdaBound{x} \AgdaBound{a}\AgdaSymbol{)} \AgdaBound{p}\AgdaSymbol{)} \AgdaSymbol{(}\AgdaBound{fn} \AgdaBound{a}\AgdaSymbol{)} \AgdaFunction{≈} \AgdaBound{fn} \AgdaBound{b}\AgdaSymbol{)} \<[68]%
\>[68]\<%
\\
\>[0]\AgdaIndent{9}{}\<[9]%
\>[9]\AgdaSymbol{;} \AgdaField{\_≈h\_} \<[19]%
\>[19]\AgdaSymbol{=} \AgdaSymbol{λ\{(}\AgdaBound{f} \AgdaInductiveConstructor{,} \AgdaSymbol{\_)} \AgdaSymbol{(}\AgdaBound{g} \AgdaInductiveConstructor{,} \AgdaSymbol{\_)} \AgdaSymbol{→}\<%
\\
\>[0]\AgdaIndent{21}{}\<[21]%
\>[21]\AgdaFunction{∀'[} \AgdaBound{a} \AgdaFunction{∶} \AgdaSymbol{\_} \AgdaFunction{]} \AgdaFunction{[} \AgdaBound{Bx} \AgdaBound{a} \AgdaFunction{]} \AgdaBound{f} \AgdaBound{a} \AgdaFunction{≈h} \AgdaBound{g} \AgdaBound{a} \AgdaSymbol{\}}\<%
\\
\>[0]\AgdaIndent{9}{}\<[9]%
\>[9]\AgdaSymbol{;} \AgdaField{refl} \<[19]%
\>[19]\AgdaSymbol{=} \AgdaSymbol{λ} \AgdaBound{a} \AgdaSymbol{→} \AgdaFunction{[} \AgdaBound{Bx} \AgdaBound{a} \AgdaFunction{]refl} \<[40]%
\>[40]\<%
\\
\>[0]\AgdaIndent{9}{}\<[9]%
\>[9]\AgdaSymbol{;} \AgdaField{sym} \<[19]%
\>[19]\AgdaSymbol{=} \AgdaSymbol{λ} \AgdaBound{f} \AgdaBound{a} \AgdaSymbol{→} \AgdaFunction{[} \AgdaBound{Bx} \AgdaBound{a} \AgdaFunction{]sym} \AgdaSymbol{(}\AgdaBound{f} \AgdaBound{a}\AgdaSymbol{)}\<%
\\
\>[0]\AgdaIndent{9}{}\<[9]%
\>[9]\AgdaSymbol{;} \AgdaField{trans} \<[19]%
\>[19]\AgdaSymbol{=} \AgdaSymbol{λ} \AgdaBound{f} \AgdaBound{g} \AgdaBound{a} \AgdaSymbol{→} \AgdaFunction{[} \AgdaBound{Bx} \AgdaBound{a} \AgdaFunction{]trans} \AgdaSymbol{(}\AgdaBound{f} \AgdaBound{a}\AgdaSymbol{)} \AgdaSymbol{(}\AgdaBound{g} \AgdaBound{a}\AgdaSymbol{)}\<%
\\
\>[0]\AgdaIndent{9}{}\<[9]%
\>[9]\AgdaSymbol{\}}\<%
\\
%
\\
\>[0]\AgdaIndent{2}{}\<[2]%
\>[2]\AgdaSymbol{;} \AgdaField{substT} \AgdaSymbol{=} \AgdaSymbol{λ} \AgdaSymbol{\{}\AgdaBound{x}\AgdaSymbol{\}} \AgdaSymbol{\{}\AgdaBound{y}\AgdaSymbol{\}} \AgdaBound{p} \<[25]%
\>[25]\<%
\\
\>[0]\AgdaIndent{15}{}\<[15]%
\>[15]\AgdaSymbol{→} \AgdaKeyword{let} \AgdaBound{y2x} \AgdaSymbol{=} \AgdaSymbol{λ} \AgdaBound{a} \AgdaSymbol{→} \AgdaFunction{[} \AgdaBound{A} \AgdaFunction{]subst} \AgdaSymbol{(}\AgdaFunction{[} \AgdaBound{Γ} \AgdaFunction{]sym} \AgdaBound{p}\AgdaSymbol{)} \AgdaBound{a} \AgdaKeyword{in}\<%
\\
\>[15]\AgdaIndent{17}{}\<[17]%
\>[17]\AgdaKeyword{let} \AgdaBound{x2y} \AgdaSymbol{=} \AgdaSymbol{λ} \AgdaBound{a} \AgdaSymbol{→} \AgdaFunction{[} \AgdaBound{A} \AgdaFunction{]subst} \AgdaBound{p} \AgdaBound{a} \AgdaKeyword{in}\<%
\\
\>[15]\AgdaIndent{17}{}\<[17]%
\>[17]\AgdaKeyword{let} \AgdaBound{p'} \AgdaSymbol{=} \AgdaSymbol{λ} \AgdaBound{a} \AgdaSymbol{→} \AgdaFunction{[} \AgdaBound{A} \AgdaFunction{]trans-refl} \AgdaKeyword{in}\<%
\\
\>[-11]\AgdaIndent{13}{}\<[13]%
\>[13]\AgdaSymbol{λ\{(}\AgdaBound{f} \AgdaInductiveConstructor{,} \AgdaBound{rsp}\AgdaSymbol{)} \AgdaSymbol{→} \<[28]%
\>[28]\<%
\\
\>[0]\AgdaIndent{15}{}\<[15]%
\>[15]\AgdaSymbol{(λ} \AgdaBound{a} \AgdaSymbol{→} \AgdaFunction{[} \AgdaBound{B} \AgdaFunction{]subst} \AgdaSymbol{(}\AgdaBound{p} \AgdaInductiveConstructor{,} \AgdaBound{p'} \AgdaBound{a}\AgdaSymbol{)} \AgdaSymbol{(}\AgdaBound{f} \AgdaSymbol{(}\AgdaBound{y2x} \AgdaBound{a}\AgdaSymbol{)))}\<%
\\
\>[0]\AgdaIndent{15}{}\<[15]%
\>[15]\AgdaInductiveConstructor{,} \<[49]%
\>[49]\<%
\\
\>[0]\AgdaIndent{15}{}\<[15]%
\>[15]\AgdaSymbol{(λ} \AgdaBound{a} \AgdaBound{b} \AgdaBound{q} \AgdaSymbol{→} \<[26]%
\>[26]\<%
\\
\>[15]\AgdaIndent{16}{}\<[16]%
\>[16]\AgdaKeyword{let} \AgdaBound{a'} \AgdaSymbol{=} \AgdaBound{y2x} \AgdaBound{a} \AgdaKeyword{in} \<[34]%
\>[34]\<%
\\
\>[15]\AgdaIndent{16}{}\<[16]%
\>[16]\AgdaKeyword{let} \AgdaBound{b'} \AgdaSymbol{=} \AgdaBound{y2x} \AgdaBound{b} \AgdaKeyword{in}\<%
\\
\>[15]\AgdaIndent{16}{}\<[16]%
\>[16]\AgdaKeyword{let} \AgdaBound{q'} \AgdaSymbol{=} \AgdaFunction{[} \AgdaBound{A} \AgdaFunction{]subst*} \AgdaSymbol{(}\AgdaFunction{[} \AgdaBound{Γ} \AgdaFunction{]sym} \AgdaBound{p}\AgdaSymbol{)} \AgdaBound{q} \AgdaKeyword{in}\<%
\\
\>[15]\AgdaIndent{16}{}\<[16]%
\>[16]\AgdaKeyword{let} \AgdaBound{H} \AgdaSymbol{=} \AgdaBound{rsp} \AgdaBound{a'} \AgdaBound{b'} \AgdaBound{q'} \AgdaKeyword{in}\<%
\\
\>[15]\AgdaIndent{16}{}\<[16]%
\>[16]\AgdaKeyword{let} \AgdaBound{r} \AgdaSymbol{:} \AgdaFunction{[} \AgdaBound{Γ} \AgdaFunction{\&} \AgdaBound{A} \AgdaFunction{]} \AgdaSymbol{(}\AgdaBound{x} \AgdaInductiveConstructor{,} \AgdaBound{b'}\AgdaSymbol{)} \AgdaFunction{≈} \AgdaSymbol{(}\AgdaBound{y} \AgdaInductiveConstructor{,} \AgdaBound{b}\AgdaSymbol{)}
                    r \AgdaSymbol{=} \AgdaSymbol{(}\AgdaBound{p} \AgdaInductiveConstructor{,} \AgdaBound{p'} \AgdaBound{b}\AgdaSymbol{)} \AgdaKeyword{in}\<%
\\
\>[15]\AgdaIndent{16}{}\<[16]%
\>[16]\AgdaKeyword{let} \AgdaBound{pre} \AgdaSymbol{=} \AgdaFunction{[} \AgdaBound{B} \AgdaFunction{]subst*} \AgdaBound{r} \AgdaBound{H} \AgdaKeyword{in}\<%
\\
\>[15]\AgdaIndent{16}{}\<[16]%
\>[16]\<%
\\
\>[15]\AgdaIndent{16}{}\<[16]%
\>[16]\AgdaFunction{[} \AgdaFunction{[} \AgdaBound{B} \AgdaFunction{]fm} \AgdaSymbol{(}\AgdaBound{y} \AgdaInductiveConstructor{,} \AgdaBound{b}\AgdaSymbol{)} \AgdaFunction{]trans} \<[41]%
\>[41]\<%
\\
\>[15]\AgdaIndent{16}{}\<[16]%
\>[16]\AgdaSymbol{(}\AgdaFunction{[} \AgdaBound{B} \AgdaFunction{]trans*} \AgdaSymbol{\_)} \<[48]%
\>[48]\<%
\\
\>[15]\AgdaIndent{16}{}\<[16]%
\>[16]\AgdaSymbol{(}\AgdaFunction{[} \AgdaFunction{[} \AgdaBound{B} \AgdaFunction{]fm} \AgdaSymbol{(}\AgdaBound{y} \AgdaInductiveConstructor{,} \AgdaBound{b}\AgdaSymbol{)} \AgdaFunction{]trans} \<[42]%
\>[42]\<%
\\
\>[15]\AgdaIndent{16}{}\<[16]%
\>[16]\AgdaFunction{[} \AgdaBound{B} \AgdaFunction{]subst-pi} \<[30]%
\>[30]\<%
\\
\>[15]\AgdaIndent{16}{}\<[16]%
\>[16]\AgdaSymbol{(}\AgdaFunction{[} \AgdaFunction{[} \AgdaBound{B} \AgdaFunction{]fm} \AgdaSymbol{(}\AgdaBound{y} \AgdaInductiveConstructor{,} \AgdaBound{b}\AgdaSymbol{)} \AgdaFunction{]trans} \<[42]%
\>[42]\<%
\\
\>[15]\AgdaIndent{16}{}\<[16]%
\>[16]\AgdaSymbol{(}\AgdaFunction{[} \AgdaFunction{[} \AgdaBound{B} \AgdaFunction{]fm} \AgdaSymbol{(}\AgdaBound{y} \AgdaInductiveConstructor{,} \AgdaBound{b}\AgdaSymbol{)} \AgdaFunction{]sym} \AgdaSymbol{(}\AgdaFunction{[} \AgdaBound{B} \AgdaFunction{]trans*} \AgdaSymbol{\_))} \<[57]%
\>[57]\<%
\\
\>[15]\AgdaIndent{16}{}\<[16]%
\>[16]\AgdaBound{pre}\AgdaSymbol{))} \<[23]%
\>[23]\<%
\\
\>[15]\AgdaIndent{16}{}\<[16]%
\>[16]\AgdaSymbol{)} \<[22]%
\>[22]\<%
\\
\>[-12]\AgdaIndent{13}{}\<[13]%
\>[13]\AgdaSymbol{\}}\<%
\\
\>[0]\AgdaIndent{2}{}\<[2]%
\>[2]\AgdaSymbol{;} \AgdaField{subst*} \AgdaSymbol{=} \AgdaSymbol{λ} \AgdaBound{\_} \AgdaBound{q} \AgdaBound{\_} \AgdaSymbol{→} \AgdaFunction{[} \AgdaBound{B} \AgdaFunction{]subst*} \AgdaSymbol{\_} \AgdaSymbol{(}\AgdaBound{q} \AgdaSymbol{\_)}\<%
\\
\>[0]\AgdaIndent{2}{}\<[2]%
\>[2]\AgdaSymbol{;} \AgdaField{refl*} \AgdaSymbol{=} \AgdaSymbol{λ} \AgdaSymbol{\{}\AgdaBound{x} \AgdaSymbol{(}\AgdaBound{f} \AgdaInductiveConstructor{,} \AgdaBound{rsp}\AgdaSymbol{)} \AgdaBound{a} \<[29]%
\>[29]\<%
\\
\>[2]\AgdaIndent{12}{}\<[12]%
\>[12]\AgdaSymbol{→} \AgdaFunction{[} \AgdaFunction{[} \AgdaBound{B} \AgdaFunction{]fm} \AgdaSymbol{\_} \AgdaFunction{]trans} \AgdaFunction{[} \AgdaBound{B} \AgdaFunction{]subst-pi} \<[47]%
\>[47]\<%
\\
\>[2]\AgdaIndent{12}{}\<[12]%
\>[12]\AgdaSymbol{(}\AgdaBound{rsp} \AgdaSymbol{(}\AgdaFunction{[} \AgdaBound{A} \AgdaFunction{]subst} \AgdaSymbol{(}\AgdaFunction{[} \AgdaBound{Γ} \AgdaFunction{]sym} \AgdaFunction{[} \AgdaBound{Γ} \AgdaFunction{]refl}\AgdaSymbol{)} \AgdaBound{a}\AgdaSymbol{)} \AgdaBound{a} \AgdaFunction{[} \AgdaBound{A} \AgdaFunction{]subst-pi'}\AgdaSymbol{)} \<[72]%
\>[72]\AgdaSymbol{\}}\<%
\\
\>[0]\AgdaIndent{2}{}\<[2]%
\>[2]\AgdaSymbol{;} \AgdaField{trans*} \AgdaSymbol{=} \AgdaSymbol{λ} \AgdaSymbol{\{(}\AgdaBound{f} \AgdaInductiveConstructor{,} \AgdaBound{rsp}\AgdaSymbol{)} \AgdaBound{a} \AgdaSymbol{→}\<%
\\
\>[0]\AgdaIndent{13}{}\<[13]%
\>[13]\AgdaFunction{[} \AgdaFunction{[} \AgdaBound{B} \AgdaFunction{]fm} \AgdaSymbol{\_} \AgdaFunction{]trans} \<[32]%
\>[32]\<%
\\
\>[0]\AgdaIndent{13}{}\<[13]%
\>[13]\AgdaSymbol{(}\AgdaFunction{[} \AgdaFunction{[} \AgdaBound{B} \AgdaFunction{]fm} \AgdaSymbol{\_} \AgdaFunction{]trans} \<[33]%
\>[33]\<%
\\
\>[0]\AgdaIndent{13}{}\<[13]%
\>[13]\AgdaSymbol{(}\AgdaFunction{[} \AgdaBound{B} \AgdaFunction{]trans*} \AgdaSymbol{\_)} \<[29]%
\>[29]\<%
\\
\>[0]\AgdaIndent{13}{}\<[13]%
\>[13]\AgdaSymbol{(}\AgdaFunction{[} \AgdaFunction{[} \AgdaBound{B} \AgdaFunction{]fm} \AgdaSymbol{\_} \AgdaFunction{]sym} \AgdaSymbol{(}\AgdaFunction{[} \AgdaFunction{[} \AgdaBound{B} \AgdaFunction{]fm} \AgdaSymbol{\_} \AgdaFunction{]trans} \<[51]%
\>[51]\<%
\\
\>[0]\AgdaIndent{13}{}\<[13]%
\>[13]\AgdaSymbol{(}\AgdaFunction{[} \AgdaBound{B} \AgdaFunction{]trans*} \AgdaSymbol{\_)} \AgdaFunction{[} \AgdaBound{B} \AgdaFunction{]subst-pi}\AgdaSymbol{)))} \<[46]%
\>[46]\<%
\\
\>[0]\AgdaIndent{13}{}\<[13]%
\>[13]\AgdaSymbol{(}\AgdaFunction{[} \AgdaBound{B} \AgdaFunction{]subst*} \AgdaSymbol{\_} \AgdaSymbol{(}\AgdaBound{rsp} \AgdaSymbol{\_} \AgdaSymbol{\_} \AgdaSymbol{(}\AgdaFunction{[} \AgdaFunction{[} \AgdaBound{A} \AgdaFunction{]fm} \AgdaSymbol{\_} \AgdaFunction{]trans} \<[57]%
\>[57]\<%
\\
\>[0]\AgdaIndent{13}{}\<[13]%
\>[13]\AgdaSymbol{(}\AgdaFunction{[} \AgdaBound{A} \AgdaFunction{]trans*} \AgdaSymbol{\_)} \AgdaFunction{[} \AgdaBound{A} \AgdaFunction{]subst-pi}\AgdaSymbol{)))} \AgdaSymbol{\}} \<[48]%
\>[48]\<%
\\
\>[0]\AgdaIndent{2}{}\<[2]%
\>[2]\AgdaSymbol{\}}\<%
\\
\end{code}

The type construction is not complete because it also includes lambda abstraction and function application which correspond to introduction and elimination rules respectively. The verification of $\beta$ and $\eta$ laws are also essential, but we can not implement it in Agda due to technical issue. Hopeful the problem can be fixed in later patches.

We also construct some basic types including natural numbers simply typed universe, equality types. $\Sigma$-types and quotient types are not discussed in the Altenkirch's original paper, which we cover here. $\Sigma$-types are intuitively simple, but the construction of it in setoid models still requires a lot of unification proofs due to the type dependency. The more important one is the construction of quotient types.


\subsection{Quotient types}

To formulate quotient types, we have to provide the encoding of equivalence in setoid model. First of all, the propositions are given by the metatheoretial propositional universe with logical equivalence:

\begin{code}
%
\\
\>\AgdaFunction{Pu} \AgdaSymbol{:} \AgdaRecord{HSetoid}\<%
\\
\>\AgdaFunction{Pu} \AgdaSymbol{=} \AgdaKeyword{record}\<%
\\
\>[0]\AgdaIndent{6}{}\<[6]%
\>[6]\AgdaSymbol{\{} \AgdaField{Carrier} \AgdaSymbol{=} \AgdaRecord{HProp}\<%
\\
\>[0]\AgdaIndent{6}{}\<[6]%
\>[6]\AgdaSymbol{;} \AgdaField{\_≈h\_} \<[16]%
\>[16]\AgdaSymbol{=} \AgdaFunction{\_⇄\_}\<%
\\
\>[0]\AgdaIndent{6}{}\<[6]%
\>[6]\AgdaSymbol{;} \AgdaField{refl} \<[16]%
\>[16]\AgdaSymbol{=} \AgdaFunction{id} \AgdaInductiveConstructor{,} \AgdaFunction{id} \<[26]%
\>[26]\<%
\\
\>[0]\AgdaIndent{6}{}\<[6]%
\>[6]\AgdaSymbol{;} \AgdaField{sym} \<[16]%
\>[16]\AgdaSymbol{=} \AgdaSymbol{λ} \AgdaSymbol{\{(}\AgdaBound{a} \AgdaInductiveConstructor{,} \AgdaBound{b}\AgdaSymbol{)} \AgdaSymbol{→(}\AgdaBound{b} \AgdaInductiveConstructor{,} \AgdaBound{a}\AgdaSymbol{)\}}\<%
\\
\>[0]\AgdaIndent{6}{}\<[6]%
\>[6]\AgdaSymbol{;} \AgdaField{trans} \<[16]%
\>[16]\AgdaSymbol{=} \AgdaSymbol{λ} \AgdaSymbol{\{(}\AgdaBound{a} \AgdaInductiveConstructor{,} \AgdaBound{b}\AgdaSymbol{)} \AgdaSymbol{(}\AgdaBound{a'} \AgdaInductiveConstructor{,} \AgdaBound{b'}\AgdaSymbol{)} \<[39]%
\>[39]\<%
\\
\>[6]\AgdaIndent{18}{}\<[18]%
\>[18]\AgdaSymbol{→} \AgdaSymbol{(}\AgdaBound{a'} \AgdaFunction{∘} \AgdaBound{a}\AgdaSymbol{)} \AgdaInductiveConstructor{,} \AgdaSymbol{(}\AgdaBound{b} \AgdaFunction{∘} \AgdaBound{b'}\AgdaSymbol{)\}}\<%
\\
\>[0]\AgdaIndent{6}{}\<[6]%
\>[6]\AgdaSymbol{\}}\<%
\\
%
\\
%
\\
\>\AgdaFunction{⟦Prop⟧} \AgdaSymbol{:} \AgdaRecord{Ty} \AgdaFunction{●}\<%
\\
\>\AgdaFunction{⟦Prop⟧} \AgdaSymbol{=} \AgdaKeyword{record} \AgdaSymbol{\{} \AgdaField{fm} \AgdaSymbol{=} \AgdaSymbol{λ} \AgdaBound{x} \AgdaSymbol{→} \AgdaFunction{Pu}\AgdaSymbol{;} \AgdaField{substT} \AgdaSymbol{=} \AgdaSymbol{λ} \AgdaBound{x'} \AgdaBound{x0} \AgdaSymbol{→} \AgdaBound{x0}\AgdaSymbol{;}\<%
\\
\>[0]\AgdaIndent{17}{}\<[17]%
\>[17]\AgdaField{subst*} \AgdaSymbol{=} \AgdaSymbol{λ} \AgdaBound{p} \AgdaBound{x'} \AgdaSymbol{→} \AgdaBound{x'}\AgdaSymbol{;} \AgdaField{refl*} \AgdaSymbol{=} \AgdaSymbol{λ} \AgdaBound{x} \AgdaBound{a} \AgdaSymbol{→} \AgdaFunction{id} \AgdaInductiveConstructor{,}\<%
\\
\>[17]\AgdaIndent{27}{}\<[27]%
\>[27]\AgdaFunction{id}\AgdaSymbol{;} \AgdaField{trans*} \AgdaSymbol{=} \AgdaSymbol{λ} \AgdaBound{\_} \AgdaSymbol{→} \AgdaFunction{id} \AgdaInductiveConstructor{,} \AgdaFunction{id} \AgdaSymbol{\}}\<%
\\
\end{code}

And then a propositional binary relation can be constructed as:

\begin{code}
\\
\>\AgdaFunction{PropRel} \<[10]%
\>[10]\AgdaSymbol{:} \AgdaSymbol{\{}\AgdaBound{Γ} \AgdaSymbol{:} \AgdaFunction{Con}\AgdaSymbol{\}(}\AgdaBound{A} \AgdaSymbol{:} \AgdaRecord{Ty} \AgdaBound{Γ}\AgdaSymbol{)} \AgdaSymbol{→} \AgdaRecord{Ty} \AgdaBound{Γ}\<%
\\
\>\AgdaFunction{PropRel} \AgdaBound{A} \AgdaSymbol{=} \AgdaBound{A} \AgdaFunction{⇒} \AgdaBound{A} \AgdaFunction{⇒} \AgdaFunction{⟦Prop⟧} \AgdaFunction{[} \AgdaInductiveConstructor{fn:} \AgdaSymbol{(λ} \AgdaBound{x} \AgdaSymbol{→} \AgdaInductiveConstructor{tt}\AgdaSymbol{)} \AgdaInductiveConstructor{resp:} \AgdaSymbol{(λ} \AgdaBound{x'} \AgdaSymbol{→} \AgdaInductiveConstructor{tt}\AgdaSymbol{)} \AgdaFunction{]T}\<%
\\
\end{code}

However practically, considering the efficiency problem of handling complicated construction in Agda, we choose to use the underlying sets straightforwardly with necessary proving of properties.


We first try to construct the predicates of defining an equivalence relation from a propositional relation, which is not as simple as we expected. We refer to Hofmann's approach in \cite{hof:95:sm}, namely by defing reflexive, symmetric and transitive closure of a propositional relation. The same concern also makes us choose the underlying set definition over the propositional construction.

\begin{code}
%
\\
\>\AgdaFunction{Closure} \AgdaSymbol{:} \AgdaSymbol{(}\AgdaBound{A} \AgdaSymbol{:} \AgdaPrimitiveType{Set}\AgdaSymbol{)(}\AgdaBound{\_∼\_} \AgdaBound{R} \AgdaSymbol{:} \AgdaBound{A} \AgdaSymbol{→} \AgdaBound{A} \AgdaSymbol{→} \AgdaRecord{HProp}\AgdaSymbol{)} \<[43]%
\>[43]\<%
\\
\>[0]\AgdaIndent{9}{}\<[9]%
\>[9]\AgdaSymbol{→} \AgdaBound{A} \AgdaSymbol{→} \AgdaBound{A} \AgdaSymbol{→} \AgdaRecord{HProp}\<%
\\
\>\AgdaFunction{Closure} \AgdaBound{A} \AgdaBound{\_∼\_} \AgdaBound{R} \AgdaBound{s} \AgdaBound{s'}\<%
\\
\>[9]\AgdaIndent{10}{}\<[10]%
\>[10]\AgdaSymbol{=} \AgdaKeyword{record}\<%
\\
\>[10]\AgdaIndent{12}{}\<[12]%
\>[12]\AgdaSymbol{\{} \AgdaField{prf} \AgdaSymbol{=} \AgdaSymbol{∀} \AgdaSymbol{(}\AgdaBound{R'} \AgdaSymbol{:} \AgdaBound{A} \AgdaSymbol{→} \AgdaBound{A} \AgdaSymbol{→} \AgdaRecord{HProp}\AgdaSymbol{)} \<[43]%
\>[43]\<%
\\
\>[12]\AgdaIndent{18}{}\<[18]%
\>[18]\AgdaSymbol{→} \AgdaSymbol{(∀} \AgdaBound{x} \AgdaBound{y} \AgdaSymbol{→} \AgdaFunction{<} \AgdaBound{R'} \AgdaBound{x} \AgdaBound{y} \AgdaFunction{>} \AgdaSymbol{→} \AgdaFunction{<} \AgdaBound{R'} \AgdaBound{y} \AgdaBound{x} \AgdaFunction{>}\AgdaSymbol{)}\<%
\\
\>[12]\AgdaIndent{18}{}\<[18]%
\>[18]\AgdaSymbol{→} \AgdaSymbol{(∀} \AgdaBound{x} \AgdaBound{y} \AgdaBound{z} \AgdaSymbol{→} \AgdaFunction{<} \AgdaBound{R'} \AgdaBound{x} \AgdaBound{y} \AgdaFunction{>} \AgdaSymbol{→} \AgdaFunction{<} \AgdaBound{R'} \AgdaBound{y} \AgdaBound{z} \AgdaFunction{>} \AgdaSymbol{→} \AgdaFunction{<} \AgdaBound{R'} \AgdaBound{x} \AgdaBound{z} \AgdaFunction{>}\AgdaSymbol{)}\<%
\\
\>[12]\AgdaIndent{18}{}\<[18]%
\>[18]\AgdaSymbol{→} \AgdaSymbol{(∀} \AgdaBound{x} \AgdaBound{y} \AgdaSymbol{→} \AgdaFunction{<} \AgdaBound{x} \AgdaBound{∼} \AgdaBound{y} \AgdaFunction{>} \AgdaSymbol{→} \AgdaFunction{<} \AgdaBound{R'} \AgdaBound{x} \AgdaBound{y} \AgdaFunction{>}\AgdaSymbol{)}\<%
\\
\>[12]\AgdaIndent{18}{}\<[18]%
\>[18]\AgdaSymbol{→} \AgdaSymbol{(∀} \AgdaBound{x} \AgdaBound{y} \AgdaSymbol{→} \AgdaFunction{<} \AgdaBound{R} \AgdaBound{x} \AgdaBound{y} \AgdaFunction{>} \AgdaSymbol{→} \AgdaFunction{<} \AgdaBound{R'} \AgdaBound{x} \AgdaBound{y} \AgdaFunction{>}\AgdaSymbol{)}\<%
\\
\>[12]\AgdaIndent{18}{}\<[18]%
\>[18]\AgdaSymbol{→} \AgdaFunction{<} \AgdaBound{R'} \AgdaBound{s} \AgdaBound{s'} \AgdaFunction{>}\<%
\\
\>[0]\AgdaIndent{12}{}\<[12]%
\>[12]\AgdaSymbol{;} \AgdaField{Uni} \AgdaSymbol{=} \AgdaBound{ext} \AgdaSymbol{(λ} \AgdaBound{R'} \AgdaSymbol{→} \AgdaBound{ext} \AgdaSymbol{(λ} \AgdaBound{x₁} \AgdaSymbol{→} \AgdaBound{ext} \<[48]%
\>[48]\<%
\\
\>[0]\AgdaIndent{20}{}\<[20]%
\>[20]\AgdaSymbol{(λ} \AgdaBound{x₂} \AgdaSymbol{→} \AgdaBound{ext} \AgdaSymbol{(λ} \AgdaBound{x₃} \AgdaSymbol{→} \AgdaBound{ext} \AgdaSymbol{(λ} \AgdaBound{x₄} \AgdaSymbol{→} \AgdaFunction{Uni} \AgdaSymbol{(}\AgdaBound{R'} \AgdaBound{s} \AgdaBound{s'}\AgdaSymbol{))))))}\<%
\\
\>[0]\AgdaIndent{12}{}\<[12]%
\>[12]\AgdaSymbol{\}}\<%
\\
\end{code}


We postulate the properties of the repsectiveness of equivalence which has to be held by the equivalence relation $\mathsf{R}$. Another property postulated is that a function $f$ which repsects $\mathsf{R}$, it also respects the closure of $\mathsf{R}$.

\begin{code}
%
\\
\>\AgdaKeyword{module} \AgdaModule{Q} \AgdaSymbol{(}\AgdaBound{Γ} \AgdaSymbol{:} \AgdaFunction{Con}\AgdaSymbol{)(}\AgdaBound{A} \AgdaSymbol{:} \AgdaRecord{Ty} \AgdaBound{Γ}\AgdaSymbol{)}\<%
\\
\>[0]\AgdaIndent{9}{}\<[9]%
\>[9]\AgdaSymbol{(}\AgdaBound{R} \AgdaSymbol{:} \AgdaSymbol{(}\AgdaBound{γ} \AgdaSymbol{:} \AgdaFunction{∣} \AgdaBound{Γ} \AgdaFunction{∣}\AgdaSymbol{)} \AgdaSymbol{→} \AgdaFunction{∣} \AgdaFunction{[} \AgdaBound{A} \AgdaFunction{]fm} \AgdaBound{γ} \AgdaFunction{∣} \AgdaSymbol{→} \AgdaFunction{∣} \AgdaFunction{[} \AgdaBound{A} \AgdaFunction{]fm} \AgdaBound{γ} \AgdaFunction{∣} \AgdaSymbol{→} \AgdaRecord{HProp}\AgdaSymbol{)}\<%
\\
\>[0]\AgdaIndent{9}{}\<[9]%
\>[9]\AgdaSymbol{(}\AgdaBound{Rresp} \AgdaSymbol{:} \<[19]%
\>[19]\AgdaSymbol{∀} \AgdaSymbol{(}\AgdaBound{γ} \AgdaSymbol{:} \AgdaFunction{∣} \AgdaBound{Γ} \AgdaFunction{∣}\AgdaSymbol{)(}\AgdaBound{x} \AgdaBound{y} \AgdaSymbol{:} \AgdaFunction{∣} \AgdaFunction{[} \AgdaBound{A} \AgdaFunction{]fm} \AgdaBound{γ} \AgdaFunction{∣}\AgdaSymbol{)(}\AgdaBound{γ'} \AgdaSymbol{:} \AgdaFunction{∣} \AgdaBound{Γ} \AgdaFunction{∣}\AgdaSymbol{)(}\AgdaBound{x'} \AgdaBound{y'} \AgdaSymbol{:} \AgdaFunction{∣} \AgdaFunction{[} \AgdaBound{A} \AgdaFunction{]fm} \AgdaBound{γ} \AgdaFunction{∣}\AgdaSymbol{)}\<%
\\
\>[9]\AgdaIndent{11}{}\<[11]%
\>[11]\AgdaSymbol{(}\AgdaBound{eq1} \AgdaSymbol{:} \AgdaFunction{<} \AgdaFunction{rel} \AgdaBound{Γ} \AgdaBound{γ} \AgdaBound{γ'} \AgdaFunction{>}\AgdaSymbol{)}\<%
\\
\>[9]\AgdaIndent{11}{}\<[11]%
\>[11]\AgdaSymbol{(}\AgdaBound{eq2} \AgdaSymbol{:} \AgdaFunction{<} \AgdaFunction{rel} \AgdaSymbol{(}\AgdaFunction{[} \AgdaBound{A} \AgdaFunction{]fm} \AgdaBound{γ}\AgdaSymbol{)} \AgdaBound{x} \AgdaBound{x'} \AgdaFunction{>}\AgdaSymbol{)}\<%
\\
\>[9]\AgdaIndent{11}{}\<[11]%
\>[11]\AgdaSymbol{(}\AgdaBound{eq3} \AgdaSymbol{:} \AgdaFunction{<} \AgdaFunction{rel} \AgdaSymbol{(}\AgdaFunction{[} \AgdaBound{A} \AgdaFunction{]fm} \AgdaBound{γ}\AgdaSymbol{)} \AgdaBound{y} \AgdaBound{y'} \AgdaFunction{>}\AgdaSymbol{)} \AgdaSymbol{→} \AgdaFunction{<} \AgdaBound{R} \AgdaBound{γ} \AgdaBound{x} \AgdaBound{y} \AgdaFunction{⇄} \AgdaBound{R} \AgdaBound{γ} \AgdaBound{x'} \AgdaBound{y'} \AgdaFunction{>}\AgdaSymbol{)}\<%
\\
\>[0]\AgdaIndent{9}{}\<[9]%
\>[9]\AgdaSymbol{(}\AgdaBound{Rresp2} \AgdaSymbol{:} \<[20]%
\>[20]\AgdaSymbol{∀} \AgdaSymbol{(}\AgdaBound{γ} \AgdaSymbol{:} \AgdaFunction{∣} \AgdaBound{Γ} \AgdaFunction{∣}\AgdaSymbol{)}\<%
\\
\>[0]\AgdaIndent{9}{}\<[9]%
\>[9]\AgdaSymbol{(}\AgdaBound{x} \AgdaBound{y} \AgdaSymbol{:} \AgdaFunction{∣} \AgdaFunction{[} \AgdaBound{A} \AgdaFunction{]fm} \AgdaBound{γ} \AgdaFunction{∣}\AgdaSymbol{)(}\AgdaBound{γ'} \AgdaSymbol{:} \AgdaFunction{∣} \AgdaBound{Γ} \AgdaFunction{∣}\AgdaSymbol{)(}\AgdaBound{x'} \AgdaBound{y'} \AgdaSymbol{:} \AgdaFunction{∣} \AgdaFunction{[} \AgdaBound{A} \AgdaFunction{]fm} \AgdaBound{γ'} \AgdaFunction{∣}\AgdaSymbol{)}\<%
\\
\>[9]\AgdaIndent{11}{}\<[11]%
\>[11]\AgdaSymbol{(}\AgdaBound{eq1} \AgdaSymbol{:} \AgdaFunction{<} \AgdaFunction{rel} \AgdaBound{Γ} \AgdaBound{γ} \AgdaBound{γ'} \AgdaFunction{>}\AgdaSymbol{)}\<%
\\
\>[9]\AgdaIndent{11}{}\<[11]%
\>[11]\AgdaSymbol{(}\AgdaBound{eq2} \AgdaSymbol{:} \AgdaFunction{<} \AgdaFunction{rel} \AgdaSymbol{(}\AgdaFunction{[} \AgdaBound{A} \AgdaFunction{]fm} \AgdaBound{γ'}\AgdaSymbol{)} \AgdaSymbol{(}\AgdaFunction{[} \AgdaBound{A} \AgdaFunction{]subst} \AgdaBound{eq1} \AgdaBound{x}\AgdaSymbol{)} \AgdaBound{x'} \AgdaFunction{>}\AgdaSymbol{)}\<%
\\
\>[9]\AgdaIndent{11}{}\<[11]%
\>[11]\AgdaSymbol{(}\AgdaBound{eq3} \AgdaSymbol{:} \AgdaFunction{<} \AgdaFunction{rel} \AgdaSymbol{(}\AgdaFunction{[} \AgdaBound{A} \AgdaFunction{]fm} \AgdaBound{γ'}\AgdaSymbol{)} \AgdaSymbol{(}\AgdaFunction{[} \AgdaBound{A} \AgdaFunction{]subst} \AgdaBound{eq1} \AgdaBound{y}\AgdaSymbol{)} \AgdaBound{y'} \AgdaFunction{>}\AgdaSymbol{)} \<[62]%
\>[62]\<%
\\
\>[9]\AgdaIndent{11}{}\<[11]%
\>[11]\AgdaSymbol{→} \AgdaFunction{<} \AgdaBound{R} \AgdaBound{γ} \AgdaBound{x} \AgdaBound{y} \AgdaFunction{⇄} \AgdaBound{R} \AgdaBound{γ'} \AgdaBound{x'} \AgdaBound{y'} \AgdaFunction{>}\AgdaSymbol{)}\<%
\\
\>[0]\AgdaIndent{9}{}\<[9]%
\>[9]\AgdaKeyword{where}\<%
\\
%
\\
\>[0]\AgdaIndent{2}{}\<[2]%
\>[2]\AgdaFunction{RC} \AgdaSymbol{:} \AgdaSymbol{(}\AgdaBound{γ} \AgdaSymbol{:} \AgdaFunction{∣} \AgdaBound{Γ} \AgdaFunction{∣}\AgdaSymbol{)} \AgdaSymbol{→} \AgdaFunction{∣} \AgdaFunction{[} \AgdaBound{A} \AgdaFunction{]fm} \AgdaBound{γ} \AgdaFunction{∣} \AgdaSymbol{→} \AgdaFunction{∣} \AgdaFunction{[} \AgdaBound{A} \AgdaFunction{]fm} \AgdaBound{γ} \AgdaFunction{∣} \AgdaSymbol{→} \AgdaRecord{HProp}\<%
\\
\>[0]\AgdaIndent{2}{}\<[2]%
\>[2]\AgdaFunction{RC} \AgdaBound{γ} \AgdaSymbol{=} \AgdaFunction{Closure} \AgdaSymbol{\_} \AgdaSymbol{(}\AgdaFunction{rel} \AgdaSymbol{(}\AgdaFunction{[} \AgdaBound{A} \AgdaFunction{]fm} \AgdaBound{γ}\AgdaSymbol{))} \AgdaSymbol{(}\AgdaBound{R} \AgdaBound{γ}\AgdaSymbol{)}\<%
\\
%
\\
\>[0]\AgdaIndent{2}{}\<[2]%
\>[2]\AgdaFunction{inc} \AgdaSymbol{:} \AgdaSymbol{\{}\AgdaBound{γ} \AgdaSymbol{:} \AgdaFunction{∣} \AgdaBound{Γ} \AgdaFunction{∣}\AgdaSymbol{\}\{}\AgdaBound{a} \AgdaBound{b} \AgdaSymbol{:} \<[27]%
\>[27]\AgdaFunction{∣} \AgdaFunction{[} \AgdaBound{A} \AgdaFunction{]fm} \AgdaBound{γ} \AgdaFunction{∣}\AgdaSymbol{\}} \AgdaSymbol{→} \AgdaFunction{<} \AgdaBound{R} \AgdaBound{γ} \AgdaBound{a} \AgdaBound{b} \AgdaFunction{>} \AgdaSymbol{→} \AgdaFunction{<} \AgdaFunction{RC} \AgdaBound{γ} \AgdaBound{a} \AgdaBound{b} \AgdaFunction{>}\<%
\\
\>[0]\AgdaIndent{2}{}\<[2]%
\>[2]\AgdaFunction{inc} \AgdaBound{r} \AgdaSymbol{=} \AgdaSymbol{λ} \AgdaBound{R'} \AgdaBound{x} \AgdaBound{x₁} \AgdaBound{x₂} \AgdaBound{x₃} \AgdaSymbol{→} \AgdaBound{x₃} \AgdaSymbol{\_} \AgdaSymbol{\_} \AgdaBound{r}\<%
\\
\>[0]\AgdaIndent{2}{}\<[2]%
\>[2]\<%
\\
\>[0]\AgdaIndent{2}{}\<[2]%
\>[2]\AgdaKeyword{postulate} \AgdaPostulate{respClosure} \AgdaSymbol{:} \AgdaSymbol{∀\{}\AgdaBound{γ} \AgdaBound{γ'} \AgdaSymbol{:} \AgdaFunction{∣} \AgdaBound{Γ} \AgdaFunction{∣}\AgdaSymbol{\}\{}\AgdaBound{a} \AgdaBound{b} \AgdaSymbol{:} \<[49]%
\>[49]\AgdaFunction{∣} \AgdaFunction{[} \AgdaBound{A} \AgdaFunction{]fm} \AgdaBound{γ} \AgdaFunction{∣}\AgdaSymbol{\}}\<%
\\
\>[2]\AgdaIndent{26}{}\<[26]%
\>[26]\AgdaSymbol{\{}\AgdaBound{f} \AgdaSymbol{:} \AgdaFunction{∣} \AgdaFunction{[} \AgdaBound{A} \AgdaFunction{]fm} \AgdaBound{γ} \AgdaFunction{∣} \AgdaSymbol{→} \AgdaFunction{∣} \AgdaFunction{[} \AgdaBound{A} \AgdaFunction{]fm} \AgdaBound{γ'} \AgdaFunction{∣}\AgdaSymbol{\}} \<[64]%
\>[64]\<%
\\
\>[0]\AgdaIndent{24}{}\<[24]%
\>[24]\AgdaSymbol{→} \AgdaSymbol{(}\AgdaFunction{<} \AgdaBound{R} \AgdaBound{γ} \AgdaBound{a} \AgdaBound{b} \AgdaFunction{>} \AgdaSymbol{→} \AgdaFunction{<} \AgdaBound{R} \AgdaBound{γ'} \AgdaSymbol{(}\AgdaBound{f} \AgdaBound{a}\AgdaSymbol{)} \AgdaSymbol{(}\AgdaBound{f} \AgdaBound{b}\AgdaSymbol{)} \AgdaFunction{>}\AgdaSymbol{)} \<[63]%
\>[63]\<%
\\
\>[0]\AgdaIndent{24}{}\<[24]%
\>[24]\AgdaSymbol{→} \AgdaSymbol{(}\AgdaFunction{<} \AgdaFunction{RC} \AgdaBound{γ} \AgdaBound{a} \AgdaBound{b} \AgdaFunction{>} \AgdaSymbol{→} \AgdaFunction{<} \AgdaFunction{RC} \AgdaBound{γ'} \AgdaSymbol{(}\AgdaBound{f} \AgdaBound{a}\AgdaSymbol{)} \AgdaSymbol{(}\AgdaBound{f} \AgdaBound{b}\AgdaSymbol{)} \AgdaFunction{>}\AgdaSymbol{)}\<%
\\
%
\\
%
\\
\>[0]\AgdaIndent{2}{}\<[2]%
\>[2]\AgdaFunction{Qsubst*} \AgdaSymbol{:} \AgdaSymbol{(}\AgdaBound{γ} \AgdaBound{γ'} \AgdaSymbol{:} \AgdaFunction{∣} \AgdaBound{Γ} \AgdaFunction{∣}\AgdaSymbol{)(}\AgdaBound{p} \AgdaSymbol{:} \AgdaFunction{[} \AgdaBound{Γ} \AgdaFunction{]} \AgdaBound{γ} \AgdaFunction{≈} \AgdaBound{γ'}\AgdaSymbol{)} \<[45]%
\>[45]\<%
\\
\>[0]\AgdaIndent{10}{}\<[10]%
\>[10]\AgdaSymbol{→} \AgdaSymbol{(}\AgdaBound{a} \AgdaBound{b} \AgdaSymbol{:} \<[20]%
\>[20]\AgdaFunction{∣} \AgdaFunction{[} \AgdaBound{A} \AgdaFunction{]fm} \AgdaBound{γ} \AgdaFunction{∣}\AgdaSymbol{)} \<[35]%
\>[35]\<%
\\
\>[0]\AgdaIndent{10}{}\<[10]%
\>[10]\AgdaSymbol{→} \AgdaFunction{<} \AgdaFunction{RC} \AgdaBound{γ} \AgdaBound{a} \AgdaBound{b} \AgdaFunction{>} \<[25]%
\>[25]\<%
\\
\>[0]\AgdaIndent{10}{}\<[10]%
\>[10]\AgdaSymbol{→} \AgdaFunction{<} \AgdaFunction{RC} \AgdaBound{γ'} \AgdaSymbol{(}\AgdaFunction{[} \AgdaBound{A} \AgdaFunction{]subst} \AgdaBound{p} \AgdaBound{a}\AgdaSymbol{)} \AgdaSymbol{(}\AgdaFunction{[} \AgdaBound{A} \AgdaFunction{]subst} \AgdaBound{p} \AgdaBound{b}\AgdaSymbol{)} \AgdaFunction{>}\<%
\\
\>[0]\AgdaIndent{2}{}\<[2]%
\>[2]\AgdaFunction{Qsubst*} \AgdaBound{γ} \AgdaBound{γ'} \AgdaBound{p} \AgdaBound{a} \AgdaBound{b} \AgdaSymbol{=} \AgdaPostulate{respClosure} \AgdaSymbol{(}\AgdaFunction{proj₁} \AgdaSymbol{(}\AgdaBound{Rresp2} \AgdaSymbol{\_} \AgdaBound{a} \AgdaBound{b} \AgdaSymbol{\_} \AgdaSymbol{\_} \AgdaSymbol{\_} \AgdaBound{p} \<[64]%
\>[64]\<%
\\
\>[0]\AgdaIndent{23}{}\<[23]%
\>[23]\AgdaFunction{[} \AgdaFunction{[} \AgdaBound{A} \AgdaFunction{]fm} \AgdaBound{γ'} \AgdaFunction{]refl} \AgdaFunction{[} \AgdaFunction{[} \AgdaBound{A} \AgdaFunction{]fm} \AgdaBound{γ'} \AgdaFunction{]refl}\AgdaSymbol{))} \<[63]%
\>[63]\<%
\\
%
\\
\>[0]\AgdaIndent{2}{}\<[2]%
\>[2]\AgdaFunction{⟦Q⟧₀} \AgdaSymbol{:} \AgdaFunction{∣} \AgdaBound{Γ} \AgdaFunction{∣} \AgdaSymbol{→} \AgdaRecord{HSetoid}\<%
\\
\>[0]\AgdaIndent{2}{}\<[2]%
\>[2]\AgdaFunction{⟦Q⟧₀} \AgdaBound{γ} \AgdaSymbol{=}\<%
\\
\>[2]\AgdaIndent{9}{}\<[9]%
\>[9]\AgdaKeyword{record}\<%
\\
\>[2]\AgdaIndent{9}{}\<[9]%
\>[9]\AgdaSymbol{\{} \AgdaField{Carrier} \AgdaSymbol{=} \AgdaFunction{∣} \AgdaFunction{[} \AgdaBound{A} \AgdaFunction{]fm} \AgdaBound{γ} \AgdaFunction{∣}\<%
\\
\>[2]\AgdaIndent{9}{}\<[9]%
\>[9]\AgdaSymbol{;} \AgdaField{\_≈h\_} \AgdaSymbol{=} \AgdaFunction{RC} \AgdaBound{γ}\<%
\\
\>[2]\AgdaIndent{9}{}\<[9]%
\>[9]\AgdaSymbol{;} \AgdaField{refl} \AgdaSymbol{=} \AgdaSymbol{λ} \AgdaSymbol{\{}\AgdaBound{x}\AgdaSymbol{\}} \AgdaBound{R'} \AgdaBound{x₁} \AgdaBound{x₂} \AgdaBound{x₃} \AgdaBound{x₄} \AgdaSymbol{→} \AgdaBound{x₃} \AgdaBound{x} \AgdaBound{x} \AgdaFunction{[} \AgdaFunction{[} \AgdaBound{A} \AgdaFunction{]fm} \AgdaBound{γ} \AgdaFunction{]refl}\<%
\\
\>[2]\AgdaIndent{9}{}\<[9]%
\>[9]\AgdaSymbol{;} \AgdaField{sym} \AgdaSymbol{=} \AgdaSymbol{λ} \AgdaSymbol{\{}\AgdaBound{x}\AgdaSymbol{\}} \AgdaSymbol{\{}\AgdaBound{y}\AgdaSymbol{\}} \AgdaBound{x≈y} \AgdaBound{R'} \AgdaBound{x₁} \AgdaBound{x₂} \AgdaBound{x₃} \AgdaBound{x₄} \AgdaSymbol{→} \AgdaBound{x₁} \AgdaBound{x} \AgdaBound{y} \AgdaSymbol{(}\AgdaBound{x≈y} \AgdaBound{R'} \AgdaBound{x₁} \AgdaBound{x₂} \AgdaBound{x₃} \AgdaBound{x₄}\AgdaSymbol{)}\<%
\\
\>[2]\AgdaIndent{9}{}\<[9]%
\>[9]\AgdaSymbol{;} \AgdaField{trans} \AgdaSymbol{=} \AgdaSymbol{λ} \AgdaSymbol{\{}\AgdaBound{x}\AgdaSymbol{\}} \AgdaSymbol{\{}\AgdaBound{y}\AgdaSymbol{\}} \AgdaSymbol{\{}\AgdaBound{z}\AgdaSymbol{\}} \AgdaBound{x≈y} \AgdaBound{y≈z} \AgdaBound{R'} \AgdaBound{x₁} \AgdaBound{x₂} \AgdaBound{x₃} \AgdaBound{x₄} \<[56]%
\>[56]\<%
\\
\>[9]\AgdaIndent{19}{}\<[19]%
\>[19]\AgdaSymbol{→} \AgdaBound{x₂} \AgdaBound{x} \AgdaBound{y} \AgdaBound{z} \AgdaSymbol{(}\AgdaBound{x≈y} \AgdaBound{R'} \AgdaBound{x₁} \AgdaBound{x₂} \AgdaBound{x₃} \AgdaBound{x₄}\AgdaSymbol{)} \AgdaSymbol{(}\AgdaBound{y≈z} \AgdaBound{R'} \AgdaBound{x₁} \AgdaBound{x₂} \AgdaBound{x₃} \AgdaBound{x₄}\AgdaSymbol{)}\<%
\\
\>[0]\AgdaIndent{9}{}\<[9]%
\>[9]\AgdaSymbol{\}}\<%
\\
%
\\
%
\\
\>[0]\AgdaIndent{2}{}\<[2]%
\>[2]\AgdaFunction{⟦Q⟧} \AgdaSymbol{:} \AgdaRecord{Ty} \AgdaBound{Γ}\<%
\\
\>[0]\AgdaIndent{2}{}\<[2]%
\>[2]\AgdaFunction{⟦Q⟧} \AgdaSymbol{=} \AgdaKeyword{record} \<[15]%
\>[15]\<%
\\
\>[2]\AgdaIndent{4}{}\<[4]%
\>[4]\AgdaSymbol{\{} \AgdaField{fm} \AgdaSymbol{=} \AgdaFunction{⟦Q⟧₀}\<%
\\
\>[2]\AgdaIndent{4}{}\<[4]%
\>[4]\AgdaSymbol{;} \AgdaField{substT} \AgdaSymbol{=} \AgdaFunction{[} \AgdaBound{A} \AgdaFunction{]subst}\<%
\\
\>[2]\AgdaIndent{4}{}\<[4]%
\>[4]\AgdaSymbol{;} \AgdaField{subst*} \AgdaSymbol{=} \AgdaSymbol{λ} \AgdaSymbol{\{}\AgdaBound{γ}\AgdaSymbol{\}} \AgdaSymbol{\{}\AgdaBound{γ'}\AgdaSymbol{\}} \AgdaBound{p} \AgdaSymbol{\{}\AgdaBound{a}\AgdaSymbol{\}} \AgdaSymbol{\{}\AgdaBound{b}\AgdaSymbol{\}} \AgdaSymbol{→} \AgdaFunction{Qsubst*} \AgdaBound{γ} \AgdaBound{γ'} \AgdaBound{p} \AgdaBound{a} \AgdaBound{b}\<%
\\
\>[2]\AgdaIndent{4}{}\<[4]%
\>[4]\AgdaSymbol{;} \AgdaField{refl*} \AgdaSymbol{=} \AgdaSymbol{λ} \AgdaBound{x} \AgdaBound{a} \AgdaBound{R'} \AgdaBound{x₁} \AgdaBound{x₂} \AgdaBound{x₃} \AgdaBound{x₄} \AgdaSymbol{→} \AgdaBound{x₃} \AgdaSymbol{\_} \AgdaSymbol{\_} \AgdaSymbol{(}\AgdaFunction{[} \AgdaBound{A} \AgdaFunction{]refl*} \AgdaSymbol{\_} \AgdaSymbol{\_)}\<%
\\
\>[2]\AgdaIndent{4}{}\<[4]%
\>[4]\AgdaSymbol{;} \AgdaField{trans*} \AgdaSymbol{=} \AgdaSymbol{λ} \AgdaBound{a} \AgdaBound{R'} \AgdaBound{x₁} \AgdaBound{x₂} \AgdaBound{x₃} \AgdaBound{x₄} \AgdaSymbol{→} \AgdaBound{x₃} \AgdaSymbol{\_} \AgdaSymbol{\_} \AgdaSymbol{(}\AgdaFunction{[} \AgdaBound{A} \AgdaFunction{]trans*} \AgdaSymbol{\_)}\<%
\\
\>[2]\AgdaIndent{4}{}\<[4]%
\>[4]\AgdaSymbol{\}}\<%
\\
\end{code}



\section{Observational equality}

%definitional distinct types

Later in in \cite{alti:ott-conf}, Altenkirch and McBride further
simplifies the setoid model by adopting McBride's heterogeneous
approach to equality. They identifies values up to observation rather than
  construction which is called observational equality. It is the
  propositional equality induced by the Setoid model.  In general we have a heterogeneous equality which
  compares terms of types which are different in construction. It only
  make sense when we can prove the types are the same. It helps us
  avoids the heavy use of $subst$ which makes formalisation and
  reasoning involved. We could simplify the setoid model by adapting this
  approach and the implementation could be easier.











\section{Groupoid model}




\section{Models of \hott}




\cite{bezem2013model}

\hott can be seen as Type Theory together with several axioms, like the univalence axiom and higher inductive types. We can postulate these axioms but then we can not keep the good computational properties. Therefore it is crucial that axioms have a computational interpretation. One solution is 









There are some models of \hott currently being studied extensively.
To interpret types as weak $\omega$-groupoids, one main problems is
the complexity of the definition of weak $\omega$-groupoid. The
coherence conditions are very difficult to specified.
It is much simpler to interpret types as Kan simplicial sets.
Voevodsky's univalent model\cite{klv:ssetmodel} is based on Kan simplicial sets. 
 Streicher wrote a concise introduction to this model \cite{DBLP:dblp_journals/japll/Streicher14}. 



\iffalse % comment out multiple line

\begin{definition}
A simplicial set $X$ is a functor from $\Delta^{op}$ to $\Set$ where
$\Delta$ is the simplex category.
\end{definition}

$\Delta^{op}$ is a category whose objects are non-empty totally ordered
finite sets. The morphisms are order-preserving functions. 
Face maps and degeneracy maps are the most important morphisms in this
category.

A simplexes is a generalisation of a triangle to arbitrary
dimensions. $3$-dimensional simplex is tetrahedron and $k+1$-simplex can
be obtained by adding one point to $k$-simplex which does not lie in the
dimension where the $k$-simplex is.

A simplicial complex is a collection of simplexes. Topologically speaking, it
is constructed by gluing n-dimensional simplexes together. 
\todo{show an example graph}

A simplicial set, therefore, can be illustrate by the same graph where
the set of points is given by $X_0$, the set of lines is $X_1$ and so
on. The graph looks very similar to n-groupoids. In fact simplical
sets can interpret types in a similar way (?).

\fi




However the simplicial set model is not constructive as Coquand showed
that it requires classical logic in an essential way \cite{TC:sset}.


The distinction of semi-simplicial set is there is only face maps
but no degeneracy maps. We can denote a semi-simplicial set as a
functor $X : \Delta_{inj} \rightarrow |Set$. The morphisms in $\Delta_{inj}$ are not only order preserving
but also injective.
A “iterated dependency” approach is believed to solve the coherence
issues.


\subsection{semi-simplicial sets}

\todo{cite https://uf-ias-2012.wikispaces.com/Semi-simplicial+types}


%Klaus and me were trying to implement semi-simplical set in Agda. 


\subsection{Cubical sets}

Someone can deduce what is a cubical set from the name literaly. It is
also a functor $S : \Box^{op} \rightarrow \Set$ (or a presheaf on the
cube category $ \Box^{op}$).

\section{Summary}

In this chapter we introduce the basic notions in \hott, discuss the
various implementations of \hott. In the next chapter we will focus on
thesyntactic implementation of \wog following
Brunerie's approach. We attempt to formalise the groupoid model of
\hott in intensional type theory, specially in Agda.







% In \itt, the uniqueness of identity types is not
% accepted in general, but derivable for types whose propositional
% equality is decidable. The homotopy interpretation fits
% nicely by provides higher levels structures which are weaker
% equivalence relation (compared to strict equality) between identity types.


However the implementation of \hott in \itt is still an open problem. We
work on defining semi-simplicial types and \wog to solve this problem. There
is also the very new model using cubical sets proposed by Bezem,
Coquand and Huber in \cite{bezem2013model}.



\hott does not only help us model type theory with a focus on the equality, but also provides mathematicians type theoretical tools to study homotopy theory.



\chapter{\og Model}
\label{wog}

\todo{Before 15th-Jan-2013}



It is possible to
define the quotient types in Homotopy Type Theory but to implement the
Homotopy Type Theory in \itt{}, it is still a difficult problem. We
work on defining semi-simplicial types and \wog{} to solve it. 



In this chapter we present an syntactic implementation of \wog{} following
Brunerie's approach. We did some contributions like adapting using
heterogeneous equality and syntactic construction of reflexivity.




\section{Quotient types in \hott}


\section{Syntax of \wog}


\section{Semantics}



\section{Higher inductive types}

%\chapter{Summary}


\subsection{Future work}

%----------------------------------------------------------------------------------------
%	THESIS CONTENT - APPENDICES
%----------------------------------------------------------------------------------------

\addtocontents{toc}{\vspace{2em}} % Add a gap in the Contents, for aesthetics

\appendix % Cue to tell LaTeX that the following 'chapters' are Appendices

% Include the appendices of the thesis as separate files from the Appendices folder
% Uncomment the lines as you write the Appendices

\chapter{Fractions reduction}

\begin{code}\label{rational-gcd}
\\
\>\AgdaFunction{GCD′→ℚ} \AgdaSymbol{:} \AgdaSymbol{∀} \AgdaBound{x} \AgdaBound{y} \AgdaBound{di} \AgdaSymbol{→} \AgdaBound{y} \AgdaFunction{≢} \AgdaNumber{0} \AgdaSymbol{→} \AgdaDatatype{C.GCD′} \AgdaBound{x} \AgdaBound{y} \AgdaBound{di} \AgdaSymbol{→} \AgdaRecord{ℚ}\<%
\\
\>\AgdaFunction{GCD′→ℚ} \AgdaSymbol{.(}\AgdaBound{q₁} \AgdaPrimitive{ℕ*} \AgdaBound{di}\AgdaSymbol{)} \AgdaSymbol{.(}\AgdaBound{q₂} \AgdaPrimitive{ℕ*} \AgdaBound{di}\AgdaSymbol{)} \AgdaBound{di} \AgdaBound{neo} \AgdaSymbol{(}\AgdaInductiveConstructor{C.gcd-*} \AgdaBound{q₁} \AgdaBound{q₂} \AgdaBound{c}\AgdaSymbol{)} \AgdaSymbol{=} \AgdaKeyword{record} \AgdaSymbol{\{} \AgdaField{numerator} \AgdaSymbol{=} \AgdaFunction{numr}\<%
\\
\>[0]\AgdaIndent{10}{}\<[10]%
\>[10]\AgdaSymbol{;} \AgdaField{denominator-1} \AgdaSymbol{=} \AgdaFunction{pred} \AgdaBound{q₂}\<%
\\
\>[0]\AgdaIndent{10}{}\<[10]%
\>[10]\AgdaSymbol{;} \AgdaField{isCoprime} \AgdaSymbol{=} \AgdaFunction{iscoprime} \AgdaSymbol{\}}\<%
\\
\>[0]\AgdaIndent{3}{}\<[3]%
\>[3]\AgdaKeyword{where}\<%
\\
\>[0]\AgdaIndent{5}{}\<[5]%
\>[5]\AgdaFunction{numr} \AgdaSymbol{=} \AgdaInductiveConstructor{ℤ.+\_} \AgdaBound{q₁}\<%
\\
\>[0]\AgdaIndent{5}{}\<[5]%
\>[5]\AgdaFunction{deno} \AgdaSymbol{=} \AgdaInductiveConstructor{suc} \AgdaSymbol{(}\AgdaFunction{pred} \AgdaBound{q₂}\AgdaSymbol{)}\<%
\\
\>[0]\AgdaIndent{5}{}\<[5]%
\>[5]\<%
\\
\>[0]\AgdaIndent{5}{}\<[5]%
\>[5]\AgdaFunction{numr≡q1} \AgdaSymbol{:} \AgdaFunction{∣} \AgdaFunction{numr} \AgdaFunction{∣} \AgdaDatatype{≡} \AgdaBound{q₁}\<%
\\
\>[0]\AgdaIndent{5}{}\<[5]%
\>[5]\AgdaFunction{numr≡q1} \AgdaSymbol{=} \AgdaInductiveConstructor{refl}\<%
\\
%
\\
\>[0]\AgdaIndent{5}{}\<[5]%
\>[5]\AgdaFunction{lzero} \AgdaSymbol{:} \AgdaSymbol{∀} \AgdaBound{x} \AgdaBound{y} \AgdaSymbol{→} \AgdaBound{x} \AgdaDatatype{≡} \AgdaNumber{0} \AgdaSymbol{→} \AgdaBound{x} \AgdaPrimitive{ℕ*} \AgdaBound{y} \AgdaDatatype{≡} \AgdaNumber{0}\<%
\\
\>[0]\AgdaIndent{5}{}\<[5]%
\>[5]\AgdaFunction{lzero} \AgdaSymbol{.}\AgdaNumber{0} \AgdaBound{y} \AgdaInductiveConstructor{refl} \AgdaSymbol{=} \AgdaInductiveConstructor{refl}\<%
\\
%
\\
\>[0]\AgdaIndent{5}{}\<[5]%
\>[5]\AgdaFunction{q2≢0} \AgdaSymbol{:} \AgdaBound{q₂} \AgdaFunction{≢} \AgdaNumber{0}\<%
\\
\>[0]\AgdaIndent{5}{}\<[5]%
\>[5]\AgdaFunction{q2≢0} \AgdaBound{qe} \AgdaSymbol{=} \AgdaBound{neo} \AgdaSymbol{(}\AgdaFunction{lzero} \AgdaBound{q₂} \AgdaBound{di} \AgdaBound{qe}\AgdaSymbol{)}\<%
\\
%
\\
\>[0]\AgdaIndent{5}{}\<[5]%
\>[5]\AgdaFunction{invsuc} \AgdaSymbol{:} \AgdaSymbol{∀} \AgdaBound{n} \AgdaSymbol{→} \AgdaBound{n} \AgdaFunction{≢} \AgdaNumber{0} \AgdaSymbol{→} \AgdaInductiveConstructor{suc} \AgdaSymbol{(}\AgdaFunction{pred} \AgdaBound{n}\AgdaSymbol{)} \AgdaDatatype{≡} \AgdaBound{n}\<%
\\
\>[0]\AgdaIndent{5}{}\<[5]%
\>[5]\AgdaFunction{invsuc} \AgdaInductiveConstructor{zero} \AgdaBound{nz} \AgdaKeyword{with} \AgdaBound{nz} \AgdaInductiveConstructor{refl}\<%
\\
\>[0]\AgdaIndent{5}{}\<[5]%
\>[5]\AgdaSymbol{...} \AgdaSymbol{|} \AgdaSymbol{()}\<%
\\
\>[0]\AgdaIndent{5}{}\<[5]%
\>[5]\AgdaFunction{invsuc} \AgdaSymbol{(}\AgdaInductiveConstructor{suc} \AgdaBound{n}\AgdaSymbol{)} \AgdaBound{nz} \AgdaSymbol{=} \AgdaInductiveConstructor{refl}\<%
\\
\>[0]\AgdaIndent{5}{}\<[5]%
\>[5]\<%
\\
\>[0]\AgdaIndent{5}{}\<[5]%
\>[5]\AgdaFunction{deno≡q2} \AgdaSymbol{:} \AgdaFunction{deno} \AgdaDatatype{≡} \AgdaBound{q₂}\<%
\\
\>[0]\AgdaIndent{5}{}\<[5]%
\>[5]\AgdaFunction{deno≡q2} \AgdaSymbol{=} \AgdaFunction{invsuc} \AgdaBound{q₂} \AgdaFunction{q2≢0}\<%
\\
\>[0]\AgdaIndent{5}{}\<[5]%
\>[5]\<%
\\
\>[0]\AgdaIndent{5}{}\<[5]%
\>[5]\AgdaFunction{transCop} \AgdaSymbol{:} \AgdaSymbol{∀} \AgdaSymbol{\{}\AgdaBound{a} \AgdaBound{b} \AgdaBound{c} \AgdaBound{d}\AgdaSymbol{\}} \AgdaSymbol{→} \AgdaBound{c} \AgdaDatatype{≡} \AgdaBound{a} \AgdaSymbol{→} \AgdaBound{d} \AgdaDatatype{≡} \AgdaBound{b} \AgdaSymbol{→} \AgdaFunction{C.Coprime} \AgdaBound{a} \AgdaBound{b} \AgdaSymbol{→} \AgdaFunction{C.Coprime} \AgdaBound{c} \AgdaBound{d}\<%
\\
\>[0]\AgdaIndent{5}{}\<[5]%
\>[5]\AgdaFunction{transCop} \AgdaInductiveConstructor{refl} \AgdaInductiveConstructor{refl} \AgdaBound{c} \AgdaSymbol{=} \AgdaBound{c}\<%
\\
%
\\
\>[0]\AgdaIndent{5}{}\<[5]%
\>[5]\AgdaFunction{copnd} \AgdaSymbol{:} \AgdaFunction{C.Coprime} \AgdaFunction{∣} \AgdaFunction{numr} \AgdaFunction{∣} \AgdaFunction{deno}\<%
\\
\>[0]\AgdaIndent{5}{}\<[5]%
\>[5]\AgdaFunction{copnd} \AgdaSymbol{=} \AgdaFunction{transCop} \AgdaFunction{numr≡q1} \AgdaFunction{deno≡q2} \AgdaBound{c}\<%
\\
%
\\
\>[0]\AgdaIndent{5}{}\<[5]%
\>[5]\AgdaFunction{witProp} \AgdaSymbol{:} \AgdaSymbol{∀} \AgdaBound{a} \AgdaBound{b} \AgdaSymbol{→} \AgdaRecord{GCD} \AgdaBound{a} \AgdaBound{b} \AgdaNumber{1} \AgdaSymbol{→} \AgdaFunction{True} \AgdaSymbol{(}\AgdaFunction{C.coprime?} \AgdaBound{a} \AgdaBound{b}\AgdaSymbol{)}\<%
\\
\>[0]\AgdaIndent{5}{}\<[5]%
\>[5]\AgdaFunction{witProp} \AgdaBound{a} \AgdaBound{b} \AgdaBound{gcd1} \AgdaKeyword{with} \AgdaFunction{gcd} \AgdaBound{a} \AgdaBound{b}\<%
\\
\>[0]\AgdaIndent{5}{}\<[5]%
\>[5]\AgdaFunction{witProp} \AgdaBound{a} \AgdaBound{b} \AgdaBound{gcd1} \AgdaSymbol{|} \AgdaInductiveConstructor{zero} \AgdaInductiveConstructor{,} \AgdaBound{y} \AgdaKeyword{with} \AgdaFunction{GCD.unique} \AgdaBound{gcd1} \AgdaBound{y}\<%
\\
\>[0]\AgdaIndent{5}{}\<[5]%
\>[5]\AgdaFunction{witProp} \AgdaBound{a} \AgdaBound{b} \AgdaBound{gcd1} \AgdaSymbol{|} \AgdaInductiveConstructor{zero} \AgdaInductiveConstructor{,} \AgdaBound{y} \AgdaSymbol{|} \AgdaSymbol{()}\<%
\\
\>[0]\AgdaIndent{5}{}\<[5]%
\>[5]\AgdaFunction{witProp} \AgdaBound{a} \AgdaBound{b} \AgdaBound{gcd1} \AgdaSymbol{|} \AgdaInductiveConstructor{suc} \AgdaInductiveConstructor{zero} \AgdaInductiveConstructor{,} \AgdaBound{y} \AgdaSymbol{=} \AgdaInductiveConstructor{tt}\<%
\\
\>[0]\AgdaIndent{5}{}\<[5]%
\>[5]\AgdaFunction{witProp} \AgdaBound{a} \AgdaBound{b} \AgdaBound{gcd1} \AgdaSymbol{|} \AgdaInductiveConstructor{suc} \AgdaSymbol{(}\AgdaInductiveConstructor{suc} \AgdaBound{n}\AgdaSymbol{)} \AgdaInductiveConstructor{,} \AgdaBound{y} \AgdaKeyword{with} \AgdaFunction{GCD.unique} \AgdaBound{gcd1} \AgdaBound{y}\<%
\\
\>[0]\AgdaIndent{5}{}\<[5]%
\>[5]\AgdaFunction{witProp} \AgdaBound{a} \AgdaBound{b} \AgdaBound{gcd1} \AgdaSymbol{|} \AgdaInductiveConstructor{suc} \AgdaSymbol{(}\AgdaInductiveConstructor{suc} \AgdaBound{n}\AgdaSymbol{)} \AgdaInductiveConstructor{,} \AgdaBound{y} \AgdaSymbol{|} \AgdaSymbol{()}\<%
\\
%
\\
\>[0]\AgdaIndent{5}{}\<[5]%
\>[5]\AgdaFunction{iscoprime} \AgdaSymbol{:} \AgdaFunction{True} \AgdaSymbol{(}\AgdaFunction{C.coprime?} \AgdaFunction{∣} \AgdaFunction{numr} \AgdaFunction{∣} \AgdaFunction{deno}\AgdaSymbol{)}\<%
\\
\>[0]\AgdaIndent{5}{}\<[5]%
\>[5]\AgdaFunction{iscoprime} \AgdaSymbol{=} \AgdaFunction{witProp} \AgdaFunction{∣} \AgdaFunction{numr} \AgdaFunction{∣} \AgdaFunction{deno} \AgdaSymbol{(}\AgdaFunction{C.coprime-gcd} \AgdaFunction{copnd}\AgdaSymbol{)}\<%
\\
\end{code}

\chapter{Set Quotients and Groupoid Quotients}


\begin{code}\>\<%
\\
%
\\
\>\AgdaKeyword{module} \AgdaModule{Quotient} \AgdaKeyword{where}\<%
\\
%
\\
\>\<\end{code}

\AgdaHide{
\begin{code}\>\<%
\\
\>\AgdaKeyword{open} \AgdaKeyword{import} \AgdaModule{Data.Product}\<%
\\
\>\AgdaKeyword{open} \AgdaKeyword{import} \AgdaModule{Function}\<%
\\
\>\AgdaKeyword{open} \AgdaKeyword{import} \AgdaModule{Level} \AgdaKeyword{using} \AgdaSymbol{(}\_⊔\_\AgdaSymbol{)}\<%
\\
\>\AgdaKeyword{import} \AgdaModule{Relation.Binary} \AgdaSymbol{as} \AgdaModule{RB}\<%
\\
\>\AgdaKeyword{open} \AgdaKeyword{import} \AgdaModule{Data.Nat} \AgdaKeyword{hiding} \AgdaSymbol{(}\_⊔\_\AgdaSymbol{)}\<%
\\
%
\\
\>\AgdaFunction{Setoid} \AgdaSymbol{=} \AgdaRecord{RB.Setoid} \AgdaPrimitive{Level.zero} \AgdaPrimitive{Level.zero}\<%
\\
%
\\
\>\AgdaKeyword{open} \AgdaKeyword{import} \AgdaModule{Relation.Binary.PropositionalEquality} \AgdaSymbol{as} \AgdaModule{PE}\<%
\\
\>[0]\AgdaIndent{2}{}\<[2]%
\>[2]\AgdaKeyword{hiding} \AgdaSymbol{(}[\_]\AgdaSymbol{)}\<%
\\
%
\\
\>\AgdaKeyword{open} \AgdaKeyword{import} \AgdaModule{ThomasProperties}\<%
\\
%
\\
%
\\
\>\AgdaFunction{\_respects\_} \AgdaSymbol{:} \AgdaSymbol{∀\{}\AgdaBound{α} \AgdaBound{β}\AgdaSymbol{\}\{}\AgdaBound{A} \AgdaSymbol{:} \AgdaPrimitiveType{Set} \AgdaBound{α}\AgdaSymbol{\}\{}\AgdaBound{B} \AgdaSymbol{:} \AgdaPrimitiveType{Set} \AgdaBound{β}\AgdaSymbol{\}(}\AgdaBound{f} \AgdaSymbol{:} \AgdaBound{A} \AgdaSymbol{→} \AgdaBound{B}\AgdaSymbol{)} \<[53]%
\>[53]\<%
\\
\>[2]\AgdaIndent{11}{}\<[11]%
\>[11]\AgdaSymbol{→} \AgdaSymbol{(}\AgdaBound{\_\textasciitilde\_} \AgdaSymbol{:} \AgdaBound{A} \AgdaSymbol{→} \AgdaBound{A} \AgdaSymbol{→} \AgdaPrimitiveType{Set} \AgdaBound{α}\AgdaSymbol{)} \AgdaSymbol{→} \AgdaPrimitiveType{Set} \AgdaSymbol{(}\AgdaBound{α} \AgdaPrimitive{⊔} \AgdaBound{β}\AgdaSymbol{)}\<%
\\
\>\AgdaBound{f} \AgdaFunction{respects} \AgdaBound{\_\textasciitilde\_} \AgdaSymbol{=} \AgdaSymbol{∀} \AgdaSymbol{\{}\AgdaBound{a} \AgdaBound{a'}\AgdaSymbol{\}} \AgdaSymbol{→} \AgdaBound{a} \AgdaBound{\textasciitilde} \AgdaBound{a'} \AgdaSymbol{→} \AgdaBound{f} \AgdaBound{a} \AgdaDatatype{≡} \AgdaBound{f} \AgdaBound{a'}\<%
\\
%
\\
\>\AgdaFunction{isProp} \AgdaSymbol{:} \AgdaSymbol{(}\AgdaBound{P} \AgdaSymbol{:} \AgdaPrimitiveType{Set}\AgdaSymbol{)} \AgdaSymbol{→} \AgdaPrimitiveType{Set}\<%
\\
\>\AgdaFunction{isProp} \AgdaBound{P} \AgdaSymbol{=} \<[12]%
\>[12]\AgdaSymbol{\{}\AgdaBound{p} \AgdaBound{p'} \AgdaSymbol{:} \AgdaBound{P}\AgdaSymbol{\}} \AgdaSymbol{→} \AgdaBound{p} \AgdaDatatype{≡} \AgdaBound{p'}\<%
\\
%
\\
\>\AgdaFunction{isSet} \AgdaSymbol{:} \AgdaSymbol{(}\AgdaBound{S} \AgdaSymbol{:} \AgdaPrimitiveType{Set}\AgdaSymbol{)} \AgdaSymbol{→} \AgdaPrimitiveType{Set}\<%
\\
\>\AgdaFunction{isSet} \AgdaBound{S} \AgdaSymbol{=} \AgdaSymbol{\{}\AgdaBound{a} \AgdaBound{b} \AgdaSymbol{:} \AgdaBound{S}\AgdaSymbol{\}} \AgdaSymbol{→} \AgdaFunction{isProp} \AgdaSymbol{(}\AgdaBound{a} \AgdaDatatype{≡} \AgdaBound{b}\AgdaSymbol{)}\<%
\\
%
\\
\>\AgdaFunction{subIrr} \AgdaSymbol{:} \AgdaSymbol{\{}\AgdaBound{S} \AgdaSymbol{:} \AgdaPrimitiveType{Set}\AgdaSymbol{\}(}\AgdaBound{isS} \AgdaSymbol{:} \AgdaFunction{isSet} \AgdaBound{S}\AgdaSymbol{)}\<%
\\
\>[0]\AgdaIndent{9}{}\<[9]%
\>[9]\AgdaSymbol{\{}\AgdaBound{A} \AgdaSymbol{:} \AgdaBound{S} \AgdaSymbol{→} \AgdaPrimitiveType{Set}\AgdaSymbol{\}\{}\AgdaBound{a} \AgdaBound{b} \AgdaSymbol{:} \AgdaBound{S}\AgdaSymbol{\}(}\AgdaBound{p} \AgdaBound{q} \AgdaSymbol{:} \AgdaBound{a} \AgdaDatatype{≡} \AgdaBound{b}\AgdaSymbol{)\{}\AgdaBound{m} \AgdaSymbol{:} \AgdaBound{A} \AgdaBound{a}\AgdaSymbol{\}}\<%
\\
\>[0]\AgdaIndent{7}{}\<[7]%
\>[7]\AgdaSymbol{→} \AgdaFunction{subst} \AgdaBound{A} \AgdaBound{p} \AgdaBound{m} \AgdaDatatype{≡} \AgdaFunction{subst} \AgdaBound{A} \AgdaBound{q} \AgdaBound{m}\<%
\\
\>\AgdaFunction{subIrr} \AgdaBound{isS} \AgdaBound{p} \AgdaBound{q} \AgdaSymbol{=} \AgdaFunction{cong} \AgdaSymbol{(λ} \AgdaBound{p'} \AgdaSymbol{→} \AgdaFunction{subst} \AgdaSymbol{\_} \AgdaBound{p'} \AgdaSymbol{\_)} \AgdaSymbol{(}\AgdaBound{isS} \AgdaSymbol{\{}p \AgdaSymbol{=} \AgdaBound{p}\AgdaSymbol{\}} \AgdaSymbol{\{}p' \AgdaSymbol{=} \AgdaBound{q}\AgdaSymbol{\})}\<%
\\
%
\\
\>\AgdaFunction{subIrr2} \AgdaSymbol{:} \AgdaSymbol{\{}\AgdaBound{S} \AgdaSymbol{:} \AgdaPrimitiveType{Set}\AgdaSymbol{\}\{}\AgdaBound{A} \AgdaSymbol{:} \AgdaPrimitiveType{Set}\AgdaSymbol{\}\{}\AgdaBound{a} \AgdaBound{b} \AgdaSymbol{:} \AgdaBound{S}\AgdaSymbol{\}(}\AgdaBound{p} \AgdaSymbol{:} \AgdaBound{a} \AgdaDatatype{≡} \AgdaBound{b}\AgdaSymbol{)\{}\AgdaBound{m} \AgdaSymbol{:} \AgdaBound{A}\AgdaSymbol{\}}\<%
\\
\>[0]\AgdaIndent{7}{}\<[7]%
\>[7]\AgdaSymbol{→} \AgdaFunction{subst} \AgdaSymbol{(λ} \AgdaBound{\_} \AgdaSymbol{→} \AgdaBound{A}\AgdaSymbol{)} \AgdaBound{p} \AgdaBound{m} \AgdaDatatype{≡} \AgdaBound{m}\<%
\\
\>\AgdaFunction{subIrr2} \AgdaInductiveConstructor{refl} \AgdaSymbol{=} \AgdaInductiveConstructor{refl}\<%
\\
%
\\
\>\<\end{code}
}

\textbf{Prequotient}

Given a setoid, we can turn it into a pre-quotient, corresponds to 
fork (or cofork) in category theory.

\begin{code}\>\<%
\\
\>\AgdaKeyword{record} \AgdaRecord{pre-Quotient} \AgdaSymbol{(}\AgdaBound{S} \AgdaSymbol{:} \AgdaFunction{Setoid}\AgdaSymbol{)} \AgdaSymbol{:} \AgdaPrimitiveType{Set₁} \AgdaKeyword{where}\<%
\\
\>\<\end{code}
\AgdaHide{
\begin{code}\>\<%
\\
\>[0]\AgdaIndent{2}{}\<[2]%
\>[2]\AgdaKeyword{open} \AgdaModule{RB.}\AgdaModule{Setoid} \AgdaBound{S} \AgdaKeyword{renaming} \AgdaSymbol{(}Carrier \AgdaSymbol{to} A\AgdaSymbol{;}
       \_≈\_ \AgdaSymbol{to} \_\textasciitilde\_ \AgdaSymbol{;} refl \AgdaSymbol{to} \textasciitilde-refl\AgdaSymbol{;} 
       sym \AgdaSymbol{to} \textasciitilde-sym\AgdaSymbol{;}
       trans \AgdaSymbol{to} \textasciitilde-trans\AgdaSymbol{)}\<%
\\
\>\<\end{code}
}
\begin{code}\>\<%
\\
\>[0]\AgdaIndent{2}{}\<[2]%
\>[2]\AgdaKeyword{field}\<%
\\
\>[2]\AgdaIndent{4}{}\<[4]%
\>[4]\AgdaField{Q} \<[8]%
\>[8]\AgdaSymbol{:} \AgdaPrimitiveType{Set}\<%
\\
\>[2]\AgdaIndent{4}{}\<[4]%
\>[4]\AgdaField{[\_]} \AgdaSymbol{:} \AgdaFunction{A} \AgdaSymbol{→} \AgdaBound{Q}\<%
\\
\>[2]\AgdaIndent{4}{}\<[4]%
\>[4]\AgdaField{[\_]⁼} \AgdaSymbol{:} \AgdaBound{[\_]} \AgdaFunction{respects} \AgdaFunction{\_\textasciitilde\_}\<%
\\
\>[2]\AgdaIndent{4}{}\<[4]%
\>[4]\AgdaField{QisSet} \AgdaSymbol{:} \AgdaFunction{isSet} \AgdaBound{Q}\<%
\\
\>\<\end{code}
\AgdaHide{
\begin{code}\>\<%
\\
\>[0]\AgdaIndent{2}{}\<[2]%
\>[2]\AgdaKeyword{open} \AgdaModule{RB.}\AgdaModule{Setoid} \AgdaBound{S} \AgdaKeyword{public} \AgdaKeyword{renaming} \<[35]%
\>[35]\<%
\\
\>[0]\AgdaIndent{7}{}\<[7]%
\>[7]\AgdaSymbol{(}Carrier \AgdaSymbol{to} A\AgdaSymbol{;} \_≈\_ \AgdaSymbol{to} \_\textasciitilde\_ 
       \AgdaSymbol{;} refl \AgdaSymbol{to} \textasciitilde-refl\AgdaSymbol{;} sym \AgdaSymbol{to} \textasciitilde-sym\AgdaSymbol{;}
       trans \AgdaSymbol{to} \textasciitilde-trans\AgdaSymbol{)}\<%
\\
\>\<\end{code}
}

\textbf{Quotient}

A prequotient with a dependent eliminator.

\begin{code}\>\<%
\\
\>\AgdaKeyword{record} \AgdaRecord{Quotient} \AgdaSymbol{\{}\AgdaBound{S} \AgdaSymbol{:} \AgdaFunction{Setoid}\AgdaSymbol{\}(}\AgdaBound{PQ} \AgdaSymbol{:} \AgdaRecord{pre-Quotient} \AgdaBound{S}\AgdaSymbol{)} \AgdaSymbol{:} \AgdaPrimitiveType{Set₁} \AgdaKeyword{where}\<%
\\
\>[0]\AgdaIndent{2}{}\<[2]%
\>[2]\AgdaKeyword{open} \AgdaModule{pre-Quotient} \AgdaBound{PQ}\<%
\\
\>[0]\AgdaIndent{2}{}\<[2]%
\>[2]\AgdaKeyword{field}\<%
\\
\>[2]\AgdaIndent{4}{}\<[4]%
\>[4]\AgdaField{qelim} \<[12]%
\>[12]\AgdaSymbol{:} \AgdaSymbol{\{}\AgdaBound{B} \AgdaSymbol{:} \AgdaFunction{Q} \AgdaSymbol{→} \AgdaPrimitiveType{Set}\AgdaSymbol{\}}\<%
\\
\>[4]\AgdaIndent{12}{}\<[12]%
\>[12]\AgdaSymbol{→} \AgdaSymbol{(}\AgdaBound{f} \AgdaSymbol{:} \AgdaSymbol{(}\AgdaBound{a} \AgdaSymbol{:} \AgdaFunction{A}\AgdaSymbol{)} \AgdaSymbol{→} \AgdaBound{B} \AgdaFunction{[} \AgdaBound{a} \AgdaFunction{]}\AgdaSymbol{)}\<%
\\
\>[4]\AgdaIndent{12}{}\<[12]%
\>[12]\AgdaSymbol{→} \AgdaSymbol{(∀} \AgdaSymbol{\{}\AgdaBound{a} \AgdaBound{a'}\AgdaSymbol{\}} \AgdaSymbol{→} \AgdaSymbol{(}\AgdaBound{p} \AgdaSymbol{:} \AgdaBound{a} \AgdaFunction{\textasciitilde} \AgdaBound{a'}\AgdaSymbol{)} \<[39]%
\>[39]\<%
\\
\>[4]\AgdaIndent{12}{}\<[12]%
\>[12]\AgdaSymbol{→} \AgdaFunction{subst} \AgdaBound{B} \AgdaFunction{[} \AgdaBound{p} \AgdaFunction{]⁼} \AgdaSymbol{(}\AgdaBound{f} \AgdaBound{a}\AgdaSymbol{)} \AgdaDatatype{≡} \AgdaBound{f} \AgdaBound{a'}\AgdaSymbol{)}\<%
\\
\>[4]\AgdaIndent{12}{}\<[12]%
\>[12]\AgdaSymbol{→} \AgdaSymbol{(}\AgdaBound{q} \AgdaSymbol{:} \AgdaFunction{Q}\AgdaSymbol{)} \AgdaSymbol{→} \AgdaBound{B} \AgdaBound{q}\<%
\\
\>[0]\AgdaIndent{4}{}\<[4]%
\>[4]\AgdaField{qelim-β} \AgdaSymbol{:} \AgdaSymbol{∀} \AgdaSymbol{\{}\AgdaBound{B} \AgdaBound{a} \AgdaBound{f}\AgdaSymbol{\}(}\AgdaBound{resp} \AgdaSymbol{:} \AgdaSymbol{(∀} \AgdaSymbol{\{}\AgdaBound{a} \AgdaBound{a'}\AgdaSymbol{\}} \AgdaSymbol{→} \AgdaSymbol{(}\AgdaBound{p} \AgdaSymbol{:} \AgdaBound{a} \AgdaFunction{\textasciitilde} \AgdaBound{a'}\AgdaSymbol{)} \<[56]%
\>[56]\<%
\\
\>[0]\AgdaIndent{12}{}\<[12]%
\>[12]\AgdaSymbol{→} \AgdaFunction{subst} \AgdaBound{B} \AgdaFunction{[} \AgdaBound{p} \AgdaFunction{]⁼} \AgdaSymbol{(}\AgdaBound{f} \AgdaBound{a}\AgdaSymbol{)} \AgdaDatatype{≡} \AgdaBound{f} \AgdaBound{a'}\AgdaSymbol{))}\<%
\\
\>[0]\AgdaIndent{12}{}\<[12]%
\>[12]\AgdaSymbol{→} \AgdaBound{qelim} \AgdaSymbol{\{}\AgdaBound{B}\AgdaSymbol{\}} \AgdaBound{f} \AgdaBound{resp} \AgdaFunction{[} \AgdaBound{a} \AgdaFunction{]} \AgdaDatatype{≡} \AgdaBound{f} \AgdaBound{a}\<%
\\
\>\<\end{code}

\textbf{Quotient (Hofmann's)}

A prequotient with a non-dependent eliminator (lifting).

\begin{code}\>\<%
\\
\>\AgdaKeyword{record} \AgdaRecord{Hof-Quotient} \AgdaSymbol{\{}\AgdaBound{S} \AgdaSymbol{:} \AgdaFunction{Setoid}\AgdaSymbol{\}}\<%
\\
\>[0]\AgdaIndent{7}{}\<[7]%
\>[7]\AgdaSymbol{(}\AgdaBound{PQ} \AgdaSymbol{:} \AgdaRecord{pre-Quotient} \AgdaBound{S}\AgdaSymbol{)} \AgdaSymbol{:} \AgdaPrimitiveType{Set₁} \AgdaKeyword{where}\<%
\\
\>[0]\AgdaIndent{2}{}\<[2]%
\>[2]\AgdaKeyword{open} \AgdaModule{pre-Quotient} \AgdaBound{PQ}\<%
\\
\>[0]\AgdaIndent{2}{}\<[2]%
\>[2]\AgdaKeyword{field}\<%
\\
\>[2]\AgdaIndent{4}{}\<[4]%
\>[4]\AgdaField{lift} \<[11]%
\>[11]\AgdaSymbol{:} \AgdaSymbol{\{}\AgdaBound{B} \AgdaSymbol{:} \AgdaPrimitiveType{Set}\AgdaSymbol{\}}\<%
\\
\>[4]\AgdaIndent{11}{}\<[11]%
\>[11]\AgdaSymbol{→} \AgdaSymbol{(}\AgdaBound{f} \AgdaSymbol{:} \AgdaFunction{A} \AgdaSymbol{→} \AgdaBound{B}\AgdaSymbol{)}\<%
\\
\>[4]\AgdaIndent{11}{}\<[11]%
\>[11]\AgdaSymbol{→} \AgdaBound{f} \AgdaFunction{respects} \AgdaFunction{\_\textasciitilde\_}\<%
\\
\>[4]\AgdaIndent{11}{}\<[11]%
\>[11]\AgdaSymbol{→} \AgdaFunction{Q} \AgdaSymbol{→} \AgdaBound{B}\<%
\\
%
\\
\>[0]\AgdaIndent{4}{}\<[4]%
\>[4]\AgdaField{lift-β} \AgdaSymbol{:} \AgdaSymbol{∀} \AgdaSymbol{\{}\AgdaBound{B} \AgdaBound{a} \AgdaBound{f}\AgdaSymbol{\}(}\AgdaBound{resp} \AgdaSymbol{:} \AgdaBound{f} \AgdaFunction{respects} \AgdaFunction{\_\textasciitilde\_}\AgdaSymbol{)} \<[46]%
\>[46]\<%
\\
\>[0]\AgdaIndent{11}{}\<[11]%
\>[11]\AgdaSymbol{→} \AgdaBound{lift} \AgdaSymbol{\{}\AgdaBound{B}\AgdaSymbol{\}} \AgdaBound{f} \AgdaBound{resp} \AgdaFunction{[} \AgdaBound{a} \AgdaFunction{]} \AgdaDatatype{≡} \AgdaBound{f} \AgdaBound{a}\<%
\\
%
\\
\>[0]\AgdaIndent{4}{}\<[4]%
\>[4]\AgdaField{qind} \<[11]%
\>[11]\AgdaSymbol{:} \AgdaSymbol{∀} \AgdaSymbol{(}\AgdaBound{P} \AgdaSymbol{:} \AgdaFunction{Q} \AgdaSymbol{→} \AgdaPrimitiveType{Set}\AgdaSymbol{)}\<%
\\
\>[0]\AgdaIndent{11}{}\<[11]%
\>[11]\AgdaSymbol{→} \AgdaSymbol{(∀\{}\AgdaBound{x}\AgdaSymbol{\}} \AgdaSymbol{→} \AgdaFunction{isProp} \AgdaSymbol{(}\AgdaBound{P} \AgdaBound{x}\AgdaSymbol{))}\<%
\\
\>[0]\AgdaIndent{11}{}\<[11]%
\>[11]\AgdaSymbol{→} \AgdaSymbol{(∀} \AgdaBound{a} \AgdaSymbol{→} \AgdaBound{P} \AgdaFunction{[} \AgdaBound{a} \AgdaFunction{]}\AgdaSymbol{)}\<%
\\
\>[0]\AgdaIndent{11}{}\<[11]%
\>[11]\AgdaSymbol{→} \AgdaSymbol{(∀} \AgdaBound{x} \AgdaSymbol{→} \AgdaBound{P} \AgdaBound{x}\AgdaSymbol{)}\<%
\\
\>\<\end{code}

\textbf{Exact quotient}

\begin{code}\>\<%
\\
\>\AgdaKeyword{record} \AgdaRecord{exact-Quotient} \AgdaSymbol{\{}\AgdaBound{S} \AgdaSymbol{:} \AgdaFunction{Setoid}\AgdaSymbol{\}}\<%
\\
\>[0]\AgdaIndent{7}{}\<[7]%
\>[7]\AgdaSymbol{(}\AgdaBound{PQ} \AgdaSymbol{:} \AgdaRecord{pre-Quotient} \AgdaBound{S}\AgdaSymbol{)} \AgdaSymbol{:} \AgdaPrimitiveType{Set₁} \AgdaKeyword{where}\<%
\\
\>[0]\AgdaIndent{2}{}\<[2]%
\>[2]\AgdaKeyword{open} \AgdaModule{pre-Quotient} \AgdaBound{PQ}\<%
\\
\>[0]\AgdaIndent{2}{}\<[2]%
\>[2]\AgdaKeyword{field}\<%
\\
\>[2]\AgdaIndent{4}{}\<[4]%
\>[4]\AgdaField{Qu} \<[10]%
\>[10]\AgdaSymbol{:} \AgdaRecord{Quotient} \AgdaBound{PQ}\<%
\\
\>[2]\AgdaIndent{4}{}\<[4]%
\>[4]\AgdaField{exact} \AgdaSymbol{:} \AgdaSymbol{∀} \AgdaSymbol{\{}\AgdaBound{a} \AgdaBound{b} \AgdaSymbol{:} \AgdaFunction{A}\AgdaSymbol{\}} \AgdaSymbol{→} \AgdaFunction{[} \AgdaBound{a} \AgdaFunction{]} \AgdaDatatype{≡} \AgdaFunction{[} \AgdaBound{b} \AgdaFunction{]} \AgdaSymbol{→} \AgdaBound{a} \AgdaFunction{\textasciitilde} \AgdaBound{b}\<%
\\
\>\<\end{code}


\textbf{Definable quotient}

\begin{code}\>\<%
\\
\>\AgdaKeyword{record} \AgdaRecord{def-Quotient} \AgdaSymbol{\{}\AgdaBound{S} \AgdaSymbol{:} \AgdaFunction{Setoid}\AgdaSymbol{\}}\<%
\\
\>[4]\AgdaIndent{7}{}\<[7]%
\>[7]\AgdaSymbol{(}\AgdaBound{PQ} \AgdaSymbol{:} \AgdaRecord{pre-Quotient} \AgdaBound{S}\AgdaSymbol{)} \AgdaSymbol{:} \AgdaPrimitiveType{Set₁} \AgdaKeyword{where}\<%
\\
\>[1]\AgdaIndent{2}{}\<[2]%
\>[2]\AgdaKeyword{open} \AgdaModule{pre-Quotient} \AgdaBound{PQ}\<%
\\
\>[0]\AgdaIndent{2}{}\<[2]%
\>[2]\AgdaKeyword{field}\<%
\\
\>[2]\AgdaIndent{4}{}\<[4]%
\>[4]\AgdaField{emb} \<[13]%
\>[13]\AgdaSymbol{:} \AgdaFunction{Q} \AgdaSymbol{→} \AgdaFunction{A}\<%
\\
\>[2]\AgdaIndent{4}{}\<[4]%
\>[4]\AgdaField{complete} \AgdaSymbol{:} \AgdaSymbol{∀} \AgdaBound{a} \AgdaSymbol{→} \AgdaBound{emb} \AgdaFunction{[} \AgdaBound{a} \AgdaFunction{]} \AgdaFunction{\textasciitilde} \AgdaBound{a}\<%
\\
\>[2]\AgdaIndent{4}{}\<[4]%
\>[4]\AgdaField{stable} \<[13]%
\>[13]\AgdaSymbol{:} \AgdaSymbol{∀} \AgdaBound{q} \AgdaSymbol{→} \AgdaFunction{[} \AgdaBound{emb} \AgdaBound{q} \AgdaFunction{]} \AgdaDatatype{≡} \AgdaBound{q}\<%
\\
\>\<\end{code}

\textbf{Proof :}\label{DQisExact} Definable quotients are exact.

\begin{code}\>\<%
\\
\>[0]\AgdaIndent{2}{}\<[2]%
\>[2]\AgdaFunction{exact} \AgdaSymbol{:} \AgdaSymbol{∀\{}\AgdaBound{a} \AgdaBound{b}\AgdaSymbol{\}} \AgdaSymbol{→} \AgdaFunction{[} \AgdaBound{a} \AgdaFunction{]} \AgdaDatatype{≡} \AgdaFunction{[} \AgdaBound{b} \AgdaFunction{]} \AgdaSymbol{→} \AgdaBound{a} \AgdaFunction{\textasciitilde} \AgdaBound{b}\<%
\\
\>[0]\AgdaIndent{2}{}\<[2]%
\>[2]\AgdaFunction{exact} \AgdaSymbol{\{}\AgdaBound{a}\AgdaSymbol{\}} \AgdaSymbol{\{}\AgdaBound{b}\AgdaSymbol{\}} \AgdaBound{p} \AgdaSymbol{=} \<[20]%
\>[20]\<%
\\
\>[2]\AgdaIndent{4}{}\<[4]%
\>[4]\AgdaFunction{\textasciitilde-trans} \AgdaSymbol{(}\AgdaFunction{\textasciitilde-sym} \AgdaSymbol{(}\AgdaFunction{complete} \AgdaBound{a}\AgdaSymbol{))} \<[33]%
\>[33]\<%
\\
\>[2]\AgdaIndent{4}{}\<[4]%
\>[4]\AgdaSymbol{(}\AgdaFunction{\textasciitilde-trans} \AgdaSymbol{(}\AgdaFunction{subst} \AgdaSymbol{(λ} \AgdaBound{x} \AgdaSymbol{→} \<[27]%
\>[27]\<%
\\
\>[2]\AgdaIndent{4}{}\<[4]%
\>[4]\AgdaFunction{emb} \AgdaFunction{[} \AgdaBound{a} \AgdaFunction{]} \AgdaFunction{\textasciitilde} \AgdaFunction{emb} \AgdaBound{x}\AgdaSymbol{)} \<[23]%
\>[23]\<%
\\
\>[2]\AgdaIndent{4}{}\<[4]%
\>[4]\AgdaBound{p} \AgdaFunction{\textasciitilde-refl}\AgdaSymbol{)} \AgdaSymbol{(}\AgdaFunction{complete} \AgdaBound{b}\AgdaSymbol{))}\<%
\\
\>\<\end{code}


\textbf{Equivalences and conversions among the quotient structures}

\AgdaHide{
\begin{code}\>\<%
\\
\>\AgdaFunction{Σeq} \AgdaSymbol{:} \AgdaSymbol{\{}\AgdaBound{A} \AgdaSymbol{:} \AgdaPrimitiveType{Set}\AgdaSymbol{\}\{}\AgdaBound{B} \AgdaSymbol{:} \AgdaBound{A} \AgdaSymbol{→} \AgdaPrimitiveType{Set}\AgdaSymbol{\}\{}\AgdaBound{a} \AgdaBound{a'} \AgdaSymbol{:} \AgdaBound{A}\AgdaSymbol{\}\{}\AgdaBound{b} \AgdaSymbol{:} \AgdaBound{B} \AgdaBound{a}\AgdaSymbol{\}\{}\AgdaBound{b'} \AgdaSymbol{:} \AgdaBound{B} \AgdaBound{a'}\AgdaSymbol{\}(}\AgdaBound{p} \AgdaSymbol{:} \AgdaBound{a} \AgdaDatatype{≡} \AgdaBound{a'}\AgdaSymbol{)} \<[71]%
\>[71]\<%
\\
\>[2]\AgdaIndent{4}{}\<[4]%
\>[4]\AgdaSymbol{→} \AgdaFunction{subst} \AgdaBound{B} \AgdaBound{p} \AgdaBound{b} \AgdaDatatype{≡} \AgdaBound{b'} \AgdaSymbol{→} \AgdaSymbol{(}\AgdaBound{a} \AgdaInductiveConstructor{,} \AgdaBound{b}\AgdaSymbol{)} \AgdaDatatype{≡} \AgdaSymbol{(}\AgdaBound{a'} \AgdaInductiveConstructor{,} \AgdaBound{b'}\AgdaSymbol{)}\<%
\\
\>\AgdaFunction{Σeq} \AgdaInductiveConstructor{refl} \AgdaInductiveConstructor{refl} \AgdaSymbol{=} \AgdaInductiveConstructor{refl}\<%
\\
%
\\
%
\\
\>\AgdaFunction{ind2dep} \AgdaSymbol{:} \AgdaSymbol{∀} \AgdaSymbol{\{}\AgdaBound{Q} \AgdaSymbol{:} \AgdaPrimitiveType{Set}\AgdaSymbol{\}\{}\AgdaBound{B} \AgdaSymbol{:} \AgdaBound{Q} \AgdaSymbol{→} \AgdaPrimitiveType{Set}\AgdaSymbol{\}}\<%
\\
\>[4]\AgdaIndent{8}{}\<[8]%
\>[8]\AgdaSymbol{→} \AgdaSymbol{(}\AgdaBound{f} \AgdaSymbol{:} \AgdaBound{Q} \AgdaSymbol{→} \AgdaRecord{Σ} \AgdaBound{Q} \AgdaBound{B}\AgdaSymbol{)}\<%
\\
\>[4]\AgdaIndent{8}{}\<[8]%
\>[8]\AgdaSymbol{→} \AgdaSymbol{(∀} \AgdaBound{q} \AgdaSymbol{→} \AgdaFunction{proj₁} \AgdaSymbol{(}\AgdaBound{f} \AgdaBound{q}\AgdaSymbol{)} \AgdaDatatype{≡} \AgdaBound{q}\AgdaSymbol{)}\<%
\\
\>[4]\AgdaIndent{8}{}\<[8]%
\>[8]\AgdaSymbol{→} \AgdaSymbol{(}\AgdaBound{q} \AgdaSymbol{:} \AgdaBound{Q}\AgdaSymbol{)} \AgdaSymbol{→} \AgdaBound{B} \AgdaBound{q}\<%
\\
\>\AgdaFunction{ind2dep} \AgdaSymbol{\{}\AgdaBound{Q}\AgdaSymbol{\}} \AgdaSymbol{\{}\AgdaBound{B}\AgdaSymbol{\}} \AgdaBound{f} \AgdaBound{id'} \AgdaBound{q} \AgdaSymbol{=} \AgdaFunction{subst} \AgdaBound{B} \AgdaSymbol{(}\AgdaBound{id'} \AgdaBound{q}\AgdaSymbol{)} \AgdaSymbol{(}\AgdaFunction{proj₂} \AgdaSymbol{(}\AgdaBound{f} \AgdaBound{q}\AgdaSymbol{))}\<%
\\
\>\<\end{code}
}

\emph{Hofmann's quotient} is equivalent to $Quotient$.

\begin{code}\>\<%
\\
\>\AgdaFunction{Hof-Quotient→Quotient} \AgdaSymbol{:} \AgdaSymbol{\{}\AgdaBound{S} \AgdaSymbol{:} \AgdaFunction{Setoid}\AgdaSymbol{\}} \AgdaSymbol{→} \AgdaSymbol{\{}\AgdaBound{PQ} \AgdaSymbol{:} \AgdaRecord{pre-Quotient} \AgdaBound{S}\AgdaSymbol{\}}\<%
\\
\>[-3]\AgdaIndent{7}{}\<[7]%
\>[7]\AgdaSymbol{→} \AgdaSymbol{(}\AgdaRecord{Hof-Quotient} \AgdaBound{PQ}\AgdaSymbol{)} \AgdaSymbol{→} \AgdaSymbol{(}\AgdaRecord{Quotient} \AgdaBound{PQ}\AgdaSymbol{)}\<%
\\
\>\AgdaFunction{Hof-Quotient→Quotient} \AgdaSymbol{\{}\AgdaBound{S}\AgdaSymbol{\}} \AgdaSymbol{\{}\AgdaBound{PQ}\AgdaSymbol{\}} \AgdaBound{QuH} \AgdaSymbol{=} \<[37]%
\>[37]\<%
\\
\>[0]\AgdaIndent{2}{}\<[2]%
\>[2]\AgdaKeyword{record} \<[9]%
\>[9]\<%
\\
\>[2]\AgdaIndent{4}{}\<[4]%
\>[4]\AgdaSymbol{\{} \AgdaField{qelim} \<[14]%
\>[14]\AgdaSymbol{=} \AgdaSymbol{λ} \AgdaSymbol{\{}\AgdaBound{B}\AgdaSymbol{\}} \AgdaBound{f} \AgdaBound{resp} \<[29]%
\>[29]\<%
\\
\>[2]\AgdaIndent{4}{}\<[4]%
\>[4]\AgdaSymbol{→} \AgdaFunction{proj₁} \AgdaSymbol{(}\AgdaFunction{qelim'} \AgdaBound{f} \AgdaBound{resp}\AgdaSymbol{)}\<%
\\
\>[2]\AgdaIndent{4}{}\<[4]%
\>[4]\AgdaSymbol{;} \AgdaField{qelim-β} \AgdaSymbol{=} \AgdaSymbol{λ} \AgdaSymbol{\{}\AgdaBound{B}\AgdaSymbol{\}} \AgdaSymbol{\{}\AgdaBound{a}\AgdaSymbol{\}} \AgdaSymbol{\{}\AgdaBound{f}\AgdaSymbol{\}} \AgdaBound{resp} \<[35]%
\>[35]\<%
\\
\>[2]\AgdaIndent{4}{}\<[4]%
\>[4]\AgdaSymbol{→} \AgdaFunction{proj₂} \AgdaSymbol{(}\AgdaFunction{qelim'} \AgdaBound{f} \AgdaBound{resp}\AgdaSymbol{)}\<%
\\
\>[2]\AgdaIndent{4}{}\<[4]%
\>[4]\AgdaSymbol{\}}\<%
\\
\>[0]\AgdaIndent{2}{}\<[2]%
\>[2]\AgdaKeyword{where}\<%
\\
\>[0]\AgdaIndent{4}{}\<[4]%
\>[4]\AgdaKeyword{open} \AgdaModule{pre-Quotient} \AgdaBound{PQ}\<%
\\
\>[0]\AgdaIndent{4}{}\<[4]%
\>[4]\AgdaKeyword{open} \AgdaModule{Hof-Quotient} \AgdaBound{QuH}\<%
\\
%
\\
\>[0]\AgdaIndent{4}{}\<[4]%
\>[4]\AgdaFunction{qelim'} \AgdaSymbol{:} \AgdaSymbol{\{}\AgdaBound{B} \AgdaSymbol{:} \AgdaFunction{Q} \AgdaSymbol{→} \AgdaPrimitiveType{Set}\AgdaSymbol{\}}\<%
\\
\>[4]\AgdaIndent{11}{}\<[11]%
\>[11]\AgdaSymbol{→} \AgdaSymbol{(}\AgdaBound{f} \AgdaSymbol{:} \AgdaSymbol{(}\AgdaBound{a} \AgdaSymbol{:} \AgdaFunction{A}\AgdaSymbol{)} \AgdaSymbol{→} \AgdaBound{B} \AgdaFunction{[} \AgdaBound{a} \AgdaFunction{]}\AgdaSymbol{)}\<%
\\
\>[4]\AgdaIndent{11}{}\<[11]%
\>[11]\AgdaSymbol{→} \AgdaSymbol{(∀} \AgdaSymbol{\{}\AgdaBound{a} \AgdaBound{a'}\AgdaSymbol{\}} \AgdaSymbol{→} \AgdaSymbol{(}\AgdaBound{p} \AgdaSymbol{:} \AgdaBound{a} \AgdaFunction{\textasciitilde} \AgdaBound{a'}\AgdaSymbol{)} \<[38]%
\>[38]\<%
\\
\>[4]\AgdaIndent{11}{}\<[11]%
\>[11]\AgdaSymbol{→} \AgdaFunction{subst} \AgdaBound{B} \AgdaFunction{[} \AgdaBound{p} \AgdaFunction{]⁼} \AgdaSymbol{(}\AgdaBound{f} \AgdaBound{a}\AgdaSymbol{)} \AgdaDatatype{≡} \AgdaBound{f} \AgdaBound{a'}\AgdaSymbol{)}\<%
\\
\>[4]\AgdaIndent{11}{}\<[11]%
\>[11]\AgdaSymbol{→} \AgdaRecord{Σ[} \AgdaBound{f\textasciicircum} \AgdaRecord{∶} \AgdaSymbol{((}\AgdaBound{q} \AgdaSymbol{:} \AgdaFunction{Q}\AgdaSymbol{)} \AgdaSymbol{→} \AgdaBound{B} \AgdaBound{q}\AgdaSymbol{)} \AgdaRecord{]} \<[39]%
\>[39]\<%
\\
\>[11]\AgdaIndent{14}{}\<[14]%
\>[14]\AgdaSymbol{(∀} \AgdaSymbol{\{}\AgdaBound{a}\AgdaSymbol{\}} \AgdaSymbol{→} \AgdaBound{f\textasciicircum} \AgdaFunction{[} \AgdaBound{a} \AgdaFunction{]} \AgdaDatatype{≡} \AgdaBound{f} \AgdaBound{a}\AgdaSymbol{)}\<%
\\
\>[-1]\AgdaIndent{4}{}\<[4]%
\>[4]\AgdaFunction{qelim'} \AgdaSymbol{\{}\AgdaBound{B}\AgdaSymbol{\}} \AgdaBound{f} \AgdaBound{resp} \AgdaSymbol{=} \<[25]%
\>[25]\AgdaFunction{f\textasciicircum} \AgdaInductiveConstructor{,} \AgdaFunction{f\textasciicircum-β}\<%
\\
\>[0]\AgdaIndent{10}{}\<[10]%
\>[10]\AgdaKeyword{where}\<%
\\
\>\<\end{code}

The dependent lifting is defined to be the second projection of a non-dependent lifting functions as follows:
 
\begin{code}\>\<%
\\
\>[10]\AgdaIndent{11}{}\<[11]%
\>[11]\AgdaFunction{f₀} \AgdaSymbol{:} \AgdaFunction{A} \AgdaSymbol{→} \AgdaRecord{Σ} \AgdaFunction{Q} \AgdaBound{B}\<%
\\
\>[10]\AgdaIndent{11}{}\<[11]%
\>[11]\AgdaFunction{f₀} \AgdaBound{a} \AgdaSymbol{=} \AgdaFunction{[} \AgdaBound{a} \AgdaFunction{]} \AgdaInductiveConstructor{,} \AgdaBound{f} \AgdaBound{a}\<%
\\
\>[-3]\AgdaIndent{4}{}\<[4]%
\>[4]\<%
\\
\>[0]\AgdaIndent{11}{}\<[11]%
\>[11]\AgdaFunction{resp₀} \AgdaSymbol{:} \AgdaFunction{f₀} \AgdaFunction{respects} \AgdaFunction{\_\textasciitilde\_}\<%
\\
\>[0]\AgdaIndent{11}{}\<[11]%
\>[11]\AgdaFunction{resp₀} \AgdaBound{p} \AgdaSymbol{=} \AgdaFunction{Σeq} \AgdaFunction{[} \AgdaBound{p} \AgdaFunction{]⁼} \AgdaSymbol{(}\AgdaBound{resp} \AgdaBound{p}\AgdaSymbol{)}\<%
\\
%
\\
%
\\
\>[0]\AgdaIndent{11}{}\<[11]%
\>[11]\AgdaFunction{f'} \AgdaSymbol{:} \AgdaFunction{Q} \AgdaSymbol{→} \AgdaRecord{Σ} \AgdaFunction{Q} \AgdaBound{B}\<%
\\
\>[0]\AgdaIndent{11}{}\<[11]%
\>[11]\AgdaFunction{f'} \AgdaSymbol{=} \AgdaFunction{lift} \AgdaFunction{f₀} \AgdaFunction{resp₀}\<%
\\
%
\\
\>[0]\AgdaIndent{11}{}\<[11]%
\>[11]\AgdaFunction{id'} \AgdaSymbol{:} \AgdaFunction{Q} \AgdaSymbol{→} \AgdaFunction{Q}\<%
\\
\>[0]\AgdaIndent{11}{}\<[11]%
\>[11]\AgdaFunction{id'} \AgdaSymbol{=} \AgdaFunction{proj₁} \AgdaFunction{∘} \AgdaFunction{f'}\<%
\\
\>[0]\AgdaIndent{11}{}\<[11]%
\>[11]\<%
\\
\>[0]\AgdaIndent{11}{}\<[11]%
\>[11]\AgdaFunction{P} \AgdaSymbol{:} \AgdaFunction{Q} \AgdaSymbol{→} \AgdaPrimitiveType{Set}\<%
\\
\>[0]\AgdaIndent{11}{}\<[11]%
\>[11]\AgdaFunction{P} \AgdaBound{q} \AgdaSymbol{=} \AgdaFunction{id'} \AgdaBound{q} \AgdaDatatype{≡} \AgdaBound{q}\<%
\\
%
\\
\>[0]\AgdaIndent{11}{}\<[11]%
\>[11]\AgdaFunction{f'-β} \AgdaSymbol{:} \AgdaSymbol{\{}\AgdaBound{a} \AgdaSymbol{:} \AgdaFunction{A}\AgdaSymbol{\}} \AgdaSymbol{→} \AgdaFunction{f'} \AgdaFunction{[} \AgdaBound{a} \AgdaFunction{]} \AgdaDatatype{≡} \AgdaFunction{[} \AgdaBound{a} \AgdaFunction{]} \AgdaInductiveConstructor{,} \AgdaBound{f} \AgdaBound{a}\<%
\\
\>[0]\AgdaIndent{11}{}\<[11]%
\>[11]\AgdaFunction{f'-β} \AgdaSymbol{=} \AgdaFunction{lift-β} \AgdaSymbol{\_}\<%
\\
%
\\
\>[0]\AgdaIndent{11}{}\<[11]%
\>[11]\AgdaFunction{f'-sound} \AgdaSymbol{:} \AgdaSymbol{∀\{}\AgdaBound{a}\AgdaSymbol{\}} \AgdaSymbol{→} \AgdaFunction{id'} \AgdaFunction{[} \AgdaBound{a} \AgdaFunction{]} \AgdaDatatype{≡} \AgdaFunction{[} \AgdaBound{a} \AgdaFunction{]}\<%
\\
\>[0]\AgdaIndent{11}{}\<[11]%
\>[11]\AgdaFunction{f'-sound} \AgdaSymbol{=} \AgdaFunction{cong} \AgdaFunction{proj₁} \AgdaFunction{f'-β}\<%
\\
%
\\
\>[0]\AgdaIndent{11}{}\<[11]%
\>[11]\AgdaFunction{f'-sound'} \AgdaSymbol{:} \AgdaSymbol{∀} \AgdaSymbol{\{}\AgdaBound{q}\AgdaSymbol{\}} \AgdaSymbol{→} \AgdaFunction{id'} \AgdaBound{q} \AgdaDatatype{≡} \AgdaBound{q}\<%
\\
\>[0]\AgdaIndent{11}{}\<[11]%
\>[11]\AgdaFunction{f'-sound'} \AgdaSymbol{\{}\AgdaBound{q}\AgdaSymbol{\}} \AgdaSymbol{=} \AgdaFunction{qind} \AgdaFunction{P} \AgdaFunction{QisSet} \<[41]%
\>[41]\<%
\\
\>[11]\AgdaIndent{27}{}\<[27]%
\>[27]\AgdaSymbol{(λ} \AgdaBound{\_} \AgdaSymbol{→} \AgdaFunction{f'-sound}\AgdaSymbol{)} \AgdaBound{q}\<%
\\
%
\\
\>[0]\AgdaIndent{11}{}\<[11]%
\>[11]\AgdaFunction{f'-sound2} \AgdaSymbol{:} \AgdaSymbol{∀} \AgdaSymbol{\{}\AgdaBound{a}\AgdaSymbol{\}} \AgdaSymbol{→} \<[31]%
\>[31]\<%
\\
\>[0]\AgdaIndent{13}{}\<[13]%
\>[13]\AgdaFunction{subst} \AgdaBound{B} \AgdaFunction{f'-sound} \AgdaSymbol{(}\AgdaFunction{proj₂} \AgdaSymbol{(}\AgdaFunction{f'} \AgdaFunction{[} \AgdaBound{a} \AgdaFunction{]}\AgdaSymbol{))} \AgdaDatatype{≡} \AgdaBound{f} \AgdaBound{a}\<%
\\
\>[0]\AgdaIndent{11}{}\<[11]%
\>[11]\AgdaFunction{f'-sound2} \AgdaSymbol{=} \AgdaFunction{cong-proj₂} \AgdaSymbol{\_} \AgdaSymbol{\_} \AgdaFunction{f'-β}\<%
\\
%
\\
\>[0]\AgdaIndent{11}{}\<[11]%
\>[11]\AgdaFunction{f\textasciicircum} \AgdaSymbol{:} \AgdaSymbol{(}\AgdaBound{q} \AgdaSymbol{:} \AgdaFunction{Q}\AgdaSymbol{)} \AgdaSymbol{→} \AgdaBound{B} \AgdaBound{q}\<%
\\
\>[0]\AgdaIndent{11}{}\<[11]%
\>[11]\AgdaFunction{f\textasciicircum} \AgdaBound{q} \AgdaSymbol{=} \AgdaFunction{subst} \AgdaBound{B} \AgdaSymbol{(}\AgdaFunction{f'-sound'}\AgdaSymbol{)} \AgdaSymbol{(}\AgdaFunction{proj₂} \AgdaSymbol{(}\AgdaFunction{f'} \AgdaBound{q}\AgdaSymbol{))}\<%
\\
\>[0]\AgdaIndent{11}{}\<[11]%
\>[11]\<%
\\
\>[0]\AgdaIndent{11}{}\<[11]%
\>[11]\AgdaFunction{f\textasciicircum-β} \AgdaSymbol{:} \AgdaSymbol{∀} \AgdaSymbol{\{}\AgdaBound{a}\AgdaSymbol{\}} \AgdaSymbol{→} \AgdaFunction{f\textasciicircum} \AgdaFunction{[} \AgdaBound{a} \AgdaFunction{]} \AgdaDatatype{≡} \AgdaBound{f} \AgdaBound{a}\<%
\\
\>[0]\AgdaIndent{11}{}\<[11]%
\>[11]\AgdaFunction{f\textasciicircum-β} \AgdaSymbol{\{}\AgdaBound{a}\AgdaSymbol{\}} \AgdaSymbol{=} \AgdaFunction{trans} \AgdaSymbol{(}\AgdaFunction{subIrr} \AgdaFunction{QisSet} \<[43]%
\>[43]\<%
\\
\>[11]\AgdaIndent{22}{}\<[22]%
\>[22]\AgdaFunction{f'-sound'} \AgdaFunction{f'-sound}\AgdaSymbol{)} \AgdaFunction{f'-sound2}\<%
\\
\>\<\end{code}

\begin{code}\>\<%
\\
\>\AgdaFunction{Quotient→Hof-Quotient} \AgdaSymbol{:} \<[24]%
\>[24]\<%
\\
\>[9]\AgdaIndent{2}{}\<[2]%
\>[2]\AgdaSymbol{\{}\AgdaBound{S} \AgdaSymbol{:} \AgdaFunction{Setoid}\AgdaSymbol{\}\{}\AgdaBound{PQ} \AgdaSymbol{:} \AgdaRecord{pre-Quotient} \AgdaBound{S}\AgdaSymbol{\}}\<%
\\
\>[0]\AgdaIndent{2}{}\<[2]%
\>[2]\AgdaSymbol{→} \AgdaSymbol{(}\AgdaRecord{Quotient} \AgdaBound{PQ}\AgdaSymbol{)}\<%
\\
\>[0]\AgdaIndent{2}{}\<[2]%
\>[2]\AgdaSymbol{→} \AgdaSymbol{(}\AgdaRecord{Hof-Quotient} \AgdaBound{PQ}\AgdaSymbol{)}\<%
\\
\>\AgdaFunction{Quotient→Hof-Quotient} \AgdaSymbol{\{}\AgdaBound{S}\AgdaSymbol{\}} \AgdaSymbol{\{}\AgdaBound{PQ}\AgdaSymbol{\}} \AgdaBound{QU} \AgdaSymbol{=}\<%
\\
\>[0]\AgdaIndent{2}{}\<[2]%
\>[2]\AgdaKeyword{record}\<%
\\
\>[0]\AgdaIndent{2}{}\<[2]%
\>[2]\AgdaSymbol{\{} \AgdaField{lift} \<[11]%
\>[11]\AgdaSymbol{=} \AgdaSymbol{λ} \AgdaBound{f} \AgdaBound{resp} \<[22]%
\>[22]\<%
\\
\>[2]\AgdaIndent{13}{}\<[13]%
\>[13]\AgdaSymbol{→} \AgdaFunction{qelim} \AgdaBound{f} \AgdaSymbol{(}\AgdaFunction{resp'} \AgdaBound{resp}\AgdaSymbol{)}\<%
\\
\>[0]\AgdaIndent{2}{}\<[2]%
\>[2]\AgdaSymbol{;} \AgdaField{lift-β} \AgdaSymbol{=} \AgdaSymbol{λ} \AgdaBound{resp} \<[20]%
\>[20]\<%
\\
\>[0]\AgdaIndent{13}{}\<[13]%
\>[13]\AgdaSymbol{→} \AgdaFunction{qelim-β} \AgdaSymbol{(}\AgdaFunction{resp'} \AgdaBound{resp}\AgdaSymbol{)}\<%
\\
\>[0]\AgdaIndent{2}{}\<[2]%
\>[2]\AgdaSymbol{;} \AgdaField{qind} \AgdaSymbol{=} \AgdaSymbol{λ} \AgdaBound{P} \AgdaBound{isP} \AgdaBound{f} \<[21]%
\>[21]\<%
\\
\>[0]\AgdaIndent{11}{}\<[11]%
\>[11]\AgdaSymbol{→} \AgdaFunction{qelim} \AgdaSymbol{\{}\AgdaBound{P}\AgdaSymbol{\}} \AgdaBound{f} \AgdaSymbol{(λ} \AgdaBound{\_} \AgdaSymbol{→} \AgdaBound{isP}\AgdaSymbol{)}\<%
\\
\>[0]\AgdaIndent{2}{}\<[2]%
\>[2]\AgdaSymbol{\}}\<%
\\
\>[0]\AgdaIndent{2}{}\<[2]%
\>[2]\AgdaKeyword{where}\<%
\\
\>[2]\AgdaIndent{4}{}\<[4]%
\>[4]\AgdaKeyword{open} \AgdaModule{pre-Quotient} \AgdaBound{PQ}\<%
\\
\>[2]\AgdaIndent{4}{}\<[4]%
\>[4]\AgdaKeyword{open} \AgdaModule{Quotient} \AgdaBound{QU}\<%
\\
%
\\
\>[2]\AgdaIndent{4}{}\<[4]%
\>[4]\AgdaFunction{resp'} \AgdaSymbol{:} \AgdaSymbol{\{}\AgdaBound{B} \AgdaSymbol{:} \AgdaPrimitiveType{Set}\AgdaSymbol{\}\{}\AgdaBound{a} \AgdaBound{a'} \AgdaSymbol{:} \AgdaFunction{A}\AgdaSymbol{\}}\<%
\\
\>[4]\AgdaIndent{10}{}\<[10]%
\>[10]\AgdaSymbol{\{}\AgdaBound{f} \AgdaSymbol{:} \AgdaFunction{A} \AgdaSymbol{→} \AgdaBound{B}\AgdaSymbol{\}}\<%
\\
\>[4]\AgdaIndent{10}{}\<[10]%
\>[10]\AgdaSymbol{(}\AgdaBound{resp} \AgdaSymbol{:} \AgdaBound{f} \AgdaFunction{respects} \AgdaFunction{\_\textasciitilde\_}\AgdaSymbol{)}\<%
\\
\>[4]\AgdaIndent{10}{}\<[10]%
\>[10]\AgdaSymbol{(}\AgdaBound{p} \AgdaSymbol{:} \AgdaBound{a} \AgdaFunction{\textasciitilde} \AgdaBound{a'}\AgdaSymbol{)}\<%
\\
\>[4]\AgdaIndent{10}{}\<[10]%
\>[10]\AgdaSymbol{→} \AgdaFunction{subst} \AgdaSymbol{(λ} \AgdaBound{\_} \AgdaSymbol{→} \AgdaBound{B}\AgdaSymbol{)} \AgdaFunction{[} \AgdaBound{p} \AgdaFunction{]⁼} \AgdaSymbol{(}\AgdaBound{f} \AgdaBound{a}\AgdaSymbol{)} \<[41]%
\>[41]\<%
\\
\>[4]\AgdaIndent{10}{}\<[10]%
\>[10]\AgdaDatatype{≡} \AgdaBound{f} \AgdaBound{a'}\<%
\\
\>[0]\AgdaIndent{4}{}\<[4]%
\>[4]\AgdaFunction{resp'} \AgdaBound{resp} \AgdaBound{p} \AgdaSymbol{=} \<[19]%
\>[19]\<%
\\
\>[0]\AgdaIndent{10}{}\<[10]%
\>[10]\AgdaFunction{trans} \AgdaSymbol{(}\AgdaFunction{subIrr2} \AgdaFunction{[} \AgdaBound{p} \AgdaFunction{]⁼}\AgdaSymbol{)} \<[33]%
\>[33]\<%
\\
\>[0]\AgdaIndent{10}{}\<[10]%
\>[10]\AgdaSymbol{(}\AgdaBound{resp} \AgdaBound{p}\AgdaSymbol{)}\<%
\\
\>\<\end{code}

\emph{Definable quotients} is the strongest and gives rise to a \emph{Quotient} (and ineed it is an \emph{Exact quotients}).

\begin{code}\>\<%
\\
\>\AgdaFunction{def-Quotient→Quotient} \AgdaSymbol{:} \<[24]%
\>[24]\<%
\\
\>[0]\AgdaIndent{2}{}\<[2]%
\>[2]\AgdaSymbol{\{}\AgdaBound{S} \AgdaSymbol{:} \AgdaFunction{Setoid}\AgdaSymbol{\}\{}\AgdaBound{PQ} \AgdaSymbol{:} \AgdaRecord{pre-Quotient} \AgdaBound{S}\AgdaSymbol{\}}\<%
\\
\>[0]\AgdaIndent{2}{}\<[2]%
\>[2]\AgdaSymbol{→} \AgdaSymbol{(}\AgdaRecord{def-Quotient} \AgdaBound{PQ}\AgdaSymbol{)} \AgdaSymbol{→} \AgdaSymbol{(}\AgdaRecord{Quotient} \AgdaBound{PQ}\AgdaSymbol{)}\<%
\\
\>\AgdaFunction{def-Quotient→Quotient} \AgdaSymbol{\{}\AgdaBound{S}\AgdaSymbol{\}} \AgdaSymbol{\{}\AgdaBound{PQ}\AgdaSymbol{\}} \AgdaBound{QuD} \AgdaSymbol{=} \<[37]%
\>[37]\<%
\\
\>[0]\AgdaIndent{2}{}\<[2]%
\>[2]\AgdaKeyword{record} \AgdaSymbol{\{} \AgdaField{qelim} \AgdaSymbol{=} \<[19]%
\>[19]\<%
\\
\>[2]\AgdaIndent{9}{}\<[9]%
\>[9]\AgdaSymbol{λ} \AgdaSymbol{\{}\AgdaBound{B}\AgdaSymbol{\}} \AgdaBound{f} \AgdaBound{resp} \AgdaBound{q} \AgdaSymbol{→} \AgdaFunction{subst} \AgdaBound{B} \AgdaSymbol{(}\AgdaFunction{stable} \AgdaBound{q}\AgdaSymbol{)} \AgdaSymbol{(}\AgdaBound{f} \AgdaSymbol{(}\AgdaFunction{emb} \AgdaBound{q}\AgdaSymbol{))}\<%
\\
\>[2]\AgdaIndent{9}{}\<[9]%
\>[9]\AgdaSymbol{;} \AgdaField{qelim-β} \AgdaSymbol{=} \<[21]%
\>[21]\<%
\\
\>[2]\AgdaIndent{9}{}\<[9]%
\>[9]\AgdaSymbol{λ} \AgdaSymbol{\{}\AgdaBound{B}\AgdaSymbol{\}} \AgdaSymbol{\{}\AgdaBound{a}\AgdaSymbol{\}} \AgdaSymbol{\{}\AgdaBound{f}\AgdaSymbol{\}} \AgdaBound{resp} \AgdaSymbol{→} \<[30]%
\>[30]\<%
\\
\>[2]\AgdaIndent{9}{}\<[9]%
\>[9]\AgdaFunction{trans} \AgdaSymbol{(}\AgdaFunction{subIrr} \AgdaFunction{QisSet} \AgdaSymbol{(}\AgdaFunction{stable} \AgdaFunction{[} \AgdaBound{a} \AgdaFunction{]}\AgdaSymbol{)} \<[45]%
\>[45]\<%
\\
\>[2]\AgdaIndent{9}{}\<[9]%
\>[9]\AgdaFunction{[} \AgdaFunction{complete} \AgdaBound{a} \AgdaFunction{]⁼}\AgdaSymbol{)} \AgdaSymbol{(}\AgdaBound{resp} \AgdaSymbol{(}\AgdaFunction{complete} \AgdaBound{a}\AgdaSymbol{))}\<%
\\
\>[0]\AgdaIndent{2}{}\<[2]%
\>[2]\AgdaSymbol{\}}\<%
\\
\>[0]\AgdaIndent{4}{}\<[4]%
\>[4]\AgdaKeyword{where}\<%
\\
\>[0]\AgdaIndent{4}{}\<[4]%
\>[4]\AgdaKeyword{open} \AgdaModule{pre-Quotient} \AgdaBound{PQ}\<%
\\
\>[0]\AgdaIndent{4}{}\<[4]%
\>[4]\AgdaKeyword{open} \AgdaModule{def-Quotient} \AgdaBound{QuD}\<%
\\
\>\<\end{code}

Remember that we have proven \emph{Definable quotients} are exact.

\begin{code}\>\<%
\\
\>\AgdaFunction{def-Quotient→exact-Quotient} \AgdaSymbol{:} \<[30]%
\>[30]\<%
\\
\>[0]\AgdaIndent{2}{}\<[2]%
\>[2]\AgdaSymbol{\{}\AgdaBound{S} \AgdaSymbol{:} \AgdaFunction{Setoid}\AgdaSymbol{\}\{}\AgdaBound{PQ} \AgdaSymbol{:} \AgdaRecord{pre-Quotient} \AgdaBound{S}\AgdaSymbol{\}}\<%
\\
\>[0]\AgdaIndent{2}{}\<[2]%
\>[2]\AgdaSymbol{→} \AgdaRecord{def-Quotient} \AgdaBound{PQ} \AgdaSymbol{→} \AgdaRecord{exact-Quotient} \AgdaBound{PQ}\<%
\\
\>\AgdaFunction{def-Quotient→exact-Quotient} \AgdaSymbol{\{}\AgdaBound{S}\AgdaSymbol{\}} \AgdaSymbol{\{}\AgdaBound{PQ}\AgdaSymbol{\}} \AgdaBound{QuD} \AgdaSymbol{=}\<%
\\
\>[0]\AgdaIndent{2}{}\<[2]%
\>[2]\AgdaKeyword{record} \AgdaSymbol{\{} \AgdaField{Qu} \AgdaSymbol{=} \AgdaFunction{def-Quotient→Quotient} \AgdaBound{QuD}\<%
\\
\>[2]\AgdaIndent{9}{}\<[9]%
\>[9]\AgdaSymbol{;} \AgdaField{exact} \AgdaSymbol{=} \AgdaFunction{exact}\<%
\\
\>[2]\AgdaIndent{9}{}\<[9]%
\>[9]\AgdaSymbol{\}}\<%
\\
\>[0]\AgdaIndent{2}{}\<[2]%
\>[2]\AgdaKeyword{where}\<%
\\
\>[0]\AgdaIndent{4}{}\<[4]%
\>[4]\AgdaKeyword{open} \AgdaModule{pre-Quotient} \AgdaBound{PQ}\<%
\\
\>[0]\AgdaIndent{4}{}\<[4]%
\>[4]\AgdaKeyword{open} \AgdaModule{def-Quotient} \AgdaBound{QuD}\<%
\\
\>\<\end{code}

\AgdaHide{
\begin{code}\>\<%
\\
\>\AgdaFunction{def-Quotient→Hof-Quotient} \<[26]%
\>[26]\<%
\\
\>[0]\AgdaIndent{2}{}\<[2]%
\>[2]\AgdaSymbol{:} \AgdaSymbol{\{}\AgdaBound{S} \AgdaSymbol{:} \AgdaFunction{Setoid}\AgdaSymbol{\}} \<[17]%
\>[17]\<%
\\
\>[0]\AgdaIndent{2}{}\<[2]%
\>[2]\AgdaSymbol{→} \AgdaSymbol{\{}\AgdaBound{PQ} \AgdaSymbol{:} \AgdaRecord{pre-Quotient} \AgdaBound{S}\AgdaSymbol{\}}\<%
\\
\>[0]\AgdaIndent{2}{}\<[2]%
\>[2]\AgdaSymbol{→} \AgdaSymbol{(}\AgdaRecord{def-Quotient} \AgdaBound{PQ}\AgdaSymbol{)} \<[22]%
\>[22]\<%
\\
\>[0]\AgdaIndent{2}{}\<[2]%
\>[2]\AgdaSymbol{→} \AgdaSymbol{(}\AgdaRecord{Hof-Quotient} \AgdaBound{PQ}\AgdaSymbol{)}\<%
\\
\>\AgdaFunction{def-Quotient→Hof-Quotient} \AgdaSymbol{\{}\AgdaBound{S}\AgdaSymbol{\}} \AgdaSymbol{\{}\AgdaBound{PQ}\AgdaSymbol{\}} \AgdaBound{QuD} \AgdaSymbol{=}\<%
\\
\>[0]\AgdaIndent{2}{}\<[2]%
\>[2]\AgdaKeyword{record} \<[9]%
\>[9]\<%
\\
\>[0]\AgdaIndent{2}{}\<[2]%
\>[2]\AgdaSymbol{\{} \AgdaField{lift} \<[11]%
\>[11]\AgdaSymbol{=} \AgdaSymbol{λ} \AgdaBound{f} \AgdaBound{\_} \AgdaSymbol{→} \AgdaBound{f} \AgdaFunction{∘} \AgdaFunction{emb}\<%
\\
\>[0]\AgdaIndent{2}{}\<[2]%
\>[2]\AgdaSymbol{;} \AgdaField{lift-β} \AgdaSymbol{=} \AgdaSymbol{λ} \AgdaBound{resp} \AgdaSymbol{→} \AgdaBound{resp} \AgdaSymbol{(}\AgdaFunction{complete} \AgdaSymbol{\_)}\<%
\\
\>[0]\AgdaIndent{2}{}\<[2]%
\>[2]\AgdaSymbol{;} \AgdaField{qind} \<[11]%
\>[11]\AgdaSymbol{=} \AgdaSymbol{λ} \AgdaBound{P} \AgdaBound{\_} \AgdaBound{f} \AgdaBound{\_} \AgdaSymbol{→} \<[25]%
\>[25]\<%
\\
\>[2]\AgdaIndent{11}{}\<[11]%
\>[11]\AgdaFunction{subst} \AgdaBound{P} \AgdaSymbol{(}\AgdaFunction{stable} \AgdaSymbol{\_)} \AgdaSymbol{(}\AgdaBound{f} \AgdaSymbol{(}\AgdaFunction{emb} \AgdaSymbol{\_))}\<%
\\
\>[0]\AgdaIndent{2}{}\<[2]%
\>[2]\AgdaSymbol{\}}\<%
\\
\>[0]\AgdaIndent{2}{}\<[2]%
\>[2]\AgdaKeyword{where}\<%
\\
\>[2]\AgdaIndent{4}{}\<[4]%
\>[4]\AgdaKeyword{open} \AgdaModule{pre-Quotient} \AgdaBound{PQ}\<%
\\
\>[2]\AgdaIndent{4}{}\<[4]%
\>[4]\AgdaKeyword{open} \AgdaModule{def-Quotient} \AgdaBound{QuD}\<%
\\
\>\<\end{code}


\begin{code}\>\<%
\\
\>\AgdaFunction{def-Quotient→Hof-Quotient'} \AgdaSymbol{:} \<[29]%
\>[29]\<%
\\
\>[0]\AgdaIndent{2}{}\<[2]%
\>[2]\AgdaSymbol{\{}\AgdaBound{S} \AgdaSymbol{:} \AgdaFunction{Setoid}\AgdaSymbol{\}\{}\AgdaBound{PQ} \AgdaSymbol{:} \AgdaRecord{pre-Quotient} \AgdaBound{S}\AgdaSymbol{\}}\<%
\\
\>[0]\AgdaIndent{2}{}\<[2]%
\>[2]\AgdaSymbol{→} \AgdaSymbol{(}\AgdaRecord{def-Quotient} \AgdaBound{PQ}\AgdaSymbol{)} \AgdaSymbol{→} \AgdaSymbol{(}\AgdaRecord{Hof-Quotient} \AgdaBound{PQ}\AgdaSymbol{)}\<%
\\
\>\AgdaFunction{def-Quotient→Hof-Quotient'} \AgdaSymbol{=} \<[29]%
\>[29]\<%
\\
\>[0]\AgdaIndent{2}{}\<[2]%
\>[2]\AgdaFunction{Quotient→Hof-Quotient} \AgdaFunction{∘} \AgdaFunction{def-Quotient→Quotient}\<%
\\
\>\<\end{code}
}

\textbf{The propositional univalence implies a quotient is always exact}\label{PUEF}

\begin{code}\>\<%
\\
%
\\
\>\AgdaFunction{Prp} \AgdaSymbol{=} \AgdaPrimitiveType{Set}\<%
\\
%
\\
\>\AgdaFunction{\_⇔\_} \AgdaSymbol{:} \AgdaSymbol{(}\AgdaBound{A} \AgdaBound{B} \AgdaSymbol{:} \AgdaFunction{Prp}\AgdaSymbol{)} \AgdaSymbol{→} \AgdaFunction{Prp}\<%
\\
\>\AgdaBound{A} \AgdaFunction{⇔} \AgdaBound{B} \AgdaSymbol{=} \AgdaSymbol{(}\AgdaBound{A} \AgdaSymbol{→} \AgdaBound{B}\AgdaSymbol{)} \AgdaFunction{×} \AgdaSymbol{(}\AgdaBound{B} \AgdaSymbol{→} \AgdaBound{A}\AgdaSymbol{)}\<%
\\
\>\<\end{code}

\begin{code}\>\<%
\\
\>\AgdaKeyword{module} \AgdaModule{PuImpEff}\<%
\\
%
\\
\>\<\end{code}

Assume we have the propositional univalence (the other part is trivial),

\begin{code}\>\<%
\\
\>[0]\AgdaIndent{2}{}\<[2]%
\>[2]\AgdaSymbol{(}\AgdaBound{PropUni₁} \AgdaSymbol{:} \AgdaSymbol{∀} \AgdaSymbol{\{}\AgdaBound{p} \AgdaBound{q} \AgdaSymbol{:} \AgdaFunction{Prp}\AgdaSymbol{\}} \AgdaSymbol{→} \AgdaSymbol{(}\AgdaBound{p} \AgdaFunction{⇔} \AgdaBound{q}\AgdaSymbol{)} \AgdaSymbol{→} \AgdaBound{p} \AgdaDatatype{≡} \AgdaBound{q}\AgdaSymbol{)}\<%
\\
\>\<\end{code}

and a quotient (Note that we postulate the lifting function for $B : Set_1$ because for convenience we did not take into account the universe levels, but it is required in the proof here)

\begin{code}\>\<%
\\
\>[0]\AgdaIndent{2}{}\<[2]%
\>[2]\AgdaSymbol{\{}\AgdaBound{S} \AgdaSymbol{:} \AgdaFunction{Setoid}\AgdaSymbol{\}}\<%
\\
\>[0]\AgdaIndent{2}{}\<[2]%
\>[2]\AgdaSymbol{\{}\AgdaBound{PQ} \AgdaSymbol{:} \AgdaRecord{pre-Quotient} \AgdaBound{S}\AgdaSymbol{\}}\<%
\\
\>[2]\AgdaIndent{4}{}\<[4]%
\>[4]\AgdaKeyword{where}\<%
\\
\>[0]\AgdaIndent{2}{}\<[2]%
\>[2]\AgdaKeyword{open} \AgdaModule{pre-Quotient} \AgdaBound{PQ}\<%
\\
\>\<\end{code}
\begin{code}\>\>[0]\AgdaIndent{2}{}\<[2]%
\>[2]\<%
\\
\>[0]\AgdaIndent{2}{}\<[2]%
\>[2]\AgdaKeyword{postulate} \<[12]%
\>[12]\<%
\\
\>[2]\AgdaIndent{4}{}\<[4]%
\>[4]\AgdaPostulate{lift₁} \AgdaSymbol{:} \AgdaSymbol{\{}\AgdaBound{B} \AgdaSymbol{:} \AgdaPrimitiveType{Set₁}\AgdaSymbol{\}} \<[24]%
\>[24]\AgdaSymbol{→} \<[26]%
\>[26]\<%
\\
\>[4]\AgdaIndent{12}{}\<[12]%
\>[12]\AgdaSymbol{(}\AgdaBound{f} \AgdaSymbol{:} \AgdaFunction{A} \AgdaSymbol{→} \AgdaBound{B}\AgdaSymbol{)} \AgdaSymbol{→} \<[26]%
\>[26]\<%
\\
\>[4]\AgdaIndent{12}{}\<[12]%
\>[12]\AgdaSymbol{(}\AgdaBound{f} \AgdaFunction{respects} \AgdaFunction{\_\textasciitilde\_}\AgdaSymbol{)} \AgdaSymbol{→} \<[31]%
\>[31]\<%
\\
\>[12]\AgdaIndent{13}{}\<[13]%
\>[13]\AgdaFunction{Q} \AgdaSymbol{→} \AgdaBound{B}\<%
\\
%
\\
\>[-1]\AgdaIndent{2}{}\<[2]%
\>[2]\AgdaKeyword{postulate}\<%
\\
\>[0]\AgdaIndent{4}{}\<[4]%
\>[4]\AgdaPostulate{lift-β₁} \AgdaSymbol{:} \AgdaSymbol{∀} \AgdaSymbol{\{}\AgdaBound{B} \AgdaBound{a} \AgdaBound{f}\AgdaSymbol{\}\{}\AgdaBound{resp} \AgdaSymbol{:} \AgdaSymbol{(}\AgdaBound{f} \AgdaFunction{respects} \AgdaFunction{\_\textasciitilde\_}\AgdaSymbol{)\}} \<[49]%
\>[49]\<%
\\
\>[4]\AgdaIndent{12}{}\<[12]%
\>[12]\AgdaSymbol{→} \AgdaPostulate{lift₁} \AgdaSymbol{\{}\AgdaBound{B}\AgdaSymbol{\}} \AgdaBound{f} \AgdaBound{resp} \AgdaFunction{[} \AgdaBound{a} \AgdaFunction{]} \<[38]%
\>[38]\AgdaDatatype{≡} \AgdaBound{f} \AgdaBound{a}\<%
\\
%
\\
\>[6]\AgdaIndent{2}{}\<[2]%
\>[2]\AgdaFunction{coerce} \AgdaSymbol{:} \AgdaSymbol{\{}\AgdaBound{A} \AgdaBound{B} \AgdaSymbol{:} \AgdaPrimitiveType{Set}\AgdaSymbol{\}} \AgdaSymbol{→} \AgdaBound{A} \AgdaDatatype{≡} \AgdaBound{B} \AgdaSymbol{→} \AgdaBound{A} \AgdaSymbol{→} \AgdaBound{B}\<%
\\
\>[0]\AgdaIndent{2}{}\<[2]%
\>[2]\AgdaFunction{coerce} \AgdaInductiveConstructor{refl} \AgdaBound{m} \AgdaSymbol{=} \AgdaBound{m}\<%
\\
%
\\
\>[0]\AgdaIndent{2}{}\<[2]%
\>[2]\AgdaFunction{exact} \AgdaSymbol{:} \AgdaSymbol{∀} \AgdaBound{a} \AgdaBound{a'} \AgdaSymbol{→} \AgdaFunction{[} \AgdaBound{a} \AgdaFunction{]} \AgdaDatatype{≡} \AgdaFunction{[} \AgdaBound{a'} \AgdaFunction{]} \AgdaSymbol{→} \AgdaBound{a} \AgdaFunction{\textasciitilde} \AgdaBound{a'}\<%
\\
\>[0]\AgdaIndent{2}{}\<[2]%
\>[2]\AgdaFunction{exact} \AgdaBound{a} \AgdaBound{a'} \AgdaBound{p} \AgdaSymbol{=} \AgdaFunction{coerce} \AgdaFunction{P\textasciicircum-β} \AgdaSymbol{(}\AgdaFunction{\textasciitilde-refl} \AgdaSymbol{\{}\AgdaBound{a}\AgdaSymbol{\})}\<%
\\
\>[2]\AgdaIndent{8}{}\<[8]%
\>[8]\AgdaKeyword{where}\<%
\\
\>[8]\AgdaIndent{10}{}\<[10]%
\>[10]\AgdaFunction{P} \AgdaSymbol{:} \AgdaFunction{A} \AgdaSymbol{→} \AgdaFunction{Prp}\<%
\\
\>[8]\AgdaIndent{10}{}\<[10]%
\>[10]\AgdaFunction{P} \AgdaBound{x} \AgdaSymbol{=} \AgdaBound{a} \AgdaFunction{\textasciitilde} \AgdaBound{x}\<%
\\
%
\\
\>[8]\AgdaIndent{10}{}\<[10]%
\>[10]\AgdaFunction{P-resp} \AgdaSymbol{:} \AgdaFunction{P} \AgdaFunction{respects} \AgdaFunction{\_\textasciitilde\_}\<%
\\
\>[8]\AgdaIndent{10}{}\<[10]%
\>[10]\AgdaFunction{P-resp} \AgdaSymbol{\{}\AgdaBound{b}\AgdaSymbol{\}} \AgdaSymbol{\{}\AgdaBound{b'}\AgdaSymbol{\}} \AgdaBound{bb'} \AgdaSymbol{=} \<[32]%
\>[32]\<%
\\
\>[10]\AgdaIndent{12}{}\<[12]%
\>[12]\AgdaBound{PropUni₁} \AgdaSymbol{((λ} \AgdaBound{ab} \AgdaSymbol{→} \AgdaFunction{\textasciitilde-trans} \AgdaBound{ab} \AgdaBound{bb'}\AgdaSymbol{)} \AgdaInductiveConstructor{,} \<[48]%
\>[48]\<%
\\
\>[10]\AgdaIndent{12}{}\<[12]%
\>[12]\AgdaSymbol{(λ} \AgdaBound{ab'} \AgdaSymbol{→} \AgdaFunction{\textasciitilde-trans} \AgdaBound{ab'} \AgdaSymbol{(}\AgdaFunction{\textasciitilde-sym} \AgdaBound{bb'}\AgdaSymbol{)))}\<%
\\
%
\\
\>[0]\AgdaIndent{10}{}\<[10]%
\>[10]\AgdaFunction{P\textasciicircum} \AgdaSymbol{:} \AgdaFunction{Q} \AgdaSymbol{→} \AgdaFunction{Prp}\<%
\\
\>[0]\AgdaIndent{10}{}\<[10]%
\>[10]\AgdaFunction{P\textasciicircum} \AgdaSymbol{=} \AgdaPostulate{lift₁} \AgdaFunction{P} \AgdaFunction{P-resp}\<%
\\
%
\\
\>[0]\AgdaIndent{10}{}\<[10]%
\>[10]\AgdaFunction{P\textasciicircum-β} \AgdaSymbol{:} \AgdaFunction{P} \AgdaBound{a} \AgdaDatatype{≡} \AgdaFunction{P} \AgdaBound{a'}\<%
\\
\>[0]\AgdaIndent{10}{}\<[10]%
\>[10]\AgdaFunction{P\textasciicircum-β} \AgdaSymbol{=} \AgdaFunction{trans} \AgdaSymbol{(}\AgdaFunction{sym} \AgdaPostulate{lift-β₁}\AgdaSymbol{)} \<[37]%
\>[37]\<%
\\
\>[10]\AgdaIndent{17}{}\<[17]%
\>[17]\AgdaSymbol{(}\AgdaFunction{trans} \AgdaSymbol{(}\AgdaFunction{cong} \AgdaFunction{P\textasciicircum} \AgdaBound{p}\AgdaSymbol{)} \AgdaPostulate{lift-β₁}\AgdaSymbol{)}\<%
\\
\>\<\end{code}


\begin{code}\>\<%
\\
%
\\
\>\AgdaKeyword{module} \AgdaModule{GroupoidQuotient} \AgdaKeyword{where}\<%
\\
%
\\
\>\AgdaKeyword{open} \AgdaKeyword{import} \AgdaModule{Relation.Binary.PropositionalEquality} \AgdaSymbol{as} \AgdaModule{PE}\<%
\\
\>[0]\AgdaIndent{2}{}\<[2]%
\>[2]\AgdaKeyword{hiding} \AgdaSymbol{(}[\_]\AgdaSymbol{)}\<%
\\
\>\AgdaKeyword{open} \AgdaKeyword{import} \AgdaModule{Quotient}\<%
\\
\>\AgdaKeyword{open} \AgdaKeyword{import} \AgdaModule{Data.Product}\<%
\\
\>\AgdaKeyword{open} \AgdaModule{≡-Reasoning}\<%
\\
\>\<\end{code}

The equivalence relation is not propositional but h-set.

\begin{code}\>\<%
\\
%
\\
\>\AgdaKeyword{record} \AgdaRecord{Groupoid} \AgdaSymbol{:} \AgdaPrimitiveType{Set₁} \AgdaKeyword{where}\<%
\\
\>[0]\AgdaIndent{2}{}\<[2]%
\>[2]\AgdaKeyword{infix} \AgdaNumber{4} \_\textasciitilde\_\<%
\\
\>[0]\AgdaIndent{2}{}\<[2]%
\>[2]\AgdaKeyword{infixl} \AgdaNumber{5} \_*\_\<%
\\
\>[0]\AgdaIndent{2}{}\<[2]%
\>[2]\AgdaKeyword{field}\<%
\\
\>[2]\AgdaIndent{4}{}\<[4]%
\>[4]\AgdaField{Carrier} \AgdaSymbol{:} \AgdaPrimitiveType{Set}\<%
\\
\>[2]\AgdaIndent{4}{}\<[4]%
\>[4]\AgdaField{\_\textasciitilde\_} \<[12]%
\>[12]\AgdaSymbol{:} \AgdaBound{Carrier} \AgdaSymbol{→} \AgdaBound{Carrier} \AgdaSymbol{→} \AgdaPrimitiveType{Set}\<%
\\
\>[2]\AgdaIndent{4}{}\<[4]%
\>[4]\AgdaField{id} \<[12]%
\>[12]\AgdaSymbol{:} \AgdaSymbol{∀\{}\AgdaBound{a}\AgdaSymbol{\}} \AgdaSymbol{→} \AgdaBound{a} \AgdaBound{\textasciitilde} \AgdaBound{a}\<%
\\
\>[2]\AgdaIndent{4}{}\<[4]%
\>[4]\AgdaField{\_⁻¹} \<[13]%
\>[13]\AgdaSymbol{:} \AgdaSymbol{∀\{}\AgdaBound{a} \AgdaBound{b}\AgdaSymbol{\}} \AgdaSymbol{→} \AgdaBound{a} \AgdaBound{\textasciitilde} \AgdaBound{b} \AgdaSymbol{→} \AgdaBound{b} \AgdaBound{\textasciitilde} \AgdaBound{a}\<%
\\
\>[2]\AgdaIndent{4}{}\<[4]%
\>[4]\AgdaField{\_*\_} \<[12]%
\>[12]\AgdaSymbol{:} \AgdaSymbol{∀\{}\AgdaBound{a} \AgdaBound{b} \AgdaBound{c}\AgdaSymbol{\}} \AgdaSymbol{→} \AgdaBound{a} \AgdaBound{\textasciitilde} \AgdaBound{b} \AgdaSymbol{→} \AgdaBound{b} \AgdaBound{\textasciitilde} \AgdaBound{c} \AgdaSymbol{→} \AgdaBound{a} \AgdaBound{\textasciitilde} \AgdaBound{c}\<%
\\
\>[2]\AgdaIndent{4}{}\<[4]%
\>[4]\AgdaField{id₁} \<[12]%
\>[12]\AgdaSymbol{:} \AgdaSymbol{∀\{}\AgdaBound{a} \AgdaBound{b}\AgdaSymbol{\}\{}\AgdaBound{p} \AgdaSymbol{:} \AgdaBound{a} \AgdaBound{\textasciitilde} \AgdaBound{b}\AgdaSymbol{\}} \AgdaSymbol{→} \AgdaBound{id} \AgdaBound{*} \AgdaBound{p} \AgdaDatatype{≡} \AgdaBound{p}\<%
\\
\>[2]\AgdaIndent{4}{}\<[4]%
\>[4]\AgdaField{id₂} \<[12]%
\>[12]\AgdaSymbol{:} \AgdaSymbol{∀\{}\AgdaBound{a} \AgdaBound{b}\AgdaSymbol{\}\{}\AgdaBound{p} \AgdaSymbol{:} \AgdaBound{a} \AgdaBound{\textasciitilde} \AgdaBound{b}\AgdaSymbol{\}} \AgdaSymbol{→} \AgdaBound{p} \AgdaBound{*} \AgdaBound{id} \AgdaDatatype{≡} \AgdaBound{p}\<%
\\
\>[2]\AgdaIndent{4}{}\<[4]%
\>[4]\AgdaField{inv₁} \<[12]%
\>[12]\AgdaSymbol{:} \AgdaSymbol{∀\{}\AgdaBound{a} \AgdaBound{b}\AgdaSymbol{\}\{}\AgdaBound{p} \AgdaSymbol{:} \AgdaBound{a} \AgdaBound{\textasciitilde} \AgdaBound{b}\AgdaSymbol{\}} \AgdaSymbol{→} \AgdaBound{p} \AgdaBound{*} \AgdaBound{p} \AgdaBound{⁻¹} \AgdaDatatype{≡} \AgdaBound{id}\<%
\\
\>[2]\AgdaIndent{4}{}\<[4]%
\>[4]\AgdaField{inv₂} \<[12]%
\>[12]\AgdaSymbol{:} \AgdaSymbol{∀\{}\AgdaBound{a} \AgdaBound{b}\AgdaSymbol{\}\{}\AgdaBound{p} \AgdaSymbol{:} \AgdaBound{a} \AgdaBound{\textasciitilde} \AgdaBound{b}\AgdaSymbol{\}} \AgdaSymbol{→} \<[35]%
\>[35]\AgdaBound{p} \AgdaBound{⁻¹} \AgdaBound{*} \AgdaBound{p} \AgdaDatatype{≡} \AgdaBound{id}\<%
\\
\>[2]\AgdaIndent{4}{}\<[4]%
\>[4]\AgdaField{assoc} \<[12]%
\>[12]\AgdaSymbol{:} \AgdaSymbol{∀\{}\AgdaBound{a} \AgdaBound{b} \AgdaBound{c} \AgdaBound{d}\AgdaSymbol{\}\{}\AgdaBound{p} \AgdaSymbol{:} \AgdaBound{a} \AgdaBound{\textasciitilde} \AgdaBound{b}\AgdaSymbol{\}\{}\AgdaBound{q} \AgdaSymbol{:} \AgdaBound{b} \AgdaBound{\textasciitilde} \AgdaBound{c}\AgdaSymbol{\}\{}\AgdaBound{r} \AgdaSymbol{:} \AgdaBound{c} \AgdaBound{\textasciitilde} \AgdaBound{d}\AgdaSymbol{\}}\<%
\\
\>[4]\AgdaIndent{12}{}\<[12]%
\>[12]\AgdaSymbol{→} \AgdaSymbol{(}\AgdaBound{p} \AgdaBound{*} \AgdaBound{q}\AgdaSymbol{)} \AgdaBound{*} \AgdaBound{r} \AgdaDatatype{≡} \AgdaBound{p} \AgdaBound{*} \AgdaSymbol{(}\AgdaBound{q} \AgdaBound{*} \AgdaBound{r}\AgdaSymbol{)}\<%
\\
\>[0]\AgdaIndent{4}{}\<[4]%
\>[4]\AgdaField{isGrp} \<[12]%
\>[12]\AgdaSymbol{:} \AgdaSymbol{∀\{}\AgdaBound{a} \AgdaBound{b}\AgdaSymbol{\}\{}\AgdaBound{p} \AgdaBound{q} \AgdaSymbol{:} \AgdaBound{a} \AgdaBound{\textasciitilde} \AgdaBound{b}\AgdaSymbol{\}} \AgdaSymbol{→} \AgdaFunction{isProp} \AgdaSymbol{(}\AgdaBound{p} \AgdaDatatype{≡} \AgdaBound{q}\AgdaSymbol{)}\<%
\\
%
\\
\>[0]\AgdaIndent{2}{}\<[2]%
\>[2]\AgdaFunction{inv₂-id₁} \AgdaSymbol{:} \AgdaSymbol{∀\{}\AgdaBound{a} \AgdaBound{b} \AgdaBound{c}\AgdaSymbol{\}\{}\AgdaBound{p} \AgdaSymbol{:} \AgdaBound{a} \AgdaFunction{\textasciitilde} \AgdaBound{b}\AgdaSymbol{\}\{}\AgdaBound{q} \AgdaSymbol{:} \AgdaBound{b} \AgdaFunction{\textasciitilde} \AgdaBound{c}\AgdaSymbol{\}} \AgdaSymbol{→} \AgdaBound{p} \AgdaFunction{⁻¹} \AgdaFunction{*} \AgdaBound{p} \AgdaFunction{*} \AgdaBound{q} \AgdaDatatype{≡} \AgdaBound{q}\<%
\\
\>[0]\AgdaIndent{2}{}\<[2]%
\>[2]\AgdaFunction{inv₂-id₁} \AgdaSymbol{=} \AgdaFunction{trans} \AgdaSymbol{(}\AgdaFunction{cong} \AgdaSymbol{(λ} \AgdaBound{x} \AgdaSymbol{→} \AgdaBound{x} \AgdaFunction{*} \AgdaSymbol{\_)} \AgdaFunction{inv₂}\AgdaSymbol{)} \AgdaFunction{id₁}\<%
\\
%
\\
\>[0]\AgdaIndent{2}{}\<[2]%
\>[2]\AgdaFunction{inv₁-id₂} \AgdaSymbol{:} \AgdaSymbol{∀\{}\AgdaBound{a} \AgdaBound{b} \AgdaBound{c}\AgdaSymbol{\}\{}\AgdaBound{p} \AgdaSymbol{:} \AgdaBound{a} \AgdaFunction{\textasciitilde} \AgdaBound{b}\AgdaSymbol{\}\{}\AgdaBound{q} \AgdaSymbol{:} \AgdaBound{b} \AgdaFunction{\textasciitilde} \AgdaBound{c}\AgdaSymbol{\}} \AgdaSymbol{→} \AgdaBound{p} \AgdaFunction{*} \AgdaSymbol{(}\AgdaBound{q} \AgdaFunction{*} \AgdaBound{q} \AgdaFunction{⁻¹}\AgdaSymbol{)} \AgdaDatatype{≡} \AgdaBound{p}\<%
\\
\>[0]\AgdaIndent{2}{}\<[2]%
\>[2]\AgdaFunction{inv₁-id₂} \AgdaSymbol{=} \AgdaFunction{trans} \AgdaSymbol{(}\AgdaFunction{cong} \AgdaSymbol{(λ} \AgdaBound{x} \AgdaSymbol{→} \AgdaSymbol{\_} \AgdaFunction{*} \AgdaBound{x}\AgdaSymbol{)} \AgdaFunction{inv₁}\AgdaSymbol{)} \AgdaFunction{id₂}\<%
\\
%
\\
\>[0]\AgdaIndent{2}{}\<[2]%
\>[2]\AgdaFunction{id-sym} \AgdaSymbol{:} \AgdaSymbol{∀} \AgdaSymbol{\{}\AgdaBound{a}\AgdaSymbol{\}} \AgdaSymbol{→} \AgdaFunction{id} \AgdaFunction{⁻¹} \AgdaDatatype{≡} \AgdaFunction{id} \AgdaSymbol{\{}\AgdaBound{a}\AgdaSymbol{\}} \<[34]%
\>[34]\<%
\\
\>[0]\AgdaIndent{2}{}\<[2]%
\>[2]\AgdaFunction{id-sym} \AgdaSymbol{=} \AgdaFunction{begin}\<%
\\
\>[2]\AgdaIndent{14}{}\<[14]%
\>[14]\AgdaFunction{id} \AgdaFunction{⁻¹} \<[27]%
\>[27]\AgdaFunction{≡⟨} \AgdaFunction{sym} \AgdaFunction{id₂} \AgdaFunction{⟩}\<%
\\
\>[2]\AgdaIndent{14}{}\<[14]%
\>[14]\AgdaFunction{id} \AgdaFunction{⁻¹} \AgdaFunction{*} \AgdaFunction{id} \<[27]%
\>[27]\AgdaFunction{≡⟨} \AgdaFunction{inv₂} \AgdaFunction{⟩}\<%
\\
\>[2]\AgdaIndent{14}{}\<[14]%
\>[14]\AgdaFunction{id} \AgdaFunction{∎}\<%
\\
%
\\
\>[0]\AgdaIndent{2}{}\<[2]%
\>[2]\AgdaFunction{sym-comp} \AgdaSymbol{:} \AgdaSymbol{∀} \AgdaSymbol{\{}\AgdaBound{a} \AgdaBound{b} \AgdaBound{c}\AgdaSymbol{\}(}\AgdaBound{p} \AgdaSymbol{:} \AgdaBound{a} \AgdaFunction{\textasciitilde} \AgdaBound{b}\AgdaSymbol{)(}\AgdaBound{q} \AgdaSymbol{:} \AgdaBound{b} \AgdaFunction{\textasciitilde} \AgdaBound{c}\AgdaSymbol{)} \AgdaSymbol{→} \AgdaSymbol{(}\AgdaBound{p} \AgdaFunction{*} \AgdaBound{q}\AgdaSymbol{)} \AgdaFunction{⁻¹} \AgdaDatatype{≡} \AgdaBound{q} \AgdaFunction{⁻¹} \AgdaFunction{*} \AgdaBound{p} \AgdaFunction{⁻¹}\<%
\\
\>[0]\AgdaIndent{2}{}\<[2]%
\>[2]\AgdaFunction{sym-comp} \AgdaBound{p} \AgdaBound{q} \AgdaSymbol{=} \AgdaFunction{begin}\<%
\\
\>[2]\AgdaIndent{14}{}\<[14]%
\>[14]\AgdaSymbol{(}\AgdaBound{p} \AgdaFunction{*} \AgdaBound{q}\AgdaSymbol{)} \AgdaFunction{⁻¹} \<[51]%
\>[51]\AgdaFunction{≡⟨} \AgdaFunction{sym} \AgdaFunction{inv₁-id₂} \AgdaFunction{⟩}\<%
\\
\>[2]\AgdaIndent{14}{}\<[14]%
\>[14]\AgdaSymbol{(}\AgdaBound{p} \AgdaFunction{*} \AgdaBound{q}\AgdaSymbol{)} \AgdaFunction{⁻¹} \AgdaFunction{*} \AgdaSymbol{(}\AgdaBound{p} \AgdaFunction{*} \AgdaBound{p} \AgdaFunction{⁻¹}\AgdaSymbol{)} \<[52]%
\>[52]\AgdaFunction{≡⟨} \AgdaFunction{cong} \AgdaSymbol{(λ} \AgdaBound{x} \AgdaSymbol{→} \AgdaSymbol{(}\AgdaBound{p} \AgdaFunction{*} \AgdaBound{q}\AgdaSymbol{)} \AgdaFunction{⁻¹} \AgdaFunction{*} \AgdaSymbol{(}\AgdaBound{x} \AgdaFunction{*} \AgdaBound{p} \AgdaFunction{⁻¹}\AgdaSymbol{))} \AgdaSymbol{(}\AgdaFunction{sym} \AgdaFunction{inv₁-id₂}\AgdaSymbol{)} \AgdaFunction{⟩}\<%
\\
\>[2]\AgdaIndent{14}{}\<[14]%
\>[14]\AgdaSymbol{(}\AgdaBound{p} \AgdaFunction{*} \AgdaBound{q}\AgdaSymbol{)} \AgdaFunction{⁻¹} \AgdaFunction{*} \AgdaSymbol{(}\AgdaBound{p} \AgdaFunction{*} \AgdaSymbol{(}\AgdaBound{q} \AgdaFunction{*} \AgdaBound{q} \AgdaFunction{⁻¹}\AgdaSymbol{)} \AgdaFunction{*} \AgdaBound{p} \AgdaFunction{⁻¹}\AgdaSymbol{)} \<[53]%
\>[53]\AgdaFunction{≡⟨} \AgdaFunction{cong} \AgdaSymbol{(λ} \AgdaBound{x} \AgdaSymbol{→} \AgdaSymbol{(}\AgdaBound{p} \AgdaFunction{*} \AgdaBound{q}\AgdaSymbol{)} \AgdaFunction{⁻¹} \AgdaFunction{*} \AgdaSymbol{(}\AgdaBound{x} \AgdaFunction{*} \AgdaBound{p} \AgdaFunction{⁻¹}\AgdaSymbol{))} \AgdaSymbol{(}\AgdaFunction{sym} \AgdaFunction{assoc}\AgdaSymbol{)} \AgdaFunction{⟩}\<%
\\
\>[2]\AgdaIndent{14}{}\<[14]%
\>[14]\AgdaSymbol{(}\AgdaBound{p} \AgdaFunction{*} \AgdaBound{q}\AgdaSymbol{)} \AgdaFunction{⁻¹} \AgdaFunction{*} \AgdaSymbol{(}\AgdaBound{p} \AgdaFunction{*} \AgdaBound{q} \AgdaFunction{*} \AgdaBound{q} \AgdaFunction{⁻¹} \AgdaFunction{*} \AgdaBound{p} \AgdaFunction{⁻¹}\AgdaSymbol{)} \<[53]%
\>[53]\AgdaFunction{≡⟨} \AgdaFunction{sym} \AgdaFunction{assoc} \AgdaFunction{⟩}\<%
\\
\>[2]\AgdaIndent{14}{}\<[14]%
\>[14]\AgdaSymbol{(}\AgdaBound{p} \AgdaFunction{*} \AgdaBound{q}\AgdaSymbol{)} \AgdaFunction{⁻¹} \AgdaFunction{*} \AgdaSymbol{(}\AgdaBound{p} \AgdaFunction{*} \AgdaBound{q} \AgdaFunction{*} \AgdaBound{q} \AgdaFunction{⁻¹}\AgdaSymbol{)} \AgdaFunction{*} \AgdaBound{p} \AgdaFunction{⁻¹} \<[53]%
\>[53]\AgdaFunction{≡⟨} \AgdaFunction{cong} \AgdaSymbol{(λ} \AgdaBound{x} \AgdaSymbol{→} \AgdaBound{x} \AgdaFunction{*} \AgdaBound{p} \AgdaFunction{⁻¹}\AgdaSymbol{)} \AgdaSymbol{(}\AgdaFunction{sym} \AgdaFunction{assoc}\AgdaSymbol{)} \AgdaFunction{⟩}\<%
\\
\>[2]\AgdaIndent{14}{}\<[14]%
\>[14]\AgdaSymbol{(}\AgdaBound{p} \AgdaFunction{*} \AgdaBound{q}\AgdaSymbol{)} \AgdaFunction{⁻¹} \AgdaFunction{*} \AgdaSymbol{(}\AgdaBound{p} \AgdaFunction{*} \AgdaBound{q}\AgdaSymbol{)} \AgdaFunction{*} \AgdaBound{q} \AgdaFunction{⁻¹} \AgdaFunction{*} \AgdaBound{p} \AgdaFunction{⁻¹} \<[53]%
\>[53]\AgdaFunction{≡⟨} \AgdaFunction{cong} \AgdaSymbol{(λ} \AgdaBound{x} \AgdaSymbol{→} \AgdaBound{x} \AgdaFunction{*} \AgdaBound{p} \AgdaFunction{⁻¹}\AgdaSymbol{)} \AgdaFunction{inv₂-id₁} \AgdaFunction{⟩}\<%
\\
\>[2]\AgdaIndent{14}{}\<[14]%
\>[14]\AgdaBound{q} \AgdaFunction{⁻¹} \AgdaFunction{*} \AgdaBound{p} \AgdaFunction{⁻¹} \AgdaFunction{∎}\<%
\\
%
\\
\>\AgdaFunction{isGroupoid} \AgdaSymbol{:} \AgdaSymbol{(}\AgdaBound{S} \AgdaSymbol{:} \AgdaPrimitiveType{Set}\AgdaSymbol{)} \AgdaSymbol{→} \AgdaPrimitiveType{Set}\<%
\\
\>\AgdaFunction{isGroupoid} \AgdaBound{S} \AgdaSymbol{=} \AgdaSymbol{\{}\AgdaBound{a} \AgdaBound{b} \AgdaSymbol{:} \AgdaBound{S}\AgdaSymbol{\}} \AgdaSymbol{→} \AgdaFunction{isSet} \AgdaSymbol{(}\AgdaBound{a} \AgdaDatatype{≡} \AgdaBound{b}\AgdaSymbol{)}\<%
\\
%
\\
\>\AgdaKeyword{record} \AgdaRecord{pre-GQuotient} \AgdaSymbol{(}\AgdaBound{G} \AgdaSymbol{:} \AgdaRecord{Groupoid}\AgdaSymbol{)} \AgdaSymbol{:} \AgdaPrimitiveType{Set₁} \AgdaKeyword{where}\<%
\\
\>[0]\AgdaIndent{2}{}\<[2]%
\>[2]\AgdaKeyword{open} \AgdaModule{Groupoid} \AgdaBound{G} \AgdaKeyword{renaming} \AgdaSymbol{(}Carrier \AgdaSymbol{to} A\AgdaSymbol{)}\<%
\\
\>[0]\AgdaIndent{2}{}\<[2]%
\>[2]\AgdaKeyword{field}\<%
\\
\>[2]\AgdaIndent{4}{}\<[4]%
\>[4]\AgdaField{Q} \<[8]%
\>[8]\AgdaSymbol{:} \AgdaPrimitiveType{Set}\<%
\\
\>[2]\AgdaIndent{4}{}\<[4]%
\>[4]\AgdaField{[\_]} \AgdaSymbol{:} \AgdaFunction{A} \AgdaSymbol{→} \AgdaBound{Q}\<%
\\
\>[2]\AgdaIndent{4}{}\<[4]%
\>[4]\AgdaField{[\_]⁼} \AgdaSymbol{:} \AgdaBound{[\_]} \AgdaFunction{respects} \AgdaFunction{\_\textasciitilde\_}\<%
\\
\>[2]\AgdaIndent{4}{}\<[4]%
\>[4]\AgdaField{QisGroupoid} \AgdaSymbol{:} \AgdaFunction{isGroupoid} \AgdaBound{Q}\<%
\\
\>[0]\AgdaIndent{2}{}\<[2]%
\>[2]\AgdaKeyword{open} \AgdaModule{Groupoid} \AgdaBound{G} \AgdaKeyword{public} \AgdaKeyword{renaming} \AgdaSymbol{(}Carrier \AgdaSymbol{to} A\AgdaSymbol{)}\<%
\\
%
\\
\>\AgdaKeyword{record} \AgdaRecord{GQuotient} \AgdaSymbol{\{}\AgdaBound{G} \AgdaSymbol{:} \AgdaRecord{Groupoid}\AgdaSymbol{\}(}\AgdaBound{PGQ} \AgdaSymbol{:} \AgdaRecord{pre-GQuotient} \AgdaBound{G}\AgdaSymbol{)} \AgdaSymbol{:} \AgdaPrimitiveType{Set₁} \AgdaKeyword{where}\<%
\\
\>[0]\AgdaIndent{2}{}\<[2]%
\>[2]\AgdaKeyword{open} \AgdaModule{pre-GQuotient} \AgdaBound{PGQ}\<%
\\
\>[0]\AgdaIndent{2}{}\<[2]%
\>[2]\AgdaKeyword{field}\<%
\\
\>[2]\AgdaIndent{4}{}\<[4]%
\>[4]\AgdaField{qelim} \<[12]%
\>[12]\AgdaSymbol{:} \AgdaSymbol{\{}\AgdaBound{B} \AgdaSymbol{:} \AgdaFunction{Q} \AgdaSymbol{→} \AgdaPrimitiveType{Set}\AgdaSymbol{\}}\<%
\\
\>[4]\AgdaIndent{12}{}\<[12]%
\>[12]\AgdaSymbol{→} \AgdaSymbol{(}\AgdaBound{f} \AgdaSymbol{:} \AgdaSymbol{(}\AgdaBound{a} \AgdaSymbol{:} \AgdaFunction{A}\AgdaSymbol{)} \AgdaSymbol{→} \AgdaBound{B} \AgdaFunction{[} \AgdaBound{a} \AgdaFunction{]}\AgdaSymbol{)}\<%
\\
\>[4]\AgdaIndent{12}{}\<[12]%
\>[12]\AgdaSymbol{→} \AgdaSymbol{(∀} \AgdaSymbol{\{}\AgdaBound{a} \AgdaBound{a'}\AgdaSymbol{\}} \AgdaSymbol{→} \AgdaSymbol{(}\AgdaBound{p} \AgdaSymbol{:} \AgdaBound{a} \AgdaFunction{\textasciitilde} \AgdaBound{a'}\AgdaSymbol{)} \AgdaSymbol{→} \AgdaFunction{subst} \AgdaBound{B} \AgdaFunction{[} \AgdaBound{p} \AgdaFunction{]⁼} \AgdaSymbol{(}\AgdaBound{f} \AgdaBound{a}\AgdaSymbol{)} \AgdaDatatype{≡} \AgdaBound{f} \AgdaBound{a'}\AgdaSymbol{)}\<%
\\
\>[4]\AgdaIndent{12}{}\<[12]%
\>[12]\AgdaSymbol{→} \AgdaSymbol{(}\AgdaBound{q} \AgdaSymbol{:} \AgdaFunction{Q}\AgdaSymbol{)} \AgdaSymbol{→} \AgdaBound{B} \AgdaBound{q}\<%
\\
\>[0]\AgdaIndent{4}{}\<[4]%
\>[4]\AgdaField{qelim-β} \AgdaSymbol{:} \AgdaSymbol{∀} \AgdaSymbol{\{}\AgdaBound{B} \AgdaBound{a} \AgdaBound{f}\AgdaSymbol{\}(}\AgdaBound{resp} \AgdaSymbol{:} \AgdaSymbol{(∀} \AgdaSymbol{\{}\AgdaBound{a} \AgdaBound{a'}\AgdaSymbol{\}} \AgdaSymbol{→} \AgdaSymbol{(}\AgdaBound{p} \AgdaSymbol{:} \AgdaBound{a} \AgdaFunction{\textasciitilde} \AgdaBound{a'}\AgdaSymbol{)} \AgdaSymbol{→} \AgdaFunction{subst} \AgdaBound{B} \AgdaFunction{[} \AgdaBound{p} \AgdaFunction{]⁼} \AgdaSymbol{(}\AgdaBound{f} \AgdaBound{a}\AgdaSymbol{)} \AgdaDatatype{≡} \AgdaBound{f} \AgdaBound{a'}\AgdaSymbol{))}\<%
\\
\>[0]\AgdaIndent{12}{}\<[12]%
\>[12]\AgdaSymbol{→} \AgdaBound{qelim} \AgdaSymbol{\{}\AgdaBound{B}\AgdaSymbol{\}} \AgdaBound{f} \AgdaBound{resp} \AgdaFunction{[} \AgdaBound{a} \AgdaFunction{]} \AgdaDatatype{≡} \AgdaBound{f} \AgdaBound{a}\<%
\\
%
\\
%
\\
%
\\
\>\AgdaKeyword{record} \AgdaRecord{ExactGQuotient} \AgdaSymbol{\{}\AgdaBound{G} \AgdaSymbol{:} \AgdaRecord{Groupoid}\AgdaSymbol{\}\{}\AgdaBound{PGQ} \AgdaSymbol{:} \AgdaRecord{pre-GQuotient} \AgdaBound{G}\AgdaSymbol{\}(}\AgdaBound{GQu} \AgdaSymbol{:} \AgdaRecord{GQuotient} \AgdaBound{PGQ}\AgdaSymbol{)} \AgdaSymbol{:} \AgdaPrimitiveType{Set₁} \AgdaKeyword{where}\<%
\\
\>[0]\AgdaIndent{2}{}\<[2]%
\>[2]\AgdaKeyword{open} \AgdaModule{pre-GQuotient} \AgdaBound{PGQ}\<%
\\
\>[0]\AgdaIndent{2}{}\<[2]%
\>[2]\AgdaKeyword{field}\<%
\\
\>[2]\AgdaIndent{4}{}\<[4]%
\>[4]\AgdaField{exact} \<[11]%
\>[11]\AgdaSymbol{:} \AgdaSymbol{∀\{}\AgdaBound{a} \AgdaBound{b} \AgdaSymbol{:} \AgdaFunction{A}\AgdaSymbol{\}} \AgdaSymbol{→} \AgdaFunction{[} \AgdaBound{a} \AgdaFunction{]} \AgdaDatatype{≡} \AgdaFunction{[} \AgdaBound{b} \AgdaFunction{]} \AgdaSymbol{→} \AgdaBound{a} \AgdaFunction{\textasciitilde} \AgdaBound{b}\<%
\\
\>[2]\AgdaIndent{4}{}\<[4]%
\>[4]\AgdaField{isInv₁} \AgdaSymbol{:} \AgdaSymbol{∀\{}\AgdaBound{a} \AgdaBound{b}\AgdaSymbol{\}\{}\AgdaBound{p} \AgdaSymbol{:} \AgdaBound{a} \AgdaFunction{\textasciitilde} \AgdaBound{b}\AgdaSymbol{\}} \AgdaSymbol{→} \AgdaBound{exact} \AgdaFunction{[} \AgdaBound{p} \AgdaFunction{]⁼} \AgdaDatatype{≡} \AgdaBound{p}\<%
\\
\>[2]\AgdaIndent{4}{}\<[4]%
\>[4]\AgdaField{isInv₂} \AgdaSymbol{:} \AgdaSymbol{∀\{}\AgdaBound{a} \AgdaBound{b}\AgdaSymbol{\}\{}\AgdaBound{p} \AgdaSymbol{:} \AgdaFunction{[} \AgdaBound{a} \AgdaFunction{]} \AgdaDatatype{≡} \AgdaFunction{[} \AgdaBound{b} \AgdaFunction{]}\AgdaSymbol{\}} \AgdaSymbol{→} \AgdaFunction{[} \AgdaBound{exact} \AgdaBound{p} \AgdaFunction{]⁼} \AgdaDatatype{≡} \AgdaBound{p}\<%
\\
%
\\
%
\\
%
\\
\>\AgdaKeyword{module} \AgdaModule{UAImpEff}\<%
\\
\>[0]\AgdaIndent{2}{}\<[2]%
\>[2]\AgdaSymbol{(}\AgdaBound{UA} \AgdaSymbol{:} \AgdaSymbol{∀} \AgdaSymbol{\{}\AgdaBound{p} \AgdaBound{q} \AgdaSymbol{:} \AgdaPrimitiveType{Set}\AgdaSymbol{\}} \AgdaSymbol{→} \AgdaSymbol{(}\AgdaBound{p} \AgdaFunction{⇔} \AgdaBound{q}\AgdaSymbol{)} \AgdaSymbol{→} \AgdaBound{p} \AgdaDatatype{≡} \AgdaBound{q}\AgdaSymbol{)}\<%
\\
\>[0]\AgdaIndent{2}{}\<[2]%
\>[2]\AgdaSymbol{(}\AgdaBound{G} \AgdaSymbol{:} \AgdaRecord{Groupoid}\AgdaSymbol{)}\<%
\\
\>[0]\AgdaIndent{2}{}\<[2]%
\>[2]\AgdaSymbol{(}\AgdaBound{PGQ} \AgdaSymbol{:} \AgdaRecord{pre-GQuotient} \AgdaBound{G}\AgdaSymbol{)}\<%
\\
\>[0]\AgdaIndent{2}{}\<[2]%
\>[2]\AgdaSymbol{(}\AgdaBound{GQu} \AgdaSymbol{:} \AgdaRecord{GQuotient} \AgdaBound{PGQ}\AgdaSymbol{)}\<%
\\
\>[2]\AgdaIndent{4}{}\<[4]%
\>[4]\AgdaKeyword{where}\<%
\\
\>[0]\AgdaIndent{2}{}\<[2]%
\>[2]\AgdaKeyword{open} \AgdaModule{pre-GQuotient} \AgdaBound{PGQ}\<%
\\
\>[0]\AgdaIndent{2}{}\<[2]%
\>[2]\AgdaKeyword{open} \AgdaModule{GQuotient} \AgdaBound{GQu}\<%
\\
\>\<\end{code}

There is also an equivalence on the arrow level $\_\sim_{1}\_$ derived from $\_\sim\_$ (higher levels are trivial)

\begin{code}\>\<%
\\
\>[0]\AgdaIndent{2}{}\<[2]%
\>[2]\AgdaFunction{\_\textasciitilde₁\_} \AgdaSymbol{:} \AgdaSymbol{∀\{}\AgdaBound{a} \AgdaBound{a'} \AgdaBound{b} \AgdaBound{b'}\AgdaSymbol{\}} \AgdaSymbol{→} \AgdaBound{a} \AgdaFunction{\textasciitilde} \AgdaBound{a'} \AgdaSymbol{→} \AgdaBound{b} \AgdaFunction{\textasciitilde} \AgdaBound{b'} \AgdaSymbol{→} \AgdaSymbol{(}\AgdaBound{a} \AgdaFunction{\textasciitilde} \AgdaBound{b} \AgdaSymbol{→} \AgdaBound{a'} \AgdaFunction{\textasciitilde} \AgdaBound{b'}\AgdaSymbol{)}\<%
\\
\>[0]\AgdaIndent{2}{}\<[2]%
\>[2]\AgdaFunction{\_\textasciitilde₁\_} \AgdaBound{p} \AgdaBound{q} \AgdaBound{r} \AgdaSymbol{=} \AgdaBound{p} \AgdaFunction{⁻¹} \AgdaFunction{*} \AgdaBound{r} \AgdaFunction{*} \AgdaBound{q}\<%
\\
%
\\
\>[0]\AgdaIndent{2}{}\<[2]%
\>[2]\AgdaFunction{\_\textasciitilde₁I\_} \AgdaSymbol{:} \AgdaSymbol{∀\{}\AgdaBound{a} \AgdaBound{a'} \AgdaBound{b} \AgdaBound{b'}\AgdaSymbol{\}} \AgdaSymbol{→} \AgdaBound{a} \AgdaFunction{\textasciitilde} \AgdaBound{a'} \AgdaSymbol{→} \AgdaBound{b} \AgdaFunction{\textasciitilde} \AgdaBound{b'} \AgdaSymbol{→} \AgdaSymbol{(}\AgdaBound{a'} \AgdaFunction{\textasciitilde} \AgdaBound{b'} \AgdaSymbol{→} \AgdaBound{a} \AgdaFunction{\textasciitilde} \AgdaBound{b}\AgdaSymbol{)}\<%
\\
\>[0]\AgdaIndent{2}{}\<[2]%
\>[2]\AgdaFunction{\_\textasciitilde₁I\_} \AgdaBound{p} \AgdaBound{q} \AgdaBound{r} \AgdaSymbol{=} \AgdaBound{p} \AgdaFunction{*} \AgdaBound{r} \AgdaFunction{*} \AgdaBound{q} \AgdaFunction{⁻¹}\<%
\\
%
\\
\>[0]\AgdaIndent{2}{}\<[2]%
\>[2]\AgdaFunction{\textasciitilde₁-equiv₁} \AgdaSymbol{:} \AgdaSymbol{∀\{}\AgdaBound{a} \AgdaBound{a'} \AgdaBound{b} \AgdaBound{b'}\AgdaSymbol{\}(}\AgdaBound{p} \AgdaSymbol{:} \AgdaBound{a} \AgdaFunction{\textasciitilde} \AgdaBound{a'}\AgdaSymbol{)(}\AgdaBound{q} \AgdaSymbol{:} \AgdaBound{b} \AgdaFunction{\textasciitilde} \AgdaBound{b'}\AgdaSymbol{)(}\AgdaBound{r} \AgdaSymbol{:} \AgdaBound{a} \AgdaFunction{\textasciitilde} \AgdaBound{b}\AgdaSymbol{)}\<%
\\
\>[2]\AgdaIndent{11}{}\<[11]%
\>[11]\AgdaSymbol{→} \AgdaSymbol{(}\AgdaBound{p} \AgdaFunction{\textasciitilde₁I} \AgdaBound{q}\AgdaSymbol{)} \AgdaSymbol{((}\AgdaBound{p} \AgdaFunction{\textasciitilde₁} \AgdaBound{q}\AgdaSymbol{)} \AgdaBound{r}\AgdaSymbol{)} \AgdaDatatype{≡} \AgdaBound{r}\<%
\\
\>[0]\AgdaIndent{2}{}\<[2]%
\>[2]\AgdaFunction{\textasciitilde₁-equiv₁} \AgdaBound{p} \AgdaBound{q} \AgdaBound{r} \AgdaSymbol{=} \AgdaFunction{begin}\<%
\\
\>[0]\AgdaIndent{13}{}\<[13]%
\>[13]\AgdaSymbol{(}\AgdaBound{p} \AgdaFunction{*} \AgdaSymbol{((}\AgdaBound{p} \AgdaFunction{⁻¹} \AgdaFunction{*} \AgdaBound{r}\AgdaSymbol{)} \AgdaFunction{*} \AgdaBound{q}\AgdaSymbol{))} \AgdaFunction{*} \AgdaBound{q} \AgdaFunction{⁻¹} \AgdaFunction{≡⟨} \AgdaFunction{cong} \AgdaSymbol{(λ} \AgdaBound{x} \AgdaSymbol{→} \AgdaBound{x} \AgdaFunction{*} \AgdaBound{q} \AgdaFunction{⁻¹}\AgdaSymbol{)} \AgdaSymbol{(}\AgdaFunction{sym} \AgdaFunction{assoc}\AgdaSymbol{)} \AgdaFunction{⟩}\<%
\\
\>[0]\AgdaIndent{13}{}\<[13]%
\>[13]\AgdaSymbol{((}\AgdaBound{p} \AgdaFunction{*} \AgdaSymbol{(}\AgdaBound{p} \AgdaFunction{⁻¹} \AgdaFunction{*} \AgdaBound{r}\AgdaSymbol{))} \AgdaFunction{*} \AgdaBound{q}\AgdaSymbol{)} \AgdaFunction{*} \AgdaBound{q} \AgdaFunction{⁻¹} \AgdaFunction{≡⟨} \AgdaFunction{cong} \AgdaSymbol{(λ} \AgdaBound{x} \AgdaSymbol{→} \AgdaBound{x} \AgdaFunction{*} \AgdaBound{q} \AgdaFunction{*} \AgdaBound{q} \AgdaFunction{⁻¹}\AgdaSymbol{)} \AgdaSymbol{(}\AgdaFunction{sym} \AgdaFunction{assoc}\AgdaSymbol{)} \AgdaFunction{⟩}\<%
\\
\>[0]\AgdaIndent{13}{}\<[13]%
\>[13]\AgdaSymbol{(((}\AgdaBound{p} \AgdaFunction{*} \AgdaBound{p} \AgdaFunction{⁻¹}\AgdaSymbol{)} \AgdaFunction{*} \AgdaBound{r}\AgdaSymbol{)} \AgdaFunction{*} \AgdaBound{q}\AgdaSymbol{)} \AgdaFunction{*} \AgdaBound{q} \AgdaFunction{⁻¹} \AgdaFunction{≡⟨} \AgdaFunction{cong} \AgdaSymbol{(λ} \AgdaBound{x} \AgdaSymbol{→} \AgdaBound{x} \AgdaFunction{*} \AgdaBound{q} \AgdaFunction{*} \AgdaBound{q} \AgdaFunction{⁻¹}\AgdaSymbol{)} \AgdaSymbol{(}\AgdaFunction{trans} \AgdaSymbol{(}\AgdaFunction{cong} \AgdaSymbol{(λ} \AgdaBound{x} \AgdaSymbol{→} \AgdaBound{x} \AgdaFunction{*} \AgdaSymbol{\_)} \AgdaFunction{inv₁}\AgdaSymbol{)} \AgdaFunction{id₁}\AgdaSymbol{)} \AgdaFunction{⟩}\<%
\\
\>[0]\AgdaIndent{13}{}\<[13]%
\>[13]\AgdaSymbol{(}\AgdaBound{r} \AgdaFunction{*} \AgdaBound{q}\AgdaSymbol{)} \AgdaFunction{*} \AgdaBound{q} \AgdaFunction{⁻¹} \<[42]%
\>[42]\AgdaFunction{≡⟨} \AgdaFunction{assoc} \AgdaFunction{⟩}\<%
\\
\>[0]\AgdaIndent{13}{}\<[13]%
\>[13]\AgdaBound{r} \AgdaFunction{*} \AgdaSymbol{(}\AgdaBound{q} \AgdaFunction{*} \AgdaBound{q} \AgdaFunction{⁻¹}\AgdaSymbol{)} \<[42]%
\>[42]\AgdaFunction{≡⟨} \AgdaFunction{trans} \AgdaSymbol{(}\AgdaFunction{cong} \AgdaSymbol{(λ} \AgdaBound{x} \AgdaSymbol{→} \AgdaBound{r} \AgdaFunction{*} \AgdaBound{x}\AgdaSymbol{)} \AgdaFunction{inv₁}\AgdaSymbol{)} \AgdaFunction{id₂} \AgdaFunction{⟩}\<%
\\
\>[0]\AgdaIndent{13}{}\<[13]%
\>[13]\AgdaBound{r} \AgdaFunction{∎}\<%
\\
%
\\
\>[0]\AgdaIndent{2}{}\<[2]%
\>[2]\AgdaFunction{\textasciitilde₁-equiv₂} \AgdaSymbol{:} \AgdaSymbol{∀\{}\AgdaBound{a} \AgdaBound{a'} \AgdaBound{b} \AgdaBound{b'}\AgdaSymbol{\}(}\AgdaBound{p} \AgdaSymbol{:} \AgdaBound{a} \AgdaFunction{\textasciitilde} \AgdaBound{a'}\AgdaSymbol{)(}\AgdaBound{q} \AgdaSymbol{:} \AgdaBound{b} \AgdaFunction{\textasciitilde} \AgdaBound{b'}\AgdaSymbol{)(}\AgdaBound{r} \AgdaSymbol{:} \AgdaBound{a'} \AgdaFunction{\textasciitilde} \AgdaBound{b'}\AgdaSymbol{)}\<%
\\
\>[0]\AgdaIndent{11}{}\<[11]%
\>[11]\AgdaSymbol{→} \AgdaSymbol{(}\AgdaBound{p} \AgdaFunction{\textasciitilde₁} \AgdaBound{q}\AgdaSymbol{)} \AgdaSymbol{((}\AgdaBound{p} \AgdaFunction{\textasciitilde₁I} \AgdaBound{q}\AgdaSymbol{)} \AgdaBound{r}\AgdaSymbol{)} \AgdaDatatype{≡} \AgdaBound{r}\<%
\\
\>[0]\AgdaIndent{2}{}\<[2]%
\>[2]\AgdaFunction{\textasciitilde₁-equiv₂} \AgdaBound{p} \AgdaBound{q} \AgdaBound{r} \AgdaSymbol{=} \AgdaFunction{begin}\<%
\\
\>[0]\AgdaIndent{13}{}\<[13]%
\>[13]\AgdaSymbol{(}\AgdaBound{p} \AgdaFunction{⁻¹} \AgdaFunction{*} \AgdaSymbol{((}\AgdaBound{p} \AgdaFunction{*} \AgdaBound{r}\AgdaSymbol{)} \AgdaFunction{*} \AgdaBound{q} \AgdaFunction{⁻¹}\AgdaSymbol{))} \AgdaFunction{*} \AgdaBound{q} \AgdaFunction{≡⟨} \AgdaFunction{cong} \AgdaSymbol{(λ} \AgdaBound{x} \AgdaSymbol{→} \AgdaBound{x} \AgdaFunction{*} \AgdaBound{q}\AgdaSymbol{)} \AgdaSymbol{(}\AgdaFunction{sym} \AgdaFunction{assoc}\AgdaSymbol{)} \AgdaFunction{⟩}\<%
\\
\>[0]\AgdaIndent{13}{}\<[13]%
\>[13]\AgdaSymbol{((}\AgdaBound{p} \AgdaFunction{⁻¹} \AgdaFunction{*} \AgdaSymbol{(}\AgdaBound{p} \AgdaFunction{*} \AgdaBound{r}\AgdaSymbol{))} \AgdaFunction{*} \AgdaBound{q} \AgdaFunction{⁻¹}\AgdaSymbol{)} \AgdaFunction{*} \AgdaBound{q} \AgdaFunction{≡⟨} \AgdaFunction{cong} \AgdaSymbol{(λ} \AgdaBound{x} \AgdaSymbol{→} \AgdaBound{x} \AgdaFunction{*} \AgdaBound{q} \<[65]%
\>[65]\AgdaFunction{⁻¹} \AgdaFunction{*} \AgdaBound{q}\AgdaSymbol{)} \AgdaSymbol{(}\AgdaFunction{sym} \AgdaFunction{assoc}\AgdaSymbol{)} \AgdaFunction{⟩}\<%
\\
\>[0]\AgdaIndent{13}{}\<[13]%
\>[13]\AgdaSymbol{(((}\AgdaBound{p} \AgdaFunction{⁻¹} \AgdaFunction{*} \AgdaBound{p}\AgdaSymbol{)} \AgdaFunction{*} \AgdaBound{r}\AgdaSymbol{)} \AgdaFunction{*} \AgdaBound{q} \AgdaFunction{⁻¹}\AgdaSymbol{)} \AgdaFunction{*} \AgdaBound{q} \AgdaFunction{≡⟨} \AgdaFunction{cong} \AgdaSymbol{(λ} \AgdaBound{x} \AgdaSymbol{→} \AgdaBound{x} \AgdaFunction{*} \AgdaBound{q} \AgdaFunction{⁻¹} \AgdaFunction{*} \AgdaBound{q}\AgdaSymbol{)} \AgdaSymbol{(}\AgdaFunction{trans} \AgdaSymbol{(}\AgdaFunction{cong} \AgdaSymbol{(λ} \AgdaBound{x} \AgdaSymbol{→} \AgdaBound{x} \AgdaFunction{*} \AgdaSymbol{\_)} \AgdaFunction{inv₂}\AgdaSymbol{)} \AgdaFunction{id₁}\AgdaSymbol{)} \AgdaFunction{⟩}\<%
\\
\>[0]\AgdaIndent{13}{}\<[13]%
\>[13]\AgdaSymbol{(}\AgdaBound{r} \AgdaFunction{*} \AgdaBound{q} \AgdaFunction{⁻¹}\AgdaSymbol{)} \AgdaFunction{*} \AgdaBound{q} \<[42]%
\>[42]\AgdaFunction{≡⟨} \AgdaFunction{assoc} \AgdaFunction{⟩}\<%
\\
\>[0]\AgdaIndent{13}{}\<[13]%
\>[13]\AgdaBound{r} \AgdaFunction{*} \AgdaSymbol{(}\AgdaBound{q} \AgdaFunction{⁻¹} \AgdaFunction{*} \AgdaBound{q}\AgdaSymbol{)} \<[42]%
\>[42]\AgdaFunction{≡⟨} \AgdaFunction{trans} \AgdaSymbol{(}\AgdaFunction{cong} \AgdaSymbol{(λ} \AgdaBound{x} \AgdaSymbol{→} \AgdaBound{r} \AgdaFunction{*} \AgdaBound{x}\AgdaSymbol{)} \AgdaFunction{inv₂}\AgdaSymbol{)} \AgdaFunction{id₂} \AgdaFunction{⟩}\<%
\\
\>[0]\AgdaIndent{13}{}\<[13]%
\>[13]\AgdaBound{r} \AgdaFunction{∎}\<%
\\
%
\\
\>[0]\AgdaIndent{2}{}\<[2]%
\>[2]\AgdaFunction{\textasciitilde₁-sound} \AgdaSymbol{:} \AgdaSymbol{∀\{}\AgdaBound{a} \AgdaBound{a'} \AgdaBound{b} \AgdaBound{b'}\AgdaSymbol{\}} \AgdaSymbol{→} \AgdaBound{a} \AgdaFunction{\textasciitilde} \AgdaBound{a'} \AgdaSymbol{→} \AgdaBound{b} \AgdaFunction{\textasciitilde} \AgdaBound{b'} \AgdaSymbol{→} \AgdaSymbol{(}\AgdaBound{a} \AgdaFunction{\textasciitilde} \AgdaBound{b}\AgdaSymbol{)} \AgdaDatatype{≡} \AgdaSymbol{(}\AgdaBound{a'} \AgdaFunction{\textasciitilde} \AgdaBound{b'}\AgdaSymbol{)}\<%
\\
\>[0]\AgdaIndent{2}{}\<[2]%
\>[2]\AgdaFunction{\textasciitilde₁-sound} \AgdaBound{p} \AgdaBound{q} \AgdaSymbol{=} \AgdaBound{UA} \AgdaSymbol{((}\AgdaBound{p} \AgdaFunction{\textasciitilde₁} \AgdaBound{q}\AgdaSymbol{)} \AgdaInductiveConstructor{,} \AgdaSymbol{(}\AgdaBound{p} \AgdaFunction{\textasciitilde₁I} \AgdaBound{q}\AgdaSymbol{))}\<%
\\
\>\<\end{code}

Functorial property

\begin{code}\>\<%
\\
\>[0]\AgdaIndent{2}{}\<[2]%
\>[2]\AgdaFunction{\_\textasciitilde₁\_-id} \AgdaSymbol{:} \AgdaSymbol{∀} \AgdaSymbol{\{}\AgdaBound{a}\AgdaSymbol{\}(}\AgdaBound{r} \AgdaSymbol{:} \AgdaBound{a} \AgdaFunction{\textasciitilde} \AgdaBound{a}\AgdaSymbol{)} \AgdaSymbol{→} \AgdaSymbol{(}\AgdaFunction{id} \AgdaFunction{\textasciitilde₁} \AgdaFunction{id}\AgdaSymbol{)} \AgdaBound{r} \AgdaDatatype{≡} \AgdaBound{r}\<%
\\
\>[0]\AgdaIndent{2}{}\<[2]%
\>[2]\AgdaFunction{\_\textasciitilde₁\_-id} \AgdaBound{r} \AgdaSymbol{=} \AgdaFunction{begin}\<%
\\
\>[2]\AgdaIndent{10}{}\<[10]%
\>[10]\AgdaFunction{id} \AgdaFunction{⁻¹} \AgdaFunction{*} \AgdaBound{r} \AgdaFunction{*} \AgdaFunction{id} \AgdaFunction{≡⟨} \AgdaFunction{id₂} \AgdaFunction{⟩}\<%
\\
\>[2]\AgdaIndent{10}{}\<[10]%
\>[10]\AgdaFunction{id} \AgdaFunction{⁻¹} \AgdaFunction{*} \AgdaBound{r} \<[25]%
\>[25]\AgdaFunction{≡⟨} \AgdaFunction{cong} \AgdaSymbol{(λ} \AgdaBound{x} \AgdaSymbol{→} \AgdaBound{x} \AgdaFunction{*} \AgdaBound{r}\AgdaSymbol{)} \AgdaFunction{id-sym} \AgdaFunction{⟩}\<%
\\
\>[2]\AgdaIndent{10}{}\<[10]%
\>[10]\AgdaFunction{id} \AgdaFunction{*} \AgdaBound{r} \<[25]%
\>[25]\AgdaFunction{≡⟨} \AgdaFunction{id₁} \AgdaFunction{⟩}\<%
\\
\>[2]\AgdaIndent{10}{}\<[10]%
\>[10]\AgdaBound{r} \AgdaFunction{∎}\<%
\\
%
\\
\>[0]\AgdaIndent{2}{}\<[2]%
\>[2]\AgdaFunction{\_\textasciitilde₁\_-comp} \AgdaSymbol{:} \AgdaSymbol{∀} \AgdaSymbol{\{}\AgdaBound{a} \AgdaBound{a'} \AgdaBound{a''} \AgdaBound{b} \AgdaBound{b'} \AgdaBound{b''}\AgdaSymbol{\}(}\AgdaBound{p} \AgdaSymbol{:} \AgdaBound{a} \AgdaFunction{\textasciitilde} \AgdaBound{a'}\AgdaSymbol{)(}\AgdaBound{q} \AgdaSymbol{:} \AgdaBound{b} \AgdaFunction{\textasciitilde} \AgdaBound{b'}\AgdaSymbol{)}\<%
\\
\>[0]\AgdaIndent{13}{}\<[13]%
\>[13]\AgdaSymbol{(}\AgdaBound{p'} \AgdaSymbol{:} \AgdaBound{a'} \AgdaFunction{\textasciitilde} \AgdaBound{a''}\AgdaSymbol{)(}\AgdaBound{q'} \AgdaSymbol{:} \AgdaBound{b'} \AgdaFunction{\textasciitilde} \AgdaBound{b''}\AgdaSymbol{)(}\AgdaBound{r} \AgdaSymbol{:} \AgdaBound{a} \AgdaFunction{\textasciitilde} \AgdaBound{b}\AgdaSymbol{)}\<%
\\
\>[0]\AgdaIndent{11}{}\<[11]%
\>[11]\AgdaSymbol{→} \AgdaSymbol{((}\AgdaBound{p} \AgdaFunction{*} \AgdaBound{p'}\AgdaSymbol{)} \AgdaFunction{\textasciitilde₁} \AgdaSymbol{(}\AgdaBound{q} \AgdaFunction{*} \AgdaBound{q'}\AgdaSymbol{))} \AgdaBound{r} \AgdaDatatype{≡} \AgdaSymbol{(}\AgdaBound{p'} \AgdaFunction{\textasciitilde₁} \AgdaBound{q'}\AgdaSymbol{)} \AgdaSymbol{((}\AgdaBound{p} \AgdaFunction{\textasciitilde₁} \AgdaBound{q}\AgdaSymbol{)} \AgdaBound{r}\AgdaSymbol{)}\<%
\\
\>[0]\AgdaIndent{2}{}\<[2]%
\>[2]\AgdaFunction{\_\textasciitilde₁\_-comp} \AgdaBound{p} \AgdaBound{q} \AgdaBound{p'} \AgdaBound{q'} \AgdaBound{r} \AgdaSymbol{=} \AgdaFunction{begin}\<%
\\
\>[2]\AgdaIndent{10}{}\<[10]%
\>[10]\AgdaSymbol{(}\AgdaBound{p} \AgdaFunction{*} \AgdaBound{p'}\AgdaSymbol{)} \AgdaFunction{⁻¹} \AgdaFunction{*} \AgdaBound{r} \AgdaFunction{*} \AgdaSymbol{(}\AgdaBound{q} \AgdaFunction{*} \AgdaBound{q'}\AgdaSymbol{)} \<[39]%
\>[39]\AgdaFunction{≡⟨} \AgdaFunction{cong} \AgdaSymbol{(λ} \AgdaBound{x} \AgdaSymbol{→} \AgdaBound{x} \AgdaFunction{*} \AgdaBound{r} \AgdaFunction{*} \AgdaSymbol{(}\AgdaBound{q} \AgdaFunction{*} \AgdaBound{q'}\AgdaSymbol{))} \AgdaSymbol{(}\AgdaFunction{sym-comp} \AgdaSymbol{\_} \AgdaSymbol{\_)} \AgdaFunction{⟩}\<%
\\
\>[2]\AgdaIndent{10}{}\<[10]%
\>[10]\AgdaSymbol{(}\AgdaBound{p'} \AgdaFunction{⁻¹} \AgdaFunction{*} \AgdaBound{p} \AgdaFunction{⁻¹}\AgdaSymbol{)} \AgdaFunction{*} \AgdaBound{r} \AgdaFunction{*} \AgdaSymbol{(}\AgdaBound{q} \AgdaFunction{*} \AgdaBound{q'}\AgdaSymbol{)} \AgdaFunction{≡⟨} \AgdaFunction{cong} \AgdaSymbol{(λ} \AgdaBound{x} \AgdaSymbol{→} \AgdaBound{x} \AgdaFunction{*} \AgdaSymbol{(}\AgdaBound{q} \AgdaFunction{*} \AgdaBound{q'}\AgdaSymbol{))} \AgdaFunction{assoc} \AgdaFunction{⟩}\<%
\\
\>[2]\AgdaIndent{10}{}\<[10]%
\>[10]\AgdaBound{p'} \AgdaFunction{⁻¹} \AgdaFunction{*} \AgdaSymbol{(}\AgdaBound{p} \AgdaFunction{⁻¹} \AgdaFunction{*} \AgdaBound{r}\AgdaSymbol{)} \AgdaFunction{*} \AgdaSymbol{(}\AgdaBound{q} \AgdaFunction{*} \AgdaBound{q'}\AgdaSymbol{)} \AgdaFunction{≡⟨} \AgdaFunction{sym} \AgdaFunction{assoc} \AgdaFunction{⟩}\<%
\\
\>[2]\AgdaIndent{10}{}\<[10]%
\>[10]\AgdaBound{p'} \AgdaFunction{⁻¹} \AgdaFunction{*} \AgdaSymbol{(}\AgdaBound{p} \AgdaFunction{⁻¹} \AgdaFunction{*} \AgdaBound{r}\AgdaSymbol{)} \AgdaFunction{*} \AgdaBound{q} \AgdaFunction{*} \AgdaBound{q'} \<[40]%
\>[40]\AgdaFunction{≡⟨} \AgdaFunction{cong} \AgdaSymbol{(λ} \AgdaBound{x} \AgdaSymbol{→} \AgdaBound{x} \AgdaFunction{*} \AgdaBound{q'}\AgdaSymbol{)} \AgdaFunction{assoc} \AgdaFunction{⟩}\<%
\\
\>[2]\AgdaIndent{10}{}\<[10]%
\>[10]\AgdaBound{p'} \AgdaFunction{⁻¹} \AgdaFunction{*} \AgdaSymbol{(}\AgdaBound{p} \AgdaFunction{⁻¹} \AgdaFunction{*} \AgdaBound{r} \AgdaFunction{*} \AgdaBound{q}\AgdaSymbol{)} \AgdaFunction{*} \AgdaBound{q'} \<[40]%
\>[40]\AgdaFunction{∎}\<%
\\
%
\\
%
\\
\>[0]\AgdaIndent{2}{}\<[2]%
\>[2]\AgdaKeyword{postulate} \<[12]%
\>[12]\<%
\\
\>[0]\AgdaIndent{4}{}\<[4]%
\>[4]\AgdaPostulate{lift₁} \AgdaSymbol{:} \AgdaSymbol{\{} \AgdaBound{B} \AgdaSymbol{:} \AgdaPrimitiveType{Set₁}\AgdaSymbol{\}} \<[24]%
\>[24]\<%
\\
\>[4]\AgdaIndent{12}{}\<[12]%
\>[12]\AgdaSymbol{(}\AgdaBound{f} \AgdaSymbol{:} \AgdaFunction{A} \AgdaSymbol{→} \AgdaBound{B}\AgdaSymbol{)}\<%
\\
\>[4]\AgdaIndent{12}{}\<[12]%
\>[12]\AgdaSymbol{(}\AgdaBound{f-resp} \AgdaSymbol{:} \AgdaBound{f} \AgdaFunction{respects} \AgdaFunction{\_\textasciitilde\_}\AgdaSymbol{)}\<%
\\
\>[-2]\AgdaIndent{10}{}\<[10]%
\>[10]\AgdaSymbol{→} \AgdaFunction{Q} \AgdaSymbol{→} \AgdaBound{B}\<%
\\
%
\\
\>[0]\AgdaIndent{2}{}\<[2]%
\>[2]\AgdaKeyword{postulate}\<%
\\
\>[2]\AgdaIndent{4}{}\<[4]%
\>[4]\AgdaPostulate{lift-β₁} \AgdaSymbol{:} \AgdaSymbol{∀} \AgdaSymbol{\{}\AgdaBound{B} \AgdaBound{a} \AgdaBound{f}\AgdaSymbol{\}}\<%
\\
\>[4]\AgdaIndent{14}{}\<[14]%
\>[14]\AgdaSymbol{\{}\AgdaBound{f-resp} \AgdaSymbol{:} \AgdaBound{f} \AgdaFunction{respects} \AgdaFunction{\_\textasciitilde\_}\AgdaSymbol{\}}\<%
\\
\>[-2]\AgdaIndent{12}{}\<[12]%
\>[12]\AgdaSymbol{→} \AgdaPostulate{lift₁} \AgdaSymbol{\{}\AgdaBound{B}\AgdaSymbol{\}} \AgdaBound{f} \AgdaBound{f-resp} \AgdaFunction{[} \AgdaBound{a} \AgdaFunction{]} \<[40]%
\>[40]\AgdaDatatype{≡} \AgdaBound{f} \AgdaBound{a}\<%
\\
%
\\
\>[0]\AgdaIndent{2}{}\<[2]%
\>[2]\AgdaFunction{exact} \AgdaSymbol{:} \AgdaSymbol{∀} \AgdaSymbol{\{}\AgdaBound{a} \AgdaBound{a'}\AgdaSymbol{\}} \AgdaSymbol{→} \AgdaFunction{[} \AgdaBound{a} \AgdaFunction{]} \AgdaDatatype{≡} \AgdaFunction{[} \AgdaBound{a'} \AgdaFunction{]} \AgdaSymbol{→} \AgdaBound{a} \AgdaFunction{\textasciitilde} \AgdaBound{a'}\<%
\\
\>[0]\AgdaIndent{2}{}\<[2]%
\>[2]\AgdaFunction{exact} \AgdaSymbol{\{}\AgdaBound{a}\AgdaSymbol{\}} \AgdaSymbol{\{}\AgdaBound{a'}\AgdaSymbol{\}} \AgdaBound{eq} \AgdaSymbol{=} \AgdaFunction{coerce} \AgdaFunction{P\textasciicircum-β} \AgdaSymbol{(}\AgdaFunction{id} \AgdaSymbol{\{}\AgdaBound{a}\AgdaSymbol{\})}\<%
\\
\>[2]\AgdaIndent{8}{}\<[8]%
\>[8]\AgdaKeyword{where}\<%
\\
\>[8]\AgdaIndent{10}{}\<[10]%
\>[10]\AgdaFunction{P} \AgdaSymbol{:} \AgdaFunction{A} \AgdaSymbol{→} \AgdaPrimitiveType{Set}\<%
\\
\>[8]\AgdaIndent{10}{}\<[10]%
\>[10]\AgdaFunction{P} \AgdaBound{x} \AgdaSymbol{=} \AgdaBound{a} \AgdaFunction{\textasciitilde} \AgdaBound{x}\<%
\\
\>\<\end{code}
Functorial
\begin{code}\>\<%
\\
\>[8]\AgdaIndent{10}{}\<[10]%
\>[10]\AgdaFunction{P₁} \AgdaSymbol{:} \AgdaSymbol{∀\{}\AgdaBound{b} \AgdaBound{c}\AgdaSymbol{\}} \AgdaSymbol{→} \AgdaBound{b} \AgdaFunction{\textasciitilde} \AgdaBound{c} \AgdaSymbol{→} \AgdaSymbol{(}\AgdaFunction{P} \AgdaBound{b} \AgdaSymbol{→} \AgdaFunction{P} \AgdaBound{c}\AgdaSymbol{)}\<%
\\
\>[8]\AgdaIndent{10}{}\<[10]%
\>[10]\AgdaFunction{P₁} \AgdaBound{bc} \AgdaBound{ab} \AgdaSymbol{=} \AgdaBound{ab} \AgdaFunction{*} \AgdaBound{bc}\<%
\\
%
\\
\>[8]\AgdaIndent{10}{}\<[10]%
\>[10]\AgdaFunction{P₁-id} \AgdaSymbol{:} \AgdaSymbol{∀\{}\AgdaBound{b}\AgdaSymbol{\}\{}\AgdaBound{r} \AgdaSymbol{:} \AgdaFunction{P} \AgdaBound{b}\AgdaSymbol{\}} \AgdaSymbol{→} \AgdaFunction{P₁} \AgdaFunction{id} \AgdaBound{r} \AgdaDatatype{≡} \AgdaBound{r}\<%
\\
\>[8]\AgdaIndent{10}{}\<[10]%
\>[10]\AgdaFunction{P₁-id} \AgdaSymbol{=} \AgdaFunction{id₂}\<%
\\
%
\\
\>[8]\AgdaIndent{10}{}\<[10]%
\>[10]\AgdaFunction{P₁-comp} \AgdaSymbol{:} \AgdaSymbol{∀\{}\AgdaBound{b} \AgdaBound{c} \AgdaBound{d}\AgdaSymbol{\}\{}\AgdaBound{p} \AgdaSymbol{:} \AgdaBound{b} \AgdaFunction{\textasciitilde} \AgdaBound{c}\AgdaSymbol{\}\{}\AgdaBound{q} \AgdaSymbol{:} \AgdaBound{c} \AgdaFunction{\textasciitilde} \AgdaBound{d}\AgdaSymbol{\}\{}\AgdaBound{r} \AgdaSymbol{:} \AgdaFunction{P} \AgdaBound{b}\AgdaSymbol{\}} \AgdaSymbol{→} \AgdaFunction{P₁} \AgdaSymbol{(}\AgdaBound{p} \AgdaFunction{*} \AgdaBound{q}\AgdaSymbol{)} \AgdaBound{r} \AgdaDatatype{≡} \AgdaFunction{P₁} \AgdaBound{q} \AgdaSymbol{(}\AgdaFunction{P₁} \AgdaBound{p} \AgdaBound{r}\AgdaSymbol{)}\<%
\\
%
\\
\>[8]\AgdaIndent{10}{}\<[10]%
\>[10]\AgdaFunction{P₁-comp} \AgdaSymbol{=} \AgdaFunction{sym} \AgdaFunction{assoc}\<%
\\
\>\<\end{code}


We didn't write the equivalence proof explicitly i.e.\ isEquivalence (P₁ x).

\begin{code}\>\<%
\\
\>[8]\AgdaIndent{10}{}\<[10]%
\>[10]\AgdaFunction{P-resp} \AgdaSymbol{:} \AgdaFunction{P} \AgdaFunction{respects} \AgdaFunction{\_\textasciitilde\_}\<%
\\
\>[8]\AgdaIndent{10}{}\<[10]%
\>[10]\AgdaFunction{P-resp} \AgdaSymbol{\{}\AgdaBound{b}\AgdaSymbol{\}} \AgdaSymbol{\{}\AgdaBound{b'}\AgdaSymbol{\}} \AgdaBound{bb'} \AgdaSymbol{=} \AgdaBound{UA} \AgdaSymbol{(}\AgdaFunction{P₁} \AgdaBound{bb'} \AgdaInductiveConstructor{,} \AgdaFunction{P₁} \AgdaSymbol{(}\AgdaBound{bb'} \AgdaFunction{⁻¹}\AgdaSymbol{))}\<%
\\
%
\\
\>[8]\AgdaIndent{10}{}\<[10]%
\>[10]\AgdaFunction{P\textasciicircum} \AgdaSymbol{:} \AgdaFunction{Q} \AgdaSymbol{→} \AgdaFunction{Prp}\<%
\\
\>[8]\AgdaIndent{10}{}\<[10]%
\>[10]\AgdaFunction{P\textasciicircum} \AgdaSymbol{=} \AgdaPostulate{lift₁} \AgdaFunction{P} \AgdaFunction{P-resp}\<%
\\
%
\\
\>[8]\AgdaIndent{10}{}\<[10]%
\>[10]\AgdaFunction{P\textasciicircum-β} \AgdaSymbol{:} \AgdaFunction{P} \AgdaBound{a} \AgdaDatatype{≡} \AgdaFunction{P} \AgdaBound{a'}\<%
\\
\>[8]\AgdaIndent{10}{}\<[10]%
\>[10]\AgdaFunction{P\textasciicircum-β} \AgdaSymbol{=} \AgdaFunction{trans} \AgdaSymbol{(}\AgdaFunction{sym} \AgdaPostulate{lift-β₁}\AgdaSymbol{)} \AgdaSymbol{(}\AgdaFunction{trans} \AgdaSymbol{(}\AgdaFunction{cong} \AgdaFunction{P\textasciicircum} \AgdaBound{eq}\AgdaSymbol{)} \AgdaPostulate{lift-β₁}\AgdaSymbol{)}\<%
\\
%
\\
%
\\
\>\AgdaComment{-- coerce (trans (sym lift-β₁) (trans (cong (lift₁ P P-resp) [ p ]⁼) lift-β₁)) (id \{a\}) ≡ p}\<%
\\
\>\AgdaComment{--  isInv₁ : ∀\{a b\}\{p : a \textasciitilde b\} → exact [ p ]⁼ ≡ p}\<%
\\
\>\AgdaComment{--  isInv₁ \{p = p\} = \{!!\}}\<%
\\
%
\\
\>\AgdaComment{--  isInv₂ : ∀\{a b\}\{p : [ a ] ≡ [ b ]\} → [ exact p ]⁼ ≡ p}\<%
\\
%
\\
%
\\
\>\<\end{code}


\chapter{Category with families of setoids}


\AgdaHide{

\begin{code}\>\<%
\\
%
\\
\>\AgdaKeyword{open} \AgdaKeyword{import} \AgdaModule{Level}\<%
\\
\>\AgdaKeyword{open} \AgdaKeyword{import} \AgdaModule{Relation.Binary.PropositionalEquality}\<%
\\
%
\\
\>\AgdaKeyword{module} \AgdaModule{hProp} \AgdaSymbol{(}\AgdaBound{ext} \AgdaSymbol{:} \AgdaFunction{Extensionality} \AgdaPrimitive{zero} \AgdaPrimitive{zero}\AgdaSymbol{)} \AgdaKeyword{where}\<%
\\
%
\\
\>\AgdaKeyword{open} \AgdaKeyword{import} \AgdaModule{Relation.Nullary}\<%
\\
\>\AgdaKeyword{open} \AgdaKeyword{import} \AgdaModule{Data.Unit}\<%
\\
\>\AgdaKeyword{open} \AgdaKeyword{import} \AgdaModule{Data.Empty}\<%
\\
\>\AgdaKeyword{open} \AgdaKeyword{import} \AgdaModule{Data.Nat}\<%
\\
\>\AgdaKeyword{open} \AgdaKeyword{import} \AgdaModule{Data.Product}\<%
\\
%
\\
\>\AgdaKeyword{infixr} \AgdaNumber{2} \_⇒\_\<%
\\
%
\\
\>\AgdaKeyword{infixr} \AgdaNumber{3} \_∧\_\<%
\\
%
\\
%
\\
\>\<\end{code}
}

A proof-irrelvant universe only contains sets with at most one inhabitant. 

\begin{code}\>\<%
\\
%
\\
\>\AgdaKeyword{record} \AgdaRecord{hProp} \AgdaSymbol{:} \AgdaPrimitiveType{Set₁} \AgdaKeyword{where}\<%
\\
\>[0]\AgdaIndent{2}{}\<[2]%
\>[2]\AgdaKeyword{constructor} \AgdaInductiveConstructor{hp}\<%
\\
\>[0]\AgdaIndent{2}{}\<[2]%
\>[2]\AgdaKeyword{field}\<%
\\
\>[2]\AgdaIndent{4}{}\<[4]%
\>[4]\AgdaField{prf} \AgdaSymbol{:} \AgdaPrimitiveType{Set}\<%
\\
\>[2]\AgdaIndent{4}{}\<[4]%
\>[4]\AgdaField{Uni} \AgdaSymbol{:} \AgdaSymbol{\{}\AgdaBound{p} \AgdaBound{q} \AgdaSymbol{:} \AgdaBound{prf}\AgdaSymbol{\}} \AgdaSymbol{→} \AgdaBound{p} \AgdaDatatype{≡} \AgdaBound{q}\<%
\\
%
\\
\>\AgdaKeyword{open} \AgdaModule{hProp} \AgdaKeyword{public} \AgdaKeyword{renaming} \AgdaSymbol{(}prf \AgdaSymbol{to} <\_>\AgdaSymbol{)}\<%
\\
%
\\
\>\<\end{code}

We can extract the proof of any propostion $A : hProp$ by using $<>$ and there is always a proof that all inhabitants of it are the same, in other words, if there is any proof of it, the proof is unique. This is not exactly the same as the $Prop$ universe in Altenkirch's approach which is judgemental. It is just a judgement whether a set behaves like a $Proposition$. The $hProp$ we define above is propositional since we can extract the proof of uniqueness.

We would like to have some basic propositions $\top$ and $\bot$. To distinguish them with the ones for non-proof irrelevant propositions which are already available in Agda library, we add a prime to all similar symbols.

\begin{code}\>\<%
\\
%
\\
\>\AgdaFunction{⊤'} \AgdaSymbol{:} \AgdaRecord{hProp}\<%
\\
\>\AgdaFunction{⊤'} \AgdaSymbol{=} \AgdaInductiveConstructor{hp} \AgdaRecord{⊤} \AgdaInductiveConstructor{refl}\<%
\\
%
\\
\>\AgdaFunction{⊥'} \AgdaSymbol{:} \AgdaRecord{hProp}\<%
\\
\>\AgdaFunction{⊥'} \AgdaSymbol{=} \AgdaInductiveConstructor{hp} \AgdaDatatype{⊥} \AgdaSymbol{(λ} \AgdaSymbol{\{}\AgdaBound{p}\AgdaSymbol{\}} \AgdaSymbol{→} \AgdaFunction{⊥-elim} \AgdaBound{p}\AgdaSymbol{)}\<%
\\
%
\\
\>\<\end{code}

We also want the universal and existential quantifier for $hProp$, namely it is closed under $\Pi$-types and $\Sigma$-types.
The universal quantifier of $hProp$ can be axiomitised but we decide to explicitly state that we 
require the functional extensionality to use this module. The reason is that functional extensionality is actually equivalent to the closure under $\Pi$-types.

\begin{code}\>\<%
\\
%
\\
\>\AgdaFunction{∀'} \AgdaSymbol{:} \AgdaSymbol{(}\AgdaBound{A} \AgdaSymbol{:} \AgdaPrimitiveType{Set}\AgdaSymbol{)(}\AgdaBound{P} \AgdaSymbol{:} \AgdaBound{A} \AgdaSymbol{→} \AgdaRecord{hProp}\AgdaSymbol{)} \AgdaSymbol{→} \AgdaRecord{hProp}\<%
\\
\>\AgdaFunction{∀'} \AgdaBound{A} \AgdaBound{P} \AgdaSymbol{=} \AgdaInductiveConstructor{hp} \AgdaSymbol{((}\AgdaBound{x} \AgdaSymbol{:} \AgdaBound{A}\AgdaSymbol{)} \AgdaSymbol{→} \AgdaFunction{<} \AgdaBound{P} \AgdaBound{x} \AgdaFunction{>}\AgdaSymbol{)} \AgdaSymbol{(}\AgdaBound{ext} \AgdaSymbol{(λ} \AgdaBound{x} \AgdaSymbol{→} \AgdaFunction{Uni} \AgdaSymbol{(}\AgdaBound{P} \AgdaBound{x}\AgdaSymbol{)))}\<%
\\
%
\\
\>\<\end{code}


\AgdaHide{
\begin{code}\>\<%
\\
%
\\
\>\AgdaFunction{sig-eq} \AgdaSymbol{:} \AgdaSymbol{\{}\AgdaBound{A} \AgdaSymbol{:} \AgdaPrimitiveType{Set}\AgdaSymbol{\}\{}\AgdaBound{B} \AgdaSymbol{:} \AgdaBound{A} \AgdaSymbol{→} \AgdaPrimitiveType{Set}\AgdaSymbol{\}\{}\AgdaBound{a} \AgdaBound{b} \AgdaSymbol{:} \AgdaBound{A}\AgdaSymbol{\}} \AgdaSymbol{→} \<[43]%
\>[43]\<%
\\
\>[4]\AgdaIndent{9}{}\<[9]%
\>[9]\AgdaSymbol{(}\AgdaBound{p} \AgdaSymbol{:} \AgdaBound{a} \AgdaDatatype{≡} \AgdaBound{b}\AgdaSymbol{)} \AgdaSymbol{→} \<[23]%
\>[23]\<%
\\
\>[4]\AgdaIndent{9}{}\<[9]%
\>[9]\AgdaSymbol{\{}\AgdaBound{c} \AgdaSymbol{:} \AgdaBound{B} \AgdaBound{a}\AgdaSymbol{\}\{}\AgdaBound{d} \AgdaSymbol{:} \AgdaBound{B} \AgdaBound{b}\AgdaSymbol{\}} \AgdaSymbol{→} \<[30]%
\>[30]\<%
\\
\>[4]\AgdaIndent{9}{}\<[9]%
\>[9]\AgdaSymbol{(}\AgdaFunction{subst} \AgdaSymbol{(λ} \AgdaBound{x} \AgdaSymbol{→} \AgdaBound{B} \AgdaBound{x}\AgdaSymbol{)} \AgdaBound{p} \AgdaBound{c} \AgdaDatatype{≡} \AgdaBound{d}\AgdaSymbol{)} \<[37]%
\>[37]\<%
\\
\>[4]\AgdaIndent{9}{}\<[9]%
\>[9]\AgdaSymbol{→} \AgdaDatatype{\_≡\_} \AgdaSymbol{\{\_\}} \AgdaSymbol{\{}\AgdaRecord{Σ} \AgdaBound{A} \AgdaBound{B}\AgdaSymbol{\}} \AgdaSymbol{(}\AgdaBound{a} \AgdaInductiveConstructor{,} \AgdaBound{c}\AgdaSymbol{)} \AgdaSymbol{(}\AgdaBound{b} \AgdaInductiveConstructor{,} \AgdaBound{d}\AgdaSymbol{)}\<%
\\
\>\AgdaFunction{sig-eq} \AgdaInductiveConstructor{refl} \AgdaInductiveConstructor{refl} \AgdaSymbol{=} \AgdaInductiveConstructor{refl}\<%
\\
%
\\
\>\<\end{code}
}

\begin{code}\>\<%
\\
%
\\
%
\\
\>\AgdaFunction{Σ'} \AgdaSymbol{:} \AgdaSymbol{(}\AgdaBound{P} \AgdaSymbol{:} \AgdaRecord{hProp}\AgdaSymbol{)(}\AgdaBound{Q} \AgdaSymbol{:} \AgdaFunction{<} \AgdaBound{P} \AgdaFunction{>} \AgdaSymbol{→} \AgdaRecord{hProp}\AgdaSymbol{)} \AgdaSymbol{→} \AgdaRecord{hProp}\<%
\\
\>\AgdaFunction{Σ'} \AgdaBound{P} \AgdaBound{Q} \AgdaSymbol{=} \AgdaInductiveConstructor{hp} \AgdaSymbol{(}\AgdaRecord{Σ} \AgdaFunction{<} \AgdaBound{P} \AgdaFunction{>} \AgdaSymbol{(λ} \AgdaBound{x} \AgdaSymbol{→} \AgdaFunction{<} \AgdaBound{Q} \AgdaBound{x} \AgdaFunction{>}\AgdaSymbol{))} \<[38]%
\>[38]\<%
\\
\>[9]\AgdaIndent{12}{}\<[12]%
\>[12]\AgdaSymbol{(λ} \AgdaSymbol{\{}\AgdaBound{p}\AgdaSymbol{\}} \AgdaSymbol{\{}\AgdaBound{q}\AgdaSymbol{\}} \AgdaSymbol{→} \<[25]%
\>[25]\<%
\\
\>[9]\AgdaIndent{12}{}\<[12]%
\>[12]\AgdaFunction{sig-eq} \AgdaSymbol{(}\AgdaFunction{Uni} \AgdaBound{P}\AgdaSymbol{)} \AgdaSymbol{(}\AgdaFunction{Uni} \AgdaSymbol{(}\AgdaBound{Q} \AgdaSymbol{(}\AgdaFunction{proj₁} \AgdaBound{q}\AgdaSymbol{))))}\<%
\\
%
\\
\>\<\end{code}

Implication and conjuction which are independent ones of them follow simply.

\begin{code}\>\<%
\\
%
\\
\>\AgdaFunction{\_⇒\_} \AgdaSymbol{:} \AgdaSymbol{(}\AgdaBound{P} \AgdaBound{Q} \AgdaSymbol{:} \AgdaRecord{hProp}\AgdaSymbol{)} \AgdaSymbol{→} \AgdaRecord{hProp}\<%
\\
\>\AgdaBound{P} \AgdaFunction{⇒} \AgdaBound{Q} \AgdaSymbol{=} \<[9]%
\>[9]\AgdaFunction{∀'} \AgdaFunction{<} \AgdaBound{P} \AgdaFunction{>} \AgdaSymbol{(λ} \AgdaBound{\_} \AgdaSymbol{→} \AgdaBound{Q}\AgdaSymbol{)}\<%
\\
%
\\
\>\AgdaFunction{\_∧\_} \AgdaSymbol{:} \AgdaSymbol{(}\AgdaBound{P} \AgdaBound{Q} \AgdaSymbol{:} \AgdaRecord{hProp}\AgdaSymbol{)} \AgdaSymbol{→} \AgdaRecord{hProp}\<%
\\
\>\AgdaBound{P} \AgdaFunction{∧} \AgdaBound{Q} \AgdaSymbol{=} \AgdaFunction{Σ'} \AgdaBound{P} \AgdaSymbol{(λ} \AgdaBound{\_} \AgdaSymbol{→} \AgdaBound{Q}\AgdaSymbol{)}\<%
\\
%
\\
\>\<\end{code}

\AgdaHide{
\begin{code}\>\<%
\\
\>\AgdaKeyword{syntax} ∀' A \AgdaSymbol{(λ} x \AgdaSymbol{→} B\AgdaSymbol{)} \AgdaSymbol{=} ∀'[ x ∶ A ] B


\AgdaKeyword{syntax} Σ' A \AgdaSymbol{(λ} x \AgdaSymbol{→} B\AgdaSymbol{)} \AgdaSymbol{=} Σ'[ x ∶ A ] B

\<\end{code}
}

As long as we have implication and conjuction, more operators on proposition can be defined, for instances negation and logical equivalence.

\begin{code}\>\<%
\\
%
\\
\>\AgdaFunction{¬} \AgdaSymbol{:} \AgdaRecord{hProp} \AgdaSymbol{→} \AgdaRecord{hProp}\<%
\\
\>\AgdaFunction{¬} \AgdaBound{P} \AgdaSymbol{=} \AgdaBound{P} \AgdaFunction{⇒} \AgdaFunction{⊥'} \<[13]%
\>[13]\<%
\\
%
\\
\>\AgdaFunction{\_↔\_} \<[6]%
\>[6]\AgdaSymbol{:} \AgdaSymbol{(}\AgdaBound{P} \AgdaBound{Q} \AgdaSymbol{:} \AgdaRecord{hProp}\AgdaSymbol{)} \AgdaSymbol{→} \AgdaRecord{hProp}\<%
\\
\>\AgdaBound{P} \AgdaFunction{↔} \AgdaBound{Q} \AgdaSymbol{=} \AgdaSymbol{(}\AgdaBound{P} \AgdaFunction{⇒} \AgdaBound{Q}\AgdaSymbol{)} \AgdaFunction{∧} \AgdaSymbol{(}\AgdaBound{Q} \AgdaFunction{⇒} \AgdaBound{P}\AgdaSymbol{)}\<%
\\
%
\\
\>\<\end{code}

\section{Category}


\AgdaHide{
\begin{code}\>\<%
\\
%
\\
\>\AgdaSymbol{\{-\#} \AgdaKeyword{OPTIONS} --type-in-type \AgdaSymbol{\#-\}}\<%
\\
%
\\
\>\AgdaKeyword{module} \AgdaModule{Category} \AgdaKeyword{where}\<%
\\
%
\\
\>\AgdaKeyword{open} \AgdaKeyword{import} \AgdaModule{Data.Product}\<%
\\
\>\AgdaKeyword{open} \AgdaKeyword{import} \AgdaModule{Relation.Binary.PropositionalEquality}\<%
\\
%
\\
\>\AgdaKeyword{open} \AgdaKeyword{import} \AgdaModule{Level}\<%
\\
%
\\
\>\<\end{code}
}

\AgdaHide{
\begin{code}\>\<%
\\
%
\\
\>\AgdaKeyword{record} \AgdaRecord{IsCategory}\<%
\\
\>[0]\AgdaIndent{2}{}\<[2]%
\>[2]\AgdaSymbol{(}\AgdaBound{obj} \<[12]%
\>[12]\AgdaSymbol{:} \AgdaPrimitiveType{Set}\AgdaSymbol{)}\<%
\\
%
\\
\>[0]\AgdaIndent{2}{}\<[2]%
\>[2]\AgdaSymbol{(}\AgdaBound{hom} \<[12]%
\>[12]\AgdaSymbol{:} \AgdaBound{obj} \AgdaSymbol{→} \AgdaBound{obj} \AgdaSymbol{→} \AgdaPrimitiveType{Set}\AgdaSymbol{)}\<%
\\
%
\\
\>[0]\AgdaIndent{2}{}\<[2]%
\>[2]\AgdaSymbol{(}\AgdaBound{id} \<[12]%
\>[12]\AgdaSymbol{:} \AgdaSymbol{∀} \AgdaBound{α} \AgdaSymbol{→} \AgdaBound{hom} \AgdaBound{α} \AgdaBound{α}\AgdaSymbol{)}\<%
\\
%
\\
\>[0]\AgdaIndent{2}{}\<[2]%
\>[2]\AgdaSymbol{(}\AgdaBound{[\_⇒\_]\_∘\_} \AgdaSymbol{:} \AgdaSymbol{∀} \AgdaBound{α} \AgdaSymbol{\{}\AgdaBound{β}\AgdaSymbol{\}} \AgdaBound{γ}\<%
\\
\>[2]\AgdaIndent{12}{}\<[12]%
\>[12]\AgdaSymbol{→} \AgdaBound{hom} \AgdaBound{β} \AgdaBound{γ}\<%
\\
\>[2]\AgdaIndent{12}{}\<[12]%
\>[12]\AgdaSymbol{→} \AgdaBound{hom} \AgdaBound{α} \AgdaBound{β}\<%
\\
\>[2]\AgdaIndent{12}{}\<[12]%
\>[12]\AgdaSymbol{→} \AgdaBound{hom} \AgdaBound{α} \AgdaBound{γ}\AgdaSymbol{)}\<%
\\
\>[0]\AgdaIndent{2}{}\<[2]%
\>[2]\AgdaSymbol{:} \AgdaPrimitiveType{Set}\<%
\\
\>[0]\AgdaIndent{2}{}\<[2]%
\>[2]\AgdaKeyword{where}\<%
\\
\>[2]\AgdaIndent{4}{}\<[4]%
\>[4]\AgdaKeyword{constructor} \AgdaInductiveConstructor{IsCatC}\<%
\\
\>[2]\AgdaIndent{4}{}\<[4]%
\>[4]\AgdaKeyword{field}\<%
\\
\>[4]\AgdaIndent{6}{}\<[6]%
\>[6]\AgdaField{id₁} \<[11]%
\>[11]\AgdaSymbol{:} \AgdaSymbol{∀} \AgdaBound{α} \AgdaBound{β} \AgdaSymbol{(}\AgdaBound{f} \AgdaSymbol{:} \AgdaBound{hom} \AgdaBound{α} \AgdaBound{β}\AgdaSymbol{)}\<%
\\
\>[6]\AgdaIndent{11}{}\<[11]%
\>[11]\AgdaSymbol{→} \AgdaBound{[} \AgdaBound{α} \AgdaBound{⇒} \AgdaBound{β} \AgdaBound{]} \AgdaBound{f} \AgdaBound{∘} \AgdaSymbol{(}\AgdaBound{id} \AgdaBound{α}\AgdaSymbol{)} \AgdaDatatype{≡} \AgdaBound{f}\<%
\\
%
\\
\>[0]\AgdaIndent{6}{}\<[6]%
\>[6]\AgdaField{id₂} \<[11]%
\>[11]\AgdaSymbol{:} \AgdaSymbol{∀} \AgdaBound{α} \AgdaBound{β} \AgdaSymbol{(}\AgdaBound{f} \AgdaSymbol{:} \AgdaBound{hom} \AgdaBound{α} \AgdaBound{β}\AgdaSymbol{)}\<%
\\
\>[0]\AgdaIndent{11}{}\<[11]%
\>[11]\AgdaSymbol{→} \AgdaBound{[} \AgdaBound{α} \AgdaBound{⇒} \AgdaBound{β} \AgdaBound{]} \AgdaSymbol{(}\AgdaBound{id} \AgdaBound{β}\AgdaSymbol{)} \AgdaBound{∘} \AgdaBound{f} \AgdaDatatype{≡} \AgdaBound{f}\<%
\\
%
\\
\>[0]\AgdaIndent{6}{}\<[6]%
\>[6]\AgdaField{comp} \AgdaSymbol{:} \AgdaSymbol{∀} \AgdaBound{α} \AgdaSymbol{\{}\AgdaBound{β} \AgdaBound{γ}\AgdaSymbol{\}} \AgdaBound{δ} \AgdaSymbol{(}\AgdaBound{f} \AgdaSymbol{:} \AgdaBound{hom} \AgdaBound{α} \AgdaBound{β}\AgdaSymbol{)} \AgdaSymbol{(}\AgdaBound{g} \AgdaSymbol{:} \AgdaBound{hom} \AgdaBound{β} \AgdaBound{γ}\AgdaSymbol{)} \AgdaSymbol{(}\AgdaBound{h} \AgdaSymbol{:} \AgdaBound{hom} \AgdaBound{γ} \AgdaBound{δ}\AgdaSymbol{)}\<%
\\
\>[0]\AgdaIndent{11}{}\<[11]%
\>[11]\AgdaSymbol{→} \AgdaBound{[} \AgdaBound{α} \AgdaBound{⇒} \AgdaBound{δ} \AgdaBound{]} \AgdaBound{[} \AgdaBound{β} \AgdaBound{⇒} \AgdaBound{δ} \AgdaBound{]} \AgdaBound{h} \AgdaBound{∘} \AgdaBound{g} \AgdaBound{∘} \AgdaBound{f} \AgdaDatatype{≡} \AgdaBound{[} \AgdaBound{α} \AgdaBound{⇒} \AgdaBound{δ} \AgdaBound{]} \AgdaBound{h} \AgdaBound{∘} \AgdaSymbol{(}\AgdaBound{[} \AgdaBound{α} \AgdaBound{⇒} \AgdaBound{γ} \AgdaBound{]} \AgdaBound{g} \AgdaBound{∘} \AgdaBound{f}\AgdaSymbol{)}\<%
\\
%
\\
%
\\
\>\<\end{code}
}


\begin{code}\>\<%
\\
\>\AgdaKeyword{record} \AgdaRecord{Category} \AgdaSymbol{:} \AgdaPrimitiveType{Set} \AgdaKeyword{where}\<%
\\
\>[0]\AgdaIndent{2}{}\<[2]%
\>[2]\AgdaKeyword{constructor} \AgdaInductiveConstructor{CatC}\<%
\\
\>[0]\AgdaIndent{2}{}\<[2]%
\>[2]\AgdaKeyword{field}\<%
\\
\>[2]\AgdaIndent{4}{}\<[4]%
\>[4]\AgdaField{obj} \<[15]%
\>[15]\AgdaSymbol{:} \AgdaPrimitiveType{Set}\<%
\\
%
\\
\>[2]\AgdaIndent{4}{}\<[4]%
\>[4]\AgdaField{hom} \<[15]%
\>[15]\AgdaSymbol{:} \AgdaBound{obj} \AgdaSymbol{→} \AgdaBound{obj} \AgdaSymbol{→} \AgdaPrimitiveType{Set}\<%
\\
%
\\
\>[2]\AgdaIndent{4}{}\<[4]%
\>[4]\AgdaField{id} \<[15]%
\>[15]\AgdaSymbol{:} \AgdaSymbol{∀} \AgdaBound{α}\<%
\\
\>[4]\AgdaIndent{15}{}\<[15]%
\>[15]\AgdaSymbol{→} \AgdaBound{hom} \AgdaBound{α} \AgdaBound{α}\<%
\\
%
\\
\>[0]\AgdaIndent{4}{}\<[4]%
\>[4]\AgdaField{[\_⇒\_]\_∘\_} \<[15]%
\>[15]\AgdaSymbol{:} \AgdaSymbol{∀} \AgdaBound{α} \AgdaSymbol{\{}\AgdaBound{β}\AgdaSymbol{\}} \AgdaBound{γ}\<%
\\
\>[0]\AgdaIndent{15}{}\<[15]%
\>[15]\AgdaSymbol{→} \AgdaBound{hom} \AgdaBound{β} \AgdaBound{γ}\<%
\\
\>[0]\AgdaIndent{15}{}\<[15]%
\>[15]\AgdaSymbol{→} \AgdaBound{hom} \AgdaBound{α} \AgdaBound{β}\<%
\\
\>[0]\AgdaIndent{15}{}\<[15]%
\>[15]\AgdaSymbol{→} \AgdaBound{hom} \AgdaBound{α} \AgdaBound{γ}\<%
\\
%
\\
\>[0]\AgdaIndent{4}{}\<[4]%
\>[4]\AgdaField{isCategory} \AgdaSymbol{:} \AgdaRecord{IsCategory} \AgdaBound{obj} \AgdaBound{hom} \AgdaBound{id} \AgdaBound{[\_⇒\_]\_∘\_}\<%
\\
%
\\
%
\\
\>\<\end{code}

\AgdaHide{
\begin{code}\>\<%
\\
%
\\
\>[0]\AgdaIndent{2}{}\<[2]%
\>[2]\AgdaKeyword{open} \AgdaModule{IsCategory} \AgdaKeyword{public}\<%
\\
%
\\
\>\<\end{code}
}

\section{Category of setoids}

\AgdaHide{
\begin{code}\>\<%
\\
%
\\
\>\AgdaSymbol{\{-\#} \AgdaKeyword{OPTIONS} --type-in-type \AgdaSymbol{\#-\}}\<%
\\
%
\\
\>\AgdaKeyword{open} \AgdaKeyword{import} \AgdaModule{Level}\<%
\\
\>\AgdaKeyword{open} \AgdaKeyword{import} \AgdaModule{Relation.Binary.PropositionalEquality} \AgdaSymbol{as} \AgdaModule{PE} \AgdaKeyword{hiding} \AgdaSymbol{(}refl\AgdaSymbol{;} sym \AgdaSymbol{;} trans\AgdaSymbol{;} isEquivalence\AgdaSymbol{)}\<%
\\
%
\\
\>\AgdaKeyword{module} \AgdaModule{CategoryOfSetoid} \<[25]%
\>[25]\AgdaSymbol{(}\AgdaBound{ext} \AgdaSymbol{:} \AgdaFunction{Extensionality} \AgdaPrimitive{zero} \AgdaPrimitive{zero}\AgdaSymbol{)} \AgdaKeyword{where}\<%
\\
%
\\
\>\AgdaKeyword{open} \AgdaKeyword{import} \AgdaModule{Cats.Category}\<%
\\
\>\AgdaKeyword{open} \AgdaKeyword{import} \AgdaModule{Function}\<%
\\
\>\AgdaKeyword{open} \AgdaKeyword{import} \AgdaModule{Relation.Binary.Core} \AgdaKeyword{using} \AgdaSymbol{(}\_⇒\_\AgdaSymbol{)}\<%
\\
\>\AgdaKeyword{open} \AgdaKeyword{import} \AgdaModule{Data.Empty}\<%
\\
\>\AgdaKeyword{open} \AgdaKeyword{import} \AgdaModule{HProp} \AgdaBound{ext} \AgdaKeyword{public}\<%
\\
\>\AgdaKeyword{open} \AgdaKeyword{import} \AgdaModule{Data.Unit}\<%
\\
%
\\
\>\<\end{code}
}

\section{Category of setoids}

\begin{code}\>\<%
\\
%
\\
\>\AgdaKeyword{record} \AgdaRecord{HSetoid} \AgdaSymbol{:} \AgdaPrimitiveType{Set₁} \AgdaKeyword{where}\<%
\\
\>[0]\AgdaIndent{2}{}\<[2]%
\>[2]\AgdaKeyword{constructor} \AgdaInductiveConstructor{\_,\_,\_,\_,\_}\<%
\\
\>[0]\AgdaIndent{2}{}\<[2]%
\>[2]\AgdaKeyword{infix} \AgdaNumber{4} \_≈h\_ \_≈\_\<%
\\
\>[0]\AgdaIndent{2}{}\<[2]%
\>[2]\AgdaKeyword{field}\<%
\\
\>[2]\AgdaIndent{4}{}\<[4]%
\>[4]\AgdaField{Carrier} \AgdaSymbol{:} \AgdaPrimitiveType{Set}\<%
\\
\>[2]\AgdaIndent{4}{}\<[4]%
\>[4]\AgdaField{\_≈h\_} \<[13]%
\>[13]\AgdaSymbol{:} \AgdaBound{Carrier} \AgdaSymbol{→} \AgdaBound{Carrier} \AgdaSymbol{→} \AgdaRecord{HProp}\<%
\\
\>[2]\AgdaIndent{4}{}\<[4]%
\>[4]\AgdaField{refl} \<[12]%
\>[12]\AgdaSymbol{:} \AgdaSymbol{\{}\AgdaBound{x} \AgdaSymbol{:} \AgdaBound{Carrier}\AgdaSymbol{\}} \AgdaSymbol{→} \AgdaFunction{<} \AgdaBound{x} \AgdaBound{≈h} \AgdaBound{x} \AgdaFunction{>}\<%
\\
\>[2]\AgdaIndent{4}{}\<[4]%
\>[4]\AgdaField{sym} \<[12]%
\>[12]\AgdaSymbol{:} \AgdaSymbol{\{}\AgdaBound{x} \AgdaBound{y} \AgdaSymbol{:} \AgdaBound{Carrier}\AgdaSymbol{\}} \AgdaSymbol{→} \AgdaFunction{<} \AgdaBound{x} \AgdaBound{≈h} \AgdaBound{y} \AgdaFunction{>} \AgdaSymbol{→} \AgdaFunction{<} \AgdaBound{y} \AgdaBound{≈h} \AgdaBound{x} \AgdaFunction{>}\<%
\\
\>[2]\AgdaIndent{4}{}\<[4]%
\>[4]\AgdaField{trans} \<[12]%
\>[12]\AgdaSymbol{:} \AgdaSymbol{\{}\AgdaBound{x} \AgdaBound{y} \AgdaBound{z} \AgdaSymbol{:} \AgdaBound{Carrier}\AgdaSymbol{\}} \AgdaSymbol{→} \AgdaFunction{<} \AgdaBound{x} \AgdaBound{≈h} \AgdaBound{y} \AgdaFunction{>} \AgdaSymbol{→} \AgdaFunction{<} \AgdaBound{y} \AgdaBound{≈h} \AgdaBound{z} \AgdaFunction{>} \AgdaSymbol{→} \AgdaFunction{<} \AgdaBound{x} \AgdaBound{≈h} \AgdaBound{z} \AgdaFunction{>}\<%
\\
\>[0]\AgdaIndent{2}{}\<[2]%
\>[2]\<%
\\
\>[0]\AgdaIndent{2}{}\<[2]%
\>[2]\AgdaFunction{\_≈\_} \AgdaSymbol{:} \AgdaFunction{Carrier} \AgdaSymbol{→} \AgdaFunction{Carrier} \AgdaSymbol{→} \AgdaPrimitiveType{Set}\<%
\\
\>[0]\AgdaIndent{2}{}\<[2]%
\>[2]\AgdaBound{a} \AgdaFunction{≈} \AgdaBound{b} \AgdaSymbol{=} \AgdaFunction{<} \AgdaBound{a} \AgdaFunction{≈h} \AgdaBound{b} \AgdaFunction{>}\<%
\\
%
\\
\>[0]\AgdaIndent{2}{}\<[2]%
\>[2]\AgdaFunction{PI} \AgdaSymbol{:} \AgdaSymbol{\{}\AgdaBound{x} \AgdaBound{y} \AgdaSymbol{:} \AgdaFunction{Carrier}\AgdaSymbol{\}\{}\AgdaBound{B} \AgdaSymbol{:} \AgdaPrimitiveType{Set}\AgdaSymbol{\}(}\AgdaBound{A} \AgdaSymbol{:} \AgdaBound{x} \AgdaFunction{≈} \AgdaBound{y} \AgdaSymbol{→} \AgdaBound{B}\AgdaSymbol{)\{}\AgdaBound{p} \AgdaBound{q} \AgdaSymbol{:} \AgdaBound{x} \AgdaFunction{≈} \AgdaBound{y}\AgdaSymbol{\}} \AgdaSymbol{→} \AgdaBound{A} \AgdaBound{p} \AgdaDatatype{≡} \AgdaBound{A} \AgdaBound{q}\<%
\\
\>[0]\AgdaIndent{2}{}\<[2]%
\>[2]\AgdaFunction{PI} \AgdaSymbol{\{}\AgdaBound{x}\AgdaSymbol{\}} \AgdaSymbol{\{}\AgdaBound{y}\AgdaSymbol{\}} \AgdaBound{A} \AgdaSymbol{\{}\AgdaBound{p}\AgdaSymbol{\}} \AgdaSymbol{\{}\AgdaBound{q}\AgdaSymbol{\}} \AgdaKeyword{with} \AgdaFunction{Uni} \AgdaSymbol{(}\AgdaBound{x} \AgdaFunction{≈h} \AgdaBound{y}\AgdaSymbol{)} \AgdaSymbol{\{}\AgdaBound{p}\AgdaSymbol{\}} \AgdaSymbol{\{}\AgdaBound{q}\AgdaSymbol{\}}\<%
\\
\>[0]\AgdaIndent{2}{}\<[2]%
\>[2]\AgdaFunction{PI} \AgdaBound{A} \AgdaSymbol{|} \AgdaInductiveConstructor{PE.refl} \AgdaSymbol{=} \AgdaInductiveConstructor{PE.refl}\<%
\\
%
\\
\>[0]\AgdaIndent{2}{}\<[2]%
\>[2]\AgdaFunction{reflexive} \AgdaSymbol{:} \AgdaDatatype{\_≡\_} \AgdaFunction{⇒} \AgdaFunction{\_≈\_}\<%
\\
\>[0]\AgdaIndent{2}{}\<[2]%
\>[2]\AgdaFunction{reflexive} \AgdaInductiveConstructor{PE.refl} \AgdaSymbol{=} \AgdaFunction{refl}\<%
\\
%
\\
\>\AgdaKeyword{open} \AgdaModule{HSetoid} \AgdaKeyword{public} \AgdaKeyword{renaming} \AgdaSymbol{(}refl \AgdaSymbol{to} [\_]refl\AgdaSymbol{;} sym \AgdaSymbol{to} [\_]sym\AgdaSymbol{;} 
  \_≈\_ \AgdaSymbol{to} [\_]\_≈\_ \AgdaSymbol{;} \_≈h\_ \AgdaSymbol{to} [\_]\_≈h\_ \AgdaSymbol{;} Carrier \AgdaSymbol{to} ∣\_∣ \AgdaSymbol{;} trans \AgdaSymbol{to} [\_]trans\AgdaSymbol{)}\<%
\\
%
\\
\>\AgdaFunction{rel} \AgdaSymbol{:} \AgdaSymbol{(}\AgdaBound{A} \AgdaSymbol{:} \AgdaRecord{HSetoid}\AgdaSymbol{)} \AgdaSymbol{→} \AgdaFunction{∣} \AgdaBound{A} \AgdaFunction{∣} \AgdaSymbol{→} \AgdaFunction{∣} \AgdaBound{A} \AgdaFunction{∣} \AgdaSymbol{→} \AgdaRecord{HProp}\<%
\\
\>\AgdaFunction{rel} \AgdaBound{A} \AgdaBound{a} \AgdaBound{b} \AgdaSymbol{=} \AgdaFunction{[} \AgdaBound{A} \AgdaFunction{]} \AgdaBound{a} \AgdaFunction{≈h} \AgdaBound{b}\<%
\\
%
\\
\>\AgdaFunction{[\_]uip} \AgdaSymbol{:} \AgdaSymbol{∀(}\AgdaBound{Γ} \AgdaSymbol{:} \AgdaRecord{HSetoid}\AgdaSymbol{)\{}\AgdaBound{a} \AgdaBound{b} \AgdaSymbol{:} \AgdaFunction{∣} \AgdaBound{Γ} \AgdaFunction{∣}\AgdaSymbol{\}\{}\AgdaBound{p} \AgdaBound{q} \AgdaSymbol{:} \AgdaFunction{[} \AgdaBound{Γ} \AgdaFunction{]} \AgdaBound{a} \AgdaFunction{≈} \AgdaBound{b}\AgdaSymbol{\}} \AgdaSymbol{→} \AgdaBound{p} \AgdaDatatype{≡} \AgdaBound{q}\<%
\\
\>\AgdaFunction{[} \AgdaBound{Γ} \AgdaFunction{]uip} \AgdaSymbol{\{}\AgdaBound{a}\AgdaSymbol{\}} \AgdaSymbol{\{}\AgdaBound{b}\AgdaSymbol{\}} \AgdaSymbol{=} \AgdaFunction{Uni} \AgdaSymbol{(}\AgdaFunction{[} \AgdaBound{Γ} \AgdaFunction{]} \AgdaBound{a} \AgdaFunction{≈h} \AgdaBound{b}\AgdaSymbol{)}\<%
\\
%
\\
\>\<\end{code}

Morphisms between HSetoids (Functors)

\begin{code}\>\<%
\\
%
\\
\>\AgdaKeyword{infix} \AgdaNumber{5} \_⇉\_\<%
\\
%
\\
\>\AgdaKeyword{record} \AgdaRecord{\_⇉\_} \AgdaSymbol{(}\AgdaBound{A} \AgdaBound{B} \AgdaSymbol{:} \AgdaRecord{HSetoid}\AgdaSymbol{)} \AgdaSymbol{:} \AgdaPrimitiveType{Set} \AgdaKeyword{where}\<%
\\
\>[0]\AgdaIndent{2}{}\<[2]%
\>[2]\AgdaKeyword{constructor} \AgdaInductiveConstructor{fn:\_resp:\_}\<%
\\
\>[0]\AgdaIndent{2}{}\<[2]%
\>[2]\AgdaKeyword{field}\<%
\\
\>[2]\AgdaIndent{4}{}\<[4]%
\>[4]\AgdaField{fn} \<[9]%
\>[9]\AgdaSymbol{:} \AgdaFunction{∣} \AgdaBound{A} \AgdaFunction{∣} \AgdaSymbol{→} \AgdaFunction{∣} \AgdaBound{B} \AgdaFunction{∣}\<%
\\
\>[2]\AgdaIndent{4}{}\<[4]%
\>[4]\AgdaField{resp} \AgdaSymbol{:} \AgdaSymbol{\{}\AgdaBound{x} \AgdaBound{y} \AgdaSymbol{:} \AgdaFunction{∣} \AgdaBound{A} \AgdaFunction{∣}\AgdaSymbol{\}} \AgdaSymbol{→} \<[27]%
\>[27]\<%
\\
\>[4]\AgdaIndent{11}{}\<[11]%
\>[11]\AgdaFunction{[} \AgdaBound{A} \AgdaFunction{]} \AgdaBound{x} \AgdaFunction{≈} \AgdaBound{y} \AgdaSymbol{→} \<[25]%
\>[25]\<%
\\
\>[4]\AgdaIndent{11}{}\<[11]%
\>[11]\AgdaFunction{[} \AgdaBound{B} \AgdaFunction{]} \AgdaBound{fn} \AgdaBound{x} \AgdaFunction{≈} \AgdaBound{fn} \AgdaBound{y}\<%
\\
\>\AgdaKeyword{open} \AgdaModule{\_⇉\_} \AgdaKeyword{public} \AgdaKeyword{renaming} \AgdaSymbol{(}fn \AgdaSymbol{to} [\_]fn \AgdaSymbol{;} resp \AgdaSymbol{to} [\_]resp\AgdaSymbol{)}\<%
\\
%
\\
\>\<\end{code}

Identity

\begin{code}\>\<%
\\
%
\\
\>\AgdaFunction{idCH} \AgdaSymbol{:} \AgdaSymbol{\{}\AgdaBound{Γ} \AgdaSymbol{:} \AgdaRecord{HSetoid}\AgdaSymbol{\}} \AgdaSymbol{→} \AgdaBound{Γ} \AgdaRecord{⇉} \AgdaBound{Γ} \<[29]%
\>[29]\<%
\\
\>\AgdaFunction{idCH} \AgdaSymbol{=} \AgdaKeyword{record} \AgdaSymbol{\{} \AgdaField{fn} \AgdaSymbol{=} \AgdaFunction{id}\AgdaSymbol{;} \AgdaField{resp} \AgdaSymbol{=} \AgdaFunction{id}\AgdaSymbol{\}}\<%
\\
%
\\
\>\<\end{code}

Composition

\begin{code}\>\<%
\\
%
\\
\>\AgdaKeyword{infixl} \AgdaNumber{5} \_∘c\_\<%
\\
%
\\
\>\AgdaFunction{\_∘c\_} \AgdaSymbol{:} \AgdaSymbol{∀\{}\AgdaBound{Γ} \AgdaBound{Δ} \AgdaBound{Z}\AgdaSymbol{\}} \AgdaSymbol{→} \AgdaBound{Δ} \AgdaRecord{⇉} \AgdaBound{Z} \AgdaSymbol{→} \AgdaBound{Γ} \AgdaRecord{⇉} \AgdaBound{Δ} \AgdaSymbol{→} \AgdaBound{Γ} \AgdaRecord{⇉} \AgdaBound{Z}\<%
\\
\>\AgdaBound{yz} \AgdaFunction{∘c} \AgdaBound{xy} \AgdaSymbol{=} \AgdaKeyword{record} \<[18]%
\>[18]\<%
\\
\>[4]\AgdaIndent{11}{}\<[11]%
\>[11]\AgdaSymbol{\{} \AgdaField{fn} \AgdaSymbol{=} \AgdaFunction{[} \AgdaBound{yz} \AgdaFunction{]fn} \AgdaFunction{∘} \AgdaFunction{[} \AgdaBound{xy} \AgdaFunction{]fn}\<%
\\
\>[4]\AgdaIndent{11}{}\<[11]%
\>[11]\AgdaSymbol{;} \AgdaField{resp} \AgdaSymbol{=} \AgdaFunction{[} \AgdaBound{yz} \AgdaFunction{]resp} \AgdaFunction{∘} \AgdaFunction{[} \AgdaBound{xy} \AgdaFunction{]resp}\<%
\\
\>[4]\AgdaIndent{11}{}\<[11]%
\>[11]\AgdaSymbol{\}}\<%
\\
\>\<\end{code}

Categorical laws

\begin{code}\>\<%
\\
%
\\
\>\AgdaFunction{id₁} \AgdaSymbol{:} \AgdaSymbol{∀} \AgdaSymbol{\{}\AgdaBound{Γ} \AgdaBound{Δ}\AgdaSymbol{\}(}\AgdaBound{ch} \AgdaSymbol{:} \AgdaBound{Γ} \AgdaRecord{⇉} \AgdaBound{Δ}\AgdaSymbol{)} \AgdaSymbol{→} \AgdaBound{ch} \AgdaFunction{∘c} \AgdaFunction{idCH} \AgdaDatatype{≡} \AgdaBound{ch}\<%
\\
\>\AgdaFunction{id₁} \AgdaBound{ch} \AgdaSymbol{=} \AgdaInductiveConstructor{PE.refl}\<%
\\
%
\\
\>\AgdaFunction{id₂} \AgdaSymbol{:} \AgdaSymbol{∀} \AgdaSymbol{\{}\AgdaBound{Γ} \AgdaBound{Δ}\AgdaSymbol{\}(}\AgdaBound{ch} \AgdaSymbol{:} \AgdaBound{Γ} \AgdaRecord{⇉} \AgdaBound{Δ}\AgdaSymbol{)} \AgdaSymbol{→} \AgdaFunction{idCH} \AgdaFunction{∘c} \AgdaBound{ch} \AgdaDatatype{≡} \AgdaBound{ch}\<%
\\
\>\AgdaFunction{id₂} \AgdaBound{ch} \AgdaSymbol{=} \AgdaInductiveConstructor{PE.refl}\<%
\\
%
\\
\>\AgdaFunction{comp} \AgdaSymbol{:} \AgdaSymbol{∀} \AgdaSymbol{\{}\AgdaBound{Γ} \AgdaBound{Δ} \AgdaBound{Φ} \AgdaBound{Ψ}\AgdaSymbol{\}} \AgdaSymbol{(}\AgdaBound{f} \AgdaSymbol{:} \AgdaBound{Γ} \AgdaRecord{⇉} \AgdaBound{Δ}\AgdaSymbol{)} \AgdaSymbol{(}\AgdaBound{g} \AgdaSymbol{:} \AgdaBound{Δ} \AgdaRecord{⇉} \AgdaBound{Φ}\AgdaSymbol{)} \AgdaSymbol{(}\AgdaBound{h} \AgdaSymbol{:} \AgdaBound{Φ} \AgdaRecord{⇉} \AgdaBound{Ψ}\AgdaSymbol{)}\<%
\\
\>[0]\AgdaIndent{7}{}\<[7]%
\>[7]\AgdaSymbol{→} \AgdaBound{h} \AgdaFunction{∘c} \AgdaBound{g} \AgdaFunction{∘c} \AgdaBound{f} \AgdaDatatype{≡} \AgdaBound{h} \AgdaFunction{∘c} \AgdaSymbol{(}\AgdaBound{g} \AgdaFunction{∘c} \AgdaBound{f}\AgdaSymbol{)}\<%
\\
\>\AgdaFunction{comp} \AgdaBound{f} \AgdaBound{g} \AgdaBound{h} \AgdaSymbol{=} \AgdaInductiveConstructor{PE.refl}\<%
\\
%
\\
\>\AgdaFunction{\_f≈\_} \AgdaSymbol{:} \<[8]%
\>[8]\AgdaSymbol{∀\{}\AgdaBound{Γ} \AgdaBound{Δ} \AgdaSymbol{:} \AgdaRecord{HSetoid}\AgdaSymbol{\}} \AgdaSymbol{→} \AgdaSymbol{(}\AgdaBound{f} \AgdaBound{g} \AgdaSymbol{:} \AgdaBound{Γ} \AgdaRecord{⇉} \AgdaBound{Δ}\AgdaSymbol{)} \AgdaSymbol{→} \AgdaRecord{HProp}\<%
\\
\>\AgdaFunction{\_f≈\_} \AgdaSymbol{\{}\AgdaBound{Γ} \AgdaInductiveConstructor{,} \AgdaBound{\_≈h\_} \AgdaInductiveConstructor{,} \AgdaBound{refl} \AgdaInductiveConstructor{,} \AgdaBound{sym} \AgdaInductiveConstructor{,} \AgdaBound{trans}\AgdaSymbol{\}} \<[37]%
\>[37]\<%
\\
\>[0]\AgdaIndent{5}{}\<[5]%
\>[5]\AgdaSymbol{\{}\AgdaBound{Δ} \AgdaInductiveConstructor{,} \AgdaBound{\_≈h₁\_} \AgdaInductiveConstructor{,} \AgdaBound{refl₁} \AgdaInductiveConstructor{,} \AgdaBound{sym₁} \AgdaInductiveConstructor{,} \AgdaBound{trans₁}\AgdaSymbol{\}}\<%
\\
\>[0]\AgdaIndent{5}{}\<[5]%
\>[5]\AgdaSymbol{(}\AgdaInductiveConstructor{fn:} \AgdaBound{fn} \AgdaInductiveConstructor{resp:} \AgdaBound{fresp}\AgdaSymbol{)} \AgdaSymbol{(}\AgdaInductiveConstructor{fn:} \AgdaBound{gn} \AgdaInductiveConstructor{resp:} \AgdaBound{gresp}\AgdaSymbol{)} \<[47]%
\>[47]\<%
\\
\>[0]\AgdaIndent{2}{}\<[2]%
\>[2]\AgdaSymbol{=} \AgdaKeyword{record} \<[11]%
\>[11]\<%
\\
\>[0]\AgdaIndent{11}{}\<[11]%
\>[11]\AgdaSymbol{\{} \AgdaField{prf} \AgdaSymbol{=} \AgdaSymbol{(}\AgdaBound{g} \AgdaSymbol{:} \AgdaBound{Γ}\AgdaSymbol{)} \AgdaSymbol{→} \AgdaFunction{<} \AgdaBound{fn} \AgdaBound{g} \AgdaBound{≈h₁} \AgdaBound{gn} \AgdaBound{g} \AgdaFunction{>}\<%
\\
\>[0]\AgdaIndent{11}{}\<[11]%
\>[11]\AgdaSymbol{;} \AgdaField{Uni} \AgdaSymbol{=} \AgdaBound{ext} \AgdaSymbol{(λ} \AgdaBound{g} \AgdaSymbol{→} \AgdaFunction{Uni} \AgdaSymbol{(}\AgdaBound{fn} \AgdaBound{g} \AgdaBound{≈h₁} \AgdaBound{gn} \AgdaBound{g}\AgdaSymbol{))}\<%
\\
\>[0]\AgdaIndent{11}{}\<[11]%
\>[11]\AgdaSymbol{\}}\<%
\\
%
\\
\>\<\end{code}

Category of Setoids

\begin{code}\>\<%
\\
%
\\
%
\\
\>\AgdaFunction{Std} \AgdaSymbol{:} \AgdaRecord{Category}\<%
\\
\>\AgdaFunction{Std} \AgdaSymbol{=} \AgdaInductiveConstructor{CatC} \AgdaRecord{HSetoid} \AgdaRecord{\_⇉\_} \AgdaSymbol{(λ} \AgdaBound{\_} \AgdaSymbol{→} \AgdaFunction{idCH}\AgdaSymbol{)} \AgdaSymbol{(λ} \AgdaBound{\_} \AgdaBound{\_} \AgdaSymbol{→} \AgdaFunction{\_∘c\_}\AgdaSymbol{)} \<[51]%
\>[51]\<%
\\
\>[0]\AgdaIndent{6}{}\<[6]%
\>[6]\AgdaSymbol{(}\AgdaInductiveConstructor{IsCatC} \AgdaSymbol{(λ} \AgdaBound{α} \AgdaBound{β} \AgdaBound{f} \AgdaSymbol{→} \AgdaInductiveConstructor{PE.refl}\AgdaSymbol{)} \<[34]%
\>[34]\<%
\\
\>[0]\AgdaIndent{6}{}\<[6]%
\>[6]\AgdaSymbol{(λ} \AgdaBound{α} \AgdaBound{β} \AgdaBound{f} \AgdaSymbol{→} \AgdaInductiveConstructor{PE.refl}\AgdaSymbol{)} \<[26]%
\>[26]\<%
\\
\>[0]\AgdaIndent{6}{}\<[6]%
\>[6]\AgdaSymbol{(λ} \AgdaBound{α} \AgdaBound{δ} \AgdaBound{f} \AgdaBound{g} \AgdaBound{h} \AgdaSymbol{→} \AgdaInductiveConstructor{PE.refl}\AgdaSymbol{))}\<%
\\
%
\\
\>\<\end{code}

Terminal object

\begin{code}\>\<%
\\
%
\\
\>\AgdaFunction{●} \AgdaSymbol{:} \AgdaRecord{HSetoid}\<%
\\
\>\AgdaFunction{●} \<[4]%
\>[4]\AgdaSymbol{=} \AgdaKeyword{record} \AgdaSymbol{\{}\<%
\\
\>[0]\AgdaIndent{6}{}\<[6]%
\>[6]\AgdaField{Carrier} \AgdaSymbol{=} \AgdaRecord{⊤}\AgdaSymbol{;}\<%
\\
\>[0]\AgdaIndent{6}{}\<[6]%
\>[6]\AgdaField{\_≈h\_} \<[14]%
\>[14]\AgdaSymbol{=} \AgdaSymbol{λ} \AgdaBound{\_} \AgdaBound{\_} \AgdaSymbol{→} \AgdaFunction{⊤'}\AgdaSymbol{;}\<%
\\
\>[0]\AgdaIndent{6}{}\<[6]%
\>[6]\AgdaField{refl} \<[14]%
\>[14]\AgdaSymbol{=} \AgdaInductiveConstructor{tt}\AgdaSymbol{;}\<%
\\
\>[0]\AgdaIndent{6}{}\<[6]%
\>[6]\AgdaField{sym} \<[14]%
\>[14]\AgdaSymbol{=} \AgdaSymbol{λ} \AgdaBound{\_} \AgdaSymbol{→} \AgdaInductiveConstructor{tt}\AgdaSymbol{;}\<%
\\
\>[0]\AgdaIndent{6}{}\<[6]%
\>[6]\AgdaField{trans} \<[14]%
\>[14]\AgdaSymbol{=} \AgdaSymbol{λ} \AgdaBound{\_} \AgdaBound{\_} \AgdaSymbol{→} \AgdaInductiveConstructor{tt} \AgdaSymbol{\}}\<%
\\
%
\\
%
\\
\>\AgdaFunction{⋆} \AgdaSymbol{:} \AgdaSymbol{\{}\AgdaBound{Δ} \AgdaSymbol{:} \AgdaRecord{HSetoid}\AgdaSymbol{\}} \AgdaSymbol{→} \AgdaBound{Δ} \AgdaRecord{⇉} \AgdaFunction{●}\<%
\\
\>\AgdaFunction{⋆} \AgdaSymbol{=} \AgdaKeyword{record} \<[11]%
\>[11]\<%
\\
\>[0]\AgdaIndent{6}{}\<[6]%
\>[6]\AgdaSymbol{\{} \AgdaField{fn} \AgdaSymbol{=} \AgdaSymbol{λ} \AgdaBound{\_} \AgdaSymbol{→} \AgdaInductiveConstructor{tt}\<%
\\
\>[0]\AgdaIndent{6}{}\<[6]%
\>[6]\AgdaSymbol{;} \AgdaField{resp} \AgdaSymbol{=} \AgdaSymbol{λ} \AgdaBound{\_} \AgdaSymbol{→} \AgdaInductiveConstructor{tt} \AgdaSymbol{\}}\<%
\\
%
\\
\>\AgdaFunction{uniqueHom} \AgdaSymbol{:} \AgdaSymbol{∀} \AgdaSymbol{(}\AgdaBound{Δ} \AgdaSymbol{:} \AgdaRecord{HSetoid}\AgdaSymbol{)} \AgdaSymbol{→} \AgdaSymbol{(}\AgdaBound{f} \AgdaSymbol{:} \AgdaBound{Δ} \AgdaRecord{⇉} \AgdaFunction{●}\AgdaSymbol{)} \AgdaSymbol{→} \AgdaBound{f} \AgdaDatatype{≡} \AgdaFunction{⋆}\<%
\\
\>\AgdaFunction{uniqueHom} \AgdaBound{Δ} \AgdaBound{f} \AgdaSymbol{=} \AgdaInductiveConstructor{PE.refl}\<%
\\
%
\\
\>\<\end{code}

\section{Categories with families of setoids}\label{cwf}


A Category with families consists of a base category and a functor
\cite{clairambault2005categories}. We firstly define the categories with
families of sets in Agda  as a guidance for the one for setoids. We
would present the setoid one here since it is relevant.

\AgdaHide{
\begin{code}\>\<%
\\
%
\\
\>\AgdaKeyword{module} \AgdaModule{CwF-setoid} \AgdaKeyword{where}\<%
\\
%
\\
\>\AgdaKeyword{open} \AgdaKeyword{import} \AgdaModule{CategoryofSetoid}\<%
\\
\>\AgdaKeyword{open} \AgdaKeyword{import} \AgdaModule{Data.Product}\<%
\\
\>\AgdaKeyword{open} \AgdaKeyword{import} \AgdaModule{Function}\<%
\\
\>\AgdaKeyword{open} \AgdaKeyword{import} \AgdaModule{Data.Nat}\<%
\\
\>\AgdaKeyword{open} \AgdaKeyword{import} \AgdaModule{Data.Empty}\<%
\\
\>\AgdaKeyword{open} \AgdaKeyword{import} \AgdaModule{Data.Unit}\<%
\\
%
\\
\>\<\end{code}
}

\section{Categories with families}

Context are interpreted as setoids

\begin{code}\>\<%
\\
\>\AgdaFunction{Con} \AgdaSymbol{=} \AgdaRecord{Setoid}\<%
\\
\>\<\end{code}

Semantic Types

\begin{code}\>\<%
\\
\>\AgdaKeyword{record} \AgdaRecord{Ty} \AgdaSymbol{(}\AgdaBound{Γ} \AgdaSymbol{:} \AgdaRecord{Setoid}\AgdaSymbol{)} \AgdaSymbol{:} \AgdaPrimitiveType{Set₁} \AgdaKeyword{where}\<%
\\
\>[0]\AgdaIndent{2}{}\<[2]%
\>[2]\AgdaKeyword{field}\<%
\\
\>[2]\AgdaIndent{4}{}\<[4]%
\>[4]\AgdaField{fm} \<[11]%
\>[11]\AgdaSymbol{:} \AgdaFunction{∣} \AgdaBound{Γ} \AgdaFunction{∣} \AgdaSymbol{→} \AgdaRecord{Setoid}\<%
\\
\>[2]\AgdaIndent{4}{}\<[4]%
\>[4]\AgdaField{substT} \AgdaSymbol{:} \AgdaSymbol{\{}\AgdaBound{x} \AgdaBound{y} \AgdaSymbol{:} \AgdaFunction{∣} \AgdaBound{Γ} \AgdaFunction{∣}\AgdaSymbol{\}} \AgdaSymbol{→} \<[29]%
\>[29]\<%
\\
\>[4]\AgdaIndent{13}{}\<[13]%
\>[13]\AgdaSymbol{.(}\AgdaFunction{[} \AgdaBound{Γ} \AgdaFunction{]} \AgdaBound{x} \AgdaFunction{≈} \AgdaBound{y}\AgdaSymbol{)} \AgdaSymbol{→}\<%
\\
\>[4]\AgdaIndent{13}{}\<[13]%
\>[13]\AgdaFunction{∣} \AgdaBound{fm} \AgdaBound{x} \AgdaFunction{∣} \AgdaSymbol{→} \<[24]%
\>[24]\<%
\\
\>[4]\AgdaIndent{13}{}\<[13]%
\>[13]\AgdaFunction{∣} \AgdaBound{fm} \AgdaBound{y} \AgdaFunction{∣}\<%
\\
\>[0]\AgdaIndent{4}{}\<[4]%
\>[4]\AgdaSymbol{.}\AgdaField{subst*} \AgdaSymbol{:} \AgdaSymbol{∀\{}\AgdaBound{x} \AgdaBound{y} \AgdaSymbol{:} \AgdaFunction{∣} \AgdaBound{Γ} \AgdaFunction{∣}\AgdaSymbol{\}}\<%
\\
\>[0]\AgdaIndent{13}{}\<[13]%
\>[13]\AgdaSymbol{(}\AgdaBound{p} \AgdaSymbol{:} \AgdaSymbol{(}\AgdaFunction{[} \AgdaBound{Γ} \AgdaFunction{]} \AgdaBound{x} \AgdaFunction{≈} \AgdaBound{y}\AgdaSymbol{))}\<%
\\
\>[0]\AgdaIndent{13}{}\<[13]%
\>[13]\AgdaSymbol{\{}\AgdaBound{a} \AgdaBound{b} \AgdaSymbol{:} \AgdaFunction{∣} \AgdaBound{fm} \AgdaBound{x} \AgdaFunction{∣}\AgdaSymbol{\}} \AgdaSymbol{→}\<%
\\
\>[0]\AgdaIndent{13}{}\<[13]%
\>[13]\AgdaSymbol{.(}\AgdaFunction{[} \AgdaBound{fm} \AgdaBound{x} \AgdaFunction{]} \AgdaBound{a} \AgdaFunction{≈} \AgdaBound{b}\AgdaSymbol{)} \AgdaSymbol{→}\<%
\\
\>[0]\AgdaIndent{13}{}\<[13]%
\>[13]\AgdaSymbol{(}\AgdaFunction{[} \AgdaBound{fm} \AgdaBound{y} \AgdaFunction{]} \AgdaBound{substT} \AgdaBound{p} \AgdaBound{a} \AgdaFunction{≈} \AgdaBound{substT} \AgdaBound{p} \AgdaBound{b}\AgdaSymbol{)}\<%
\\
%
\\
\>[0]\AgdaIndent{4}{}\<[4]%
\>[4]\AgdaSymbol{.}\AgdaField{refl*} \<[12]%
\>[12]\AgdaSymbol{:} \AgdaSymbol{∀\{}\AgdaBound{x} \AgdaSymbol{:} \AgdaFunction{∣} \AgdaBound{Γ} \AgdaFunction{∣}\AgdaSymbol{\}\{}\AgdaBound{a} \AgdaSymbol{:} \AgdaFunction{∣} \AgdaBound{fm} \AgdaBound{x} \AgdaFunction{∣}\AgdaSymbol{\}} \AgdaSymbol{→} \<[43]%
\>[43]\<%
\\
\>[0]\AgdaIndent{13}{}\<[13]%
\>[13]\AgdaFunction{[} \AgdaBound{fm} \AgdaBound{x} \AgdaFunction{]} \AgdaBound{substT} \AgdaSymbol{(}\AgdaFunction{[} \AgdaBound{Γ} \AgdaFunction{]refl}\AgdaSymbol{)} \AgdaBound{a} \AgdaFunction{≈} \AgdaBound{a}\<%
\\
\>[0]\AgdaIndent{4}{}\<[4]%
\>[4]\AgdaSymbol{.}\AgdaField{trans*} \AgdaSymbol{:} \AgdaSymbol{∀\{}\AgdaBound{x} \AgdaBound{y} \AgdaBound{z} \AgdaSymbol{:} \AgdaFunction{∣} \AgdaBound{Γ} \AgdaFunction{∣}\AgdaSymbol{\}}\<%
\\
\>[0]\AgdaIndent{13}{}\<[13]%
\>[13]\AgdaSymbol{\{}\AgdaBound{p} \AgdaSymbol{:} \AgdaFunction{[} \AgdaBound{Γ} \AgdaFunction{]} \AgdaBound{x} \AgdaFunction{≈} \AgdaBound{y}\AgdaSymbol{\}}\<%
\\
\>[0]\AgdaIndent{13}{}\<[13]%
\>[13]\AgdaSymbol{\{}\AgdaBound{q} \AgdaSymbol{:} \AgdaFunction{[} \AgdaBound{Γ} \AgdaFunction{]} \AgdaBound{y} \AgdaFunction{≈} \AgdaBound{z}\AgdaSymbol{\}}\<%
\\
\>[0]\AgdaIndent{13}{}\<[13]%
\>[13]\AgdaSymbol{(}\AgdaBound{a} \AgdaSymbol{:} \AgdaFunction{∣} \AgdaBound{fm} \AgdaBound{x} \AgdaFunction{∣}\AgdaSymbol{)} \AgdaSymbol{→} \<[30]%
\>[30]\<%
\\
\>[0]\AgdaIndent{13}{}\<[13]%
\>[13]\AgdaFunction{[} \AgdaBound{fm} \AgdaBound{z} \AgdaFunction{]} \AgdaBound{substT} \AgdaBound{q} \AgdaSymbol{(}\AgdaBound{substT} \AgdaBound{p} \AgdaBound{a}\AgdaSymbol{)} \<[44]%
\>[44]\<%
\\
\>[0]\AgdaIndent{12}{}\<[12]%
\>[12]\AgdaFunction{≈} \AgdaBound{substT} \AgdaSymbol{(}\AgdaFunction{[} \AgdaBound{Γ} \AgdaFunction{]trans} \AgdaBound{p} \AgdaBound{q}\AgdaSymbol{)} \AgdaBound{a}\<%
\\
%
\\
\>[0]\AgdaIndent{2}{}\<[2]%
\>[2]\AgdaSymbol{.}\AgdaFunction{tr*} \AgdaSymbol{:} \AgdaSymbol{∀\{}\AgdaBound{x} \AgdaBound{y} \AgdaSymbol{:} \AgdaFunction{∣} \AgdaBound{Γ} \AgdaFunction{∣}\AgdaSymbol{\}}\<%
\\
\>[2]\AgdaIndent{10}{}\<[10]%
\>[10]\AgdaSymbol{\{}\AgdaBound{p} \AgdaSymbol{:} \AgdaFunction{[} \AgdaBound{Γ} \AgdaFunction{]} \AgdaBound{y} \AgdaFunction{≈} \AgdaBound{x}\AgdaSymbol{\}}\<%
\\
\>[2]\AgdaIndent{10}{}\<[10]%
\>[10]\AgdaSymbol{\{}\AgdaBound{q} \AgdaSymbol{:} \AgdaFunction{[} \AgdaBound{Γ} \AgdaFunction{]} \AgdaBound{x} \AgdaFunction{≈} \AgdaBound{y}\AgdaSymbol{\}}\<%
\\
\>[2]\AgdaIndent{10}{}\<[10]%
\>[10]\AgdaSymbol{\{}\AgdaBound{a} \AgdaSymbol{:} \AgdaFunction{∣} \AgdaFunction{fm} \AgdaBound{x} \AgdaFunction{∣}\AgdaSymbol{\}} \AgdaSymbol{→} \<[27]%
\>[27]\<%
\\
\>[2]\AgdaIndent{10}{}\<[10]%
\>[10]\AgdaFunction{[} \AgdaFunction{fm} \AgdaBound{x} \AgdaFunction{]} \AgdaFunction{substT} \AgdaBound{p} \AgdaSymbol{(}\AgdaFunction{substT} \AgdaBound{q} \AgdaBound{a}\AgdaSymbol{)} \AgdaFunction{≈} \AgdaBound{a}\<%
\\
\>[0]\AgdaIndent{2}{}\<[2]%
\>[2]\AgdaFunction{tr*} \AgdaSymbol{=} \AgdaFunction{[} \AgdaFunction{fm} \AgdaSymbol{\_} \AgdaFunction{]trans} \AgdaSymbol{(}\AgdaFunction{trans*} \AgdaSymbol{\_)} \AgdaFunction{refl*}\<%
\\
%
\\
\>[0]\AgdaIndent{2}{}\<[2]%
\>[2]\AgdaFunction{substT-inv} \AgdaSymbol{:} \AgdaSymbol{\{}\AgdaBound{x} \AgdaBound{y} \AgdaSymbol{:} \AgdaFunction{∣} \AgdaBound{Γ} \AgdaFunction{∣}\AgdaSymbol{\}} \AgdaSymbol{→} \<[31]%
\>[31]\<%
\\
\>[2]\AgdaIndent{13}{}\<[13]%
\>[13]\AgdaSymbol{.(}\AgdaFunction{[} \AgdaBound{Γ} \AgdaFunction{]} \AgdaBound{x} \AgdaFunction{≈} \AgdaBound{y}\AgdaSymbol{)} \AgdaSymbol{→}\<%
\\
\>[2]\AgdaIndent{13}{}\<[13]%
\>[13]\AgdaFunction{∣} \AgdaFunction{fm} \AgdaBound{y} \AgdaFunction{∣} \AgdaSymbol{→} \<[24]%
\>[24]\<%
\\
\>[2]\AgdaIndent{13}{}\<[13]%
\>[13]\AgdaFunction{∣} \AgdaFunction{fm} \AgdaBound{x} \AgdaFunction{∣}\<%
\\
\>[0]\AgdaIndent{2}{}\<[2]%
\>[2]\AgdaFunction{substT-inv} \AgdaBound{p} \AgdaBound{y} \AgdaSymbol{=} \AgdaFunction{substT} \AgdaSymbol{(}\AgdaFunction{[} \AgdaBound{Γ} \AgdaFunction{]sym} \AgdaBound{p}\AgdaSymbol{)} \AgdaBound{y}\<%
\\
\>\<\end{code}

\AgdaHide{
\begin{code}\>\<%
\\
\>[0]\AgdaIndent{2}{}\<[2]%
\>[2]\AgdaFunction{subst-mir1} \AgdaSymbol{:} \AgdaSymbol{∀\{}\AgdaBound{x} \AgdaBound{y} \AgdaSymbol{:} \AgdaFunction{∣} \AgdaBound{Γ} \AgdaFunction{∣}\AgdaSymbol{\}\{}\AgdaBound{P} \AgdaSymbol{:} \AgdaPrimitiveType{Set}\AgdaSymbol{\}}\<%
\\
\>[2]\AgdaIndent{14}{}\<[14]%
\>[14]\AgdaSymbol{\{}\AgdaBound{p} \AgdaSymbol{:} \AgdaFunction{[} \AgdaBound{Γ} \AgdaFunction{]} \AgdaBound{x} \AgdaFunction{≈} \AgdaBound{y}\AgdaSymbol{\}\{}\AgdaBound{q} \AgdaSymbol{:} \AgdaFunction{[} \AgdaBound{Γ} \AgdaFunction{]} \AgdaBound{y} \AgdaFunction{≈} \AgdaBound{x}\AgdaSymbol{\}}\<%
\\
\>[2]\AgdaIndent{14}{}\<[14]%
\>[14]\AgdaSymbol{\{}\AgdaBound{a} \AgdaSymbol{:} \AgdaFunction{∣} \AgdaFunction{fm} \AgdaBound{x} \AgdaFunction{∣}\AgdaSymbol{\}\{}\AgdaBound{b} \AgdaSymbol{:} \AgdaFunction{∣} \AgdaFunction{fm} \AgdaBound{y} \AgdaFunction{∣}\AgdaSymbol{\}} \AgdaSymbol{→} \<[45]%
\>[45]\<%
\\
\>[2]\AgdaIndent{14}{}\<[14]%
\>[14]\AgdaSymbol{(.(}\AgdaFunction{[} \AgdaFunction{fm} \AgdaBound{y} \AgdaFunction{]} \AgdaFunction{substT} \AgdaBound{p} \AgdaBound{a} \AgdaFunction{≈} \AgdaBound{b}\AgdaSymbol{)} \AgdaSymbol{→} \AgdaBound{P}\AgdaSymbol{)} \<[47]%
\>[47]\<%
\\
\>[-1]\AgdaIndent{13}{}\<[13]%
\>[13]\AgdaSymbol{→} \AgdaSymbol{(.(}\AgdaFunction{[} \AgdaFunction{fm} \AgdaBound{x} \AgdaFunction{]} \AgdaBound{a} \AgdaFunction{≈} \AgdaFunction{substT} \AgdaBound{q} \AgdaBound{b}\AgdaSymbol{)} \AgdaSymbol{→} \AgdaBound{P}\AgdaSymbol{)}\<%
\\
\>[0]\AgdaIndent{2}{}\<[2]%
\>[2]\AgdaFunction{subst-mir1} \AgdaBound{eq} \AgdaBound{p} \AgdaSymbol{=} \AgdaBound{eq} \AgdaSymbol{(}\AgdaFunction{[} \AgdaFunction{fm} \AgdaSymbol{\_} \AgdaFunction{]trans} \AgdaSymbol{(}\AgdaFunction{subst*} \AgdaSymbol{\_} \AgdaBound{p}\AgdaSymbol{)} \AgdaFunction{tr*}\AgdaSymbol{)}\<%
\\
%
\\
\>\AgdaKeyword{open} \AgdaModule{Ty} \AgdaKeyword{public} \<[15]%
\>[15]\<%
\\
\>[0]\AgdaIndent{2}{}\<[2]%
\>[2]\AgdaKeyword{renaming} \AgdaSymbol{(}substT \AgdaSymbol{to} [\_]subst\AgdaSymbol{;} substT-inv \AgdaSymbol{to} [\_]subst-inv
           \AgdaSymbol{;} subst* \AgdaSymbol{to} [\_]subst*\AgdaSymbol{;} fm \AgdaSymbol{to} [\_]fm \AgdaSymbol{;}
            refl* \AgdaSymbol{to} [\_]refl* \AgdaSymbol{;} trans* \AgdaSymbol{to} [\_]trans*\AgdaSymbol{;} tr* \AgdaSymbol{to} [\_]tr*\AgdaSymbol{)}\<%
\\
%
\\
\>\<\end{code}
}

Type substitution

\begin{code}\>\<%
\\
\>\AgdaFunction{\_[\_]T} \AgdaSymbol{:} \AgdaSymbol{∀} \AgdaSymbol{\{}\AgdaBound{Γ} \AgdaBound{Δ} \AgdaSymbol{:} \AgdaRecord{Setoid}\AgdaSymbol{\}} \AgdaSymbol{→} \AgdaRecord{Ty} \AgdaBound{Δ} \AgdaSymbol{→} \AgdaBound{Γ} \AgdaRecord{⇉} \AgdaBound{Δ} \AgdaSymbol{→} \AgdaRecord{Ty} \AgdaBound{Γ}\<%
\\
\>\AgdaFunction{\_[\_]T} \AgdaSymbol{\{}\AgdaBound{Γ}\AgdaSymbol{\}} \AgdaSymbol{\{}\AgdaBound{Δ}\AgdaSymbol{\}} \AgdaBound{A} \AgdaBound{f}\<%
\\
\>[2]\AgdaIndent{5}{}\<[5]%
\>[5]\AgdaSymbol{=} \AgdaKeyword{record}\<%
\\
\>[2]\AgdaIndent{5}{}\<[5]%
\>[5]\AgdaSymbol{\{} \AgdaField{fm} \<[14]%
\>[14]\AgdaSymbol{=} \AgdaSymbol{λ} \AgdaBound{x} \AgdaSymbol{→} \AgdaFunction{fm} \AgdaSymbol{(}\AgdaFunction{fn} \AgdaBound{x}\AgdaSymbol{)}\<%
\\
\>[2]\AgdaIndent{5}{}\<[5]%
\>[5]\AgdaSymbol{;} \AgdaField{substT} \AgdaSymbol{=} \AgdaSymbol{λ} \AgdaBound{p} \AgdaSymbol{→} \AgdaFunction{substT} \AgdaSymbol{\_}\<%
\\
\>[2]\AgdaIndent{5}{}\<[5]%
\>[5]\AgdaSymbol{;} \AgdaField{subst*} \AgdaSymbol{=} \AgdaSymbol{λ} \AgdaBound{p} \AgdaSymbol{→} \AgdaFunction{subst*} \AgdaSymbol{(}\AgdaFunction{resp} \AgdaBound{p}\AgdaSymbol{)}\<%
\\
\>[2]\AgdaIndent{5}{}\<[5]%
\>[5]\AgdaSymbol{;} \AgdaField{refl*} \<[14]%
\>[14]\AgdaSymbol{=} \AgdaFunction{refl*}\<%
\\
\>[2]\AgdaIndent{5}{}\<[5]%
\>[5]\AgdaSymbol{;} \AgdaField{trans*} \AgdaSymbol{=} \AgdaFunction{trans*}\<%
\\
\>[2]\AgdaIndent{5}{}\<[5]%
\>[5]\AgdaSymbol{\}}\<%
\\
\>[2]\AgdaIndent{5}{}\<[5]%
\>[5]\AgdaKeyword{where} \<[11]%
\>[11]\<%
\\
\>[5]\AgdaIndent{7}{}\<[7]%
\>[7]\AgdaKeyword{open} \AgdaModule{Ty} \AgdaBound{A}\<%
\\
\>[5]\AgdaIndent{7}{}\<[7]%
\>[7]\AgdaKeyword{open} \AgdaModule{\_⇉\_} \AgdaBound{f}\<%
\\
%
\\
%
\\
\>\<\end{code}


Semantic Terms

\begin{code}\>\<%
\\
\>\AgdaKeyword{record} \AgdaRecord{Tm} \AgdaSymbol{\{}\AgdaBound{Γ} \AgdaSymbol{:} \AgdaFunction{Con}\AgdaSymbol{\}(}\AgdaBound{A} \AgdaSymbol{:} \AgdaRecord{Ty} \AgdaBound{Γ}\AgdaSymbol{)} \AgdaSymbol{:} \AgdaPrimitiveType{Set} \AgdaKeyword{where}\<%
\\
\>[0]\AgdaIndent{2}{}\<[2]%
\>[2]\AgdaKeyword{constructor} \AgdaInductiveConstructor{tm:\_resp:\_}\<%
\\
\>[0]\AgdaIndent{2}{}\<[2]%
\>[2]\AgdaKeyword{field}\<%
\\
\>[2]\AgdaIndent{4}{}\<[4]%
\>[4]\AgdaField{tm} \<[10]%
\>[10]\AgdaSymbol{:} \AgdaSymbol{(}\AgdaBound{x} \AgdaSymbol{:} \AgdaFunction{∣} \AgdaBound{Γ} \AgdaFunction{∣}\AgdaSymbol{)} \AgdaSymbol{→} \AgdaFunction{∣} \AgdaFunction{[} \AgdaBound{A} \AgdaFunction{]fm} \AgdaBound{x} \AgdaFunction{∣}\<%
\\
\>[2]\AgdaIndent{4}{}\<[4]%
\>[4]\AgdaSymbol{.}\AgdaField{respt} \AgdaSymbol{:} \AgdaSymbol{∀} \AgdaSymbol{\{}\AgdaBound{x} \AgdaBound{y} \AgdaSymbol{:} \AgdaFunction{∣} \AgdaBound{Γ} \AgdaFunction{∣}\AgdaSymbol{\}} \AgdaSymbol{→} \<[31]%
\>[31]\<%
\\
\>[4]\AgdaIndent{10}{}\<[10]%
\>[10]\AgdaSymbol{(}\AgdaBound{p} \AgdaSymbol{:} \AgdaFunction{[} \AgdaBound{Γ} \AgdaFunction{]} \AgdaBound{x} \AgdaFunction{≈} \AgdaBound{y}\AgdaSymbol{)} \AgdaSymbol{→} \<[30]%
\>[30]\<%
\\
\>[4]\AgdaIndent{10}{}\<[10]%
\>[10]\AgdaFunction{[} \AgdaFunction{[} \AgdaBound{A} \AgdaFunction{]fm} \AgdaBound{y} \AgdaFunction{]} \AgdaFunction{[} \AgdaBound{A} \AgdaFunction{]subst} \AgdaBound{p} \AgdaSymbol{(}\AgdaBound{tm} \AgdaBound{x}\AgdaSymbol{)} \AgdaFunction{≈} \AgdaBound{tm} \AgdaBound{y}\<%
\\
%
\\
\>\AgdaKeyword{open} \AgdaModule{Tm} \AgdaKeyword{public} \AgdaKeyword{renaming} \AgdaSymbol{(}tm \AgdaSymbol{to} [\_]tm \AgdaSymbol{;} respt \AgdaSymbol{to} [\_]respt\AgdaSymbol{)}\<%
\\
\>\<\end{code}

Term substitution

\begin{code}\>\<%
\\
\>\AgdaFunction{\_[\_]m} \AgdaSymbol{:} \AgdaSymbol{∀} \AgdaSymbol{\{}\AgdaBound{Γ} \AgdaBound{Δ} \AgdaSymbol{:} \AgdaFunction{Con}\AgdaSymbol{\}\{}\AgdaBound{A} \AgdaSymbol{:} \AgdaRecord{Ty} \AgdaBound{Δ}\AgdaSymbol{\}} \AgdaSymbol{→} \AgdaRecord{Tm} \AgdaBound{A} \<[39]%
\>[39]\<%
\\
\>[0]\AgdaIndent{6}{}\<[6]%
\>[6]\AgdaSymbol{→} \AgdaSymbol{(}\AgdaBound{f} \AgdaSymbol{:} \AgdaBound{Γ} \AgdaRecord{⇉} \AgdaBound{Δ}\AgdaSymbol{)} \AgdaSymbol{→} \AgdaRecord{Tm} \AgdaSymbol{(}\AgdaBound{A} \AgdaFunction{[} \AgdaBound{f} \AgdaFunction{]T}\AgdaSymbol{)}\<%
\\
\>\AgdaFunction{\_[\_]m} \AgdaBound{t} \AgdaBound{f} \AgdaSymbol{=} \AgdaKeyword{record} \<[19]%
\>[19]\<%
\\
\>[0]\AgdaIndent{10}{}\<[10]%
\>[10]\AgdaSymbol{\{} \AgdaField{tm} \AgdaSymbol{=} \AgdaFunction{[} \AgdaBound{t} \AgdaFunction{]tm} \AgdaFunction{∘} \AgdaFunction{[} \AgdaBound{f} \AgdaFunction{]fn}\<%
\\
\>[0]\AgdaIndent{10}{}\<[10]%
\>[10]\AgdaSymbol{;} \AgdaField{respt} \AgdaSymbol{=} \AgdaFunction{[} \AgdaBound{t} \AgdaFunction{]respt} \AgdaFunction{∘} \AgdaFunction{[} \AgdaBound{f} \AgdaFunction{]resp} \<[43]%
\>[43]\<%
\\
\>[0]\AgdaIndent{10}{}\<[10]%
\>[10]\AgdaSymbol{\}}\<%
\\
\>\<\end{code}

Context comprehension

\begin{code}\>\<%
\\
\>\AgdaFunction{\_\&\_} \AgdaSymbol{:} \AgdaSymbol{(}\AgdaBound{Γ} \AgdaSymbol{:} \AgdaRecord{Setoid}\AgdaSymbol{)} \AgdaSymbol{→} \AgdaRecord{Ty} \AgdaBound{Γ} \AgdaSymbol{→} \AgdaRecord{Setoid}\<%
\\
\>\AgdaBound{Γ} \AgdaFunction{\&} \AgdaBound{A} \AgdaSymbol{=} \<[8]%
\>[8]\<%
\\
\>[0]\AgdaIndent{2}{}\<[2]%
\>[2]\AgdaKeyword{record} \AgdaSymbol{\{} \AgdaField{Carrier} \AgdaSymbol{=} \AgdaRecord{Σ[} \AgdaBound{x} \AgdaRecord{∶} \AgdaFunction{∣} \AgdaBound{Γ} \AgdaFunction{∣} \AgdaRecord{]} \AgdaFunction{∣} \AgdaFunction{fm} \AgdaBound{x} \AgdaFunction{∣} \<[45]%
\>[45]\<%
\\
\>[0]\AgdaIndent{9}{}\<[9]%
\>[9]\AgdaSymbol{;} \AgdaField{\_≈\_} \AgdaSymbol{=} \AgdaSymbol{λ\{(}\AgdaBound{x} \AgdaInductiveConstructor{,} \AgdaBound{a}\AgdaSymbol{)} \AgdaSymbol{(}\AgdaBound{y} \AgdaInductiveConstructor{,} \AgdaBound{b}\AgdaSymbol{)} \AgdaSymbol{→} \<[37]%
\>[37]\<%
\\
\>[9]\AgdaIndent{13}{}\<[13]%
\>[13]\AgdaRecord{Σ[} \AgdaBound{p} \AgdaRecord{∶} \AgdaBound{x} \AgdaFunction{≈} \AgdaBound{y} \AgdaRecord{]} \AgdaFunction{[} \AgdaFunction{fm} \AgdaBound{y} \AgdaFunction{]} \AgdaSymbol{(}\AgdaFunction{substT} \AgdaBound{p} \AgdaBound{a}\AgdaSymbol{)} \AgdaFunction{≈} \AgdaBound{b} \AgdaSymbol{\}}\<%
\\
\>[0]\AgdaIndent{9}{}\<[9]%
\>[9]\AgdaSymbol{;} \AgdaField{refl} \AgdaSymbol{=} \AgdaFunction{refl} \AgdaInductiveConstructor{,} \AgdaFunction{refl*}\<%
\\
\>[0]\AgdaIndent{9}{}\<[9]%
\>[9]\AgdaSymbol{;} \AgdaField{sym} \AgdaSymbol{=} \<[18]%
\>[18]\AgdaSymbol{λ} \AgdaSymbol{\{(}\AgdaBound{p} \AgdaInductiveConstructor{,} \AgdaBound{q}\AgdaSymbol{)} \AgdaSymbol{→} \AgdaSymbol{(}\AgdaFunction{sym} \AgdaBound{p}\AgdaSymbol{)} \AgdaInductiveConstructor{,} \<[41]%
\>[41]\<%
\\
\>[9]\AgdaIndent{13}{}\<[13]%
\>[13]\AgdaFunction{[} \AgdaFunction{fm} \AgdaSymbol{\_} \AgdaFunction{]trans} \AgdaSymbol{(}\AgdaFunction{subst*} \AgdaSymbol{\_} \AgdaSymbol{(}\AgdaFunction{[} \AgdaFunction{fm} \AgdaSymbol{\_} \AgdaFunction{]sym} \AgdaBound{q}\AgdaSymbol{))} \AgdaFunction{tr*} \AgdaSymbol{\}}\<%
\\
\>[0]\AgdaIndent{9}{}\<[9]%
\>[9]\AgdaSymbol{;} \AgdaField{trans} \AgdaSymbol{=} \AgdaSymbol{λ} \AgdaSymbol{\{(}\AgdaBound{p} \AgdaInductiveConstructor{,} \AgdaBound{q}\AgdaSymbol{)} \AgdaSymbol{(}\AgdaBound{m} \AgdaInductiveConstructor{,} \AgdaBound{n}\AgdaSymbol{)} \AgdaSymbol{→} \AgdaFunction{trans} \AgdaBound{p} \AgdaBound{m} \AgdaInductiveConstructor{,}\<%
\\
\>[0]\AgdaIndent{19}{}\<[19]%
\>[19]\AgdaFunction{[} \AgdaFunction{fm} \AgdaSymbol{\_} \AgdaFunction{]trans} \AgdaSymbol{(}\AgdaFunction{[} \AgdaFunction{fm} \AgdaSymbol{\_} \AgdaFunction{]trans} \<[48]%
\>[48]\<%
\\
\>[0]\AgdaIndent{19}{}\<[19]%
\>[19]\AgdaSymbol{(}\AgdaFunction{[} \AgdaFunction{fm} \AgdaSymbol{\_} \AgdaFunction{]sym} \AgdaSymbol{(}\AgdaFunction{trans*} \AgdaSymbol{\_))} \AgdaSymbol{(}\AgdaFunction{subst*} \AgdaSymbol{\_} \AgdaBound{q}\AgdaSymbol{))} \AgdaBound{n}\AgdaSymbol{\}}\<%
\\
\>[0]\AgdaIndent{9}{}\<[9]%
\>[9]\AgdaSymbol{\}}\<%
\\
\>[0]\AgdaIndent{4}{}\<[4]%
\>[4]\AgdaKeyword{where}\<%
\\
\>[4]\AgdaIndent{6}{}\<[6]%
\>[6]\AgdaKeyword{open} \AgdaModule{Setoid} \AgdaBound{Γ}\<%
\\
\>[4]\AgdaIndent{6}{}\<[6]%
\>[6]\AgdaKeyword{open} \AgdaModule{Ty} \AgdaBound{A}\<%
\\
%
\\
%
\\
\>\AgdaKeyword{infixl} \AgdaNumber{5} \_\&\_\<%
\\
%
\\
\>\AgdaFunction{fst\&} \AgdaSymbol{:} \AgdaSymbol{\{}\AgdaBound{Γ} \AgdaSymbol{:} \AgdaFunction{Con}\AgdaSymbol{\}\{}\AgdaBound{A} \AgdaSymbol{:} \AgdaRecord{Ty} \AgdaBound{Γ}\AgdaSymbol{\}} \AgdaSymbol{→} \AgdaBound{Γ} \AgdaFunction{\&} \AgdaBound{A} \AgdaRecord{⇉} \AgdaBound{Γ}\<%
\\
\>\AgdaFunction{fst\&} \AgdaSymbol{=} \AgdaKeyword{record} \<[14]%
\>[14]\<%
\\
\>[6]\AgdaIndent{9}{}\<[9]%
\>[9]\AgdaSymbol{\{} \AgdaField{fn} \AgdaSymbol{=} \AgdaFunction{proj₁}\<%
\\
\>[6]\AgdaIndent{9}{}\<[9]%
\>[9]\AgdaSymbol{;} \AgdaField{resp} \AgdaSymbol{=} \AgdaFunction{proj₁}\<%
\\
\>[6]\AgdaIndent{9}{}\<[9]%
\>[9]\AgdaSymbol{\}}\<%
\\
\>\<\end{code}

Pairing operation

\begin{code}\>\<%
\\
\>\AgdaFunction{\_,,\_} \AgdaSymbol{:} \AgdaSymbol{\{}\AgdaBound{Γ} \AgdaBound{Δ} \AgdaSymbol{:} \AgdaFunction{Con}\AgdaSymbol{\}\{}\AgdaBound{A} \AgdaSymbol{:} \AgdaRecord{Ty} \AgdaBound{Δ}\AgdaSymbol{\}(}\AgdaBound{f} \AgdaSymbol{:} \AgdaBound{Γ} \AgdaRecord{⇉} \AgdaBound{Δ}\AgdaSymbol{)} \<[40]%
\>[40]\<%
\\
\>[-2]\AgdaIndent{5}{}\<[5]%
\>[5]\AgdaSymbol{→} \AgdaSymbol{(}\AgdaRecord{Tm} \AgdaSymbol{(}\AgdaBound{A} \AgdaFunction{[} \AgdaBound{f} \AgdaFunction{]T}\AgdaSymbol{))} \AgdaSymbol{→} \AgdaBound{Γ} \AgdaRecord{⇉} \AgdaSymbol{(}\AgdaBound{Δ} \AgdaFunction{\&} \AgdaBound{A}\AgdaSymbol{)}\<%
\\
\>\AgdaBound{f} \AgdaFunction{,,} \AgdaBound{t} \AgdaSymbol{=} \AgdaKeyword{record} \<[16]%
\>[16]\<%
\\
\>[0]\AgdaIndent{9}{}\<[9]%
\>[9]\AgdaSymbol{\{} \AgdaField{fn} \AgdaSymbol{=} \AgdaFunction{<} \AgdaFunction{[} \AgdaBound{f} \AgdaFunction{]fn} \AgdaFunction{,} \AgdaFunction{[} \AgdaBound{t} \AgdaFunction{]tm} \AgdaFunction{>}\<%
\\
\>[0]\AgdaIndent{9}{}\<[9]%
\>[9]\AgdaSymbol{;} \AgdaField{resp} \AgdaSymbol{=} \AgdaFunction{<} \AgdaFunction{[} \AgdaBound{f} \AgdaFunction{]resp} \AgdaFunction{,} \AgdaFunction{[} \AgdaBound{t} \AgdaFunction{]respt} \AgdaFunction{>}\<%
\\
\>[0]\AgdaIndent{9}{}\<[9]%
\>[9]\AgdaSymbol{\}}\<%
\\
\>\<\end{code}

Projections

\begin{code}\>\<%
\\
\>\AgdaFunction{fst} \AgdaSymbol{:} \AgdaSymbol{\{}\AgdaBound{Γ} \AgdaBound{Δ} \AgdaSymbol{:} \AgdaFunction{Con}\AgdaSymbol{\}\{}\AgdaBound{A} \AgdaSymbol{:} \AgdaRecord{Ty} \AgdaBound{Δ}\AgdaSymbol{\}} \AgdaSymbol{→} \AgdaBound{Γ} \AgdaRecord{⇉} \AgdaSymbol{(}\AgdaBound{Δ} \AgdaFunction{\&} \AgdaBound{A}\AgdaSymbol{)} \AgdaSymbol{→} \AgdaBound{Γ} \AgdaRecord{⇉} \AgdaBound{Δ}\<%
\\
\>\AgdaFunction{fst} \AgdaBound{f} \AgdaSymbol{=} \AgdaKeyword{record} \<[15]%
\>[15]\<%
\\
\>[0]\AgdaIndent{8}{}\<[8]%
\>[8]\AgdaSymbol{\{} \AgdaField{fn} \AgdaSymbol{=} \AgdaFunction{proj₁} \AgdaFunction{∘} \AgdaFunction{[} \AgdaBound{f} \AgdaFunction{]fn}\<%
\\
\>[0]\AgdaIndent{8}{}\<[8]%
\>[8]\AgdaSymbol{;} \AgdaField{resp} \AgdaSymbol{=} \AgdaFunction{proj₁} \AgdaFunction{∘} \AgdaFunction{[} \AgdaBound{f} \AgdaFunction{]resp} \<[35]%
\>[35]\<%
\\
\>[0]\AgdaIndent{8}{}\<[8]%
\>[8]\AgdaSymbol{\}}\<%
\\
%
\\
\>\AgdaFunction{snd} \AgdaSymbol{:} \AgdaSymbol{\{}\AgdaBound{Γ} \AgdaBound{Δ} \AgdaSymbol{:} \AgdaFunction{Con}\AgdaSymbol{\}\{}\AgdaBound{A} \AgdaSymbol{:} \AgdaRecord{Ty} \AgdaBound{Δ}\AgdaSymbol{\}} \AgdaSymbol{→} \AgdaSymbol{(}\AgdaBound{f} \AgdaSymbol{:} \AgdaBound{Γ} \AgdaRecord{⇉} \AgdaSymbol{(}\AgdaBound{Δ} \AgdaFunction{\&} \AgdaBound{A}\AgdaSymbol{))}\<%
\\
\>[0]\AgdaIndent{4}{}\<[4]%
\>[4]\AgdaSymbol{→} \AgdaRecord{Tm} \AgdaSymbol{(}\AgdaBound{A} \AgdaFunction{[} \AgdaFunction{fst} \AgdaSymbol{\{}A \AgdaSymbol{=} \AgdaBound{A}\AgdaSymbol{\}} \AgdaBound{f} \AgdaFunction{]T}\AgdaSymbol{)}\<%
\\
\>\AgdaFunction{snd} \AgdaBound{f} \AgdaSymbol{=} \AgdaKeyword{record} \<[15]%
\>[15]\<%
\\
\>[0]\AgdaIndent{8}{}\<[8]%
\>[8]\AgdaSymbol{\{} \AgdaField{tm} \AgdaSymbol{=} \AgdaFunction{proj₂} \AgdaFunction{∘} \AgdaFunction{[} \AgdaBound{f} \AgdaFunction{]fn}\<%
\\
\>[0]\AgdaIndent{8}{}\<[8]%
\>[8]\AgdaSymbol{;} \AgdaField{respt} \AgdaSymbol{=} \AgdaFunction{proj₂} \AgdaFunction{∘} \AgdaFunction{[} \AgdaBound{f} \AgdaFunction{]resp} \<[36]%
\>[36]\<%
\\
\>[0]\AgdaIndent{8}{}\<[8]%
\>[8]\AgdaSymbol{\}}\<%
\\
%
\\
\>\AgdaFunction{\_\textasciicircum\_} \AgdaSymbol{:} \AgdaSymbol{\{}\AgdaBound{Γ} \AgdaBound{Δ} \AgdaSymbol{:} \AgdaFunction{Con}\AgdaSymbol{\}(}\AgdaBound{f} \AgdaSymbol{:} \AgdaBound{Γ} \AgdaRecord{⇉} \AgdaBound{Δ}\AgdaSymbol{)(}\AgdaBound{A} \AgdaSymbol{:} \AgdaRecord{Ty} \AgdaBound{Δ}\AgdaSymbol{)} \<[39]%
\>[39]\<%
\\
\>[0]\AgdaIndent{4}{}\<[4]%
\>[4]\AgdaSymbol{→} \AgdaBound{Γ} \AgdaFunction{\&} \AgdaBound{A} \AgdaFunction{[} \AgdaBound{f} \AgdaFunction{]T} \AgdaRecord{⇉} \AgdaBound{Δ} \AgdaFunction{\&} \AgdaBound{A}\<%
\\
\>\AgdaBound{f} \AgdaFunction{\textasciicircum} \AgdaBound{A} \AgdaSymbol{=} \AgdaKeyword{record} \<[15]%
\>[15]\<%
\\
\>[0]\AgdaIndent{8}{}\<[8]%
\>[8]\AgdaSymbol{\{} \AgdaField{fn} \AgdaSymbol{=} \AgdaFunction{<} \AgdaFunction{[} \AgdaBound{f} \AgdaFunction{]fn} \AgdaFunction{∘} \AgdaFunction{proj₁} \AgdaFunction{,} \AgdaFunction{proj₂} \AgdaFunction{>}\<%
\\
\>[0]\AgdaIndent{8}{}\<[8]%
\>[8]\AgdaSymbol{;} \AgdaField{resp} \AgdaSymbol{=} \AgdaFunction{<} \AgdaFunction{[} \AgdaBound{f} \AgdaFunction{]resp} \AgdaFunction{∘} \AgdaFunction{proj₁} \AgdaFunction{,} \AgdaFunction{proj₂} \AgdaFunction{>}\<%
\\
\>[0]\AgdaIndent{8}{}\<[8]%
\>[8]\AgdaSymbol{\}}\<%
\\
\>\<\end{code}

$\Pi$-types (object level)

\begin{code}\>\<%
\\
\>\AgdaFunction{Π} \AgdaSymbol{:} \AgdaSymbol{\{}\AgdaBound{Γ} \AgdaSymbol{:} \AgdaRecord{Setoid}\AgdaSymbol{\}(}\AgdaBound{A} \AgdaSymbol{:} \AgdaRecord{Ty} \AgdaBound{Γ}\AgdaSymbol{)(}\AgdaBound{B} \AgdaSymbol{:} \AgdaRecord{Ty} \AgdaSymbol{(}\AgdaBound{Γ} \AgdaFunction{\&} \AgdaBound{A}\AgdaSymbol{))} \AgdaSymbol{→} \AgdaRecord{Ty} \AgdaBound{Γ}\<%
\\
\>\AgdaFunction{Π} \AgdaSymbol{\{}\AgdaBound{Γ}\AgdaSymbol{\}} \AgdaBound{A} \AgdaBound{B} \AgdaSymbol{=} \AgdaKeyword{record} \<[19]%
\>[19]\<%
\\
\>[0]\AgdaIndent{2}{}\<[2]%
\>[2]\AgdaSymbol{\{} \AgdaField{fm} \AgdaSymbol{=} \AgdaSymbol{λ} \AgdaBound{x} \AgdaSymbol{→} \AgdaKeyword{let} \AgdaBound{Ax} \AgdaSymbol{=} \AgdaFunction{[} \AgdaBound{A} \AgdaFunction{]fm} \AgdaBound{x} \AgdaKeyword{in}\<%
\\
\>[0]\AgdaIndent{15}{}\<[15]%
\>[15]\AgdaKeyword{let} \AgdaBound{Bx} \AgdaSymbol{=} \AgdaSymbol{λ} \AgdaBound{a} \AgdaSymbol{→} \AgdaFunction{[} \AgdaBound{B} \AgdaFunction{]fm} \AgdaSymbol{(}\AgdaBound{x} \AgdaInductiveConstructor{,} \AgdaBound{a}\AgdaSymbol{)} \AgdaKeyword{in}\<%
\\
\>[0]\AgdaIndent{9}{}\<[9]%
\>[9]\AgdaKeyword{record}\<%
\\
\>[0]\AgdaIndent{9}{}\<[9]%
\>[9]\AgdaSymbol{\{} \AgdaField{Carrier} \AgdaSymbol{=} \AgdaRecord{Subset} \AgdaSymbol{((}\AgdaBound{a} \AgdaSymbol{:} \AgdaFunction{∣} \AgdaBound{Ax} \AgdaFunction{∣}\AgdaSymbol{)} \AgdaSymbol{→} \AgdaFunction{∣} \AgdaBound{Bx} \AgdaBound{a} \AgdaFunction{∣}\AgdaSymbol{)} \AgdaSymbol{(λ} \AgdaBound{fn} \AgdaSymbol{→} \<[62]%
\>[62]\<%
\\
\>[9]\AgdaIndent{21}{}\<[21]%
\>[21]\AgdaSymbol{(}\AgdaBound{a} \AgdaBound{b} \AgdaSymbol{:} \AgdaFunction{∣} \AgdaBound{Ax} \AgdaFunction{∣}\AgdaSymbol{)}\<%
\\
\>[9]\AgdaIndent{21}{}\<[21]%
\>[21]\AgdaSymbol{(}\AgdaBound{p} \AgdaSymbol{:} \AgdaFunction{[} \AgdaBound{Ax} \AgdaFunction{]} \AgdaBound{a} \AgdaFunction{≈} \AgdaBound{b}\AgdaSymbol{)} \AgdaSymbol{→}\<%
\\
\>[9]\AgdaIndent{21}{}\<[21]%
\>[21]\AgdaFunction{[} \AgdaBound{Bx} \AgdaBound{b} \AgdaFunction{]} \AgdaFunction{[} \AgdaBound{B} \AgdaFunction{]subst} \AgdaSymbol{(}\AgdaFunction{[} \AgdaBound{Γ} \AgdaFunction{]refl} \AgdaInductiveConstructor{,} \<[54]%
\>[54]\<%
\\
\>[9]\AgdaIndent{21}{}\<[21]%
\>[21]\AgdaFunction{[} \AgdaBound{Ax} \AgdaFunction{]trans} \AgdaFunction{[} \AgdaBound{A} \AgdaFunction{]refl*} \AgdaBound{p}\AgdaSymbol{)} \AgdaSymbol{(}\AgdaBound{fn} \AgdaBound{a}\AgdaSymbol{)} \AgdaFunction{≈} \AgdaBound{fn} \AgdaBound{b}\AgdaSymbol{)}\<%
\\
%
\\
%
\\
\>[0]\AgdaIndent{9}{}\<[9]%
\>[9]\AgdaSymbol{;} \AgdaField{\_≈\_} \<[18]%
\>[18]\AgdaSymbol{=} \AgdaSymbol{λ\{(}\AgdaBound{f} \AgdaInductiveConstructor{,} \AgdaSymbol{\_)} \AgdaSymbol{(}\AgdaBound{g} \AgdaInductiveConstructor{,} \AgdaSymbol{\_)} \AgdaSymbol{→} \AgdaSymbol{∀} \AgdaBound{a} \AgdaSymbol{→} \AgdaFunction{[} \AgdaBound{Bx} \AgdaBound{a} \AgdaFunction{]} \AgdaBound{f} \AgdaBound{a} \AgdaFunction{≈} \AgdaBound{g} \AgdaBound{a} \AgdaSymbol{\}}\<%
\\
\>[0]\AgdaIndent{9}{}\<[9]%
\>[9]\AgdaSymbol{;} \AgdaField{refl} \<[19]%
\>[19]\AgdaSymbol{=} \AgdaSymbol{λ} \AgdaBound{a} \AgdaSymbol{→} \AgdaFunction{[} \AgdaBound{Bx} \AgdaSymbol{\_} \AgdaFunction{]refl} \<[40]%
\>[40]\<%
\\
\>[0]\AgdaIndent{9}{}\<[9]%
\>[9]\AgdaSymbol{;} \AgdaField{sym} \<[19]%
\>[19]\AgdaSymbol{=} \AgdaSymbol{λ} \AgdaBound{f} \AgdaBound{a} \AgdaSymbol{→} \AgdaFunction{[} \AgdaBound{Bx} \AgdaSymbol{\_} \AgdaFunction{]sym} \AgdaSymbol{(}\AgdaBound{f} \AgdaBound{a}\AgdaSymbol{)}\<%
\\
\>[0]\AgdaIndent{9}{}\<[9]%
\>[9]\AgdaSymbol{;} \AgdaField{trans} \<[19]%
\>[19]\AgdaSymbol{=} \AgdaSymbol{λ} \AgdaBound{f} \AgdaBound{g} \AgdaBound{a} \AgdaSymbol{→} \AgdaFunction{[} \AgdaBound{Bx} \AgdaSymbol{\_} \AgdaFunction{]trans} \AgdaSymbol{(}\AgdaBound{f} \AgdaBound{a}\AgdaSymbol{)} \AgdaSymbol{(}\AgdaBound{g} \AgdaBound{a}\AgdaSymbol{)}\<%
\\
\>[0]\AgdaIndent{9}{}\<[9]%
\>[9]\AgdaSymbol{\}}\<%
\\
%
\\
\>[0]\AgdaIndent{2}{}\<[2]%
\>[2]\AgdaSymbol{;} \AgdaField{substT} \AgdaSymbol{=} \AgdaSymbol{λ} \AgdaSymbol{\{}\AgdaBound{x}\AgdaSymbol{\}} \AgdaSymbol{\{}\AgdaBound{y}\AgdaSymbol{\}} \AgdaBound{p} \AgdaSymbol{→} \AgdaSymbol{λ} \AgdaSymbol{\{(}\AgdaBound{f} \AgdaInductiveConstructor{,} \AgdaBound{rsp}\AgdaSymbol{)} \AgdaSymbol{→}\<%
\\
\>[0]\AgdaIndent{19}{}\<[19]%
\>[19]\AgdaKeyword{let} \AgdaBound{y2x} \AgdaSymbol{=} \AgdaSymbol{λ} \AgdaBound{a} \AgdaSymbol{→} \AgdaFunction{[} \AgdaBound{A} \AgdaFunction{]subst} \AgdaSymbol{(}\AgdaFunction{[} \AgdaBound{Γ} \AgdaFunction{]sym} \AgdaBound{p}\AgdaSymbol{)} \AgdaBound{a} \AgdaKeyword{in}\<%
\\
\>[0]\AgdaIndent{19}{}\<[19]%
\>[19]\AgdaKeyword{let} \AgdaBound{x2y} \AgdaSymbol{=} \AgdaSymbol{λ} \AgdaBound{a} \AgdaSymbol{→} \AgdaFunction{[} \AgdaBound{A} \AgdaFunction{]subst} \AgdaBound{p} \AgdaBound{a} \AgdaKeyword{in}\<%
\\
\>[0]\AgdaIndent{13}{}\<[13]%
\>[13]\AgdaSymbol{(λ} \AgdaBound{a} \AgdaSymbol{→} \AgdaFunction{[} \AgdaBound{B} \AgdaFunction{]subst} \AgdaSymbol{(}\AgdaBound{p} \AgdaInductiveConstructor{,} \AgdaFunction{[} \AgdaBound{A} \AgdaFunction{]tr*}\AgdaSymbol{)} \<[46]%
\>[46]\<%
\\
\>[0]\AgdaIndent{13}{}\<[13]%
\>[13]\AgdaSymbol{(}\AgdaBound{f} \AgdaSymbol{(}\AgdaBound{y2x} \AgdaBound{a}\AgdaSymbol{)))} \AgdaInductiveConstructor{,} \<[28]%
\>[28]\<%
\\
\>[0]\AgdaIndent{13}{}\<[13]%
\>[13]\AgdaSymbol{(λ} \AgdaBound{a} \AgdaBound{b} \AgdaBound{q} \AgdaSymbol{→}\<%
\\
\>[13]\AgdaIndent{16}{}\<[16]%
\>[16]\AgdaKeyword{let} \AgdaBound{a'} \AgdaSymbol{=} \AgdaBound{y2x} \AgdaBound{a} \AgdaKeyword{in} \<[34]%
\>[34]\<%
\\
\>[13]\AgdaIndent{16}{}\<[16]%
\>[16]\AgdaKeyword{let} \AgdaBound{b'} \AgdaSymbol{=} \AgdaBound{y2x} \AgdaBound{b} \AgdaKeyword{in}\<%
\\
\>[13]\AgdaIndent{16}{}\<[16]%
\>[16]\AgdaKeyword{let} \AgdaBound{q'} \AgdaSymbol{=} \AgdaFunction{[} \AgdaBound{A} \AgdaFunction{]subst*} \AgdaSymbol{(}\AgdaFunction{[} \AgdaBound{Γ} \AgdaFunction{]sym} \AgdaBound{p}\AgdaSymbol{)} \AgdaBound{q} \AgdaKeyword{in}\<%
\\
\>[13]\AgdaIndent{16}{}\<[16]%
\>[16]\AgdaKeyword{let} \AgdaBound{H} \AgdaSymbol{=} \AgdaBound{rsp} \AgdaBound{a'} \AgdaBound{b'} \AgdaSymbol{(}\AgdaFunction{[} \AgdaBound{A} \AgdaFunction{]subst*} \AgdaSymbol{(}\AgdaFunction{[} \AgdaBound{Γ} \AgdaFunction{]sym} \AgdaBound{p}\AgdaSymbol{)} \AgdaBound{q}\AgdaSymbol{)} \AgdaKeyword{in}\<%
\\
\>[13]\AgdaIndent{16}{}\<[16]%
\>[16]\AgdaKeyword{let} \AgdaBound{r} \AgdaSymbol{:} \AgdaFunction{[} \AgdaBound{Γ} \AgdaFunction{\&} \AgdaBound{A} \AgdaFunction{]} \AgdaSymbol{(}\AgdaBound{x} \AgdaInductiveConstructor{,} \AgdaBound{b'}\AgdaSymbol{)} \AgdaFunction{≈} \AgdaSymbol{(}\AgdaBound{y} \AgdaInductiveConstructor{,} \AgdaBound{b}\AgdaSymbol{)}
                    r \AgdaSymbol{=} \AgdaSymbol{(}\AgdaBound{p} \AgdaInductiveConstructor{,} \AgdaFunction{[} \AgdaBound{A} \AgdaFunction{]tr*}\AgdaSymbol{)} \AgdaKeyword{in}\<%
\\
\>[13]\AgdaIndent{16}{}\<[16]%
\>[16]\AgdaKeyword{let} \AgdaBound{pre} \AgdaSymbol{=} \AgdaFunction{[} \AgdaBound{B} \AgdaFunction{]subst*} \AgdaBound{r} \<[40]%
\>[40]\<%
\\
\>[16]\AgdaIndent{24}{}\<[24]%
\>[24]\AgdaSymbol{(}\AgdaBound{rsp} \AgdaBound{a'} \AgdaBound{b'} \AgdaSymbol{(}\AgdaFunction{[} \AgdaBound{A} \AgdaFunction{]subst*} \AgdaSymbol{(}\AgdaFunction{[} \AgdaBound{Γ} \AgdaFunction{]sym} \AgdaBound{p}\AgdaSymbol{)} \AgdaBound{q}\AgdaSymbol{))} \AgdaKeyword{in} \<[68]%
\>[68]\<%
\\
\>[0]\AgdaIndent{16}{}\<[16]%
\>[16]\AgdaFunction{[} \AgdaFunction{[} \AgdaBound{B} \AgdaFunction{]fm} \AgdaSymbol{(}\AgdaBound{y} \AgdaInductiveConstructor{,} \AgdaBound{b}\AgdaSymbol{)} \AgdaFunction{]trans} \<[41]%
\>[41]\<%
\\
\>[0]\AgdaIndent{16}{}\<[16]%
\>[16]\AgdaSymbol{(}\AgdaFunction{[} \AgdaBound{B} \AgdaFunction{]trans*} \AgdaSymbol{\_)} \<[32]%
\>[32]\<%
\\
\>[0]\AgdaIndent{16}{}\<[16]%
\>[16]\AgdaSymbol{(}\AgdaFunction{[} \AgdaFunction{[} \AgdaBound{B} \AgdaFunction{]fm} \AgdaSymbol{(}\AgdaBound{y} \AgdaInductiveConstructor{,} \AgdaBound{b}\AgdaSymbol{)} \AgdaFunction{]trans} \<[42]%
\>[42]\<%
\\
\>[0]\AgdaIndent{16}{}\<[16]%
\>[16]\AgdaSymbol{(}\AgdaFunction{[} \AgdaFunction{[} \AgdaBound{B} \AgdaFunction{]fm} \AgdaSymbol{(}\AgdaBound{y} \AgdaInductiveConstructor{,} \AgdaBound{b}\AgdaSymbol{)} \AgdaFunction{]sym} \AgdaSymbol{(}\AgdaFunction{[} \AgdaBound{B} \AgdaFunction{]trans*} \AgdaSymbol{\_))} \<[57]%
\>[57]\<%
\\
\>[0]\AgdaIndent{16}{}\<[16]%
\>[16]\AgdaBound{pre}\AgdaSymbol{))\}}\<%
\\
%
\\
%
\\
\>[0]\AgdaIndent{2}{}\<[2]%
\>[2]\AgdaSymbol{;} \AgdaField{subst*} \AgdaSymbol{=} \AgdaSymbol{λ} \AgdaBound{\_} \AgdaBound{q} \AgdaBound{\_} \AgdaSymbol{→} \AgdaFunction{[} \AgdaBound{B} \AgdaFunction{]subst*} \AgdaSymbol{\_} \AgdaSymbol{(}\AgdaBound{q} \AgdaSymbol{\_)}\<%
\\
\>[0]\AgdaIndent{2}{}\<[2]%
\>[2]\AgdaSymbol{;} \AgdaField{refl*} \AgdaSymbol{=} \AgdaSymbol{λ} \AgdaSymbol{\{}\AgdaBound{x}\AgdaSymbol{\}} \AgdaSymbol{\{}\AgdaBound{a}\AgdaSymbol{\}} \AgdaBound{ax} \<[25]%
\>[25]\<%
\\
\>[2]\AgdaIndent{10}{}\<[10]%
\>[10]\AgdaSymbol{→} \AgdaKeyword{let} \AgdaBound{rsp} \AgdaSymbol{=} \AgdaFunction{prj₂} \AgdaBound{a} \AgdaKeyword{in} \AgdaSymbol{(}\AgdaBound{rsp} \AgdaSymbol{\_} \AgdaSymbol{\_} \AgdaFunction{[} \AgdaBound{A} \AgdaFunction{]refl*}\AgdaSymbol{)}\<%
\\
\>[0]\AgdaIndent{2}{}\<[2]%
\>[2]\AgdaSymbol{;} \AgdaField{trans*} \AgdaSymbol{=} \<[14]%
\>[14]\AgdaSymbol{λ} \AgdaSymbol{\{(}\AgdaBound{f} \AgdaInductiveConstructor{,} \AgdaBound{rsp}\AgdaSymbol{)} \AgdaBound{a} \AgdaSymbol{→}\<%
\\
\>[0]\AgdaIndent{13}{}\<[13]%
\>[13]\AgdaFunction{[} \AgdaFunction{[} \AgdaBound{B} \AgdaFunction{]fm} \AgdaSymbol{\_} \AgdaFunction{]trans} \<[32]%
\>[32]\<%
\\
\>[0]\AgdaIndent{13}{}\<[13]%
\>[13]\AgdaSymbol{(}\AgdaFunction{[} \AgdaFunction{[} \AgdaBound{B} \AgdaFunction{]fm} \AgdaSymbol{\_} \AgdaFunction{]trans} \<[33]%
\>[33]\<%
\\
\>[0]\AgdaIndent{13}{}\<[13]%
\>[13]\AgdaSymbol{(}\AgdaFunction{[} \AgdaBound{B} \AgdaFunction{]trans*} \AgdaSymbol{\_)} \<[29]%
\>[29]\<%
\\
\>[0]\AgdaIndent{13}{}\<[13]%
\>[13]\AgdaSymbol{(}\AgdaFunction{[} \AgdaFunction{[} \AgdaBound{B} \AgdaFunction{]fm} \AgdaSymbol{\_} \AgdaFunction{]sym} \AgdaSymbol{(}\AgdaFunction{[} \AgdaBound{B} \AgdaFunction{]trans*} \AgdaSymbol{\_)))}\<%
\\
\>[0]\AgdaIndent{13}{}\<[13]%
\>[13]\AgdaSymbol{(}\AgdaFunction{[} \AgdaBound{B} \AgdaFunction{]subst*} \AgdaSymbol{\_} \AgdaSymbol{(}\AgdaBound{rsp} \AgdaSymbol{\_} \AgdaSymbol{\_} \AgdaSymbol{(}\AgdaFunction{[} \AgdaBound{A} \AgdaFunction{]trans*} \AgdaSymbol{\_)} \AgdaSymbol{))} \AgdaSymbol{\}}\<%
\\
\>[0]\AgdaIndent{2}{}\<[2]%
\>[2]\AgdaSymbol{\}}\<%
\\
%
\\
%
\\
\>\AgdaFunction{lam} \AgdaSymbol{:} \AgdaSymbol{\{}\AgdaBound{Γ} \AgdaSymbol{:} \AgdaFunction{Con}\AgdaSymbol{\}\{}\AgdaBound{A} \AgdaSymbol{:} \AgdaRecord{Ty} \AgdaBound{Γ}\AgdaSymbol{\}\{}\AgdaBound{B} \AgdaSymbol{:} \AgdaRecord{Ty} \AgdaSymbol{(}\AgdaBound{Γ} \AgdaFunction{\&} \AgdaBound{A}\AgdaSymbol{)\}} \AgdaSymbol{→} \AgdaRecord{Tm} \AgdaBound{B} \AgdaSymbol{→} \AgdaRecord{Tm} \AgdaSymbol{(}\AgdaFunction{Π} \AgdaBound{A} \AgdaBound{B}\AgdaSymbol{)}\<%
\\
\>\AgdaFunction{lam} \AgdaSymbol{\{}\AgdaBound{Γ}\AgdaSymbol{\}} \AgdaSymbol{\{}\AgdaBound{A}\AgdaSymbol{\}} \AgdaSymbol{(}\AgdaInductiveConstructor{tm:} \AgdaBound{tm} \AgdaInductiveConstructor{resp:} \AgdaBound{respt}\AgdaSymbol{)} \AgdaSymbol{=} \<[35]%
\>[35]\<%
\\
\>[0]\AgdaIndent{2}{}\<[2]%
\>[2]\AgdaKeyword{record} \AgdaSymbol{\{} \AgdaField{tm} \AgdaSymbol{=} \AgdaSymbol{λ} \AgdaBound{x} \AgdaSymbol{→} \AgdaSymbol{(λ} \AgdaBound{a} \AgdaSymbol{→} \AgdaBound{tm} \AgdaSymbol{(}\AgdaBound{x} \AgdaInductiveConstructor{,} \AgdaBound{a}\AgdaSymbol{))} \<[41]%
\>[41]\<%
\\
\>[2]\AgdaIndent{16}{}\<[16]%
\>[16]\AgdaInductiveConstructor{,} \AgdaSymbol{(λ} \AgdaBound{a} \AgdaBound{b} \AgdaBound{p} \AgdaSymbol{→} \AgdaBound{respt} \AgdaSymbol{(}\AgdaFunction{[} \AgdaBound{Γ} \AgdaFunction{]refl} \AgdaInductiveConstructor{,} \<[48]%
\>[48]\<%
\\
\>[2]\AgdaIndent{16}{}\<[16]%
\>[16]\AgdaFunction{[} \AgdaFunction{[} \AgdaBound{A} \AgdaFunction{]fm} \AgdaBound{x} \AgdaFunction{]trans} \AgdaFunction{[} \AgdaBound{A} \AgdaFunction{]refl*} \AgdaBound{p}\AgdaSymbol{))}\<%
\\
\>[5]\AgdaIndent{9}{}\<[9]%
\>[9]\AgdaSymbol{;} \AgdaField{respt} \AgdaSymbol{=} \AgdaSymbol{λ} \AgdaBound{p} \AgdaBound{\_} \AgdaSymbol{→} \AgdaBound{respt} \AgdaSymbol{(}\AgdaBound{p} \AgdaInductiveConstructor{,} \AgdaFunction{[} \AgdaBound{A} \AgdaFunction{]tr*}\AgdaSymbol{)}\<%
\\
\>[0]\AgdaIndent{9}{}\<[9]%
\>[9]\AgdaSymbol{\}}\<%
\\
%
\\
%
\\
\>\AgdaFunction{app} \AgdaSymbol{:} \AgdaSymbol{\{}\AgdaBound{Γ} \AgdaSymbol{:} \AgdaFunction{Con}\AgdaSymbol{\}\{}\AgdaBound{A} \AgdaSymbol{:} \AgdaRecord{Ty} \AgdaBound{Γ}\AgdaSymbol{\}\{}\AgdaBound{B} \AgdaSymbol{:} \AgdaRecord{Ty} \AgdaSymbol{(}\AgdaBound{Γ} \AgdaFunction{\&} \AgdaBound{A}\AgdaSymbol{)\}} \AgdaSymbol{→} \AgdaRecord{Tm} \AgdaSymbol{(}\AgdaFunction{Π} \AgdaBound{A} \AgdaBound{B}\AgdaSymbol{)} \AgdaSymbol{→} \AgdaRecord{Tm} \AgdaBound{B}\<%
\\
\>\AgdaFunction{app} \AgdaSymbol{\{}\AgdaBound{Γ}\AgdaSymbol{\}} \AgdaSymbol{\{}\AgdaBound{A}\AgdaSymbol{\}} \AgdaSymbol{\{}\AgdaBound{B}\AgdaSymbol{\}} \AgdaSymbol{(}\AgdaInductiveConstructor{tm:} \AgdaBound{tm} \AgdaInductiveConstructor{resp:} \AgdaBound{respt}\AgdaSymbol{)} \AgdaSymbol{=} \<[39]%
\>[39]\<%
\\
\>[0]\AgdaIndent{2}{}\<[2]%
\>[2]\AgdaKeyword{record} \AgdaSymbol{\{} \AgdaField{tm} \AgdaSymbol{=} \AgdaSymbol{λ} \AgdaSymbol{\{(}\AgdaBound{x} \AgdaInductiveConstructor{,} \AgdaBound{a}\AgdaSymbol{)} \AgdaSymbol{→} \AgdaFunction{prj₁} \AgdaSymbol{(}\AgdaBound{tm} \AgdaBound{x}\AgdaSymbol{)} \AgdaBound{a}\AgdaSymbol{\}}\<%
\\
\>[0]\AgdaIndent{9}{}\<[9]%
\>[9]\AgdaSymbol{;} \AgdaField{respt} \AgdaSymbol{=} \AgdaSymbol{λ} \AgdaSymbol{\{}\AgdaBound{x}\AgdaSymbol{\}} \AgdaSymbol{\{}\AgdaBound{y}\AgdaSymbol{\}} \AgdaSymbol{→} \AgdaSymbol{λ} \AgdaSymbol{\{(}\AgdaBound{p} \AgdaInductiveConstructor{,} \AgdaBound{tr}\AgdaSymbol{)} \AgdaSymbol{→} \<[45]%
\>[45]\<%
\\
\>[9]\AgdaIndent{17}{}\<[17]%
\>[17]\AgdaKeyword{let} \AgdaBound{fresp} \AgdaSymbol{=} \AgdaFunction{prj₂} \AgdaSymbol{(}\AgdaBound{tm} \AgdaSymbol{(}\AgdaFunction{proj₁} \AgdaBound{x}\AgdaSymbol{))} \AgdaKeyword{in}\<%
\\
\>[9]\AgdaIndent{17}{}\<[17]%
\>[17]\AgdaFunction{[} \AgdaFunction{[} \AgdaBound{B} \AgdaFunction{]fm} \AgdaSymbol{\_} \AgdaFunction{]trans} \<[36]%
\>[36]\<%
\\
\>[9]\AgdaIndent{17}{}\<[17]%
\>[17]\AgdaSymbol{(}\AgdaFunction{[} \AgdaBound{B} \AgdaFunction{]subst*} \AgdaSymbol{(}\AgdaBound{p} \AgdaInductiveConstructor{,} \AgdaBound{tr}\AgdaSymbol{)} \<[39]%
\>[39]\<%
\\
\>[9]\AgdaIndent{17}{}\<[17]%
\>[17]\AgdaSymbol{(}\AgdaFunction{[} \AgdaFunction{[} \AgdaBound{B} \AgdaFunction{]fm} \AgdaSymbol{\_} \AgdaFunction{]sym} \AgdaFunction{[} \AgdaBound{B} \AgdaFunction{]refl*}\AgdaSymbol{))}\<%
\\
\>[9]\AgdaIndent{17}{}\<[17]%
\>[17]\AgdaSymbol{(}\AgdaFunction{[} \AgdaFunction{[} \AgdaBound{B} \AgdaFunction{]fm} \AgdaSymbol{\_} \AgdaFunction{]trans}\<%
\\
\>[9]\AgdaIndent{17}{}\<[17]%
\>[17]\AgdaSymbol{(}\AgdaFunction{[} \AgdaBound{B} \AgdaFunction{]trans*} \AgdaSymbol{\{}p \AgdaSymbol{=} \AgdaSymbol{(}\AgdaFunction{[} \AgdaBound{Γ} \AgdaFunction{]refl} \AgdaInductiveConstructor{,} \AgdaFunction{[} \AgdaBound{A} \AgdaFunction{]refl*}\AgdaSymbol{)\}} \AgdaSymbol{\_)}\<%
\\
\>[0]\AgdaIndent{8}{}\<[8]%
\>[8]\<%
\\
\>[0]\AgdaIndent{17}{}\<[17]%
\>[17]\AgdaSymbol{(}\AgdaFunction{[} \AgdaFunction{[} \AgdaBound{B} \AgdaFunction{]fm} \AgdaSymbol{\_} \AgdaFunction{]trans} \<[37]%
\>[37]\<%
\\
\>[0]\AgdaIndent{17}{}\<[17]%
\>[17]\AgdaSymbol{(}\AgdaFunction{[} \AgdaFunction{[} \AgdaBound{B} \AgdaFunction{]fm} \AgdaSymbol{\_} \AgdaFunction{]sym} \<[35]%
\>[35]\<%
\\
\>[0]\AgdaIndent{17}{}\<[17]%
\>[17]\AgdaSymbol{(}\AgdaFunction{[} \AgdaBound{B} \AgdaFunction{]trans*} \AgdaSymbol{\{}q \AgdaSymbol{=} \AgdaSymbol{(}\AgdaBound{p} \AgdaInductiveConstructor{,} \AgdaFunction{[} \AgdaBound{A} \AgdaFunction{]tr*}\AgdaSymbol{)\}} \AgdaSymbol{\_))}\<%
\\
\>[0]\AgdaIndent{17}{}\<[17]%
\>[17]\AgdaSymbol{(}\AgdaFunction{[} \AgdaFunction{[} \AgdaBound{B} \AgdaFunction{]fm} \AgdaSymbol{\_} \AgdaFunction{]trans} \<[37]%
\>[37]\<%
\\
\>[0]\AgdaIndent{17}{}\<[17]%
\>[17]\AgdaSymbol{(}\AgdaFunction{[} \AgdaBound{B} \AgdaFunction{]subst*} \AgdaSymbol{\_} \AgdaSymbol{(}\AgdaBound{fresp} \AgdaSymbol{\_} \AgdaSymbol{\_} \<[43]%
\>[43]\<%
\\
\>[17]\AgdaIndent{19}{}\<[19]%
\>[19]\AgdaSymbol{(}\AgdaFunction{[} \AgdaFunction{[} \AgdaBound{A} \AgdaFunction{]fm} \AgdaSymbol{\_} \AgdaFunction{]trans} \AgdaSymbol{(}\AgdaFunction{[} \AgdaFunction{[} \AgdaBound{A} \AgdaFunction{]fm} \AgdaSymbol{\_} \AgdaFunction{]sym} \AgdaFunction{[} \AgdaBound{A} \AgdaFunction{]tr*}\AgdaSymbol{)} \<[67]%
\>[67]\<%
\\
\>[17]\AgdaIndent{19}{}\<[19]%
\>[19]\AgdaSymbol{(}\AgdaFunction{[} \AgdaBound{A} \AgdaFunction{]subst*} \AgdaSymbol{(}\AgdaFunction{[} \AgdaBound{Γ} \AgdaFunction{]sym} \AgdaBound{p}\AgdaSymbol{)} \AgdaBound{tr}\AgdaSymbol{))))}\<%
\\
\>[0]\AgdaIndent{17}{}\<[17]%
\>[17]\AgdaSymbol{(}\AgdaBound{respt} \AgdaBound{p} \AgdaSymbol{\_))))} \AgdaSymbol{\}}\<%
\\
\>[0]\AgdaIndent{9}{}\<[9]%
\>[9]\AgdaSymbol{\}}\<%
\\
%
\\
\>\AgdaFunction{\_⇒\_} \AgdaSymbol{:} \AgdaSymbol{\{}\AgdaBound{Γ} \AgdaSymbol{:} \AgdaFunction{Con}\AgdaSymbol{\}(}\AgdaBound{A} \AgdaBound{B} \AgdaSymbol{:} \AgdaRecord{Ty} \AgdaBound{Γ}\AgdaSymbol{)} \AgdaSymbol{→} \AgdaRecord{Ty} \AgdaBound{Γ}\<%
\\
\>\AgdaBound{A} \AgdaFunction{⇒} \AgdaBound{B} \AgdaSymbol{=} \AgdaFunction{Π} \AgdaBound{A} \AgdaSymbol{(}\AgdaBound{B} \AgdaFunction{[} \AgdaFunction{fst\&} \AgdaSymbol{\{}A \AgdaSymbol{=} \AgdaBound{A}\AgdaSymbol{\}} \AgdaFunction{]T}\AgdaSymbol{)}\<%
\\
%
\\
\>\AgdaKeyword{infixr} \AgdaNumber{6} \_⇒\_\<%
\\
\>\<\end{code}

Simpler definition for functions

\begin{code}\>\<%
\\
\>\AgdaFunction{[\_,\_]\_⇒fm\_} \AgdaSymbol{:} \AgdaSymbol{(}\AgdaBound{Γ} \AgdaSymbol{:} \AgdaFunction{Con}\AgdaSymbol{)(}\AgdaBound{x} \AgdaSymbol{:} \AgdaFunction{∣} \AgdaBound{Γ} \AgdaFunction{∣}\AgdaSymbol{)} \<[34]%
\>[34]\<%
\\
\>[9]\AgdaIndent{11}{}\<[11]%
\>[11]\AgdaSymbol{→} \AgdaRecord{Setoid} \AgdaSymbol{→} \AgdaRecord{Setoid} \AgdaSymbol{→} \AgdaRecord{Setoid}\<%
\\
\>\AgdaFunction{[} \AgdaBound{Γ} \AgdaFunction{,} \AgdaBound{x} \AgdaFunction{]} \AgdaBound{Ax} \AgdaFunction{⇒fm} \AgdaBound{Bx} \<[20]%
\>[20]\<%
\\
\>[0]\AgdaIndent{2}{}\<[2]%
\>[2]\AgdaSymbol{=} \AgdaKeyword{record}\<%
\\
\>[0]\AgdaIndent{6}{}\<[6]%
\>[6]\AgdaSymbol{\{} \AgdaField{Carrier} \AgdaSymbol{=} \AgdaRecord{Σ[} \AgdaBound{fn} \AgdaRecord{∶} \AgdaSymbol{(}\AgdaFunction{∣} \AgdaBound{Ax} \AgdaFunction{∣} \AgdaSymbol{→} \AgdaFunction{∣} \AgdaBound{Bx} \AgdaFunction{∣}\AgdaSymbol{)} \AgdaRecord{]} \AgdaSymbol{((}\AgdaBound{a} \AgdaBound{b} \AgdaSymbol{:} \AgdaFunction{∣} \AgdaBound{Ax} \AgdaFunction{∣}\AgdaSymbol{)}\<%
\\
\>[6]\AgdaIndent{18}{}\<[18]%
\>[18]\AgdaSymbol{(}\AgdaBound{p} \AgdaSymbol{:} \AgdaFunction{[} \AgdaBound{Ax} \AgdaFunction{]} \AgdaBound{a} \AgdaFunction{≈} \AgdaBound{b}\AgdaSymbol{)} \AgdaSymbol{→} \AgdaFunction{[} \AgdaBound{Bx} \AgdaFunction{]} \AgdaBound{fn} \AgdaBound{a} \AgdaFunction{≈} \AgdaBound{fn} \AgdaBound{b}\AgdaSymbol{)}\<%
\\
\>[0]\AgdaIndent{6}{}\<[6]%
\>[6]\AgdaSymbol{;} \AgdaField{\_≈\_} \<[15]%
\>[15]\AgdaSymbol{=} \AgdaSymbol{λ\{(}\AgdaBound{f} \AgdaInductiveConstructor{,} \AgdaSymbol{\_)} \AgdaSymbol{(}\AgdaBound{g} \AgdaInductiveConstructor{,} \AgdaSymbol{\_)} \AgdaSymbol{→} \AgdaSymbol{∀} \AgdaBound{a} \AgdaSymbol{→} \AgdaFunction{[} \AgdaBound{Bx} \AgdaFunction{]} \AgdaBound{f} \AgdaBound{a} \AgdaFunction{≈} \AgdaBound{g} \AgdaBound{a} \AgdaSymbol{\}}\<%
\\
\>[0]\AgdaIndent{6}{}\<[6]%
\>[6]\AgdaSymbol{;} \AgdaField{refl} \<[16]%
\>[16]\AgdaSymbol{=} \AgdaSymbol{λ} \AgdaBound{\_} \AgdaSymbol{→} \AgdaFunction{[} \AgdaBound{Bx} \AgdaFunction{]refl} \<[35]%
\>[35]\<%
\\
\>[0]\AgdaIndent{6}{}\<[6]%
\>[6]\AgdaSymbol{;} \AgdaField{sym} \<[16]%
\>[16]\AgdaSymbol{=} \AgdaSymbol{λ} \AgdaBound{f} \AgdaBound{a} \AgdaSymbol{→} \AgdaFunction{[} \AgdaBound{Bx} \AgdaFunction{]sym} \AgdaSymbol{(}\AgdaBound{f} \AgdaBound{a}\AgdaSymbol{)}\<%
\\
\>[0]\AgdaIndent{6}{}\<[6]%
\>[6]\AgdaSymbol{;} \AgdaField{trans} \<[16]%
\>[16]\AgdaSymbol{=} \AgdaSymbol{λ} \AgdaBound{f} \AgdaBound{g} \AgdaBound{a} \AgdaSymbol{→} \AgdaFunction{[} \AgdaBound{Bx} \AgdaFunction{]trans} \AgdaSymbol{(}\AgdaBound{f} \AgdaBound{a}\AgdaSymbol{)} \AgdaSymbol{(}\AgdaBound{g} \AgdaBound{a}\AgdaSymbol{)}\<%
\\
\>[0]\AgdaIndent{6}{}\<[6]%
\>[6]\AgdaSymbol{\}}\<%
\\
\>\<\end{code}

$\Sigma$-types (object level)

\begin{code}\>\<%
\\
%
\\
\>\AgdaFunction{Σ'} \AgdaSymbol{:} \AgdaSymbol{\{}\AgdaBound{Γ} \AgdaSymbol{:} \AgdaFunction{Con}\AgdaSymbol{\}(}\AgdaBound{A} \AgdaSymbol{:} \AgdaRecord{Ty} \AgdaBound{Γ}\AgdaSymbol{)(}\AgdaBound{B} \AgdaSymbol{:} \AgdaRecord{Ty} \AgdaSymbol{(}\AgdaBound{Γ} \AgdaFunction{\&} \AgdaBound{A}\AgdaSymbol{))} \AgdaSymbol{→} \AgdaRecord{Ty} \AgdaBound{Γ}\<%
\\
\>\AgdaFunction{Σ'} \AgdaSymbol{\{}\AgdaBound{Γ}\AgdaSymbol{\}} \AgdaBound{A} \AgdaBound{B} \AgdaSymbol{=} \AgdaKeyword{record} \<[20]%
\>[20]\<%
\\
\>[6]\AgdaIndent{8}{}\<[8]%
\>[8]\AgdaSymbol{\{} \AgdaField{fm} \AgdaSymbol{=} \AgdaSymbol{λ} \AgdaBound{x} \AgdaSymbol{→} \AgdaKeyword{let} \AgdaBound{Ax} \AgdaSymbol{=} \AgdaFunction{[} \AgdaBound{A} \AgdaFunction{]fm} \AgdaBound{x} \AgdaKeyword{in}\<%
\\
\>[8]\AgdaIndent{15}{}\<[15]%
\>[15]\AgdaKeyword{let} \AgdaBound{Bx} \AgdaSymbol{=} \AgdaSymbol{λ} \AgdaBound{a} \AgdaSymbol{→} \AgdaFunction{[} \AgdaBound{B} \AgdaFunction{]fm} \AgdaSymbol{(}\AgdaBound{x} \AgdaInductiveConstructor{,} \AgdaBound{a}\AgdaSymbol{)} \AgdaKeyword{in}\<%
\\
\>[-2]\AgdaIndent{9}{}\<[9]%
\>[9]\AgdaKeyword{record}\<%
\\
\>[0]\AgdaIndent{11}{}\<[11]%
\>[11]\AgdaSymbol{\{} \AgdaField{Carrier} \AgdaSymbol{=} \AgdaRecord{Σ[} \AgdaBound{a} \AgdaRecord{∶} \AgdaFunction{∣} \AgdaBound{Ax} \AgdaFunction{∣} \AgdaRecord{]} \AgdaFunction{∣} \AgdaBound{Bx} \AgdaBound{a} \AgdaFunction{∣}\<%
\\
%
\\
\>[0]\AgdaIndent{11}{}\<[11]%
\>[11]\AgdaSymbol{;} \AgdaField{\_≈\_} \<[20]%
\>[20]\AgdaSymbol{=} \AgdaSymbol{λ\{(}\AgdaBound{a₁} \AgdaInductiveConstructor{,} \AgdaBound{b₁}\AgdaSymbol{)} \AgdaSymbol{(}\AgdaBound{a₂} \AgdaInductiveConstructor{,} \AgdaBound{b₂}\AgdaSymbol{)} \AgdaSymbol{→} \<[46]%
\>[46]\<%
\\
\>[11]\AgdaIndent{22}{}\<[22]%
\>[22]\AgdaRecord{Subset} \AgdaSymbol{(}\AgdaFunction{[} \AgdaBound{Ax} \AgdaFunction{]} \AgdaBound{a₁} \AgdaFunction{≈} \AgdaBound{a₂}\AgdaSymbol{)} \<[46]%
\>[46]\<%
\\
\>[11]\AgdaIndent{22}{}\<[22]%
\>[22]\AgdaSymbol{(λ} \AgdaBound{eq₁} \AgdaSymbol{→} \AgdaFunction{[} \AgdaBound{Bx} \AgdaSymbol{\_} \AgdaFunction{]} \AgdaFunction{[} \AgdaBound{B} \AgdaFunction{]subst} \<[51]%
\>[51]\<%
\\
\>[22]\AgdaIndent{23}{}\<[23]%
\>[23]\AgdaSymbol{(}\AgdaFunction{[} \AgdaBound{Γ} \AgdaFunction{]refl} \AgdaInductiveConstructor{,} \AgdaFunction{[} \AgdaFunction{[} \AgdaBound{A} \AgdaFunction{]fm} \AgdaBound{x} \AgdaFunction{]trans} \<[55]%
\>[55]\<%
\\
\>[22]\AgdaIndent{23}{}\<[23]%
\>[23]\AgdaFunction{[} \AgdaBound{A} \AgdaFunction{]refl*} \AgdaBound{eq₁}\AgdaSymbol{)} \AgdaBound{b₁} \AgdaFunction{≈} \AgdaBound{b₂}\AgdaSymbol{)}\<%
\\
\>[-10]\AgdaIndent{11}{}\<[11]%
\>[11]\AgdaSymbol{\}}\<%
\\
%
\\
\>[0]\AgdaIndent{11}{}\<[11]%
\>[11]\AgdaSymbol{;} \AgdaField{refl} \<[21]%
\>[21]\AgdaSymbol{=} \AgdaSymbol{λ} \AgdaSymbol{\{}\AgdaBound{t}\AgdaSymbol{\}} \AgdaSymbol{→} \AgdaFunction{[} \AgdaBound{Ax} \AgdaFunction{]refl} \AgdaInductiveConstructor{,} \AgdaFunction{[} \AgdaBound{B} \AgdaFunction{]refl*}\<%
\\
%
\\
\>[0]\AgdaIndent{11}{}\<[11]%
\>[11]\AgdaSymbol{;} \AgdaField{sym} \<[21]%
\>[21]\AgdaSymbol{=} \AgdaSymbol{λ} \AgdaSymbol{\{(}\AgdaBound{p} \AgdaInductiveConstructor{,} \AgdaBound{q}\AgdaSymbol{)} \AgdaSymbol{→} \AgdaSymbol{(}\AgdaFunction{[} \AgdaBound{Ax} \AgdaFunction{]sym} \AgdaBound{p}\AgdaSymbol{)} \AgdaInductiveConstructor{,} \<[52]%
\>[52]\<%
\\
\>[11]\AgdaIndent{23}{}\<[23]%
\>[23]\AgdaFunction{[} \AgdaBound{Bx} \AgdaSymbol{\_} \AgdaFunction{]trans} \AgdaSymbol{(}\AgdaFunction{[} \AgdaBound{B} \AgdaFunction{]subst*} \AgdaSymbol{\_} \<[52]%
\>[52]\<%
\\
\>[11]\AgdaIndent{23}{}\<[23]%
\>[23]\AgdaSymbol{(}\AgdaFunction{[} \AgdaBound{Bx} \AgdaSymbol{\_} \AgdaFunction{]sym} \AgdaBound{q}\AgdaSymbol{))} \AgdaFunction{[} \AgdaBound{B} \AgdaFunction{]tr*}\AgdaSymbol{\}}\<%
\\
%
\\
\>[0]\AgdaIndent{11}{}\<[11]%
\>[11]\AgdaSymbol{;} \AgdaField{trans} \<[21]%
\>[21]\AgdaSymbol{=} \AgdaSymbol{λ} \AgdaSymbol{\{(}\AgdaBound{p} \AgdaInductiveConstructor{,} \AgdaBound{q}\AgdaSymbol{)} \AgdaSymbol{(}\AgdaBound{r} \AgdaInductiveConstructor{,} \AgdaBound{s}\AgdaSymbol{)} \AgdaSymbol{→} \AgdaSymbol{(}\AgdaFunction{[} \AgdaBound{Ax} \AgdaFunction{]trans} \AgdaBound{p} \AgdaBound{r}\AgdaSymbol{)} \AgdaInductiveConstructor{,}\<%
\\
\>[0]\AgdaIndent{23}{}\<[23]%
\>[23]\AgdaFunction{[} \AgdaBound{Bx} \AgdaSymbol{\_} \AgdaFunction{]trans} \AgdaSymbol{(}\AgdaFunction{[} \AgdaBound{Bx} \AgdaSymbol{\_} \AgdaFunction{]trans} \<[52]%
\>[52]\<%
\\
\>[0]\AgdaIndent{23}{}\<[23]%
\>[23]\AgdaSymbol{(}\AgdaFunction{[} \AgdaBound{Bx} \AgdaSymbol{\_} \AgdaFunction{]sym} \AgdaSymbol{(}\AgdaFunction{[} \AgdaBound{B} \AgdaFunction{]trans*} \AgdaSymbol{\_))} \<[53]%
\>[53]\<%
\\
\>[0]\AgdaIndent{23}{}\<[23]%
\>[23]\AgdaSymbol{(}\AgdaFunction{[} \AgdaBound{B} \AgdaFunction{]subst*} \AgdaSymbol{\_} \AgdaBound{q}\AgdaSymbol{))} \AgdaBound{s}\AgdaSymbol{\}}\<%
\\
\>[0]\AgdaIndent{11}{}\<[11]%
\>[11]\AgdaSymbol{\}}\<%
\\
%
\\
\>[0]\AgdaIndent{8}{}\<[8]%
\>[8]\AgdaSymbol{;} \AgdaField{substT} \AgdaSymbol{=} \AgdaSymbol{λ} \AgdaBound{x≈y} \AgdaSymbol{→} \AgdaSymbol{λ} \AgdaSymbol{\{(}\AgdaBound{p} \AgdaInductiveConstructor{,} \AgdaBound{q}\AgdaSymbol{)} \AgdaSymbol{→} \<[40]%
\>[40]\<%
\\
\>[8]\AgdaIndent{19}{}\<[19]%
\>[19]\AgdaSymbol{(}\AgdaFunction{[} \AgdaBound{A} \AgdaFunction{]subst} \AgdaBound{x≈y} \AgdaBound{p}\AgdaSymbol{)} \AgdaInductiveConstructor{,} \AgdaFunction{[} \AgdaBound{B} \AgdaFunction{]subst} \AgdaSymbol{(}\AgdaBound{x≈y} \AgdaInductiveConstructor{,} \<[58]%
\>[58]\<%
\\
\>[8]\AgdaIndent{19}{}\<[19]%
\>[19]\AgdaFunction{[} \AgdaFunction{[} \AgdaBound{A} \AgdaFunction{]fm} \AgdaSymbol{\_} \AgdaFunction{]refl}\AgdaSymbol{)} \AgdaBound{q}\AgdaSymbol{\}}\<%
\\
%
\\
\>[0]\AgdaIndent{8}{}\<[8]%
\>[8]\AgdaSymbol{;} \AgdaField{subst*} \AgdaSymbol{=} \AgdaSymbol{λ} \AgdaBound{x≈y} \AgdaSymbol{→} \AgdaSymbol{λ} \AgdaSymbol{\{(}\AgdaBound{p} \AgdaInductiveConstructor{,} \AgdaBound{q}\AgdaSymbol{)} \AgdaSymbol{→} \AgdaFunction{[} \AgdaBound{A} \AgdaFunction{]subst*} \AgdaBound{x≈y} \AgdaBound{p} \AgdaInductiveConstructor{,} \<[60]%
\>[60]\<%
\\
\>[0]\AgdaIndent{19}{}\<[19]%
\>[19]\AgdaFunction{[} \AgdaFunction{[} \AgdaBound{B} \AgdaFunction{]fm} \AgdaSymbol{\_} \AgdaFunction{]trans} \AgdaSymbol{(}\AgdaFunction{[} \AgdaFunction{[} \AgdaBound{B} \AgdaFunction{]fm} \AgdaSymbol{\_} \AgdaFunction{]trans} \<[58]%
\>[58]\<%
\\
\>[0]\AgdaIndent{19}{}\<[19]%
\>[19]\AgdaSymbol{(}\AgdaFunction{[} \AgdaBound{B} \AgdaFunction{]trans*} \AgdaSymbol{\_)}\<%
\\
\>[0]\AgdaIndent{19}{}\<[19]%
\>[19]\AgdaSymbol{(}\AgdaFunction{[} \AgdaFunction{[} \AgdaBound{B} \AgdaFunction{]fm} \AgdaSymbol{\_} \AgdaFunction{]sym} \AgdaSymbol{(}\AgdaFunction{[} \AgdaBound{B} \AgdaFunction{]trans*} \AgdaSymbol{\_)))} \<[55]%
\>[55]\<%
\\
\>[0]\AgdaIndent{19}{}\<[19]%
\>[19]\AgdaSymbol{(}\AgdaFunction{[} \AgdaBound{B} \AgdaFunction{]subst*} \AgdaSymbol{(}\AgdaBound{x≈y} \AgdaInductiveConstructor{,} \AgdaFunction{[} \AgdaFunction{[} \AgdaBound{A} \AgdaFunction{]fm} \AgdaSymbol{\_} \AgdaFunction{]refl}\AgdaSymbol{)} \AgdaBound{q}\AgdaSymbol{)} \AgdaSymbol{\}}\<%
\\
\>[0]\AgdaIndent{8}{}\<[8]%
\>[8]\AgdaSymbol{;} \AgdaField{refl*} \AgdaSymbol{=} \AgdaSymbol{λ} \AgdaSymbol{\{}\AgdaBound{x}\AgdaSymbol{\}} \AgdaSymbol{\{}\AgdaBound{a}\AgdaSymbol{\}} \AgdaSymbol{→} \<[30]%
\>[30]\<%
\\
\>[0]\AgdaIndent{18}{}\<[18]%
\>[18]\AgdaKeyword{let} \AgdaSymbol{(}\AgdaBound{p} \AgdaInductiveConstructor{,} \AgdaBound{q}\AgdaSymbol{)} \AgdaSymbol{=} \AgdaBound{a} \AgdaKeyword{in} \AgdaFunction{[} \AgdaBound{A} \AgdaFunction{]refl*} \AgdaInductiveConstructor{,} \AgdaFunction{[} \AgdaBound{B} \AgdaFunction{]tr*}\<%
\\
\>[0]\AgdaIndent{8}{}\<[8]%
\>[8]\AgdaSymbol{;} \AgdaField{trans*} \AgdaSymbol{=} \<[20]%
\>[20]\AgdaSymbol{λ} \AgdaSymbol{\{(}\AgdaBound{p} \AgdaInductiveConstructor{,} \AgdaBound{q}\AgdaSymbol{)} \<[32]%
\>[32]\AgdaSymbol{→} \AgdaSymbol{(}\AgdaFunction{[} \AgdaBound{A} \AgdaFunction{]trans*} \AgdaSymbol{\_)} \AgdaInductiveConstructor{,} \<[52]%
\>[52]\<%
\\
\>[0]\AgdaIndent{20}{}\<[20]%
\>[20]\AgdaSymbol{(}\AgdaFunction{[} \AgdaFunction{[} \AgdaBound{B} \AgdaFunction{]fm} \AgdaSymbol{\_} \AgdaFunction{]trans}\<%
\\
\>[0]\AgdaIndent{20}{}\<[20]%
\>[20]\AgdaSymbol{(}\AgdaFunction{[} \AgdaBound{B} \AgdaFunction{]trans*} \AgdaSymbol{\_)} \AgdaSymbol{(}\AgdaFunction{[} \AgdaBound{B} \AgdaFunction{]trans*} \AgdaSymbol{\_))} \AgdaSymbol{\}}\<%
\\
\>[0]\AgdaIndent{8}{}\<[8]%
\>[8]\AgdaSymbol{\}}\<%
\\
%
\\
\>\<\end{code}

Binary relation

\begin{code}\>\<%
\\
\>\AgdaFunction{Rel} \AgdaSymbol{:} \AgdaSymbol{\{}\AgdaBound{Γ} \AgdaSymbol{:} \AgdaFunction{Con}\AgdaSymbol{\}} \AgdaSymbol{→} \AgdaRecord{Ty} \AgdaBound{Γ} \AgdaSymbol{→} \AgdaPrimitiveType{Set₁}\<%
\\
\>\AgdaFunction{Rel} \AgdaSymbol{\{}\AgdaBound{Γ}\AgdaSymbol{\}} \AgdaBound{A} \AgdaSymbol{=} \AgdaRecord{Ty} \AgdaSymbol{(}\AgdaBound{Γ} \AgdaFunction{\&} \AgdaBound{A} \AgdaFunction{\&} \AgdaBound{A} \AgdaFunction{[} \AgdaFunction{fst\&} \AgdaSymbol{\{}A \AgdaSymbol{=} \AgdaBound{A}\AgdaSymbol{\}} \AgdaFunction{]T}\AgdaSymbol{)}\<%
\\
\>\<\end{code}

Natural numbers

Axiom: irrelevant:

\begin{code}\>\<%
\\
\>\AgdaKeyword{postulate}\<%
\\
\>[0]\AgdaIndent{4}{}\<[4]%
\>[4]\AgdaSymbol{.}\AgdaPostulate{irrelevant} \AgdaSymbol{:} \AgdaSymbol{\{}\AgdaBound{A} \AgdaSymbol{:} \AgdaPrimitiveType{Set}\AgdaSymbol{\}} \AgdaSymbol{→} \AgdaSymbol{.}\AgdaBound{A} \AgdaSymbol{→} \AgdaBound{A}\<%
\\
\>\<\end{code}


\begin{code}\>\<%
\\
\>\AgdaKeyword{module} \AgdaModule{Natural} \AgdaSymbol{(}\AgdaBound{Γ} \AgdaSymbol{:} \AgdaFunction{Con}\AgdaSymbol{)} \AgdaKeyword{where}\<%
\\
%
\\
\>[0]\AgdaIndent{2}{}\<[2]%
\>[2]\AgdaFunction{\_≈nat\_} \AgdaSymbol{:} \AgdaDatatype{ℕ} \AgdaSymbol{→} \AgdaDatatype{ℕ} \AgdaSymbol{→} \AgdaPrimitiveType{Set}\<%
\\
\>[0]\AgdaIndent{2}{}\<[2]%
\>[2]\AgdaInductiveConstructor{zero} \AgdaFunction{≈nat} \AgdaInductiveConstructor{zero} \AgdaSymbol{=} \AgdaRecord{⊤}\<%
\\
\>[0]\AgdaIndent{2}{}\<[2]%
\>[2]\AgdaInductiveConstructor{zero} \AgdaFunction{≈nat} \AgdaInductiveConstructor{suc} \AgdaBound{n} \AgdaSymbol{=} \AgdaDatatype{⊥}\<%
\\
\>[0]\AgdaIndent{2}{}\<[2]%
\>[2]\AgdaInductiveConstructor{suc} \AgdaBound{m} \AgdaFunction{≈nat} \AgdaInductiveConstructor{zero} \AgdaSymbol{=} \AgdaDatatype{⊥}\<%
\\
\>[0]\AgdaIndent{2}{}\<[2]%
\>[2]\AgdaInductiveConstructor{suc} \AgdaBound{m} \AgdaFunction{≈nat} \AgdaInductiveConstructor{suc} \AgdaBound{n} \AgdaSymbol{=} \AgdaBound{m} \AgdaFunction{≈nat} \AgdaBound{n}\<%
\\
\>[0]\AgdaIndent{2}{}\<[2]%
\>[2]\<%
\\
\>[0]\AgdaIndent{2}{}\<[2]%
\>[2]\AgdaFunction{reflNat} \AgdaSymbol{:} \AgdaSymbol{\{}\AgdaBound{x} \AgdaSymbol{:} \AgdaDatatype{ℕ}\AgdaSymbol{\}} \AgdaSymbol{→} \AgdaBound{x} \AgdaFunction{≈nat} \AgdaBound{x}\<%
\\
\>[0]\AgdaIndent{2}{}\<[2]%
\>[2]\AgdaFunction{reflNat} \AgdaSymbol{\{}\AgdaInductiveConstructor{zero}\AgdaSymbol{\}} \AgdaSymbol{=} \AgdaInductiveConstructor{tt}\<%
\\
\>[0]\AgdaIndent{2}{}\<[2]%
\>[2]\AgdaFunction{reflNat} \AgdaSymbol{\{}\AgdaInductiveConstructor{suc} \AgdaBound{n}\AgdaSymbol{\}} \AgdaSymbol{=} \AgdaFunction{reflNat} \AgdaSymbol{\{}\AgdaBound{n}\AgdaSymbol{\}}\<%
\\
%
\\
\>[0]\AgdaIndent{2}{}\<[2]%
\>[2]\AgdaFunction{symNat} \AgdaSymbol{:} \AgdaSymbol{\{}\AgdaBound{x} \AgdaBound{y} \AgdaSymbol{:} \AgdaDatatype{ℕ}\AgdaSymbol{\}} \AgdaSymbol{→} \AgdaBound{x} \AgdaFunction{≈nat} \AgdaBound{y} \AgdaSymbol{→} \AgdaBound{y} \AgdaFunction{≈nat} \AgdaBound{x}\<%
\\
\>[0]\AgdaIndent{2}{}\<[2]%
\>[2]\AgdaFunction{symNat} \AgdaSymbol{\{}\AgdaInductiveConstructor{zero}\AgdaSymbol{\}} \AgdaSymbol{\{}\AgdaInductiveConstructor{zero}\AgdaSymbol{\}} \AgdaBound{eq} \AgdaSymbol{=} \AgdaInductiveConstructor{tt}\<%
\\
\>[0]\AgdaIndent{2}{}\<[2]%
\>[2]\AgdaFunction{symNat} \AgdaSymbol{\{}\AgdaInductiveConstructor{zero}\AgdaSymbol{\}} \AgdaSymbol{\{}\AgdaInductiveConstructor{suc} \AgdaSymbol{\_\}} \AgdaBound{eq} \AgdaSymbol{=} \AgdaBound{eq}\<%
\\
\>[0]\AgdaIndent{2}{}\<[2]%
\>[2]\AgdaFunction{symNat} \AgdaSymbol{\{}\AgdaInductiveConstructor{suc} \AgdaSymbol{\_\}} \AgdaSymbol{\{}\AgdaInductiveConstructor{zero}\AgdaSymbol{\}} \AgdaBound{eq} \AgdaSymbol{=} \AgdaBound{eq}\<%
\\
\>[0]\AgdaIndent{2}{}\<[2]%
\>[2]\AgdaFunction{symNat} \AgdaSymbol{\{}\AgdaInductiveConstructor{suc} \AgdaBound{x}\AgdaSymbol{\}} \AgdaSymbol{\{}\AgdaInductiveConstructor{suc} \AgdaBound{y}\AgdaSymbol{\}} \AgdaBound{eq} \AgdaSymbol{=} \AgdaFunction{symNat} \AgdaSymbol{\{}\AgdaBound{x}\AgdaSymbol{\}} \AgdaSymbol{\{}\AgdaBound{y}\AgdaSymbol{\}} \AgdaBound{eq}\<%
\\
%
\\
\>[0]\AgdaIndent{2}{}\<[2]%
\>[2]\AgdaFunction{transNat} \AgdaSymbol{:} \AgdaSymbol{\{}\AgdaBound{x} \AgdaBound{y} \AgdaBound{z} \AgdaSymbol{:} \AgdaDatatype{ℕ}\AgdaSymbol{\}}\<%
\\
\>[2]\AgdaIndent{11}{}\<[11]%
\>[11]\AgdaSymbol{→} \AgdaBound{x} \AgdaFunction{≈nat} \AgdaBound{y} \AgdaSymbol{→} \AgdaBound{y} \AgdaFunction{≈nat} \AgdaBound{z} \AgdaSymbol{→} \AgdaBound{x} \AgdaFunction{≈nat} \AgdaBound{z}\<%
\\
\>[0]\AgdaIndent{2}{}\<[2]%
\>[2]\AgdaFunction{transNat} \AgdaSymbol{\{}\AgdaInductiveConstructor{zero}\AgdaSymbol{\}} \AgdaSymbol{\{}\AgdaInductiveConstructor{zero}\AgdaSymbol{\}} \AgdaBound{xy} \AgdaBound{yz} \AgdaSymbol{=} \AgdaBound{yz}\<%
\\
\>[0]\AgdaIndent{2}{}\<[2]%
\>[2]\AgdaFunction{transNat} \AgdaSymbol{\{}\AgdaInductiveConstructor{zero}\AgdaSymbol{\}} \AgdaSymbol{\{}\AgdaInductiveConstructor{suc} \AgdaSymbol{\_\}} \AgdaSymbol{()} \AgdaBound{yz}\<%
\\
\>[0]\AgdaIndent{2}{}\<[2]%
\>[2]\AgdaFunction{transNat} \AgdaSymbol{\{}\AgdaInductiveConstructor{suc} \AgdaSymbol{\_\}} \AgdaSymbol{\{}\AgdaInductiveConstructor{zero}\AgdaSymbol{\}} \AgdaSymbol{()} \AgdaBound{yz}\<%
\\
\>[0]\AgdaIndent{2}{}\<[2]%
\>[2]\AgdaFunction{transNat} \AgdaSymbol{\{}\AgdaInductiveConstructor{suc} \AgdaSymbol{\_\}} \AgdaSymbol{\{}\AgdaInductiveConstructor{suc} \AgdaSymbol{\_\}} \AgdaSymbol{\{}\AgdaInductiveConstructor{zero}\AgdaSymbol{\}} \AgdaBound{xy} \AgdaBound{yz} \AgdaSymbol{=} \AgdaBound{yz}\<%
\\
\>[0]\AgdaIndent{2}{}\<[2]%
\>[2]\AgdaFunction{transNat} \AgdaSymbol{\{}\AgdaInductiveConstructor{suc} \AgdaBound{x}\AgdaSymbol{\}} \AgdaSymbol{\{}\AgdaInductiveConstructor{suc} \AgdaBound{y}\AgdaSymbol{\}} \AgdaSymbol{\{}\AgdaInductiveConstructor{suc} \AgdaBound{z}\AgdaSymbol{\}} \AgdaBound{xy} \AgdaBound{yz} \AgdaSymbol{=}\<%
\\
\>[2]\AgdaIndent{11}{}\<[11]%
\>[11]\AgdaFunction{transNat} \AgdaSymbol{\{}\AgdaBound{x}\AgdaSymbol{\}} \AgdaSymbol{\{}\AgdaBound{y}\AgdaSymbol{\}} \AgdaSymbol{\{}\AgdaBound{z}\AgdaSymbol{\}} \AgdaBound{xy} \AgdaBound{yz}\<%
\\
\>[6]\AgdaIndent{3}{}\<[3]%
\>[3]\<%
\\
%
\\
\>[0]\AgdaIndent{2}{}\<[2]%
\>[2]\AgdaFunction{⟦Nat⟧} \AgdaSymbol{:} \AgdaRecord{Ty} \AgdaBound{Γ}\<%
\\
\>[0]\AgdaIndent{2}{}\<[2]%
\>[2]\AgdaFunction{⟦Nat⟧} \AgdaSymbol{=} \AgdaKeyword{record} \<[17]%
\>[17]\<%
\\
\>[2]\AgdaIndent{4}{}\<[4]%
\>[4]\AgdaSymbol{\{} \AgdaField{fm} \AgdaSymbol{=} \AgdaSymbol{λ} \AgdaBound{γ} \AgdaSymbol{→} \AgdaKeyword{record}\<%
\\
\>[4]\AgdaIndent{9}{}\<[9]%
\>[9]\AgdaSymbol{\{} \AgdaField{Carrier} \AgdaSymbol{=} \AgdaDatatype{ℕ}\<%
\\
\>[4]\AgdaIndent{9}{}\<[9]%
\>[9]\AgdaSymbol{;} \AgdaField{\_≈\_} \AgdaSymbol{=} \AgdaFunction{\_≈nat\_}\<%
\\
\>[4]\AgdaIndent{9}{}\<[9]%
\>[9]\AgdaSymbol{;} \AgdaField{refl} \AgdaSymbol{=} \AgdaSymbol{λ} \AgdaSymbol{\{}\AgdaBound{n}\AgdaSymbol{\}} \AgdaSymbol{→} \AgdaFunction{reflNat} \AgdaSymbol{\{}\AgdaBound{n}\AgdaSymbol{\}}\<%
\\
\>[4]\AgdaIndent{9}{}\<[9]%
\>[9]\AgdaSymbol{;} \AgdaField{sym} \AgdaSymbol{=} \AgdaSymbol{λ} \AgdaSymbol{\{}\AgdaBound{x}\AgdaSymbol{\}} \AgdaSymbol{\{}\AgdaBound{y}\AgdaSymbol{\}} \AgdaSymbol{→} \AgdaFunction{symNat} \AgdaSymbol{\{}\AgdaBound{x}\AgdaSymbol{\}} \AgdaSymbol{\{}\AgdaBound{y}\AgdaSymbol{\}}\<%
\\
\>[4]\AgdaIndent{9}{}\<[9]%
\>[9]\AgdaSymbol{;} \AgdaField{trans} \AgdaSymbol{=} \AgdaSymbol{λ} \AgdaSymbol{\{}\AgdaBound{x}\AgdaSymbol{\}} \AgdaSymbol{\{}\AgdaBound{y}\AgdaSymbol{\}} \AgdaSymbol{\{}\AgdaBound{z}\AgdaSymbol{\}} \AgdaSymbol{→} \AgdaFunction{transNat} \AgdaSymbol{\{}\AgdaBound{x}\AgdaSymbol{\}} \AgdaSymbol{\{}\AgdaBound{y}\AgdaSymbol{\}} \AgdaSymbol{\{}\AgdaBound{z}\AgdaSymbol{\}}\<%
\\
\>[4]\AgdaIndent{9}{}\<[9]%
\>[9]\AgdaSymbol{\}}\<%
\\
\>[0]\AgdaIndent{4}{}\<[4]%
\>[4]\AgdaSymbol{;} \AgdaField{substT} \AgdaSymbol{=} \AgdaSymbol{λ} \AgdaBound{\_} \AgdaBound{n} \AgdaSymbol{→} \AgdaBound{n}\<%
\\
\>[0]\AgdaIndent{4}{}\<[4]%
\>[4]\AgdaSymbol{;} \AgdaField{subst*} \AgdaSymbol{=} \AgdaSymbol{λ} \AgdaBound{\_} \AgdaBound{x} \AgdaSymbol{→} \AgdaPostulate{irrelevant} \AgdaBound{x}\<%
\\
\>[0]\AgdaIndent{4}{}\<[4]%
\>[4]\AgdaSymbol{;} \AgdaField{refl*} \AgdaSymbol{=} \AgdaSymbol{λ} \AgdaSymbol{\{}\AgdaBound{x}\AgdaSymbol{\}} \AgdaSymbol{\{}\AgdaBound{a}\AgdaSymbol{\}} \AgdaSymbol{→} \AgdaFunction{reflNat} \AgdaSymbol{\{}\AgdaBound{a}\AgdaSymbol{\}}\<%
\\
\>[0]\AgdaIndent{4}{}\<[4]%
\>[4]\AgdaSymbol{;} \AgdaField{trans*} \AgdaSymbol{=} \AgdaSymbol{λ} \AgdaBound{a} \AgdaSymbol{→} \AgdaFunction{reflNat} \AgdaSymbol{\{}\AgdaBound{a}\AgdaSymbol{\}} \<[33]%
\>[33]\<%
\\
\>[0]\AgdaIndent{4}{}\<[4]%
\>[4]\AgdaSymbol{\}}\<%
\\
%
\\
\>[0]\AgdaIndent{2}{}\<[2]%
\>[2]\AgdaFunction{⟦0⟧} \AgdaSymbol{:} \AgdaRecord{Tm} \AgdaFunction{⟦Nat⟧}\<%
\\
\>[0]\AgdaIndent{2}{}\<[2]%
\>[2]\AgdaFunction{⟦0⟧} \AgdaSymbol{=} \AgdaKeyword{record}\<%
\\
\>[2]\AgdaIndent{6}{}\<[6]%
\>[6]\AgdaSymbol{\{} \AgdaField{tm} \AgdaSymbol{=} \AgdaSymbol{λ} \AgdaBound{\_} \AgdaSymbol{→} \AgdaNumber{0}\<%
\\
\>[2]\AgdaIndent{6}{}\<[6]%
\>[6]\AgdaSymbol{;} \AgdaField{respt} \AgdaSymbol{=} \AgdaSymbol{λ} \AgdaBound{p} \AgdaSymbol{→} \AgdaInductiveConstructor{tt}\<%
\\
\>[2]\AgdaIndent{6}{}\<[6]%
\>[6]\AgdaSymbol{\}}\<%
\\
%
\\
\>[0]\AgdaIndent{2}{}\<[2]%
\>[2]\AgdaFunction{⟦s⟧} \AgdaSymbol{:} \AgdaRecord{Tm} \AgdaFunction{⟦Nat⟧} \AgdaSymbol{→} \AgdaRecord{Tm} \AgdaFunction{⟦Nat⟧}\<%
\\
\>[0]\AgdaIndent{2}{}\<[2]%
\>[2]\AgdaFunction{⟦s⟧} \AgdaSymbol{(}\AgdaInductiveConstructor{tm:} \AgdaBound{t} \AgdaInductiveConstructor{resp:} \AgdaBound{respt}\AgdaSymbol{)} \<[26]%
\>[26]\<%
\\
\>[2]\AgdaIndent{6}{}\<[6]%
\>[6]\AgdaSymbol{=} \AgdaKeyword{record}\<%
\\
\>[2]\AgdaIndent{6}{}\<[6]%
\>[6]\AgdaSymbol{\{} \AgdaField{tm} \AgdaSymbol{=} \AgdaInductiveConstructor{suc} \AgdaFunction{∘} \AgdaBound{t}\<%
\\
\>[2]\AgdaIndent{6}{}\<[6]%
\>[6]\AgdaSymbol{;} \AgdaField{respt} \AgdaSymbol{=} \AgdaBound{respt}\<%
\\
\>[2]\AgdaIndent{6}{}\<[6]%
\>[6]\AgdaSymbol{\}}\<%
\\
\>\<\end{code}

Simply typed universe

\AgdaHide{
\begin{code}\>\<%
\\
\>[0]\AgdaIndent{2}{}\<[2]%
\>[2]\AgdaKeyword{data} \AgdaDatatype{⟦U⟧⁰} \AgdaSymbol{:} \AgdaPrimitiveType{Set} \AgdaKeyword{where}\<%
\\
\>[0]\AgdaIndent{4}{}\<[4]%
\>[4]\AgdaInductiveConstructor{nat} \AgdaSymbol{:} \AgdaDatatype{⟦U⟧⁰}\<%
\\
\>[0]\AgdaIndent{4}{}\<[4]%
\>[4]\AgdaInductiveConstructor{arr<\_,\_>} \AgdaSymbol{:} \AgdaSymbol{(}\AgdaBound{a} \AgdaBound{b} \AgdaSymbol{:} \AgdaDatatype{⟦U⟧⁰}\AgdaSymbol{)} \AgdaSymbol{→} \AgdaDatatype{⟦U⟧⁰}\<%
\\
%
\\
\>[0]\AgdaIndent{2}{}\<[2]%
\>[2]\AgdaFunction{\_\textasciitilde⟦U⟧\_} \AgdaSymbol{:} \AgdaDatatype{⟦U⟧⁰} \AgdaSymbol{→} \AgdaDatatype{⟦U⟧⁰} \AgdaSymbol{→} \AgdaPrimitiveType{Set} \AgdaComment{-- HProp}\<%
\\
\>[0]\AgdaIndent{2}{}\<[2]%
\>[2]\AgdaInductiveConstructor{nat} \AgdaFunction{\textasciitilde⟦U⟧} \AgdaInductiveConstructor{nat} \AgdaSymbol{=} \AgdaRecord{⊤}\<%
\\
\>[0]\AgdaIndent{2}{}\<[2]%
\>[2]\AgdaInductiveConstructor{nat} \AgdaFunction{\textasciitilde⟦U⟧} \AgdaInductiveConstructor{arr<} \AgdaBound{a} \AgdaInductiveConstructor{,} \AgdaBound{b} \AgdaInductiveConstructor{>} \AgdaSymbol{=} \AgdaDatatype{⊥}\<%
\\
\>[0]\AgdaIndent{2}{}\<[2]%
\>[2]\AgdaInductiveConstructor{arr<} \AgdaBound{a} \AgdaInductiveConstructor{,} \AgdaBound{b} \AgdaInductiveConstructor{>} \AgdaFunction{\textasciitilde⟦U⟧} \AgdaInductiveConstructor{nat} \AgdaSymbol{=} \AgdaDatatype{⊥}\<%
\\
\>[0]\AgdaIndent{2}{}\<[2]%
\>[2]\AgdaInductiveConstructor{arr<} \AgdaBound{a} \AgdaInductiveConstructor{,} \AgdaBound{b} \AgdaInductiveConstructor{>} \AgdaFunction{\textasciitilde⟦U⟧} \AgdaInductiveConstructor{arr<} \AgdaBound{a'} \AgdaInductiveConstructor{,} \AgdaBound{b'} \AgdaInductiveConstructor{>} \AgdaSymbol{=} \AgdaBound{a} \AgdaFunction{\textasciitilde⟦U⟧} \AgdaBound{a'} \AgdaFunction{×} \AgdaBound{b} \AgdaFunction{\textasciitilde⟦U⟧} \AgdaBound{b'}\<%
\\
%
\\
\>[0]\AgdaIndent{2}{}\<[2]%
\>[2]\AgdaFunction{reflU} \AgdaSymbol{:} \<[11]%
\>[11]\AgdaSymbol{\{}\AgdaBound{x} \AgdaSymbol{:} \AgdaDatatype{⟦U⟧⁰}\AgdaSymbol{\}} \AgdaSymbol{→} \AgdaBound{x} \AgdaFunction{\textasciitilde⟦U⟧} \AgdaBound{x}\<%
\\
\>[0]\AgdaIndent{2}{}\<[2]%
\>[2]\AgdaFunction{reflU} \AgdaSymbol{\{}\AgdaInductiveConstructor{nat}\AgdaSymbol{\}} \AgdaSymbol{=} \AgdaInductiveConstructor{tt}\<%
\\
\>[0]\AgdaIndent{2}{}\<[2]%
\>[2]\AgdaFunction{reflU} \AgdaSymbol{\{}\AgdaInductiveConstructor{arr<} \AgdaBound{a} \AgdaInductiveConstructor{,} \AgdaBound{b} \AgdaInductiveConstructor{>}\AgdaSymbol{\}} \AgdaSymbol{=} \AgdaFunction{reflU} \AgdaSymbol{\{}\AgdaBound{a}\AgdaSymbol{\}} \AgdaInductiveConstructor{,} \AgdaFunction{reflU} \AgdaSymbol{\{}\AgdaBound{b}\AgdaSymbol{\}}\<%
\\
%
\\
\>[0]\AgdaIndent{2}{}\<[2]%
\>[2]\AgdaFunction{symU} \AgdaSymbol{:} \AgdaSymbol{\{}\AgdaBound{x} \AgdaBound{y} \AgdaSymbol{:} \AgdaDatatype{⟦U⟧⁰}\AgdaSymbol{\}} \AgdaSymbol{→} \AgdaBound{x} \AgdaFunction{\textasciitilde⟦U⟧} \AgdaBound{y} \AgdaSymbol{→} \AgdaBound{y} \AgdaFunction{\textasciitilde⟦U⟧} \AgdaBound{x}\<%
\\
\>[0]\AgdaIndent{2}{}\<[2]%
\>[2]\AgdaFunction{symU} \AgdaSymbol{\{}\AgdaInductiveConstructor{nat}\AgdaSymbol{\}} \AgdaSymbol{\{}\AgdaInductiveConstructor{nat}\AgdaSymbol{\}} \AgdaBound{eq} \AgdaSymbol{=} \AgdaInductiveConstructor{tt}\<%
\\
\>[0]\AgdaIndent{2}{}\<[2]%
\>[2]\AgdaFunction{symU} \AgdaSymbol{\{}\AgdaInductiveConstructor{nat}\AgdaSymbol{\}} \AgdaSymbol{\{}\AgdaInductiveConstructor{arr<} \AgdaBound{a} \AgdaInductiveConstructor{,} \AgdaBound{b} \AgdaInductiveConstructor{>}\AgdaSymbol{\}} \AgdaBound{eq} \AgdaSymbol{=} \AgdaBound{eq}\<%
\\
\>[0]\AgdaIndent{2}{}\<[2]%
\>[2]\AgdaFunction{symU} \AgdaSymbol{\{}\AgdaInductiveConstructor{arr<} \AgdaBound{a} \AgdaInductiveConstructor{,} \AgdaBound{b} \AgdaInductiveConstructor{>}\AgdaSymbol{\}} \AgdaSymbol{\{}\AgdaInductiveConstructor{nat}\AgdaSymbol{\}} \AgdaBound{eq} \AgdaSymbol{=} \AgdaBound{eq}\<%
\\
\>[0]\AgdaIndent{2}{}\<[2]%
\>[2]\AgdaFunction{symU} \AgdaSymbol{\{}\AgdaInductiveConstructor{arr<} \AgdaBound{a} \AgdaInductiveConstructor{,} \AgdaBound{b} \AgdaInductiveConstructor{>}\AgdaSymbol{\}} \AgdaSymbol{\{}\AgdaInductiveConstructor{arr<} \AgdaBound{a'} \AgdaInductiveConstructor{,} \AgdaBound{b'} \AgdaInductiveConstructor{>}\AgdaSymbol{\}} \AgdaSymbol{(}\AgdaBound{p} \AgdaInductiveConstructor{,} \AgdaBound{q}\AgdaSymbol{)} \AgdaSymbol{=} \AgdaSymbol{(}\AgdaFunction{symU} \AgdaSymbol{\{}\AgdaBound{a}\AgdaSymbol{\}} \AgdaSymbol{\{}\AgdaBound{a'}\AgdaSymbol{\}} \AgdaBound{p}\AgdaSymbol{)} \<[67]%
\>[67]\<%
\\
\>[2]\AgdaIndent{47}{}\<[47]%
\>[47]\AgdaInductiveConstructor{,} \AgdaSymbol{(}\AgdaFunction{symU} \AgdaSymbol{\{}\AgdaBound{b}\AgdaSymbol{\}} \AgdaSymbol{\{}\AgdaBound{b'}\AgdaSymbol{\}} \AgdaBound{q}\AgdaSymbol{)}\<%
\\
%
\\
\>[0]\AgdaIndent{2}{}\<[2]%
\>[2]\AgdaFunction{transU} \AgdaSymbol{:} \AgdaSymbol{\{}\AgdaBound{x} \AgdaBound{y} \AgdaBound{z} \AgdaSymbol{:} \AgdaDatatype{⟦U⟧⁰}\AgdaSymbol{\}} \AgdaSymbol{→} \AgdaBound{x} \AgdaFunction{\textasciitilde⟦U⟧} \AgdaBound{y} \AgdaSymbol{→} \AgdaBound{y} \AgdaFunction{\textasciitilde⟦U⟧} \AgdaBound{z} \AgdaSymbol{→} \AgdaBound{x} \AgdaFunction{\textasciitilde⟦U⟧} \AgdaBound{z}\<%
\\
\>[0]\AgdaIndent{2}{}\<[2]%
\>[2]\AgdaFunction{transU} \AgdaSymbol{\{}\AgdaInductiveConstructor{nat}\AgdaSymbol{\}} \AgdaSymbol{\{}\AgdaInductiveConstructor{nat}\AgdaSymbol{\}} \AgdaBound{eq1} \AgdaBound{eq2} \AgdaSymbol{=} \AgdaBound{eq2}\<%
\\
\>[0]\AgdaIndent{2}{}\<[2]%
\>[2]\AgdaFunction{transU} \AgdaSymbol{\{}\AgdaInductiveConstructor{nat}\AgdaSymbol{\}} \AgdaSymbol{\{}\AgdaInductiveConstructor{arr<} \AgdaBound{a} \AgdaInductiveConstructor{,} \AgdaBound{b} \AgdaInductiveConstructor{>}\AgdaSymbol{\}} \AgdaSymbol{()} \AgdaBound{eq2}\<%
\\
\>[0]\AgdaIndent{2}{}\<[2]%
\>[2]\AgdaFunction{transU} \AgdaSymbol{\{}\AgdaInductiveConstructor{arr<} \AgdaBound{a} \AgdaInductiveConstructor{,} \AgdaBound{b} \AgdaInductiveConstructor{>}\AgdaSymbol{\}} \AgdaSymbol{\{}\AgdaInductiveConstructor{nat}\AgdaSymbol{\}} \AgdaSymbol{()} \AgdaBound{eq2}\<%
\\
\>[0]\AgdaIndent{2}{}\<[2]%
\>[2]\AgdaFunction{transU} \AgdaSymbol{\{}\AgdaInductiveConstructor{arr<} \AgdaBound{a} \AgdaInductiveConstructor{,} \AgdaBound{b} \AgdaInductiveConstructor{>}\AgdaSymbol{\}} \AgdaSymbol{\{}\AgdaInductiveConstructor{arr<} \AgdaBound{a'} \AgdaInductiveConstructor{,} \AgdaBound{b'} \AgdaInductiveConstructor{>}\AgdaSymbol{\}} \AgdaSymbol{\{}\AgdaInductiveConstructor{nat}\AgdaSymbol{\}} \AgdaBound{eq1} \AgdaBound{eq2} \AgdaSymbol{=} \AgdaBound{eq2}\<%
\\
\>[0]\AgdaIndent{2}{}\<[2]%
\>[2]\AgdaFunction{transU} \AgdaSymbol{\{}\AgdaInductiveConstructor{arr<} \AgdaBound{a} \AgdaInductiveConstructor{,} \AgdaBound{b} \AgdaInductiveConstructor{>}\AgdaSymbol{\}} \AgdaSymbol{\{}\AgdaInductiveConstructor{arr<} \AgdaBound{a'} \AgdaInductiveConstructor{,} \AgdaBound{b'} \AgdaInductiveConstructor{>}\AgdaSymbol{\}} \AgdaSymbol{\{}\AgdaInductiveConstructor{arr<} \AgdaBound{a0} \AgdaInductiveConstructor{,} \AgdaBound{b0} \AgdaInductiveConstructor{>}\AgdaSymbol{\}} \AgdaSymbol{(}\AgdaBound{p1} \AgdaInductiveConstructor{,} \AgdaBound{q1}\AgdaSymbol{)} \<[68]%
\>[68]\<%
\\
\>[2]\AgdaIndent{9}{}\<[9]%
\>[9]\AgdaSymbol{(}\AgdaBound{p2} \AgdaInductiveConstructor{,} \AgdaBound{q2}\AgdaSymbol{)} \AgdaSymbol{=} \AgdaSymbol{(}\AgdaFunction{transU} \AgdaSymbol{\{}\AgdaBound{a}\AgdaSymbol{\}} \AgdaSymbol{\{}\AgdaBound{a'}\AgdaSymbol{\}} \AgdaSymbol{\{}\AgdaBound{a0}\AgdaSymbol{\}} \AgdaBound{p1} \AgdaBound{p2}\AgdaSymbol{)} \<[50]%
\>[50]\<%
\\
\>[2]\AgdaIndent{9}{}\<[9]%
\>[9]\AgdaInductiveConstructor{,} \AgdaFunction{transU} \AgdaSymbol{\{}\AgdaBound{b}\AgdaSymbol{\}} \AgdaSymbol{\{}\AgdaBound{b'}\AgdaSymbol{\}} \AgdaSymbol{\{}\AgdaBound{b0}\AgdaSymbol{\}} \AgdaBound{q1} \AgdaBound{q2}\<%
\\
%
\\
\>[0]\AgdaIndent{2}{}\<[2]%
\>[2]\AgdaFunction{⟦U⟧} \AgdaSymbol{:} \AgdaRecord{Ty} \AgdaBound{Γ}\<%
\\
\>[0]\AgdaIndent{2}{}\<[2]%
\>[2]\AgdaFunction{⟦U⟧} \AgdaSymbol{=} \AgdaKeyword{record} \<[15]%
\>[15]\<%
\\
\>[2]\AgdaIndent{4}{}\<[4]%
\>[4]\AgdaSymbol{\{} \AgdaField{fm} \AgdaSymbol{=} \AgdaSymbol{λ} \AgdaBound{γ} \AgdaSymbol{→} \AgdaKeyword{record}\<%
\\
\>[4]\AgdaIndent{9}{}\<[9]%
\>[9]\AgdaSymbol{\{} \AgdaField{Carrier} \AgdaSymbol{=} \AgdaDatatype{⟦U⟧⁰}\<%
\\
\>[4]\AgdaIndent{9}{}\<[9]%
\>[9]\AgdaSymbol{;} \AgdaField{\_≈\_} \AgdaSymbol{=} \AgdaFunction{\_\textasciitilde⟦U⟧\_}\<%
\\
\>[4]\AgdaIndent{9}{}\<[9]%
\>[9]\AgdaSymbol{;} \AgdaField{refl} \AgdaSymbol{=} \AgdaSymbol{λ} \AgdaSymbol{\{}\AgdaBound{x}\AgdaSymbol{\}} \AgdaSymbol{→} \AgdaFunction{reflU} \AgdaSymbol{\{}\AgdaBound{x}\AgdaSymbol{\}}\<%
\\
\>[4]\AgdaIndent{9}{}\<[9]%
\>[9]\AgdaSymbol{;} \AgdaField{sym} \AgdaSymbol{=} \AgdaSymbol{λ} \AgdaSymbol{\{}\AgdaBound{x}\AgdaSymbol{\}} \AgdaSymbol{\{}\AgdaBound{y}\AgdaSymbol{\}} \AgdaSymbol{→} \AgdaFunction{symU} \AgdaSymbol{\{}\AgdaBound{x}\AgdaSymbol{\}} \AgdaSymbol{\{}\AgdaBound{y}\AgdaSymbol{\}}\<%
\\
\>[4]\AgdaIndent{9}{}\<[9]%
\>[9]\AgdaSymbol{;} \AgdaField{trans} \AgdaSymbol{=} \AgdaSymbol{λ} \AgdaSymbol{\{}\AgdaBound{x}\AgdaSymbol{\}} \AgdaSymbol{\{}\AgdaBound{y}\AgdaSymbol{\}} \AgdaSymbol{\{}\AgdaBound{z}\AgdaSymbol{\}} \AgdaSymbol{→} \AgdaFunction{transU} \AgdaSymbol{\{}\AgdaBound{x}\AgdaSymbol{\}} \AgdaSymbol{\{}\AgdaBound{y}\AgdaSymbol{\}} \AgdaSymbol{\{}\AgdaBound{z}\AgdaSymbol{\}}\<%
\\
\>[4]\AgdaIndent{9}{}\<[9]%
\>[9]\AgdaSymbol{\}}\<%
\\
\>[0]\AgdaIndent{4}{}\<[4]%
\>[4]\AgdaSymbol{;} \AgdaField{substT} \AgdaSymbol{=} \AgdaSymbol{λ} \AgdaBound{\_} \AgdaBound{x} \AgdaSymbol{→} \AgdaBound{x}\<%
\\
\>[0]\AgdaIndent{4}{}\<[4]%
\>[4]\AgdaSymbol{;} \AgdaField{subst*} \AgdaSymbol{=} \<[16]%
\>[16]\AgdaSymbol{λ} \AgdaBound{\_} \AgdaBound{x} \AgdaSymbol{→} \AgdaPostulate{irrelevant} \AgdaBound{x}\<%
\\
\>[0]\AgdaIndent{4}{}\<[4]%
\>[4]\AgdaSymbol{;} \AgdaField{refl*} \AgdaSymbol{=} \AgdaSymbol{λ} \AgdaSymbol{\{}\AgdaBound{x}\AgdaSymbol{\}} \AgdaSymbol{\{}\AgdaBound{a}\AgdaSymbol{\}} \AgdaSymbol{→} \AgdaFunction{reflU} \AgdaSymbol{\{}\AgdaBound{a}\AgdaSymbol{\}}\<%
\\
\>[0]\AgdaIndent{4}{}\<[4]%
\>[4]\AgdaSymbol{;} \AgdaField{trans*} \AgdaSymbol{=} \AgdaSymbol{λ} \AgdaBound{a} \AgdaSymbol{→} \AgdaFunction{reflU} \AgdaSymbol{\{}\AgdaBound{a}\AgdaSymbol{\}}\<%
\\
\>[0]\AgdaIndent{4}{}\<[4]%
\>[4]\AgdaSymbol{\}}\<%
\\
%
\\
\>[0]\AgdaIndent{2}{}\<[2]%
\>[2]\AgdaFunction{elfm} \AgdaSymbol{:} \AgdaRecord{Σ} \AgdaFunction{∣} \AgdaBound{Γ} \AgdaFunction{∣} \AgdaSymbol{(λ} \AgdaBound{x} \AgdaSymbol{→} \AgdaDatatype{⟦U⟧⁰}\AgdaSymbol{)} \AgdaSymbol{→} \AgdaRecord{Setoid}\<%
\\
\>[0]\AgdaIndent{2}{}\<[2]%
\>[2]\AgdaFunction{elfm} \AgdaSymbol{(}\AgdaBound{γ} \AgdaInductiveConstructor{,} \AgdaInductiveConstructor{nat}\AgdaSymbol{)} \AgdaSymbol{=} \AgdaFunction{[} \AgdaFunction{⟦Nat⟧} \AgdaFunction{]fm} \AgdaBound{γ}\<%
\\
\>[0]\AgdaIndent{2}{}\<[2]%
\>[2]\AgdaFunction{elfm} \AgdaSymbol{(}\AgdaBound{γ} \AgdaInductiveConstructor{,} \AgdaInductiveConstructor{arr<} \AgdaBound{a} \AgdaInductiveConstructor{,} \AgdaBound{b} \AgdaInductiveConstructor{>}\AgdaSymbol{)} \AgdaSymbol{=} \AgdaFunction{[} \AgdaBound{Γ} \AgdaFunction{,} \AgdaBound{γ} \AgdaFunction{]} \AgdaFunction{elfm} \AgdaSymbol{(}\AgdaBound{γ} \AgdaInductiveConstructor{,} \AgdaBound{a}\AgdaSymbol{)} \AgdaFunction{⇒fm} \AgdaFunction{elfm} \AgdaSymbol{(}\AgdaBound{γ} \AgdaInductiveConstructor{,} \AgdaBound{b}\AgdaSymbol{)}\<%
\\
\>\<\end{code}
}

\section{Examples of types}

\AgdaHide{
\begin{code}\>\<%
\\
%
\\
\>\AgdaSymbol{\{-\#} \AgdaKeyword{OPTIONS} --type-in-type \AgdaSymbol{\#-\}}\<%
\\
%
\\
\>\AgdaKeyword{import} \AgdaModule{Level}\<%
\\
\>\AgdaKeyword{open} \AgdaKeyword{import} \AgdaModule{Relation.Binary.PropositionalEquality} \AgdaSymbol{as} \AgdaModule{PE} \AgdaKeyword{hiding} \AgdaSymbol{(}refl \AgdaSymbol{;} sym \AgdaSymbol{;} trans\AgdaSymbol{;} isEquivalence\AgdaSymbol{;} [\_]\AgdaSymbol{)}\<%
\\
%
\\
\>\AgdaKeyword{module} \AgdaModule{CwF-ctd} \AgdaSymbol{(}\AgdaBound{ext} \AgdaSymbol{:} \AgdaFunction{Extensionality} \AgdaPrimitive{Level.zero} \AgdaPrimitive{Level.zero}\AgdaSymbol{)} \AgdaKeyword{where}\<%
\\
%
\\
\>\AgdaKeyword{open} \AgdaKeyword{import} \AgdaModule{Data.Unit}\<%
\\
\>\AgdaKeyword{open} \AgdaKeyword{import} \AgdaModule{Function}\<%
\\
\>\AgdaKeyword{open} \AgdaKeyword{import} \AgdaModule{Data.Product}\<%
\\
%
\\
\>\AgdaKeyword{open} \AgdaKeyword{import} \AgdaModule{CwF-setoidwo} \AgdaBound{ext} \AgdaKeyword{public}\<%
\\
%
\\
\>\AgdaKeyword{open} \AgdaKeyword{import} \AgdaModule{Data.Nat}\<%
\\
%
\\
\>\<\end{code}
}

Binary relation

\begin{code}\>\<%
\\
\>\AgdaFunction{Rel} \AgdaSymbol{:} \AgdaSymbol{\{}\AgdaBound{Γ} \AgdaSymbol{:} \AgdaFunction{Con}\AgdaSymbol{\}} \AgdaSymbol{→} \AgdaRecord{Ty} \AgdaBound{Γ} \AgdaSymbol{→} \AgdaPrimitiveType{Set₁}\<%
\\
\>\AgdaFunction{Rel} \AgdaSymbol{\{}\AgdaBound{Γ}\AgdaSymbol{\}} \AgdaBound{A} \AgdaSymbol{=} \AgdaRecord{Ty} \AgdaSymbol{(}\AgdaBound{Γ} \AgdaFunction{\&} \AgdaBound{A} \AgdaFunction{\&} \AgdaBound{A} \AgdaFunction{[} \AgdaFunction{fst\&} \AgdaSymbol{\{}A \AgdaSymbol{=} \AgdaBound{A}\AgdaSymbol{\}} \AgdaFunction{]T}\AgdaSymbol{)}\<%
\\
\>\<\end{code}

Natural numbers

\begin{code}\>\<%
\\
\>\AgdaKeyword{module} \AgdaModule{Natural} \AgdaSymbol{(}\AgdaBound{Γ} \AgdaSymbol{:} \AgdaFunction{Con}\AgdaSymbol{)} \AgdaKeyword{where}\<%
\\
%
\\
\>[0]\AgdaIndent{2}{}\<[2]%
\>[2]\AgdaFunction{\_≈nat\_} \AgdaSymbol{:} \AgdaDatatype{ℕ} \AgdaSymbol{→} \AgdaDatatype{ℕ} \AgdaSymbol{→} \AgdaRecord{HProp}\<%
\\
\>[0]\AgdaIndent{2}{}\<[2]%
\>[2]\AgdaInductiveConstructor{zero} \AgdaFunction{≈nat} \AgdaInductiveConstructor{zero} \AgdaSymbol{=} \AgdaFunction{⊤'}\<%
\\
\>[0]\AgdaIndent{2}{}\<[2]%
\>[2]\AgdaInductiveConstructor{zero} \AgdaFunction{≈nat} \AgdaInductiveConstructor{suc} \AgdaBound{n} \AgdaSymbol{=} \AgdaFunction{⊥'}\<%
\\
\>[0]\AgdaIndent{2}{}\<[2]%
\>[2]\AgdaInductiveConstructor{suc} \AgdaBound{m} \AgdaFunction{≈nat} \AgdaInductiveConstructor{zero} \AgdaSymbol{=} \AgdaFunction{⊥'}\<%
\\
\>[0]\AgdaIndent{2}{}\<[2]%
\>[2]\AgdaInductiveConstructor{suc} \AgdaBound{m} \AgdaFunction{≈nat} \AgdaInductiveConstructor{suc} \AgdaBound{n} \AgdaSymbol{=} \AgdaBound{m} \AgdaFunction{≈nat} \AgdaBound{n}\<%
\\
\>[0]\AgdaIndent{2}{}\<[2]%
\>[2]\<%
\\
\>[0]\AgdaIndent{2}{}\<[2]%
\>[2]\AgdaFunction{reflNat} \AgdaSymbol{:} \AgdaSymbol{\{}\AgdaBound{x} \AgdaSymbol{:} \AgdaDatatype{ℕ}\AgdaSymbol{\}} \AgdaSymbol{→} \AgdaFunction{<} \AgdaBound{x} \AgdaFunction{≈nat} \AgdaBound{x} \AgdaFunction{>} \<[35]%
\>[35]\<%
\\
\>[0]\AgdaIndent{2}{}\<[2]%
\>[2]\AgdaFunction{reflNat} \AgdaSymbol{\{}\AgdaInductiveConstructor{zero}\AgdaSymbol{\}} \AgdaSymbol{=} \AgdaInductiveConstructor{tt}\<%
\\
\>[0]\AgdaIndent{2}{}\<[2]%
\>[2]\AgdaFunction{reflNat} \AgdaSymbol{\{}\AgdaInductiveConstructor{suc} \AgdaBound{n}\AgdaSymbol{\}} \AgdaSymbol{=} \AgdaFunction{reflNat} \AgdaSymbol{\{}\AgdaBound{n}\AgdaSymbol{\}}\<%
\\
%
\\
\>[0]\AgdaIndent{2}{}\<[2]%
\>[2]\AgdaFunction{symNat} \AgdaSymbol{:} \AgdaSymbol{\{}\AgdaBound{x} \AgdaBound{y} \AgdaSymbol{:} \AgdaDatatype{ℕ}\AgdaSymbol{\}} \AgdaSymbol{→} \AgdaFunction{<} \AgdaBound{x} \AgdaFunction{≈nat} \AgdaBound{y} \AgdaFunction{>} \AgdaSymbol{→} \AgdaFunction{<} \AgdaBound{y} \AgdaFunction{≈nat} \AgdaBound{x} \AgdaFunction{>}\<%
\\
\>[0]\AgdaIndent{2}{}\<[2]%
\>[2]\AgdaFunction{symNat} \AgdaSymbol{\{}\AgdaInductiveConstructor{zero}\AgdaSymbol{\}} \AgdaSymbol{\{}\AgdaInductiveConstructor{zero}\AgdaSymbol{\}} \AgdaBound{eq} \AgdaSymbol{=} \AgdaInductiveConstructor{tt}\<%
\\
\>[0]\AgdaIndent{2}{}\<[2]%
\>[2]\AgdaFunction{symNat} \AgdaSymbol{\{}\AgdaInductiveConstructor{zero}\AgdaSymbol{\}} \AgdaSymbol{\{}\AgdaInductiveConstructor{suc} \AgdaSymbol{\_\}} \AgdaBound{eq} \AgdaSymbol{=} \AgdaBound{eq}\<%
\\
\>[0]\AgdaIndent{2}{}\<[2]%
\>[2]\AgdaFunction{symNat} \AgdaSymbol{\{}\AgdaInductiveConstructor{suc} \AgdaSymbol{\_\}} \AgdaSymbol{\{}\AgdaInductiveConstructor{zero}\AgdaSymbol{\}} \AgdaBound{eq} \AgdaSymbol{=} \AgdaBound{eq}\<%
\\
\>[0]\AgdaIndent{2}{}\<[2]%
\>[2]\AgdaFunction{symNat} \AgdaSymbol{\{}\AgdaInductiveConstructor{suc} \AgdaBound{x}\AgdaSymbol{\}} \AgdaSymbol{\{}\AgdaInductiveConstructor{suc} \AgdaBound{y}\AgdaSymbol{\}} \AgdaBound{eq} \AgdaSymbol{=} \AgdaFunction{symNat} \AgdaSymbol{\{}\AgdaBound{x}\AgdaSymbol{\}} \AgdaSymbol{\{}\AgdaBound{y}\AgdaSymbol{\}} \AgdaBound{eq}\<%
\\
%
\\
\>[0]\AgdaIndent{2}{}\<[2]%
\>[2]\AgdaFunction{transNat} \AgdaSymbol{:} \AgdaSymbol{\{}\AgdaBound{x} \AgdaBound{y} \AgdaBound{z} \AgdaSymbol{:} \AgdaDatatype{ℕ}\AgdaSymbol{\}} \AgdaSymbol{→} \AgdaFunction{<} \AgdaBound{x} \AgdaFunction{≈nat} \AgdaBound{y} \AgdaFunction{>} \AgdaSymbol{→} \AgdaFunction{<} \AgdaBound{y} \AgdaFunction{≈nat} \AgdaBound{z} \AgdaFunction{>} \AgdaSymbol{→} \AgdaFunction{<} \AgdaBound{x} \AgdaFunction{≈nat} \AgdaBound{z} \AgdaFunction{>}\<%
\\
\>[0]\AgdaIndent{2}{}\<[2]%
\>[2]\AgdaFunction{transNat} \AgdaSymbol{\{}\AgdaInductiveConstructor{zero}\AgdaSymbol{\}} \AgdaSymbol{\{}\AgdaInductiveConstructor{zero}\AgdaSymbol{\}} \AgdaBound{xy} \AgdaBound{yz} \AgdaSymbol{=} \AgdaBound{yz}\<%
\\
\>[0]\AgdaIndent{2}{}\<[2]%
\>[2]\AgdaFunction{transNat} \AgdaSymbol{\{}\AgdaInductiveConstructor{zero}\AgdaSymbol{\}} \AgdaSymbol{\{}\AgdaInductiveConstructor{suc} \AgdaSymbol{\_\}} \AgdaSymbol{()} \AgdaBound{yz}\<%
\\
\>[0]\AgdaIndent{2}{}\<[2]%
\>[2]\AgdaFunction{transNat} \AgdaSymbol{\{}\AgdaInductiveConstructor{suc} \AgdaSymbol{\_\}} \AgdaSymbol{\{}\AgdaInductiveConstructor{zero}\AgdaSymbol{\}} \AgdaSymbol{()} \AgdaBound{yz}\<%
\\
\>[0]\AgdaIndent{2}{}\<[2]%
\>[2]\AgdaFunction{transNat} \AgdaSymbol{\{}\AgdaInductiveConstructor{suc} \AgdaSymbol{\_\}} \AgdaSymbol{\{}\AgdaInductiveConstructor{suc} \AgdaSymbol{\_\}} \AgdaSymbol{\{}\AgdaInductiveConstructor{zero}\AgdaSymbol{\}} \AgdaBound{xy} \AgdaBound{yz} \AgdaSymbol{=} \AgdaBound{yz}\<%
\\
\>[0]\AgdaIndent{2}{}\<[2]%
\>[2]\AgdaFunction{transNat} \AgdaSymbol{\{}\AgdaInductiveConstructor{suc} \AgdaBound{x}\AgdaSymbol{\}} \AgdaSymbol{\{}\AgdaInductiveConstructor{suc} \AgdaBound{y}\AgdaSymbol{\}} \AgdaSymbol{\{}\AgdaInductiveConstructor{suc} \AgdaBound{z}\AgdaSymbol{\}} \AgdaBound{xy} \AgdaBound{yz} \AgdaSymbol{=} \AgdaFunction{transNat} \AgdaSymbol{\{}\AgdaBound{x}\AgdaSymbol{\}} \AgdaSymbol{\{}\AgdaBound{y}\AgdaSymbol{\}} \AgdaSymbol{\{}\AgdaBound{z}\AgdaSymbol{\}} \AgdaBound{xy} \AgdaBound{yz}\<%
\\
%
\\
\>[0]\AgdaIndent{2}{}\<[2]%
\>[2]\AgdaFunction{⟦Nat⟧} \AgdaSymbol{:} \AgdaRecord{Ty} \AgdaBound{Γ}\<%
\\
\>[0]\AgdaIndent{2}{}\<[2]%
\>[2]\AgdaFunction{⟦Nat⟧} \AgdaSymbol{=} \AgdaKeyword{record} \<[17]%
\>[17]\<%
\\
\>[2]\AgdaIndent{4}{}\<[4]%
\>[4]\AgdaSymbol{\{} \AgdaField{fm} \AgdaSymbol{=} \AgdaSymbol{λ} \AgdaBound{γ} \AgdaSymbol{→} \AgdaKeyword{record}\<%
\\
\>[4]\AgdaIndent{9}{}\<[9]%
\>[9]\AgdaSymbol{\{} \AgdaField{Carrier} \AgdaSymbol{=} \AgdaDatatype{ℕ}\<%
\\
\>[4]\AgdaIndent{9}{}\<[9]%
\>[9]\AgdaSymbol{;} \AgdaField{\_≈h\_} \AgdaSymbol{=} \AgdaFunction{\_≈nat\_}\<%
\\
\>[4]\AgdaIndent{9}{}\<[9]%
\>[9]\AgdaSymbol{;} \AgdaField{refl} \AgdaSymbol{=} \AgdaSymbol{λ} \AgdaSymbol{\{}\AgdaBound{n}\AgdaSymbol{\}} \AgdaSymbol{→} \AgdaFunction{reflNat} \AgdaSymbol{\{}\AgdaBound{n}\AgdaSymbol{\}}\<%
\\
\>[4]\AgdaIndent{9}{}\<[9]%
\>[9]\AgdaSymbol{;} \AgdaField{sym} \AgdaSymbol{=} \AgdaSymbol{λ} \AgdaSymbol{\{}\AgdaBound{x}\AgdaSymbol{\}} \AgdaSymbol{\{}\AgdaBound{y}\AgdaSymbol{\}} \AgdaSymbol{→} \AgdaFunction{symNat} \AgdaSymbol{\{}\AgdaBound{x}\AgdaSymbol{\}} \AgdaSymbol{\{}\AgdaBound{y}\AgdaSymbol{\}}\<%
\\
\>[4]\AgdaIndent{9}{}\<[9]%
\>[9]\AgdaSymbol{;} \AgdaField{trans} \AgdaSymbol{=} \AgdaSymbol{λ} \AgdaSymbol{\{}\AgdaBound{x}\AgdaSymbol{\}} \AgdaSymbol{\{}\AgdaBound{y}\AgdaSymbol{\}} \AgdaSymbol{\{}\AgdaBound{z}\AgdaSymbol{\}} \AgdaSymbol{→} \AgdaFunction{transNat} \AgdaSymbol{\{}\AgdaBound{x}\AgdaSymbol{\}} \AgdaSymbol{\{}\AgdaBound{y}\AgdaSymbol{\}} \AgdaSymbol{\{}\AgdaBound{z}\AgdaSymbol{\}}\<%
\\
\>[4]\AgdaIndent{9}{}\<[9]%
\>[9]\AgdaSymbol{\}}\<%
\\
\>[0]\AgdaIndent{4}{}\<[4]%
\>[4]\AgdaSymbol{;} \AgdaField{substT} \AgdaSymbol{=} \AgdaSymbol{λ} \AgdaBound{\_} \AgdaSymbol{→} \AgdaFunction{id}\<%
\\
\>[0]\AgdaIndent{4}{}\<[4]%
\>[4]\AgdaSymbol{;} \AgdaField{subst*} \AgdaSymbol{=} \AgdaSymbol{λ} \AgdaBound{\_} \AgdaSymbol{→} \AgdaFunction{id}\<%
\\
\>[0]\AgdaIndent{4}{}\<[4]%
\>[4]\AgdaSymbol{;} \AgdaField{refl*} \AgdaSymbol{=} \AgdaSymbol{λ} \AgdaBound{x} \AgdaBound{a} \AgdaSymbol{→} \AgdaFunction{reflNat} \AgdaSymbol{\{}\AgdaBound{a}\AgdaSymbol{\}}\<%
\\
\>[0]\AgdaIndent{4}{}\<[4]%
\>[4]\AgdaSymbol{;} \AgdaField{trans*} \AgdaSymbol{=} \AgdaSymbol{λ} \AgdaBound{a} \AgdaSymbol{→} \AgdaFunction{reflNat} \AgdaSymbol{\{}\AgdaBound{a}\AgdaSymbol{\}} \<[33]%
\>[33]\<%
\\
\>[0]\AgdaIndent{4}{}\<[4]%
\>[4]\AgdaSymbol{\}}\<%
\\
%
\\
\>[0]\AgdaIndent{2}{}\<[2]%
\>[2]\AgdaFunction{⟦0⟧} \AgdaSymbol{:} \AgdaRecord{Tm} \AgdaFunction{⟦Nat⟧}\<%
\\
\>[0]\AgdaIndent{2}{}\<[2]%
\>[2]\AgdaFunction{⟦0⟧} \AgdaSymbol{=} \AgdaKeyword{record}\<%
\\
\>[2]\AgdaIndent{6}{}\<[6]%
\>[6]\AgdaSymbol{\{} \AgdaField{tm} \AgdaSymbol{=} \AgdaSymbol{λ} \AgdaBound{\_} \AgdaSymbol{→} \AgdaNumber{0}\<%
\\
\>[2]\AgdaIndent{6}{}\<[6]%
\>[6]\AgdaSymbol{;} \AgdaField{respt} \AgdaSymbol{=} \AgdaSymbol{λ} \AgdaBound{p} \AgdaSymbol{→} \AgdaInductiveConstructor{tt}\<%
\\
\>[2]\AgdaIndent{6}{}\<[6]%
\>[6]\AgdaSymbol{\}}\<%
\\
%
\\
\>[0]\AgdaIndent{2}{}\<[2]%
\>[2]\AgdaFunction{⟦s⟧} \AgdaSymbol{:} \AgdaRecord{Tm} \AgdaFunction{⟦Nat⟧} \AgdaSymbol{→} \AgdaRecord{Tm} \AgdaFunction{⟦Nat⟧}\<%
\\
\>[0]\AgdaIndent{2}{}\<[2]%
\>[2]\AgdaFunction{⟦s⟧} \AgdaSymbol{(}\AgdaInductiveConstructor{tm:} \AgdaBound{t} \AgdaInductiveConstructor{resp:} \AgdaBound{respt}\AgdaSymbol{)} \<[26]%
\>[26]\<%
\\
\>[2]\AgdaIndent{6}{}\<[6]%
\>[6]\AgdaSymbol{=} \AgdaKeyword{record}\<%
\\
\>[2]\AgdaIndent{6}{}\<[6]%
\>[6]\AgdaSymbol{\{} \AgdaField{tm} \AgdaSymbol{=} \AgdaInductiveConstructor{suc} \AgdaFunction{∘} \AgdaBound{t}\<%
\\
\>[2]\AgdaIndent{6}{}\<[6]%
\>[6]\AgdaSymbol{;} \AgdaField{respt} \AgdaSymbol{=} \AgdaBound{respt}\<%
\\
\>[2]\AgdaIndent{6}{}\<[6]%
\>[6]\AgdaSymbol{\}}\<%
\\
\>\<\end{code}

Simply typed universe

\AgdaHide{
\begin{code}\>\<%
\\
\>[0]\AgdaIndent{2}{}\<[2]%
\>[2]\AgdaKeyword{data} \AgdaDatatype{⟦U⟧⁰} \AgdaSymbol{:} \AgdaPrimitiveType{Set} \AgdaKeyword{where}\<%
\\
\>[0]\AgdaIndent{4}{}\<[4]%
\>[4]\AgdaInductiveConstructor{nat} \AgdaSymbol{:} \AgdaDatatype{⟦U⟧⁰}\<%
\\
\>[0]\AgdaIndent{4}{}\<[4]%
\>[4]\AgdaInductiveConstructor{arr<\_,\_>} \AgdaSymbol{:} \AgdaSymbol{(}\AgdaBound{a} \AgdaBound{b} \AgdaSymbol{:} \AgdaDatatype{⟦U⟧⁰}\AgdaSymbol{)} \AgdaSymbol{→} \AgdaDatatype{⟦U⟧⁰}\<%
\\
%
\\
\>[0]\AgdaIndent{2}{}\<[2]%
\>[2]\AgdaFunction{\_\textasciitilde⟦U⟧\_} \AgdaSymbol{:} \AgdaDatatype{⟦U⟧⁰} \AgdaSymbol{→} \AgdaDatatype{⟦U⟧⁰} \AgdaSymbol{→} \AgdaRecord{HProp}\<%
\\
\>[0]\AgdaIndent{2}{}\<[2]%
\>[2]\AgdaInductiveConstructor{nat} \AgdaFunction{\textasciitilde⟦U⟧} \AgdaInductiveConstructor{nat} \AgdaSymbol{=} \AgdaFunction{⊤'}\<%
\\
\>[0]\AgdaIndent{2}{}\<[2]%
\>[2]\AgdaInductiveConstructor{nat} \AgdaFunction{\textasciitilde⟦U⟧} \AgdaInductiveConstructor{arr<} \AgdaBound{a} \AgdaInductiveConstructor{,} \AgdaBound{b} \AgdaInductiveConstructor{>} \AgdaSymbol{=} \AgdaFunction{⊥'}\<%
\\
\>[0]\AgdaIndent{2}{}\<[2]%
\>[2]\AgdaInductiveConstructor{arr<} \AgdaBound{a} \AgdaInductiveConstructor{,} \AgdaBound{b} \AgdaInductiveConstructor{>} \AgdaFunction{\textasciitilde⟦U⟧} \AgdaInductiveConstructor{nat} \AgdaSymbol{=} \AgdaFunction{⊥'}\<%
\\
\>[0]\AgdaIndent{2}{}\<[2]%
\>[2]\AgdaInductiveConstructor{arr<} \AgdaBound{a} \AgdaInductiveConstructor{,} \AgdaBound{b} \AgdaInductiveConstructor{>} \AgdaFunction{\textasciitilde⟦U⟧} \AgdaInductiveConstructor{arr<} \AgdaBound{a'} \AgdaInductiveConstructor{,} \AgdaBound{b'} \AgdaInductiveConstructor{>} \AgdaSymbol{=} \AgdaBound{a} \AgdaFunction{\textasciitilde⟦U⟧} \AgdaBound{a'} \AgdaFunction{∧} \AgdaBound{b} \AgdaFunction{\textasciitilde⟦U⟧} \AgdaBound{b'}\<%
\\
%
\\
\>[0]\AgdaIndent{2}{}\<[2]%
\>[2]\AgdaFunction{reflU} \AgdaSymbol{:} \<[11]%
\>[11]\AgdaSymbol{\{}\AgdaBound{x} \AgdaSymbol{:} \AgdaDatatype{⟦U⟧⁰}\AgdaSymbol{\}} \AgdaSymbol{→} \AgdaFunction{<} \AgdaBound{x} \AgdaFunction{\textasciitilde⟦U⟧} \AgdaBound{x} \AgdaFunction{>}\<%
\\
\>[0]\AgdaIndent{2}{}\<[2]%
\>[2]\AgdaFunction{reflU} \AgdaSymbol{\{}\AgdaInductiveConstructor{nat}\AgdaSymbol{\}} \AgdaSymbol{=} \AgdaInductiveConstructor{tt}\<%
\\
\>[0]\AgdaIndent{2}{}\<[2]%
\>[2]\AgdaFunction{reflU} \AgdaSymbol{\{}\AgdaInductiveConstructor{arr<} \AgdaBound{a} \AgdaInductiveConstructor{,} \AgdaBound{b} \AgdaInductiveConstructor{>}\AgdaSymbol{\}} \AgdaSymbol{=} \AgdaFunction{reflU} \AgdaSymbol{\{}\AgdaBound{a}\AgdaSymbol{\}} \AgdaInductiveConstructor{,} \AgdaFunction{reflU} \AgdaSymbol{\{}\AgdaBound{b}\AgdaSymbol{\}}\<%
\\
%
\\
\>[0]\AgdaIndent{2}{}\<[2]%
\>[2]\AgdaFunction{symU} \AgdaSymbol{:} \AgdaSymbol{\{}\AgdaBound{x} \AgdaBound{y} \AgdaSymbol{:} \AgdaDatatype{⟦U⟧⁰}\AgdaSymbol{\}} \AgdaSymbol{→} \AgdaFunction{<} \AgdaBound{x} \AgdaFunction{\textasciitilde⟦U⟧} \AgdaBound{y} \AgdaFunction{>} \AgdaSymbol{→} \AgdaFunction{<} \AgdaBound{y} \AgdaFunction{\textasciitilde⟦U⟧} \AgdaBound{x} \AgdaFunction{>}\<%
\\
\>[0]\AgdaIndent{2}{}\<[2]%
\>[2]\AgdaFunction{symU} \AgdaSymbol{\{}\AgdaInductiveConstructor{nat}\AgdaSymbol{\}} \AgdaSymbol{\{}\AgdaInductiveConstructor{nat}\AgdaSymbol{\}} \AgdaBound{eq} \AgdaSymbol{=} \AgdaInductiveConstructor{tt}\<%
\\
\>[0]\AgdaIndent{2}{}\<[2]%
\>[2]\AgdaFunction{symU} \AgdaSymbol{\{}\AgdaInductiveConstructor{nat}\AgdaSymbol{\}} \AgdaSymbol{\{}\AgdaInductiveConstructor{arr<} \AgdaBound{a} \AgdaInductiveConstructor{,} \AgdaBound{b} \AgdaInductiveConstructor{>}\AgdaSymbol{\}} \AgdaBound{eq} \AgdaSymbol{=} \AgdaBound{eq}\<%
\\
\>[0]\AgdaIndent{2}{}\<[2]%
\>[2]\AgdaFunction{symU} \AgdaSymbol{\{}\AgdaInductiveConstructor{arr<} \AgdaBound{a} \AgdaInductiveConstructor{,} \AgdaBound{b} \AgdaInductiveConstructor{>}\AgdaSymbol{\}} \AgdaSymbol{\{}\AgdaInductiveConstructor{nat}\AgdaSymbol{\}} \AgdaBound{eq} \AgdaSymbol{=} \AgdaBound{eq}\<%
\\
\>[0]\AgdaIndent{2}{}\<[2]%
\>[2]\AgdaFunction{symU} \AgdaSymbol{\{}\AgdaInductiveConstructor{arr<} \AgdaBound{a} \AgdaInductiveConstructor{,} \AgdaBound{b} \AgdaInductiveConstructor{>}\AgdaSymbol{\}} \AgdaSymbol{\{}\AgdaInductiveConstructor{arr<} \AgdaBound{a'} \AgdaInductiveConstructor{,} \AgdaBound{b'} \AgdaInductiveConstructor{>}\AgdaSymbol{\}} \AgdaSymbol{(}\AgdaBound{p} \AgdaInductiveConstructor{,} \AgdaBound{q}\AgdaSymbol{)} \AgdaSymbol{=} \AgdaSymbol{(}\AgdaFunction{symU} \AgdaSymbol{\{}\AgdaBound{a}\AgdaSymbol{\}} \AgdaSymbol{\{}\AgdaBound{a'}\AgdaSymbol{\}} \AgdaBound{p}\AgdaSymbol{)} \<[67]%
\>[67]\<%
\\
\>[2]\AgdaIndent{47}{}\<[47]%
\>[47]\AgdaInductiveConstructor{,} \AgdaSymbol{(}\AgdaFunction{symU} \AgdaSymbol{\{}\AgdaBound{b}\AgdaSymbol{\}} \AgdaSymbol{\{}\AgdaBound{b'}\AgdaSymbol{\}} \AgdaBound{q}\AgdaSymbol{)}\<%
\\
%
\\
\>[0]\AgdaIndent{2}{}\<[2]%
\>[2]\AgdaFunction{transU} \AgdaSymbol{:} \AgdaSymbol{\{}\AgdaBound{x} \AgdaBound{y} \AgdaBound{z} \AgdaSymbol{:} \AgdaDatatype{⟦U⟧⁰}\AgdaSymbol{\}} \AgdaSymbol{→} \AgdaFunction{<} \AgdaBound{x} \AgdaFunction{\textasciitilde⟦U⟧} \AgdaBound{y} \AgdaFunction{>} \AgdaSymbol{→} \AgdaFunction{<} \AgdaBound{y} \AgdaFunction{\textasciitilde⟦U⟧} \AgdaBound{z} \AgdaFunction{>} \AgdaSymbol{→} \AgdaFunction{<} \AgdaBound{x} \AgdaFunction{\textasciitilde⟦U⟧} \AgdaBound{z} \AgdaFunction{>}\<%
\\
\>[0]\AgdaIndent{2}{}\<[2]%
\>[2]\AgdaFunction{transU} \AgdaSymbol{\{}\AgdaInductiveConstructor{nat}\AgdaSymbol{\}} \AgdaSymbol{\{}\AgdaInductiveConstructor{nat}\AgdaSymbol{\}} \AgdaBound{eq1} \AgdaBound{eq2} \AgdaSymbol{=} \AgdaBound{eq2}\<%
\\
\>[0]\AgdaIndent{2}{}\<[2]%
\>[2]\AgdaFunction{transU} \AgdaSymbol{\{}\AgdaInductiveConstructor{nat}\AgdaSymbol{\}} \AgdaSymbol{\{}\AgdaInductiveConstructor{arr<} \AgdaBound{a} \AgdaInductiveConstructor{,} \AgdaBound{b} \AgdaInductiveConstructor{>}\AgdaSymbol{\}} \AgdaSymbol{()} \AgdaBound{eq2}\<%
\\
\>[0]\AgdaIndent{2}{}\<[2]%
\>[2]\AgdaFunction{transU} \AgdaSymbol{\{}\AgdaInductiveConstructor{arr<} \AgdaBound{a} \AgdaInductiveConstructor{,} \AgdaBound{b} \AgdaInductiveConstructor{>}\AgdaSymbol{\}} \AgdaSymbol{\{}\AgdaInductiveConstructor{nat}\AgdaSymbol{\}} \AgdaSymbol{()} \AgdaBound{eq2}\<%
\\
\>[0]\AgdaIndent{2}{}\<[2]%
\>[2]\AgdaFunction{transU} \AgdaSymbol{\{}\AgdaInductiveConstructor{arr<} \AgdaBound{a} \AgdaInductiveConstructor{,} \AgdaBound{b} \AgdaInductiveConstructor{>}\AgdaSymbol{\}} \AgdaSymbol{\{}\AgdaInductiveConstructor{arr<} \AgdaBound{a'} \AgdaInductiveConstructor{,} \AgdaBound{b'} \AgdaInductiveConstructor{>}\AgdaSymbol{\}} \AgdaSymbol{\{}\AgdaInductiveConstructor{nat}\AgdaSymbol{\}} \AgdaBound{eq1} \AgdaBound{eq2} \AgdaSymbol{=} \AgdaBound{eq2}\<%
\\
\>[0]\AgdaIndent{2}{}\<[2]%
\>[2]\AgdaFunction{transU} \AgdaSymbol{\{}\AgdaInductiveConstructor{arr<} \AgdaBound{a} \AgdaInductiveConstructor{,} \AgdaBound{b} \AgdaInductiveConstructor{>}\AgdaSymbol{\}} \AgdaSymbol{\{}\AgdaInductiveConstructor{arr<} \AgdaBound{a'} \AgdaInductiveConstructor{,} \AgdaBound{b'} \AgdaInductiveConstructor{>}\AgdaSymbol{\}} \AgdaSymbol{\{}\AgdaInductiveConstructor{arr<} \AgdaBound{a0} \AgdaInductiveConstructor{,} \AgdaBound{b0} \AgdaInductiveConstructor{>}\AgdaSymbol{\}} \AgdaSymbol{(}\AgdaBound{p1} \AgdaInductiveConstructor{,} \AgdaBound{q1}\AgdaSymbol{)} \<[68]%
\>[68]\<%
\\
\>[2]\AgdaIndent{9}{}\<[9]%
\>[9]\AgdaSymbol{(}\AgdaBound{p2} \AgdaInductiveConstructor{,} \AgdaBound{q2}\AgdaSymbol{)} \AgdaSymbol{=} \AgdaSymbol{(}\AgdaFunction{transU} \AgdaSymbol{\{}\AgdaBound{a}\AgdaSymbol{\}} \AgdaSymbol{\{}\AgdaBound{a'}\AgdaSymbol{\}} \AgdaSymbol{\{}\AgdaBound{a0}\AgdaSymbol{\}} \AgdaBound{p1} \AgdaBound{p2}\AgdaSymbol{)} \<[50]%
\>[50]\<%
\\
\>[2]\AgdaIndent{9}{}\<[9]%
\>[9]\AgdaInductiveConstructor{,} \AgdaFunction{transU} \AgdaSymbol{\{}\AgdaBound{b}\AgdaSymbol{\}} \AgdaSymbol{\{}\AgdaBound{b'}\AgdaSymbol{\}} \AgdaSymbol{\{}\AgdaBound{b0}\AgdaSymbol{\}} \AgdaBound{q1} \AgdaBound{q2}\<%
\\
%
\\
\>[0]\AgdaIndent{2}{}\<[2]%
\>[2]\AgdaFunction{⟦U⟧} \AgdaSymbol{:} \AgdaRecord{Ty} \AgdaBound{Γ}\<%
\\
\>[0]\AgdaIndent{2}{}\<[2]%
\>[2]\AgdaFunction{⟦U⟧} \AgdaSymbol{=} \AgdaKeyword{record} \<[15]%
\>[15]\<%
\\
\>[2]\AgdaIndent{4}{}\<[4]%
\>[4]\AgdaSymbol{\{} \AgdaField{fm} \AgdaSymbol{=} \AgdaSymbol{λ} \AgdaBound{γ} \AgdaSymbol{→} \AgdaKeyword{record}\<%
\\
\>[4]\AgdaIndent{9}{}\<[9]%
\>[9]\AgdaSymbol{\{} \AgdaField{Carrier} \AgdaSymbol{=} \AgdaDatatype{⟦U⟧⁰}\<%
\\
\>[4]\AgdaIndent{9}{}\<[9]%
\>[9]\AgdaSymbol{;} \AgdaField{\_≈h\_} \AgdaSymbol{=} \AgdaFunction{\_\textasciitilde⟦U⟧\_}\<%
\\
\>[4]\AgdaIndent{9}{}\<[9]%
\>[9]\AgdaSymbol{;} \AgdaField{refl} \AgdaSymbol{=} \AgdaSymbol{λ} \AgdaSymbol{\{}\AgdaBound{x}\AgdaSymbol{\}} \AgdaSymbol{→} \AgdaFunction{reflU} \AgdaSymbol{\{}\AgdaBound{x}\AgdaSymbol{\}}\<%
\\
\>[4]\AgdaIndent{9}{}\<[9]%
\>[9]\AgdaSymbol{;} \AgdaField{sym} \AgdaSymbol{=} \AgdaSymbol{λ} \AgdaSymbol{\{}\AgdaBound{x}\AgdaSymbol{\}} \AgdaSymbol{\{}\AgdaBound{y}\AgdaSymbol{\}} \AgdaSymbol{→} \AgdaFunction{symU} \AgdaSymbol{\{}\AgdaBound{x}\AgdaSymbol{\}} \AgdaSymbol{\{}\AgdaBound{y}\AgdaSymbol{\}}\<%
\\
\>[4]\AgdaIndent{9}{}\<[9]%
\>[9]\AgdaSymbol{;} \AgdaField{trans} \AgdaSymbol{=} \AgdaSymbol{λ} \AgdaSymbol{\{}\AgdaBound{x}\AgdaSymbol{\}} \AgdaSymbol{\{}\AgdaBound{y}\AgdaSymbol{\}} \AgdaSymbol{\{}\AgdaBound{z}\AgdaSymbol{\}} \AgdaSymbol{→} \AgdaFunction{transU} \AgdaSymbol{\{}\AgdaBound{x}\AgdaSymbol{\}} \AgdaSymbol{\{}\AgdaBound{y}\AgdaSymbol{\}} \AgdaSymbol{\{}\AgdaBound{z}\AgdaSymbol{\}}\<%
\\
\>[4]\AgdaIndent{9}{}\<[9]%
\>[9]\AgdaSymbol{\}}\<%
\\
\>[0]\AgdaIndent{4}{}\<[4]%
\>[4]\AgdaSymbol{;} \AgdaField{substT} \AgdaSymbol{=} \AgdaSymbol{λ} \AgdaBound{\_} \AgdaSymbol{→} \AgdaFunction{id}\<%
\\
\>[0]\AgdaIndent{4}{}\<[4]%
\>[4]\AgdaSymbol{;} \AgdaField{subst*} \AgdaSymbol{=} \AgdaSymbol{λ} \AgdaBound{\_} \AgdaSymbol{→} \AgdaFunction{id}\<%
\\
\>[0]\AgdaIndent{4}{}\<[4]%
\>[4]\AgdaSymbol{;} \AgdaField{refl*} \AgdaSymbol{=} \AgdaSymbol{λ} \AgdaBound{x} \AgdaBound{a} \AgdaSymbol{→} \AgdaFunction{reflU} \AgdaSymbol{\{}\AgdaBound{a}\AgdaSymbol{\}}\<%
\\
\>[0]\AgdaIndent{4}{}\<[4]%
\>[4]\AgdaSymbol{;} \AgdaField{trans*} \AgdaSymbol{=} \AgdaSymbol{λ} \AgdaBound{a} \AgdaSymbol{→} \AgdaFunction{reflU} \AgdaSymbol{\{}\AgdaBound{a}\AgdaSymbol{\}}\<%
\\
\>[0]\AgdaIndent{4}{}\<[4]%
\>[4]\AgdaSymbol{\}}\<%
\\
%
\\
\>[0]\AgdaIndent{2}{}\<[2]%
\>[2]\AgdaFunction{elfm} \AgdaSymbol{:} \AgdaRecord{Σ} \AgdaFunction{∣} \AgdaBound{Γ} \AgdaFunction{∣} \AgdaSymbol{(λ} \AgdaBound{x} \AgdaSymbol{→} \AgdaDatatype{⟦U⟧⁰}\AgdaSymbol{)} \AgdaSymbol{→} \AgdaRecord{HSetoid}\<%
\\
\>[0]\AgdaIndent{2}{}\<[2]%
\>[2]\AgdaFunction{elfm} \AgdaSymbol{(}\AgdaBound{γ} \AgdaInductiveConstructor{,} \AgdaInductiveConstructor{nat}\AgdaSymbol{)} \AgdaSymbol{=} \AgdaFunction{[} \AgdaFunction{⟦Nat⟧} \AgdaFunction{]fm} \AgdaBound{γ}\<%
\\
\>[0]\AgdaIndent{2}{}\<[2]%
\>[2]\AgdaFunction{elfm} \AgdaSymbol{(}\AgdaBound{γ} \AgdaInductiveConstructor{,} \AgdaInductiveConstructor{arr<} \AgdaBound{a} \AgdaInductiveConstructor{,} \AgdaBound{b} \AgdaInductiveConstructor{>}\AgdaSymbol{)} \AgdaSymbol{=} \AgdaFunction{[} \AgdaBound{Γ} \AgdaFunction{,} \AgdaBound{γ} \AgdaFunction{]} \AgdaFunction{elfm} \AgdaSymbol{(}\AgdaBound{γ} \AgdaInductiveConstructor{,} \AgdaBound{a}\AgdaSymbol{)} \AgdaFunction{⇒fm} \AgdaFunction{elfm} \AgdaSymbol{(}\AgdaBound{γ} \AgdaInductiveConstructor{,} \AgdaBound{b}\AgdaSymbol{)}\<%
\\
\>\<\end{code}
}

\AgdaHide{
\begin{code}\>\<%
\\
%
\\
\>\AgdaComment{\{- To do : To find the way to extract the substT from ->

  elsubstT : \{x y : Σ ∣ Γ ∣ (λ x' → ⟦U⟧⁰)\} →
      Σ < [ Γ ] proj₁ x ≈h proj₁ y > (λ x' → < proj₂ x \textasciitilde⟦U⟧ proj₂ y >) →
      ∣ elfm x ∣ → ∣ elfm y ∣
  elsubstT \{\_ , nat\} \{\_ , nat\} \_ x' = x'
  elsubstT \{\_ , nat\} \{\_ , arr< a , b >\} (p , ()) x'
  elsubstT \{\_ , arr< a , b >\} \{\_ , nat\} (p , ()) x'
  elsubstT \{γ , arr< a , b >\} \{γ' , arr< a' , b' >\} (p , qa , qb) (s1 , s2) = 
   \{!!\}

  ⟦El⟧ : Ty (Γ \& ⟦U⟧)
  ⟦El⟧ = record 
       \{ fm = elfm
       ; substT = elsubstT
       ; subst* = \{!!\}
       ; refl* = \{!!\}
       ; trans* = \{!!\} 
       \}

-\}}\<%
\\
\>\<\end{code}
}

The equality type

\begin{code}\>\<%
\\
\>\AgdaKeyword{module} \AgdaModule{Equality-Type} \AgdaSymbol{(}\AgdaBound{Γ} \AgdaSymbol{:} \AgdaFunction{Con}\AgdaSymbol{)(}\AgdaBound{A} \AgdaSymbol{:} \AgdaRecord{Ty} \AgdaBound{Γ}\AgdaSymbol{)} \AgdaKeyword{where}\<%
\\
%
\\
\>[0]\AgdaIndent{2}{}\<[2]%
\>[2]\AgdaFunction{⟦Id⟧} \AgdaSymbol{:} \AgdaFunction{Rel} \AgdaBound{A}\<%
\\
\>[0]\AgdaIndent{2}{}\<[2]%
\>[2]\AgdaFunction{⟦Id⟧} \AgdaSymbol{=} \AgdaKeyword{record} \<[16]%
\>[16]\<%
\\
\>[2]\AgdaIndent{4}{}\<[4]%
\>[4]\AgdaSymbol{\{} \AgdaField{fm} \AgdaSymbol{=} \AgdaSymbol{λ} \AgdaSymbol{\{((}\AgdaBound{x} \AgdaInductiveConstructor{,} \AgdaBound{a}\AgdaSymbol{)} \AgdaInductiveConstructor{,} \AgdaBound{b}\AgdaSymbol{)} \AgdaSymbol{→} \<[30]%
\>[30]\<%
\\
\>[4]\AgdaIndent{13}{}\<[13]%
\>[13]\AgdaKeyword{record} \<[20]%
\>[20]\<%
\\
\>[4]\AgdaIndent{13}{}\<[13]%
\>[13]\AgdaSymbol{\{} \AgdaField{Carrier} \AgdaSymbol{=} \AgdaFunction{[} \AgdaFunction{[} \AgdaBound{A} \AgdaFunction{]fm} \AgdaBound{x} \AgdaFunction{]} \AgdaBound{a} \AgdaFunction{≈} \AgdaBound{b}\<%
\\
\>[4]\AgdaIndent{13}{}\<[13]%
\>[13]\AgdaSymbol{;} \AgdaField{\_≈h\_} \AgdaSymbol{=} \AgdaSymbol{λ} \AgdaBound{\_} \AgdaBound{\_} \AgdaSymbol{→} \AgdaKeyword{record} \AgdaSymbol{\{} \AgdaField{prf} \AgdaSymbol{=} \AgdaRecord{⊤} \AgdaSymbol{;} \AgdaField{Uni} \AgdaSymbol{=} \AgdaInductiveConstructor{PE.refl} \AgdaSymbol{\}}\<%
\\
\>[4]\AgdaIndent{13}{}\<[13]%
\>[13]\AgdaSymbol{;} \AgdaField{refl} \AgdaSymbol{=} \AgdaInductiveConstructor{tt} \<[25]%
\>[25]\<%
\\
\>[4]\AgdaIndent{13}{}\<[13]%
\>[13]\AgdaSymbol{;} \AgdaField{sym} \AgdaSymbol{=} \AgdaFunction{id}\<%
\\
\>[4]\AgdaIndent{13}{}\<[13]%
\>[13]\AgdaSymbol{;} \AgdaField{trans} \AgdaSymbol{=} \AgdaSymbol{λ} \AgdaBound{\_} \AgdaBound{\_} \AgdaSymbol{→} \AgdaInductiveConstructor{tt}\<%
\\
\>[4]\AgdaIndent{13}{}\<[13]%
\>[13]\AgdaSymbol{\}}\<%
\\
\>[4]\AgdaIndent{13}{}\<[13]%
\>[13]\AgdaSymbol{\}}\<%
\\
\>[0]\AgdaIndent{4}{}\<[4]%
\>[4]\AgdaSymbol{;} \AgdaField{substT} \AgdaSymbol{=} \AgdaSymbol{λ} \AgdaSymbol{\{((}\AgdaBound{x} \AgdaInductiveConstructor{,} \AgdaBound{a}\AgdaSymbol{)} \AgdaInductiveConstructor{,} \AgdaBound{b}\AgdaSymbol{)} \AgdaBound{x0} \AgdaSymbol{→} \<[37]%
\>[37]\<%
\\
\>[0]\AgdaIndent{15}{}\<[15]%
\>[15]\AgdaFunction{[} \AgdaFunction{[} \AgdaBound{A} \AgdaFunction{]fm} \AgdaSymbol{\_} \AgdaFunction{]trans} \<[34]%
\>[34]\<%
\\
\>[0]\AgdaIndent{15}{}\<[15]%
\>[15]\AgdaSymbol{(}\AgdaFunction{[} \AgdaFunction{[} \AgdaBound{A} \AgdaFunction{]fm} \AgdaSymbol{\_} \AgdaFunction{]sym} \AgdaBound{a}\AgdaSymbol{)} \<[36]%
\>[36]\<%
\\
\>[0]\AgdaIndent{15}{}\<[15]%
\>[15]\AgdaSymbol{(}\AgdaFunction{[} \AgdaFunction{[} \AgdaBound{A} \AgdaFunction{]fm} \AgdaSymbol{\_} \AgdaFunction{]trans} \<[35]%
\>[35]\<%
\\
\>[0]\AgdaIndent{15}{}\<[15]%
\>[15]\AgdaSymbol{(}\AgdaFunction{[} \AgdaBound{A} \AgdaFunction{]subst*} \AgdaSymbol{\_} \AgdaBound{x0}\AgdaSymbol{)} \AgdaBound{b}\AgdaSymbol{)} \<[37]%
\>[37]\<%
\\
\>[0]\AgdaIndent{15}{}\<[15]%
\>[15]\AgdaSymbol{\}}\<%
\\
\>[0]\AgdaIndent{4}{}\<[4]%
\>[4]\AgdaSymbol{;} \AgdaField{subst*} \AgdaSymbol{=} \AgdaSymbol{λ} \AgdaBound{\_} \AgdaBound{\_} \AgdaSymbol{→} \AgdaInductiveConstructor{tt}\<%
\\
\>[0]\AgdaIndent{4}{}\<[4]%
\>[4]\AgdaSymbol{;} \AgdaField{refl*} \AgdaSymbol{=} \AgdaSymbol{λ} \AgdaBound{\_} \AgdaBound{\_} \AgdaSymbol{→} \AgdaInductiveConstructor{tt}\<%
\\
\>[0]\AgdaIndent{4}{}\<[4]%
\>[4]\AgdaSymbol{;} \AgdaField{trans*} \AgdaSymbol{=} \AgdaSymbol{λ} \AgdaBound{\_} \AgdaSymbol{→} \AgdaInductiveConstructor{tt} \<[24]%
\>[24]\<%
\\
\>[0]\AgdaIndent{4}{}\<[4]%
\>[4]\AgdaSymbol{\}}\<%
\\
%
\\
\>[0]\AgdaIndent{2}{}\<[2]%
\>[2]\AgdaFunction{⟦refl⟧⁰} \AgdaSymbol{:} \AgdaRecord{Tm} \AgdaSymbol{\{}\AgdaBound{Γ} \AgdaFunction{\&} \AgdaBound{A}\AgdaSymbol{\}} \AgdaSymbol{(}\AgdaFunction{⟦Id⟧} \AgdaFunction{[} \AgdaKeyword{record} \AgdaSymbol{\{} \AgdaField{fn} \AgdaSymbol{=} \AgdaSymbol{λ} \AgdaBound{x'} \AgdaSymbol{→} \AgdaBound{x'} \AgdaInductiveConstructor{,} \AgdaFunction{proj₂} \AgdaBound{x'} \<[66]%
\>[66]\<%
\\
\>[0]\AgdaIndent{23}{}\<[23]%
\>[23]\AgdaSymbol{;} \AgdaField{resp} \AgdaSymbol{=} \AgdaSymbol{λ} \AgdaBound{x'} \AgdaSymbol{→} \AgdaBound{x'} \AgdaInductiveConstructor{,} \AgdaFunction{proj₂} \AgdaBound{x'} \AgdaSymbol{\}} \AgdaFunction{]T}\AgdaSymbol{)} \<[59]%
\>[59]\<%
\\
\>[0]\AgdaIndent{2}{}\<[2]%
\>[2]\AgdaFunction{⟦refl⟧⁰} \AgdaSymbol{=} \AgdaKeyword{record}\<%
\\
\>[0]\AgdaIndent{11}{}\<[11]%
\>[11]\AgdaSymbol{\{} \AgdaField{tm} \AgdaSymbol{=} \AgdaSymbol{λ} \AgdaSymbol{\{(}\AgdaBound{x} \AgdaInductiveConstructor{,} \AgdaBound{a}\AgdaSymbol{)} \AgdaSymbol{→} \AgdaFunction{[} \AgdaFunction{[} \AgdaBound{A} \AgdaFunction{]fm} \AgdaBound{x} \AgdaFunction{]refl} \AgdaSymbol{\{}\AgdaBound{a}\AgdaSymbol{\}} \AgdaSymbol{\}}\<%
\\
\>[0]\AgdaIndent{11}{}\<[11]%
\>[11]\AgdaSymbol{;} \AgdaField{respt} \AgdaSymbol{=} \AgdaSymbol{λ} \AgdaBound{p} \AgdaSymbol{→} \AgdaInductiveConstructor{tt}\<%
\\
\>[0]\AgdaIndent{11}{}\<[11]%
\>[11]\AgdaSymbol{\}}\<%
\\
%
\\
\>[0]\AgdaIndent{2}{}\<[2]%
\>[2]\AgdaFunction{⟦refl⟧} \AgdaSymbol{=} \<[12]%
\>[12]\AgdaFunction{lam} \AgdaSymbol{\{}\AgdaBound{Γ}\AgdaSymbol{\}} \AgdaSymbol{\{}\AgdaBound{A}\AgdaSymbol{\}} \AgdaFunction{⟦refl⟧⁰}\<%
\\
\>\<\end{code}

Subst using equality types

\begin{code}\>\<%
\\
%
\\
\>[0]\AgdaIndent{2}{}\<[2]%
\>[2]\AgdaKeyword{module} \AgdaModule{substIn} \AgdaSymbol{(}\AgdaBound{B} \AgdaSymbol{:} \AgdaRecord{Ty} \AgdaSymbol{(}\AgdaBound{Γ} \AgdaFunction{\&} \AgdaBound{A}\AgdaSymbol{))} \AgdaKeyword{where}\<%
\\
\>[0]\AgdaIndent{2}{}\<[2]%
\>[2]\<%
\\
\>[2]\AgdaIndent{4}{}\<[4]%
\>[4]\AgdaFunction{⟦subst⟧⁰} \AgdaSymbol{:} \AgdaRecord{Tm} \AgdaSymbol{\{}\AgdaBound{Γ} \AgdaFunction{\&} \AgdaBound{A} \AgdaFunction{\&} \AgdaSymbol{(}\AgdaBound{A} \AgdaFunction{[} \AgdaFunction{fst\&} \AgdaSymbol{\{}A \AgdaSymbol{=} \AgdaBound{A}\AgdaSymbol{\}} \AgdaFunction{]T}\AgdaSymbol{)} \<[49]%
\>[49]\<%
\\
\>[4]\AgdaIndent{15}{}\<[15]%
\>[15]\AgdaFunction{\&} \AgdaFunction{⟦Id⟧} \AgdaFunction{\&} \AgdaBound{B} \AgdaFunction{[} \AgdaFunction{fst\&} \AgdaSymbol{\{}A \AgdaSymbol{=} \AgdaBound{A} \AgdaFunction{[} \AgdaFunction{fst\&} \AgdaSymbol{\{}A \AgdaSymbol{=} \AgdaBound{A}\AgdaSymbol{\}} \AgdaFunction{]T}\AgdaSymbol{\}} \<[60]%
\>[60]\AgdaFunction{]T} \<[63]%
\>[63]\<%
\\
\>[4]\AgdaIndent{15}{}\<[15]%
\>[15]\AgdaFunction{[} \AgdaFunction{fst\&} \AgdaSymbol{\{}A \AgdaSymbol{=} \AgdaFunction{⟦Id⟧}\AgdaSymbol{\}} \AgdaFunction{]T}\AgdaSymbol{\}} \<[37]%
\>[37]\<%
\\
\>[-2]\AgdaIndent{13}{}\<[13]%
\>[13]\AgdaSymbol{(}\AgdaBound{B} \AgdaFunction{[} \AgdaKeyword{record} \AgdaSymbol{\{} \AgdaField{fn} \AgdaSymbol{=} \AgdaSymbol{λ} \AgdaBound{x} \AgdaSymbol{→} \AgdaSymbol{(}\AgdaFunction{proj₁} \AgdaSymbol{(}\AgdaFunction{proj₁} \AgdaSymbol{(}\AgdaFunction{proj₁} \AgdaSymbol{(}\AgdaFunction{proj₁} \AgdaBound{x}\AgdaSymbol{))))} \<[72]%
\>[72]\<%
\\
\>[0]\AgdaIndent{13}{}\<[13]%
\>[13]\AgdaInductiveConstructor{,} \AgdaSymbol{(}\AgdaFunction{proj₂} \AgdaSymbol{(}\AgdaFunction{proj₁} \AgdaSymbol{(}\AgdaFunction{proj₁} \AgdaBound{x}\AgdaSymbol{)))} \<[41]%
\>[41]\<%
\\
\>[0]\AgdaIndent{13}{}\<[13]%
\>[13]\AgdaSymbol{;} \AgdaField{resp} \AgdaSymbol{=} \AgdaSymbol{λ} \AgdaBound{x} \AgdaSymbol{→} \AgdaFunction{proj₁} \AgdaSymbol{(}\AgdaFunction{proj₁} \AgdaSymbol{(}\AgdaFunction{proj₁} \AgdaSymbol{(}\AgdaFunction{proj₁} \AgdaBound{x}\AgdaSymbol{)))} \<[60]%
\>[60]\<%
\\
\>[0]\AgdaIndent{13}{}\<[13]%
\>[13]\AgdaInductiveConstructor{,} \AgdaFunction{proj₂} \AgdaSymbol{(}\AgdaFunction{proj₁} \AgdaSymbol{(}\AgdaFunction{proj₁} \AgdaBound{x}\AgdaSymbol{))} \AgdaSymbol{\}} \AgdaFunction{]T}\AgdaSymbol{)}\<%
\\
%
\\
\>[0]\AgdaIndent{4}{}\<[4]%
\>[4]\AgdaFunction{⟦subst⟧⁰} \AgdaSymbol{=} \AgdaKeyword{record}\<%
\\
\>[0]\AgdaIndent{11}{}\<[11]%
\>[11]\AgdaSymbol{\{} \AgdaField{tm} \AgdaSymbol{=} \AgdaSymbol{λ} \AgdaSymbol{\{((((}\AgdaBound{x} \AgdaInductiveConstructor{,} \AgdaBound{a}\AgdaSymbol{)} \AgdaInductiveConstructor{,} \AgdaBound{b}\AgdaSymbol{)} \AgdaInductiveConstructor{,} \AgdaBound{p}\AgdaSymbol{)} \AgdaInductiveConstructor{,} \AgdaBound{PA}\AgdaSymbol{)} \AgdaSymbol{→} \AgdaFunction{[} \AgdaBound{B} \AgdaFunction{]subst} \<[61]%
\>[61]\<%
\\
\>[11]\AgdaIndent{18}{}\<[18]%
\>[18]\AgdaSymbol{(}\AgdaFunction{[} \AgdaBound{Γ} \AgdaFunction{]refl} \AgdaInductiveConstructor{,} \AgdaFunction{[} \AgdaFunction{[} \AgdaBound{A} \AgdaFunction{]fm} \AgdaSymbol{\_} \AgdaFunction{]trans} \<[50]%
\>[50]\<%
\\
\>[11]\AgdaIndent{18}{}\<[18]%
\>[18]\AgdaSymbol{(}\AgdaFunction{[} \AgdaBound{A} \AgdaFunction{]refl*} \AgdaSymbol{\_} \AgdaSymbol{\_)} \AgdaBound{p}\AgdaSymbol{)} \AgdaBound{PA} \AgdaSymbol{\}}\<%
\\
\>[0]\AgdaIndent{11}{}\<[11]%
\>[11]\AgdaSymbol{;} \AgdaField{respt} \AgdaSymbol{=} \AgdaSymbol{λ} \AgdaSymbol{\{((((}\AgdaBound{m} \AgdaInductiveConstructor{,} \AgdaBound{a}\AgdaSymbol{)} \AgdaInductiveConstructor{,} \AgdaBound{b}\AgdaSymbol{)} \AgdaInductiveConstructor{,} \AgdaBound{p}\AgdaSymbol{)} \AgdaInductiveConstructor{,} \AgdaBound{PA}\AgdaSymbol{)} \AgdaSymbol{→} \<[53]%
\>[53]\<%
\\
\>[0]\AgdaIndent{13}{}\<[13]%
\>[13]\AgdaFunction{[} \AgdaFunction{[} \AgdaBound{B} \AgdaFunction{]fm} \AgdaSymbol{\_} \AgdaFunction{]trans} \<[32]%
\>[32]\<%
\\
\>[0]\AgdaIndent{13}{}\<[13]%
\>[13]\AgdaSymbol{(}\AgdaFunction{[} \AgdaBound{B} \AgdaFunction{]trans*} \AgdaSymbol{\_)} \<[29]%
\>[29]\<%
\\
\>[13]\AgdaIndent{14}{}\<[14]%
\>[14]\AgdaSymbol{(}\AgdaFunction{[} \AgdaFunction{[} \AgdaBound{B} \AgdaFunction{]fm} \AgdaSymbol{\_} \AgdaFunction{]trans} \<[34]%
\>[34]\<%
\\
\>[0]\AgdaIndent{13}{}\<[13]%
\>[13]\AgdaFunction{[} \AgdaBound{B} \AgdaFunction{]subst-pi} \<[27]%
\>[27]\<%
\\
\>[0]\AgdaIndent{13}{}\<[13]%
\>[13]\AgdaSymbol{(}\AgdaFunction{[} \AgdaFunction{[} \AgdaBound{B} \AgdaFunction{]fm} \AgdaSymbol{\_} \AgdaFunction{]trans} \<[33]%
\>[33]\<%
\\
\>[0]\AgdaIndent{13}{}\<[13]%
\>[13]\AgdaSymbol{(}\AgdaFunction{[} \AgdaFunction{[} \AgdaBound{B} \AgdaFunction{]fm} \AgdaSymbol{\_} \AgdaFunction{]sym} \AgdaSymbol{(}\AgdaFunction{[} \AgdaBound{B} \AgdaFunction{]trans*} \AgdaSymbol{\_))}\<%
\\
\>[0]\AgdaIndent{13}{}\<[13]%
\>[13]\AgdaSymbol{(}\AgdaFunction{[} \AgdaBound{B} \AgdaFunction{]subst*} \AgdaSymbol{\_} \AgdaBound{PA}\AgdaSymbol{)} \AgdaSymbol{))} \AgdaSymbol{\}}\<%
\\
\>[0]\AgdaIndent{11}{}\<[11]%
\>[11]\AgdaSymbol{\}}\<%
\\
%
\\
\>[0]\AgdaIndent{4}{}\<[4]%
\>[4]\AgdaFunction{⟦subst⟧} \AgdaSymbol{=} \AgdaFunction{lam} \AgdaSymbol{\{}A \AgdaSymbol{=} \AgdaSymbol{(}\AgdaBound{A} \AgdaFunction{[} \AgdaFunction{fst\&} \AgdaSymbol{\{}A \AgdaSymbol{=} \AgdaBound{A}\AgdaSymbol{\}} \AgdaFunction{]T}\AgdaSymbol{)\}} \<[46]%
\>[46]\<%
\\
\>[4]\AgdaIndent{14}{}\<[14]%
\>[14]\AgdaSymbol{(}\AgdaFunction{lam} \AgdaSymbol{\{}A \AgdaSymbol{=} \AgdaFunction{⟦Id⟧}\AgdaSymbol{\}} \<[30]%
\>[30]\<%
\\
\>[4]\AgdaIndent{14}{}\<[14]%
\>[14]\AgdaSymbol{(}\AgdaFunction{lam} \AgdaSymbol{\{}A \AgdaSymbol{=} \AgdaBound{B} \AgdaFunction{[} \AgdaFunction{fst\&} \AgdaSymbol{\{}A \AgdaSymbol{=} \AgdaBound{A} \AgdaFunction{[} \AgdaFunction{fst\&} \AgdaSymbol{\{}A \AgdaSymbol{=} \AgdaBound{A}\AgdaSymbol{\}} \AgdaFunction{]T}\AgdaSymbol{\}}\AgdaFunction{]T} \<[61]%
\>[61]\<%
\\
\>[14]\AgdaIndent{15}{}\<[15]%
\>[15]\AgdaFunction{[} \AgdaFunction{fst\&} \AgdaSymbol{\{}A \AgdaSymbol{=} \AgdaFunction{⟦Id⟧}\AgdaSymbol{\}} \AgdaFunction{]T}\AgdaSymbol{\}} \AgdaFunction{⟦subst⟧⁰}\AgdaSymbol{))}\<%
\\
%
\\
\>\<\end{code}


\AgdaHide{
\begin{code}\>\<%
\\
\>\AgdaKeyword{import} \AgdaModule{Level}\<%
\\
\>\AgdaKeyword{open} \AgdaKeyword{import} \AgdaModule{Relation.Binary.PropositionalEquality} \AgdaSymbol{as} \AgdaModule{PE} \AgdaKeyword{hiding} \AgdaSymbol{(}refl \AgdaSymbol{;} sym \AgdaSymbol{;} trans\AgdaSymbol{;} isEquivalence\AgdaSymbol{;} [\_]\AgdaSymbol{)}\<%
\\
%
\\
\>\AgdaKeyword{module} \AgdaModule{CwF-quotient} \AgdaSymbol{(}\AgdaBound{ext} \AgdaSymbol{:} \AgdaFunction{Extensionality} \AgdaPrimitive{Level.zero} \AgdaPrimitive{Level.zero}\AgdaSymbol{)} \AgdaKeyword{where}\<%
\\
%
\\
\>\AgdaKeyword{open} \AgdaKeyword{import} \AgdaModule{Data.Unit}\<%
\\
\>\AgdaKeyword{open} \AgdaKeyword{import} \AgdaModule{Function}\<%
\\
\>\AgdaKeyword{open} \AgdaKeyword{import} \AgdaModule{Data.Product}\<%
\\
\>\AgdaKeyword{open} \AgdaKeyword{import} \AgdaModule{Data.Bool}\<%
\\
\>\AgdaKeyword{open} \AgdaKeyword{import} \AgdaModule{CategoryofSetoid}\<%
\\
\>\AgdaKeyword{open} \AgdaKeyword{import} \AgdaModule{CwF-setoid}\<%
\\
%
\\
\>\<\end{code}
}

Quotient types

\begin{code}\>\<%
\\
\>\AgdaKeyword{module} \AgdaModule{Q} \AgdaSymbol{(}\AgdaBound{Γ} \AgdaSymbol{:} \AgdaFunction{Con}\AgdaSymbol{)(}\AgdaBound{A} \AgdaSymbol{:} \AgdaRecord{Ty} \AgdaBound{Γ}\AgdaSymbol{)}\<%
\\
%
\\
\>[0]\AgdaIndent{9}{}\<[9]%
\>[9]\AgdaSymbol{(}\AgdaBound{R} \AgdaSymbol{:} \AgdaSymbol{(}\AgdaBound{γ} \AgdaSymbol{:} \AgdaFunction{∣} \AgdaBound{Γ} \AgdaFunction{∣}\AgdaSymbol{)} \AgdaSymbol{→} \AgdaFunction{∣} \AgdaFunction{[} \AgdaBound{A} \AgdaFunction{]fm} \AgdaBound{γ} \AgdaFunction{∣} \AgdaSymbol{→} \AgdaFunction{∣} \AgdaFunction{[} \AgdaBound{A} \AgdaFunction{]fm} \AgdaBound{γ} \AgdaFunction{∣} \AgdaSymbol{→} \AgdaPrimitiveType{Set}\AgdaSymbol{)}\<%
\\
%
\\
\>[0]\AgdaIndent{9}{}\<[9]%
\>[9]\AgdaSymbol{.(}\AgdaBound{Rrespt} \AgdaSymbol{:} \AgdaSymbol{∀\{}\AgdaBound{γ} \AgdaBound{γ'} \AgdaSymbol{:} \AgdaFunction{∣} \AgdaBound{Γ} \AgdaFunction{∣}\AgdaSymbol{\}}\<%
\\
\>[9]\AgdaIndent{19}{}\<[19]%
\>[19]\AgdaSymbol{(}\AgdaBound{p} \AgdaSymbol{:} \AgdaFunction{[} \AgdaBound{Γ} \AgdaFunction{]} \AgdaBound{γ} \AgdaFunction{≈} \AgdaBound{γ'}\AgdaSymbol{)}\<%
\\
\>[9]\AgdaIndent{19}{}\<[19]%
\>[19]\AgdaSymbol{(}\AgdaBound{a} \AgdaBound{b} \AgdaSymbol{:} \AgdaFunction{∣} \AgdaFunction{[} \AgdaBound{A} \AgdaFunction{]fm} \AgdaBound{γ} \AgdaFunction{∣}\AgdaSymbol{)} \AgdaSymbol{→}\<%
\\
\>[9]\AgdaIndent{19}{}\<[19]%
\>[19]\AgdaSymbol{.(}\AgdaBound{R} \AgdaBound{γ} \AgdaBound{a} \AgdaBound{b}\AgdaSymbol{)} \AgdaSymbol{→} \<[32]%
\>[32]\<%
\\
\>[9]\AgdaIndent{19}{}\<[19]%
\>[19]\AgdaBound{R} \AgdaBound{γ'} \AgdaSymbol{(}\AgdaFunction{[} \AgdaBound{A} \AgdaFunction{]subst} \AgdaBound{p} \AgdaBound{a}\AgdaSymbol{)} \AgdaSymbol{(}\AgdaFunction{[} \AgdaBound{A} \AgdaFunction{]subst} \AgdaBound{p} \AgdaBound{b}\AgdaSymbol{))}\<%
\\
%
\\
\>[0]\AgdaIndent{9}{}\<[9]%
\>[9]\AgdaSymbol{.(}\AgdaBound{Rrsp} \AgdaSymbol{:} \AgdaSymbol{∀} \AgdaSymbol{\{}\AgdaBound{γ} \AgdaBound{a} \AgdaBound{b}\AgdaSymbol{\}} \AgdaSymbol{→} \AgdaSymbol{.(}\AgdaFunction{[} \AgdaFunction{[} \AgdaBound{A} \AgdaFunction{]fm} \AgdaBound{γ} \AgdaFunction{]} \AgdaBound{a} \AgdaFunction{≈} \AgdaBound{b}\AgdaSymbol{)} \AgdaSymbol{→} \AgdaBound{R} \AgdaBound{γ} \AgdaBound{a} \AgdaBound{b}\AgdaSymbol{)}\<%
\\
%
\\
\>[0]\AgdaIndent{9}{}\<[9]%
\>[9]\AgdaSymbol{.(}\AgdaBound{Rref} \AgdaSymbol{:} \AgdaSymbol{∀} \AgdaSymbol{\{}\AgdaBound{γ} \AgdaBound{a}\AgdaSymbol{\}} \AgdaSymbol{→} \AgdaBound{R} \AgdaBound{γ} \AgdaBound{a} \AgdaBound{a}\AgdaSymbol{)}\<%
\\
\>[0]\AgdaIndent{9}{}\<[9]%
\>[9]\AgdaSymbol{.(}\AgdaBound{Rsym} \AgdaSymbol{:} \AgdaSymbol{(∀} \AgdaSymbol{\{}\AgdaBound{γ} \AgdaBound{a} \AgdaBound{b}\AgdaSymbol{\}} \AgdaSymbol{→} \AgdaSymbol{.(}\AgdaBound{R} \AgdaBound{γ} \AgdaBound{a} \AgdaBound{b}\AgdaSymbol{)} \AgdaSymbol{→} \AgdaBound{R} \AgdaBound{γ} \AgdaBound{b} \AgdaBound{a}\AgdaSymbol{))}\<%
\\
\>[0]\AgdaIndent{9}{}\<[9]%
\>[9]\AgdaSymbol{.(}\AgdaBound{Rtrn} \AgdaSymbol{:} \<[19]%
\>[19]\AgdaSymbol{(∀} \AgdaSymbol{\{}\AgdaBound{γ} \AgdaBound{a} \AgdaBound{b} \AgdaBound{c}\AgdaSymbol{\}} \AgdaSymbol{→} \AgdaSymbol{.(}\AgdaBound{R} \AgdaBound{γ} \AgdaBound{a} \AgdaBound{b}\AgdaSymbol{)}\<%
\\
\>[9]\AgdaIndent{17}{}\<[17]%
\>[17]\AgdaSymbol{→} \<[20]%
\>[20]\AgdaSymbol{.(}\AgdaBound{R} \AgdaBound{γ} \AgdaBound{b} \AgdaBound{c}\AgdaSymbol{)} \AgdaSymbol{→} \AgdaBound{R} \AgdaBound{γ} \AgdaBound{a} \AgdaBound{c}\AgdaSymbol{))}\<%
\\
\>[0]\AgdaIndent{9}{}\<[9]%
\>[9]\AgdaKeyword{where}\<%
\\
%
\\
\>[0]\AgdaIndent{2}{}\<[2]%
\>[2]\AgdaFunction{⟦Q⟧₀} \AgdaSymbol{:} \AgdaFunction{∣} \AgdaBound{Γ} \AgdaFunction{∣} \AgdaSymbol{→} \AgdaRecord{Setoid}\<%
\\
\>[0]\AgdaIndent{2}{}\<[2]%
\>[2]\AgdaFunction{⟦Q⟧₀} \AgdaBound{γ} \AgdaSymbol{=} \AgdaKeyword{record}\<%
\\
\>[2]\AgdaIndent{9}{}\<[9]%
\>[9]\AgdaSymbol{\{} \AgdaField{Carrier} \AgdaSymbol{=} \AgdaFunction{∣} \AgdaFunction{[} \AgdaBound{A} \AgdaFunction{]fm} \AgdaBound{γ} \AgdaFunction{∣}\<%
\\
\>[2]\AgdaIndent{9}{}\<[9]%
\>[9]\AgdaSymbol{;} \AgdaField{\_≈\_} \AgdaSymbol{=} \AgdaBound{R} \AgdaBound{γ}\<%
\\
\>[2]\AgdaIndent{9}{}\<[9]%
\>[9]\AgdaSymbol{;} \AgdaField{refl} \AgdaSymbol{=} \AgdaBound{Rref}\<%
\\
\>[2]\AgdaIndent{9}{}\<[9]%
\>[9]\AgdaSymbol{;} \AgdaField{sym} \AgdaSymbol{=} \AgdaBound{Rsym}\<%
\\
\>[2]\AgdaIndent{9}{}\<[9]%
\>[9]\AgdaSymbol{;} \AgdaField{trans} \AgdaSymbol{=} \AgdaBound{Rtrn}\<%
\\
\>[2]\AgdaIndent{9}{}\<[9]%
\>[9]\AgdaSymbol{\}}\<%
\\
%
\\
%
\\
\>[0]\AgdaIndent{2}{}\<[2]%
\>[2]\AgdaFunction{⟦Q⟧} \AgdaSymbol{:} \AgdaRecord{Ty} \AgdaBound{Γ}\<%
\\
\>[0]\AgdaIndent{2}{}\<[2]%
\>[2]\AgdaFunction{⟦Q⟧} \AgdaSymbol{=} \AgdaKeyword{record} \<[15]%
\>[15]\<%
\\
\>[2]\AgdaIndent{4}{}\<[4]%
\>[4]\AgdaSymbol{\{} \AgdaField{fm} \AgdaSymbol{=} \AgdaFunction{⟦Q⟧₀}\<%
\\
\>[2]\AgdaIndent{4}{}\<[4]%
\>[4]\AgdaSymbol{;} \AgdaField{substT} \AgdaSymbol{=} \AgdaFunction{[} \AgdaBound{A} \AgdaFunction{]subst}\<%
\\
\>[2]\AgdaIndent{4}{}\<[4]%
\>[4]\AgdaSymbol{;} \AgdaField{subst*} \AgdaSymbol{=} \AgdaSymbol{λ} \AgdaBound{p} \AgdaBound{q} \AgdaSymbol{→} \AgdaBound{Rrespt} \AgdaBound{p} \AgdaSymbol{\_} \AgdaSymbol{\_} \AgdaBound{q}\<%
\\
\>[2]\AgdaIndent{4}{}\<[4]%
\>[4]\AgdaSymbol{;} \AgdaField{refl*} \AgdaSymbol{=} \AgdaBound{Rrsp} \AgdaFunction{[} \AgdaBound{A} \AgdaFunction{]refl*}\<%
\\
\>[2]\AgdaIndent{4}{}\<[4]%
\>[4]\AgdaSymbol{;} \AgdaField{trans*} \AgdaSymbol{=} \AgdaSymbol{λ} \AgdaBound{a} \AgdaSymbol{→} \AgdaBound{Rrsp} \AgdaSymbol{(}\AgdaFunction{[} \AgdaBound{A} \AgdaFunction{]trans*} \AgdaSymbol{\_)}\<%
\\
\>[2]\AgdaIndent{4}{}\<[4]%
\>[4]\AgdaSymbol{\}}\<%
\\
%
\\
\>[0]\AgdaIndent{2}{}\<[2]%
\>[2]\AgdaFunction{⟦[\_]⟧} \AgdaSymbol{:} \AgdaRecord{Tm} \AgdaBound{A} \AgdaSymbol{→} \AgdaRecord{Tm} \AgdaFunction{⟦Q⟧}\<%
\\
\>[0]\AgdaIndent{2}{}\<[2]%
\>[2]\AgdaFunction{⟦[} \AgdaBound{x} \AgdaFunction{]⟧} \AgdaSymbol{=} \AgdaKeyword{record}\<%
\\
\>[2]\AgdaIndent{11}{}\<[11]%
\>[11]\AgdaSymbol{\{} \AgdaField{tm} \AgdaSymbol{=} \AgdaFunction{[} \AgdaBound{x} \AgdaFunction{]tm}\<%
\\
\>[2]\AgdaIndent{11}{}\<[11]%
\>[11]\AgdaSymbol{;} \AgdaField{respt} \AgdaSymbol{=} \AgdaSymbol{λ} \AgdaBound{p} \AgdaSymbol{→} \AgdaBound{Rrsp} \AgdaSymbol{(}\AgdaFunction{[} \AgdaBound{x} \AgdaFunction{]respt} \AgdaBound{p}\AgdaSymbol{)}\<%
\\
\>[2]\AgdaIndent{11}{}\<[11]%
\>[11]\AgdaSymbol{\}}\<%
\\
\>[0]\AgdaIndent{2}{}\<[2]%
\>[2]\<%
\\
\>[0]\AgdaIndent{2}{}\<[2]%
\>[2]\AgdaFunction{⟦[\_]⟧'} \AgdaSymbol{:} \AgdaRecord{Tm} \AgdaSymbol{(}\AgdaBound{A} \AgdaFunction{⇒} \AgdaFunction{⟦Q⟧}\AgdaSymbol{)}\<%
\\
\>[0]\AgdaIndent{2}{}\<[2]%
\>[2]\AgdaFunction{⟦[\_]⟧'} \AgdaSymbol{=} \AgdaKeyword{record}\<%
\\
\>[2]\AgdaIndent{11}{}\<[11]%
\>[11]\AgdaSymbol{\{} \AgdaField{tm} \AgdaSymbol{=} \AgdaSymbol{λ} \AgdaBound{x} \AgdaSymbol{→} \AgdaSymbol{(λ} \AgdaBound{a} \AgdaSymbol{→} \AgdaBound{a}\AgdaSymbol{)} \AgdaInductiveConstructor{,} \<[36]%
\>[36]\<%
\\
\>[11]\AgdaIndent{18}{}\<[18]%
\>[18]\AgdaSymbol{(λ} \AgdaBound{a} \AgdaBound{b} \AgdaBound{p} \AgdaSymbol{→} \<[29]%
\>[29]\<%
\\
\>[11]\AgdaIndent{18}{}\<[18]%
\>[18]\AgdaBound{Rrsp} \AgdaSymbol{(}\AgdaFunction{[} \AgdaFunction{[} \AgdaBound{A} \AgdaFunction{]fm} \AgdaSymbol{\_} \AgdaFunction{]trans} \AgdaFunction{[} \AgdaBound{A} \AgdaFunction{]refl*} \AgdaBound{p}\AgdaSymbol{))}\<%
\\
\>[0]\AgdaIndent{11}{}\<[11]%
\>[11]\AgdaSymbol{;} \AgdaField{respt} \AgdaSymbol{=} \AgdaSymbol{λ} \AgdaBound{p} \AgdaBound{a} \AgdaSymbol{→} \AgdaBound{Rrsp} \AgdaFunction{[} \AgdaBound{A} \AgdaFunction{]tr*}\<%
\\
\>[0]\AgdaIndent{11}{}\<[11]%
\>[11]\AgdaSymbol{\}}\<%
\\
\>\<\end{code}

\begin{code}\>\>[0]\AgdaIndent{2}{}\<[2]%
\>[2]\<%
\\
\>[0]\AgdaIndent{2}{}\<[2]%
\>[2]\AgdaSymbol{.}\AgdaFunction{Q-Ax} \AgdaSymbol{:} \AgdaSymbol{∀} \AgdaBound{γ} \AgdaBound{a} \AgdaBound{b} \AgdaSymbol{→} \AgdaFunction{[} \AgdaFunction{[} \AgdaBound{A} \AgdaFunction{]fm} \AgdaBound{γ} \AgdaFunction{]} \AgdaBound{a} \AgdaFunction{≈} \AgdaBound{b} \AgdaSymbol{→} \AgdaFunction{[} \AgdaFunction{[} \AgdaFunction{⟦Q⟧} \AgdaFunction{]fm} \AgdaSymbol{\_} \AgdaFunction{]} \AgdaBound{a} \AgdaFunction{≈} \AgdaBound{b}\<%
\\
\>[0]\AgdaIndent{2}{}\<[2]%
\>[2]\AgdaFunction{Q-Ax} \AgdaBound{γ} \AgdaBound{a} \AgdaBound{b} \AgdaSymbol{=} \AgdaBound{Rrsp}\<%
\\
%
\\
\>[0]\AgdaIndent{2}{}\<[2]%
\>[2]\AgdaFunction{Q-elim} \AgdaSymbol{:} \AgdaSymbol{(}\AgdaBound{B} \AgdaSymbol{:} \AgdaRecord{Ty} \AgdaBound{Γ}\AgdaSymbol{)} \<[22]%
\>[22]\<%
\\
\>[2]\AgdaIndent{9}{}\<[9]%
\>[9]\AgdaSymbol{→} \AgdaSymbol{(}\AgdaBound{f} \AgdaSymbol{:} \AgdaRecord{Tm} \AgdaSymbol{(}\AgdaBound{A} \AgdaFunction{⇒} \AgdaBound{B}\AgdaSymbol{))}\<%
\\
\>[2]\AgdaIndent{9}{}\<[9]%
\>[9]\AgdaSymbol{→} \AgdaSymbol{(}\AgdaBound{frespR} \AgdaSymbol{:} \AgdaSymbol{∀} \AgdaBound{γ} \AgdaBound{a} \AgdaBound{b}\<%
\\
\>[9]\AgdaIndent{18}{}\<[18]%
\>[18]\AgdaSymbol{→} \AgdaSymbol{(}\AgdaBound{R} \AgdaBound{γ} \AgdaBound{a} \AgdaBound{b}\AgdaSymbol{)}\<%
\\
\>[9]\AgdaIndent{18}{}\<[18]%
\>[18]\AgdaSymbol{→} \AgdaFunction{[} \AgdaFunction{[} \AgdaBound{B} \AgdaFunction{]fm} \AgdaBound{γ} \AgdaFunction{]} \AgdaFunction{prj₁} \AgdaSymbol{(}\AgdaFunction{[} \AgdaBound{f} \AgdaFunction{]tm} \AgdaBound{γ}\AgdaSymbol{)} \AgdaBound{a} \<[53]%
\>[53]\<%
\\
\>[18]\AgdaIndent{20}{}\<[20]%
\>[20]\AgdaFunction{≈} \<[23]%
\>[23]\AgdaFunction{prj₁} \AgdaSymbol{(}\AgdaFunction{[} \AgdaBound{f} \AgdaFunction{]tm} \AgdaBound{γ}\AgdaSymbol{)} \AgdaBound{b}\AgdaSymbol{)}\<%
\\
\>[-7]\AgdaIndent{9}{}\<[9]%
\>[9]\AgdaSymbol{→} \AgdaRecord{Tm} \AgdaSymbol{(}\AgdaFunction{⟦Q⟧} \AgdaFunction{⇒} \AgdaBound{B}\AgdaSymbol{)}\<%
\\
\>[0]\AgdaIndent{2}{}\<[2]%
\>[2]\AgdaFunction{Q-elim} \AgdaBound{B} \AgdaBound{f} \AgdaBound{frespR} \AgdaSymbol{=} \AgdaKeyword{record}\<%
\\
\>[2]\AgdaIndent{11}{}\<[11]%
\>[11]\AgdaSymbol{\{} \AgdaField{tm} \AgdaSymbol{=} \AgdaSymbol{λ} \AgdaBound{γ} \AgdaSymbol{→} \AgdaFunction{prj₁} \AgdaSymbol{(}\AgdaFunction{[} \AgdaBound{f} \AgdaFunction{]tm} \AgdaBound{γ}\AgdaSymbol{)} \AgdaInductiveConstructor{,} \AgdaSymbol{(λ} \AgdaBound{a} \AgdaBound{b} \AgdaBound{p} \AgdaSymbol{→} \<[54]%
\>[54]\<%
\\
\>[11]\AgdaIndent{18}{}\<[18]%
\>[18]\AgdaFunction{[} \AgdaFunction{[} \AgdaBound{B} \AgdaFunction{]fm} \AgdaSymbol{\_} \AgdaFunction{]trans} \AgdaFunction{[} \AgdaBound{B} \AgdaFunction{]refl*} \AgdaSymbol{(}\AgdaBound{frespR} \AgdaSymbol{\_} \AgdaSymbol{\_} \AgdaSymbol{\_} \AgdaBound{p}\AgdaSymbol{))}\<%
\\
\>[0]\AgdaIndent{11}{}\<[11]%
\>[11]\AgdaSymbol{;} \AgdaField{respt} \AgdaSymbol{=} \AgdaSymbol{λ} \AgdaSymbol{\{}\AgdaBound{γ}\AgdaSymbol{\}} \AgdaSymbol{\{}\AgdaBound{γ'}\AgdaSymbol{\}} \AgdaBound{p} \AgdaBound{a} \AgdaSymbol{→} \AgdaFunction{[} \AgdaBound{f} \AgdaFunction{]respt} \AgdaBound{p} \AgdaBound{a}\<%
\\
\>[0]\AgdaIndent{11}{}\<[11]%
\>[11]\AgdaSymbol{\}}\<%
\\
%
\\
\>[0]\AgdaIndent{2}{}\<[2]%
\>[2]\<%
\\
\>[0]\AgdaIndent{2}{}\<[2]%
\>[2]\AgdaFunction{substQ} \AgdaSymbol{:} \AgdaSymbol{(}\AgdaBound{Γ} \AgdaFunction{\&} \AgdaBound{A}\AgdaSymbol{)} \AgdaRecord{⇉} \AgdaSymbol{(}\AgdaBound{Γ} \AgdaFunction{\&} \AgdaFunction{⟦Q⟧}\AgdaSymbol{)}\<%
\\
\>[0]\AgdaIndent{2}{}\<[2]%
\>[2]\AgdaFunction{substQ} \AgdaSymbol{=} \AgdaKeyword{record}\<%
\\
\>[2]\AgdaIndent{11}{}\<[11]%
\>[11]\AgdaSymbol{\{} \AgdaField{fn} \AgdaSymbol{=} \AgdaSymbol{λ} \AgdaSymbol{\{(}\AgdaBound{x} \AgdaInductiveConstructor{,} \AgdaBound{a}\AgdaSymbol{)} \AgdaSymbol{→} \AgdaBound{x} \AgdaInductiveConstructor{,} \AgdaBound{a}\AgdaSymbol{\}}\<%
\\
\>[2]\AgdaIndent{11}{}\<[11]%
\>[11]\AgdaSymbol{;} \AgdaField{resp} \AgdaSymbol{=} \AgdaSymbol{λ\{} \AgdaSymbol{(}\AgdaBound{p} \AgdaInductiveConstructor{,} \AgdaBound{q}\AgdaSymbol{)} \AgdaSymbol{→} \AgdaBound{p} \AgdaInductiveConstructor{,} \AgdaSymbol{(}\AgdaBound{Rrsp} \AgdaBound{q}\AgdaSymbol{)\}}\<%
\\
\>[2]\AgdaIndent{11}{}\<[11]%
\>[11]\AgdaSymbol{\}}\<%
\\
%
\\
\>[0]\AgdaIndent{2}{}\<[2]%
\>[2]\AgdaFunction{Q-ind} \AgdaSymbol{:} \AgdaSymbol{(}\AgdaBound{P} \AgdaSymbol{:} \AgdaRecord{Ty} \AgdaSymbol{(}\AgdaBound{Γ} \AgdaFunction{\&} \AgdaFunction{⟦Q⟧}\AgdaSymbol{))}\<%
\\
\>[0]\AgdaIndent{9}{}\<[9]%
\>[9]\AgdaSymbol{→} \AgdaSymbol{(}\AgdaBound{isProp} \AgdaSymbol{:} \AgdaSymbol{∀} \AgdaSymbol{\{}\AgdaBound{x} \AgdaBound{a}\AgdaSymbol{\}} \AgdaSymbol{(}\AgdaBound{r} \AgdaBound{s} \AgdaSymbol{:} \AgdaFunction{∣} \AgdaFunction{[} \AgdaBound{P} \AgdaFunction{]fm} \AgdaSymbol{(}\AgdaBound{x} \AgdaInductiveConstructor{,} \AgdaBound{a}\AgdaSymbol{)} \AgdaFunction{∣}\AgdaSymbol{)} \AgdaSymbol{→} \<[59]%
\>[59]\<%
\\
\>[9]\AgdaIndent{22}{}\<[22]%
\>[22]\AgdaFunction{[} \AgdaFunction{[} \AgdaBound{P} \AgdaFunction{]fm} \AgdaSymbol{(}\AgdaBound{x} \AgdaInductiveConstructor{,} \AgdaBound{a}\AgdaSymbol{)} \AgdaFunction{]} \AgdaBound{r} \AgdaFunction{≈} \AgdaBound{s} \AgdaSymbol{)}\<%
\\
\>[0]\AgdaIndent{9}{}\<[9]%
\>[9]\AgdaSymbol{→} \AgdaSymbol{(}\AgdaBound{h} \AgdaSymbol{:} \AgdaRecord{Tm} \AgdaSymbol{(}\AgdaFunction{Π} \AgdaBound{A} \AgdaSymbol{(}\AgdaBound{P} \AgdaFunction{[} \AgdaFunction{substQ} \AgdaFunction{]T}\AgdaSymbol{)))}\<%
\\
\>[0]\AgdaIndent{9}{}\<[9]%
\>[9]\AgdaSymbol{→} \AgdaRecord{Tm} \AgdaSymbol{(}\AgdaFunction{Π} \AgdaFunction{⟦Q⟧} \AgdaBound{P}\AgdaSymbol{)}\<%
\\
\>[0]\AgdaIndent{2}{}\<[2]%
\>[2]\AgdaFunction{Q-ind} \AgdaBound{P} \AgdaBound{isProp} \AgdaBound{h} \AgdaSymbol{=} \AgdaKeyword{record}\<%
\\
\>[0]\AgdaIndent{11}{}\<[11]%
\>[11]\AgdaSymbol{\{} \AgdaField{tm} \AgdaSymbol{=} \AgdaSymbol{λ} \AgdaBound{x} \AgdaSymbol{→} \AgdaSymbol{(}\AgdaFunction{prj₁} \AgdaSymbol{(}\AgdaFunction{[} \AgdaBound{h} \AgdaFunction{]tm} \AgdaBound{x}\AgdaSymbol{))} \AgdaInductiveConstructor{,} \<[45]%
\>[45]\<%
\\
\>[11]\AgdaIndent{18}{}\<[18]%
\>[18]\AgdaSymbol{(λ} \AgdaBound{a} \AgdaBound{b} \AgdaBound{p} \AgdaSymbol{→} \AgdaBound{isProp} \AgdaSymbol{\{}\AgdaBound{x}\AgdaSymbol{\}} \AgdaSymbol{\{}\AgdaBound{b}\AgdaSymbol{\}} \AgdaSymbol{\_} \AgdaSymbol{\_)}\<%
\\
\>[0]\AgdaIndent{11}{}\<[11]%
\>[11]\AgdaSymbol{;} \AgdaField{respt} \AgdaSymbol{=} \AgdaFunction{[} \AgdaBound{h} \AgdaFunction{]respt}\<%
\\
\>[0]\AgdaIndent{11}{}\<[11]%
\>[11]\AgdaSymbol{\}}\<%
\\
\>\<\end{code}







\addtocontents{toc}{\vspace{2em}} % Add a gap in the Contents, for aesthetics

\backmatter

%----------------------------------------------------------------------------------------
%	BIBLIOGRAPHY
%----------------------------------------------------------------------------------------

\label{Bibliography}

\lhead{\emph{Bibliography}} % Change the page header to say "Bibliography"

\bibliographystyle{unsrtnat} % Use the "unsrtnat" BibTeX style for formatting the Bibliography

\bibliography{my}


\end{document}
